

\newpage
\subsubsection{ESP.VCMULAS.S16.QACC.H}\label{instruction:ESPVCMULASS16QACC_H}
\textbf{Syntax}\\
ESP.VCMULAS.S16.QACC.H qx,qy,sat

\textbf{Description}\\
This instruction performs multiplication of 16-bit vectors and stores the result in \lstinline{QACC_H}.

\textbf{Operation}\\
\begin{lstlisting}
    add_0[64:0] = QACC_H[ 63:  0] + (qx[ 79: 64]*qy[ 79: 64] - qx[95:80]*qy[95:80])
    add_1[64:0] = QACC_H[127: 64] + (qx[ 95: 80]*qy[ 79: 64] + qx[79:64]*qy[95:80])
    add_2[64:0] = QACC_H[191:128] + (qx[111:96]*qy[111:96] - qx[127:112]*qy[127:112])
    add_3[64:0] = QACC_H[255:192] + (qx[127:112]*qy[111:96] - qx[111:96]*qy[127:112])
    QACC_H[ 63:  0] = min(max(add_0[64:0], -2^{63}), 2^{63}-1)
    QACC_H[127: 64] = min(max(add_1[64:0], -2^{63}), 2^{63}-1)
    QACC_H[191:128] = min(max(add_2[64:0], -2^{63}), 2^{63}-1)
    QACC_H[255:192] = min(max(add_3[64:0], -2^{63}), 2^{63}-1)
\end{lstlisting}

\textbf{Schedule}\\
\begin{longtable}[c]{|l|l|l|}
    \hline
    \rowcolor{lightgray}
    \textbf{Operands} & \textbf{Use} & \textbf{Def} \\\hline
    \endfirsthead
    \hline
    \rowcolor{lightgray}
     \textbf{Operands} & \textbf{Use} & \textbf{Def} \\\hline
    \endhead
    \hline
    \endfoot

     qx & 1 & - \\\hline
     qy & 1 & - \\\hline
     QACC\_H & 2 & 2

\end{longtable}

\textbf{Format}\\

\begin{figure*}[!htbp]
 {\setlength{\tabcolsep}{1pt}
 \begin{center}
\begin{tabular}{@{}F@{}F@{}T@{}I@{}T@{}O}
\instbitrange{31}{29}&
\instbitrange{28}{26}&
\instbitrange{25}{17}&
\instbit{16}&
\instbitrange{15}{7}&
\instbitrange{6}{0} \\
\hline
\multicolumn{1}{|c|}{{\scriptsize qx[2:0]}} &
\multicolumn{1}{c|}{{\scriptsize qy[2:0]}} &
\multicolumn{1}{c|}{1X1X0XX11} &
\multicolumn{1}{c|}{{\scriptsize sat[0]}} &
\multicolumn{1}{c|}{101XX0X0X} &
\multicolumn{1}{c|}{0011011}\\
\hline
3& 3& 9& 1& 9& 7 \\
\end{tabular}
 \end{center}
 }
 \vspace{-0.1in}
 \caption[]{ESP.VCMULAS.S16.QACC.H} \end{figure*}


\newpage
\subsubsection{ESP.VCMULAS.S16.QACC.H.LD.IP}\label{instruction:ESPVCMULASS16QACC_HLDIP}
\textbf{Syntax}\\
ESP.VCMULAS.S16.QACC.H.LD.IP qu,rs1,imm16,qx,qy,sat

\textbf{Description}\\
This instruction performs multiplication of 16-bit vectors and stores the result in \lstinline{QACC_H}.

During the operation, it loads the 16-byte data from the memory into the register \lstinline{qu}. After the access is completed, the value in the register \lstinline{rs1} is incremented by a 4-bit sign-extended constant in the instruction code segment, left-shifted by 4.

imm16 starts at -128, ends at 112, with a step of 16.

\textbf{Operation}\\
\begin{lstlisting}
    add_0[64:0] = QACC_H[63:0] + (qx[79:64]*qy[79:64] - qx[95:80]*qy[95:80])
    add_1[64:0] = QACC_H[127:64] + (qx[95:80]*qy[79:64] + qx[79:64]*qy[95:80])
    add_2[64:0] = QACC_H[191:128] + (qx[111:96]*qy[111:96] - qx[127:112]*qy[127:112])
    add_3[64:0] = QACC_H[255:192] + (qx[127:112]*qy[111:96] - qx[111:96]*qy[127:112])

    QACC_H[ 63:  0] = min(max(add_0[64:0], -2^{63}), 2^{63}-1)
    QACC_H[127: 64] = min(max(add_1[64:0], -2^{63}), 2^{63}-1)
    QACC_H[191:128] = min(max(add_2[64:0], -2^{63}), 2^{63}-1)
    QACC_H[255:192] = min(max(add_3[64:0], -2^{63}), 2^{63}-1)
    qu[127:0] = load128(rs1[31:0])
    rs1[31:0] = rs1[31:0] + imm16
\end{lstlisting}

\textbf{Schedule}\\
\begin{longtable}[c]{|l|l|l|}
    \hline
    \rowcolor{lightgray}
    \textbf{Operands} & \textbf{Use} & \textbf{Def} \\\hline
    \endfirsthead
    \hline
    \rowcolor{lightgray}
     \textbf{Operands} & \textbf{Use} & \textbf{Def} \\\hline
    \endhead
    \hline
    \endfoot

    qu & - & 2 \\\hline
    rs1 & 1 & 1 \\\hline
     qx & 1 & - \\\hline
     qy & 1 & - \\\hline
     QACC\_H & 2 & 2

\end{longtable}

\textbf{Format}\\

\begin{figure*}[!htbp]
 {\setlength{\tabcolsep}{1pt}
 \begin{center}
\begin{tabular}{@{}F@{}F@{}W@{}I@{}Y@{}I@{}F@{}W@{}F@{}F@{}O}
\instbitrange{31}{29}&
\instbitrange{28}{26}&
\instbitrange{25}{24}&
\instbit{23}&
\instbitrange{22}{19}&
\instbit{18}&
\instbitrange{17}{15}&
\instbitrange{14}{13}&
\instbitrange{12}{10}&
\instbitrange{9}{7}&
\instbitrange{6}{0} \\
\hline
\multicolumn{1}{|c|}{{\scriptsize qx[2:0]}} &
\multicolumn{1}{c|}{{\scriptsize qy[2:0]}} &
\multicolumn{1}{c|}{11} &
\multicolumn{1}{c|}{{\scriptsize sat[0]}} &
\multicolumn{1}{c|}{{\scriptsize imm16[3:0]}} &
\multicolumn{1}{c|}{{\scriptsize rs1[4]}} &
\multicolumn{1}{c|}{{\scriptsize rs1[2:0]}} &
\multicolumn{1}{c|}{01} &
\multicolumn{1}{c|}{{\scriptsize qu[2:0]}} &
\multicolumn{1}{c|}{1X1} &
\multicolumn{1}{c|}{0011111}\\
\hline
3& 3& 2& 1& 4& 1& 3& 2& 3& 3& 7 \\
\end{tabular}
 \end{center}
 }
 \vspace{-0.1in}
 \caption[]{ESP.VCMULAS.S16.QACC.H.LD.IP} \end{figure*}


\newpage
\subsubsection{ESP.VCMULAS.S16.QACC.H.LD.XP}\label{instruction:ESPVCMULASS16QACC_HLDXP}
\textbf{Syntax}\\
ESP.VCMULAS.S16.QACC.H.LD.XP qu,rs1,rs2,qx,qy,sat

\textbf{Description}\\
This instruction performs multiplication of 16-bit vectors and stores the result in \lstinline{QACC_H}.

During the operation, it loads the 16-byte data from the memory into the register \lstinline{qu}. After the access is completed, the value in the register \lstinline{rs1} is incremented by the value of the register \lstinline{rs2}.

\textbf{Operation}\\
\begin{lstlisting}
    add_0[64:0] = QACC_H[63:0] + (qx[79:64]*qy[79:64] - qx[95:80]*qy[95:80])
    add_1[64:0] = QACC_H[127:64] + (qx[95:80]*qy[79:64] + qx[79:64]*qy[95:80])
    add_2[64:0] = QACC_H[191:128] + (qx[111:96]*qy[111:96] - qx[127:112]*qy[127:112])
    add_3[64:0] = QACC_H[255:192] + (qx[127:112]*qy[111:96] - qx[111:96]*qy[127:112])

    QACC_H[ 63:  0] = min(max(add_0[64:0], -2^{63}), 2^{63}-1)
    QACC_H[127: 64] = min(max(add_1[64:0], -2^{63}), 2^{63}-1)
    QACC_H[191:128] = min(max(add_2[64:0], -2^{63}), 2^{63}-1)
    QACC_H[255:192] = min(max(add_3[64:0], -2^{63}), 2^{63}-1)

    qu[127:0] = load128(rs1[31:0])
    rs1[31:0] = rs1[31:0] + rs2[31:0]
\end{lstlisting}

\textbf{Schedule}\\
\begin{longtable}[c]{|l|l|l|}
    \hline
    \rowcolor{lightgray}
    \textbf{Operands} & \textbf{Use} & \textbf{Def} \\\hline
    \endfirsthead
    \hline
    \rowcolor{lightgray}
     \textbf{Operands} & \textbf{Use} & \textbf{Def} \\\hline
    \endhead
    \hline
    \endfoot

    qu & - & 2 \\\hline
    rs1 & 1 & 1 \\\hline
    rs2 & 1 & - \\\hline
     qx & 1 & - \\\hline
     qy & 1 & - \\\hline
     QACC\_H & 2 & 2

\end{longtable}

\textbf{Format}\\

\begin{figure*}[!htbp]
 {\setlength{\tabcolsep}{1pt}
 \begin{center}
\begin{tabular}{@{}F@{}F@{}W@{}I@{}F@{}I@{}I@{}F@{}W@{}F@{}F@{}O}
\instbitrange{31}{29}&
\instbitrange{28}{26}&
\instbitrange{25}{24}&
\instbit{23}&
\instbitrange{22}{20}&
\instbit{19}&
\instbit{18}&
\instbitrange{17}{15}&
\instbitrange{14}{13}&
\instbitrange{12}{10}&
\instbitrange{9}{7}&
\instbitrange{6}{0} \\
\hline
\multicolumn{1}{|c|}{{\scriptsize qx[2:0]}} &
\multicolumn{1}{c|}{{\scriptsize qy[2:0]}} &
\multicolumn{1}{c|}{11} &
\multicolumn{1}{c|}{{\scriptsize rs2[4]}} &
\multicolumn{1}{c|}{{\scriptsize rs2[2:0]}} &
\multicolumn{1}{c|}{{\scriptsize sat[0]}} &
\multicolumn{1}{c|}{{\scriptsize rs1[4]}} &
\multicolumn{1}{c|}{{\scriptsize rs1[2:0]}} &
\multicolumn{1}{c|}{10} &
\multicolumn{1}{c|}{{\scriptsize qu[2:0]}} &
\multicolumn{1}{c|}{111} &
\multicolumn{1}{c|}{0011111}\\
\hline
3& 3& 2& 1& 3& 1& 1& 3& 2& 3& 3& 7 \\
\end{tabular}
 \end{center}
 }
 \vspace{-0.1in}
 \caption[]{ESP.VCMULAS.S16.QACC.H.LD.XP} \end{figure*}


\newpage
\subsubsection{ESP.VCMULAS.S16.QACC.L}\label{instruction:ESPVCMULASS16QACC_L}
\textbf{Syntax}\\
ESP.VCMULAS.S16.QACC.L qx,qy,sat

\textbf{Description}\\
This instruction performs multiplication of 16-bit vectors and stores the result in \lstinline{QACC_L}.

\textbf{Operation}\\
\begin{lstlisting}
    add_0[64:0] = QACC_L[63:0] + (qx[15:0]*qy[15:0] - qx[31:16]*qy[31:16])
    add_1[64:0] = QACC_L[127:64] + (qx[31:16]*qy[15:0] + qx[15:0]*qy[31:16])
    add_2[64:0] = QACC_L[191:128] + (qx[47:32]*qy[47:32] - qx[63:48]*qy[63:48])
    add_3[64:0] = QACC_L[255:192] + (qx[63:48]*qy[47:32] - qx[47:32]*qy[63:48])

    QACC_L[ 63:  0] = min(max(add_0[64:0], -2^{63}), 2^{63}-1)
    QACC_L[127: 64] = min(max(add_1[64:0], -2^{63}), 2^{63}-1)
    QACC_L[191:128] = min(max(add_2[64:0], -2^{63}), 2^{63}-1)
    QACC_L[255:192] = min(max(add_3[64:0], -2^{63}), 2^{63}-1)
\end{lstlisting}

\textbf{Schedule}\\
\begin{longtable}[c]{|l|l|l|}
    \hline
    \rowcolor{lightgray}
    \textbf{Operands} & \textbf{Use} & \textbf{Def} \\\hline
    \endfirsthead
    \hline
    \rowcolor{lightgray}
     \textbf{Operands} & \textbf{Use} & \textbf{Def} \\\hline
    \endhead
    \hline
    \endfoot


     qx & 1 & - \\\hline
     qy & 1 & - \\\hline
     QACC\_L & 2 & 2

\end{longtable}

\textbf{Format}\\

\begin{figure*}[!htbp]
 {\setlength{\tabcolsep}{1pt}
 \begin{center}
\begin{tabular}{@{}F@{}F@{}T@{}I@{}T@{}O}
\instbitrange{31}{29}&
\instbitrange{28}{26}&
\instbitrange{25}{17}&
\instbit{16}&
\instbitrange{15}{7}&
\instbitrange{6}{0} \\
\hline
\multicolumn{1}{|c|}{{\scriptsize qx[2:0]}} &
\multicolumn{1}{c|}{{\scriptsize qy[2:0]}} &
\multicolumn{1}{c|}{1X1X0XX01} &
\multicolumn{1}{c|}{{\scriptsize sat[0]}} &
\multicolumn{1}{c|}{101XX0X0X} &
\multicolumn{1}{c|}{0011011}\\
\hline
3& 3& 9& 1& 9& 7 \\
\end{tabular}
 \end{center}
 }
 \vspace{-0.1in}
 \caption[]{ESP.VCMULAS.S16.QACC.L} \end{figure*}


\newpage
\subsubsection{ESP.VCMULAS.S16.QACC.L.LD.IP}\label{instruction:ESPVCMULASS16QACC_LLDIP}
\textbf{Syntax}\\
ESP.VCMULAS.S16.QACC.L.LD.IP qu,rs1,imm16,qx,qy,sat

\textbf{Description}\\
This instruction performs multiplication of 16-bit vectors and stores the result in \lstinline{QACC_L}.

During the operation, it loads the 16-byte data from the memory into the register \lstinline{qu}. After the access is completed, the value in the register \lstinline{rs1} is incremented by a 4-bit sign-extended constant in the instruction code segment, left-shifted by 4.

imm16 starts at -128, ends at 112, with a step of 16.

\textbf{Operation}\\
\begin{lstlisting}
    add_0[64:0] = QACC_L[63:0] + (qx[15:0]*qy[15:0] - qx[31:16]*qy[31:16])
    add_1[64:0] = QACC_L[127:64] + (qx[31:16]*qy[15:0] + qx[15:0]*qy[31:16])
    add_2[64:0] = QACC_L[191:128] + (qx[47:32]*qy[47:32] - qx[63:48]*qy[63:48])
    add_3[64:0] = QACC_L[255:192] + (qx[63:48]*qy[47:32] - qx[47:32]*qy[63:48])

    QACC_L[ 63:  0] = min(max(add_0[64:0], -2^{63}), 2^{63}-1)
    QACC_L[127: 64] = min(max(add_1[64:0], -2^{63}), 2^{63}-1)
    QACC_L[191:128] = min(max(add_2[64:0], -2^{63}), 2^{63}-1)
    QACC_L[255:192] = min(max(add_3[64:0], -2^{63}), 2^{63}-1)

    qu[127:0] = load128(rs1[31:0])
    rs1[31:0] = rs1[31:0] + imm16
\end{lstlisting}

\textbf{Schedule}\\
\begin{longtable}[c]{|l|l|l|}
    \hline
    \rowcolor{lightgray}
    \textbf{Operands} & \textbf{Use} & \textbf{Def} \\\hline
    \endfirsthead
    \hline
    \rowcolor{lightgray}
     \textbf{Operands} & \textbf{Use} & \textbf{Def} \\\hline
    \endhead
    \hline
    \endfoot

    qu & - & 2 \\\hline
    rs1 & 1 & 1 \\\hline
     qx & 1 & - \\\hline
     qy & 1 & - \\\hline
     QACC\_L & 2 & 2

\end{longtable}

\textbf{Format}\\

\begin{figure*}[!htbp]
 {\setlength{\tabcolsep}{1pt}
 \begin{center}
\begin{tabular}{@{}F@{}F@{}W@{}I@{}Y@{}I@{}F@{}W@{}F@{}F@{}O}
\instbitrange{31}{29}&
\instbitrange{28}{26}&
\instbitrange{25}{24}&
\instbit{23}&
\instbitrange{22}{19}&
\instbit{18}&
\instbitrange{17}{15}&
\instbitrange{14}{13}&
\instbitrange{12}{10}&
\instbitrange{9}{7}&
\instbitrange{6}{0} \\
\hline
\multicolumn{1}{|c|}{{\scriptsize qx[2:0]}} &
\multicolumn{1}{c|}{{\scriptsize qy[2:0]}} &
\multicolumn{1}{c|}{01} &
\multicolumn{1}{c|}{{\scriptsize sat[0]}} &
\multicolumn{1}{c|}{{\scriptsize imm16[3:0]}} &
\multicolumn{1}{c|}{{\scriptsize rs1[4]}} &
\multicolumn{1}{c|}{{\scriptsize rs1[2:0]}} &
\multicolumn{1}{c|}{01} &
\multicolumn{1}{c|}{{\scriptsize qu[2:0]}} &
\multicolumn{1}{c|}{1X1} &
\multicolumn{1}{c|}{0011111}\\
\hline
3& 3& 2& 1& 4& 1& 3& 2& 3& 3& 7 \\
\end{tabular}
 \end{center}
 }
 \vspace{-0.1in}
 \caption[]{ESP.VCMULAS.S16.QACC.L.LD.IP} \end{figure*}


\newpage
\subsubsection{ESP.VCMULAS.S16.QACC.L.LD.XP}\label{instruction:ESPVCMULASS16QACC_LLDXP}
\textbf{Syntax}\\
ESP.VCMULAS.S16.QACC.L.LD.XP qu,rs1,rs2,qx,qy,sat

\textbf{Description}\\
This instruction performs multiplication of 16-bit vectors and stores the result in \lstinline{QACC_L}.

During the operation, it loads the 16-byte data from the memory into the register \lstinline{qu}. After the access is completed, the value in the register \lstinline{rs1} is incremented by the value of the register \lstinline{rs2}.

\textbf{Operation}\\
\begin{lstlisting}
    add_0[64:0] = QACC_L[63:0] + (qx[15:0]*qy[15:0] - qx[31:16]*qy[31:16])
    add_1[64:0] = QACC_L[127:64] + (qx[31:16]*qy[15:0] + qx[15:0]*qy[31:16])
    add_2[64:0] = QACC_L[191:128] + (qx[47:32]*qy[47:32] - qx[63:48]*qy[63:48])
    add_3[64:0] = QACC_L[255:192] + (qx[63:48]*qy[47:32] - qx[47:32]*qy[63:48])

    QACC_L[ 63:  0] = min(max(add_0[64:0], -2^{63}), 2^{63}-1)
    QACC_L[127: 64] = min(max(add_1[64:0], -2^{63}), 2^{63}-1)
    QACC_L[191:128] = min(max(add_2[64:0], -2^{63}), 2^{63}-1)
    QACC_L[255:192] = min(max(add_3[64:0], -2^{63}), 2^{63}-1)

    qu[127:0] = load128(rs1[31:0])
    rs1[31:0] = rs1[31:0] + rs2[31:0]
\end{lstlisting}

\textbf{Schedule}\\
\begin{longtable}[c]{|l|l|l|}
    \hline
    \rowcolor{lightgray}
    \textbf{Operands} & \textbf{Use} & \textbf{Def} \\\hline
    \endfirsthead
    \hline
    \rowcolor{lightgray}
     \textbf{Operands} & \textbf{Use} & \textbf{Def} \\\hline
    \endhead
    \hline
    \endfoot

    qu & - & 2 \\\hline
    rs1 & 1 & 1 \\\hline
    rs1 & 1 & - \\\hline
     qx & 1 & - \\\hline
     qy & 1 & - \\\hline
     QACC\_L & 2 & 2

\end{longtable}

\textbf{Format}\\

\begin{figure*}[!htbp]
 {\setlength{\tabcolsep}{1pt}
 \begin{center}
\begin{tabular}{@{}F@{}F@{}W@{}I@{}F@{}I@{}I@{}F@{}W@{}F@{}F@{}O}
\instbitrange{31}{29}&
\instbitrange{28}{26}&
\instbitrange{25}{24}&
\instbit{23}&
\instbitrange{22}{20}&
\instbit{19}&
\instbit{18}&
\instbitrange{17}{15}&
\instbitrange{14}{13}&
\instbitrange{12}{10}&
\instbitrange{9}{7}&
\instbitrange{6}{0} \\
\hline
\multicolumn{1}{|c|}{{\scriptsize qx[2:0]}} &
\multicolumn{1}{c|}{{\scriptsize qy[2:0]}} &
\multicolumn{1}{c|}{01} &
\multicolumn{1}{c|}{{\scriptsize rs2[4]}} &
\multicolumn{1}{c|}{{\scriptsize rs2[2:0]}} &
\multicolumn{1}{c|}{{\scriptsize sat[0]}} &
\multicolumn{1}{c|}{{\scriptsize rs1[4]}} &
\multicolumn{1}{c|}{{\scriptsize rs1[2:0]}} &
\multicolumn{1}{c|}{10} &
\multicolumn{1}{c|}{{\scriptsize qu[2:0]}} &
\multicolumn{1}{c|}{111} &
\multicolumn{1}{c|}{0011111}\\
\hline
3& 3& 2& 1& 3& 1& 1& 3& 2& 3& 3& 7 \\
\end{tabular}
 \end{center}
 }
 \vspace{-0.1in}
 \caption[]{ESP.VCMULAS.S16.QACC.L.LD.XP} \end{figure*}


\newpage
\subsubsection{ESP.VCMULAS.S8.QACC.H}\label{instruction:ESPVCMULASS8QACC_H}
\textbf{Syntax}\\
ESP.VCMULAS.S8.QACC.H qx,qy,sat

\textbf{Description}\\
This instruction performs multiplication of 8-bit vectors and stores the result in \lstinline{QACC_H}.

\textbf{Operation}\\
\begin{lstlisting}
    add_0[32:0] = QACC_H[ 31:  0] + (qx[ 71: 64]*qy[ 71: 64] - qx[ 79: 72]*qy[ 79: 72])
    add_1[32:0] = QACC_H[ 63: 32] + (qx[ 79: 72]*qy[ 71: 64] - qx[ 71: 64]*qy[ 79: 72])

    add_2[32:0] = QACC_H[ 95: 64] + (qx[ 87: 80]*qy[ 87: 80] - qx[ 95: 88]*qy[ 95: 88])
    add_3[32:0] = QACC_H[127: 96] + (qx[ 95: 88]*qy[ 87: 80] - qx[ 87: 80]*qy[ 95: 88])

    add_4[32:0] = QACC_H[159:128] + (qx[103: 96]*qy[103: 96] - qx[111:104]*qy[111:104])
    add_5[32:0] = QACC_H[191:160] + (qx[111:104]*qy[103: 96] - qx[103: 96]*qy[111:104])

    add_6[32:0] = QACC_H[223:192] + (qx[119:112]*qy[119:112] - qx[127:120]*qy[127:120])
    add_7[32:0] = QACC_H[255:224] + (qx[127:120]*qy[119:112] - qx[119:112]*qy[127:120])

    QACC_H[ 31:  0] = min(max(add_0[32:0], -2^{31}), 2^{31}-1)
    QACC_H[ 63: 32] = min(max(add_1[32:0], -2^{31}), 2^{31}-1)
    QACC_H[ 95: 64] = min(max(add_2[32:0], -2^{31}), 2^{31}-1)
    QACC_H[127: 96] = min(max(add_3[32:0], -2^{31}), 2^{31}-1)
    QACC_H[159:128] = min(max(add_4[32:0], -2^{31}), 2^{31}-1)
    QACC_H[191:160] = min(max(add_5[32:0], -2^{31}), 2^{31}-1)
    QACC_H[223:192] = min(max(add_6[32:0], -2^{31}), 2^{31}-1)
    QACC_H[255:224] = min(max(add_7[32:0], -2^{31}), 2^{31}-1)

\end{lstlisting}

\textbf{Schedule}\\
\begin{longtable}[c]{|l|l|l|}
    \hline
    \rowcolor{lightgray}
    \textbf{Operands} & \textbf{Use} & \textbf{Def} \\\hline
    \endfirsthead
    \hline
    \rowcolor{lightgray}
     \textbf{Operands} & \textbf{Use} & \textbf{Def} \\\hline
    \endhead
    \hline
    \endfoot

     qx & 1 & - \\\hline
     qy & 1 & - \\\hline
     QACC\_H & 2 & 2

\end{longtable}

\textbf{Format}\\

\begin{figure*}[!htbp]
 {\setlength{\tabcolsep}{1pt}
 \begin{center}
\begin{tabular}{@{}F@{}F@{}T@{}I@{}T@{}O}
\instbitrange{31}{29}&
\instbitrange{28}{26}&
\instbitrange{25}{17}&
\instbit{16}&
\instbitrange{15}{7}&
\instbitrange{6}{0} \\
\hline
\multicolumn{1}{|c|}{{\scriptsize qx[2:0]}} &
\multicolumn{1}{c|}{{\scriptsize qy[2:0]}} &
\multicolumn{1}{c|}{1X1X0XX10} &
\multicolumn{1}{c|}{{\scriptsize sat[0]}} &
\multicolumn{1}{c|}{101XX0X0X} &
\multicolumn{1}{c|}{0011011}\\
\hline
3& 3& 9& 1& 9& 7 \\
\end{tabular}
 \end{center}
 }
 \vspace{-0.1in}
 \caption[]{ESP.VCMULAS.S8.QACC.H} \end{figure*}


\newpage
\subsubsection{ESP.VCMULAS.S8.QACC.H.LD.IP}\label{instruction:ESPVCMULASS8QACC_HLDIP}
\textbf{Syntax}\\
ESP.VCMULAS.S8.QACC.H.LD.IP qu,rs1,imm16,qx,qy,sat

\textbf{Description}\\
This instruction performs multiplication of 8-bit vectors and stores the result in \lstinline{QACC_H}.

During the operation, it loads the 16-byte data from the memory into the register \lstinline{qu}. After the access is completed, the value in the register \lstinline{rs1} is incremented by a 4-bit sign-extended constant in the instruction code segment, left-shifted by 4.

imm16 starts at -128, ends at 112, with a step of 16.

\textbf{Operation}\\
\begin{lstlisting}
    add_0[32:0] = QACC_H[ 31:  0] + (qx[ 71: 64]*qy[ 71: 64] - qx[ 79: 72]*qy[ 79: 72])
    add_1[32:0] = QACC_H[ 63: 32] + (qx[ 79: 72]*qy[ 71: 64] - qx[ 71: 64]*qy[ 79: 72])

    add_2[32:0] = QACC_H[ 95: 64] + (qx[ 87: 80]*qy[ 87: 80] - qx[ 95: 88]*qy[ 95: 88])
    add_3[32:0] = QACC_H[127: 96] + (qx[ 95: 88]*qy[ 87: 80] - qx[ 87: 80]*qy[ 95: 88])

    add_4[32:0] = QACC_H[159:128] + (qx[103: 96]*qy[103: 96] - qx[111:104]*qy[111:104])
    add_5[32:0] = QACC_H[191:160] + (qx[111:104]*qy[103: 96] - qx[103: 96]*qy[111:104])

    add_6[32:0] = QACC_H[223:192] + (qx[119:112]*qy[119:112] - qx[127:120]*qy[127:120])
    add_7[32:0] = QACC_H[255:224] + (qx[127:120]*qy[119:112] - qx[119:112]*qy[127:120])

    QACC_H[ 31:  0] = min(max(add_0[32:0], -2^{31}), 2^{31}-1)
    QACC_H[ 63: 32] = min(max(add_1[32:0], -2^{31}), 2^{31}-1)
    QACC_H[ 95: 64] = min(max(add_2[32:0], -2^{31}), 2^{31}-1)
    QACC_H[127: 96] = min(max(add_3[32:0], -2^{31}), 2^{31}-1)
    QACC_H[159:128] = min(max(add_4[32:0], -2^{31}), 2^{31}-1)
    QACC_H[191:160] = min(max(add_5[32:0], -2^{31}), 2^{31}-1)
    QACC_H[223:192] = min(max(add_6[32:0], -2^{31}), 2^{31}-1)
    QACC_H[255:224] = min(max(add_7[32:0], -2^{31}), 2^{31}-1)

    qu[127:0] = load128(rs1[31:0])
    rs1[31:0] = rs1[31:0] + signed32(imm16)
\end{lstlisting}

\textbf{Schedule}\\
\begin{longtable}[c]{|l|l|l|}
    \hline
    \rowcolor{lightgray}
    \textbf{Operands} & \textbf{Use} & \textbf{Def} \\\hline
    \endfirsthead
    \hline
    \rowcolor{lightgray}
     \textbf{Operands} & \textbf{Use} & \textbf{Def} \\\hline
    \endhead
    \hline
    \endfoot

    qu & - & 2 \\\hline
    rs1 & 1 & 1 \\\hline
     qx & 1 & - \\\hline
     qy & 1 & - \\\hline
     QACC\_H & 2 & 2

\end{longtable}

\textbf{Format}\\

\begin{figure*}[!htbp]
 {\setlength{\tabcolsep}{1pt}
 \begin{center}
\begin{tabular}{@{}F@{}F@{}W@{}I@{}Y@{}I@{}F@{}W@{}F@{}F@{}O}
\instbitrange{31}{29}&
\instbitrange{28}{26}&
\instbitrange{25}{24}&
\instbit{23}&
\instbitrange{22}{19}&
\instbit{18}&
\instbitrange{17}{15}&
\instbitrange{14}{13}&
\instbitrange{12}{10}&
\instbitrange{9}{7}&
\instbitrange{6}{0} \\
\hline
\multicolumn{1}{|c|}{{\scriptsize qx[2:0]}} &
\multicolumn{1}{c|}{{\scriptsize qy[2:0]}} &
\multicolumn{1}{c|}{10} &
\multicolumn{1}{c|}{{\scriptsize sat[0]}} &
\multicolumn{1}{c|}{{\scriptsize imm16[3:0]}} &
\multicolumn{1}{c|}{{\scriptsize rs1[4]}} &
\multicolumn{1}{c|}{{\scriptsize rs1[2:0]}} &
\multicolumn{1}{c|}{01} &
\multicolumn{1}{c|}{{\scriptsize qu[2:0]}} &
\multicolumn{1}{c|}{1X1} &
\multicolumn{1}{c|}{0011111}\\
\hline
3& 3& 2& 1& 4& 1& 3& 2& 3& 3& 7 \\
\end{tabular}
 \end{center}
 }
 \vspace{-0.1in}
 \caption[]{ESP.VCMULAS.S8.QACC.H.LD.IP} \end{figure*}


\newpage
\subsubsection{ESP.VCMULAS.S8.QACC.H.LD.XP}\label{instruction:ESPVCMULASS8QACC_HLDXP}
\textbf{Syntax}\\
ESP.VCMULAS.S8.QACC.H.LD.XP qu,rs1,rs2,qx,qy,sat

\textbf{Description}\\
This instruction performs multiplication of 8-bit vectors and stores the result in \lstinline{QACC_H}.

During the operation, it loads the 16-byte data from the memory into the register \lstinline{qu}. After the access is completed, the value in the register \lstinline{rs1} is incremented by the value of the register \lstinline{rs2}.

\textbf{Operation}\\
\begin{lstlisting}
    add_0[32:0] = QACC_H[ 31:  0] + (qx[ 71: 64]*qy[ 71: 64] - qx[ 79: 72]*qy[ 79: 72])
    add_1[32:0] = QACC_H[ 63: 32] + (qx[ 79: 72]*qy[ 71: 64] - qx[ 71: 64]*qy[ 79: 72])

    add_2[32:0] = QACC_H[ 95: 64] + (qx[ 87: 80]*qy[ 87: 80] - qx[ 95: 88]*qy[ 95: 88])
    add_3[32:0] = QACC_H[127: 96] + (qx[ 95: 88]*qy[ 87: 80] - qx[ 87: 80]*qy[ 95: 88])

    add_4[32:0] = QACC_H[159:128] + (qx[103: 96]*qy[103: 96] - qx[111:104]*qy[111:104])
    add_5[32:0] = QACC_H[191:160] + (qx[111:104]*qy[103: 96] - qx[103: 96]*qy[111:104])

    add_6[32:0] = QACC_H[223:192] + (qx[119:112]*qy[119:112] - qx[127:120]*qy[127:120])
    add_7[32:0] = QACC_H[255:224] + (qx[127:120]*qy[119:112] - qx[119:112]*qy[127:120])

    QACC_H[ 31:  0] = min(max(add_0[32:0], -2^{31}), 2^{31}-1)
    QACC_H[ 63: 32] = min(max(add_1[32:0], -2^{31}), 2^{31}-1)
    QACC_H[ 95: 64] = min(max(add_2[32:0], -2^{31}), 2^{31}-1)
    QACC_H[127: 96] = min(max(add_3[32:0], -2^{31}), 2^{31}-1)
    QACC_H[159:128] = min(max(add_4[32:0], -2^{31}), 2^{31}-1)
    QACC_H[191:160] = min(max(add_5[32:0], -2^{31}), 2^{31}-1)
    QACC_H[223:192] = min(max(add_6[32:0], -2^{31}), 2^{31}-1)
    QACC_H[255:224] = min(max(add_7[32:0], -2^{31}), 2^{31}-1)

    qu[127:0] = load128(rs1[31:0])
    rs1[31:0] = rs1[31:0] + rs2[31:0]
\end{lstlisting}

\textbf{Schedule}\\
\begin{longtable}[c]{|l|l|l|}
    \hline
    \rowcolor{lightgray}
    \textbf{Operands} & \textbf{Use} & \textbf{Def} \\\hline
    \endfirsthead
    \hline
    \rowcolor{lightgray}
     \textbf{Operands} & \textbf{Use} & \textbf{Def} \\\hline
    \endhead
    \hline
    \endfoot

    qu & - & 2 \\\hline
    rs1 & 1 & 1 \\\hline
    rs2 & 1 & - \\\hline
     qx & 1 & - \\\hline
     qy & 1 & - \\\hline
     QACC\_H & 2 & 2

\end{longtable}

\textbf{Format}\\

\begin{figure*}[!htbp]
 {\setlength{\tabcolsep}{1pt}
 \begin{center}
\begin{tabular}{@{}F@{}F@{}W@{}I@{}F@{}I@{}I@{}F@{}W@{}F@{}F@{}O}
\instbitrange{31}{29}&
\instbitrange{28}{26}&
\instbitrange{25}{24}&
\instbit{23}&
\instbitrange{22}{20}&
\instbit{19}&
\instbit{18}&
\instbitrange{17}{15}&
\instbitrange{14}{13}&
\instbitrange{12}{10}&
\instbitrange{9}{7}&
\instbitrange{6}{0} \\
\hline
\multicolumn{1}{|c|}{{\scriptsize qx[2:0]}} &
\multicolumn{1}{c|}{{\scriptsize qy[2:0]}} &
\multicolumn{1}{c|}{10} &
\multicolumn{1}{c|}{{\scriptsize rs2[4]}} &
\multicolumn{1}{c|}{{\scriptsize rs2[2:0]}} &
\multicolumn{1}{c|}{{\scriptsize sat[0]}} &
\multicolumn{1}{c|}{{\scriptsize rs1[4]}} &
\multicolumn{1}{c|}{{\scriptsize rs1[2:0]}} &
\multicolumn{1}{c|}{10} &
\multicolumn{1}{c|}{{\scriptsize qu[2:0]}} &
\multicolumn{1}{c|}{111} &
\multicolumn{1}{c|}{0011111}\\
\hline
3& 3& 2& 1& 3& 1& 1& 3& 2& 3& 3& 7 \\
\end{tabular}
 \end{center}
 }
 \vspace{-0.1in}
 \caption[]{ESP.VCMULAS.S8.QACC.H.LD.XP} \end{figure*}


\newpage
\subsubsection{ESP.VCMULAS.S8.QACC.L}\label{instruction:ESPVCMULASS8QACC_L}
\textbf{Syntax}\\
ESP.VCMULAS.S8.QACC.L qx,qy,sat

\textbf{Description}\\
This instruction performs multiplication of 8-bit vectors and stores the result in \lstinline{QACC_L}.

\textbf{Operation}\\
\begin{lstlisting}
    add_0[32:0] = QACC_L[ 31:  0] + (qx[  7:  0]*qy[  7:  0] - qx[ 15:  8]*qy[ 15:  8])
    add_1[32:0] = QACC_L[ 63: 32] + (qx[ 15:  8]*qy[  7:  0] - qx[  7:  0]*qy[ 15:  8])

    add_2[32:0] = QACC_L[ 95: 64] + (qx[ 23: 16]*qy[ 23: 16] - qx[ 31: 24]*qy[ 31: 24])
    add_3[32:0] = QACC_L[127: 96] + (qx[ 31: 24]*qy[ 23: 16] - qx[ 23: 16]*qy[ 31: 24])

    add_4[32:0] = QACC_L[159:128] + (qx[ 39: 32]*qy[ 39: 32] - qx[ 47: 40]*qy[ 47: 40])
    add_5[32:0] = QACC_L[191:160] + (qx[ 47: 40]*qy[ 39: 32] - qx[ 39: 32]*qy[ 47: 40])

    add_6[32:0] = QACC_L[223:192] + (qx[ 55: 48]*qy[ 55: 48] - qx[ 63: 56]*qy[ 63: 56])
    add_7[32:0] = QACC_L[255:224] + (qx[ 63: 56]*qy[ 55: 48] - qx[ 55: 48]*qy[ 63: 56])

    QACC_L[ 31:  0] = min(max(add_0[32:0], -2^{31}), 2^{31}-1)
    QACC_L[ 63: 32] = min(max(add_1[32:0], -2^{31}), 2^{31}-1)
    QACC_L[ 95: 64] = min(max(add_2[32:0], -2^{31}), 2^{31}-1)
    QACC_L[127: 96] = min(max(add_3[32:0], -2^{31}), 2^{31}-1)
    QACC_L[159:128] = min(max(add_4[32:0], -2^{31}), 2^{31}-1)
    QACC_L[191:160] = min(max(add_5[32:0], -2^{31}), 2^{31}-1)
    QACC_L[223:192] = min(max(add_6[32:0], -2^{31}), 2^{31}-1)
    QACC_L[255:224] = min(max(add_7[32:0], -2^{31}), 2^{31}-1)

\end{lstlisting}

\textbf{Schedule}\\
\begin{longtable}[c]{|l|l|l|}
    \hline
    \rowcolor{lightgray}
    \textbf{Operands} & \textbf{Use} & \textbf{Def} \\\hline
    \endfirsthead
    \hline
    \rowcolor{lightgray}
     \textbf{Operands} & \textbf{Use} & \textbf{Def} \\\hline
    \endhead
    \hline
    \endfoot

     qx & 1 & - \\\hline
     qy & 1 & - \\\hline
     QACC\_L & 2 & 2

\end{longtable}

\textbf{Format}\\

\begin{figure*}[!htbp]
 {\setlength{\tabcolsep}{1pt}
 \begin{center}
\begin{tabular}{@{}F@{}F@{}T@{}I@{}T@{}O}
\instbitrange{31}{29}&
\instbitrange{28}{26}&
\instbitrange{25}{17}&
\instbit{16}&
\instbitrange{15}{7}&
\instbitrange{6}{0} \\
\hline
\multicolumn{1}{|c|}{{\scriptsize qx[2:0]}} &
\multicolumn{1}{c|}{{\scriptsize qy[2:0]}} &
\multicolumn{1}{c|}{1X1X0XX00} &
\multicolumn{1}{c|}{{\scriptsize sat[0]}} &
\multicolumn{1}{c|}{101XX0X0X} &
\multicolumn{1}{c|}{0011011}\\
\hline
3& 3& 9& 1& 9& 7 \\
\end{tabular}
 \end{center}
 }
 \vspace{-0.1in}
 \caption[]{ESP.VCMULAS.S8.QACC.L} \end{figure*}


\newpage
\subsubsection{ESP.VCMULAS.S8.QACC.L.LD.IP}\label{instruction:ESPVCMULASS8QACC_LLDIP}
\textbf{Syntax}\\
ESP.VCMULAS.S8.QACC.L.LD.IP qu,rs1,imm16,qx,qy,sat

\textbf{Description}\\
This instruction performs multiplication of 8-bit vectors and stores the result in \lstinline{QACC_L}.

During the operation, it loads the 16-byte data from the memory into the register \lstinline{qu}. After the access is completed, the value in the register \lstinline{rs1} is incremented by a 4-bit sign-extended constant in the instruction code segment, left-shifted by 4.

imm16 starts at -128, ends at 112, with a step of 16.

\textbf{Operation}\\
\begin{lstlisting}
    add_0[32:0] = QACC_L[ 31:  0] + (qx[  7:  0]*qy[  7:  0] - qx[ 15:  8]*qy[ 15:  8])
    add_1[32:0] = QACC_L[ 63: 32] + (qx[ 15:  8]*qy[  7:  0] - qx[  7:  0]*qy[ 15:  8])

    add_2[32:0] = QACC_L[ 95: 64] + (qx[ 23: 16]*qy[ 23: 16] - qx[ 31: 24]*qy[ 31: 24])
    add_3[32:0] = QACC_L[127: 96] + (qx[ 31: 24]*qy[ 23: 16] - qx[ 23: 16]*qy[ 31: 24])

    add_4[32:0] = QACC_L[159:128] + (qx[ 39: 32]*qy[ 39: 32] - qx[ 47: 40]*qy[ 47: 40])
    add_5[32:0] = QACC_L[191:160] + (qx[ 47: 40]*qy[ 39: 32] - qx[ 39: 32]*qy[ 47: 40])

    add_6[32:0] = QACC_L[223:192] + (qx[ 55: 48]*qy[ 55: 48] - qx[ 63: 56]*qy[ 63: 56])
    add_7[32:0] = QACC_L[255:224] + (qx[ 63: 56]*qy[ 55: 48] - qx[ 55: 48]*qy[ 63: 56])

    QACC_L[ 31:  0] = min(max(add_0[32:0], -2^{31}), 2^{31}-1)
    QACC_L[ 63: 32] = min(max(add_1[32:0], -2^{31}), 2^{31}-1)
    QACC_L[ 95: 64] = min(max(add_2[32:0], -2^{31}), 2^{31}-1)
    QACC_L[127: 96] = min(max(add_3[32:0], -2^{31}), 2^{31}-1)
    QACC_L[159:128] = min(max(add_4[32:0], -2^{31}), 2^{31}-1)
    QACC_L[191:160] = min(max(add_5[32:0], -2^{31}), 2^{31}-1)
    QACC_L[223:192] = min(max(add_6[32:0], -2^{31}), 2^{31}-1)
    QACC_L[255:224] = min(max(add_7[32:0], -2^{31}), 2^{31}-1)

    qu[127:0] = load128(rs1[31:0])
    rs1[31:0] = rs1[31:0] + imm16
\end{lstlisting}

\textbf{Schedule}\\
\begin{longtable}[c]{|l|l|l|}
    \hline
    \rowcolor{lightgray}
    \textbf{Operands} & \textbf{Use} & \textbf{Def} \\\hline
    \endfirsthead
    \hline
    \rowcolor{lightgray}
     \textbf{Operands} & \textbf{Use} & \textbf{Def} \\\hline
    \endhead
    \hline
    \endfoot

    qu & - & 2 \\\hline
    rs1 & 1 & 1 \\\hline
     qx & 1 & - \\\hline
     qy & 1 & - \\\hline
     QACC\_L & 2 & 2

\end{longtable}

\textbf{Format}\\

\begin{figure*}[!htbp]
 {\setlength{\tabcolsep}{1pt}
 \begin{center}
\begin{tabular}{@{}F@{}F@{}W@{}I@{}Y@{}I@{}F@{}W@{}F@{}F@{}O}
\instbitrange{31}{29}&
\instbitrange{28}{26}&
\instbitrange{25}{24}&
\instbit{23}&
\instbitrange{22}{19}&
\instbit{18}&
\instbitrange{17}{15}&
\instbitrange{14}{13}&
\instbitrange{12}{10}&
\instbitrange{9}{7}&
\instbitrange{6}{0} \\
\hline
\multicolumn{1}{|c|}{{\scriptsize qx[2:0]}} &
\multicolumn{1}{c|}{{\scriptsize qy[2:0]}} &
\multicolumn{1}{c|}{00} &
\multicolumn{1}{c|}{{\scriptsize sat[0]}} &
\multicolumn{1}{c|}{{\scriptsize imm16[3:0]}} &
\multicolumn{1}{c|}{{\scriptsize rs1[4]}} &
\multicolumn{1}{c|}{{\scriptsize rs1[2:0]}} &
\multicolumn{1}{c|}{01} &
\multicolumn{1}{c|}{{\scriptsize qu[2:0]}} &
\multicolumn{1}{c|}{1X1} &
\multicolumn{1}{c|}{0011111}\\
\hline
3& 3& 2& 1& 4& 1& 3& 2& 3& 3& 7 \\
\end{tabular}
 \end{center}
 }
 \vspace{-0.1in}
 \caption[]{ESP.VCMULAS.S8.QACC.L.LD.IP} \end{figure*}


\newpage
\subsubsection{ESP.VCMULAS.S8.QACC.L.LD.XP}\label{instruction:ESPVCMULASS8QACC_LLDXP}
\textbf{Syntax}\\
ESP.VCMULAS.S8.QACC.L.LD.XP qu,rs1,rs2,qx,qy,sat

\textbf{Description}\\
This instruction performs multiplication of 8-bit vectors and stores the result in \lstinline{QACC_L}.

During the operation, it loads the 16-byte data from the memory into the register \lstinline{qu}. After the access is completed, the value in the register \lstinline{rs1} is incremented by the value of the register \lstinline{rs2}.


\textbf{Operation}\\
\begin{lstlisting}
    add_0[32:0] = QACC_L[ 31:  0] + (qx[  7:  0]*qy[  7:  0] - qx[ 15:  8]*qy[ 15:  8])
    add_1[32:0] = QACC_L[ 63: 32] + (qx[ 15:  8]*qy[  7:  0] - qx[  7:  0]*qy[ 15:  8])

    add_2[32:0] = QACC_L[ 95: 64] + (qx[ 23: 16]*qy[ 23: 16] - qx[ 31: 24]*qy[ 31: 24])
    add_3[32:0] = QACC_L[127: 96] + (qx[ 31: 24]*qy[ 23: 16] - qx[ 23: 16]*qy[ 31: 24])

    add_4[32:0] = QACC_L[159:128] + (qx[ 39: 32]*qy[ 39: 32] - qx[ 47: 40]*qy[ 47: 40])
    add_5[32:0] = QACC_L[191:160] + (qx[ 47: 40]*qy[ 39: 32] - qx[ 39: 32]*qy[ 47: 40])

    add_6[32:0] = QACC_L[223:192] + (qx[ 55: 48]*qy[ 55: 48] - qx[ 63: 56]*qy[ 63: 56])
    add_7[32:0] = QACC_L[255:224] + (qx[ 63: 56]*qy[ 55: 48] - qx[ 55: 48]*qy[ 63: 56])

    QACC_L[ 31:  0] = min(max(add_0[32:0], -2^{31}), 2^{31}-1)
    QACC_L[ 63: 32] = min(max(add_1[32:0], -2^{31}), 2^{31}-1)
    QACC_L[ 95: 64] = min(max(add_2[32:0], -2^{31}), 2^{31}-1)
    QACC_L[127: 96] = min(max(add_3[32:0], -2^{31}), 2^{31}-1)
    QACC_L[159:128] = min(max(add_4[32:0], -2^{31}), 2^{31}-1)
    QACC_L[191:160] = min(max(add_5[32:0], -2^{31}), 2^{31}-1)
    QACC_L[223:192] = min(max(add_6[32:0], -2^{31}), 2^{31}-1)
    QACC_L[255:224] = min(max(add_7[32:0], -2^{31}), 2^{31}-1)

    qu[127:0] = load128(rs1[31:0])
    rs1[31:0] = rs1[31:0] + rs2[31:0]
\end{lstlisting}

\textbf{Schedule}\\
\begin{longtable}[c]{|l|l|l|}
    \hline
    \rowcolor{lightgray}
    \textbf{Operands} & \textbf{Use} & \textbf{Def} \\\hline
    \endfirsthead
    \hline
    \rowcolor{lightgray}
     \textbf{Operands} & \textbf{Use} & \textbf{Def} \\\hline
    \endhead
    \hline
    \endfoot

    qu & - & 2 \\\hline
    rs1 & 1 & 1 \\\hline
    rs2 & 1 & - \\\hline
     qx & 1 & - \\\hline
     qy & 1 & - \\\hline
     QACC\_L & 2 & 2

\end{longtable}

\textbf{Format}\\

\begin{figure*}[!htbp]
 {\setlength{\tabcolsep}{1pt}
 \begin{center}
\begin{tabular}{@{}F@{}F@{}W@{}I@{}F@{}I@{}I@{}F@{}W@{}F@{}F@{}O}
\instbitrange{31}{29}&
\instbitrange{28}{26}&
\instbitrange{25}{24}&
\instbit{23}&
\instbitrange{22}{20}&
\instbit{19}&
\instbit{18}&
\instbitrange{17}{15}&
\instbitrange{14}{13}&
\instbitrange{12}{10}&
\instbitrange{9}{7}&
\instbitrange{6}{0} \\
\hline
\multicolumn{1}{|c|}{{\scriptsize qx[2:0]}} &
\multicolumn{1}{c|}{{\scriptsize qy[2:0]}} &
\multicolumn{1}{c|}{00} &
\multicolumn{1}{c|}{{\scriptsize rs2[4]}} &
\multicolumn{1}{c|}{{\scriptsize rs2[2:0]}} &
\multicolumn{1}{c|}{{\scriptsize sat[0]}} &
\multicolumn{1}{c|}{{\scriptsize rs1[4]}} &
\multicolumn{1}{c|}{{\scriptsize rs1[2:0]}} &
\multicolumn{1}{c|}{10} &
\multicolumn{1}{c|}{{\scriptsize qu[2:0]}} &
\multicolumn{1}{c|}{111} &
\multicolumn{1}{c|}{0011111}\\
\hline
3& 3& 2& 1& 3& 1& 1& 3& 2& 3& 3& 7 \\
\end{tabular}
 \end{center}
 }
 \vspace{-0.1in}
 \caption[]{ESP.VCMULAS.S8.QACC.L.LD.XP} \end{figure*}


\newpage
\subsubsection{ESP.VMULAS.S16.QACC}\label{instruction:ESPVMULASS16QACC}
\textbf{Syntax}\\
ESP.VMULAS.S16.QACC qx,qy,sat

\textbf{Description}\\
This instruction divides registers \lstinline{qx} and \lstinline{qy} into 8 data segments of 16 bits each. Then, the signed multiplication result of the 8 sets of segments is added to the corresponding 64-bit data segment in the special registers \lstinline{QACC_H} and \lstinline{QACC_L} respectively. The calculated result is saturated to a 64-bit signed number and then stored in the corresponding 64-bit data segment in \lstinline{QACC_H} and \lstinline{QACC_L}.

\textbf{Operation}\\
\begin{lstlisting}
QACC_L[63 :  0] = min(max(QACC_L[63 :  0] + qx[15 :  0] * qy[15 :  0], -2^63), 2^63-1)
QACC_L[127: 64] = min(max(QACC_L[127: 64] + qx[31 : 16] * qy[31 : 16], -2^63), 2^63-1)
QACC_L[191:128] = min(max(QACC_L[191:128] + qx[47 : 32] * qy[47 : 32], -2^63), 2^63-1)
QACC_L[255:192] = min(max(QACC_L[255:192] + qx[63 : 48] * qy[63 : 48], -2^63), 2^63-1)
QACC_H[63 :  0] = min(max(QACC_H[63 :  0] + qx[79 : 64] * qy[79 : 64], -2^63), 2^63-1)
QACC_H[127: 64] = min(max(QACC_H[127: 64] + qx[95 : 80] * qy[95 : 80], -2^63), 2^63-1)
QACC_H[191:128] = min(max(QACC_H[191:128] + qx[111: 96] * qy[111: 96], -2^63), 2^63-1)
QACC_H[255:192] = min(max(QACC_H[255:192] + qx[127:112] * qy[127:112], -2^63), 2^63-1)
\end{lstlisting}

\textbf{Schedule}\\
\begin{longtable}[c]{|l|l|l|}
    \hline
    \rowcolor{lightgray}
    \textbf{Operands} & \textbf{Use} & \textbf{Def} \\\hline
    \endfirsthead
    \hline
    \rowcolor{lightgray}
     \textbf{Operands} & \textbf{Use} & \textbf{Def} \\\hline
    \endhead
    \hline
    \endfoot


     qx & 1 & - \\\hline
     qy & 1 & - \\\hline
     QACC\_H & 2 & 2 \\\hline
     QACC\_L & 2 & 2

\end{longtable}

\textbf{Format}\\

\begin{figure*}[!htbp]
 {\setlength{\tabcolsep}{1pt}
 \begin{center}
\begin{tabular}{@{}F@{}F@{}T@{}I@{}E@{}O}
\instbitrange{31}{29}&
\instbitrange{28}{26}&
\instbitrange{25}{16}&
\instbit{15}&
\instbitrange{14}{7}&
\instbitrange{6}{0} \\
\hline
\multicolumn{1}{|c|}{{\scriptsize qx[2:0]}} &
\multicolumn{1}{c|}{{\scriptsize qy[2:0]}} &
\multicolumn{1}{c|}{1X101XX111} &
\multicolumn{1}{c|}{{\scriptsize sat[0]}} &
\multicolumn{1}{c|}{01XX0X0X} &
\multicolumn{1}{c|}{0011011}\\
\hline
3& 3& 10& 1& 8& 7 \\
\end{tabular}
 \end{center}
 }
 \vspace{-0.1in}
 \caption[]{ESP.VMULAS.S16.QACC} \end{figure*}


\newpage
\subsubsection{ESP.VMULAS.S16.QACC.LD.IP}\label{instruction:ESPVMULASS16QACCLDIP}
\textbf{Syntax}\\
ESP.VMULAS.S16.QACC.LD.IP qu,rs1,imm16,qx,qy,sat

\textbf{Description}\\
This instruction divides registers \lstinline{qx} and \lstinline{qy} into 8 data segments of 16 bits each. Then, the signed multiplication result of the 8 sets of segments is added to the corresponding 64-bit data segment in special registers \lstinline{QACC_H} and \lstinline{QACC_L} respectively. The calculated result is saturated to a 64-bit signed number and then stored in the corresponding 64-bit data segment in \lstinline{QACC_H} and \lstinline{QACC_L}.

During the operation, it loads the 16-byte data from the memory into the register \lstinline{qu}. After the access is completed, the value in the register \lstinline{rs1} is incremented by the 4-bit sign-extended constant in the instruction code segment, which is left-shifted by 4.

The value of imm starts at -128, ends at 112, and steps by 16.

\textbf{Operation}\\
\begin{lstlisting}
QACC_L[63 :  0] = min(max(QACC_L[63 :  0] + qx[15 :  0] * qy[15 :  0], -2^63), 2^63-1)
QACC_L[127: 64] = min(max(QACC_L[127: 64] + qx[31 : 16] * qy[31 : 16], -2^63), 2^63-1)
QACC_L[191:128] = min(max(QACC_L[191:128] + qx[47 : 32] * qy[47 : 32], -2^63), 2^63-1)
QACC_L[255:192] = min(max(QACC_L[255:192] + qx[63 : 48] * qy[63 : 48], -2^63), 2^63-1)
QACC_H[63 :  0] = min(max(QACC_H[63 :  0] + qx[79 : 64] * qy[79 : 64], -2^63), 2^63-1)
QACC_H[127: 64] = min(max(QACC_H[127: 64] + qx[95 : 80] * qy[95 : 80], -2^63), 2^63-1)
QACC_H[191:128] = min(max(QACC_H[191:128] + qx[111: 96] * qy[111: 96], -2^63), 2^63-1)
QACC_H[255:192] = min(max(QACC_H[255:192] + qx[127:112] * qy[127:112], -2^63), 2^63-1)

qu[127:0] = load128(rs1[31:0])
rs1[31:0] = rs1[31:0] + imm
\end{lstlisting}

\textbf{Schedule}\\
\begin{longtable}[c]{|l|l|l|}
    \hline
    \rowcolor{lightgray}
    \textbf{Operands} & \textbf{Use} & \textbf{Def} \\\hline
    \endfirsthead
    \hline
    \rowcolor{lightgray}
     \textbf{Operands} & \textbf{Use} & \textbf{Def} \\\hline
    \endhead
    \hline
    \endfoot

    qu & - & 2 \\\hline
    rs1 & 1 & 1 \\\hline
     qx & 1 & - \\\hline
     qy & 1 & - \\\hline
     QACC\_H & 2 & 2 \\\hline
     QACC\_L & 2 & 2

\end{longtable}

\textbf{Format}\\

\begin{figure*}[!htbp]
 {\setlength{\tabcolsep}{1pt}
 \begin{center}
\begin{tabular}{@{}F@{}F@{}Y@{}I@{}W@{}I@{}F@{}W@{}F@{}W@{}I@{}O}
\instbitrange{31}{29}&
\instbitrange{28}{26}&
\instbitrange{25}{22}&
\instbit{21}&
\instbitrange{20}{19}&
\instbit{18}&
\instbitrange{17}{15}&
\instbitrange{14}{13}&
\instbitrange{12}{10}&
\instbitrange{9}{8}&
\instbit{7}&
\instbitrange{6}{0} \\
\hline
\multicolumn{1}{|c|}{{\scriptsize qx[2:0]}} &
\multicolumn{1}{c|}{{\scriptsize qy[2:0]}} &
\multicolumn{1}{c|}{1011} &
\multicolumn{1}{c|}{{\scriptsize sat[0]}} &
\multicolumn{1}{c|}{{\scriptsize imm16[3:2]}} &
\multicolumn{1}{c|}{{\scriptsize rs1[4]}} &
\multicolumn{1}{c|}{{\scriptsize rs1[2:0]}} &
\multicolumn{1}{c|}{11} &
\multicolumn{1}{c|}{{\scriptsize qu[2:0]}} &
\multicolumn{1}{c|}{{\scriptsize imm16[1:0]}} &
\multicolumn{1}{c|}{0} &
\multicolumn{1}{c|}{0011111}\\
\hline
3& 3& 4& 1& 2& 1& 3& 2& 3& 2& 1& 7 \\
\end{tabular}
 \end{center}
 }
 \vspace{-0.1in}
 \caption[]{ESP.VMULAS.S16.QACC.LD.IP} \end{figure*}


\newpage
\subsubsection{ESP.VMULAS.S16.QACC.LD.XP}\label{instruction:ESPVMULASS16QACCLDXP}
\textbf{Syntax}\\
ESP.VMULAS.S16.QACC.LD.XP qu,rs1,rs2,qx,qy,sat

\textbf{Description}\\
This instruction divides registers \lstinline{qx} and \lstinline{qy} into 8 data segments of 16 bits each. Then, the signed multiplication result of the 8 sets of segments is added to the corresponding 64-bit data segment in special registers \lstinline{QACC_H} and \lstinline{QACC_L} respectively. The calculated result is saturated to a 64-bit signed number and then stored in the corresponding 64-bit data segment in \lstinline{QACC_H} and \lstinline{QACC_L}.

During the operation, it loads the 16-byte data from the memory into the register \lstinline{qu}. After the access is completed, the value in the register \lstinline{rs1} is incremented by the value of the register \lstinline{rs2}.

\textbf{Operation}\\
\begin{lstlisting}
QACC_L[63 :  0] = min(max(QACC_L[63 :  0] + qx[15 :  0] * qy[15 :  0], -2^63), 2^63-1)
QACC_L[127: 64] = min(max(QACC_L[127: 64] + qx[31 : 16] * qy[31 : 16], -2^63), 2^63-1)
QACC_L[191:128] = min(max(QACC_L[191:128] + qx[47 : 32] * qy[47 : 32], -2^63), 2^63-1)
QACC_L[255:192] = min(max(QACC_L[255:192] + qx[63 : 48] * qy[63 : 48], -2^63), 2^63-1)
QACC_H[63 :  0] = min(max(QACC_H[63 :  0] + qx[79 : 64] * qy[79 : 64], -2^63), 2^63-1)
QACC_H[127: 64] = min(max(QACC_H[127: 64] + qx[95 : 80] * qy[95 : 80], -2^63), 2^63-1)
QACC_H[191:128] = min(max(QACC_H[191:128] + qx[111: 96] * qy[111: 96], -2^63), 2^63-1)
QACC_H[255:192] = min(max(QACC_H[255:192] + qx[127:112] * qy[127:112], -2^63), 2^63-1)

qu[127:0] = load128(rs1[31:0])
rs1[31:0] = rs1[31:0] + rs2[31:0]
\end{lstlisting}

\textbf{Schedule}\\
\begin{longtable}[c]{|l|l|l|}
    \hline
    \rowcolor{lightgray}
    \textbf{Operands} & \textbf{Use} & \textbf{Def} \\\hline
    \endfirsthead
    \hline
    \rowcolor{lightgray}
     \textbf{Operands} & \textbf{Use} & \textbf{Def} \\\hline
    \endhead
    \hline
    \endfoot

    qu & - & 2 \\\hline
    rs1 & 1 & 1 \\\hline
    rs2 & 1 & - \\\hline
     qx & 1 & - \\\hline
     qy & 1 & - \\\hline
     QACC\_H & 2 & 2 \\\hline
     QACC\_L & 2 & 2

\end{longtable}

\textbf{Format}\\

\begin{figure*}[!htbp]
 {\setlength{\tabcolsep}{1pt}
 \begin{center}
\begin{tabular}{@{}F@{}F@{}W@{}I@{}F@{}I@{}I@{}F@{}W@{}F@{}I@{}I@{}I@{}O}
\instbitrange{31}{29}&
\instbitrange{28}{26}&
\instbitrange{25}{24}&
\instbit{23}&
\instbitrange{22}{20}&
\instbit{19}&
\instbit{18}&
\instbitrange{17}{15}&
\instbitrange{14}{13}&
\instbitrange{12}{10}&
\instbit{9}&
\instbit{8}&
\instbit{7}&
\instbitrange{6}{0} \\
\hline
\multicolumn{1}{|c|}{{\scriptsize qx[2:0]}} &
\multicolumn{1}{c|}{{\scriptsize qy[2:0]}} &
\multicolumn{1}{c|}{10} &
\multicolumn{1}{c|}{{\scriptsize rs2[4]}} &
\multicolumn{1}{c|}{{\scriptsize rs2[2:0]}} &
\multicolumn{1}{c|}{1} &
\multicolumn{1}{c|}{{\scriptsize rs1[4]}} &
\multicolumn{1}{c|}{{\scriptsize rs1[2:0]}} &
\multicolumn{1}{c|}{01} &
\multicolumn{1}{c|}{{\scriptsize qu[2:0]}} &
\multicolumn{1}{c|}{1} &
\multicolumn{1}{c|}{{\scriptsize sat[0]}} &
\multicolumn{1}{c|}{0} &
\multicolumn{1}{c|}{0011111}\\
\hline
3& 3& 2& 1& 3& 1& 1& 3& 2& 3& 1& 1& 1& 7 \\
\end{tabular}
 \end{center}
 }
 \vspace{-0.1in}
 \caption[]{ESP.VMULAS.S16.QACC.LD.XP} \end{figure*}


\newpage
\subsubsection{ESP.VMULAS.S16.QACC.ST.IP}\label{instruction:ESPVMULASS16QACCSTIP}
\textbf{Syntax}\\
ESP.VMULAS.S16.QACC.ST.IP qu,rs1,imm16,qx,qy,sat

\textbf{Description}\\
This instruction divides registers \lstinline{qx} and \lstinline{qy} into 8 data segments of 16 bits each. Then, the signed multiplication results of the 8 sets of segments are added to the corresponding 64-bit data segments in special registers \lstinline{QACC_H} and \lstinline{QACC_L} respectively. The calculated result is saturated to a 64-bit signed number and then stored in the corresponding 64-bit data segment in \lstinline{QACC_H} and \lstinline{QACC_L}.

During the operation, the value of the register \lstinline{qu} is stored in memory. After the access is completed, the value in the register \lstinline{rs1} is incremented by a 4-bit sign-extended constant in the instruction code segment, which is left-shifted by 4.

The imm starts at -128, ends at 112, and steps by 16.

\textbf{Operation}\\
\begin{lstlisting}
QACC_L[63 :  0] = min(max(QACC_L[63 :  0] + qx[15 :  0] * qy[15 :  0], -2^63), 2^63-1)
QACC_L[127: 64] = min(max(QACC_L[127: 64] + qx[31 : 16] * qy[31 : 16], -2^63), 2^63-1)
QACC_L[191:128] = min(max(QACC_L[191:128] + qx[47 : 32] * qy[47 : 32], -2^63), 2^63-1)
QACC_L[255:192] = min(max(QACC_L[255:192] + qx[63 : 48] * qy[63 : 48], -2^63), 2^63-1)
QACC_H[63 :  0] = min(max(QACC_H[63 :  0] + qx[79 : 64] * qy[79 : 64], -2^63), 2^63-1)
QACC_H[127: 64] = min(max(QACC_H[127: 64] + qx[95 : 80] * qy[95 : 80], -2^63), 2^63-1)
QACC_H[191:128] = min(max(QACC_H[191:128] + qx[111: 96] * qy[111: 96], -2^63), 2^63-1)
QACC_H[255:192] = min(max(QACC_H[255:192] + qx[127:112] * qy[127:112], -2^63), 2^63-1)

qu[127:0] => store128(rs1[31:0])
rs1[31:0] = rs1[31:0] + imm
\end{lstlisting}

\textbf{Schedule}\\
\begin{longtable}[c]{|l|l|l|}
    \hline
    \rowcolor{lightgray}
    \textbf{Operands} & \textbf{Use} & \textbf{Def} \\\hline
    \endfirsthead
    \hline
    \rowcolor{lightgray}
     \textbf{Operands} & \textbf{Use} & \textbf{Def} \\\hline
    \endhead
    \hline
    \endfoot

    qu & 2 & - \\\hline
    rs1 & 1 & 1 \\\hline

     qx & 1 & - \\\hline
     qy & 1 & - \\\hline
     QACC\_H & 2 & 2 \\\hline
     QACC\_L & 2 & 2

\end{longtable}

\textbf{Format}\\

\begin{figure*}[!htbp]
 {\setlength{\tabcolsep}{1pt}
 \begin{center}
\begin{tabular}{@{}F@{}F@{}Y@{}I@{}W@{}I@{}F@{}W@{}F@{}W@{}I@{}O}
\instbitrange{31}{29}&
\instbitrange{28}{26}&
\instbitrange{25}{22}&
\instbit{21}&
\instbitrange{20}{19}&
\instbit{18}&
\instbitrange{17}{15}&
\instbitrange{14}{13}&
\instbitrange{12}{10}&
\instbitrange{9}{8}&
\instbit{7}&
\instbitrange{6}{0} \\
\hline
\multicolumn{1}{|c|}{{\scriptsize qx[2:0]}} &
\multicolumn{1}{c|}{{\scriptsize qy[2:0]}} &
\multicolumn{1}{c|}{1111} &
\multicolumn{1}{c|}{{\scriptsize sat[0]}} &
\multicolumn{1}{c|}{{\scriptsize imm16[3:2]}} &
\multicolumn{1}{c|}{{\scriptsize rs1[4]}} &
\multicolumn{1}{c|}{{\scriptsize rs1[2:0]}} &
\multicolumn{1}{c|}{11} &
\multicolumn{1}{c|}{{\scriptsize qu[2:0]}} &
\multicolumn{1}{c|}{{\scriptsize imm16[1:0]}} &
\multicolumn{1}{c|}{0} &
\multicolumn{1}{c|}{0011111}\\
\hline
3& 3& 4& 1& 2& 1& 3& 2& 3& 2& 1& 7 \\
\end{tabular}
 \end{center}
 }
 \vspace{-0.1in}
 \caption[]{ESP.VMULAS.S16.QACC.ST.IP} \end{figure*}


\newpage
\subsubsection{ESP.VMULAS.S16.QACC.ST.XP}\label{instruction:ESPVMULASS16QACCSTXP}
\textbf{Syntax}\\
ESP.VMULAS.S16.QACC.ST.XP qu,rs1,rs2,qx,qy,sat

\textbf{Description}\\
This instruction divides registers \lstinline{qx} and \lstinline{qy} into 8 data segments of 16 bits each. Then, the signed multiplication results of the 8 sets of segments are added to the corresponding 64-bit data segments in special registers \lstinline{QACC_H} and \lstinline{QACC_L} respectively. The calculated results are saturated to a 64-bit signed number and then stored in the corresponding 64-bit data segments in \lstinline{QACC_H} and \lstinline{QACC_L}.

During the operation, the value of the register \lstinline{qu} is stored in memory. After the access is completed, the value in the register \lstinline{rs1} is incremented by the value of the register \lstinline{rs2}.

\textbf{Operation}\\
\begin{lstlisting}
QACC_L[63 :  0] = min(max(QACC_L[63 :  0] + qx[15 :  0] * qy[15 :  0], -2^63), 2^63-1)
QACC_L[127: 64] = min(max(QACC_L[127: 64] + qx[31 : 16] * qy[31 : 16], -2^63), 2^63-1)
QACC_L[191:128] = min(max(QACC_L[191:128] + qx[47 : 32] * qy[47 : 32], -2^63), 2^63-1)
QACC_L[255:192] = min(max(QACC_L[255:192] + qx[63 : 48] * qy[63 : 48], -2^63), 2^63-1)
QACC_H[63 :  0] = min(max(QACC_H[63 :  0] + qx[79 : 64] * qy[79 : 64], -2^63), 2^63-1)
QACC_H[127: 64] = min(max(QACC_H[127: 64] + qx[95 : 80] * qy[95 : 80], -2^63), 2^63-1)
QACC_H[191:128] = min(max(QACC_H[191:128] + qx[111: 96] * qy[111: 96], -2^63), 2^63-1)
QACC_H[255:192] = min(max(QACC_H[255:192] + qx[127:112] * qy[127:112], -2^63), 2^63-1)

qu[127:0] => store128(rs1[31:0])
rs1[31:0] = rs1[31:0] + rs2[31:0]
\end{lstlisting}

\textbf{Schedule}\\
\begin{longtable}[c]{|l|l|l|}
    \hline
    \rowcolor{lightgray}
    \textbf{Operands} & \textbf{Use} & \textbf{Def} \\\hline
    \endfirsthead
    \hline
    \rowcolor{lightgray}
     \textbf{Operands} & \textbf{Use} & \textbf{Def} \\\hline
    \endhead
    \hline
    \endfoot

    qu & 2 & - \\\hline
    rs1 & 1 & 1 \\\hline
    rs2 & 1 & - \\\hline
     qx & 1 & - \\\hline
     qy & 1 & - \\\hline
     QACC\_H & 2 & 2 \\\hline
     QACC\_L & 2 & 2

\end{longtable}

\textbf{Format}\\

\begin{figure*}[!htbp]
 {\setlength{\tabcolsep}{1pt}
 \begin{center}
\begin{tabular}{@{}F@{}F@{}W@{}I@{}F@{}I@{}I@{}F@{}W@{}F@{}I@{}I@{}I@{}O}
\instbitrange{31}{29}&
\instbitrange{28}{26}&
\instbitrange{25}{24}&
\instbit{23}&
\instbitrange{22}{20}&
\instbit{19}&
\instbit{18}&
\instbitrange{17}{15}&
\instbitrange{14}{13}&
\instbitrange{12}{10}&
\instbit{9}&
\instbit{8}&
\instbit{7}&
\instbitrange{6}{0} \\
\hline
\multicolumn{1}{|c|}{{\scriptsize qx[2:0]}} &
\multicolumn{1}{c|}{{\scriptsize qy[2:0]}} &
\multicolumn{1}{c|}{11} &
\multicolumn{1}{c|}{{\scriptsize rs2[4]}} &
\multicolumn{1}{c|}{{\scriptsize rs2[2:0]}} &
\multicolumn{1}{c|}{1} &
\multicolumn{1}{c|}{{\scriptsize rs1[4]}} &
\multicolumn{1}{c|}{{\scriptsize rs1[2:0]}} &
\multicolumn{1}{c|}{01} &
\multicolumn{1}{c|}{{\scriptsize qu[2:0]}} &
\multicolumn{1}{c|}{1} &
\multicolumn{1}{c|}{{\scriptsize sat[0]}} &
\multicolumn{1}{c|}{0} &
\multicolumn{1}{c|}{0011111}\\
\hline
3& 3& 2& 1& 3& 1& 1& 3& 2& 3& 1& 1& 1& 7 \\
\end{tabular}
 \end{center}
 }
 \vspace{-0.1in}
 \caption[]{ESP.VMULAS.S16.QACC.ST.XP} \end{figure*}


\newpage
\subsubsection{ESP.VMULAS.S16.XACC}\label{instruction:ESPVMULASS16XACC}
\textbf{Syntax}\\
ESP.VMULAS.S16.XACC qx,qy,sat

\textbf{Description}\\
This instruction divides registers \lstinline{qx} and \lstinline{qy} into 8 data segments of 16 bits each. Then, it performs a signed multiply-accumulate operation on each of the 8 sets of segments. The accumulated result is added to the value in the special register \lstinline{XACC}. Then, the resulting sum is saturated and stored in \lstinline{XACC}.

\textbf{Operation}\\
\begin{lstlisting}
    res[39 :  0] = qx[15 :  0] * qy[15 :  0] + qx[31 : 16] * qy[31 : 16] + … + qx[127:112] * qy[127:112]
    XACC = min(max(XACC + res[39 :  0], -2^39), 2^39-1)
\end{lstlisting}

\textbf{Schedule}\\
\begin{longtable}[c]{|l|l|l|}
    \hline
    \rowcolor{lightgray}
    \textbf{Operands} & \textbf{Use} & \textbf{Def} \\\hline
    \endfirsthead
    \hline
    \rowcolor{lightgray}
     \textbf{Operands} & \textbf{Use} & \textbf{Def} \\\hline
    \endhead
    \hline
    \endfoot


     qx & 1 & - \\\hline
     qy & 1 & - \\\hline
     XACC & 3 & 3

\end{longtable}

\textbf{Format}\\

\begin{figure*}[!htbp]
 {\setlength{\tabcolsep}{1pt}
 \begin{center}
\begin{tabular}{@{}F@{}F@{}T@{}I@{}E@{}O}
\instbitrange{31}{29}&
\instbitrange{28}{26}&
\instbitrange{25}{16}&
\instbit{15}&
\instbitrange{14}{7}&
\instbitrange{6}{0} \\
\hline
\multicolumn{1}{|c|}{{\scriptsize qx[2:0]}} &
\multicolumn{1}{c|}{{\scriptsize qy[2:0]}} &
\multicolumn{1}{c|}{1X101XX011} &
\multicolumn{1}{c|}{{\scriptsize sat[0]}} &
\multicolumn{1}{c|}{01XX0X0X} &
\multicolumn{1}{c|}{0011011}\\
\hline
3& 3& 10& 1& 8& 7 \\
\end{tabular}
 \end{center}
 }
 \vspace{-0.1in}
 \caption[]{ESP.VMULAS.S16.XACC} \end{figure*}


\newpage
\subsubsection{ESP.VMULAS.S16.XACC.LD.IP}\label{instruction:ESPVMULASS16XACCLDIP}
\textbf{Syntax}\\
ESP.VMULAS.S16.XACC.LD.IP qu,rs1,imm16,qx,qy,sat

\textbf{Description}\\
This instruction divides registers \lstinline{qx} and \lstinline{qy} into 8 data segments of 16 bits each. Then, it performs a signed multiply-accumulate operation on the 8 sets of segments respectively. The accumulated result is added to the value in the special register \lstinline{XACC}. Then, the sum obtained is saturated and stored in \lstinline{XACC}.

During the operation, the 16-byte data is loaded from the memory to the register \lstinline{qu}. After the access is completed, the value in the register \lstinline{rs1} is incremented by a 4-bit sign-extended constant in the instruction code segment, which is left-shifted by 4.

The value of imm starts at -128, ends at 112, and steps by 16.

\textbf{Operation}\\
\begin{lstlisting}
    res[39 :  0] = qx[15 :  0] * qy[15 :  0] + qx[31 : 16] * qy[31 : 16] + … + qx[127:112] * qy[127:112]
    XACC = min(max(XACC + res[39 :  0], -2^39), 2^39-1)
    qu[127:0] = load128(rs1[31:0])
    rs1[31:0] = rs1[31:0] + imm
\end{lstlisting}

\textbf{Schedule}\\
\begin{longtable}[c]{|l|l|l|}
    \hline
    \rowcolor{lightgray}
    \textbf{Operands} & \textbf{Use} & \textbf{Def} \\\hline
    \endfirsthead
    \hline
    \rowcolor{lightgray}
     \textbf{Operands} & \textbf{Use} & \textbf{Def} \\\hline
    \endhead
    \hline
    \endfoot

    qu & - & 2 \\\hline
    rs1 & 1 & 1 \\\hline
     qx & 1 & - \\\hline
     qy & 1 & - \\\hline
     XACC & 3 & 3

\end{longtable}

\textbf{Format}\\

\begin{figure*}[!htbp]
 {\setlength{\tabcolsep}{1pt}
 \begin{center}
\begin{tabular}{@{}F@{}F@{}Y@{}I@{}W@{}I@{}F@{}W@{}F@{}W@{}I@{}O}
\instbitrange{31}{29}&
\instbitrange{28}{26}&
\instbitrange{25}{22}&
\instbit{21}&
\instbitrange{20}{19}&
\instbit{18}&
\instbitrange{17}{15}&
\instbitrange{14}{13}&
\instbitrange{12}{10}&
\instbitrange{9}{8}&
\instbit{7}&
\instbitrange{6}{0} \\
\hline
\multicolumn{1}{|c|}{{\scriptsize qx[2:0]}} &
\multicolumn{1}{c|}{{\scriptsize qy[2:0]}} &
\multicolumn{1}{c|}{0011} &
\multicolumn{1}{c|}{{\scriptsize sat[0]}} &
\multicolumn{1}{c|}{{\scriptsize imm16[3:2]}} &
\multicolumn{1}{c|}{{\scriptsize rs1[4]}} &
\multicolumn{1}{c|}{{\scriptsize rs1[2:0]}} &
\multicolumn{1}{c|}{11} &
\multicolumn{1}{c|}{{\scriptsize qu[2:0]}} &
\multicolumn{1}{c|}{{\scriptsize imm16[1:0]}} &
\multicolumn{1}{c|}{0} &
\multicolumn{1}{c|}{0011111}\\
\hline
3& 3& 4& 1& 2& 1& 3& 2& 3& 2& 1& 7 \\
\end{tabular}
 \end{center}
 }
 \vspace{-0.1in}
 \caption[]{ESP.VMULAS.S16.XACC.LD.IP} \end{figure*}


\newpage
\subsubsection{ESP.VMULAS.S16.XACC.LD.XP}\label{instruction:ESPVMULASS16XACCLDXP}
\textbf{Syntax}\\
ESP.VMULAS.S16.XACC.LD.XP qu,rs1,rs2,qx,qy,sat

\textbf{Description}\\
This instruction divides registers \lstinline{qx} and \lstinline{qy} into 8 data segments of 16 bits each. Then, it performs a signed multiply-accumulate operation on each of the 8 segment sets. The accumulated result is added to the value in the special register \lstinline{XACC}. The resulting sum is then saturated and stored in \lstinline{XACC}.

During the operation, the 16-byte data is loaded from the memory into the register \lstinline{qu}. After the access is completed, the value in the register \lstinline{rs1} is incremented by the value of the register \lstinline{rs2}.

\textbf{Operation}\\
\begin{lstlisting}
    res[39 :  0] = qx[15 :  0] * qy[15 :  0] + qx[31 : 16] * qy[31 : 16] + … + qx[127:112] * qy[127:112]
    XACC = min(max(XACC + res[39 :  0], -2^39), 2^39-1)
    qu[127:0] = load128(rs1[31:0])
    rs1[31:0] = rs1[31:0] + rs2[31:0]
\end{lstlisting}

\textbf{Schedule}\\
\begin{longtable}[c]{|l|l|l|}
    \hline
    \rowcolor{lightgray}
    \textbf{Operands} & \textbf{Use} & \textbf{Def} \\\hline
    \endfirsthead
    \hline
    \rowcolor{lightgray}
     \textbf{Operands} & \textbf{Use} & \textbf{Def} \\\hline
    \endhead
    \hline
    \endfoot

    qu & - & 2 \\\hline
    rs1 & 1 & 1 \\\hline
    rs2 & 1 & - \\\hline
     qx & 1 & - \\\hline
     qy & 1 & - \\\hline
     XACC & 3 & 3
\end{longtable}

\textbf{Format}\\

\begin{figure*}[!htbp]
 {\setlength{\tabcolsep}{1pt}
 \begin{center}
\begin{tabular}{@{}F@{}F@{}W@{}I@{}F@{}I@{}I@{}F@{}W@{}F@{}I@{}I@{}I@{}O}
\instbitrange{31}{29}&
\instbitrange{28}{26}&
\instbitrange{25}{24}&
\instbit{23}&
\instbitrange{22}{20}&
\instbit{19}&
\instbit{18}&
\instbitrange{17}{15}&
\instbitrange{14}{13}&
\instbitrange{12}{10}&
\instbit{9}&
\instbit{8}&
\instbit{7}&
\instbitrange{6}{0} \\
\hline
\multicolumn{1}{|c|}{{\scriptsize qx[2:0]}} &
\multicolumn{1}{c|}{{\scriptsize qy[2:0]}} &
\multicolumn{1}{c|}{00} &
\multicolumn{1}{c|}{{\scriptsize rs2[4]}} &
\multicolumn{1}{c|}{{\scriptsize rs2[2:0]}} &
\multicolumn{1}{c|}{1} &
\multicolumn{1}{c|}{{\scriptsize rs1[4]}} &
\multicolumn{1}{c|}{{\scriptsize rs1[2:0]}} &
\multicolumn{1}{c|}{01} &
\multicolumn{1}{c|}{{\scriptsize qu[2:0]}} &
\multicolumn{1}{c|}{1} &
\multicolumn{1}{c|}{{\scriptsize sat[0]}} &
\multicolumn{1}{c|}{0} &
\multicolumn{1}{c|}{0011111}\\
\hline
3& 3& 2& 1& 3& 1& 1& 3& 2& 3& 1& 1& 1& 7 \\
\end{tabular}
 \end{center}
 }
 \vspace{-0.1in}
 \caption[]{ESP.VMULAS.S16.XACC.LD.XP} \end{figure*}


\newpage
\subsubsection{ESP.VMULAS.S16.XACC.ST.IP}\label{instruction:ESPVMULASS16XACCSTIP}
\textbf{Syntax}\\
ESP.VMULAS.S16.XACC.ST.IP qu,rs1,imm16,qx,qy,sat

\textbf{Description}\\
This instruction divides registers \lstinline{qx} and \lstinline{qy} into 8 data segments of 16 bits each. Then, it performs a signed multiply-accumulate operation on the 8 sets of segments respectively. The accumulated result is added to the value in the special register \lstinline{XACC}. Then, the sum obtained is saturated and then stored in \lstinline{XACC}.

During the operation, the value of the register \lstinline{qu} is stored in the memory. After the access is completed, the value in the register \lstinline{rs1} is incremented by a 4-bit sign-extended constant in the instruction code segment, which is left-shifted by 4.

The value of imm starts at -128, ends at 112, and steps by 16.

\textbf{Operation}\\
\begin{lstlisting}
    res[39 :  0] = qx[15 :  0] * qy[15 :  0] + qx[31 : 16] * qy[31 : 16] + … + qx[127:112] * qy[127:112]
    XACC = min(max(XACC + res[39 :  0], -2^39), 2^39-1)
    qu[127:0] => store128(rs1[31:0])
    rs1[31:0] = rs1[31:0] + imm
\end{lstlisting}

\textbf{Schedule}\\
\begin{longtable}[c]{|l|l|l|}
    \hline
    \rowcolor{lightgray}
    \textbf{Operands} & \textbf{Use} & \textbf{Def} \\\hline
    \endfirsthead
    \hline
    \rowcolor{lightgray}
     \textbf{Operands} & \textbf{Use} & \textbf{Def} \\\hline
    \endhead
    \hline
    \endfoot

    qu & 2 & - \\\hline
    rs1 & 1 & 1 \\\hline

     qx & 1 & - \\\hline
     qy & 1 & - \\\hline
     XACC & 3 & 3

\end{longtable}

\textbf{Format}\\

\begin{figure*}[!htbp]
 {\setlength{\tabcolsep}{1pt}
 \begin{center}
\begin{tabular}{@{}F@{}F@{}Y@{}I@{}W@{}I@{}F@{}W@{}F@{}W@{}I@{}O}
\instbitrange{31}{29}&
\instbitrange{28}{26}&
\instbitrange{25}{22}&
\instbit{21}&
\instbitrange{20}{19}&
\instbit{18}&
\instbitrange{17}{15}&
\instbitrange{14}{13}&
\instbitrange{12}{10}&
\instbitrange{9}{8}&
\instbit{7}&
\instbitrange{6}{0} \\
\hline
\multicolumn{1}{|c|}{{\scriptsize qx[2:0]}} &
\multicolumn{1}{c|}{{\scriptsize qy[2:0]}} &
\multicolumn{1}{c|}{0111} &
\multicolumn{1}{c|}{{\scriptsize sat[0]}} &
\multicolumn{1}{c|}{{\scriptsize imm16[3:2]}} &
\multicolumn{1}{c|}{{\scriptsize rs1[4]}} &
\multicolumn{1}{c|}{{\scriptsize rs1[2:0]}} &
\multicolumn{1}{c|}{11} &
\multicolumn{1}{c|}{{\scriptsize qu[2:0]}} &
\multicolumn{1}{c|}{{\scriptsize imm16[1:0]}} &
\multicolumn{1}{c|}{0} &
\multicolumn{1}{c|}{0011111}\\
\hline
3& 3& 4& 1& 2& 1& 3& 2& 3& 2& 1& 7 \\
\end{tabular}
 \end{center}
 }
 \vspace{-0.1in}
 \caption[]{ESP.VMULAS.S16.XACC.ST.IP} \end{figure*}


\newpage
\subsubsection{ESP.VMULAS.S16.XACC.ST.XP}\label{instruction:ESPVMULASS16XACCSTXP}
\textbf{Syntax}\\
ESP.VMULAS.S16.XACC.ST.XP qu,rs1,rs2,qx,qy,sat

\textbf{Description}\\
This instruction divides registers \lstinline{qx} and \lstinline{qy} into 8 data segments of 16 bits each. Then, it performs a signed multiply-accumulate operation on each of the 8 segment sets. The accumulated result is added to the value in the special register \lstinline{XACC}. The resulting sum is then saturated and stored in \lstinline{XACC}.

During the operation, the value of the register \lstinline{qu} is stored in memory. After the access is completed, the value in the register \lstinline{rs1} is incremented by the value of the register \lstinline{rs2}.

\textbf{Operation}\\
\begin{lstlisting}
    res[39 :  0] = qx[15 :  0] * qy[15 :  0] + qx[31 : 16] * qy[31 : 16] + … + qx[127:112] * qy[127:112]
    XACC = min(max(XACC + res[39 :  0], -2^39), 2^39-1)
    qu[127:0] => store128(rs1[31:0])
    rs1[31:0] = rs1[31:0] + rs2[31:0]
\end{lstlisting}

\textbf{Schedule}\\
\begin{longtable}[c]{|l|l|l|}
    \hline
    \rowcolor{lightgray}
    \textbf{Operands} & \textbf{Use} & \textbf{Def} \\\hline
    \endfirsthead
    \hline
    \rowcolor{lightgray}
     \textbf{Operands} & \textbf{Use} & \textbf{Def} \\\hline
    \endhead
    \hline
    \endfoot

    qu & 2 & - \\\hline
    rs1 & 1 & 1 \\\hline
    rs2 & 1 & - \\\hline
     qx & 1 & - \\\hline
     qy & 1 & - \\\hline
     XACC & 3 & 3

\end{longtable}

\textbf{Format}\\

\begin{figure*}[!htbp]
 {\setlength{\tabcolsep}{1pt}
 \begin{center}
\begin{tabular}{@{}F@{}F@{}W@{}I@{}F@{}I@{}I@{}F@{}W@{}F@{}I@{}I@{}I@{}O}
\instbitrange{31}{29}&
\instbitrange{28}{26}&
\instbitrange{25}{24}&
\instbit{23}&
\instbitrange{22}{20}&
\instbit{19}&
\instbit{18}&
\instbitrange{17}{15}&
\instbitrange{14}{13}&
\instbitrange{12}{10}&
\instbit{9}&
\instbit{8}&
\instbit{7}&
\instbitrange{6}{0} \\
\hline
\multicolumn{1}{|c|}{{\scriptsize qx[2:0]}} &
\multicolumn{1}{c|}{{\scriptsize qy[2:0]}} &
\multicolumn{1}{c|}{01} &
\multicolumn{1}{c|}{{\scriptsize rs2[4]}} &
\multicolumn{1}{c|}{{\scriptsize rs2[2:0]}} &
\multicolumn{1}{c|}{1} &
\multicolumn{1}{c|}{{\scriptsize rs1[4]}} &
\multicolumn{1}{c|}{{\scriptsize rs1[2:0]}} &
\multicolumn{1}{c|}{01} &
\multicolumn{1}{c|}{{\scriptsize qu[2:0]}} &
\multicolumn{1}{c|}{1} &
\multicolumn{1}{c|}{{\scriptsize sat[0]}} &
\multicolumn{1}{c|}{0} &
\multicolumn{1}{c|}{0011111}\\
\hline
3& 3& 2& 1& 3& 1& 1& 3& 2& 3& 1& 1& 1& 7 \\
\end{tabular}
 \end{center}
 }
 \vspace{-0.1in}
 \caption[]{ESP.VMULAS.S16.XACC.ST.XP} \end{figure*}


\newpage
\subsubsection{ESP.VMULAS.S8.QACC}\label{instruction:ESPVMULASS8QACC}
\textbf{Syntax}\\
ESP.VMULAS.S8.QACC qx,qy,sat

\textbf{Description}\\
This instruction divides registers \lstinline{qx} and \lstinline{qy} into 16 data segments of 8 bits each. Then, the signed multiplication results of the 16 sets of segments are added to the corresponding 32-bit data segments in special registers \lstinline{QACC_H} and \lstinline{QACC_L} respectively. The calculated results are saturated to 32-bit signed numbers and then stored in the corresponding 32-bit data segments in \lstinline{QACC_H} and \lstinline{QACC_L}.

\textbf{Operation}\\
\begin{lstlisting}
QACC_L[31 :  0] = min(max(QACC_L[31 :  0] + qx[7  :  0] * qy[7  :  0], -2^31), 2^31-1)
QACC_L[63 : 32] = min(max(QACC_L[63 : 32] + qx[15 :  8] * qy[15 :  8], -2^31), 2^31-1)
QACC_L[95 : 64] = min(max(QACC_L[95 : 64] + qx[23 : 16] * qy[23 : 16], -2^31), 2^31-1)
QACC_L[127: 96] = min(max(QACC_L[127: 96] + qx[31 : 24] * qy[31 : 24], -2^31), 2^31-1)
QACC_L[159:128] = min(max(QACC_L[159:128] + qx[39 : 32] * qy[39 : 32], -2^31), 2^31-1)
QACC_L[191:160] = min(max(QACC_L[191:160] + qx[47 : 40] * qy[47 : 40], -2^31), 2^31-1)
QACC_L[223:192] = min(max(QACC_L[223:192] + qx[55 : 48] * qy[55 : 48], -2^31), 2^31-1)
QACC_L[255:224] = min(max(QACC_L[255:224] + qx[63 : 56] * qy[63 : 56], -2^31), 2^31-1)
QACC_H[31 :  0] = min(max(QACC_H[31 :  0] + qx[71 : 64] * qy[71 : 64], -2^31), 2^31-1)
QACC_H[63 : 32] = min(max(QACC_H[63 : 32] + qx[79 : 72] * qy[79 : 72], -2^31), 2^31-1)
QACC_H[95 : 64] = min(max(QACC_H[95 : 64] + qx[87 : 80] * qy[87 : 80], -2^31), 2^31-1)
QACC_H[127: 96] = min(max(QACC_H[127: 96] + qx[95 : 88] * qy[95 : 88], -2^31), 2^31-1)
QACC_H[159:128] = min(max(QACC_H[159:128] + qx[103: 96] * qy[103: 96], -2^31), 2^31-1)
QACC_H[191:160] = min(max(QACC_H[191:160] + qx[111:104] * qy[111:104], -2^31), 2^31-1)
QACC_H[223:192] = min(max(QACC_H[223:192] + qx[119:112] * qy[119:112], -2^31), 2^31-1)
QACC_H[255:224] = min(max(QACC_H[255:224] + qx[127:120] * qy[127:120], -2^31), 2^31-1)

\end{lstlisting}

\textbf{Schedule}\\
\begin{longtable}[c]{|l|l|l|}
    \hline
    \rowcolor{lightgray}
    \textbf{Operands} & \textbf{Use} & \textbf{Def} \\\hline
    \endfirsthead
    \hline
    \rowcolor{lightgray}
     \textbf{Operands} & \textbf{Use} & \textbf{Def} \\\hline
    \endhead
    \hline
    \endfoot

     qx & 1 & - \\\hline
     qy & 1 & - \\\hline
     QACC\_H & 2 & 2 \\\hline
     QACC\_L & 2 & 2

\end{longtable}

\textbf{Format}\\

\begin{figure*}[!htbp]
 {\setlength{\tabcolsep}{1pt}
 \begin{center}
\begin{tabular}{@{}F@{}F@{}T@{}I@{}E@{}O}
\instbitrange{31}{29}&
\instbitrange{28}{26}&
\instbitrange{25}{16}&
\instbit{15}&
\instbitrange{14}{7}&
\instbitrange{6}{0} \\
\hline
\multicolumn{1}{|c|}{{\scriptsize qx[2:0]}} &
\multicolumn{1}{c|}{{\scriptsize qy[2:0]}} &
\multicolumn{1}{c|}{1X101XX101} &
\multicolumn{1}{c|}{{\scriptsize sat[0]}} &
\multicolumn{1}{c|}{01XX0X0X} &
\multicolumn{1}{c|}{0011011}\\
\hline
3& 3& 10& 1& 8& 7 \\
\end{tabular}
 \end{center}
 }
 \vspace{-0.1in}
 \caption[]{ESP.VMULAS.S8.QACC} \end{figure*}


\newpage
\subsubsection{ESP.VMULAS.S8.QACC.LD.IP}\label{instruction:ESPVMULASS8QACCLDIP}
\textbf{Syntax}\\
ESP.VMULAS.S8.QACC.LD.IP qu,rs1,imm16,qx,qy,sat

\textbf{Description}\\
This instruction divides registers \lstinline{qx} and \lstinline{qy} into 16 data segments of 8 bits each. Then, the signed multiplication results of the 16 sets of segments are added to the corresponding 32-bit data segments in special registers \lstinline{QACC_H} and \lstinline{QACC_L} respectively. The calculated result is saturated to a 32-bit signed number and then stored in the corresponding 32-bit data segment in \lstinline{QACC_H} and \lstinline{QACC_L}.

During the operation, the 16-byte data is loaded from the memory to the register \lstinline{qu}. After the access is completed, the value in the register \lstinline{rs1} is incremented by the 4-bit sign-extended constant in the instruction code segment, which is left-shifted by 4.

imm starts at -128, ends at 112, with a step of 16.

\textbf{Operation}\\
\begin{lstlisting}
QACC_L[31 :  0] = min(max(QACC_L[31 :  0] + qx[7  :  0] * qy[7  :  0], -2^31), 2^31-1)
QACC_L[63 : 32] = min(max(QACC_L[63 : 32] + qx[15 :  8] * qy[15 :  8], -2^31), 2^31-1)
QACC_L[95 : 64] = min(max(QACC_L[95 : 64] + qx[23 : 16] * qy[23 : 16], -2^31), 2^31-1)
QACC_L[127: 96] = min(max(QACC_L[127: 96] + qx[31 : 24] * qy[31 : 24], -2^31), 2^31-1)
QACC_L[159:128] = min(max(QACC_L[159:128] + qx[39 : 32] * qy[39 : 32], -2^31), 2^31-1)
QACC_L[191:160] = min(max(QACC_L[191:160] + qx[47 : 40] * qy[47 : 40], -2^31), 2^31-1)
QACC_L[223:192] = min(max(QACC_L[223:192] + qx[55 : 48] * qy[55 : 48], -2^31), 2^31-1)
QACC_L[255:224] = min(max(QACC_L[255:224] + qx[63 : 56] * qy[63 : 56], -2^31), 2^31-1)
QACC_H[31 :  0] = min(max(QACC_H[31 :  0] + qx[71 : 64] * qy[71 : 64], -2^31), 2^31-1)
QACC_H[63 : 32] = min(max(QACC_H[63 : 32] + qx[79 : 72] * qy[79 : 72], -2^31), 2^31-1)
QACC_H[95 : 64] = min(max(QACC_H[95 : 64] + qx[87 : 80] * qy[87 : 80], -2^31), 2^31-1)
QACC_H[127: 96] = min(max(QACC_H[127: 96] + qx[95 : 88] * qy[95 : 88], -2^31), 2^31-1)
QACC_H[159:128] = min(max(QACC_H[159:128] + qx[103: 96] * qy[103: 96], -2^31), 2^31-1)
QACC_H[191:160] = min(max(QACC_H[191:160] + qx[111:104] * qy[111:104], -2^31), 2^31-1)
QACC_H[223:192] = min(max(QACC_H[223:192] + qx[119:112] * qy[119:112], -2^31), 2^31-1)
QACC_H[255:224] = min(max(QACC_H[255:224] + qx[127:120] * qy[127:120], -2^31), 2^31-1)

qu[127:0] = load128(rs1[31:0])
rs1[31:0] = rs1[31:0] + imm
\end{lstlisting}

\textbf{Schedule}\\
\begin{longtable}[c]{|l|l|l|}
    \hline
    \rowcolor{lightgray}
    \textbf{Operands} & \textbf{Use} & \textbf{Def} \\\hline
    \endfirsthead
    \hline
    \rowcolor{lightgray}
     \textbf{Operands} & \textbf{Use} & \textbf{Def} \\\hline
    \endhead
    \hline
    \endfoot

    qu & - & 2 \\\hline
    rs1 & 1 & 1 \\\hline
     qx & 1 & - \\\hline
     qy & 1 & - \\\hline
     QACC\_H & 2 & 2 \\\hline
     QACC\_L & 2 & 2

\end{longtable}

\textbf{Format}\\

\begin{figure*}[!htbp]
 {\setlength{\tabcolsep}{1pt}
 \begin{center}
\begin{tabular}{@{}F@{}F@{}Y@{}I@{}W@{}I@{}F@{}W@{}F@{}W@{}I@{}O}
\instbitrange{31}{29}&
\instbitrange{28}{26}&
\instbitrange{25}{22}&
\instbit{21}&
\instbitrange{20}{19}&
\instbit{18}&
\instbitrange{17}{15}&
\instbitrange{14}{13}&
\instbitrange{12}{10}&
\instbitrange{9}{8}&
\instbit{7}&
\instbitrange{6}{0} \\
\hline
\multicolumn{1}{|c|}{{\scriptsize qx[2:0]}} &
\multicolumn{1}{c|}{{\scriptsize qy[2:0]}} &
\multicolumn{1}{c|}{1001} &
\multicolumn{1}{c|}{{\scriptsize sat[0]}} &
\multicolumn{1}{c|}{{\scriptsize imm16[3:2]}} &
\multicolumn{1}{c|}{{\scriptsize rs1[4]}} &
\multicolumn{1}{c|}{{\scriptsize rs1[2:0]}} &
\multicolumn{1}{c|}{11} &
\multicolumn{1}{c|}{{\scriptsize qu[2:0]}} &
\multicolumn{1}{c|}{{\scriptsize imm16[1:0]}} &
\multicolumn{1}{c|}{0} &
\multicolumn{1}{c|}{0011111}\\
\hline
3& 3& 4& 1& 2& 1& 3& 2& 3& 2& 1& 7 \\
\end{tabular}
 \end{center}
 }
 \vspace{-0.1in}
 \caption[]{ESP.VMULAS.S8.QACC.LD.IP} \end{figure*}


\newpage
\subsubsection{ESP.VMULAS.S8.QACC.LD.XP}\label{instruction:ESPVMULASS8QACCLDXP}
\textbf{Syntax}\\
ESP.VMULAS.S8.QACC.LD.XP qu,rs1,rs2,qx,qy,sat

\textbf{Description}\\
This instruction divides registers \lstinline{qx} and \lstinline{qy} into 16 data segments of 8 bits each. Then, the signed multiplication results of the 16 sets of segments are added to the corresponding 32-bit data segments in special registers \lstinline{QACC_H} and \lstinline{QACC_L} respectively. The calculated results are saturated to 32-bit signed numbers and then stored in the corresponding 32-bit data segments in \lstinline{QACC_H} and \lstinline{QACC_L}.

During the operation, the 16-byte data is loaded from the memory into the register \lstinline{qu}. After the access is completed, the value in the register \lstinline{rs1} is incremented by the value of the register \lstinline{rs2}.

\textbf{Operation}\\
\begin{lstlisting}
QACC_L[31 :  0] = min(max(QACC_L[31 :  0] + qx[7  :  0] * qy[7  :  0], -2^31), 2^31-1)
QACC_L[63 : 32] = min(max(QACC_L[63 : 32] + qx[15 :  8] * qy[15 :  8], -2^31), 2^31-1)
QACC_L[95 : 64] = min(max(QACC_L[95 : 64] + qx[23 : 16] * qy[23 : 16], -2^31), 2^31-1)
QACC_L[127: 96] = min(max(QACC_L[127: 96] + qx[31 : 24] * qy[31 : 24], -2^31), 2^31-1)
QACC_L[159:128] = min(max(QACC_L[159:128] + qx[39 : 32] * qy[39 : 32], -2^31), 2^31-1)
QACC_L[191:160] = min(max(QACC_L[191:160] + qx[47 : 40] * qy[47 : 40], -2^31), 2^31-1)
QACC_L[223:192] = min(max(QACC_L[223:192] + qx[55 : 48] * qy[55 : 48], -2^31), 2^31-1)
QACC_L[255:224] = min(max(QACC_L[255:224] + qx[63 : 56] * qy[63 : 56], -2^31), 2^31-1)
QACC_H[31 :  0] = min(max(QACC_H[31 :  0] + qx[71 : 64] * qy[71 : 64], -2^31), 2^31-1)
QACC_H[63 : 32] = min(max(QACC_H[63 : 32] + qx[79 : 72] * qy[79 : 72], -2^31), 2^31-1)
QACC_H[95 : 64] = min(max(QACC_H[95 : 64] + qx[87 : 80] * qy[87 : 80], -2^31), 2^31-1)
QACC_H[127: 96] = min(max(QACC_H[127: 96] + qx[95 : 88] * qy[95 : 88], -2^31), 2^31-1)
QACC_H[159:128] = min(max(QACC_H[159:128] + qx[103: 96] * qy[103: 96], -2^31), 2^31-1)
QACC_H[191:160] = min(max(QACC_H[191:160] + qx[111:104] * qy[111:104], -2^31), 2^31-1)
QACC_H[223:192] = min(max(QACC_H[223:192] + qx[119:112] * qy[119:112], -2^31), 2^31-1)
QACC_H[255:224] = min(max(QACC_H[255:224] + qx[127:120] * qy[127:120], -2^31), 2^31-1)

qu[127:0] = load128(rs1[31:0])
rs1[31:0] = rs1[31:0] + rs2[31:0]
\end{lstlisting}

\textbf{Schedule}\\
\begin{longtable}[c]{|l|l|l|}
    \hline
    \rowcolor{lightgray}
    \textbf{Operands} & \textbf{Use} & \textbf{Def} \\\hline
    \endfirsthead
    \hline
    \rowcolor{lightgray}
     \textbf{Operands} & \textbf{Use} & \textbf{Def} \\\hline
    \endhead
    \hline
    \endfoot

    qu & - & 2 \\\hline
    rs1 & 1 & 1 \\\hline
    rs2 & 1 & - \\\hline
     qx & 1 & - \\\hline
     qy & 1 & - \\\hline
     QACC\_H & 2 & 2 \\\hline
     QACC\_L & 2 & 2

\end{longtable}

\textbf{Format}\\

\begin{figure*}[!htbp]
 {\setlength{\tabcolsep}{1pt}
 \begin{center}
\begin{tabular}{@{}F@{}F@{}W@{}I@{}F@{}I@{}I@{}F@{}W@{}F@{}I@{}I@{}I@{}O}
\instbitrange{31}{29}&
\instbitrange{28}{26}&
\instbitrange{25}{24}&
\instbit{23}&
\instbitrange{22}{20}&
\instbit{19}&
\instbit{18}&
\instbitrange{17}{15}&
\instbitrange{14}{13}&
\instbitrange{12}{10}&
\instbit{9}&
\instbit{8}&
\instbit{7}&
\instbitrange{6}{0} \\
\hline
\multicolumn{1}{|c|}{{\scriptsize qx[2:0]}} &
\multicolumn{1}{c|}{{\scriptsize qy[2:0]}} &
\multicolumn{1}{c|}{10} &
\multicolumn{1}{c|}{{\scriptsize rs2[4]}} &
\multicolumn{1}{c|}{{\scriptsize rs2[2:0]}} &
\multicolumn{1}{c|}{0} &
\multicolumn{1}{c|}{{\scriptsize rs1[4]}} &
\multicolumn{1}{c|}{{\scriptsize rs1[2:0]}} &
\multicolumn{1}{c|}{01} &
\multicolumn{1}{c|}{{\scriptsize qu[2:0]}} &
\multicolumn{1}{c|}{1} &
\multicolumn{1}{c|}{{\scriptsize sat[0]}} &
\multicolumn{1}{c|}{0} &
\multicolumn{1}{c|}{0011111}\\
\hline
3& 3& 2& 1& 3& 1& 1& 3& 2& 3& 1& 1& 1& 7 \\
\end{tabular}
 \end{center}
 }
 \vspace{-0.1in}
 \caption[]{ESP.VMULAS.S8.QACC.LD.XP} \end{figure*}


\newpage
\subsubsection{ESP.VMULAS.S8.QACC.ST.IP}\label{instruction:ESPVMULASS8QACCSTIP}
\textbf{Syntax}\\
ESP.VMULAS.S8.QACC.ST.IP qu,rs1,imm16,qx,qy,sat

\textbf{Description}\\
This instruction divides registers \lstinline{qx} and \lstinline{qy} into 16 data segments of 8 bits each. Then, the signed multiplication results of the 16 sets of segments are added to the corresponding 32-bit data segments in special registers \lstinline{QACC_H} and \lstinline{QACC_L} respectively. The calculated result is saturated to a 32-bit signed number and then stored in the corresponding 32-bit data segment in \lstinline{QACC_H} and \lstinline{QACC_L}.

During the operation, the value of the register \lstinline{qu} is stored in memory. After the access is completed, the value in the register \lstinline{rs1} is incremented by a 4-bit sign-extended constant in the instruction code segment, which is left-shifted by 4.

The value of imm starts at -128, ends at 112, and steps in increments of 16.

\textbf{Operation}\\
\begin{lstlisting}
QACC_L[31 :  0] = min(max(QACC_L[31 :  0] + qx[7  :  0] * qy[7  :  0], -2^31), 2^31-1)
QACC_L[63 : 32] = min(max(QACC_L[63 : 32] + qx[15 :  8] * qy[15 :  8], -2^31), 2^31-1)
QACC_L[95 : 64] = min(max(QACC_L[95 : 64] + qx[23 : 16] * qy[23 : 16], -2^31), 2^31-1)
QACC_L[127: 96] = min(max(QACC_L[127: 96] + qx[31 : 24] * qy[31 : 24], -2^31), 2^31-1)
QACC_L[159:128] = min(max(QACC_L[159:128] + qx[39 : 32] * qy[39 : 32], -2^31), 2^31-1)
QACC_L[191:160] = min(max(QACC_L[191:160] + qx[47 : 40] * qy[47 : 40], -2^31), 2^31-1)
QACC_L[223:192] = min(max(QACC_L[223:192] + qx[55 : 48] * qy[55 : 48], -2^31), 2^31-1)
QACC_L[255:224] = min(max(QACC_L[255:224] + qx[63 : 56] * qy[63 : 56], -2^31), 2^31-1)
QACC_H[31 :  0] = min(max(QACC_H[31 :  0] + qx[71 : 64] * qy[71 : 64], -2^31), 2^31-1)
QACC_H[63 : 32] = min(max(QACC_H[63 : 32] + qx[79 : 72] * qy[79 : 72], -2^31), 2^31-1)
QACC_H[95 : 64] = min(max(QACC_H[95 : 64] + qx[87 : 80] * qy[87 : 80], -2^31), 2^31-1)
QACC_H[127: 96] = min(max(QACC_H[127: 96] + qx[95 : 88] * qy[95 : 88], -2^31), 2^31-1)
QACC_H[159:128] = min(max(QACC_H[159:128] + qx[103: 96] * qy[103: 96], -2^31), 2^31-1)
QACC_H[191:160] = min(max(QACC_H[191:160] + qx[111:104] * qy[111:104], -2^31), 2^31-1)
QACC_H[223:192] = min(max(QACC_H[223:192] + qx[119:112] * qy[119:112], -2^31), 2^31-1)
QACC_H[255:224] = min(max(QACC_H[255:224] + qx[127:120] * qy[127:120], -2^31), 2^31-1)

qu[127:0] => store128(rs1[31:0])
rs1[31:0] = rs1[31:0] + imm
\end{lstlisting}

\textbf{Schedule}\\
\begin{longtable}[c]{|l|l|l|}
    \hline
    \rowcolor{lightgray}
    \textbf{Operands} & \textbf{Use} & \textbf{Def} \\\hline
    \endfirsthead
    \hline
    \rowcolor{lightgray}
     \textbf{Operands} & \textbf{Use} & \textbf{Def} \\\hline
    \endhead
    \hline
    \endfoot

    qu & 2 & - \\\hline
    rs2 & 1 & - \\\hline
     qx & 1 & - \\\hline
     qy & 1 & - \\\hline
     QACC\_H & 2 & 2 \\\hline
     QACC\_L & 2 & 2

\end{longtable}

\textbf{Format}\\

\begin{figure*}[!htbp]
 {\setlength{\tabcolsep}{1pt}
 \begin{center}
\begin{tabular}{@{}F@{}F@{}Y@{}I@{}W@{}I@{}F@{}W@{}F@{}W@{}I@{}O}
\instbitrange{31}{29}&
\instbitrange{28}{26}&
\instbitrange{25}{22}&
\instbit{21}&
\instbitrange{20}{19}&
\instbit{18}&
\instbitrange{17}{15}&
\instbitrange{14}{13}&
\instbitrange{12}{10}&
\instbitrange{9}{8}&
\instbit{7}&
\instbitrange{6}{0} \\
\hline
\multicolumn{1}{|c|}{{\scriptsize qx[2:0]}} &
\multicolumn{1}{c|}{{\scriptsize qy[2:0]}} &
\multicolumn{1}{c|}{1101} &
\multicolumn{1}{c|}{{\scriptsize sat[0]}} &
\multicolumn{1}{c|}{{\scriptsize imm16[3:2]}} &
\multicolumn{1}{c|}{{\scriptsize rs1[4]}} &
\multicolumn{1}{c|}{{\scriptsize rs1[2:0]}} &
\multicolumn{1}{c|}{11} &
\multicolumn{1}{c|}{{\scriptsize qu[2:0]}} &
\multicolumn{1}{c|}{{\scriptsize imm16[1:0]}} &
\multicolumn{1}{c|}{0} &
\multicolumn{1}{c|}{0011111}\\
\hline
3& 3& 4& 1& 2& 1& 3& 2& 3& 2& 1& 7 \\
\end{tabular}
 \end{center}
 }
 \vspace{-0.1in}
 \caption[]{ESP.VMULAS.S8.QACC.ST.IP} \end{figure*}


\newpage
\subsubsection{ESP.VMULAS.S8.QACC.ST.XP}\label{instruction:ESPVMULASS8QACCSTXP}
\textbf{Syntax}\\
ESP.VMULAS.S8.QACC.ST.XP qu,rs1,rs2,qx,qy,sat

\textbf{Description}\\
This instruction divides registers \lstinline{qx} and \lstinline{qy} into 16 data segments of 8 bits each. Then, the signed multiplication results of the 16 sets of segments are added to the corresponding 32-bit data segments in special registers \lstinline{QACC_H} and \lstinline{QACC_L} respectively. The calculated results are saturated to 32-bit signed numbers and then stored in the corresponding 32-bit data segments in \lstinline{QACC_H} and \lstinline{QACC_L}.

During the operation, the value of the register \lstinline{qu} is stored in memory. After the access is completed, the value in the register \lstinline{rs1} is incremented by the value of the register \lstinline{rs2}.

\textbf{Operation}\\
\begin{lstlisting}
QACC_L[31 :  0] = min(max(QACC_L[31 :  0] + qx[7  :  0] * qy[7  :  0], -2^31), 2^31-1)
QACC_L[63 : 32] = min(max(QACC_L[63 : 32] + qx[15 :  8] * qy[15 :  8], -2^31), 2^31-1)
QACC_L[95 : 64] = min(max(QACC_L[95 : 64] + qx[23 : 16] * qy[23 : 16], -2^31), 2^31-1)
QACC_L[127: 96] = min(max(QACC_L[127: 96] + qx[31 : 24] * qy[31 : 24], -2^31), 2^31-1)
QACC_L[159:128] = min(max(QACC_L[159:128] + qx[39 : 32] * qy[39 : 32], -2^31), 2^31-1)
QACC_L[191:160] = min(max(QACC_L[191:160] + qx[47 : 40] * qy[47 : 40], -2^31), 2^31-1)
QACC_L[223:192] = min(max(QACC_L[223:192] + qx[55 : 48] * qy[55 : 48], -2^31), 2^31-1)
QACC_L[255:224] = min(max(QACC_L[255:224] + qx[63 : 56] * qy[63 : 56], -2^31), 2^31-1)
QACC_H[31 :  0] = min(max(QACC_H[31 :  0] + qx[71 : 64] * qy[71 : 64], -2^31), 2^31-1)
QACC_H[63 : 32] = min(max(QACC_H[63 : 32] + qx[79 : 72] * qy[79 : 72], -2^31), 2^31-1)
QACC_H[95 : 64] = min(max(QACC_H[95 : 64] + qx[87 : 80] * qy[87 : 80], -2^31), 2^31-1)
QACC_H[127: 96] = min(max(QACC_H[127: 96] + qx[95 : 88] * qy[95 : 88], -2^31), 2^31-1)
QACC_H[159:128] = min(max(QACC_H[159:128] + qx[103: 96] * qy[103: 96], -2^31), 2^31-1)
QACC_H[191:160] = min(max(QACC_H[191:160] + qx[111:104] * qy[111:104], -2^31), 2^31-1)
QACC_H[223:192] = min(max(QACC_H[223:192] + qx[119:112] * qy[119:112], -2^31), 2^31-1)
QACC_H[255:224] = min(max(QACC_H[255:224] + qx[127:120] * qy[127:120], -2^31), 2^31-1)

qu[127:0] => store128(rs1[31:0])
rs1[31:0] = rs1[31:0] + rs2[31:0]
\end{lstlisting}

\textbf{Schedule}\\
\begin{longtable}[c]{|l|l|l|}
    \hline
    \rowcolor{lightgray}
    \textbf{Operands} & \textbf{Use} & \textbf{Def} \\\hline
    \endfirsthead
    \hline
    \rowcolor{lightgray}
     \textbf{Operands} & \textbf{Use} & \textbf{Def} \\\hline
    \endhead
    \hline
    \endfoot

    qu & 2 & - \\\hline
    rs1 & 1 & 1 \\\hline
    rs2 & 1 & - \\\hline
     qx & 1 & - \\\hline
     qy & 1 & - \\\hline
     QACC\_H & 2 & 2 \\\hline
     QACC\_L & 2 & 2

\end{longtable}

\textbf{Format}\\

\begin{figure*}[!htbp]
 {\setlength{\tabcolsep}{1pt}
 \begin{center}
\begin{tabular}{@{}F@{}F@{}W@{}I@{}F@{}I@{}I@{}F@{}W@{}F@{}I@{}I@{}I@{}O}
\instbitrange{31}{29}&
\instbitrange{28}{26}&
\instbitrange{25}{24}&
\instbit{23}&
\instbitrange{22}{20}&
\instbit{19}&
\instbit{18}&
\instbitrange{17}{15}&
\instbitrange{14}{13}&
\instbitrange{12}{10}&
\instbit{9}&
\instbit{8}&
\instbit{7}&
\instbitrange{6}{0} \\
\hline
\multicolumn{1}{|c|}{{\scriptsize qx[2:0]}} &
\multicolumn{1}{c|}{{\scriptsize qy[2:0]}} &
\multicolumn{1}{c|}{11} &
\multicolumn{1}{c|}{{\scriptsize rs2[4]}} &
\multicolumn{1}{c|}{{\scriptsize rs2[2:0]}} &
\multicolumn{1}{c|}{0} &
\multicolumn{1}{c|}{{\scriptsize rs1[4]}} &
\multicolumn{1}{c|}{{\scriptsize rs1[2:0]}} &
\multicolumn{1}{c|}{01} &
\multicolumn{1}{c|}{{\scriptsize qu[2:0]}} &
\multicolumn{1}{c|}{1} &
\multicolumn{1}{c|}{{\scriptsize sat[0]}} &
\multicolumn{1}{c|}{0} &
\multicolumn{1}{c|}{0011111}\\
\hline
3& 3& 2& 1& 3& 1& 1& 3& 2& 3& 1& 1& 1& 7 \\
\end{tabular}
 \end{center}
 }
 \vspace{-0.1in}
 \caption[]{ESP.VMULAS.S8.QACC.ST.XP} \end{figure*}


\newpage
\subsubsection{ESP.VMULAS.S8.XACC}\label{instruction:ESPVMULASS8XACC}
\textbf{Syntax}\\
ESP.VMULAS.S8.XACC qx,qy,sat

\textbf{Description}\\
This instruction divides registers \lstinline{qx} and \lstinline{qy} into 16 data segments of 8 bits each. Then, it performs a signed multiply-accumulate operation on each of the 16 sets of segments. The accumulated result is added to the value in the special register \lstinline{XACC}. The resulting sum is then saturated and stored in \lstinline{XACC}.\\
\textbf{Operation}
\begin{lstlisting}
  add0[15:0] = qx[  7:  0] * qy[  7:  0]
  add1[15:0] = qx[ 15:  8] * qy[ 15:  8]
  ...
  add15[15:0] = qx[127:120] * qy[127:120]
  sum[40:0] = XACC[39:0] + add0[15:0] + add1[15:0] + ... + add15[15:0]

  XACC[39:0] = min(max(sum[40:0], -2^{39}), 2^{39}-1)
\end{lstlisting}

\textbf{Schedule}\\
\begin{longtable}[c]{|l|l|l|}
    \hline
    \rowcolor{lightgray}
    \textbf{Operands} & \textbf{Use} & \textbf{Def} \\\hline
    \endfirsthead
    \hline
    \rowcolor{lightgray}
     \textbf{Operands} & \textbf{Use} & \textbf{Def} \\\hline
    \endhead
    \hline
    \endfoot


     qx & 1 & - \\\hline
     qy & 1 & - \\\hline
     XACC & 3 & 3

\end{longtable}

\textbf{Format}\\

\begin{figure*}[!htbp]
 {\setlength{\tabcolsep}{1pt}
 \begin{center}
\begin{tabular}{@{}F@{}F@{}T@{}I@{}E@{}O}
\instbitrange{31}{29}&
\instbitrange{28}{26}&
\instbitrange{25}{16}&
\instbit{15}&
\instbitrange{14}{7}&
\instbitrange{6}{0} \\
\hline
\multicolumn{1}{|c|}{{\scriptsize qx[2:0]}} &
\multicolumn{1}{c|}{{\scriptsize qy[2:0]}} &
\multicolumn{1}{c|}{1X101XX001} &
\multicolumn{1}{c|}{{\scriptsize sat[0]}} &
\multicolumn{1}{c|}{01XX0X0X} &
\multicolumn{1}{c|}{0011011}\\
\hline
3& 3& 10& 1& 8& 7 \\
\end{tabular}
 \end{center}
 }
 \vspace{-0.1in}
 \caption[]{ESP.VMULAS.S8.XACC} \end{figure*}


\newpage
\subsubsection{ESP.VMULAS.S8.XACC.LD.IP}\label{instruction:ESPVMULASS8XACCLDIP}
\textbf{Syntax}\\
ESP.VMULAS.S8.XACC.LD.IP qu,rs1,imm16,qx,qy,sat

\textbf{Description}\\
This instruction divides registers \lstinline{qx} and \lstinline{qy} into 16 data segments of 8 bits each. Then, it performs a signed multiply-accumulate operation on each of the 16 segment sets. The accumulated result is added to the value in the special register \lstinline{XACC}. The resulting sum is then saturated and stored in \lstinline{XACC}. \\ During the operation, the 16-byte data is loaded from the memory into the register \lstinline{qu}. After the access is completed, the value in the register \lstinline{rs1} is incremented by a 4-bit sign-extended constant from the instruction code segment, which is left-shifted by 4.\\

The imm starts at -128, ends at 112, and steps in increments of 16.

\textbf{Operation}
\begin{lstlisting}
  add0[15:0] = qx[  7:  0] * qy[  7:  0]
  add1[15:0] = qx[ 15:  8] * qy[ 15:  8]
  ...
  add15[15:0] = qx[127:120] * qy[127:120]
  sum[40:0] = XACC[39:0] + add0[15:0] + add1[15:0] + ... + add15[15:0]

  XACC[39:0] = min(max(sum[40:0], -2^{39}), 2^{39}-1)
  qu[127:0] = load128(rs1[31:0])
  rs1[31:0] = rs1[31:0] + imm
\end{lstlisting}

\textbf{Schedule}\\
\begin{longtable}[c]{|l|l|l|}
    \hline
    \rowcolor{lightgray}
    \textbf{Operands} & \textbf{Use} & \textbf{Def} \\\hline
    \endfirsthead
    \hline
    \rowcolor{lightgray}
     \textbf{Operands} & \textbf{Use} & \textbf{Def} \\\hline
    \endhead
    \hline
    \endfoot

    qu & - & 2 \\\hline
    rs1 & 1 & 1 \\\hline
     qx & 1 & - \\\hline
     qy & 1 & - \\\hline
     XACC & 3 & 3


\end{longtable}

\textbf{Format}\\

\begin{figure*}[!htbp]
 {\setlength{\tabcolsep}{1pt}
 \begin{center}
\begin{tabular}{@{}F@{}F@{}Y@{}I@{}W@{}I@{}F@{}W@{}F@{}W@{}I@{}O}
\instbitrange{31}{29}&
\instbitrange{28}{26}&
\instbitrange{25}{22}&
\instbit{21}&
\instbitrange{20}{19}&
\instbit{18}&
\instbitrange{17}{15}&
\instbitrange{14}{13}&
\instbitrange{12}{10}&
\instbitrange{9}{8}&
\instbit{7}&
\instbitrange{6}{0} \\
\hline
\multicolumn{1}{|c|}{{\scriptsize qx[2:0]}} &
\multicolumn{1}{c|}{{\scriptsize qy[2:0]}} &
\multicolumn{1}{c|}{0001} &
\multicolumn{1}{c|}{{\scriptsize sat[0]}} &
\multicolumn{1}{c|}{{\scriptsize imm16[3:2]}} &
\multicolumn{1}{c|}{{\scriptsize rs1[4]}} &
\multicolumn{1}{c|}{{\scriptsize rs1[2:0]}} &
\multicolumn{1}{c|}{11} &
\multicolumn{1}{c|}{{\scriptsize qu[2:0]}} &
\multicolumn{1}{c|}{{\scriptsize imm16[1:0]}} &
\multicolumn{1}{c|}{0} &
\multicolumn{1}{c|}{0011111}\\
\hline
3& 3& 4& 1& 2& 1& 3& 2& 3& 2& 1& 7 \\
\end{tabular}
 \end{center}
 }
 \vspace{-0.1in}
 \caption[]{ESP.VMULAS.S8.XACC.LD.IP} \end{figure*}


\newpage
\subsubsection{ESP.VMULAS.S8.XACC.LD.XP}\label{instruction:ESPVMULASS8XACCLDXP}
\textbf{Syntax}\\
ESP.VMULAS.S8.XACC.LD.XP qu,rs1,rs2,qx,qy,sat

\textbf{Description}\\

This instruction divides registers \lstinline{qx} and \lstinline{qy} into 16 data segments of 8 bits each. Then, it performs a signed multiply-accumulate operation on each of the 16 segments. The accumulated result is added to the value in the special register \lstinline{XACC}. The resulting sum is then saturated and stored in \lstinline{XACC}. \\ During the operation, the 16-byte data is loaded from the memory into the register \lstinline{qu}. After the access is completed, the value in the register \lstinline{rs1} is incremented by the value of the register \lstinline{rs2}.\\
\textbf{Operation}
\begin{lstlisting}
  add0[15:0] = qx[  7:  0] * qy[  7:  0]
  add1[15:0] = qx[ 15:  8] * qy[ 15:  8]
  ...
  add15[15:0] = qx[127:120] * qy[127:120]
  sum[40:0] = XACC[39:0] + add0[15:0] + add1[15:0] + ... + add15[15:0]

  XACC[39:0] = min(max(sum[40:0], -2^{39}), 2^{39}-1)
  qu[127:0] = load128(rs1[31:0])
  rs1[31:0] = rs1[31:0] + rs2[31:0]
\end{lstlisting}

\textbf{Schedule}\\
\begin{longtable}[c]{|l|l|l|}
    \hline
    \rowcolor{lightgray}
    \textbf{Operands} & \textbf{Use} & \textbf{Def} \\\hline
    \endfirsthead
    \hline
    \rowcolor{lightgray}
     \textbf{Operands} & \textbf{Use} & \textbf{Def} \\\hline
    \endhead
    \hline
    \endfoot

    qu & - & 2 \\\hline
    rs1 & 1 & 1 \\\hline
    rs2 & 1 & - \\\hline
     qx & 1 & - \\\hline
     qy & 1 & - \\\hline
     QACC\_H & 2 & 2 \\\hline
     QACC\_L & 2 & 2

\end{longtable}

\textbf{Format}\\

\begin{figure*}[!htbp]
 {\setlength{\tabcolsep}{1pt}
 \begin{center}
\begin{tabular}{@{}F@{}F@{}W@{}I@{}F@{}I@{}I@{}F@{}W@{}F@{}I@{}I@{}I@{}O}
\instbitrange{31}{29}&
\instbitrange{28}{26}&
\instbitrange{25}{24}&
\instbit{23}&
\instbitrange{22}{20}&
\instbit{19}&
\instbit{18}&
\instbitrange{17}{15}&
\instbitrange{14}{13}&
\instbitrange{12}{10}&
\instbit{9}&
\instbit{8}&
\instbit{7}&
\instbitrange{6}{0} \\
\hline
\multicolumn{1}{|c|}{{\scriptsize qx[2:0]}} &
\multicolumn{1}{c|}{{\scriptsize qy[2:0]}} &
\multicolumn{1}{c|}{00} &
\multicolumn{1}{c|}{{\scriptsize rs2[4]}} &
\multicolumn{1}{c|}{{\scriptsize rs2[2:0]}} &
\multicolumn{1}{c|}{0} &
\multicolumn{1}{c|}{{\scriptsize rs1[4]}} &
\multicolumn{1}{c|}{{\scriptsize rs1[2:0]}} &
\multicolumn{1}{c|}{01} &
\multicolumn{1}{c|}{{\scriptsize qu[2:0]}} &
\multicolumn{1}{c|}{1} &
\multicolumn{1}{c|}{{\scriptsize sat[0]}} &
\multicolumn{1}{c|}{0} &
\multicolumn{1}{c|}{0011111}\\
\hline
3& 3& 2& 1& 3& 1& 1& 3& 2& 3& 1& 1& 1& 7 \\
\end{tabular}
 \end{center}
 }
 \vspace{-0.1in}
 \caption[]{ESP.VMULAS.S8.XACC.LD.XP} \end{figure*}


\newpage
\subsubsection{ESP.VMULAS.S8.XACC.ST.IP}\label{instruction:ESPVMULASS8XACCSTIP}
\textbf{Syntax}\\
ESP.VMULAS.S8.XACC.ST.IP qu,rs1,imm16,qx,qy,sat

\textbf{Description}\\
This instruction divides registers \lstinline{qx} and \lstinline{qy} into 16 data segments of 8 bits each. Then, it performs a signed multiply-accumulate operation on each of the 16 segments. The accumulated result is added to the value in the special register \lstinline{XACC}. The resulting sum is then saturated and stored in \lstinline{XACC}. \\ During the operation, the register \lstinline{qu} is stored to memory. After the access is completed, the value in the register \lstinline{rs1} is incremented by a 4-bit sign-extended constant from the instruction code segment, which is left-shifted by 4.\\

The imm starts at -128, ends at 112, and steps in increments of 16.

\textbf{Operation}
\begin{lstlisting}
  add0[15:0] = qx[  7:  0] * qy[  7:  0]
  add1[15:0] = qx[ 15:  8] * qy[ 15:  8]
  ...
  add15[15:0] = qx[127:120] * qy[127:120]
  sum[40:0] = XACC[39:0] + add0[15:0] + add1[15:0] + ... + add15[15:0]

  XACC[39:0] = min(max(sum[40:0], -2^{39}), 2^{39}-1)
  qu[127:0] => store128(rs1[31:0])
  rs1[31:0] = rs1[31:0] + imm
\end{lstlisting}

\textbf{Schedule}\\
\begin{longtable}[c]{|l|l|l|}
    \hline
    \rowcolor{lightgray}
    \textbf{Operands} & \textbf{Use} & \textbf{Def} \\\hline
    \endfirsthead
    \hline
    \rowcolor{lightgray}
     \textbf{Operands} & \textbf{Use} & \textbf{Def} \\\hline
    \endhead
    \hline
    \endfoot

    qu & 2 & 1 \\\hline
    rs1 & 1 & 1 \\\hline
     qx & 1 & - \\\hline
     qy & 1 & - \\\hline
     XACC & 3 & 3

\end{longtable}

\textbf{Format}\\

\begin{figure*}[!htbp]
 {\setlength{\tabcolsep}{1pt}
 \begin{center}
\begin{tabular}{@{}F@{}F@{}Y@{}I@{}W@{}I@{}F@{}W@{}F@{}W@{}I@{}O}
\instbitrange{31}{29}&
\instbitrange{28}{26}&
\instbitrange{25}{22}&
\instbit{21}&
\instbitrange{20}{19}&
\instbit{18}&
\instbitrange{17}{15}&
\instbitrange{14}{13}&
\instbitrange{12}{10}&
\instbitrange{9}{8}&
\instbit{7}&
\instbitrange{6}{0} \\
\hline
\multicolumn{1}{|c|}{{\scriptsize qx[2:0]}} &
\multicolumn{1}{c|}{{\scriptsize qy[2:0]}} &
\multicolumn{1}{c|}{0101} &
\multicolumn{1}{c|}{{\scriptsize sat[0]}} &
\multicolumn{1}{c|}{{\scriptsize imm16[3:2]}} &
\multicolumn{1}{c|}{{\scriptsize rs1[4]}} &
\multicolumn{1}{c|}{{\scriptsize rs1[2:0]}} &
\multicolumn{1}{c|}{11} &
\multicolumn{1}{c|}{{\scriptsize qu[2:0]}} &
\multicolumn{1}{c|}{{\scriptsize imm16[1:0]}} &
\multicolumn{1}{c|}{0} &
\multicolumn{1}{c|}{0011111}\\
\hline
3& 3& 4& 1& 2& 1& 3& 2& 3& 2& 1& 7 \\
\end{tabular}
 \end{center}
 }
 \vspace{-0.1in}
 \caption[]{ESP.VMULAS.S8.XACC.ST.IP} \end{figure*}


\newpage
\subsubsection{ESP.VMULAS.S8.XACC.ST.XP}\label{instruction:ESPVMULASS8XACCSTXP}
\textbf{Syntax}\\
ESP.VMULAS.S8.XACC.ST.XP qu,rs1,rs2,qx,qy,sat

\textbf{Description}\\
This instruction divides registers \lstinline{qx} and \lstinline{qy} into 16 data segments of 8 bits each. Then, it performs a signed multiply-accumulate operation on each of the 16 sets of segments. The accumulated result is added to the value in the special register \lstinline{XACC}. The resulting sum is then saturated and stored in \lstinline{XACC}. \\ During the operation, the value in the register \lstinline{qu} is stored to memory. After the access is completed, the value in the register \lstinline{rs1} is incremented by the value of the register \lstinline{rs2}.\\
\textbf{Operation}
\begin{lstlisting}
  add0[15:0] = qx[  7:  0] * qy[  7:  0]
  add1[15:0] = qx[ 15:  8] * qy[ 15:  8]
  ...
  add15[15:0] = qx[127:120] * qy[127:120]
  sum[40:0] = XACC[39:0] + add0[15:0] + add1[15:0] + ... + add15[15:0]

  XACC[39:0] = min(max(sum[40:0], -2^{39}), 2^{39}-1)
  qu[127:0] => store128(rs1[31:0])
  rs1[31:0] = rs1[31:0] + rs2[31:0]
\end{lstlisting}

\textbf{Schedule}\\
\begin{longtable}[c]{|l|l|l|}
    \hline
    \rowcolor{lightgray}
    \textbf{Operands} & \textbf{Use} & \textbf{Def} \\\hline
    \endfirsthead
    \hline
    \rowcolor{lightgray}
     \textbf{Operands} & \textbf{Use} & \textbf{Def} \\\hline
    \endhead
    \hline
    \endfoot

    qu & - & 2 \\\hline
    rs1 & 1 & 1 \\\hline
    rs2 & 1 & - \\\hline
     qx & 1 & - \\\hline
     qy & 1 & - \\\hline
     XACC & 3 & 3

\end{longtable}

\textbf{Format}\\

\begin{figure*}[!htbp]
 {\setlength{\tabcolsep}{1pt}
 \begin{center}
\begin{tabular}{@{}F@{}F@{}W@{}I@{}F@{}I@{}I@{}F@{}W@{}F@{}I@{}I@{}I@{}O}
\instbitrange{31}{29}&
\instbitrange{28}{26}&
\instbitrange{25}{24}&
\instbit{23}&
\instbitrange{22}{20}&
\instbit{19}&
\instbit{18}&
\instbitrange{17}{15}&
\instbitrange{14}{13}&
\instbitrange{12}{10}&
\instbit{9}&
\instbit{8}&
\instbit{7}&
\instbitrange{6}{0} \\
\hline
\multicolumn{1}{|c|}{{\scriptsize qx[2:0]}} &
\multicolumn{1}{c|}{{\scriptsize qy[2:0]}} &
\multicolumn{1}{c|}{01} &
\multicolumn{1}{c|}{{\scriptsize rs2[4]}} &
\multicolumn{1}{c|}{{\scriptsize rs2[2:0]}} &
\multicolumn{1}{c|}{0} &
\multicolumn{1}{c|}{{\scriptsize rs1[4]}} &
\multicolumn{1}{c|}{{\scriptsize rs1[2:0]}} &
\multicolumn{1}{c|}{01} &
\multicolumn{1}{c|}{{\scriptsize qu[2:0]}} &
\multicolumn{1}{c|}{1} &
\multicolumn{1}{c|}{{\scriptsize sat[0]}} &
\multicolumn{1}{c|}{0} &
\multicolumn{1}{c|}{0011111}\\
\hline
3& 3& 2& 1& 3& 1& 1& 3& 2& 3& 1& 1& 1& 7 \\
\end{tabular}
 \end{center}
 }
 \vspace{-0.1in}
 \caption[]{ESP.VMULAS.S8.XACC.ST.XP} \end{figure*}


\newpage
\subsubsection{ESP.VMULAS.U16.QACC}\label{instruction:ESPVMULASU16QACC}
\textbf{Syntax}\\
ESP.VMULAS.U16.QACC qx,qy,sat

\textbf{Description}\\
This instruction divides registers \lstinline{qx} and \lstinline{qy} into 8 data segments of 16 bits each. Then, the unsigned multiplication results of the 8 sets of segments are added to the corresponding 64-bit data segments in special registers \lstinline{QACC_H} and \lstinline{QACC_L} respectively. The calculated results are saturated to 64-bit signed numbers and then stored in the corresponding 64-bit data segments in \lstinline{QACC_H} and \lstinline{QACC_L}.

\textbf{Operation}\\
\begin{lstlisting}
QACC_L[63 :  0] = min(QACC_L[63 :  0] + qx[15 :  0] * qy[15 :  0], 2^63-1)
QACC_L[127: 64] = min(QACC_L[127: 64] + qx[31 : 16] * qy[31 : 16], 2^63-1)
QACC_L[191:128] = min(QACC_L[191:128] + qx[47 : 32] * qy[47 : 32], 2^63-1)
QACC_L[255:192] = min(QACC_L[255:192] + qx[63 : 48] * qy[63 : 48], 2^63-1)
QACC_H[63 :  0] = min(QACC_H[63 :  0] + qx[79 : 64] * qy[79 : 64], 2^63-1)
QACC_H[127: 64] = min(QACC_H[127: 64] + qx[95 : 80] * qy[95 : 80], 2^63-1)
QACC_H[191:128] = min(QACC_H[191:128] + qx[111: 96] * qy[111: 96], 2^63-1)
QACC_H[255:192] = min(QACC_H[255:192] + qx[127:112] * qy[127:112], 2^63-1)
\end{lstlisting}

\textbf{Schedule}\\
\begin{longtable}[c]{|l|l|l|}
    \hline
    \rowcolor{lightgray}
    \textbf{Operands} & \textbf{Use} & \textbf{Def} \\\hline
    \endfirsthead
    \hline
    \rowcolor{lightgray}
     \textbf{Operands} & \textbf{Use} & \textbf{Def} \\\hline
    \endhead
    \hline
    \endfoot

     qx & 1 & - \\\hline
     qy & 1 & - \\\hline
     QACC\_H & 2 & 2 \\\hline
     QACC\_L & 2 & 2

\end{longtable}

\textbf{Format}\\

\begin{figure*}[!htbp]
 {\setlength{\tabcolsep}{1pt}
 \begin{center}
\begin{tabular}{@{}F@{}F@{}T@{}I@{}E@{}O}
\instbitrange{31}{29}&
\instbitrange{28}{26}&
\instbitrange{25}{16}&
\instbit{15}&
\instbitrange{14}{7}&
\instbitrange{6}{0} \\
\hline
\multicolumn{1}{|c|}{{\scriptsize qx[2:0]}} &
\multicolumn{1}{c|}{{\scriptsize qy[2:0]}} &
\multicolumn{1}{c|}{1X101XX110} &
\multicolumn{1}{c|}{{\scriptsize sat[0]}} &
\multicolumn{1}{c|}{01XX0X0X} &
\multicolumn{1}{c|}{0011011}\\
\hline
3& 3& 10& 1& 8& 7 \\
\end{tabular}
 \end{center}
 }
 \vspace{-0.1in}
 \caption[]{ESP.VMULAS.U16.QACC} \end{figure*}


\newpage
\subsubsection{ESP.VMULAS.U16.QACC.LD.IP}\label{instruction:ESPVMULASU16QACCLDIP}
\textbf{Syntax}\\
ESP.VMULAS.U16.QACC.LD.IP qu,rs1,imm16,qx,qy,sat

\textbf{Description}\\
This instruction divides registers \lstinline{qx} and \lstinline{qy} into 8 data segments of 16 bits each. Then, the unsigned multiplication results of the 8 sets of segments are added to the corresponding 64-bit data segments in special registers \lstinline{QACC_H} and \lstinline{QACC_L} respectively. The calculated result is saturated to a 64-bit signed number and then stored in the corresponding 64-bit data segment in \lstinline{QACC_H} and \lstinline{QACC_L}.

During the operation, the 16-byte data is loaded from the memory into the register \lstinline{qu}. After the access is completed, the value in the register \lstinline{rs1} is incremented by the 4-bit sign-extended constant in the instruction code segment, which is left-shifted by 4.

The range of imm starts at -128, ends at 112, and steps by 16.

\textbf{Operation}\\
\begin{lstlisting}
QACC_L[63 :  0] = min(QACC_L[63 :  0] + qx[15 :  0] * qy[15 :  0], 2^63-1)
QACC_L[127: 64] = min(QACC_L[127: 64] + qx[31 : 16] * qy[31 : 16], 2^63-1)
QACC_L[191:128] = min(QACC_L[191:128] + qx[47 : 32] * qy[47 : 32], 2^63-1)
QACC_L[255:192] = min(QACC_L[255:192] + qx[63 : 48] * qy[63 : 48], 2^63-1)
QACC_H[63 :  0] = min(QACC_H[63 :  0] + qx[79 : 64] * qy[79 : 64], 2^63-1)
QACC_H[127: 64] = min(QACC_H[127: 64] + qx[95 : 80] * qy[95 : 80], 2^63-1)
QACC_H[191:128] = min(QACC_H[191:128] + qx[111: 96] * qy[111: 96], 2^63-1)
QACC_H[255:192] = min(QACC_H[255:192] + qx[127:112] * qy[127:112], 2^63-1)

qu[127:0] = load128(rs1[31:0])
rs1[31:0] = rs1[31:0] + imm
\end{lstlisting}

\textbf{Schedule}\\
\begin{longtable}[c]{|l|l|l|}
    \hline
    \rowcolor{lightgray}
    \textbf{Operands} & \textbf{Use} & \textbf{Def} \\\hline
    \endfirsthead
    \hline
    \rowcolor{lightgray}
     \textbf{Operands} & \textbf{Use} & \textbf{Def} \\\hline
    \endhead
    \hline
    \endfoot

    qu & - & 2 \\\hline
    rs1 & 1 & 1 \\\hline
     qx & 1 & - \\\hline
     qy & 1 & - \\\hline
     QACC\_H & 2 & 2 \\\hline
     QACC\_L & 2 & 2

\end{longtable}

\textbf{Format}\\

\begin{figure*}[!htbp]
 {\setlength{\tabcolsep}{1pt}
 \begin{center}
\begin{tabular}{@{}F@{}F@{}Y@{}I@{}W@{}I@{}F@{}W@{}F@{}W@{}I@{}O}
\instbitrange{31}{29}&
\instbitrange{28}{26}&
\instbitrange{25}{22}&
\instbit{21}&
\instbitrange{20}{19}&
\instbit{18}&
\instbitrange{17}{15}&
\instbitrange{14}{13}&
\instbitrange{12}{10}&
\instbitrange{9}{8}&
\instbit{7}&
\instbitrange{6}{0} \\
\hline
\multicolumn{1}{|c|}{{\scriptsize qx[2:0]}} &
\multicolumn{1}{c|}{{\scriptsize qy[2:0]}} &
\multicolumn{1}{c|}{1010} &
\multicolumn{1}{c|}{{\scriptsize sat[0]}} &
\multicolumn{1}{c|}{{\scriptsize imm16[3:2]}} &
\multicolumn{1}{c|}{{\scriptsize rs1[4]}} &
\multicolumn{1}{c|}{{\scriptsize rs1[2:0]}} &
\multicolumn{1}{c|}{11} &
\multicolumn{1}{c|}{{\scriptsize qu[2:0]}} &
\multicolumn{1}{c|}{{\scriptsize imm16[1:0]}} &
\multicolumn{1}{c|}{0} &
\multicolumn{1}{c|}{0011111}\\
\hline
3& 3& 4& 1& 2& 1& 3& 2& 3& 2& 1& 7 \\
\end{tabular}
 \end{center}
 }
 \vspace{-0.1in}
 \caption[]{ESP.VMULAS.U16.QACC.LD.IP} \end{figure*}


\newpage
\subsubsection{ESP.VMULAS.U16.QACC.LD.XP}\label{instruction:ESPVMULASU16QACCLDXP}
\textbf{Syntax}\\
ESP.VMULAS.U16.QACC.LD.XP qu,rs1,rs2,qx,qy,sat

\textbf{Description}\\

This instruction divides registers \lstinline{qx} and \lstinline{qy} into 8 data segments of 16 bits each. Then, it performs the unsigned multiplication of 8 sets of segments and adds the result to the corresponding 64-bit data segment in special registers \lstinline{QACC_H} and \lstinline{QACC_L} respectively. The calculated result is saturated to a 64-bit signed number and then stored in the corresponding 64-bit data segment in \lstinline{QACC_H} and \lstinline{QACC_L}.

During the operation, it loads the 16-byte data from the memory into the register \lstinline{qu}. After the access is completed, the value in the register \lstinline{rs1} is incremented by the value of the register \lstinline{rs2}.

\textbf{Operation}\\
\begin{lstlisting}
QACC_L[63 :  0] = min(QACC_L[63 :  0] + qx[15 :  0] * qy[15 :  0], 2^63-1)
QACC_L[127: 64] = min(QACC_L[127: 64] + qx[31 : 16] * qy[31 : 16], 2^63-1)
QACC_L[191:128] = min(QACC_L[191:128] + qx[47 : 32] * qy[47 : 32], 2^63-1)
QACC_L[255:192] = min(QACC_L[255:192] + qx[63 : 48] * qy[63 : 48], 2^63-1)
QACC_H[63 :  0] = min(QACC_H[63 :  0] + qx[79 : 64] * qy[79 : 64], 2^63-1)
QACC_H[127: 64] = min(QACC_H[127: 64] + qx[95 : 80] * qy[95 : 80], 2^63-1)
QACC_H[191:128] = min(QACC_H[191:128] + qx[111: 96] * qy[111: 96], 2^63-1)
QACC_H[255:192] = min(QACC_H[255:192] + qx[127:112] * qy[127:112], 2^63-1)
qu[127:0] = load128(rs1[31:0])
rs1[31:0] = rs1[31:0] + rs2[31:0]
\end{lstlisting}

\textbf{Schedule}\\
\begin{longtable}[c]{|l|l|l|}
    \hline
    \rowcolor{lightgray}
    \textbf{Operands} & \textbf{Use} & \textbf{Def} \\\hline
    \endfirsthead
    \hline
    \rowcolor{lightgray}
     \textbf{Operands} & \textbf{Use} & \textbf{Def} \\\hline
    \endhead
    \hline
    \endfoot

    qu & - & 2 \\\hline
    rs1 & 1 & 1 \\\hline
    rs2 & 1 & - \\\hline
     qx & 1 & - \\\hline
     qy & 1 & - \\\hline
     QACC\_H & 2 & 2 \\\hline
     QACC\_L & 2 & 2

\end{longtable}

\textbf{Format}\\

\begin{figure*}[!htbp]
 {\setlength{\tabcolsep}{1pt}
 \begin{center}
\begin{tabular}{@{}F@{}F@{}W@{}I@{}F@{}I@{}I@{}F@{}W@{}F@{}I@{}I@{}I@{}O}
\instbitrange{31}{29}&
\instbitrange{28}{26}&
\instbitrange{25}{24}&
\instbit{23}&
\instbitrange{22}{20}&
\instbit{19}&
\instbit{18}&
\instbitrange{17}{15}&
\instbitrange{14}{13}&
\instbitrange{12}{10}&
\instbit{9}&
\instbit{8}&
\instbit{7}&
\instbitrange{6}{0} \\
\hline
\multicolumn{1}{|c|}{{\scriptsize qx[2:0]}} &
\multicolumn{1}{c|}{{\scriptsize qy[2:0]}} &
\multicolumn{1}{c|}{10} &
\multicolumn{1}{c|}{{\scriptsize rs2[4]}} &
\multicolumn{1}{c|}{{\scriptsize rs2[2:0]}} &
\multicolumn{1}{c|}{1} &
\multicolumn{1}{c|}{{\scriptsize rs1[4]}} &
\multicolumn{1}{c|}{{\scriptsize rs1[2:0]}} &
\multicolumn{1}{c|}{01} &
\multicolumn{1}{c|}{{\scriptsize qu[2:0]}} &
\multicolumn{1}{c|}{0} &
\multicolumn{1}{c|}{{\scriptsize sat[0]}} &
\multicolumn{1}{c|}{0} &
\multicolumn{1}{c|}{0011111}\\
\hline
3& 3& 2& 1& 3& 1& 1& 3& 2& 3& 1& 1& 1& 7 \\
\end{tabular}
 \end{center}
 }
 \vspace{-0.1in}
 \caption[]{ESP.VMULAS.U16.QACC.LD.XP} \end{figure*}


\newpage
\subsubsection{ESP.VMULAS.U16.QACC.ST.IP}\label{instruction:ESPVMULASU16QACCSTIP}
\textbf{Syntax}\\
ESP.VMULAS.U16.QACC.ST.IP qu,rs1,imm16,qx,qy,sat

\textbf{Description}\\
This instruction divides registers \lstinline{qx} and \lstinline{qy} into 8 data segments of 16 bits each. Then, the unsigned multiplication results of 8 sets of segments are added to the corresponding 64-bit data segments in special registers \lstinline{QACC_H} and \lstinline{QACC_L} respectively. The calculated result is saturated to a 64-bit signed number and then stored in the corresponding 64-bit data segments in \lstinline{QACC_H} and \lstinline{QACC_L}.

During the operation, the register \lstinline{qu} is stored to memory. After the access is completed, the value in the register \lstinline{rs1} is incremented by a 4-bit sign-extended constant in the instruction code segment, left-shifted by 4.

imm starts at -128, ends at 112, with a step of 16.

\textbf{Operation}\\
\begin{lstlisting}
QACC_L[63 :  0] = min(QACC_L[63 :  0] + qx[15 :  0] * qy[15 :  0], 2^63-1)
QACC_L[127: 64] = min(QACC_L[127: 64] + qx[31 : 16] * qy[31 : 16], 2^63-1)
QACC_L[191:128] = min(QACC_L[191:128] + qx[47 : 32] * qy[47 : 32], 2^63-1)
QACC_L[255:192] = min(QACC_L[255:192] + qx[63 : 48] * qy[63 : 48], 2^63-1)
QACC_H[63 :  0] = min(QACC_H[63 :  0] + qx[79 : 64] * qy[79 : 64], 2^63-1)
QACC_H[127: 64] = min(QACC_H[127: 64] + qx[95 : 80] * qy[95 : 80], 2^63-1)
QACC_H[191:128] = min(QACC_H[191:128] + qx[111: 96] * qy[111: 96], 2^63-1)
QACC_H[255:192] = min(QACC_H[255:192] + qx[127:112] * qy[127:112], 2^63-1)

qu[127:0] => store128(rs1[31:0])
rs1[31:0] = rs1[31:0] + imm
\end{lstlisting}

\textbf{Schedule}\\
\begin{longtable}[c]{|l|l|l|}
    \hline
    \rowcolor{lightgray}
    \textbf{Operands} & \textbf{Use} & \textbf{Def} \\\hline
    \endfirsthead
    \hline
    \rowcolor{lightgray}
     \textbf{Operands} & \textbf{Use} & \textbf{Def} \\\hline
    \endhead
    \hline
    \endfoot

    qu & 2 & - \\\hline
    rs1 & 1 & 1 \\\hline
     qx & 1 & - \\\hline
     qy & 1 & - \\\hline
     QACC\_H & 2 & 2 \\\hline
     QACC\_L & 2 & 2

\end{longtable}

\textbf{Format}\\

\begin{figure*}[!htbp]
 {\setlength{\tabcolsep}{1pt}
 \begin{center}
\begin{tabular}{@{}F@{}F@{}Y@{}I@{}W@{}I@{}F@{}W@{}F@{}W@{}I@{}O}
\instbitrange{31}{29}&
\instbitrange{28}{26}&
\instbitrange{25}{22}&
\instbit{21}&
\instbitrange{20}{19}&
\instbit{18}&
\instbitrange{17}{15}&
\instbitrange{14}{13}&
\instbitrange{12}{10}&
\instbitrange{9}{8}&
\instbit{7}&
\instbitrange{6}{0} \\
\hline
\multicolumn{1}{|c|}{{\scriptsize qx[2:0]}} &
\multicolumn{1}{c|}{{\scriptsize qy[2:0]}} &
\multicolumn{1}{c|}{1110} &
\multicolumn{1}{c|}{{\scriptsize sat[0]}} &
\multicolumn{1}{c|}{{\scriptsize imm16[3:2]}} &
\multicolumn{1}{c|}{{\scriptsize rs1[4]}} &
\multicolumn{1}{c|}{{\scriptsize rs1[2:0]}} &
\multicolumn{1}{c|}{11} &
\multicolumn{1}{c|}{{\scriptsize qu[2:0]}} &
\multicolumn{1}{c|}{{\scriptsize imm16[1:0]}} &
\multicolumn{1}{c|}{0} &
\multicolumn{1}{c|}{0011111}\\
\hline
3& 3& 4& 1& 2& 1& 3& 2& 3& 2& 1& 7 \\
\end{tabular}
 \end{center}
 }
 \vspace{-0.1in}
 \caption[]{ESP.VMULAS.U16.QACC.ST.IP} \end{figure*}


\newpage
\subsubsection{ESP.VMULAS.U16.QACC.ST.XP}\label{instruction:ESPVMULASU16QACCSTXP}
\textbf{Syntax}\\
ESP.VMULAS.U16.QACC.ST.XP qu,rs1,rs2,qx,qy,sat

\textbf{Description}\\
This instruction divides registers \lstinline{qx} and \lstinline{qy} into 8 data segments of 16 bits each. Then, the unsigned multiplication results of the 8 sets of segments are added to the corresponding 64-bit data segments in special registers \lstinline{QACC_H} and \lstinline{QACC_L} respectively. The calculated result is saturated to a 64-bit signed number and then stored in the corresponding 64-bit data segment in \lstinline{QACC_H} and \lstinline{QACC_L}.

During the operation, the register \lstinline{qu} is stored to memory. After the access is completed, the value in the register \lstinline{rs1} is incremented by the value of the register \lstinline{rs2}.

\textbf{Operation}\\
\begin{lstlisting}
QACC_L[63 :  0] = min(QACC_L[63 :  0] + qx[15 :  0] * qy[15 :  0], 2^63-1)
QACC_L[127: 64] = min(QACC_L[127: 64] + qx[31 : 16] * qy[31 : 16], 2^63-1)
QACC_L[191:128] = min(QACC_L[191:128] + qx[47 : 32] * qy[47 : 32], 2^63-1)
QACC_L[255:192] = min(QACC_L[255:192] + qx[63 : 48] * qy[63 : 48], 2^63-1)
QACC_H[63 :  0] = min(QACC_H[63 :  0] + qx[79 : 64] * qy[79 : 64], 2^63-1)
QACC_H[127: 64] = min(QACC_H[127: 64] + qx[95 : 80] * qy[95 : 80], 2^63-1)
QACC_H[191:128] = min(QACC_H[191:128] + qx[111: 96] * qy[111: 96], 2^63-1)
QACC_H[255:192] = min(QACC_H[255:192] + qx[127:112] * qy[127:112], 2^63-1)
qu[127:0] => store128(rs1[31:0])
rs1[31:0] = rs1[31:0] + rs2[31:0]
\end{lstlisting}

\textbf{Schedule}\\
\begin{longtable}[c]{|l|l|l|}
    \hline
    \rowcolor{lightgray}
    \textbf{Operands} & \textbf{Use} & \textbf{Def} \\\hline
    \endfirsthead
    \hline
    \rowcolor{lightgray}
     \textbf{Operands} & \textbf{Use} & \textbf{Def} \\\hline
    \endhead
    \hline
    \endfoot

    qu & 2 & - \\\hline
    rs1 & 1 & 1 \\\hline
    rs2 & 1 & - \\\hline
     qx & 1 & - \\\hline
     qy & 1 & - \\\hline
     QACC\_H & 2 & 2 \\\hline
     QACC\_L & 2 & 2

\end{longtable}


\textbf{Format}\\

\begin{figure*}[!htbp]
 {\setlength{\tabcolsep}{1pt}
 \begin{center}
\begin{tabular}{@{}F@{}F@{}W@{}I@{}F@{}I@{}I@{}F@{}W@{}F@{}I@{}I@{}I@{}O}
\instbitrange{31}{29}&
\instbitrange{28}{26}&
\instbitrange{25}{24}&
\instbit{23}&
\instbitrange{22}{20}&
\instbit{19}&
\instbit{18}&
\instbitrange{17}{15}&
\instbitrange{14}{13}&
\instbitrange{12}{10}&
\instbit{9}&
\instbit{8}&
\instbit{7}&
\instbitrange{6}{0} \\
\hline
\multicolumn{1}{|c|}{{\scriptsize qx[2:0]}} &
\multicolumn{1}{c|}{{\scriptsize qy[2:0]}} &
\multicolumn{1}{c|}{11} &
\multicolumn{1}{c|}{{\scriptsize rs2[4]}} &
\multicolumn{1}{c|}{{\scriptsize rs2[2:0]}} &
\multicolumn{1}{c|}{1} &
\multicolumn{1}{c|}{{\scriptsize rs1[4]}} &
\multicolumn{1}{c|}{{\scriptsize rs1[2:0]}} &
\multicolumn{1}{c|}{01} &
\multicolumn{1}{c|}{{\scriptsize qu[2:0]}} &
\multicolumn{1}{c|}{0} &
\multicolumn{1}{c|}{{\scriptsize sat[0]}} &
\multicolumn{1}{c|}{0} &
\multicolumn{1}{c|}{0011111}\\
\hline
3& 3& 2& 1& 3& 1& 1& 3& 2& 3& 1& 1& 1& 7 \\
\end{tabular}
 \end{center}
 }
 \vspace{-0.1in}
 \caption[]{ESP.VMULAS.U16.QACC.ST.XP} \end{figure*}


\newpage
\subsubsection{ESP.VMULAS.U16.XACC}\label{instruction:ESPVMULASU16XACC}
\textbf{Syntax}\\
ESP.VMULAS.U16.XACC qx,qy,sat

\textbf{Description}\\
This instruction divides registers \lstinline{qx} and \lstinline{qy} into 8 data segments of 16 bits each. Then, it performs an unsigned multiply-accumulate operation on each of the 8 sets of segments. The accumulated result is added to the value in the special register \lstinline{XACC}. The resulting sum is then saturated and stored in \lstinline{XACC}.\\
\textbf{Operation}
\begin{lstlisting}
  add0[31:0] = qx[ 15:  0] * qy[ 15:  0]
  add1[31:0] = qx[ 31: 16] * qy[ 31: 16]
  ...
  add7[31:0] = qx[127:112] * qy[127:112]
  sum[40:0] = XACC[39:0] + add0[31:0] + add1[31:0] + ... + add7[31:0]

  XACC[39:0] = min(max(sum[40:0], 0), 2^{40}-1)
\end{lstlisting}

\textbf{Schedule}\\
\begin{longtable}[c]{|l|l|l|}
    \hline
    \rowcolor{lightgray}
    \textbf{Operands} & \textbf{Use} & \textbf{Def} \\\hline
    \endfirsthead
    \hline
    \rowcolor{lightgray}
     \textbf{Operands} & \textbf{Use} & \textbf{Def} \\\hline
    \endhead
    \hline
    \endfoot

     qx & 1 & - \\\hline
     qy & 1 & - \\\hline
     XACC & 3 & 3

\end{longtable}

\textbf{Format}\\

\begin{figure*}[!htbp]
 {\setlength{\tabcolsep}{1pt}
 \begin{center}
\begin{tabular}{@{}F@{}F@{}T@{}I@{}E@{}O}
\instbitrange{31}{29}&
\instbitrange{28}{26}&
\instbitrange{25}{16}&
\instbit{15}&
\instbitrange{14}{7}&
\instbitrange{6}{0} \\
\hline
\multicolumn{1}{|c|}{{\scriptsize qx[2:0]}} &
\multicolumn{1}{c|}{{\scriptsize qy[2:0]}} &
\multicolumn{1}{c|}{1X101XX010} &
\multicolumn{1}{c|}{{\scriptsize sat[0]}} &
\multicolumn{1}{c|}{01XX0X0X} &
\multicolumn{1}{c|}{0011011}\\
\hline
3& 3& 10& 1& 8& 7 \\
\end{tabular}
 \end{center}
 }
 \vspace{-0.1in}
 \caption[]{ESP.VMULAS.U16.XACC} \end{figure*}


\newpage
\subsubsection{ESP.VMULAS.U16.XACC.LD.IP}\label{instruction:ESPVMULASU16XACCLDIP}
\textbf{Syntax}\\
ESP.VMULAS.U16.XACC.LD.IP qu,rs1,imm16,qx,qy,sat

\textbf{Description}\\
This instruction divides registers \lstinline{qx} and \lstinline{qy} into 8 data segments of 16 bits each. Then, it performs an unsigned multiply-accumulate operation on each of the 8 segment sets. The accumulated result is added to the value in the special register \lstinline{XACC}. The resulting sum is then saturated and stored in \lstinline{XACC}. \\ During the operation, the 16-byte data is loaded from the memory into the register \lstinline{qu}. After the access is completed, the value in the register \lstinline{rs1} is incremented by a 4-bit sign-extended constant from the instruction code segment, which is left-shifted by 4.\\

imm starts at -128, ends at 112, and steps by 16.

\textbf{Operation}
\begin{lstlisting}
  add0[31:0] = qx[ 15:  0] * qy[ 15:  0]
  add1[31:0] = qx[ 31: 16] * qy[ 31: 16]
  ...
  add7[31:0] = qx[127:112] * qy[127:112]
  sum[40:0] = XACC[39:0] + add0[31:0] + add1[31:0] + ... + add7[31:0]
  XACC[40:0] = min(max(sum[40:0], 0), 2^{40}-1)

  qu[127:0] = load128(rs1[31:0])
  rs1[31:0] = rs1[31:0] + imm
\end{lstlisting}

\textbf{Schedule}\\
\begin{longtable}[c]{|l|l|l|}
    \hline
    \rowcolor{lightgray}
    \textbf{Operands} & \textbf{Use} & \textbf{Def} \\\hline
    \endfirsthead
    \hline
    \rowcolor{lightgray}
     \textbf{Operands} & \textbf{Use} & \textbf{Def} \\\hline
    \endhead
    \hline
    \endfoot

    qu & - & 2 \\\hline
    rs1 & 1 & 1 \\\hline
     qx & 1 & - \\\hline
     qy & 1 & - \\\hline
     XACC & 3 & 3

\end{longtable}

\textbf{Format}\\

\begin{figure*}[!htbp]
 {\setlength{\tabcolsep}{1pt}
 \begin{center}
\begin{tabular}{@{}F@{}F@{}Y@{}I@{}W@{}I@{}F@{}W@{}F@{}W@{}I@{}O}
\instbitrange{31}{29}&
\instbitrange{28}{26}&
\instbitrange{25}{22}&
\instbit{21}&
\instbitrange{20}{19}&
\instbit{18}&
\instbitrange{17}{15}&
\instbitrange{14}{13}&
\instbitrange{12}{10}&
\instbitrange{9}{8}&
\instbit{7}&
\instbitrange{6}{0} \\
\hline
\multicolumn{1}{|c|}{{\scriptsize qx[2:0]}} &
\multicolumn{1}{c|}{{\scriptsize qy[2:0]}} &
\multicolumn{1}{c|}{0010} &
\multicolumn{1}{c|}{{\scriptsize sat[0]}} &
\multicolumn{1}{c|}{{\scriptsize imm16[3:2]}} &
\multicolumn{1}{c|}{{\scriptsize rs1[4]}} &
\multicolumn{1}{c|}{{\scriptsize rs1[2:0]}} &
\multicolumn{1}{c|}{11} &
\multicolumn{1}{c|}{{\scriptsize qu[2:0]}} &
\multicolumn{1}{c|}{{\scriptsize imm16[1:0]}} &
\multicolumn{1}{c|}{0} &
\multicolumn{1}{c|}{0011111}\\
\hline
3& 3& 4& 1& 2& 1& 3& 2& 3& 2& 1& 7 \\
\end{tabular}
 \end{center}
 }
 \vspace{-0.1in}
 \caption[]{ESP.VMULAS.U16.XACC.LD.IP} \end{figure*}


\newpage
\subsubsection{ESP.VMULAS.U16.XACC.LD.XP}\label{instruction:ESPVMULASU16XACCLDXP}
\textbf{Syntax}\\
ESP.VMULAS.U16.XACC.LD.XP qu,rs1,rs2,qx,qy,sat

\textbf{Description}\\
This instruction divides registers \lstinline{qx} and \lstinline{qy} into 8 data segments of 16 bits each. Then, it performs an unsigned multiply-accumulate operation on each of the 8 segment sets. The accumulated result is added to the value in the special register \lstinline{XACC}. The resulting sum is then saturated and stored in \lstinline{XACC}. \\ During the operation, it loads the 16-byte data from the memory into the register \lstinline{qu}. After the access is completed, the value in the register \lstinline{rs1} is incremented by the value of the register \lstinline{rs2}.\\

\textbf{Operation}
\begin{lstlisting}
  add0[31:0] = qx[ 15:  0] * qy[ 15:  0]
  add1[31:0] = qx[ 31: 16] * qy[ 31: 16]
  ...
  add7[31:0] = qx[127:112] * qy[127:112]
  sum[40:0] = XACC[39:0] + add0[31:0] + add1[31:0] + ... + add7[31:0]
  XACC[40:0] = min(max(sum[40:0], 0), 2^{40}-1)

  qu[127:0] = load128(rs1[31:0])
  rs1[31:0] = rs1[31:0] + rs2[31:0]
\end{lstlisting}

\textbf{Schedule}\\
\begin{longtable}[c]{|l|l|l|}
    \hline
    \rowcolor{lightgray}
    \textbf{Operands} & \textbf{Use} & \textbf{Def} \\\hline
    \endfirsthead
    \hline
    \rowcolor{lightgray}
     \textbf{Operands} & \textbf{Use} & \textbf{Def} \\\hline
    \endhead
    \hline
    \endfoot

    qu & - & 2 \\\hline
    rs1 & 1 & 1 \\\hline
    rs2 & 1 & - \\\hline
     qx & 1 & - \\\hline
     qy & 1 & - \\\hline
     XACC & 3 & 3

\end{longtable}

\textbf{Format}\\

\begin{figure*}[!htbp]
 {\setlength{\tabcolsep}{1pt}
 \begin{center}
\begin{tabular}{@{}F@{}F@{}W@{}I@{}F@{}I@{}I@{}F@{}W@{}F@{}I@{}I@{}I@{}O}
\instbitrange{31}{29}&
\instbitrange{28}{26}&
\instbitrange{25}{24}&
\instbit{23}&
\instbitrange{22}{20}&
\instbit{19}&
\instbit{18}&
\instbitrange{17}{15}&
\instbitrange{14}{13}&
\instbitrange{12}{10}&
\instbit{9}&
\instbit{8}&
\instbit{7}&
\instbitrange{6}{0} \\
\hline
\multicolumn{1}{|c|}{{\scriptsize qx[2:0]}} &
\multicolumn{1}{c|}{{\scriptsize qy[2:0]}} &
\multicolumn{1}{c|}{00} &
\multicolumn{1}{c|}{{\scriptsize rs2[4]}} &
\multicolumn{1}{c|}{{\scriptsize rs2[2:0]}} &
\multicolumn{1}{c|}{1} &
\multicolumn{1}{c|}{{\scriptsize rs1[4]}} &
\multicolumn{1}{c|}{{\scriptsize rs1[2:0]}} &
\multicolumn{1}{c|}{01} &
\multicolumn{1}{c|}{{\scriptsize qu[2:0]}} &
\multicolumn{1}{c|}{0} &
\multicolumn{1}{c|}{{\scriptsize sat[0]}} &
\multicolumn{1}{c|}{0} &
\multicolumn{1}{c|}{0011111}\\
\hline
3& 3& 2& 1& 3& 1& 1& 3& 2& 3& 1& 1& 1& 7 \\
\end{tabular}
 \end{center}
 }
 \vspace{-0.1in}
 \caption[]{ESP.VMULAS.U16.XACC.LD.XP} \end{figure*}


\newpage
\subsubsection{ESP.VMULAS.U16.XACC.ST.IP}\label{instruction:ESPVMULASU16XACCSTIP}
\textbf{Syntax}\\
ESP.VMULAS.U16.XACC.ST.IP qu,rs1,imm16,qx,qy,sat

\textbf{Description}\\
This instruction divides registers \lstinline{qx} and \lstinline{qy} into 8 data segments of 16 bits each. Then, it performs an unsigned multiply-accumulate operation on the 8 sets of segments respectively. The accumulated result is added to the value in the special register \lstinline{XACC}. Then, the sum obtained is saturated and stored in \lstinline{XACC}. \\ During the operation, the register \lstinline{qu} is stored to memory. After the access is completed, the value in the register \lstinline{rs1} is incremented by a 4-bit sign-extended constant in the instruction code segment, which is left-shifted by 4.\\

imm starts at -128, ends at 112, and steps by 16.

\textbf{Operation}
\begin{lstlisting}
  add0[31:0] = qx[ 15:  0] * qy[ 15:  0]
  add1[31:0] = qx[ 31: 16] * qy[ 31: 16]
  ...
  add7[31:0] = qx[127:112] * qy[127:112]
  sum[40:0] = XACC[39:0] + add0[31:0] + add1[31:0] + ... + add7[31:0]
  XACC[40:0] = min(max(sum[40:0], 0), 2^{40}-1)

  qu[127:0] => store128(rs1[31:0])
  rs1[31:0] = rs1[31:0] + imm
\end{lstlisting}

\textbf{Schedule}\\
\begin{longtable}[c]{|l|l|l|}
    \hline
    \rowcolor{lightgray}
    \textbf{Operands} & \textbf{Use} & \textbf{Def} \\\hline
    \endfirsthead
    \hline
    \rowcolor{lightgray}
     \textbf{Operands} & \textbf{Use} & \textbf{Def} \\\hline
    \endhead
    \hline
    \endfoot

    qu & 2 & - \\\hline
    rs1 & 1 & 1 \\\hline
     qx & 1 & - \\\hline
     qy & 1 & - \\\hline
     XACC & 3 & 3

\end{longtable}

\textbf{Format}\\

\begin{figure*}[!htbp]
 {\setlength{\tabcolsep}{1pt}
 \begin{center}
\begin{tabular}{@{}F@{}F@{}Y@{}I@{}W@{}I@{}F@{}W@{}F@{}W@{}I@{}O}
\instbitrange{31}{29}&
\instbitrange{28}{26}&
\instbitrange{25}{22}&
\instbit{21}&
\instbitrange{20}{19}&
\instbit{18}&
\instbitrange{17}{15}&
\instbitrange{14}{13}&
\instbitrange{12}{10}&
\instbitrange{9}{8}&
\instbit{7}&
\instbitrange{6}{0} \\
\hline
\multicolumn{1}{|c|}{{\scriptsize qx[2:0]}} &
\multicolumn{1}{c|}{{\scriptsize qy[2:0]}} &
\multicolumn{1}{c|}{0110} &
\multicolumn{1}{c|}{{\scriptsize sat[0]}} &
\multicolumn{1}{c|}{{\scriptsize imm16[3:2]}} &
\multicolumn{1}{c|}{{\scriptsize rs1[4]}} &
\multicolumn{1}{c|}{{\scriptsize rs1[2:0]}} &
\multicolumn{1}{c|}{11} &
\multicolumn{1}{c|}{{\scriptsize qu[2:0]}} &
\multicolumn{1}{c|}{{\scriptsize imm16[1:0]}} &
\multicolumn{1}{c|}{0} &
\multicolumn{1}{c|}{0011111}\\
\hline
3& 3& 4& 1& 2& 1& 3& 2& 3& 2& 1& 7 \\
\end{tabular}
 \end{center}
 }
 \vspace{-0.1in}
 \caption[]{ESP.VMULAS.U16.XACC.ST.IP} \end{figure*}


\newpage
\subsubsection{ESP.VMULAS.U16.XACC.ST.XP}\label{instruction:ESPVMULASU16XACCSTXP}
\textbf{Syntax}\\
ESP.VMULAS.U16.XACC.ST.XP qu,rs1,rs2,qx,qy,sat

\textbf{Description}\\
This instruction divides registers \lstinline{qx} and \lstinline{qy} into 8 data segments of 16 bits each. Then, it performs an unsigned multiply-accumulate operation on each of the 8 segment sets. The accumulated result is added to the value in the special register \lstinline{XACC}. The resulting sum is then saturated and stored in \lstinline{XACC}. \\ During the operation, the register \lstinline{qu} is stored to memory. After the access is completed, the value in the register \lstinline{rs1} is incremented by the value of the register \lstinline{rs2}.\\

\textbf{Operation}
\begin{lstlisting}
  add0[31:0] = qx[ 15:  0] * qy[ 15:  0]
  add1[31:0] = qx[ 31: 16] * qy[ 31: 16]
  ...
  add7[31:0] = qx[127:112] * qy[127:112]
  sum[40:0] = XACC[39:0] + add0[31:0] + add1[31:0] + ... + add7[31:0]
  XACC[40:0] = min(max(sum[40:0], 0), 2^{40}-1)

  qu[127:0] => store128(rs1[31:0])
  rs1[31:0] = rs1[31:0] + rs2[31:0]
\end{lstlisting}

\textbf{Schedule}\\
\begin{longtable}[c]{|l|l|l|}
    \hline
    \rowcolor{lightgray}
    \textbf{Operands} & \textbf{Use} & \textbf{Def} \\\hline
    \endfirsthead
    \hline
    \rowcolor{lightgray}
     \textbf{Operands} & \textbf{Use} & \textbf{Def} \\\hline
    \endhead
    \hline
    \endfoot

    qu & 2 & - \\\hline
    rs1 & 1 & 1 \\\hline
    rs2 & 1 & - \\\hline
     qx & 1 & - \\\hline
     qy & 1 & - \\\hline
     XACC & 3 & 3

\end{longtable}

\textbf{Format}\\

\begin{figure*}[!htbp]
 {\setlength{\tabcolsep}{1pt}
 \begin{center}
\begin{tabular}{@{}F@{}F@{}W@{}I@{}F@{}I@{}I@{}F@{}W@{}F@{}I@{}I@{}I@{}O}
\instbitrange{31}{29}&
\instbitrange{28}{26}&
\instbitrange{25}{24}&
\instbit{23}&
\instbitrange{22}{20}&
\instbit{19}&
\instbit{18}&
\instbitrange{17}{15}&
\instbitrange{14}{13}&
\instbitrange{12}{10}&
\instbit{9}&
\instbit{8}&
\instbit{7}&
\instbitrange{6}{0} \\
\hline
\multicolumn{1}{|c|}{{\scriptsize qx[2:0]}} &
\multicolumn{1}{c|}{{\scriptsize qy[2:0]}} &
\multicolumn{1}{c|}{01} &
\multicolumn{1}{c|}{{\scriptsize rs2[4]}} &
\multicolumn{1}{c|}{{\scriptsize rs2[2:0]}} &
\multicolumn{1}{c|}{1} &
\multicolumn{1}{c|}{{\scriptsize rs1[4]}} &
\multicolumn{1}{c|}{{\scriptsize rs1[2:0]}} &
\multicolumn{1}{c|}{01} &
\multicolumn{1}{c|}{{\scriptsize qu[2:0]}} &
\multicolumn{1}{c|}{0} &
\multicolumn{1}{c|}{{\scriptsize sat[0]}} &
\multicolumn{1}{c|}{0} &
\multicolumn{1}{c|}{0011111}\\
\hline
3& 3& 2& 1& 3& 1& 1& 3& 2& 3& 1& 1& 1& 7 \\
\end{tabular}
 \end{center}
 }
 \vspace{-0.1in}
 \caption[]{ESP.VMULAS.U16.XACC.ST.XP} \end{figure*}


\newpage
\subsubsection{ESP.VMULAS.U8.QACC}\label{instruction:ESPVMULASU8QACC}
\textbf{Syntax}\\
ESP.VMULAS.U8.QACC qx,qy,sat

\textbf{Description}\\
This instruction divides registers \lstinline{qx} and \lstinline{qy} into 16 data segments of 8 bits each. Then, it performs the unsigned multiplication of the 16 sets of segments and adds the result to the corresponding 32-bit data segment in special registers \lstinline{QACC_H} and \lstinline{QACC_L} respectively. The calculated result is saturated to a 32-bit unsigned number and then stored in the corresponding 32-bit data segment in \lstinline{QACC_H} and \lstinline{QACC_L}.\\

\textbf{Operation}\\

\begin{lstlisting}
QACC_L[31 :  0] = min(QACC_L[31 :  0] + qx[7  :  0] * qy[7  :  0], 2^31-1)
QACC_L[63 : 32] = min(QACC_L[63 : 32] + qx[15 :  8] * qy[15 :  8], 2^31-1)
QACC_L[95 : 64] = min(QACC_L[95 : 64] + qx[23 : 16] * qy[23 : 16], 2^31-1)
QACC_L[127: 96] = min(QACC_L[127: 96] + qx[31 : 24] * qy[31 : 24], 2^31-1)
QACC_L[159:128] = min(QACC_L[159:128] + qx[39 : 32] * qy[39 : 32], 2^31-1)
QACC_L[191:160] = min(QACC_L[191:160] + qx[47 : 40] * qy[47 : 40], 2^31-1)
QACC_L[223:192] = min(QACC_L[223:192] + qx[55 : 48] * qy[55 : 48], 2^31-1)
QACC_L[255:224] = min(QACC_L[255:224] + qx[63 : 56] * qy[63 : 56], 2^31-1)
QACC_H[31 :  0] = min(QACC_H[31 :  0] + qx[71 : 64] * qy[71 : 64], 2^31-1)
QACC_H[63 : 32] = min(QACC_H[63 : 32] + qx[79 : 72] * qy[79 : 72], 2^31-1)
QACC_H[95 : 64] = min(QACC_H[95 : 64] + qx[87 : 80] * qy[87 : 80], 2^31-1)
QACC_H[127: 96] = min(QACC_H[127: 96] + qx[95 : 88] * qy[95 : 88], 2^31-1)
QACC_H[159:128] = min(QACC_H[159:128] + qx[103: 96] * qy[103: 96], 2^31-1)
QACC_H[191:160] = min(QACC_H[191:160] + qx[111:104] * qy[111:104], 2^31-1)
QACC_H[223:192] = min(QACC_H[223:192] + qx[119:112] * qy[119:112], 2^31-1)
QACC_H[255:224] = min(QACC_H[255:224] + qx[127:120] * qy[127:120], 2^31-1)

\end{lstlisting}

\textbf{Schedule}\\
\begin{longtable}[c]{|l|l|l|}
    \hline
    \rowcolor{lightgray}
    \textbf{Operands} & \textbf{Use} & \textbf{Def} \\\hline
    \endfirsthead
    \hline
    \rowcolor{lightgray}
     \textbf{Operands} & \textbf{Use} & \textbf{Def} \\\hline
    \endhead
    \hline
    \endfoot

     qx & 1 & - \\\hline
     qy & 1 & - \\\hline
     QACC\_H & 2 & 2 \\\hline
     QACC\_L & 2 & 2

\end{longtable}

\textbf{Format}\\

\begin{figure*}[!htbp]
 {\setlength{\tabcolsep}{1pt}
 \begin{center}
\begin{tabular}{@{}F@{}F@{}T@{}I@{}E@{}O}
\instbitrange{31}{29}&
\instbitrange{28}{26}&
\instbitrange{25}{16}&
\instbit{15}&
\instbitrange{14}{7}&
\instbitrange{6}{0} \\
\hline
\multicolumn{1}{|c|}{{\scriptsize qx[2:0]}} &
\multicolumn{1}{c|}{{\scriptsize qy[2:0]}} &
\multicolumn{1}{c|}{1X101XX100} &
\multicolumn{1}{c|}{{\scriptsize sat[0]}} &
\multicolumn{1}{c|}{01XX0X0X} &
\multicolumn{1}{c|}{0011011}\\
\hline
3& 3& 10& 1& 8& 7 \\
\end{tabular}
 \end{center}
 }
 \vspace{-0.1in}
 \caption[]{ESP.VMULAS.U8.QACC} \end{figure*}


\newpage
\subsubsection{ESP.VMULAS.U8.QACC.LD.IP}\label{instruction:ESPVMULASU8QACCLDIP}
\textbf{Syntax}\\
ESP.VMULAS.U8.QACC.LD.IP qu,rs1,imm16,qx,qy,sat

\textbf{Description}\\

This instruction divides registers \lstinline{qx} and \lstinline{qy} into 16 data segments of 8 bits each. Then, the unsigned multiplication results of the 16 sets of segments are added to the corresponding 32-bit data segments in special registers \lstinline{QACC_H} and \lstinline{QACC_L} respectively. The calculated result is saturated to a 32-bit unsigned number and then stored in the corresponding 32-bit data segment in \lstinline{QACC_H} and \lstinline{QACC_L}.\\

During the operation, the 16-byte data is loaded from the memory to the register \lstinline{qu}. After the access is completed, the value in the register \lstinline{rs1} is incremented by the 4-bit sign-extended constant in the instruction code segment, which is left-shifted by 4.

The value of imm starts at -128, ends at 112, and steps by 16.

\textbf{Operation}\\

\begin{lstlisting}
QACC_L[31 :  0] = min(QACC_L[31 :  0] + qx[7  :  0] * qy[7  :  0], 2^31-1)
QACC_L[63 : 32] = min(QACC_L[63 : 32] + qx[15 :  8] * qy[15 :  8], 2^31-1)
QACC_L[95 : 64] = min(QACC_L[95 : 64] + qx[23 : 16] * qy[23 : 16], 2^31-1)
QACC_L[127: 96] = min(QACC_L[127: 96] + qx[31 : 24] * qy[31 : 24], 2^31-1)
QACC_L[159:128] = min(QACC_L[159:128] + qx[39 : 32] * qy[39 : 32], 2^31-1)
QACC_L[191:160] = min(QACC_L[191:160] + qx[47 : 40] * qy[47 : 40], 2^31-1)
QACC_L[223:192] = min(QACC_L[223:192] + qx[55 : 48] * qy[55 : 48], 2^31-1)
QACC_L[255:224] = min(QACC_L[255:224] + qx[63 : 56] * qy[63 : 56], 2^31-1)
QACC_H[31 :  0] = min(QACC_H[31 :  0] + qx[71 : 64] * qy[71 : 64], 2^31-1)
QACC_H[63 : 32] = min(QACC_H[63 : 32] + qx[79 : 72] * qy[79 : 72], 2^31-1)
QACC_H[95 : 64] = min(QACC_H[95 : 64] + qx[87 : 80] * qy[87 : 80], 2^31-1)
QACC_H[127: 96] = min(QACC_H[127: 96] + qx[95 : 88] * qy[95 : 88], 2^31-1)
QACC_H[159:128] = min(QACC_H[159:128] + qx[103: 96] * qy[103: 96], 2^31-1)
QACC_H[191:160] = min(QACC_H[191:160] + qx[111:104] * qy[111:104], 2^31-1)
QACC_H[223:192] = min(QACC_H[223:192] + qx[119:112] * qy[119:112], 2^31-1)
QACC_H[255:224] = min(QACC_H[255:224] + qx[127:120] * qy[127:120], 2^31-1)

qu[127:0] = load128(rs1[31:0])
rs1[31:0] = rs1[31:0] + imm
\end{lstlisting}

\textbf{Schedule}\\
\begin{longtable}[c]{|l|l|l|}
    \hline
    \rowcolor{lightgray}
    \textbf{Operands} & \textbf{Use} & \textbf{Def} \\\hline
    \endfirsthead
    \hline
    \rowcolor{lightgray}
     \textbf{Operands} & \textbf{Use} & \textbf{Def} \\\hline
    \endhead
    \hline
    \endfoot

    qu & - & 2 \\\hline
    rs1 & 1 & 1 \\\hline
     qx & 1 & - \\\hline
     qy & 1 & - \\\hline
     QACC\_H & 2 & 2 \\\hline
     QACC\_L & 2 & 2

\end{longtable}

\textbf{Format}\\

\begin{figure*}[!htbp]
 {\setlength{\tabcolsep}{1pt}
 \begin{center}
\begin{tabular}{@{}F@{}F@{}Y@{}I@{}W@{}I@{}F@{}W@{}F@{}W@{}I@{}O}
\instbitrange{31}{29}&
\instbitrange{28}{26}&
\instbitrange{25}{22}&
\instbit{21}&
\instbitrange{20}{19}&
\instbit{18}&
\instbitrange{17}{15}&
\instbitrange{14}{13}&
\instbitrange{12}{10}&
\instbitrange{9}{8}&
\instbit{7}&
\instbitrange{6}{0} \\
\hline
\multicolumn{1}{|c|}{{\scriptsize qx[2:0]}} &
\multicolumn{1}{c|}{{\scriptsize qy[2:0]}} &
\multicolumn{1}{c|}{1000} &
\multicolumn{1}{c|}{{\scriptsize sat[0]}} &
\multicolumn{1}{c|}{{\scriptsize imm16[3:2]}} &
\multicolumn{1}{c|}{{\scriptsize rs1[4]}} &
\multicolumn{1}{c|}{{\scriptsize rs1[2:0]}} &
\multicolumn{1}{c|}{11} &
\multicolumn{1}{c|}{{\scriptsize qu[2:0]}} &
\multicolumn{1}{c|}{{\scriptsize imm16[1:0]}} &
\multicolumn{1}{c|}{0} &
\multicolumn{1}{c|}{0011111}\\
\hline
3& 3& 4& 1& 2& 1& 3& 2& 3& 2& 1& 7 \\
\end{tabular}
 \end{center}
 }
 \vspace{-0.1in}
 \caption[]{ESP.VMULAS.U8.QACC.LD.IP} \end{figure*}


\newpage
\subsubsection{ESP.VMULAS.U8.QACC.LD.XP}\label{instruction:ESPVMULASU8QACCLDXP}
\textbf{Syntax}\\
ESP.VMULAS.U8.QACC.LD.XP qu,rs1,rs2,qx,qy,sat

\textbf{Description}\\

This instruction divides registers \lstinline{qx} and \lstinline{qy} into 16 data segments of 8 bits each. Then, the unsigned multiplication results of the 16 sets of segments are added to the corresponding 32-bit data segments in special registers \lstinline{QACC_H} and \lstinline{QACC_L} respectively. The calculated result is saturated to a 32-bit unsigned number and then stored in the corresponding 32-bit data segment in \lstinline{QACC_H} and \lstinline{QACC_L}.\\

During the operation, the 16-byte data is loaded from the memory to the register \lstinline{qu}. After the access is completed, the value in the register \lstinline{rs1} is incremented by the value of the register \lstinline{rs2}.

\textbf{Operation}\\

\begin{lstlisting}
QACC_L[31 :  0] = min(QACC_L[31 :  0] + qx[7  :  0] * qy[7  :  0], 2^31-1)
QACC_L[63 : 32] = min(QACC_L[63 : 32] + qx[15 :  8] * qy[15 :  8], 2^31-1)
QACC_L[95 : 64] = min(QACC_L[95 : 64] + qx[23 : 16] * qy[23 : 16], 2^31-1)
QACC_L[127: 96] = min(QACC_L[127: 96] + qx[31 : 24] * qy[31 : 24], 2^31-1)
QACC_L[159:128] = min(QACC_L[159:128] + qx[39 : 32] * qy[39 : 32], 2^31-1)
QACC_L[191:160] = min(QACC_L[191:160] + qx[47 : 40] * qy[47 : 40], 2^31-1)
QACC_L[223:192] = min(QACC_L[223:192] + qx[55 : 48] * qy[55 : 48], 2^31-1)
QACC_L[255:224] = min(QACC_L[255:224] + qx[63 : 56] * qy[63 : 56], 2^31-1)
QACC_H[31 :  0] = min(QACC_H[31 :  0] + qx[71 : 64] * qy[71 : 64], 2^31-1)
QACC_H[63 : 32] = min(QACC_H[63 : 32] + qx[79 : 72] * qy[79 : 72], 2^31-1)
QACC_H[95 : 64] = min(QACC_H[95 : 64] + qx[87 : 80] * qy[87 : 80], 2^31-1)
QACC_H[127: 96] = min(QACC_H[127: 96] + qx[95 : 88] * qy[95 : 88], 2^31-1)
QACC_H[159:128] = min(QACC_H[159:128] + qx[103: 96] * qy[103: 96], 2^31-1)
QACC_H[191:160] = min(QACC_H[191:160] + qx[111:104] * qy[111:104], 2^31-1)
QACC_H[223:192] = min(QACC_H[223:192] + qx[119:112] * qy[119:112], 2^31-1)
QACC_H[255:224] = min(QACC_H[255:224] + qx[127:120] * qy[127:120], 2^31-1)

qu[127:0] = load128(rs1[31:0])
rs1[31:0] = rs1[31:0] + rs2[31:0]
\end{lstlisting}

\textbf{Schedule}\\
\begin{longtable}[c]{|l|l|l|}
    \hline
    \rowcolor{lightgray}
    \textbf{Operands} & \textbf{Use} & \textbf{Def} \\\hline
    \endfirsthead
    \hline
    \rowcolor{lightgray}
     \textbf{Operands} & \textbf{Use} & \textbf{Def} \\\hline
    \endhead
    \hline
    \endfoot

    qu & - & 2 \\\hline
    rs1 & 1 & 1 \\\hline
    rs2 & 1 & - \\\hline
     qx & 1 & - \\\hline
     qy & 1 & - \\\hline
     QACC\_H & 2 & 2 \\\hline
     QACC\_L & 2 & 2

\end{longtable}

\textbf{Format}\\

\begin{figure*}[!htbp]
 {\setlength{\tabcolsep}{1pt}
 \begin{center}
\begin{tabular}{@{}F@{}F@{}W@{}I@{}F@{}I@{}I@{}F@{}W@{}F@{}I@{}I@{}I@{}O}
\instbitrange{31}{29}&
\instbitrange{28}{26}&
\instbitrange{25}{24}&
\instbit{23}&
\instbitrange{22}{20}&
\instbit{19}&
\instbit{18}&
\instbitrange{17}{15}&
\instbitrange{14}{13}&
\instbitrange{12}{10}&
\instbit{9}&
\instbit{8}&
\instbit{7}&
\instbitrange{6}{0} \\
\hline
\multicolumn{1}{|c|}{{\scriptsize qx[2:0]}} &
\multicolumn{1}{c|}{{\scriptsize qy[2:0]}} &
\multicolumn{1}{c|}{10} &
\multicolumn{1}{c|}{{\scriptsize rs2[4]}} &
\multicolumn{1}{c|}{{\scriptsize rs2[2:0]}} &
\multicolumn{1}{c|}{0} &
\multicolumn{1}{c|}{{\scriptsize rs1[4]}} &
\multicolumn{1}{c|}{{\scriptsize rs1[2:0]}} &
\multicolumn{1}{c|}{01} &
\multicolumn{1}{c|}{{\scriptsize qu[2:0]}} &
\multicolumn{1}{c|}{0} &
\multicolumn{1}{c|}{{\scriptsize sat[0]}} &
\multicolumn{1}{c|}{0} &
\multicolumn{1}{c|}{0011111}\\
\hline
3& 3& 2& 1& 3& 1& 1& 3& 2& 3& 1& 1& 1& 7 \\
\end{tabular}
 \end{center}
 }
 \vspace{-0.1in}
 \caption[]{ESP.VMULAS.U8.QACC.LD.XP} \end{figure*}


\newpage
\subsubsection{ESP.VMULAS.U8.QACC.ST.IP}\label{instruction:ESPVMULASU8QACCSTIP}
\textbf{Syntax}\\
ESP.VMULAS.U8.QACC.ST.IP qu,rs1,imm16,qx,qy,sat

\textbf{Description}\\
This instruction divides registers \lstinline{qx} and \lstinline{qy} into 16 data segments of 8 bits each. Then, the unsigned multiplication results of the 16 sets of segments are added to the corresponding 32-bit data segments in special registers \lstinline{QACC_H} and \lstinline{QACC_L} respectively. The calculated result is saturated to a 32-bit unsigned number and then stored in the corresponding 32-bit data segment in \lstinline{QACC_H} and \lstinline{QACC_L}.\\

During the operation, the register \lstinline{qu} is stored to memory. After the access is completed, the value in the register \lstinline{rs1} is incremented by the 4-bit sign-extended constant in the instruction code segment, which is left-shifted by 4.

The imm starts at -128, ends at 112, and steps by 16.

\textbf{Operation}\\

\begin{lstlisting}
QACC_L[31 :  0] = min(QACC_L[31 :  0] + qx[7  :  0] * qy[7  :  0], 2^31-1)
QACC_L[63 : 32] = min(QACC_L[63 : 32] + qx[15 :  8] * qy[15 :  8], 2^31-1)
QACC_L[95 : 64] = min(QACC_L[95 : 64] + qx[23 : 16] * qy[23 : 16], 2^31-1)
QACC_L[127: 96] = min(QACC_L[127: 96] + qx[31 : 24] * qy[31 : 24], 2^31-1)
QACC_L[159:128] = min(QACC_L[159:128] + qx[39 : 32] * qy[39 : 32], 2^31-1)
QACC_L[191:160] = min(QACC_L[191:160] + qx[47 : 40] * qy[47 : 40], 2^31-1)
QACC_L[223:192] = min(QACC_L[223:192] + qx[55 : 48] * qy[55 : 48], 2^31-1)
QACC_L[255:224] = min(QACC_L[255:224] + qx[63 : 56] * qy[63 : 56], 2^31-1)
QACC_H[31 :  0] = min(QACC_H[31 :  0] + qx[71 : 64] * qy[71 : 64], 2^31-1)
QACC_H[63 : 32] = min(QACC_H[63 : 32] + qx[79 : 72] * qy[79 : 72], 2^31-1)
QACC_H[95 : 64] = min(QACC_H[95 : 64] + qx[87 : 80] * qy[87 : 80], 2^31-1)
QACC_H[127: 96] = min(QACC_H[127: 96] + qx[95 : 88] * qy[95 : 88], 2^31-1)
QACC_H[159:128] = min(QACC_H[159:128] + qx[103: 96] * qy[103: 96], 2^31-1)
QACC_H[191:160] = min(QACC_H[191:160] + qx[111:104] * qy[111:104], 2^31-1)
QACC_H[223:192] = min(QACC_H[223:192] + qx[119:112] * qy[119:112], 2^31-1)
QACC_H[255:224] = min(QACC_H[255:224] + qx[127:120] * qy[127:120], 2^31-1)

qu[127:0] => store128(rs1[31:0])
rs1[31:0] = rs1[31:0] + imm
\end{lstlisting}

\textbf{Schedule}\\
\begin{longtable}[c]{|l|l|l|}
    \hline
    \rowcolor{lightgray}
    \textbf{Operands} & \textbf{Use} & \textbf{Def} \\\hline
    \endfirsthead
    \hline
    \rowcolor{lightgray}
     \textbf{Operands} & \textbf{Use} & \textbf{Def} \\\hline
    \endhead
    \hline
    \endfoot

    qu & 2 & - \\\hline
    rs1 & 1 & 1 \\\hline
     qx & 1 & - \\\hline
     qy & 1 & - \\\hline
     QACC\_H & 2 & 2 \\\hline
     QACC\_L & 2 & 2

\end{longtable}

\textbf{Format}\\

\begin{figure*}[!htbp]
 {\setlength{\tabcolsep}{1pt}
 \begin{center}
\begin{tabular}{@{}F@{}F@{}Y@{}I@{}W@{}I@{}F@{}W@{}F@{}W@{}I@{}O}
\instbitrange{31}{29}&
\instbitrange{28}{26}&
\instbitrange{25}{22}&
\instbit{21}&
\instbitrange{20}{19}&
\instbit{18}&
\instbitrange{17}{15}&
\instbitrange{14}{13}&
\instbitrange{12}{10}&
\instbitrange{9}{8}&
\instbit{7}&
\instbitrange{6}{0} \\
\hline
\multicolumn{1}{|c|}{{\scriptsize qx[2:0]}} &
\multicolumn{1}{c|}{{\scriptsize qy[2:0]}} &
\multicolumn{1}{c|}{1100} &
\multicolumn{1}{c|}{{\scriptsize sat[0]}} &
\multicolumn{1}{c|}{{\scriptsize imm16[3:2]}} &
\multicolumn{1}{c|}{{\scriptsize rs1[4]}} &
\multicolumn{1}{c|}{{\scriptsize rs1[2:0]}} &
\multicolumn{1}{c|}{11} &
\multicolumn{1}{c|}{{\scriptsize qu[2:0]}} &
\multicolumn{1}{c|}{{\scriptsize imm16[1:0]}} &
\multicolumn{1}{c|}{0} &
\multicolumn{1}{c|}{0011111}\\
\hline
3& 3& 4& 1& 2& 1& 3& 2& 3& 2& 1& 7 \\
\end{tabular}
 \end{center}
 }
 \vspace{-0.1in}
 \caption[]{ESP.VMULAS.U8.QACC.ST.IP} \end{figure*}


\newpage
\subsubsection{ESP.VMULAS.U8.QACC.ST.XP}\label{instruction:ESPVMULASU8QACCSTXP}
\textbf{Syntax}\\
ESP.VMULAS.U8.QACC.ST.XP qu,rs1,rs2,qx,qy,sat

\textbf{Description}\\
This instruction divides registers \lstinline{qx} and \lstinline{qy} into 16 data segments of 8 bits each. Then, the unsigned multiplication results of the 16 sets of segments are added to the corresponding 32-bit data segments in special registers \lstinline{QACC_H} and \lstinline{QACC_L} respectively. The calculated result is saturated to a 32-bit unsigned number and then stored in the corresponding 32-bit data segment in \lstinline{QACC_H} and \lstinline{QACC_L}.\\

During the operation, the register \lstinline{qu} is stored to memory. After the access is completed, the value in the register \lstinline{rs1} is incremented by the value of the register \lstinline{rs2}.

\textbf{Operation}\\

\begin{lstlisting}
QACC_L[31 :  0] = min(QACC_L[31 :  0] + qx[7  :  0] * qy[7  :  0], 2^31-1)
QACC_L[63 : 32] = min(QACC_L[63 : 32] + qx[15 :  8] * qy[15 :  8], 2^31-1)
QACC_L[95 : 64] = min(QACC_L[95 : 64] + qx[23 : 16] * qy[23 : 16], 2^31-1)
QACC_L[127: 96] = min(QACC_L[127: 96] + qx[31 : 24] * qy[31 : 24], 2^31-1)
QACC_L[159:128] = min(QACC_L[159:128] + qx[39 : 32] * qy[39 : 32], 2^31-1)
QACC_L[191:160] = min(QACC_L[191:160] + qx[47 : 40] * qy[47 : 40], 2^31-1)
QACC_L[223:192] = min(QACC_L[223:192] + qx[55 : 48] * qy[55 : 48], 2^31-1)
QACC_L[255:224] = min(QACC_L[255:224] + qx[63 : 56] * qy[63 : 56], 2^31-1)
QACC_H[31 :  0] = min(QACC_H[31 :  0] + qx[71 : 64] * qy[71 : 64], 2^31-1)
QACC_H[63 : 32] = min(QACC_H[63 : 32] + qx[79 : 72] * qy[79 : 72], 2^31-1)
QACC_H[95 : 64] = min(QACC_H[95 : 64] + qx[87 : 80] * qy[87 : 80], 2^31-1)
QACC_H[127: 96] = min(QACC_H[127: 96] + qx[95 : 88] * qy[95 : 88], 2^31-1)
QACC_H[159:128] = min(QACC_H[159:128] + qx[103: 96] * qy[103: 96], 2^31-1)
QACC_H[191:160] = min(QACC_H[191:160] + qx[111:104] * qy[111:104], 2^31-1)
QACC_H[223:192] = min(QACC_H[223:192] + qx[119:112] * qy[119:112], 2^31-1)
QACC_H[255:224] = min(QACC_H[255:224] + qx[127:120] * qy[127:120], 2^31-1)

qu[127:0] => store128(rs1[31:0])
rs1[31:0] = rs1[31:0] + rs2[31:0]
\end{lstlisting}

\textbf{Schedule}\\
\begin{longtable}[c]{|l|l|l|}
    \hline
    \rowcolor{lightgray}
    \textbf{Operands} & \textbf{Use} & \textbf{Def} \\\hline
    \endfirsthead
    \hline
    \rowcolor{lightgray}
     \textbf{Operands} & \textbf{Use} & \textbf{Def} \\\hline
    \endhead
    \hline
    \endfoot

    qu & 2 & - \\\hline
    rs1 & 1 & 1 \\\hline
    rs2 & 1 & - \\\hline
     qx & 1 & - \\\hline
     qy & 1 & - \\\hline
     QACC\_H & 2 & 2 \\\hline
     QACC\_L & 2 & 2

\end{longtable}

\textbf{Format}\\

\begin{figure*}[!htbp]
 {\setlength{\tabcolsep}{1pt}
 \begin{center}
\begin{tabular}{@{}F@{}F@{}W@{}I@{}F@{}I@{}I@{}F@{}W@{}F@{}I@{}I@{}I@{}O}
\instbitrange{31}{29}&
\instbitrange{28}{26}&
\instbitrange{25}{24}&
\instbit{23}&
\instbitrange{22}{20}&
\instbit{19}&
\instbit{18}&
\instbitrange{17}{15}&
\instbitrange{14}{13}&
\instbitrange{12}{10}&
\instbit{9}&
\instbit{8}&
\instbit{7}&
\instbitrange{6}{0} \\
\hline
\multicolumn{1}{|c|}{{\scriptsize qx[2:0]}} &
\multicolumn{1}{c|}{{\scriptsize qy[2:0]}} &
\multicolumn{1}{c|}{11} &
\multicolumn{1}{c|}{{\scriptsize rs2[4]}} &
\multicolumn{1}{c|}{{\scriptsize rs2[2:0]}} &
\multicolumn{1}{c|}{0} &
\multicolumn{1}{c|}{{\scriptsize rs1[4]}} &
\multicolumn{1}{c|}{{\scriptsize rs1[2:0]}} &
\multicolumn{1}{c|}{01} &
\multicolumn{1}{c|}{{\scriptsize qu[2:0]}} &
\multicolumn{1}{c|}{0} &
\multicolumn{1}{c|}{{\scriptsize sat[0]}} &
\multicolumn{1}{c|}{0} &
\multicolumn{1}{c|}{0011111}\\
\hline
3& 3& 2& 1& 3& 1& 1& 3& 2& 3& 1& 1& 1& 7 \\
\end{tabular}
 \end{center}
 }
 \vspace{-0.1in}
 \caption[]{ESP.VMULAS.U8.QACC.ST.XP} \end{figure*}


\newpage
\subsubsection{ESP.VMULAS.U8.XACC}\label{instruction:ESPVMULASU8XACC}
\textbf{Syntax}\\
ESP.VMULAS.U8.XACC qx,qy,sat

\textbf{Description}\\
This instruction divides registers \lstinline{qx} and \lstinline{qy} into 16 data segments of 8 bits each. Then, it performs an unsigned multiply-accumulate operation on each of the 16 sets of segments. The accumulated result is added to the value in the special register \lstinline{XACC}. The resulting sum is then saturated and stored in \lstinline{XACC}.\\

\textbf{Operation}
\begin{lstlisting}
  add0[15:0] = qx[  7:  0] * qy[  7:  0]
  add1[15:0] = qx[ 15:  8] * qy[ 15:  8]
  ...
  add15[15:0] = qx[127:120] * qy[127:120]
  sum[40:0] = XACC[39:0] + add0[15:0] + add1[15:0] + ... + add15[15:0]

  XACC[40:0] = min(max(sum[40:0], 0), 2^{40}-1)
\end{lstlisting}

\textbf{Schedule}\\
\begin{longtable}[c]{|l|l|l|}
    \hline
    \rowcolor{lightgray}
    \textbf{Operands} & \textbf{Use} & \textbf{Def} \\\hline
    \endfirsthead
    \hline
    \rowcolor{lightgray}
     \textbf{Operands} & \textbf{Use} & \textbf{Def} \\\hline
    \endhead
    \hline
    \endfoot

     qx & 1 & - \\\hline
     qy & 1 & - \\\hline
     XACC & 3 & 3


\end{longtable}

\textbf{Format}\\

\begin{figure*}[!htbp]
 {\setlength{\tabcolsep}{1pt}
 \begin{center}
\begin{tabular}{@{}F@{}F@{}T@{}I@{}E@{}O}
\instbitrange{31}{29}&
\instbitrange{28}{26}&
\instbitrange{25}{16}&
\instbit{15}&
\instbitrange{14}{7}&
\instbitrange{6}{0} \\
\hline
\multicolumn{1}{|c|}{{\scriptsize qx[2:0]}} &
\multicolumn{1}{c|}{{\scriptsize qy[2:0]}} &
\multicolumn{1}{c|}{1X101XX000} &
\multicolumn{1}{c|}{{\scriptsize sat[0]}} &
\multicolumn{1}{c|}{01XX0X0X} &
\multicolumn{1}{c|}{0011011}\\
\hline
3& 3& 10& 1& 8& 7 \\
\end{tabular}
 \end{center}
 }
 \vspace{-0.1in}
 \caption[]{ESP.VMULAS.U8.XACC} \end{figure*}


\newpage
\subsubsection{ESP.VMULAS.U8.XACC.LD.IP}\label{instruction:ESPVMULASU8XACCLDIP}
\textbf{Syntax}\\
ESP.VMULAS.U8.XACC.LD.IP qu,rs1,imm16,qx,qy,sat

\textbf{Description}\\
This instruction divides registers \lstinline{qx} and \lstinline{qy} into 16 data segments of 8 bits each. Then, it performs an unsigned multiply-accumulate operation on each of the 16 segment sets. The accumulated result is added to the value in the special register \lstinline{XACC}. The resulting sum is then saturated and stored in \lstinline{XACC}. \\ During the operation, the 16-byte data is loaded from the memory into the register \lstinline{qu}. After the access is completed, the value in the register \lstinline{rs1} is incremented by the 4-bit sign-extended constant in the instruction code segment, which is left-shifted by 4.\\

imm starts at -128, ends at 112, and steps by 16.

\textbf{Operation}
\begin{lstlisting}
  add0[15:0] = qx[  7:  0] * qy[  7:  0]
  add1[15:0] = qx[ 15:  8] * qy[ 15:  8]
  ...
  add15[15:0] = qx[127:120] * qy[127:120]
  sum[40:0] = XACC[39:0] + add0[15:0] + add1[15:0] + ... + add15[15:0]
  XACC[39:0] = min(max(sum[40:0], 0), 2^{40}-1)

  qu[127:0] = load128(rs1[31:0])
  rs1[31:0] = rs1[31:0] + imm
\end{lstlisting}

\textbf{Schedule}\\
\begin{longtable}[c]{|l|l|l|}
    \hline
    \rowcolor{lightgray}
    \textbf{Operands} & \textbf{Use} & \textbf{Def} \\\hline
    \endfirsthead
    \hline
    \rowcolor{lightgray}
     \textbf{Operands} & \textbf{Use} & \textbf{Def} \\\hline
    \endhead
    \hline
    \endfoot

    qu & - & 2 \\\hline
    rs1 & 1 & 1 \\\hline
     qx & 1 & - \\\hline
     qy & 1 & - \\\hline
     XACC & 3 & 3

\end{longtable}

\textbf{Format}\\

\begin{figure*}[!htbp]
 {\setlength{\tabcolsep}{1pt}
 \begin{center}
\begin{tabular}{@{}F@{}F@{}Y@{}I@{}W@{}I@{}F@{}W@{}F@{}W@{}I@{}O}
\instbitrange{31}{29}&
\instbitrange{28}{26}&
\instbitrange{25}{22}&
\instbit{21}&
\instbitrange{20}{19}&
\instbit{18}&
\instbitrange{17}{15}&
\instbitrange{14}{13}&
\instbitrange{12}{10}&
\instbitrange{9}{8}&
\instbit{7}&
\instbitrange{6}{0} \\
\hline
\multicolumn{1}{|c|}{{\scriptsize qx[2:0]}} &
\multicolumn{1}{c|}{{\scriptsize qy[2:0]}} &
\multicolumn{1}{c|}{0000} &
\multicolumn{1}{c|}{{\scriptsize sat[0]}} &
\multicolumn{1}{c|}{{\scriptsize imm16[3:2]}} &
\multicolumn{1}{c|}{{\scriptsize rs1[4]}} &
\multicolumn{1}{c|}{{\scriptsize rs1[2:0]}} &
\multicolumn{1}{c|}{11} &
\multicolumn{1}{c|}{{\scriptsize qu[2:0]}} &
\multicolumn{1}{c|}{{\scriptsize imm16[1:0]}} &
\multicolumn{1}{c|}{0} &
\multicolumn{1}{c|}{0011111}\\
\hline
3& 3& 4& 1& 2& 1& 3& 2& 3& 2& 1& 7 \\
\end{tabular}
 \end{center}
 }
 \vspace{-0.1in}
 \caption[]{ESP.VMULAS.U8.XACC.LD.IP} \end{figure*}


\newpage
\subsubsection{ESP.VMULAS.U8.XACC.LD.XP}\label{instruction:ESPVMULASU8XACCLDXP}
\textbf{Syntax}\\
ESP.VMULAS.U8.XACC.LD.XP qu,rs1,rs2,qx,qy,sat

\textbf{Description}\\
This instruction divides registers \lstinline{qx} and \lstinline{qy} into 16 data segments of 8 bits each. Then, it performs an unsigned multiply-accumulate operation on each of the 16 sets of segments. The accumulated result is added to the value in the special register \lstinline{XACC}. The resulting sum is then saturated and stored in \lstinline{XACC}.\\
During the operation, the 16-byte data is loaded from the memory into the register \lstinline{qu}. After the access is completed, the value in the register \lstinline{rs1} is incremented by the value in the register \lstinline{rs2}.\\
\textbf{Operation}
\begin{lstlisting}
  add0[15:0] = qx[  7:  0] * qy[  7:  0]
  add1[15:0] = qx[ 15:  8] * qy[ 15:  8]
  ...
  add15[15:0] = qx[127:120] * qy[127:120]
  sum[40:0] = XACC[39:0] + add0[15:0] + add1[15:0] + ... + add15[15:0]
  XACC[39:0] = min(max(sum[40:0], 0), 2^{40}-1)

  qu[127:0] = load128(rs1[31:0])
  rs1[31:0] = rs1[31:0] + rs2[31:0]
\end{lstlisting}

\textbf{Schedule}\\
\begin{longtable}[c]{|l|l|l|}
    \hline
    \rowcolor{lightgray}
    \textbf{Operands} & \textbf{Use} & \textbf{Def} \\\hline
    \endfirsthead
    \hline
    \rowcolor{lightgray}
     \textbf{Operands} & \textbf{Use} & \textbf{Def} \\\hline
    \endhead
    \hline
    \endfoot

    qu & - & 2 \\\hline
    rs1 & 1 & 1 \\\hline
    rs2 & 1 & - \\\hline
     qx & 1 & - \\\hline
     qy & 1 & - \\\hline
     XACC & 3 & 3

\end{longtable}

\textbf{Format}\\

\begin{figure*}[!htbp]
 {\setlength{\tabcolsep}{1pt}
 \begin{center}
\begin{tabular}{@{}F@{}F@{}W@{}I@{}F@{}I@{}I@{}F@{}W@{}F@{}I@{}I@{}I@{}O}
\instbitrange{31}{29}&
\instbitrange{28}{26}&
\instbitrange{25}{24}&
\instbit{23}&
\instbitrange{22}{20}&
\instbit{19}&
\instbit{18}&
\instbitrange{17}{15}&
\instbitrange{14}{13}&
\instbitrange{12}{10}&
\instbit{9}&
\instbit{8}&
\instbit{7}&
\instbitrange{6}{0} \\
\hline
\multicolumn{1}{|c|}{{\scriptsize qx[2:0]}} &
\multicolumn{1}{c|}{{\scriptsize qy[2:0]}} &
\multicolumn{1}{c|}{00} &
\multicolumn{1}{c|}{{\scriptsize rs2[4]}} &
\multicolumn{1}{c|}{{\scriptsize rs2[2:0]}} &
\multicolumn{1}{c|}{0} &
\multicolumn{1}{c|}{{\scriptsize rs1[4]}} &
\multicolumn{1}{c|}{{\scriptsize rs1[2:0]}} &
\multicolumn{1}{c|}{01} &
\multicolumn{1}{c|}{{\scriptsize qu[2:0]}} &
\multicolumn{1}{c|}{0} &
\multicolumn{1}{c|}{{\scriptsize sat[0]}} &
\multicolumn{1}{c|}{0} &
\multicolumn{1}{c|}{0011111}\\
\hline
3& 3& 2& 1& 3& 1& 1& 3& 2& 3& 1& 1& 1& 7 \\
\end{tabular}
 \end{center}
 }
 \vspace{-0.1in}
 \caption[]{ESP.VMULAS.U8.XACC.LD.XP} \end{figure*}


\newpage
\subsubsection{ESP.VMULAS.U8.XACC.ST.IP}\label{instruction:ESPVMULASU8XACCSTIP}
\textbf{Syntax}\\
ESP.VMULAS.U8.XACC.ST.IP qu,rs1,imm16,qx,qy,sat

\textbf{Description}\\
This instruction divides registers \lstinline{qx} and \lstinline{qy} into 16 data segments of 8 bits each. Then, it performs an unsigned multiply-accumulate operation on each of the 16 segment sets. The accumulated result is added to the value in the special register \lstinline{XACC}. The resulting sum is then saturated and stored in \lstinline{XACC}. \\ During the operation, the value of the register \lstinline{qu} is stored to memory. After the access is completed, the value in the register \lstinline{rs1} is incremented by a 4-bit sign-extended constant from the instruction code segment, which is left-shifted by 4.\\

imm starts at -128, ends at 112, and steps by 16.

\textbf{Operation}
\begin{lstlisting}
  add0[15:0] = qx[  7:  0] * qy[  7:  0]
  add1[15:0] = qx[ 15:  8] * qy[ 15:  8]
  ...
  add15[15:0] = qx[127:120] * qy[127:120]
  sum[40:0] = XACC[39:0] + add0[15:0] + add1[15:0] + ... + add15[15:0]
  XACC[39:0] = min(max(sum[40:0], 0), 2^{40}-1)

  qu[127:0] => store128(rs1[31:0])
  rs1[31:0] = rs1[31:0] + imm
\end{lstlisting}

\textbf{Schedule}\\
\begin{longtable}[c]{|l|l|l|}
    \hline
    \rowcolor{lightgray}
    \textbf{Operands} & \textbf{Use} & \textbf{Def} \\\hline
    \endfirsthead
    \hline
    \rowcolor{lightgray}
     \textbf{Operands} & \textbf{Use} & \textbf{Def} \\\hline
    \endhead
    \hline
    \endfoot

    qu & 2 & - \\\hline
    rs1 & 1 & 1 \\\hline
     qx & 1 & - \\\hline
     qy & 1 & - \\\hline
     XACC & 3 & 3

\end{longtable}

\textbf{Format}\\

\begin{figure*}[!htbp]
 {\setlength{\tabcolsep}{1pt}
 \begin{center}
\begin{tabular}{@{}F@{}F@{}Y@{}I@{}W@{}I@{}F@{}W@{}F@{}W@{}I@{}O}
\instbitrange{31}{29}&
\instbitrange{28}{26}&
\instbitrange{25}{22}&
\instbit{21}&
\instbitrange{20}{19}&
\instbit{18}&
\instbitrange{17}{15}&
\instbitrange{14}{13}&
\instbitrange{12}{10}&
\instbitrange{9}{8}&
\instbit{7}&
\instbitrange{6}{0} \\
\hline
\multicolumn{1}{|c|}{{\scriptsize qx[2:0]}} &
\multicolumn{1}{c|}{{\scriptsize qy[2:0]}} &
\multicolumn{1}{c|}{0100} &
\multicolumn{1}{c|}{{\scriptsize sat[0]}} &
\multicolumn{1}{c|}{{\scriptsize imm16[3:2]}} &
\multicolumn{1}{c|}{{\scriptsize rs1[4]}} &
\multicolumn{1}{c|}{{\scriptsize rs1[2:0]}} &
\multicolumn{1}{c|}{11} &
\multicolumn{1}{c|}{{\scriptsize qu[2:0]}} &
\multicolumn{1}{c|}{{\scriptsize imm16[1:0]}} &
\multicolumn{1}{c|}{0} &
\multicolumn{1}{c|}{0011111}\\
\hline
3& 3& 4& 1& 2& 1& 3& 2& 3& 2& 1& 7 \\
\end{tabular}
 \end{center}
 }
 \vspace{-0.1in}
 \caption[]{ESP.VMULAS.U8.XACC.ST.IP} \end{figure*}


\newpage
\subsubsection{ESP.VMULAS.U8.XACC.ST.XP}\label{instruction:ESPVMULASU8XACCSTXP}
\textbf{Syntax}\\
ESP.VMULAS.U8.XACC.ST.XP qu,rs1,rs2,qx,qy,sat

\textbf{Description}\\

This instruction divides registers \lstinline{qx} and \lstinline{qy} into 16 data segments of 8 bits each. Then, it performs an unsigned multiply-accumulate operation on each of the 16 segments. The accumulated result is added to the value in the special register \lstinline{XACC}. The resulting sum is then saturated and stored in \lstinline{XACC}. \\ During the operation, the value of the register \lstinline{qu} is stored to memory. After the access is completed, the value in the register \lstinline{rs1} is incremented by the value of the register \lstinline{rs2}.\\
\textbf{Operation}
\begin{lstlisting}
  add0[15:0] = qx[  7:  0] * qy[  7:  0]
  add1[15:0] = qx[ 15:  8] * qy[ 15:  8]
  ...
  add15[15:0] = qx[127:120] * qy[127:120]
  sum[40:0] = XACC[39:0] + add0[15:0] + add1[15:0] + ... + add15[15:0]
  XACC[39:0] = min(max(sum[40:0], 0), 2^{40}-1)

  qu[127:0] => store128(rs1[31:0])
  rs1[31:0] = rs1[31:0] + rs2[31:0]
\end{lstlisting}

\textbf{Schedule}\\
\begin{longtable}[c]{|l|l|l|}
    \hline
    \rowcolor{lightgray}
    \textbf{Operands} & \textbf{Use} & \textbf{Def} \\\hline
    \endfirsthead
    \hline
    \rowcolor{lightgray}
     \textbf{Operands} & \textbf{Use} & \textbf{Def} \\\hline
    \endhead
    \hline
    \endfoot

    qu & 2 & - \\\hline
    rs1 & 1 & 1 \\\hline
    rs2 & 1 & - \\\hline
     qx & 1 & - \\\hline
     qy & 1 & - \\\hline
     XACC & 3 & 3

\end{longtable}

\textbf{Format}\\

\begin{figure*}[!htbp]
 {\setlength{\tabcolsep}{1pt}
 \begin{center}
\begin{tabular}{@{}F@{}F@{}W@{}I@{}F@{}I@{}I@{}F@{}W@{}F@{}I@{}I@{}I@{}O}
\instbitrange{31}{29}&
\instbitrange{28}{26}&
\instbitrange{25}{24}&
\instbit{23}&
\instbitrange{22}{20}&
\instbit{19}&
\instbit{18}&
\instbitrange{17}{15}&
\instbitrange{14}{13}&
\instbitrange{12}{10}&
\instbit{9}&
\instbit{8}&
\instbit{7}&
\instbitrange{6}{0} \\
\hline
\multicolumn{1}{|c|}{{\scriptsize qx[2:0]}} &
\multicolumn{1}{c|}{{\scriptsize qy[2:0]}} &
\multicolumn{1}{c|}{01} &
\multicolumn{1}{c|}{{\scriptsize rs2[4]}} &
\multicolumn{1}{c|}{{\scriptsize rs2[2:0]}} &
\multicolumn{1}{c|}{0} &
\multicolumn{1}{c|}{{\scriptsize rs1[4]}} &
\multicolumn{1}{c|}{{\scriptsize rs1[2:0]}} &
\multicolumn{1}{c|}{01} &
\multicolumn{1}{c|}{{\scriptsize qu[2:0]}} &
\multicolumn{1}{c|}{0} &
\multicolumn{1}{c|}{{\scriptsize sat[0]}} &
\multicolumn{1}{c|}{0} &
\multicolumn{1}{c|}{0011111}\\
\hline
3& 3& 2& 1& 3& 1& 1& 3& 2& 3& 1& 1& 1& 7 \\
\end{tabular}
 \end{center}
 }
 \vspace{-0.1in}
 \caption[]{ESP.VMULAS.U8.XACC.ST.XP} \end{figure*}


\newpage
\subsubsection{ESP.VMULAS.S16.QACC.LDBC.INCP}\label{instruction:ESPVMULASS16QACCLDBCINCP}
\textbf{Syntax}\\
ESP.VMULAS.S16.QACC.LDBC.INCP qu,rs1,qx,qy,sat

\textbf{Description}\\
This instruction divides registers \lstinline{qx} and \lstinline{qy} into 8 data segments of 16 bits each. Then, the signed multiplication results of the 8 sets of segments are added to the corresponding 64-bit data segments in special registers \lstinline{QACC_H} and \lstinline{QACC_L} respectively. The calculated result is saturated to a 64-bit signed number and then stored in the corresponding 64-bit data segment in \lstinline{QACC_H} and \lstinline{QACC_L}.\\
During the operation, the 2-byte data is loaded from the memory and broadcast to the register \lstinline{qu}. After the access is completed, the value in the register \lstinline{rs1} is incremented by 2.\\
\textbf{Operation}
\begin{lstlisting}
  QACC_L[ 63:  0] = min(max(QACC_L[ 63:  0] + qx[ 15:  0] * qy[ 15:  0], -2^{63}), 2^{63}-1)
  QACC_L[127: 64] = min(max(QACC_L[127: 64] + qx[ 31: 16] * qy[ 31: 16], -2^{63}), 2^{63}-1)
  QACC_L[191:128] = min(max(QACC_L[191:128] + qx[ 47: 32] * qy[ 47: 32], -2^{63}), 2^{63}-1)
  QACC_L[255:192] = min(max(QACC_L[255:192] + qx[ 63: 48] * qy[ 63: 48], -2^{63}), 2^{63}-1)
  QACC_H[ 63:  0] = min(max(QACC_H[ 63:  0] + qx[ 79: 64] * qy[ 79: 64], -2^{63}), 2^{63}-1)
  QACC_H[127: 64] = min(max(QACC_H[127: 64] + qx[ 95: 80] * qy[ 95: 80], -2^{63}), 2^{63}-1)
  QACC_H[191:128] = min(max(QACC_H[191:128] + qx[111: 96] * qy[111: 96], -2^{63}), 2^{63}-1)
  QACC_H[255:192] = min(max(QACC_H[255:192] + qx[127:112] * qy[127:112], -2^{63}), 2^{63}-1)

  qu[127:0] = {8{load16(rs1[31:0])}}
  rs1[31:0] = rs1[31:0] + 2
\end{lstlisting}

\textbf{Schedule}\\
\begin{longtable}[c]{|l|l|l|}
    \hline
    \rowcolor{lightgray}
    \textbf{Operands} & \textbf{Use} & \textbf{Def} \\\hline
    \endfirsthead
    \hline
    \rowcolor{lightgray}
     \textbf{Operands} & \textbf{Use} & \textbf{Def} \\\hline
    \endhead
    \hline
    \endfoot

    qu & - & 2 \\\hline
    rs1 & 1 & 1 \\\hline
     qx & 1 & - \\\hline
     qy & 1 & - \\\hline
     QACC\_H & 2 & 2 \\\hline
     QACC\_L & 2 & 2

\end{longtable}

\textbf{Format}\\

\begin{figure*}[!htbp]
 {\setlength{\tabcolsep}{1pt}
 \begin{center}
\begin{tabular}{@{}F@{}F@{}Y@{}I@{}W@{}I@{}F@{}W@{}F@{}F@{}O}
\instbitrange{31}{29}&
\instbitrange{28}{26}&
\instbitrange{25}{22}&
\instbit{21}&
\instbitrange{20}{19}&
\instbit{18}&
\instbitrange{17}{15}&
\instbitrange{14}{13}&
\instbitrange{12}{10}&
\instbitrange{9}{7}&
\instbitrange{6}{0} \\
\hline
\multicolumn{1}{|c|}{{\scriptsize qx[2:0]}} &
\multicolumn{1}{c|}{{\scriptsize qy[2:0]}} &
\multicolumn{1}{c|}{XX11} &
\multicolumn{1}{c|}{{\scriptsize sat[0]}} &
\multicolumn{1}{c|}{XX} &
\multicolumn{1}{c|}{{\scriptsize rs1[4]}} &
\multicolumn{1}{c|}{{\scriptsize rs1[2:0]}} &
\multicolumn{1}{c|}{11} &
\multicolumn{1}{c|}{{\scriptsize qu[2:0]}} &
\multicolumn{1}{c|}{X01} &
\multicolumn{1}{c|}{0011111}\\
\hline
3& 3& 4& 1& 2& 1& 3& 2& 3& 3& 7 \\
\end{tabular}
 \end{center}
 }
 \vspace{-0.1in}
 \caption[]{ESP.VMULAS.S16.QACC.LDBC.INCP} \end{figure*}


\newpage
\subsubsection{ESP.VMULAS.S8.QACC.LDBC.INCP}\label{instruction:ESPVMULASS8QACCLDBCINCP}
\textbf{Syntax}\\
ESP.VMULAS.S8.QACC.LDBC.INCP qu,rs1,qx,qy,sat

\textbf{Description}\\
This instruction divides registers \lstinline{qx} and \lstinline{qy} into 16 data segments of 8 bits each. Then, the signed multiplication results of the 16 sets of segments are added to the corresponding 32-bit data segments in special registers \lstinline{QACC_H} and \lstinline{QACC_L} respectively. The calculated result is saturated to a 32-bit signed number and then stored in the corresponding 32-bit data segment in \lstinline{QACC_H} and \lstinline{QACC_L}.\\
During the operation, the 1-byte data is loaded from the memory and broadcast to the register \lstinline{qu}. After the access is completed, the value in the register \lstinline{rs1} is incremented by 1.\\
\textbf{Operation}
\begin{lstlisting}
  QACC_L[ 31:  0] = min(max(QACC_L[ 31:  0] + qx[  7:  0] * qy[  7:  0], -2^{31}), 2^{31}-1)
  QACC_L[ 63: 32] = min(max(QACC_L[ 63: 32] + qx[ 15:  8] * qy[ 15:  8], -2^{31}), 2^{31}-1)
  QACC_L[ 95: 64] = min(max(QACC_L[ 95: 64] + qx[ 23: 16] * qy[ 23: 16], -2^{31}), 2^{31}-1)
  QACC_L[127: 96] = min(max(QACC_L[127: 96] + qx[ 31: 24] * qy[ 31: 24], -2^{31}), 2^{31}-1)
  QACC_L[159:128] = min(max(QACC_L[159:128] + qx[ 39: 32] * qy[ 39: 32], -2^{31}), 2^{31}-1)
  QACC_L[191:160] = min(max(QACC_L[191:160] + qx[ 47: 40] * qy[ 47: 40], -2^{31}), 2^{31}-1)
  QACC_L[223:192] = min(max(QACC_L[223:192] + qx[ 55: 48] * qy[ 55: 48], -2^{31}), 2^{31}-1)
  QACC_L[255:224] = min(max(QACC_L[255:224] + qx[ 63: 56] * qy[ 63: 56], -2^{31}), 2^{31}-1)
  QACC_H[ 31:  0] = min(max(QACC_H[ 31:  0] + qx[ 71: 64] * qy[ 71: 64], -2^{31}), 2^{31}-1)
  QACC_H[ 63: 32] = min(max(QACC_H[ 63: 32] + qx[ 79: 72] * qy[ 79: 72], -2^{31}), 2^{31}-1)
  QACC_H[ 95: 64] = min(max(QACC_H[ 95: 64] + qx[ 87: 80] * qy[ 87: 80], -2^{31}), 2^{31}-1)
  QACC_H[127: 96] = min(max(QACC_H[127: 96] + qx[ 95: 88] * qy[ 95: 88], -2^{31}), 2^{31}-1)
  QACC_H[159:128] = min(max(QACC_H[159:128] + qx[103: 96] * qy[103: 96], -2^{31}), 2^{31}-1)
  QACC_H[191:160] = min(max(QACC_H[191:160] + qx[111:104] * qy[111:104], -2^{31}), 2^{31}-1)
  QACC_H[223:192] = min(max(QACC_H[223:192] + qx[119:112] * qy[119:112], -2^{31}), 2^{31}-1)
  QACC_H[255:224] = min(max(QACC_H[255:224] + qx[127:120] * qy[127:120], -2^{31}), 2^{31}-1)


  qu[127:0] = {16{load8(rs1[31:0])}}
  rs1[31:0] = rs1[31:0] + 1
\end{lstlisting}

\textbf{Schedule}\\
\begin{longtable}[c]{|l|l|l|}
    \hline
    \rowcolor{lightgray}
    \textbf{Operands} & \textbf{Use} & \textbf{Def} \\\hline
    \endfirsthead
    \hline
    \rowcolor{lightgray}
     \textbf{Operands} & \textbf{Use} & \textbf{Def} \\\hline
    \endhead
    \hline
    \endfoot

    qu & - & 2 \\\hline
    rs1 & 1 & 1 \\\hline
     qx & 1 & - \\\hline
     qy & 1 & - \\\hline
     QACC\_H & 2 & 2 \\\hline
     QACC\_L & 2 & 2

\end{longtable}

\textbf{Format}\\

\begin{figure*}[!htbp]
 {\setlength{\tabcolsep}{1pt}
 \begin{center}
\begin{tabular}{@{}F@{}F@{}Y@{}I@{}W@{}I@{}F@{}W@{}F@{}F@{}O}
\instbitrange{31}{29}&
\instbitrange{28}{26}&
\instbitrange{25}{22}&
\instbit{21}&
\instbitrange{20}{19}&
\instbit{18}&
\instbitrange{17}{15}&
\instbitrange{14}{13}&
\instbitrange{12}{10}&
\instbitrange{9}{7}&
\instbitrange{6}{0} \\
\hline
\multicolumn{1}{|c|}{{\scriptsize qx[2:0]}} &
\multicolumn{1}{c|}{{\scriptsize qy[2:0]}} &
\multicolumn{1}{c|}{XX01} &
\multicolumn{1}{c|}{{\scriptsize sat[0]}} &
\multicolumn{1}{c|}{XX} &
\multicolumn{1}{c|}{{\scriptsize rs1[4]}} &
\multicolumn{1}{c|}{{\scriptsize rs1[2:0]}} &
\multicolumn{1}{c|}{11} &
\multicolumn{1}{c|}{{\scriptsize qu[2:0]}} &
\multicolumn{1}{c|}{X01} &
\multicolumn{1}{c|}{0011111}\\
\hline
3& 3& 4& 1& 2& 1& 3& 2& 3& 3& 7 \\
\end{tabular}
 \end{center}
 }
 \vspace{-0.1in}
 \caption[]{ESP.VMULAS.S8.QACC.LDBC.INCP} \end{figure*}


\newpage
\subsubsection{ESP.VMULAS.U16.QACC.LDBC.INCP}\label{instruction:ESPVMULASU16QACCLDBCINCP}
\textbf{Syntax}\\
ESP.VMULAS.U16.QACC.LDBC.INCP qu,rs1,qx,qy,sat

\textbf{Description}\\
This instruction divides registers \lstinline{qx} and \lstinline{qy} into 8 data segments of 16 bits each. Then, it performs the unsigned multiplication of 8 sets of segments and adds the result to the corresponding 64-bit data segment in special registers \lstinline{QACC_H} and \lstinline{QACC_L} respectively. The calculated result is saturated to a 64-bit unsigned number and then stored in the corresponding 64-bit data segment in \lstinline{QACC_H} and \lstinline{QACC_L}.\\
During the operation, it loads the 2-byte data from the memory and broadcasts it to the register \lstinline{qu}. After the access is completed, the value in the register \lstinline{rs1} is incremented by 2.\\
\textbf{Operation}
\begin{lstlisting}
  QACC_L[ 63:  0] = min(QACC_L[ 63:  0] + qx[ 15:  0] * qy[ 15:  0], 2^{64}-1)
  QACC_L[127: 64] = min(QACC_L[127: 64] + qx[ 31: 16] * qy[ 31: 16], 2^{64}-1)
  QACC_L[191:128] = min(QACC_L[191:128] + qx[ 47: 32] * qy[ 47: 32], 2^{64}-1)
  QACC_L[255:192] = min(QACC_L[255:192] + qx[ 63: 48] * qy[ 63: 48], 2^{64}-1)
  QACC_H[ 63:  0] = min(QACC_H[ 63:  0] + qx[ 79: 64] * qy[ 79: 64], 2^{64}-1)
  QACC_H[127: 64] = min(QACC_H[127: 64] + qx[ 95: 80] * qy[ 95: 80], 2^{64}-1)
  QACC_H[191:128] = min(QACC_H[191:128] + qx[111: 96] * qy[111: 96], 2^{64}-1)
  QACC_H[255:192] = min(QACC_H[255:192] + qx[127:112] * qy[127:112], 2^{64}-1)

  qu[127:0] = {8{load16(rs1[31:0])}}
  rs1[31:0] = rs1[31:0] + 2
\end{lstlisting}

\textbf{Schedule}\\
\begin{longtable}[c]{|l|l|l|}
    \hline
    \rowcolor{lightgray}
    \textbf{Operands} & \textbf{Use} & \textbf{Def} \\\hline
    \endfirsthead
    \hline
    \rowcolor{lightgray}
     \textbf{Operands} & \textbf{Use} & \textbf{Def} \\\hline
    \endhead
    \hline
    \endfoot

    qu & - & 2 \\\hline
    rs1 & 1 & 1 \\\hline
     qx & 1 & - \\\hline
     qy & 1 & - \\\hline
     QACC\_H & 2 & 2 \\\hline
     QACC\_L & 2 & 2

\end{longtable}

\textbf{Format}\\

\begin{figure*}[!htbp]
 {\setlength{\tabcolsep}{1pt}
 \begin{center}
\begin{tabular}{@{}F@{}F@{}Y@{}I@{}W@{}I@{}F@{}W@{}F@{}F@{}O}
\instbitrange{31}{29}&
\instbitrange{28}{26}&
\instbitrange{25}{22}&
\instbit{21}&
\instbitrange{20}{19}&
\instbit{18}&
\instbitrange{17}{15}&
\instbitrange{14}{13}&
\instbitrange{12}{10}&
\instbitrange{9}{7}&
\instbitrange{6}{0} \\
\hline
\multicolumn{1}{|c|}{{\scriptsize qx[2:0]}} &
\multicolumn{1}{c|}{{\scriptsize qy[2:0]}} &
\multicolumn{1}{c|}{XX10} &
\multicolumn{1}{c|}{{\scriptsize sat[0]}} &
\multicolumn{1}{c|}{XX} &
\multicolumn{1}{c|}{{\scriptsize rs1[4]}} &
\multicolumn{1}{c|}{{\scriptsize rs1[2:0]}} &
\multicolumn{1}{c|}{11} &
\multicolumn{1}{c|}{{\scriptsize qu[2:0]}} &
\multicolumn{1}{c|}{X01} &
\multicolumn{1}{c|}{0011111}\\
\hline
3& 3& 4& 1& 2& 1& 3& 2& 3& 3& 7 \\
\end{tabular}
 \end{center}
 }
 \vspace{-0.1in}
 \caption[]{ESP.VMULAS.U16.QACC.LDBC.INCP} \end{figure*}


\newpage
\subsubsection{ESP.VMULAS.U8.QACC.LDBC.INCP}\label{instruction:ESPVMULASU8QACCLDBCINCP}
\textbf{Syntax}\\
ESP.VMULAS.U8.QACC.LDBC.INCP qu,rs1,qx,qy,sat

\textbf{Description}\\
This instruction divides registers \lstinline{qx} and \lstinline{qy} into 16 data segments of 8 bits each. Then, the unsigned multiplication results of the 16 sets of segments are added to the corresponding 32-bit data segments in special registers \lstinline{QACC_H} and \lstinline{QACC_L} respectively. The calculated result is saturated to a 32-bit unsigned number and then stored in the corresponding 32-bit data segment in \lstinline{QACC_H} and \lstinline{QACC_L}.\\
During the operation, the 1-byte data is loaded from the memory and broadcast to the register \lstinline{qu}. After the access is completed, the value in the register \lstinline{rs1} is incremented by 1.\\
\textbf{Operation}
\begin{lstlisting}
  QACC_L[ 31:  0] = min(QACC_L[ 31:  0] + qx[  7:  0] * qy[  7:  0], 2^{32}-1)
  QACC_L[ 63: 32] = min(QACC_L[ 63: 32] + qx[ 15:  8] * qy[ 15:  8], 2^{32}-1)
  QACC_L[ 95: 64] = min(QACC_L[ 95: 64] + qx[ 23: 16] * qy[ 23: 16], 2^{32}-1)
  QACC_L[127: 96] = min(QACC_L[127: 96] + qx[ 31: 24] * qy[ 31: 24], 2^{32}-1)
  QACC_L[159:128] = min(QACC_L[159:128] + qx[ 39: 32] * qy[ 39: 32], 2^{32}-1)
  QACC_L[191:160] = min(QACC_L[191:160] + qx[ 47: 40] * qy[ 47: 40], 2^{32}-1)
  QACC_L[223:192] = min(QACC_L[223:192] + qx[ 55: 48] * qy[ 55: 48], 2^{32}-1)
  QACC_L[255:224] = min(QACC_L[255:224] + qx[ 63: 56] * qy[ 63: 56], 2^{32}-1)
  QACC_H[ 31:  0] = min(QACC_H[ 31:  0] + qx[ 71: 64] * qy[ 71: 64], 2^{32}-1)
  QACC_H[ 63: 32] = min(QACC_H[ 63: 32] + qx[ 79: 72] * qy[ 79: 72], 2^{32}-1)
  QACC_H[ 95: 64] = min(QACC_H[ 95: 64] + qx[ 87: 80] * qy[ 87: 80], 2^{32}-1)
  QACC_H[127: 96] = min(QACC_H[127: 96] + qx[ 95: 88] * qy[ 95: 88], 2^{32}-1)
  QACC_H[159:128] = min(QACC_H[159:128] + qx[103: 96] * qy[103: 96], 2^{32}-1)
  QACC_H[191:160] = min(QACC_H[191:160] + qx[111:104] * qy[111:104], 2^{32}-1)
  QACC_H[223:192] = min(QACC_H[223:192] + qx[119:112] * qy[119:112], 2^{32}-1)
  QACC_H[255:224] = min(QACC_H[255:224] + qx[127:120] * qy[127:120], 2^{32}-1)


  qu[127:0] = {16{load8(rs1[31:0])}}
  rs1[31:0] = rs1[31:0] + 1
\end{lstlisting}

\textbf{Schedule}\\
\begin{longtable}[c]{|l|l|l|}
    \hline
    \rowcolor{lightgray}
    \textbf{Operands} & \textbf{Use} & \textbf{Def} \\\hline
    \endfirsthead
    \hline
    \rowcolor{lightgray}
     \textbf{Operands} & \textbf{Use} & \textbf{Def} \\\hline
    \endhead
    \hline
    \endfoot

    qu & - & 2 \\\hline
    rs1 & 1 & 1 \\\hline
     qx & 1 & - \\\hline
     qy & 1 & - \\\hline
     QACC\_H & 2 & 2 \\\hline
     QACC\_L & 2 & 2

\end{longtable}

\textbf{Format}\\

\begin{figure*}[!htbp]
 {\setlength{\tabcolsep}{1pt}
 \begin{center}
\begin{tabular}{@{}F@{}F@{}Y@{}I@{}W@{}I@{}F@{}W@{}F@{}F@{}O}
\instbitrange{31}{29}&
\instbitrange{28}{26}&
\instbitrange{25}{22}&
\instbit{21}&
\instbitrange{20}{19}&
\instbit{18}&
\instbitrange{17}{15}&
\instbitrange{14}{13}&
\instbitrange{12}{10}&
\instbitrange{9}{7}&
\instbitrange{6}{0} \\
\hline
\multicolumn{1}{|c|}{{\scriptsize qx[2:0]}} &
\multicolumn{1}{c|}{{\scriptsize qy[2:0]}} &
\multicolumn{1}{c|}{XX00} &
\multicolumn{1}{c|}{{\scriptsize sat[0]}} &
\multicolumn{1}{c|}{XX} &
\multicolumn{1}{c|}{{\scriptsize rs1[4]}} &
\multicolumn{1}{c|}{{\scriptsize rs1[2:0]}} &
\multicolumn{1}{c|}{11} &
\multicolumn{1}{c|}{{\scriptsize qu[2:0]}} &
\multicolumn{1}{c|}{X01} &
\multicolumn{1}{c|}{0011111}\\
\hline
3& 3& 4& 1& 2& 1& 3& 2& 3& 3& 7 \\
\end{tabular}
 \end{center}
 }
 \vspace{-0.1in}
 \caption[]{ESP.VMULAS.U8.QACC.LDBC.INCP} \end{figure*}


\newpage
\subsubsection{ESP.VSMULAS.S16.QACC}\label{instruction:ESPVSMULASS16QACC}
\textbf{Syntax}\\
ESP.VSMULAS.S16.QACC qx,qy,sel16,sat

\textbf{Description}\\
This instruction selects one out of the eight 16-bit data segments in the register \lstinline{qy} according to the immediate number \lstinline{sel16} and performs a signed multiplication on it and the eight 16-bit data segments in the register \lstinline{qx} respectively. The eight results obtained are added to the corresponding 64-bit data segments in the special registers \lstinline{QACC_H} and \lstinline{QACC_L} respectively. Then, each result is saturated to a 64-bit signed number and then stored in the corresponding 64-bit data segment in \lstinline{QACC_H} and \lstinline{QACC_L}.\\
\textbf{Operation}
\begin{lstlisting}
  temp[15:0] = qy[sel16*16+15:sel16*16]
  QACC_L[ 63:  0] = min(max(QACC_L[ 63:  0] + qx[ 15:  0] * temp[15:0], -2^{63}), 2^{63}-1)
  QACC_L[127: 64] = min(max(QACC_L[127: 64] + qx[ 31: 16] * temp[15:0], -2^{63}), 2^{63}-1)
  QACC_L[191:128] = min(max(QACC_L[191:128] + qx[ 47: 32] * temp[15:0], -2^{63}), 2^{63}-1)
  QACC_L[255:192] = min(max(QACC_L[255:192] + qx[ 63: 48] * temp[15:0], -2^{63}), 2^{63}-1)
  QACC_H[ 63:  0] = min(max(QACC_H[ 63:  0] + qx[ 79: 64] * temp[15:0], -2^{63}), 2^{63}-1)
  QACC_H[127: 64] = min(max(QACC_H[127: 64] + qx[ 95: 80] * temp[15:0], -2^{63}), 2^{63}-1)
  QACC_H[191:128] = min(max(QACC_H[191:128] + qx[111: 96] * temp[15:0], -2^{63}), 2^{63}-1)
  QACC_H[255:192] = min(max(QACC_H[255:192] + qx[127:112] * temp[15:0], -2^{63}), 2^{63}-1)
\end{lstlisting}

\textbf{Schedule}\\
\begin{longtable}[c]{|l|l|l|}
    \hline
    \rowcolor{lightgray}
    \textbf{Operands} & \textbf{Use} & \textbf{Def} \\\hline
    \endfirsthead
    \hline
    \rowcolor{lightgray}
     \textbf{Operands} & \textbf{Use} & \textbf{Def} \\\hline
    \endhead
    \hline
    \endfoot

     qx & 1 & - \\\hline
     qy & 1 & - \\\hline
     QACC\_H & 2 & 2 \\\hline
     QACC\_L & 2 & 2

\end{longtable}

\textbf{Format}\\

\begin{figure*}[!htbp]
 {\setlength{\tabcolsep}{1pt}
 \begin{center}
\begin{tabular}{@{}F@{}F@{}V@{}I@{}I@{}Y@{}E@{}O}
\instbitrange{31}{29}&
\instbitrange{28}{26}&
\instbitrange{25}{21}&
\instbit{20}&
\instbit{19}&
\instbitrange{18}{15}&
\instbitrange{14}{7}&
\instbitrange{6}{0} \\
\hline
\multicolumn{1}{|c|}{{\scriptsize qx[2:0]}} &
\multicolumn{1}{c|}{{\scriptsize qy[2:0]}} &
\multicolumn{1}{c|}{1X111} &
\multicolumn{1}{c|}{{\scriptsize sat[0]}} &
\multicolumn{1}{c|}{X} &
\multicolumn{1}{c|}{{\scriptsize sel16[3:0]}} &
\multicolumn{1}{c|}{01XX1X0X} &
\multicolumn{1}{c|}{0011011}\\
\hline
3& 3& 5& 1& 1& 4& 8& 7 \\
\end{tabular}
 \end{center}
 }
 \vspace{-0.1in}
 \caption[]{ESP.VSMULAS.S16.QACC} \end{figure*}


\newpage
\subsubsection{ESP.VSMULAS.S16.QACC.LD.INCP}\label{instruction:ESPVSMULASS16QACCLDINCP}
\textbf{Syntax}\\
ESP.VSMULAS.S16.QACC.LD.INCP qu,rs1,qx,qy,sel16,sat

\textbf{Description}\\
This instruction selects one out of the eight 16-bit data segments in the register \lstinline{qy} according to the immediate number \lstinline{sel16} and performs a signed multiplication on it and the eight 16-bit data segments in the register \lstinline{qx} respectively. The eight results obtained are added to the corresponding 64-bit data segments in the special registers \lstinline{QACC_H} and \lstinline{QACC_L} respectively. Then, each result is saturated to a 64-bit signed number and then stored in the corresponding 64-bit data segment in \lstinline{QACC_H} and \lstinline{QACC_L}.\\
During the operation, the 16-byte data is loaded from the memory to the register \lstinline{qu}. After the access is completed, the value in the register \lstinline{rs1} is incremented by 16.\\
\textbf{Operation}
\begin{lstlisting}
  temp[15:0] = qy[sel16*16+15:sel16*16]
  QACC_L[ 63:  0] = min(max(QACC_L[ 63:  0] + qx[ 15:  0] * temp[15:0], -2^{63}), 2^{63}-1)
  QACC_L[127: 64] = min(max(QACC_L[127: 64] + qx[ 31: 16] * temp[15:0], -2^{63}), 2^{63}-1)
  QACC_L[191:128] = min(max(QACC_L[191:128] + qx[ 47: 32] * temp[15:0], -2^{63}), 2^{63}-1)
  QACC_L[255:192] = min(max(QACC_L[255:192] + qx[ 63: 48] * temp[15:0], -2^{63}), 2^{63}-1)
  QACC_H[ 63:  0] = min(max(QACC_H[ 63:  0] + qx[ 79: 64] * temp[15:0], -2^{63}), 2^{63}-1)
  QACC_H[127: 64] = min(max(QACC_H[127: 64] + qx[ 95: 80] * temp[15:0], -2^{63}), 2^{63}-1)
  QACC_H[191:128] = min(max(QACC_H[191:128] + qx[111: 96] * temp[15:0], -2^{63}), 2^{63}-1)
  QACC_H[255:192] = min(max(QACC_H[255:192] + qx[127:112] * temp[15:0], -2^{63}), 2^{63}-1)

  qu[127:0] = load128(rs1[31:0])
  rs1[31:0] = rs1[31:0] + 16
\end{lstlisting}

\textbf{Schedule}\\
\begin{longtable}[c]{|l|l|l|}
    \hline
    \rowcolor{lightgray}
    \textbf{Operands} & \textbf{Use} & \textbf{Def} \\\hline
    \endfirsthead
    \hline
    \rowcolor{lightgray}
     \textbf{Operands} & \textbf{Use} & \textbf{Def} \\\hline
    \endhead
    \hline
    \endfoot

    qu & - & 2 \\\hline
    rs1 & 1 & 1 \\\hline
     qx & 1 & - \\\hline
     qy & 1 & - \\\hline
     QACC\_H & 2 & 2 \\\hline
     QACC\_L & 2 & 2

\end{longtable}

\textbf{Format}\\

\begin{figure*}[!htbp]
 {\setlength{\tabcolsep}{1pt}
 \begin{center}
\begin{tabular}{@{}F@{}F@{}W@{}I@{}Y@{}I@{}F@{}W@{}F@{}F@{}O}
\instbitrange{31}{29}&
\instbitrange{28}{26}&
\instbitrange{25}{24}&
\instbit{23}&
\instbitrange{22}{19}&
\instbit{18}&
\instbitrange{17}{15}&
\instbitrange{14}{13}&
\instbitrange{12}{10}&
\instbitrange{9}{7}&
\instbitrange{6}{0} \\
\hline
\multicolumn{1}{|c|}{{\scriptsize qx[2:0]}} &
\multicolumn{1}{c|}{{\scriptsize qy[2:0]}} &
\multicolumn{1}{c|}{11} &
\multicolumn{1}{c|}{{\scriptsize sat[0]}} &
\multicolumn{1}{c|}{{\scriptsize sel16[3:0]}} &
\multicolumn{1}{c|}{{\scriptsize rs1[4]}} &
\multicolumn{1}{c|}{{\scriptsize rs1[2:0]}} &
\multicolumn{1}{c|}{11} &
\multicolumn{1}{c|}{{\scriptsize qu[2:0]}} &
\multicolumn{1}{c|}{111} &
\multicolumn{1}{c|}{0011111}\\
\hline
3& 3& 2& 1& 4& 1& 3& 2& 3& 3& 7 \\
\end{tabular}
 \end{center}
 }
 \vspace{-0.1in}
 \caption[]{ESP.VSMULAS.S16.QACC.LD.INCP} \end{figure*}


\newpage
\subsubsection{ESP.VSMULAS.S8.QACC}\label{instruction:ESPVSMULASS8QACC}
\textbf{Syntax}\\
ESP.VSMULAS.S8.QACC qx,qy,sel16,sat

\textbf{Description}\\
This instruction selects one out of the 16 8-bit data segments in the register \lstinline{qy} according to the immediate number \lstinline{sel16} and performs a signed multiplication on it and the 16 8-bit data segments in the register \lstinline{qx} respectively. The 16 results obtained are added to the corresponding 32-bit data segments in the special registers \lstinline{QACC_H} and \lstinline{QACC_L} respectively. Then, each result is saturated to a 32-bit signed number and then stored in the corresponding 32-bit data segment in \lstinline{QACC_H} and \lstinline{QACC_L}.\\
\textbf{Operation}
\begin{lstlisting}
  temp[7:0] = qy[sel16*8+7:sel16*8]
  QACC_L[ 31:  0] = min(max(QACC_L[ 31:  0] + qx[  7:  0] * temp[7:0], -2^{31}), 2^{31}-1)
  QACC_L[ 63: 32] = min(max(QACC_L[ 63: 32] + qx[ 15:  8] * temp[7:0], -2^{31}), 2^{31}-1)
  QACC_L[ 95: 64] = min(max(QACC_L[ 95: 64] + qx[ 23: 16] * temp[7:0], -2^{31}), 2^{31}-1)
  QACC_L[127: 96] = min(max(QACC_L[128: 96] + qx[ 31: 24] * temp[7:0], -2^{31}), 2^{31}-1)
  QACC_L[159:128] = min(max(QACC_L[159:128] + qx[ 39: 32] * temp[7:0], -2^{31}), 2^{31}-1)
  QACC_L[191:160] = min(max(QACC_L[191:160] + qx[ 47: 40] * temp[7:0], -2^{31}), 2^{31}-1)
  QACC_L[223:192] = min(max(QACC_L[223:192] + qx[ 55: 48] * temp[7:0], -2^{31}), 2^{31}-1)
  QACC_L[255:224] = min(max(QACC_L[255:224] + qx[ 63: 56] * temp[7:0], -2^{31}), 2^{31}-1)
  QACC_H[ 31:  0] = min(max(QACC_H[ 31:  0] + qx[ 71: 64] * temp[7:0], -2^{31}), 2^{31}-1)
  QACC_H[ 63: 32] = min(max(QACC_H[ 63: 32] + qx[ 79: 72] * temp[7:0], -2^{31}), 2^{31}-1)
  QACC_H[ 95: 64] = min(max(QACC_H[ 95: 64] + qx[ 87: 80] * temp[7:0], -2^{31}), 2^{31}-1)
  QACC_H[127: 96] = min(max(QACC_H[127: 96] + qx[ 95: 88] * temp[7:0], -2^{31}), 2^{31}-1)
  QACC_H[159:128] = min(max(QACC_H[159:128] + qx[103: 96] * temp[7:0], -2^{31}), 2^{31}-1)
  QACC_H[191:160] = min(max(QACC_H[191:160] + qx[111:104] * temp[7:0], -2^{31}), 2^{31}-1)
  QACC_H[223:192] = min(max(QACC_H[223:192] + qx[119:112] * temp[7:0], -2^{31}), 2^{31}-1)
  QACC_H[255:224] = min(max(QACC_H[255:224] + qx[127:120] * temp[7:0], -2^{31}), 2^{31}-1)

\end{lstlisting}

\textbf{Schedule}\\
\begin{longtable}[c]{|l|l|l|}
    \hline
    \rowcolor{lightgray}
    \textbf{Operands} & \textbf{Use} & \textbf{Def} \\\hline
    \endfirsthead
    \hline
    \rowcolor{lightgray}
     \textbf{Operands} & \textbf{Use} & \textbf{Def} \\\hline
    \endhead
    \hline
    \endfoot

     qx & 1 & - \\\hline
     qy & 1 & - \\\hline
     QACC\_H & 2 & 2 \\\hline
     QACC\_L & 2 & 2

\end{longtable}

\textbf{Format}\\

\begin{figure*}[!htbp]
 {\setlength{\tabcolsep}{1pt}
 \begin{center}
\begin{tabular}{@{}F@{}F@{}V@{}I@{}I@{}Y@{}E@{}O}
\instbitrange{31}{29}&
\instbitrange{28}{26}&
\instbitrange{25}{21}&
\instbit{20}&
\instbit{19}&
\instbitrange{18}{15}&
\instbitrange{14}{7}&
\instbitrange{6}{0} \\
\hline
\multicolumn{1}{|c|}{{\scriptsize qx[2:0]}} &
\multicolumn{1}{c|}{{\scriptsize qy[2:0]}} &
\multicolumn{1}{c|}{1X101} &
\multicolumn{1}{c|}{{\scriptsize sat[0]}} &
\multicolumn{1}{c|}{X} &
\multicolumn{1}{c|}{{\scriptsize sel16[3:0]}} &
\multicolumn{1}{c|}{01XX1X0X} &
\multicolumn{1}{c|}{0011011}\\
\hline
3& 3& 5& 1& 1& 4& 8& 7 \\
\end{tabular}
 \end{center}
 }
 \vspace{-0.1in}
 \caption[]{ESP.VSMULAS.S8.QACC} \end{figure*}


\newpage
\subsubsection{ESP.VSMULAS.S8.QACC.LD.INCP}\label{instruction:ESPVSMULASS8QACCLDINCP}
\textbf{Syntax}\\
ESP.VSMULAS.S8.QACC.LD.INCP qu,rs1,qx,qy,sel16,sat

\textbf{Description}\\
This instruction selects one out of the 16 8-bit data segments in the register \lstinline{qy} according to the immediate number \lstinline{sel16} and performs a signed multiplication on it and the 16 8-bit data segments in the register \lstinline{qx} respectively. The 16 results obtained are added to the corresponding 32-bit data segments in the special registers \lstinline{QACC_H} and \lstinline{QACC_L} respectively. Then, the result is saturated to a 32-bit signed number and stored in the corresponding 32-bit data segment in \lstinline{QACC_H} and \lstinline{QACC_L}.\\
During the operation, the 16-byte data is loaded from the memory into the register \lstinline{qu}. After the access is completed, the value in the register \lstinline{rs1} is incremented by 16.\\
\textbf{Operation}
\begin{lstlisting}
  temp[7:0] = qy[sel16*8+7:sel16*8]
  QACC_L[ 31:  0] = min(max(QACC_L[ 31:  0] + qx[  7:  0] * temp[7:0], -2^{31}), 2^{31}-1)
  QACC_L[ 63: 32] = min(max(QACC_L[ 63: 32] + qx[ 15:  8] * temp[7:0], -2^{31}), 2^{31}-1)
  QACC_L[ 95: 64] = min(max(QACC_L[ 95: 64] + qx[ 23: 16] * temp[7:0], -2^{31}), 2^{31}-1)
  QACC_L[127: 96] = min(max(QACC_L[128: 96] + qx[ 31: 24] * temp[7:0], -2^{31}), 2^{31}-1)
  QACC_L[159:128] = min(max(QACC_L[159:128] + qx[ 39: 32] * temp[7:0], -2^{31}), 2^{31}-1)
  QACC_L[191:160] = min(max(QACC_L[191:160] + qx[ 47: 40] * temp[7:0], -2^{31}), 2^{31}-1)
  QACC_L[223:192] = min(max(QACC_L[223:192] + qx[ 55: 48] * temp[7:0], -2^{31}), 2^{31}-1)
  QACC_L[255:224] = min(max(QACC_L[255:224] + qx[ 63: 56] * temp[7:0], -2^{31}), 2^{31}-1)
  QACC_H[ 31:  0] = min(max(QACC_H[ 31:  0] + qx[ 71: 64] * temp[7:0], -2^{31}), 2^{31}-1)
  QACC_H[ 63: 32] = min(max(QACC_H[ 63: 32] + qx[ 79: 72] * temp[7:0], -2^{31}), 2^{31}-1)
  QACC_H[ 95: 64] = min(max(QACC_H[ 95: 64] + qx[ 87: 80] * temp[7:0], -2^{31}), 2^{31}-1)
  QACC_H[127: 96] = min(max(QACC_H[127: 96] + qx[ 95: 88] * temp[7:0], -2^{31}), 2^{31}-1)
  QACC_H[159:128] = min(max(QACC_H[159:128] + qx[103: 96] * temp[7:0], -2^{31}), 2^{31}-1)
  QACC_H[191:160] = min(max(QACC_H[191:160] + qx[111:104] * temp[7:0], -2^{31}), 2^{31}-1)
  QACC_H[223:192] = min(max(QACC_H[223:192] + qx[119:112] * temp[7:0], -2^{31}), 2^{31}-1)
  QACC_H[255:224] = min(max(QACC_H[255:224] + qx[127:120] * temp[7:0], -2^{31}), 2^{31}-1)

  qu[127:0] = load128(rs1[31:0])
  rs1[31:0] = rs1[31:0] + 16
\end{lstlisting}

\textbf{Schedule}\\
\begin{longtable}[c]{|l|l|l|}
    \hline
    \rowcolor{lightgray}
    \textbf{Operands} & \textbf{Use} & \textbf{Def} \\\hline
    \endfirsthead
    \hline
    \rowcolor{lightgray}
     \textbf{Operands} & \textbf{Use} & \textbf{Def} \\\hline
    \endhead
    \hline
    \endfoot

    qu & - & 2 \\\hline
    rs1 & 1 & 1 \\\hline
     qx & 1 & - \\\hline
     qy & 1 & - \\\hline
     QACC\_H & 2 & 2 \\\hline
     QACC\_L & 2 & 2

\end{longtable}

\textbf{Format}\\

\begin{figure*}[!htbp]
 {\setlength{\tabcolsep}{1pt}
 \begin{center}
\begin{tabular}{@{}F@{}F@{}W@{}I@{}Y@{}I@{}F@{}W@{}F@{}F@{}O}
\instbitrange{31}{29}&
\instbitrange{28}{26}&
\instbitrange{25}{24}&
\instbit{23}&
\instbitrange{22}{19}&
\instbit{18}&
\instbitrange{17}{15}&
\instbitrange{14}{13}&
\instbitrange{12}{10}&
\instbitrange{9}{7}&
\instbitrange{6}{0} \\
\hline
\multicolumn{1}{|c|}{{\scriptsize qx[2:0]}} &
\multicolumn{1}{c|}{{\scriptsize qy[2:0]}} &
\multicolumn{1}{c|}{01} &
\multicolumn{1}{c|}{{\scriptsize sat[0]}} &
\multicolumn{1}{c|}{{\scriptsize sel16[3:0]}} &
\multicolumn{1}{c|}{{\scriptsize rs1[4]}} &
\multicolumn{1}{c|}{{\scriptsize rs1[2:0]}} &
\multicolumn{1}{c|}{11} &
\multicolumn{1}{c|}{{\scriptsize qu[2:0]}} &
\multicolumn{1}{c|}{111} &
\multicolumn{1}{c|}{0011111}\\
\hline
3& 3& 2& 1& 4& 1& 3& 2& 3& 3& 7 \\
\end{tabular}
 \end{center}
 }
 \vspace{-0.1in}
 \caption[]{ESP.VSMULAS.S8.QACC.LD.INCP} \end{figure*}


\newpage
\subsubsection{ESP.VSMULAS.U16.QACC}\label{instruction:ESPVSMULASU16QACC}
\textbf{Syntax}\\
ESP.VSMULAS.U16.QACC qx,qy,sel16,sat

\textbf{Description}\\
This instruction selects one out of the eight 16-bit data segments in the register \lstinline{qy} according to the immediate number \lstinline{sel16} and performs an unsigned multiplication on it and the eight 16-bit data segments in the register \lstinline{qx} respectively. The eight results obtained are added to the corresponding 64-bit data segments in the special registers \lstinline{QACC_H} and \lstinline{QACC_L} respectively. Then, each result is saturated to a 64-bit unsigned number and then stored in the corresponding 64-bit data segment in \lstinline{QACC_H} and \lstinline{QACC_L}.\\
\textbf{Operation}
\begin{lstlisting}
  temp[15:0] = qy[sel16*16+15:sel16*16]
  QACC_L[ 63:  0] = min(QACC_L[ 63:  0] + qx[ 15:  0] * temp[15:0], 2^{64}-1)
  QACC_L[127: 64] = min(QACC_L[127: 64] + qx[ 31: 16] * temp[15:0], 2^{64}-1)
  QACC_L[191:128] = min(QACC_L[191:128] + qx[ 47: 32] * temp[15:0], 2^{64}-1)
  QACC_L[255:192] = min(QACC_L[255:192] + qx[ 63: 48] * temp[15:0], 2^{64}-1)
  QACC_H[ 63:  0] = min(QACC_H[ 63:  0] + qx[ 79: 64] * temp[15:0], 2^{64}-1)
  QACC_H[127: 64] = min(QACC_H[127: 64] + qx[ 95: 80] * temp[15:0], 2^{64}-1)
  QACC_H[191:128] = min(QACC_H[191:128] + qx[111: 96] * temp[15:0], 2^{64}-1)
  QACC_H[255:192] = min(QACC_H[255:192] + qx[127:112] * temp[15:0], 2^{64}-1)
\end{lstlisting}

\textbf{Schedule}\\
\begin{longtable}[c]{|l|l|l|}
    \hline
    \rowcolor{lightgray}
    \textbf{Operands} & \textbf{Use} & \textbf{Def} \\\hline
    \endfirsthead
    \hline
    \rowcolor{lightgray}
     \textbf{Operands} & \textbf{Use} & \textbf{Def} \\\hline
    \endhead
    \hline
    \endfoot

     qx & 1 & - \\\hline
     qy & 1 & - \\\hline
     QACC\_H & 2 & 2 \\\hline
     QACC\_L & 2 & 2

\end{longtable}

\textbf{Format}\\

\begin{figure*}[!htbp]
 {\setlength{\tabcolsep}{1pt}
 \begin{center}
\begin{tabular}{@{}F@{}F@{}V@{}I@{}I@{}Y@{}E@{}O}
\instbitrange{31}{29}&
\instbitrange{28}{26}&
\instbitrange{25}{21}&
\instbit{20}&
\instbit{19}&
\instbitrange{18}{15}&
\instbitrange{14}{7}&
\instbitrange{6}{0} \\
\hline
\multicolumn{1}{|c|}{{\scriptsize qx[2:0]}} &
\multicolumn{1}{c|}{{\scriptsize qy[2:0]}} &
\multicolumn{1}{c|}{1X110} &
\multicolumn{1}{c|}{{\scriptsize sat[0]}} &
\multicolumn{1}{c|}{X} &
\multicolumn{1}{c|}{{\scriptsize sel16[3:0]}} &
\multicolumn{1}{c|}{01XX1X0X} &
\multicolumn{1}{c|}{0011011}\\
\hline
3& 3& 5& 1& 1& 4& 8& 7 \\
\end{tabular}
 \end{center}
 }
 \vspace{-0.1in}
 \caption[]{ESP.VSMULAS.U16.QACC} \end{figure*}


\newpage
\subsubsection{ESP.VSMULAS.U16.QACC.LD.INCP}\label{instruction:ESPVSMULASU16QACCLDINCP}
\textbf{Syntax}\\
ESP.VSMULAS.U16.QACC.LD.INCP qu,rs1,qx,qy,sel16,sat

\textbf{Description}\\
This instruction selects one out of the eight 16-bit data segments in the register \lstinline{qy} according to the immediate number \lstinline{sel16} and performs an unsigned multiplication on it and the eight 16-bit data segments in the register \lstinline{qx} respectively. The 8 results obtained are added to the corresponding 64-bit data segments in the special registers \lstinline{QACC_H} and \lstinline{QACC_L} respectively. Then, the result is saturated to a 64-bit unsigned number and stored in the corresponding 64-bit data segment in \lstinline{QACC_H} and \lstinline{QACC_L}.\\
During the operation, the 16-byte data is loaded from the memory to the register \lstinline{qu}. After the access is completed, the value in the register \lstinline{rs1} is incremented by 16.\\
\textbf{Operation}
\begin{lstlisting}
  temp[15:0] = qy[sel16*16+15:sel16*16]
  QACC_L[ 63:  0] = min(QACC_L[ 63:  0] + qx[ 15:  0] * temp[15:0], 2^{64}-1)
  QACC_L[127: 64] = min(QACC_L[127: 64] + qx[ 31: 16] * temp[15:0], 2^{64}-1)
  QACC_L[191:128] = min(QACC_L[191:128] + qx[ 47: 32] * temp[15:0], 2^{64}-1)
  QACC_L[255:192] = min(QACC_L[255:192] + qx[ 63: 48] * temp[15:0], 2^{64}-1)
  QACC_H[ 63:  0] = min(QACC_H[ 63:  0] + qx[ 79: 64] * temp[15:0], 2^{64}-1)
  QACC_H[127: 64] = min(QACC_H[127: 64] + qx[ 95: 80] * temp[15:0], 2^{64}-1)
  QACC_H[191:128] = min(QACC_H[191:128] + qx[111: 96] * temp[15:0], 2^{64}-1)
  QACC_H[255:192] = min(QACC_H[255:192] + qx[127:112] * temp[15:0], 2^{64}-1)

  qu[127:0] = load128(rs1[31:0])
  rs1[31:0] = rs1[31:0] + 16
\end{lstlisting}

\textbf{Schedule}\\
\begin{longtable}[c]{|l|l|l|}
    \hline
    \rowcolor{lightgray}
    \textbf{Operands} & \textbf{Use} & \textbf{Def} \\\hline
    \endfirsthead
    \hline
    \rowcolor{lightgray}
     \textbf{Operands} & \textbf{Use} & \textbf{Def} \\\hline
    \endhead
    \hline
    \endfoot

    qu & - & 2 \\\hline
    rs1 & 1 & 1 \\\hline
     qx & 1 & - \\\hline
     qy & 1 & - \\\hline
     QACC\_H & 2 & 2 \\\hline
     QACC\_L & 2 & 2

\end{longtable}

\textbf{Format}\\

\begin{figure*}[!htbp]
 {\setlength{\tabcolsep}{1pt}
 \begin{center}
\begin{tabular}{@{}F@{}F@{}W@{}I@{}Y@{}I@{}F@{}W@{}F@{}F@{}O}
\instbitrange{31}{29}&
\instbitrange{28}{26}&
\instbitrange{25}{24}&
\instbit{23}&
\instbitrange{22}{19}&
\instbit{18}&
\instbitrange{17}{15}&
\instbitrange{14}{13}&
\instbitrange{12}{10}&
\instbitrange{9}{7}&
\instbitrange{6}{0} \\
\hline
\multicolumn{1}{|c|}{{\scriptsize qx[2:0]}} &
\multicolumn{1}{c|}{{\scriptsize qy[2:0]}} &
\multicolumn{1}{c|}{10} &
\multicolumn{1}{c|}{{\scriptsize sat[0]}} &
\multicolumn{1}{c|}{{\scriptsize sel16[3:0]}} &
\multicolumn{1}{c|}{{\scriptsize rs1[4]}} &
\multicolumn{1}{c|}{{\scriptsize rs1[2:0]}} &
\multicolumn{1}{c|}{11} &
\multicolumn{1}{c|}{{\scriptsize qu[2:0]}} &
\multicolumn{1}{c|}{111} &
\multicolumn{1}{c|}{0011111}\\
\hline
3& 3& 2& 1& 4& 1& 3& 2& 3& 3& 7 \\
\end{tabular}
 \end{center}
 }
 \vspace{-0.1in}
 \caption[]{ESP.VSMULAS.U16.QACC.LD.INCP} \end{figure*}


\newpage
\subsubsection{ESP.VSMULAS.U8.QACC}\label{instruction:ESPVSMULASU8QACC}
\textbf{Syntax}\\
ESP.VSMULAS.U8.QACC qx,qy,sel16,sat

\textbf{Description}\\
This instruction selects one out of the 16 8-bit data segments in the register \lstinline{qy} according to the immediate number \lstinline{sel16} and performs an unsigned multiplication on it and the 16 8-bit data segments in the register \lstinline{qx} respectively. The 16 results obtained are added to the corresponding 32-bit data segments in the special registers \lstinline{QACC_H} and \lstinline{QACC_L} respectively. Then, each result is saturated to a 32-bit unsigned number and then stored in the corresponding 32-bit data segment in \lstinline{QACC_H} and \lstinline{QACC_L}.\\
\textbf{Operation}
\begin{lstlisting}
  temp[7:0] = qy[sel16*8+7:sel16*8]
  QACC_L[ 31:  0] = min(QACC_L[ 31:  0] + qx[  7:  0] * temp[7:0] , 2^{32}-1)
  QACC_L[ 63: 32] = min(QACC_L[ 63: 32] + qx[ 15:  8] * temp[7:0] , 2^{32}-1)
  QACC_L[ 95: 64] = min(QACC_L[ 95: 64] + qx[ 23: 16] * temp[7:0] , 2^{32}-1)
  QACC_L[127: 96] = min(QACC_L[128: 96] + qx[ 31: 24] * temp[7:0] , 2^{32}-1)
  QACC_L[159:128] = min(QACC_L[159:128] + qx[ 39: 32] * temp[7:0] , 2^{32}-1)
  QACC_L[191:160] = min(QACC_L[191:160] + qx[ 47: 40] * temp[7:0] , 2^{32}-1)
  QACC_L[223:192] = min(QACC_L[223:192] + qx[ 55: 48] * temp[7:0] , 2^{32}-1)
  QACC_L[255:224] = min(QACC_L[255:224] + qx[ 63: 56] * temp[7:0] , 2^{32}-1)
  QACC_H[ 31:  0] = min(QACC_H[ 31:  0] + qx[ 71: 64] * temp[7:0] , 2^{32}-1)
  QACC_H[ 63: 32] = min(QACC_H[ 63: 32] + qx[ 79: 72] * temp[7:0] , 2^{32}-1)
  QACC_H[ 95: 64] = min(QACC_H[ 95: 64] + qx[ 87: 80] * temp[7:0] , 2^{32}-1)
  QACC_H[127: 96] = min(QACC_H[127: 96] + qx[ 95: 88] * temp[7:0] , 2^{32}-1)
  QACC_H[159:128] = min(QACC_H[159:128] + qx[103: 96] * temp[7:0] , 2^{32}-1)
  QACC_H[191:160] = min(QACC_H[191:160] + qx[111:104] * temp[7:0] , 2^{32}-1)
  QACC_H[223:192] = min(QACC_H[223:192] + qx[119:112] * temp[7:0] , 2^{32}-1)
  QACC_H[255:224] = min(QACC_H[255:224] + qx[127:120] * temp[7:0] , 2^{32}-1)

\end{lstlisting}

\textbf{Schedule}\\
\begin{longtable}[c]{|l|l|l|}
    \hline
    \rowcolor{lightgray}
    \textbf{Operands} & \textbf{Use} & \textbf{Def} \\\hline
    \endfirsthead
    \hline
    \rowcolor{lightgray}
     \textbf{Operands} & \textbf{Use} & \textbf{Def} \\\hline
    \endhead
    \hline
    \endfoot

     qx & 1 & - \\\hline
     qy & 1 & - \\\hline
     QACC\_H & 2 & 2 \\\hline
     QACC\_L & 2 & 2

\end{longtable}

\textbf{Format}\\

\begin{figure*}[!htbp]
 {\setlength{\tabcolsep}{1pt}
 \begin{center}
\begin{tabular}{@{}F@{}F@{}V@{}I@{}I@{}Y@{}E@{}O}
\instbitrange{31}{29}&
\instbitrange{28}{26}&
\instbitrange{25}{21}&
\instbit{20}&
\instbit{19}&
\instbitrange{18}{15}&
\instbitrange{14}{7}&
\instbitrange{6}{0} \\
\hline
\multicolumn{1}{|c|}{{\scriptsize qx[2:0]}} &
\multicolumn{1}{c|}{{\scriptsize qy[2:0]}} &
\multicolumn{1}{c|}{1X100} &
\multicolumn{1}{c|}{{\scriptsize sat[0]}} &
\multicolumn{1}{c|}{X} &
\multicolumn{1}{c|}{{\scriptsize sel16[3:0]}} &
\multicolumn{1}{c|}{01XX1X0X} &
\multicolumn{1}{c|}{0011011}\\
\hline
3& 3& 5& 1& 1& 4& 8& 7 \\
\end{tabular}
 \end{center}
 }
 \vspace{-0.1in}
 \caption[]{ESP.VSMULAS.U8.QACC} \end{figure*}


\newpage
\subsubsection{ESP.VSMULAS.U8.QACC.LD.INCP}\label{instruction:ESPVSMULASU8QACCLDINCP}
\textbf{Syntax}\\
ESP.VSMULAS.U8.QACC.LD.INCP qu,rs1,qx,qy,sel16,sat

\textbf{Description}\\
This instruction selects one out of the 16 8-bit data segments in the register \lstinline{qy} according to the immediate number \lstinline{sel16} and performs an unsigned multiplication on it and the 16 8-bit data segments in the register \lstinline{qx} respectively. The 16 results obtained are added to the corresponding 32-bit data segments in the special registers \lstinline{QACC_H} and \lstinline{QACC_L} respectively. Then, each result is saturated to a 32-bit unsigned number and then stored in the corresponding 32-bit data segment in \lstinline{QACC_H} and \lstinline{QACC_L}.\\
During the operation, the 16-byte data is loaded from the memory into the register \lstinline{qu}. After the access is completed, the value in the register \lstinline{rs1} is incremented by 16.\\
\textbf{Operation}
\begin{lstlisting}
  temp[7:0] = qy[sel16*8+7:sel16*8]
  QACC_L[ 31:  0] = min(QACC_L[ 31:  0] + qx[  7:  0] * temp[7:0] , 2^{32}-1)
  QACC_L[ 63: 32] = min(QACC_L[ 63: 32] + qx[ 15:  8] * temp[7:0] , 2^{32}-1)
  QACC_L[ 95: 64] = min(QACC_L[ 95: 64] + qx[ 23: 16] * temp[7:0] , 2^{32}-1)
  QACC_L[127: 96] = min(QACC_L[128: 96] + qx[ 31: 24] * temp[7:0] , 2^{32}-1)
  QACC_L[159:128] = min(QACC_L[159:128] + qx[ 39: 32] * temp[7:0] , 2^{32}-1)
  QACC_L[191:160] = min(QACC_L[191:160] + qx[ 47: 40] * temp[7:0] , 2^{32}-1)
  QACC_L[223:192] = min(QACC_L[223:192] + qx[ 55: 48] * temp[7:0] , 2^{32}-1)
  QACC_L[255:224] = min(QACC_L[255:224] + qx[ 63: 56] * temp[7:0] , 2^{32}-1)
  QACC_H[ 31:  0] = min(QACC_H[ 31:  0] + qx[ 71: 64] * temp[7:0] , 2^{32}-1)
  QACC_H[ 63: 32] = min(QACC_H[ 63: 32] + qx[ 79: 72] * temp[7:0] , 2^{32}-1)
  QACC_H[ 95: 64] = min(QACC_H[ 95: 64] + qx[ 87: 80] * temp[7:0] , 2^{32}-1)
  QACC_H[127: 96] = min(QACC_H[127: 96] + qx[ 95: 88] * temp[7:0] , 2^{32}-1)
  QACC_H[159:128] = min(QACC_H[159:128] + qx[103: 96] * temp[7:0] , 2^{32}-1)
  QACC_H[191:160] = min(QACC_H[191:160] + qx[111:104] * temp[7:0] , 2^{32}-1)
  QACC_H[223:192] = min(QACC_H[223:192] + qx[119:112] * temp[7:0] , 2^{32}-1)
  QACC_H[255:224] = min(QACC_H[255:224] + qx[127:120] * temp[7:0] , 2^{32}-1)

  qu[127:0] = load128(rs1[31:0])
  rs1[31:0] = rs1[31:0] + 16
\end{lstlisting}

\textbf{Schedule}\\
\begin{longtable}[c]{|l|l|l|}
    \hline
    \rowcolor{lightgray}
    \textbf{Operands} & \textbf{Use} & \textbf{Def} \\\hline
    \endfirsthead
    \hline
    \rowcolor{lightgray}
     \textbf{Operands} & \textbf{Use} & \textbf{Def} \\\hline
    \endhead
    \hline
    \endfoot

    qu & - & 2 \\\hline
    rs1 & 1 & 1 \\\hline
     qx & 1 & - \\\hline
     qy & 1 & - \\\hline
     QACC\_H & 2 & 2 \\\hline
     QACC\_L & 2 & 2

\end{longtable}

\textbf{Format}\\

\begin{figure*}[!htbp]
 {\setlength{\tabcolsep}{1pt}
 \begin{center}
\begin{tabular}{@{}F@{}F@{}W@{}I@{}Y@{}I@{}F@{}W@{}F@{}F@{}O}
\instbitrange{31}{29}&
\instbitrange{28}{26}&
\instbitrange{25}{24}&
\instbit{23}&
\instbitrange{22}{19}&
\instbit{18}&
\instbitrange{17}{15}&
\instbitrange{14}{13}&
\instbitrange{12}{10}&
\instbitrange{9}{7}&
\instbitrange{6}{0} \\
\hline
\multicolumn{1}{|c|}{{\scriptsize qx[2:0]}} &
\multicolumn{1}{c|}{{\scriptsize qy[2:0]}} &
\multicolumn{1}{c|}{00} &
\multicolumn{1}{c|}{{\scriptsize sat[0]}} &
\multicolumn{1}{c|}{{\scriptsize sel16[3:0]}} &
\multicolumn{1}{c|}{{\scriptsize rs1[4]}} &
\multicolumn{1}{c|}{{\scriptsize rs1[2:0]}} &
\multicolumn{1}{c|}{11} &
\multicolumn{1}{c|}{{\scriptsize qu[2:0]}} &
\multicolumn{1}{c|}{111} &
\multicolumn{1}{c|}{0011111}\\
\hline
3& 3& 2& 1& 4& 1& 3& 2& 3& 3& 7 \\
\end{tabular}
 \end{center}
 }
 \vspace{-0.1in}
 \caption[]{ESP.VSMULAS.U8.QACC.LD.INCP} \end{figure*}


\newpage
\subsubsection{ESP.CMUL.S16}\label{instruction:ESPCMULS16}
\textbf{Syntax}\\
ESP.CMUL.S16 qz,qx,qy,sel4,sat,rm

\textbf{Description}\\

This instruction performs a 16-bit signed complex multiplication. The range of the immediate number \lstinline{sel4} is 0 \verb+~+ 3, which specifies the 32 bits in the two QR registers \lstinline{qx} and \lstinline{qy} for complex multiplication. The real and imaginary parts of the complex numbers are stored in the upper 16 bits and lower 16 bits of the 32 bits, respectively. The calculated real part and imaginary part results are stored in the corresponding 32 bits of the register \lstinline{qz}.\\

\textbf{Operation}
\begin{lstlisting}
if sel4 == 0:
  qz[ 15:  0] = (qx[ 15:  0] * qy[ 15:  0] - qx[ 31: 16] * qy[ 31: 16]) >> SAR[5:0]
  qz[ 31: 16] = (qx[ 15:  0] * qy[ 31: 16] + qx[ 31: 16] * qy[ 15:  0]) >> SAR[5:0]
  qz[ 47: 32] = (qx[ 47: 32] * qy[ 47: 32] - qx[ 63: 48] * qy[ 63: 48]) >> SAR[5:0]
  qz[ 63: 48] = (qx[ 47: 32] * qy[ 63: 48] + qx[ 63: 48] * qy[ 47: 32]) >> SAR[5:0]
if sel4 == 1:
  qz[ 79: 64] = (qx[ 79: 64] * qy[ 79: 64] - qx[ 95: 80] * qy[ 95: 80]) >> SAR[5:0]
  qz[ 95: 80] = (qx[ 79: 64] * qy[ 95: 80] + qx[ 95: 80] * qy[ 79: 64]) >> SAR[5:0]
  qz[111: 96] = (qx[111: 96] * qy[111: 96] - qx[127:112] * qy[127:112]) >> SAR[5:0]
  qz[127:112] = (qx[111: 96] * qy[127:112] + qx[127:112] * qy[111: 96]) >> SAR[5:0]
if sel4 == 2:
  qz[ 15:  0] = (qx[ 15:  0] * qy[ 15:  0] + qx[ 31: 16] * qy[ 31: 16]) >> SAR[5:0]
  qz[ 31: 16] = (qx[ 15:  0] * qy[ 31: 16] - qx[ 31: 16] * qy[ 15:  0]) >> SAR[5:0]
  qz[ 47: 32] = (qx[ 47: 32] * qy[ 47: 32] + qx[ 63: 48] * qy[ 63: 48]) >> SAR[5:0]
  qz[ 63: 48] = (qx[ 47: 32] * qy[ 63: 48] - qx[ 63: 48] * qy[ 47: 32]) >> SAR[5:0]
if sel4 == 3:
  qz[ 79: 64] = (qx[ 79: 64] * qy[ 79: 64] + qx[ 95: 80] * qy[ 95: 80]) >> SAR[5:0]
  qz[ 95: 80] = (qx[ 79: 64] * qy[ 95: 80] - qx[ 95: 80] * qy[ 79: 64]) >> SAR[5:0]
  qz[111: 96] = (qx[111: 96] * qy[111: 96] + qx[127:112] * qy[127:112]) >> SAR[5:0]
  qz[127:112] = (qx[111: 96] * qy[127:112] - qx[127:112] * qy[111: 96]) >> SAR[5:0]
\end{lstlisting}

\textbf{Schedule}\\
\begin{longtable}[c]{|l|l|l|}
    \hline
    \rowcolor{lightgray}
    \textbf{Operands} & \textbf{Use} & \textbf{Def} \\\hline
    \endfirsthead
    \hline
    \rowcolor{lightgray}
     \textbf{Operands} & \textbf{Use} & \textbf{Def} \\\hline
    \endhead
    \hline
    \endfoot


    qz & - & 2 \\\hline
     qx & 1 & - \\\hline
     qy & 1 & -


\end{longtable}

\textbf{Format}\\

\begin{figure*}[!htbp]
 {\setlength{\tabcolsep}{1pt}
 \begin{center}
\begin{tabular}{@{}F@{}F@{}E@{}I@{}W@{}W@{}F@{}F@{}O}
\instbitrange{31}{29}&
\instbitrange{28}{26}&
\instbitrange{25}{18}&
\instbit{17}&
\instbitrange{16}{15}&
\instbitrange{14}{13}&
\instbitrange{12}{10}&
\instbitrange{9}{7}&
\instbitrange{6}{0} \\
\hline
\multicolumn{1}{|c|}{{\scriptsize qx[2:0]}} &
\multicolumn{1}{c|}{{\scriptsize qy[2:0]}} &
\multicolumn{1}{c|}{111000X1} &
\multicolumn{1}{c|}{{\scriptsize sat[0]}} &
\multicolumn{1}{c|}{{\scriptsize sel4[1:0]}} &
\multicolumn{1}{c|}{00} &
\multicolumn{1}{c|}{{\scriptsize rm[2:0]}} &
\multicolumn{1}{c|}{{\scriptsize qz[2:0]}} &
\multicolumn{1}{c|}{0011111}\\
\hline
3& 3& 8& 1& 2& 2& 3& 3& 7 \\
\end{tabular}
 \end{center}
 }
 \vspace{-0.1in}
 \caption[]{ESP.CMUL.S16} \end{figure*}


\newpage
\subsubsection{ESP.CMUL.S16.LD.INCP}\label{instruction:ESPCMULS16LDINCP}
\textbf{Syntax}\\
ESP.CMUL.S16.LD.INCP qu,rs1,qz,qx,qy,sel4,sat,rm

\textbf{Description}\\
This instruction performs a 16-bit signed complex multiplication. The range of the immediate number \lstinline{sel4} is 0 \verb+~+ 3, which specifies the 32 bits in the two QR registers \lstinline{qx} and \lstinline{qy} for complex multiplication. The real and imaginary parts of complex numbers are stored in the upper 16 bits and lower 16 bits of the 32 bits respectively. The calculated real part and imaginary part results are stored in the corresponding 32 bits of the register \lstinline{qz}.\\
Meanwhile, it also loads 16-byte data from the memory into the register \lstinline{qu}. After the access, the value in the register \lstinline{rs1} is incremented by 16.\\

\textbf{Operation}
\begin{lstlisting}
if sel4 == 0:
  qz[ 15:  0] = (qx[ 15:  0] * qy[ 15:  0] - qx[ 31: 16] * qy[ 31: 16]) >> SAR[5:0]
  qz[ 31: 16] = (qx[ 15:  0] * qy[ 31: 16] + qx[ 31: 16] * qy[ 15:  0]) >> SAR[5:0]
  qz[ 47: 32] = (qx[ 47: 32] * qy[ 47: 32] - qx[ 63: 48] * qy[ 63: 48]) >> SAR[5:0]
  qz[ 63: 48] = (qx[ 47: 32] * qy[ 63: 48] + qx[ 63: 48] * qy[ 47: 32]) >> SAR[5:0]
if sel4 == 1:
  qz[ 79: 64] = (qx[ 79: 64] * qy[ 79: 64] - qx[ 95: 80] * qy[ 95: 80]) >> SAR[5:0]
  qz[ 95: 80] = (qx[ 79: 64] * qy[ 95: 80] + qx[ 95: 80] * qy[ 79: 64]) >> SAR[5:0]
  qz[111: 96] = (qx[111: 96] * qy[111: 96] - qx[127:112] * qy[127:112]) >> SAR[5:0]
  qz[127:112] = (qx[111: 96] * qy[127:112] + qx[127:112] * qy[111: 96]) >> SAR[5:0]
if sel4 == 2:
  qz[ 15:  0] = (qx[ 15:  0] * qy[ 15:  0] + qx[ 31: 16] * qy[ 31: 16]) >> SAR[5:0]
  qz[ 31: 16] = (qx[ 15:  0] * qy[ 31: 16] - qx[ 31: 16] * qy[ 15:  0]) >> SAR[5:0]
  qz[ 47: 32] = (qx[ 47: 32] * qy[ 47: 32] + qx[ 63: 48] * qy[ 63: 48]) >> SAR[5:0]
  qz[ 63: 48] = (qx[ 47: 32] * qy[ 63: 48] - qx[ 63: 48] * qy[ 47: 32]) >> SAR[5:0]
if sel4 == 3:
  qz[ 79: 64] = (qx[ 79: 64] * qy[ 79: 64] + qx[ 95: 80] * qy[ 95: 80]) >> SAR[5:0]
  qz[ 95: 80] = (qx[ 79: 64] * qy[ 95: 80] - qx[ 95: 80] * qy[ 79: 64]) >> SAR[5:0]
  qz[111: 96] = (qx[111: 96] * qy[111: 96] + qx[127:112] * qy[127:112]) >> SAR[5:0]
  qz[127:112] = (qx[111: 96] * qy[127:112] - qx[127:112] * qy[111: 96]) >> SAR[5:0]

  qu[127:0] = load128(rs1[31:0])
  rs1[31:0] = rs1[31:0] + 16
\end{lstlisting}

\textbf{Schedule}\\
\begin{longtable}[c]{|l|l|l|}
    \hline
    \rowcolor{lightgray}
    \textbf{Operands} & \textbf{Use} & \textbf{Def} \\\hline
    \endfirsthead
    \hline
    \rowcolor{lightgray}
     \textbf{Operands} & \textbf{Use} & \textbf{Def} \\\hline
    \endhead
    \hline
    \endfoot


    qz & - & 2 \\\hline
     qx & 1 & - \\\hline
     qy & 1 & - \\\hline
     qu & - & 2 \\\hline
     rs1 & 1 & 1

\end{longtable}

\textbf{Format}\\

\begin{figure*}[!htbp]
 {\setlength{\tabcolsep}{1pt}
 \begin{center}
\begin{tabular}{@{}F@{}F@{}F@{}I@{}W@{}I@{}I@{}F@{}W@{}F@{}F@{}O}
\instbitrange{31}{29}&
\instbitrange{28}{26}&
\instbitrange{25}{23}&
\instbit{22}&
\instbitrange{21}{20}&
\instbit{19}&
\instbit{18}&
\instbitrange{17}{15}&
\instbitrange{14}{13}&
\instbitrange{12}{10}&
\instbitrange{9}{7}&
\instbitrange{6}{0} \\
\hline
\multicolumn{1}{|c|}{{\scriptsize qx[2:0]}} &
\multicolumn{1}{c|}{{\scriptsize qy[2:0]}} &
\multicolumn{1}{c|}{011} &
\multicolumn{1}{c|}{{\scriptsize sat[0]}} &
\multicolumn{1}{c|}{{\scriptsize sel4[1:0]}} &
\multicolumn{1}{c|}{{\scriptsize rm[2]}} &
\multicolumn{1}{c|}{{\scriptsize rs1[4]}} &
\multicolumn{1}{c|}{{\scriptsize rs1[2:0]}} &
\multicolumn{1}{c|}{{\scriptsize rm[1:0]}} &
\multicolumn{1}{c|}{{\scriptsize qu[2:0]}} &
\multicolumn{1}{c|}{{\scriptsize qz[2:0]}} &
\multicolumn{1}{c|}{1111011}\\
\hline
3& 3& 3& 1& 2& 1& 1& 3& 2& 3& 3& 7 \\
\end{tabular}
 \end{center}
 }
 \vspace{-0.1in}
 \caption[]{ESP.CMUL.S16.LD.INCP} \end{figure*}


\newpage
\subsubsection{ESP.CMUL.S16.ST.INCP}\label{instruction:ESPCMULS16STINCP}
\textbf{Syntax}\\
ESP.CMUL.S16.ST.INCP qu,rs1,qz,qx,qy,sel4,sat,rm

\textbf{Description}\\
This instruction performs a 16-bit signed complex multiplication. The range of the immediate number \lstinline{sel4} is 0 \verb+~+ 3, which specifies the 32 bits in the two QR registers \lstinline{qx} and \lstinline{qy} for complex multiplication. The real and imaginary parts of complex numbers are stored in the upper 16 bits and lower 16 bits of the 32 bits, respectively. The calculated real part and imaginary part results are stored in the corresponding 32 bits of the \lstinline{qz} register.

Meanwhile, it also stores the 16-byte data of \lstinline{qu} in memory. After the access, the value in the \lstinline{rs1} register is incremented by 16.\\
\textbf{Operation}
\begin{lstlisting}
if sel4 == 0:
  qz[ 15:  0] = (qx[ 15:  0] * qy[ 15:  0] - qx[ 31: 16] * qy[ 31: 16]) >> SAR[5:0]
  qz[ 31: 16] = (qx[ 15:  0] * qy[ 31: 16] + qx[ 31: 16] * qy[ 15:  0]) >> SAR[5:0]
  qz[ 47: 32] = (qx[ 47: 32] * qy[ 47: 32] - qx[ 63: 48] * qy[ 63: 48]) >> SAR[5:0]
  qz[ 63: 48] = (qx[ 47: 32] * qy[ 63: 48] + qx[ 63: 48] * qy[ 47: 32]) >> SAR[5:0]
if sel4 == 1:
  qz[ 79: 64] = (qx[ 79: 64] * qy[ 79: 64] - qx[ 95: 80] * qy[ 95: 80]) >> SAR[5:0]
  qz[ 95: 80] = (qx[ 79: 64] * qy[ 95: 80] + qx[ 95: 80] * qy[ 79: 64]) >> SAR[5:0]
  qz[111: 96] = (qx[111: 96] * qy[111: 96] - qx[127:112] * qy[127:112]) >> SAR[5:0]
  qz[127:112] = (qx[111: 96] * qy[127:112] + qx[127:112] * qy[111: 96]) >> SAR[5:0]
if sel4 == 2:
  qz[ 15:  0] = (qx[ 15:  0] * qy[ 15:  0] + qx[ 31: 16] * qy[ 31: 16]) >> SAR[5:0]
  qz[ 31: 16] = (qx[ 15:  0] * qy[ 31: 16] - qx[ 31: 16] * qy[ 15:  0]) >> SAR[5:0]
  qz[ 47: 32] = (qx[ 47: 32] * qy[ 47: 32] + qx[ 63: 48] * qy[ 63: 48]) >> SAR[5:0]
  qz[ 63: 48] = (qx[ 47: 32] * qy[ 63: 48] - qx[ 63: 48] * qy[ 47: 32]) >> SAR[5:0]
if sel4 == 3:
  qz[ 79: 64] = (qx[ 79: 64] * qy[ 79: 64] + qx[ 95: 80] * qy[ 95: 80]) >> SAR[5:0]
  qz[ 95: 80] = (qx[ 79: 64] * qy[ 95: 80] - qx[ 95: 80] * qy[ 79: 64]) >> SAR[5:0]
  qz[111: 96] = (qx[111: 96] * qy[111: 96] + qx[127:112] * qy[127:112]) >> SAR[5:0]
  qz[127:112] = (qx[111: 96] * qy[127:112] - qx[127:112] * qy[111: 96]) >> SAR[5:0]

  qu[127:0] => store128(rs1[31:0])
  rs1[31:0] = rs1[31:0] + 16
\end{lstlisting}

\textbf{Schedule}\\
\begin{longtable}[c]{|l|l|l|}
    \hline
    \rowcolor{lightgray}
    \textbf{Operands} & \textbf{Use} & \textbf{Def} \\\hline
    \endfirsthead
    \hline
    \rowcolor{lightgray}
     \textbf{Operands} & \textbf{Use} & \textbf{Def} \\\hline
    \endhead
    \hline
    \endfoot


    qz & - & 2 \\\hline
     qx & 1 & - \\\hline
     qy & 1 & - \\\hline
     qu & 2 & - \\\hline
     rs1 & 1 & 1

\end{longtable}

\textbf{Format}\\

\begin{figure*}[!htbp]
 {\setlength{\tabcolsep}{1pt}
 \begin{center}
\begin{tabular}{@{}F@{}F@{}F@{}I@{}W@{}I@{}I@{}F@{}W@{}F@{}F@{}O}
\instbitrange{31}{29}&
\instbitrange{28}{26}&
\instbitrange{25}{23}&
\instbit{22}&
\instbitrange{21}{20}&
\instbit{19}&
\instbit{18}&
\instbitrange{17}{15}&
\instbitrange{14}{13}&
\instbitrange{12}{10}&
\instbitrange{9}{7}&
\instbitrange{6}{0} \\
\hline
\multicolumn{1}{|c|}{{\scriptsize qx[2:0]}} &
\multicolumn{1}{c|}{{\scriptsize qy[2:0]}} &
\multicolumn{1}{c|}{111} &
\multicolumn{1}{c|}{{\scriptsize sat[0]}} &
\multicolumn{1}{c|}{{\scriptsize sel4[1:0]}} &
\multicolumn{1}{c|}{{\scriptsize rm[2]}} &
\multicolumn{1}{c|}{{\scriptsize rs1[4]}} &
\multicolumn{1}{c|}{{\scriptsize rs1[2:0]}} &
\multicolumn{1}{c|}{{\scriptsize rm[1:0]}} &
\multicolumn{1}{c|}{{\scriptsize qu[2:0]}} &
\multicolumn{1}{c|}{{\scriptsize qz[2:0]}} &
\multicolumn{1}{c|}{1111011}\\
\hline
3& 3& 3& 1& 2& 1& 1& 3& 2& 3& 3& 7 \\
\end{tabular}
 \end{center}
 }
 \vspace{-0.1in}
 \caption[]{ESP.CMUL.S16.ST.INCP} \end{figure*}


\newpage
\subsubsection{ESP.CMUL.S8}\label{instruction:ESPCMULS8}
\textbf{Syntax}\\
ESP.CMUL.S8 qz,qx,qy,sel4,sat,rm

\textbf{Description}\\
This instruction performs an 8-bit signed complex multiplication. The range of the immediate number \lstinline{sel4} is 0 \verb+~+ 3, which specifies the 16 bits in the two QR registers \lstinline{qx} and \lstinline{qy} for complex multiplication. The real and imaginary parts of complex numbers are stored in the upper 8 bits and lower 8 bits of the 16 bits respectively. The calculated real part and imaginary part results are stored in the corresponding 16 bits of the register \lstinline{qz}.

\textbf{Operation}\\
\begin{lstlisting}
if sel4 == 0:
  qz[  7:  0] = (qx[  7:  0] * qy[  7:  0] - qx[ 15:  8] * qy[ 15:  8]) >> SAR[5:0]
  qz[ 15:  8] = (qx[  7:  0] * qy[ 15:  8] + qx[ 15:  8] * qy[  7:  0]) >> SAR[5:0]
  qz[ 23: 16] = (qx[ 23: 16] * qy[ 23: 16] - qx[ 31: 16] * qy[ 31: 24]) >> SAR[5:0]
  qz[ 31: 24] = (qx[ 23: 16] * qy[ 31: 24] + qx[ 31: 16] * qy[ 23: 16]) >> SAR[5:0]
  qz[ 39: 32] = (qx[ 39: 32] * qy[ 39: 32] - qx[ 47: 40] * qy[ 47: 40]) >> SAR[5:0]
  qz[47 : 40] = (qx[39 : 32] * qy[47 : 40] + qx[47 : 40] * qy[39 : 32]) >> SAR[5:0]
  qz[55 : 48] = (qx[55 : 48] * qy[55 : 48] - qx[63 : 56] * qy[63 : 56]) >> SAR[5:0]
  qz[63 : 56] = (qx[55 : 48] * qy[63 : 56] + qx[63 : 56] * qy[55 : 48]) >> SAR[5:0]
else if sel4 == 1:
  qz[71 : 64] = (qx[71 : 64] * qy[71 : 64] - qx[79 : 72] * qy[79 : 72]) >> SAR[5:0]
  qz[79 : 72] = (qx[71 : 64] * qy[79 : 72] + qx[79 : 72] * qy[71 : 64]) >> SAR[5:0]
  qz[87 : 80] = (qx[87 : 80] * qy[87 : 80] - qx[95 : 88] * qy[95 : 88]) >> SAR[5:0]
  qz[95 : 88] = (qx[87 : 80] * qy[95 : 88] + qx[95 : 88] * qy[87 : 80]) >> SAR[5:0]
  qz[103: 96] = (qx[103: 96] * qy[103: 96] - qx[111:104] * qy[111:104]) >> SAR[5:0]
  qz[111:104] = (qx[103: 96] * qy[111:104] + qx[111:104] * qy[103: 96]) >> SAR[5:0]
  qz[119:112] = (qx[119:112] * qy[119:112] - qx[127:120] * qy[127:120]) >> SAR[5:0]
  qz[127:120] = (qx[119:112] * qy[127:120] + qx[127:120] * qy[119:112]) >> SAR[5:0]
else if sel4 == 2:
  qz[7  :  0] = (qx[7  :  0] * qy[7  :  0] + qx[15 :  8] * qy[15 :  8]) >> SAR[5:0]
  qz[15 :  8] = (qx[7  :  0] * qy[15 :  8] - qx[15 :  8] * qy[7  :  0]) >> SAR[5:0]
  qz[23 : 16] = (qx[23 : 16] * qy[23 : 16] + qx[31 : 24] * qy[31 : 24]) >> SAR[5:0]
  qz[31 : 24] = (qx[23 : 16] * qy[31 : 24] - qx[31 : 24] * qy[23 : 16]) >> SAR[5:0]
  qz[39 : 32] = (qx[39 : 32] * qy[39 : 32] + qx[47 : 40] * qy[47 : 40]) >> SAR[5:0]
  qz[47 : 40] = (qx[39 : 32] * qy[47 : 40] - qx[47 : 40] * qy[39 : 32]) >> SAR[5:0]
  qz[55 : 48] = (qx[55 : 48] * qy[55 : 48] + qx[63 : 56] * qy[63 : 56]) >> SAR[5:0]
  qz[63 : 56] = (qx[55 : 48] * qy[63 : 56] - qx[63 : 56] * qy[55 : 48]) >> SAR[5:0]
else if sel4 == 3:
  qz[71 : 64] = (qx[71 : 64] * qy[71 : 64] + qx[79 : 72] * qy[79 : 72]) >> SAR[5:0]
  qz[79 : 72] = (qx[71 : 64] * qy[79 : 72] - qx[79 : 72] * qy[71 : 64]) >> SAR[5:0]
  qz[87 : 80] = (qx[87 : 80] * qy[87 : 80] + qx[95 : 88] * qy[95 : 88]) >> SAR[5:0]
  qz[95 : 88] = (qx[87 : 80] * qy[95 : 88] - qx[95 : 88] * qy[87 : 80]) >> SAR[5:0]
  qz[103: 96] = (qx[103: 96] * qy[103: 96] + qx[111:104] * qy[111:104]) >> SAR[5:0]
  qz[111:104] = (qx[103: 96] * qy[111:104] - qx[111:104] * qy[103: 96]) >> SAR[5:0]
  qz[119:112] = (qx[119:112] * qy[119:112] + qx[127:120] * qy[127:120]) >> SAR[5:0]
  qz[127:120] = (qx[119:112] * qy[127:120] - qx[127:120] * qy[119:112]) >> SAR[5:0]

\end{lstlisting}

\textbf{Schedule}\\
\begin{longtable}[c]{|l|l|l|}
    \hline
    \rowcolor{lightgray}
    \textbf{Operands} & \textbf{Use} & \textbf{Def} \\\hline
    \endfirsthead
    \hline
    \rowcolor{lightgray}
     \textbf{Operands} & \textbf{Use} & \textbf{Def} \\\hline
    \endhead
    \hline
    \endfoot


    qz & - & 2 \\\hline
     qx & 1 & - \\\hline
     qy & 1 & -
\end{longtable}

\textbf{Format}\\

\begin{figure*}[!htbp]
 {\setlength{\tabcolsep}{1pt}
 \begin{center}
\begin{tabular}{@{}F@{}F@{}E@{}I@{}W@{}W@{}F@{}F@{}O}
\instbitrange{31}{29}&
\instbitrange{28}{26}&
\instbitrange{25}{18}&
\instbit{17}&
\instbitrange{16}{15}&
\instbitrange{14}{13}&
\instbitrange{12}{10}&
\instbitrange{9}{7}&
\instbitrange{6}{0} \\
\hline
\multicolumn{1}{|c|}{{\scriptsize qx[2:0]}} &
\multicolumn{1}{c|}{{\scriptsize qy[2:0]}} &
\multicolumn{1}{c|}{110000X1} &
\multicolumn{1}{c|}{{\scriptsize sat[0]}} &
\multicolumn{1}{c|}{{\scriptsize sel4[1:0]}} &
\multicolumn{1}{c|}{00} &
\multicolumn{1}{c|}{{\scriptsize rm[2:0]}} &
\multicolumn{1}{c|}{{\scriptsize qz[2:0]}} &
\multicolumn{1}{c|}{0011111}\\
\hline
3& 3& 8& 1& 2& 2& 3& 3& 7 \\
\end{tabular}
 \end{center}
 }
 \vspace{-0.1in}
 \caption[]{ESP.CMUL.S8} \end{figure*}


\newpage
\subsubsection{ESP.CMUL.S8.LD.INCP}\label{instruction:ESPCMULS8LDINCP}
\textbf{Syntax}\\
ESP.CMUL.S8.LD.INCP qu,rs1,qz,qx,qy,sel4,sat,rm

\textbf{Description}\\
This instruction performs an 8-bit signed complex multiplication. The range of the immediate number \lstinline{sel4} is 0 \verb+~+ 3, which specifies the 16 bits in the two QR registers \lstinline{qx} and \lstinline{qy} for complex multiplication. The real and imaginary parts of complex numbers are stored in the upper 8 bits and lower 8 bits of the 16 bits respectively. The calculated real part and imaginary part results are stored in the corresponding 16 bits of the register \lstinline{qz}.

Meanwhile, it also loads 16-byte data from the memory into the register \lstinline{qu}. After the access, the value in the register \lstinline{rs1} is incremented by 16.\\
\textbf{Operation}\\
\begin{lstlisting}
if sel4 == 0:
  qz[  7:  0] = (qx[  7:  0] * qy[  7:  0] - qx[ 15:  8] * qy[ 15:  8]) >> SAR[5  :  0]
  qz[ 15:  8] = (qx[  7:  0] * qy[ 15:  8] + qx[ 15:  8] * qy[  7:  0]) >> SAR[5  :  0]
  qz[ 23: 16] = (qx[ 23: 16] * qy[ 23: 16] - qx[ 31: 16] * qy[ 31: 24]) >> SAR[5  :  0]
  qz[ 31: 24] = (qx[ 23: 16] * qy[ 31: 24] + qx[ 31: 16] * qy[ 23: 16]) >> SAR[5  :  0]
  qz[ 39: 32] = (qx[ 39: 32] * qy[ 39: 32] - qx[ 47: 40] * qy[ 47: 40]) >> SAR[5  :  0]
  qz[47 : 40] = (qx[39 : 32] * qy[47 : 40] + qx[47 : 40] * qy[39 : 32]) >> SAR[5  :  0]
  qz[55 : 48] = (qx[55 : 48] * qy[55 : 48] - qx[63 : 56] * qy[63 : 56]) >> SAR[5  :  0]
  qz[63 : 56] = (qx[55 : 48] * qy[63 : 56] + qx[63 : 56] * qy[55 : 48]) >> SAR[5  :  0]
else if sel4 == 1:
  qz[71 : 64] = (qx[71 : 64] * qy[71 : 64] - qx[79 : 72] * qy[79 : 72]) >> SAR[5  :  0]
  qz[79 : 72] = (qx[71 : 64] * qy[79 : 72] + qx[79 : 72] * qy[71 : 64]) >> SAR[5  :  0]
  qz[87 : 80] = (qx[87 : 80] * qy[87 : 80] - qx[95 : 88] * qy[95 : 88]) >> SAR[5  :  0]
  qz[95 : 88] = (qx[87 : 80] * qy[95 : 88] + qx[95 : 88] * qy[87 : 80]) >> SAR[5  :  0]
  qz[103: 96] = (qx[103: 96] * qy[103: 96] - qx[111:104] * qy[111:104]) >> SAR[5  :  0]
  qz[111:104] = (qx[103: 96] * qy[111:104] + qx[111:104] * qy[103: 96]) >> SAR[5  :  0]
  qz[119:112] = (qx[119:112] * qy[119:112] - qx[127:120] * qy[127:120]) >> SAR[5  :  0]
  qz[127:120] = (qx[119:112] * qy[127:120] + qx[127:120] * qy[119:112]) >> SAR[5  :  0]
else if sel4 == 2:
  qz[7  :  0] = (qx[7  :  0] * qy[7  :  0] + qx[15 :  8] * qy[15 :  8]) >> SAR[5  :  0]
  qz[15 :  8] = (qx[7  :  0] * qy[15 :  8] - qx[15 :  8] * qy[7  :  0]) >> SAR[5  :  0]
  qz[23 : 16] = (qx[23 : 16] * qy[23 : 16] + qx[31 : 24] * qy[31 : 24]) >> SAR[5  :  0]
  qz[31 : 24] = (qx[23 : 16] * qy[31 : 24] - qx[31 : 24] * qy[23 : 16]) >> SAR[5  :  0]
  qz[39 : 32] = (qx[39 : 32] * qy[39 : 32] + qx[47 : 40] * qy[47 : 40]) >> SAR[5  :  0]
  qz[47 : 40] = (qx[39 : 32] * qy[47 : 40] - qx[47 : 40] * qy[39 : 32]) >> SAR[5  :  0]
  qz[55 : 48] = (qx[55 : 48] * qy[55 : 48] + qx[63 : 56] * qy[63 : 56]) >> SAR[5  :  0]
  qz[63 : 56] = (qx[55 : 48] * qy[63 : 56] - qx[63 : 56] * qy[55 : 48]) >> SAR[5  :  0]
else if sel4 == 3:
  qz[71 : 64] = (qx[71 : 64] * qy[71 : 64] + qx[79 : 72] * qy[79 : 72]) >> SAR[5  :  0]
  qz[79 : 72] = (qx[71 : 64] * qy[79 : 72] - qx[79 : 72] * qy[71 : 64]) >> SAR[5  :  0]
  qz[87 : 80] = (qx[87 : 80] * qy[87 : 80] + qx[95 : 88] * qy[95 : 88]) >> SAR[5  :  0]
  qz[95 : 88] = (qx[87 : 80] * qy[95 : 88] - qx[95 : 88] * qy[87 : 80]) >> SAR[5  :  0]
  qz[103: 96] = (qx[103: 96] * qy[103: 96] + qx[111:104] * qy[111:104]) >> SAR[5  :  0]
  qz[111:104] = (qx[103: 96] * qy[111:104] - qx[111:104] * qy[103: 96]) >> SAR[5  :  0]
  qz[119:112] = (qx[119:112] * qy[119:112] + qx[127:120] * qy[127:120]) >> SAR[5  :  0]
  qz[127:120] = (qx[119:112] * qy[127:120] - qx[127:120] * qy[119:112]) >> SAR[5  :  0]

qu[127:0] = load128(rs1[31:0])
rs1[31:0] = rs1[31:0] + 16

\end{lstlisting}

\textbf{Schedule}\\
\begin{longtable}[c]{|l|l|l|}
    \hline
    \rowcolor{lightgray}
    \textbf{Operands} & \textbf{Use} & \textbf{Def} \\\hline
    \endfirsthead
    \hline
    \rowcolor{lightgray}
     \textbf{Operands} & \textbf{Use} & \textbf{Def} \\\hline
    \endhead
    \hline
    \endfoot


    qz & - & 2 \\\hline
     qx & 1 & - \\\hline
     qy & 1 & - \\\hline
     qu & - & 2 \\\hline
     rs1 & 1 & 1

\end{longtable}

\textbf{Format}\\

\begin{figure*}[!htbp]
 {\setlength{\tabcolsep}{1pt}
 \begin{center}
\begin{tabular}{@{}F@{}F@{}F@{}I@{}W@{}I@{}I@{}F@{}W@{}F@{}F@{}O}
\instbitrange{31}{29}&
\instbitrange{28}{26}&
\instbitrange{25}{23}&
\instbit{22}&
\instbitrange{21}{20}&
\instbit{19}&
\instbit{18}&
\instbitrange{17}{15}&
\instbitrange{14}{13}&
\instbitrange{12}{10}&
\instbitrange{9}{7}&
\instbitrange{6}{0} \\
\hline
\multicolumn{1}{|c|}{{\scriptsize qx[2:0]}} &
\multicolumn{1}{c|}{{\scriptsize qy[2:0]}} &
\multicolumn{1}{c|}{001} &
\multicolumn{1}{c|}{{\scriptsize sat[0]}} &
\multicolumn{1}{c|}{{\scriptsize sel4[1:0]}} &
\multicolumn{1}{c|}{{\scriptsize rm[2]}} &
\multicolumn{1}{c|}{{\scriptsize rs1[4]}} &
\multicolumn{1}{c|}{{\scriptsize rs1[2:0]}} &
\multicolumn{1}{c|}{{\scriptsize rm[1:0]}} &
\multicolumn{1}{c|}{{\scriptsize qu[2:0]}} &
\multicolumn{1}{c|}{{\scriptsize qz[2:0]}} &
\multicolumn{1}{c|}{1111011}\\
\hline
3& 3& 3& 1& 2& 1& 1& 3& 2& 3& 3& 7 \\
\end{tabular}
 \end{center}
 }
 \vspace{-0.1in}
 \caption[]{ESP.CMUL.S8.LD.INCP} \end{figure*}


\newpage
\subsubsection{ESP.CMUL.S8.ST.INCP}\label{instruction:ESPCMULS8STINCP}
\textbf{Syntax}\\
ESP.CMUL.S8.ST.INCP qu,rs1,qz,qx,qy,sel4,sat,rm

\textbf{Description}\\
This instruction performs an 8-bit signed complex multiplication. The range of the immediate number \lstinline{sel4} is 0 \verb+~+ 3, which specifies the 16 bits in the two QR registers \lstinline{qx} and \lstinline{qy} for complex multiplication. The real and imaginary parts of complex numbers are stored in the upper 8 bits and lower 8 bits of the 16 bits respectively. The calculated real part and imaginary part results are stored in the corresponding 16 bits of the register \lstinline{qz}.

Meanwhile, it also stores the 16-byte data of \lstinline{qu} in memory. After the access, the value in the register \lstinline{rs1} is incremented by 16.

\textbf{Operation}\\
\begin{lstlisting}
if sel4 == 0:
  qz[  7:  0] = (qx[  7:  0] * qy[  7:  0] - qx[ 15:  8] * qy[ 15:  8]) >> SAR[5  :  0]
  qz[ 15:  8] = (qx[  7:  0] * qy[ 15:  8] + qx[ 15:  8] * qy[  7:  0]) >> SAR[5  :  0]
  qz[ 23: 16] = (qx[ 23: 16] * qy[ 23: 16] - qx[ 31: 16] * qy[ 31: 24]) >> SAR[5  :  0]
  qz[ 31: 24] = (qx[ 23: 16] * qy[ 31: 24] + qx[ 31: 16] * qy[ 23: 16]) >> SAR[5  :  0]
  qz[ 39: 32] = (qx[ 39: 32] * qy[ 39: 32] - qx[ 47: 40] * qy[ 47: 40]) >> SAR[5  :  0]
  qz[47 : 40] = (qx[39 : 32] * qy[47 : 40] + qx[47 : 40] * qy[39 : 32]) >> SAR[5  :  0]
  qz[55 : 48] = (qx[55 : 48] * qy[55 : 48] - qx[63 : 56] * qy[63 : 56]) >> SAR[5  :  0]
  qz[63 : 56] = (qx[55 : 48] * qy[63 : 56] + qx[63 : 56] * qy[55 : 48]) >> SAR[5  :  0]
else if sel4 == 1:
  qz[71 : 64] = (qx[71 : 64] * qy[71 : 64] - qx[79 : 72] * qy[79 : 72]) >> SAR[5  :  0]
  qz[79 : 72] = (qx[71 : 64] * qy[79 : 72] + qx[79 : 72] * qy[71 : 64]) >> SAR[5  :  0]
  qz[87 : 80] = (qx[87 : 80] * qy[87 : 80] - qx[95 : 88] * qy[95 : 88]) >> SAR[5  :  0]
  qz[95 : 88] = (qx[87 : 80] * qy[95 : 88] + qx[95 : 88] * qy[87 : 80]) >> SAR[5  :  0]
  qz[103: 96] = (qx[103: 96] * qy[103: 96] - qx[111:104] * qy[111:104]) >> SAR[5  :  0]
  qz[111:104] = (qx[103: 96] * qy[111:104] + qx[111:104] * qy[103: 96]) >> SAR[5  :  0]
  qz[119:112] = (qx[119:112] * qy[119:112] - qx[127:120] * qy[127:120]) >> SAR[5  :  0]
  qz[127:120] = (qx[119:112] * qy[127:120] + qx[127:120] * qy[119:112]) >> SAR[5  :  0]
else if sel4 == 2:
  qz[7  :  0] = (qx[7  :  0] * qy[7  :  0] + qx[15 :  8] * qy[15 :  8]) >> SAR[5  :  0]
  qz[15 :  8] = (qx[7  :  0] * qy[15 :  8] - qx[15 :  8] * qy[7  :  0]) >> SAR[5  :  0]
  qz[23 : 16] = (qx[23 : 16] * qy[23 : 16] + qx[31 : 24] * qy[31 : 24]) >> SAR[5  :  0]
  qz[31 : 24] = (qx[23 : 16] * qy[31 : 24] - qx[31 : 24] * qy[23 : 16]) >> SAR[5 :  0]
  qz[39 : 32] = (qx[39 : 32] * qy[39 : 32] + qx[47 : 40] * qy[47 : 40]) >> SAR[5  :  0]
  qz[47 : 40] = (qx[39 : 32] * qy[47 : 40] - qx[47 : 40] * qy[39 : 32]) >> SAR[5  :  0]
  qz[55 : 48] = (qx[55 : 48] * qy[55 : 48] + qx[63 : 56] * qy[63 : 56]) >> SAR[5  :  0]
  qz[63 : 56] = (qx[55 : 48] * qy[63 : 56] - qx[63 : 56] * qy[55 : 48]) >> SAR[5  :  0]
else if sel4 == 3:
  qz[71 : 64] = (qx[71 : 64] * qy[71 : 64] + qx[79 : 72] * qy[79 : 72]) >> SAR[5  :  0]
  qz[79 : 72] = (qx[71 : 64] * qy[79 : 72] - qx[79 : 72] * qy[71 : 64]) >> SAR[5  :  0]
  qz[87 : 80] = (qx[87 : 80] * qy[87 : 80] + qx[95 : 88] * qy[95 : 88]) >> SAR[5  :  0]
  qz[95 : 88] = (qx[87 : 80] * qy[95 : 88] - qx[95 : 88] * qy[87 : 80]) >> SAR[5  :  0]
  qz[103: 96] = (qx[103: 96] * qy[103: 96] + qx[111:104] * qy[111:104]) >> SAR[5  :  0]
  qz[111:104] = (qx[103: 96] * qy[111:104] - qx[111:104] * qy[103: 96]) >> SAR[5  :  0]
  qz[119:112] = (qx[119:112] * qy[119:112] + qx[127:120] * qy[127:120]) >> SAR[5  :  0]
  qz[127:120] = (qx[119:112] * qy[127:120] - qx[127:120] * qy[119:112]) >> SAR[5  :  0]

  qu[127:0] => store128(rs1[31:0])
  rs1[31:0] = rs1[31:0] + 16
\end{lstlisting}


\textbf{Schedule}\\
\begin{longtable}[c]{|l|l|l|}
    \hline
    \rowcolor{lightgray}
    \textbf{Operands} & \textbf{Use} & \textbf{Def} \\\hline
    \endfirsthead
    \hline
    \rowcolor{lightgray}
     \textbf{Operands} & \textbf{Use} & \textbf{Def} \\\hline
    \endhead
    \hline
    \endfoot


    qz & - & 2 \\\hline
     qx & 1 & - \\\hline
     qy & 1 & - \\\hline
     qu & 2 & - \\\hline
     rs1 & 1 & 1

\end{longtable}

\textbf{Format}\\

\begin{figure*}[!htbp]
 {\setlength{\tabcolsep}{1pt}
 \begin{center}
\begin{tabular}{@{}F@{}F@{}F@{}I@{}W@{}I@{}I@{}F@{}W@{}F@{}F@{}O}
\instbitrange{31}{29}&
\instbitrange{28}{26}&
\instbitrange{25}{23}&
\instbit{22}&
\instbitrange{21}{20}&
\instbit{19}&
\instbit{18}&
\instbitrange{17}{15}&
\instbitrange{14}{13}&
\instbitrange{12}{10}&
\instbitrange{9}{7}&
\instbitrange{6}{0} \\
\hline
\multicolumn{1}{|c|}{{\scriptsize qx[2:0]}} &
\multicolumn{1}{c|}{{\scriptsize qy[2:0]}} &
\multicolumn{1}{c|}{101} &
\multicolumn{1}{c|}{{\scriptsize sat[0]}} &
\multicolumn{1}{c|}{{\scriptsize sel4[1:0]}} &
\multicolumn{1}{c|}{{\scriptsize rm[2]}} &
\multicolumn{1}{c|}{{\scriptsize rs1[4]}} &
\multicolumn{1}{c|}{{\scriptsize rs1[2:0]}} &
\multicolumn{1}{c|}{{\scriptsize rm[1:0]}} &
\multicolumn{1}{c|}{{\scriptsize qu[2:0]}} &
\multicolumn{1}{c|}{{\scriptsize qz[2:0]}} &
\multicolumn{1}{c|}{1111011}\\
\hline
3& 3& 3& 1& 2& 1& 1& 3& 2& 3& 3& 7 \\
\end{tabular}
 \end{center}
 }
 \vspace{-0.1in}
 \caption[]{ESP.CMUL.S8.ST.INCP} \end{figure*}


\newpage
\subsubsection{ESP.CMUL.U16}\label{instruction:ESPCMULU16}
\textbf{Syntax}\\
ESP.CMUL.U16 qz,qx,qy,sel4,sat,rm

\textbf{Description}\\
This instruction performs a 16-bit unsigned complex multiplication. The range of the immediate number \lstinline{sel4} is 0 \verb+~+ 3, which specifies the 32 bits in the two QR registers \lstinline{qx} and \lstinline{qy} for complex multiplication. The real and imaginary parts of the complex numbers are stored in the upper 16 bits and lower 16 bits of the 32 bits, respectively. The calculated real part and imaginary part results are stored in the corresponding 32 bits of the register \lstinline{qz}.

\textbf{Operation}\\
\begin{lstlisting}
if sel4 == 0:
  qz[ 15:  0] = (qx[ 15:  0] * qy[ 15:  0] - qx[ 31: 16] * qy[ 31: 16]) >> SAR[5:0]
  qz[ 31: 16] = (qx[ 15:  0] * qy[ 31: 16] + qx[ 31: 16] * qy[ 15:  0]) >> SAR[5:0]
  qz[ 47: 32] = (qx[ 47: 32] * qy[ 47: 32] - qx[ 63: 48] * qy[ 63: 48]) >> SAR[5:0]
  qz[ 63: 48] = (qx[ 47: 32] * qy[ 63: 48] + qx[ 63: 48] * qy[ 47: 32]) >> SAR[5:0]
if sel4 == 1:
  qz[ 79: 64] = (qx[ 79: 64] * qy[ 79: 64] - qx[ 95: 80] * qy[ 95: 80]) >> SAR[5:0]
  qz[ 95: 80] = (qx[ 79: 64] * qy[ 95: 80] + qx[ 95: 80] * qy[ 79: 64]) >> SAR[5:0]
  qz[111: 96] = (qx[111: 96] * qy[111: 96] - qx[127:112] * qy[127:112]) >> SAR[5:0]
  qz[127:112] = (qx[111: 96] * qy[127:112] + qx[127:112] * qy[111: 96]) >> SAR[5:0]
if sel4 == 2:
  qz[ 15:  0] = (qx[ 15:  0] * qy[ 15:  0] + qx[ 31: 16] * qy[ 31: 16]) >> SAR[5:0]
  qz[ 31: 16] = (qx[ 15:  0] * qy[ 31: 16] - qx[ 31: 16] * qy[ 15:  0]) >> SAR[5:0]
  qz[ 47: 32] = (qx[ 47: 32] * qy[ 47: 32] + qx[ 63: 48] * qy[ 63: 48]) >> SAR[5:0]
  qz[ 63: 48] = (qx[ 47: 32] * qy[ 63: 48] - qx[ 63: 48] * qy[ 47: 32]) >> SAR[5:0]
if sel4 == 3:
  qz[ 79: 64] = (qx[ 79: 64] * qy[ 79: 64] + qx[ 95: 80] * qy[ 95: 80]) >> SAR[5:0]
  qz[ 95: 80] = (qx[ 79: 64] * qy[ 95: 80] - qx[ 95: 80] * qy[ 79: 64]) >> SAR[5:0]
  qz[111: 96] = (qx[111: 96] * qy[111: 96] + qx[127:112] * qy[127:112]) >> SAR[5:0]
  qz[127:112] = (qx[111: 96] * qy[127:112] - qx[127:112] * qy[111: 96]) >> SAR[5:0]
\end{lstlisting}



\textbf{Schedule}\\
\begin{longtable}[c]{|l|l|l|}
    \hline
    \rowcolor{lightgray}
    \textbf{Operands} & \textbf{Use} & \textbf{Def} \\\hline
    \endfirsthead
    \hline
    \rowcolor{lightgray}
     \textbf{Operands} & \textbf{Use} & \textbf{Def} \\\hline
    \endhead
    \hline
    \endfoot


    qz & - & 2 \\\hline
     qx & 1 & - \\\hline
     qy & 1 & -

\end{longtable}

\textbf{Format}\\

\begin{figure*}[!htbp]
 {\setlength{\tabcolsep}{1pt}
 \begin{center}
\begin{tabular}{@{}F@{}F@{}E@{}I@{}W@{}W@{}F@{}F@{}O}
\instbitrange{31}{29}&
\instbitrange{28}{26}&
\instbitrange{25}{18}&
\instbit{17}&
\instbitrange{16}{15}&
\instbitrange{14}{13}&
\instbitrange{12}{10}&
\instbitrange{9}{7}&
\instbitrange{6}{0} \\
\hline
\multicolumn{1}{|c|}{{\scriptsize qx[2:0]}} &
\multicolumn{1}{c|}{{\scriptsize qy[2:0]}} &
\multicolumn{1}{c|}{111000X0} &
\multicolumn{1}{c|}{{\scriptsize sat[0]}} &
\multicolumn{1}{c|}{{\scriptsize sel4[1:0]}} &
\multicolumn{1}{c|}{00} &
\multicolumn{1}{c|}{{\scriptsize rm[2:0]}} &
\multicolumn{1}{c|}{{\scriptsize qz[2:0]}} &
\multicolumn{1}{c|}{0011111}\\
\hline
3& 3& 8& 1& 2& 2& 3& 3& 7 \\
\end{tabular}
 \end{center}
 }
 \vspace{-0.1in}
 \caption[]{ESP.CMUL.U16} \end{figure*}


\newpage
\subsubsection{ESP.CMUL.U16.LD.INCP}\label{instruction:ESPCMULU16LDINCP}
\textbf{Syntax}\\
ESP.CMUL.U16.LD.INCP qu,rs1,qz,qx,qy,sel4,sat,rm

\textbf{Description}\\
This instruction performs a 16-bit unsigned complex multiplication. The range of the immediate number \lstinline{sel4} is 0 \verb+~+ 3, which specifies the 32 bits in the two QR registers \lstinline{qx} and \lstinline{qy} for complex multiplication. The real and imaginary parts of complex numbers are stored in the upper 16 bits and lower 16 bits of the 32 bits respectively. The calculated real part and imaginary part results are stored in the corresponding 32 bits of the register \lstinline{qz}.\\
Meanwhile, it also loads 16-byte data from the memory into the register \lstinline{qu}. After the access, the value in the register \lstinline{rs1} is incremented by 16.

\textbf{Operation}\\
\begin{lstlisting}
if sel4 == 0:
  qz[ 15:  0] = (qx[ 15:  0] * qy[ 15:  0] - qx[ 31: 16] * qy[ 31: 16]) >> SAR[5:0]
  qz[ 31: 16] = (qx[ 15:  0] * qy[ 31: 16] + qx[ 31: 16] * qy[ 15:  0]) >> SAR[5:0]
  qz[ 47: 32] = (qx[ 47: 32] * qy[ 47: 32] - qx[ 63: 48] * qy[ 63: 48]) >> SAR[5:0]
  qz[ 63: 48] = (qx[ 47: 32] * qy[ 63: 48] + qx[ 63: 48] * qy[ 47: 32]) >> SAR[5:0]
if sel4 == 1:
  qz[ 79: 64] = (qx[ 79: 64] * qy[ 79: 64] - qx[ 95: 80] * qy[ 95: 80]) >> SAR[5:0]
  qz[ 95: 80] = (qx[ 79: 64] * qy[ 95: 80] + qx[ 95: 80] * qy[ 79: 64]) >> SAR[5:0]
  qz[111: 96] = (qx[111: 96] * qy[111: 96] - qx[127:112] * qy[127:112]) >> SAR[5:0]
  qz[127:112] = (qx[111: 96] * qy[127:112] + qx[127:112] * qy[111: 96]) >> SAR[5:0]
if sel4 == 2:
  qz[ 15:  0] = (qx[ 15:  0] * qy[ 15:  0] + qx[ 31: 16] * qy[ 31: 16]) >> SAR[5:0]
  qz[ 31: 16] = (qx[ 15:  0] * qy[ 31: 16] - qx[ 31: 16] * qy[ 15:  0]) >> SAR[5:0]
  qz[ 47: 32] = (qx[ 47: 32] * qy[ 47: 32] + qx[ 63: 48] * qy[ 63: 48]) >> SAR[5:0]
  qz[ 63: 48] = (qx[ 47: 32] * qy[ 63: 48] - qx[ 63: 48] * qy[ 47: 32]) >> SAR[5:0]
if sel4 == 3:
  qz[ 79: 64] = (qx[ 79: 64] * qy[ 79: 64] + qx[ 95: 80] * qy[ 95: 80]) >> SAR[5:0]
  qz[ 95: 80] = (qx[ 79: 64] * qy[ 95: 80] - qx[ 95: 80] * qy[ 79: 64]) >> SAR[5:0]
  qz[111: 96] = (qx[111: 96] * qy[111: 96] + qx[127:112] * qy[127:112]) >> SAR[5:0]
  qz[127:112] = (qx[111: 96] * qy[127:112] - qx[127:112] * qy[111: 96]) >> SAR[5:0]

  qu[127:0] = load128(rs1[31:0])
  rs1[31:0] = rs1[31:0] + 16
\end{lstlisting}

\textbf{Schedule}\\
\begin{longtable}[c]{|l|l|l|}
    \hline
    \rowcolor{lightgray}
    \textbf{Operands} & \textbf{Use} & \textbf{Def} \\\hline
    \endfirsthead
    \hline
    \rowcolor{lightgray}
     \textbf{Operands} & \textbf{Use} & \textbf{Def} \\\hline
    \endhead
    \hline
    \endfoot


    qz & - & 2 \\\hline
     qx & 1 & - \\\hline
     qy & 1 & - \\\hline
     qu & - & 2 \\\hline
     rs1 & 1 & 1

\end{longtable}

\textbf{Format}\\

\begin{figure*}[!htbp]
 {\setlength{\tabcolsep}{1pt}
 \begin{center}
\begin{tabular}{@{}F@{}F@{}F@{}I@{}W@{}I@{}I@{}F@{}W@{}F@{}F@{}O}
\instbitrange{31}{29}&
\instbitrange{28}{26}&
\instbitrange{25}{23}&
\instbit{22}&
\instbitrange{21}{20}&
\instbit{19}&
\instbit{18}&
\instbitrange{17}{15}&
\instbitrange{14}{13}&
\instbitrange{12}{10}&
\instbitrange{9}{7}&
\instbitrange{6}{0} \\
\hline
\multicolumn{1}{|c|}{{\scriptsize qx[2:0]}} &
\multicolumn{1}{c|}{{\scriptsize qy[2:0]}} &
\multicolumn{1}{c|}{010} &
\multicolumn{1}{c|}{{\scriptsize sat[0]}} &
\multicolumn{1}{c|}{{\scriptsize sel4[1:0]}} &
\multicolumn{1}{c|}{{\scriptsize rm[2]}} &
\multicolumn{1}{c|}{{\scriptsize rs1[4]}} &
\multicolumn{1}{c|}{{\scriptsize rs1[2:0]}} &
\multicolumn{1}{c|}{{\scriptsize rm[1:0]}} &
\multicolumn{1}{c|}{{\scriptsize qu[2:0]}} &
\multicolumn{1}{c|}{{\scriptsize qz[2:0]}} &
\multicolumn{1}{c|}{1111011}\\
\hline
3& 3& 3& 1& 2& 1& 1& 3& 2& 3& 3& 7 \\
\end{tabular}
 \end{center}
 }
 \vspace{-0.1in}
 \caption[]{ESP.CMUL.U16.LD.INCP} \end{figure*}


\newpage
\subsubsection{ESP.CMUL.U16.ST.INCP}\label{instruction:ESPCMULU16STINCP}
\textbf{Syntax}\\
ESP.CMUL.U16.ST.INCP qu,rs1,qz,qx,qy,sel4,sat,rm

\textbf{Description}\\
This instruction performs a 16-bit unsigned complex multiplication. The range of the immediate number \lstinline{sel4} is 0 \verb+~+ 3, which specifies the 32 bits in the two QR registers \lstinline{qx} and \lstinline{qy} for complex multiplication. The real and imaginary parts of complex numbers are stored in the upper 16 bits and lower 16 bits of the 32 bits, respectively. The calculated real part and imaginary part results are stored in the corresponding 32 bits of the register \lstinline{qz}.

Meanwhile, it also stores the 16-byte data of \lstinline{qu} in memory. After the access, the value in the register \lstinline{rs1} is incremented by 16.

\textbf{Operation}\\
\begin{lstlisting}
if sel4 == 0:
  qz[ 15:  0] = (qx[ 15:  0] * qy[ 15:  0] - qx[ 31: 16] * qy[ 31: 16]) >> SAR[5:0]
  qz[ 31: 16] = (qx[ 15:  0] * qy[ 31: 16] + qx[ 31: 16] * qy[ 15:  0]) >> SAR[5:0]
  qz[ 47: 32] = (qx[ 47: 32] * qy[ 47: 32] - qx[ 63: 48] * qy[ 63: 48]) >> SAR[5:0]
  qz[ 63: 48] = (qx[ 47: 32] * qy[ 63: 48] + qx[ 63: 48] * qy[ 47: 32]) >> SAR[5:0]
if sel4 == 1:
  qz[ 79: 64] = (qx[ 79: 64] * qy[ 79: 64] - qx[ 95: 80] * qy[ 95: 80]) >> SAR[5:0]
  qz[ 95: 80] = (qx[ 79: 64] * qy[ 95: 80] + qx[ 95: 80] * qy[ 79: 64]) >> SAR[5:0]
  qz[111: 96] = (qx[111: 96] * qy[111: 96] - qx[127:112] * qy[127:112]) >> SAR[5:0]
  qz[127:112] = (qx[111: 96] * qy[127:112] + qx[127:112] * qy[111: 96]) >> SAR[5:0]
if sel4 == 2:
  qz[ 15:  0] = (qx[ 15:  0] * qy[ 15:  0] + qx[ 31: 16] * qy[ 31: 16]) >> SAR[5:0]
  qz[ 31: 16] = (qx[ 15:  0] * qy[ 31: 16] - qx[ 31: 16] * qy[ 15:  0]) >> SAR[5:0]
  qz[ 47: 32] = (qx[ 47: 32] * qy[ 47: 32] + qx[ 63: 48] * qy[ 63: 48]) >> SAR[5:0]
  qz[ 63: 48] = (qx[ 47: 32] * qy[ 63: 48] - qx[ 63: 48] * qy[ 47: 32]) >> SAR[5:0]
if sel4 == 3:
  qz[ 79: 64] = (qx[ 79: 64] * qy[ 79: 64] + qx[ 95: 80] * qy[ 95: 80]) >> SAR[5:0]
  qz[ 95: 80] = (qx[ 79: 64] * qy[ 95: 80] - qx[ 95: 80] * qy[ 79: 64]) >> SAR[5:0]
  qz[111: 96] = (qx[111: 96] * qy[111: 96] + qx[127:112] * qy[127:112]) >> SAR[5:0]
  qz[127:112] = (qx[111: 96] * qy[127:112] - qx[127:112] * qy[111: 96]) >> SAR[5:0]

  qu[127:0] => store128(rs1[31:0])
  rs1[31:0] = rs1[31:0] + 16
\end{lstlisting}

\textbf{Schedule}\\
\begin{longtable}[c]{|l|l|l|}
    \hline
    \rowcolor{lightgray}
    \textbf{Operands} & \textbf{Use} & \textbf{Def} \\\hline
    \endfirsthead
    \hline
    \rowcolor{lightgray}
     \textbf{Operands} & \textbf{Use} & \textbf{Def} \\\hline
    \endhead
    \hline
    \endfoot


    qz & - & 2 \\\hline
     qx & 1 & - \\\hline
     qy & 1 & - \\\hline
     qu & 2 & - \\\hline
     rs1 & 1 & 1

\end{longtable}

\textbf{Format}\\

\begin{figure*}[!htbp]
 {\setlength{\tabcolsep}{1pt}
 \begin{center}
\begin{tabular}{@{}F@{}F@{}F@{}I@{}W@{}I@{}I@{}F@{}W@{}F@{}F@{}O}
\instbitrange{31}{29}&
\instbitrange{28}{26}&
\instbitrange{25}{23}&
\instbit{22}&
\instbitrange{21}{20}&
\instbit{19}&
\instbit{18}&
\instbitrange{17}{15}&
\instbitrange{14}{13}&
\instbitrange{12}{10}&
\instbitrange{9}{7}&
\instbitrange{6}{0} \\
\hline
\multicolumn{1}{|c|}{{\scriptsize qx[2:0]}} &
\multicolumn{1}{c|}{{\scriptsize qy[2:0]}} &
\multicolumn{1}{c|}{110} &
\multicolumn{1}{c|}{{\scriptsize sat[0]}} &
\multicolumn{1}{c|}{{\scriptsize sel4[1:0]}} &
\multicolumn{1}{c|}{{\scriptsize rm[2]}} &
\multicolumn{1}{c|}{{\scriptsize rs1[4]}} &
\multicolumn{1}{c|}{{\scriptsize rs1[2:0]}} &
\multicolumn{1}{c|}{{\scriptsize rm[1:0]}} &
\multicolumn{1}{c|}{{\scriptsize qu[2:0]}} &
\multicolumn{1}{c|}{{\scriptsize qz[2:0]}} &
\multicolumn{1}{c|}{1111011}\\
\hline
3& 3& 3& 1& 2& 1& 1& 3& 2& 3& 3& 7 \\
\end{tabular}
 \end{center}
 }
 \vspace{-0.1in}
 \caption[]{ESP.CMUL.U16.ST.INCP} \end{figure*}


\newpage
\subsubsection{ESP.CMUL.U8}\label{instruction:ESPCMULU8}
\textbf{Syntax}\\
ESP.CMUL.U8 qz,qx,qy,sel4,sat,rm

\textbf{Description}\\
This instruction performs an 8-bit unsigned complex multiplication. The range of the immediate number \lstinline{sel4} is 0 \verb+~+ 3, which specifies the 16 bits in the two QR registers \lstinline{qx} and \lstinline{qy} for complex multiplication. The real and imaginary parts of complex numbers are stored in the upper 8 bits and lower 8 bits of the 16 bits respectively. The calculated real part and imaginary part results are stored in the corresponding 16 bits of the register \lstinline{qz}.

\textbf{Operation}\\
\begin{lstlisting}
if sel4 == 0:
  qz[  7:  0] = (qx[  7:  0] * qy[  7:  0] - qx[ 15:  8] * qy[ 15:  8]) >> SAR[5  :  0]
  qz[ 15:  8] = (qx[  7:  0] * qy[ 15:  8] + qx[ 15:  8] * qy[  7:  0]) >> SAR[5  :  0]
  qz[ 23: 16] = (qx[ 23: 16] * qy[ 23: 16] - qx[ 31: 16] * qy[ 31: 24]) >> SAR[5  :  0]
  qz[ 31: 24] = (qx[ 23: 16] * qy[ 31: 24] + qx[ 31: 16] * qy[ 23: 16]) >> SAR[5  :  0]
  qz[ 39: 32] = (qx[ 39: 32] * qy[ 39: 32] - qx[ 47: 40] * qy[ 47: 40]) >> SAR[5  :  0]
  qz[47 : 40] = (qx[39 : 32] * qy[47 : 40] + qx[47 : 40] * qy[39 : 32]) >> SAR[5  :  0]
  qz[55 : 48] = (qx[55 : 48] * qy[55 : 48] - qx[63 : 56] * qy[63 : 56]) >> SAR[5  :  0]
  qz[63 : 56] = (qx[55 : 48] * qy[63 : 56] + qx[63 : 56] * qy[55 : 48]) >> SAR[5  :  0]
else if sel4 == 1:
  qz[71 : 64] = (qx[71 : 64] * qy[71 : 64] - qx[79 : 72] * qy[79 : 72]) >> SAR[5  :  0]
  qz[79 : 72] = (qx[71 : 64] * qy[79 : 72] + qx[79 : 72] * qy[71 : 64]) >> SAR[5  :  0]
  qz[87 : 80] = (qx[87 : 80] * qy[87 : 80] - qx[95 : 88] * qy[95 : 88]) >> SAR[5  :  0]
  qz[95 : 88] = (qx[87 : 80] * qy[95 : 88] + qx[95 : 88] * qy[87 : 80]) >> SAR[5  :  0]
  qz[103: 96] = (qx[103: 96] * qy[103: 96] - qx[111:104] * qy[111:104]) >> SAR[5  :  0]
  qz[111:104] = (qx[103: 96] * qy[111:104] + qx[111:104] * qy[103: 96]) >> SAR[5  :  0]
  qz[119:112] = (qx[119:112] * qy[119:112] - qx[127:120] * qy[127:120]) >> SAR[5  :  0]
  qz[127:120] = (qx[119:112] * qy[127:120] + qx[127:120] * qy[119:112]) >> SAR[5  :  0]
else if sel4 == 2:
  qz[7  :  0] = (qx[7  :  0] * qy[7  :  0] + qx[15 :  8] * qy[15 :  8]) >> SAR[5  :  0]
  qz[15 :  8] = (qx[7  :  0] * qy[15 :  8] - qx[15 :  8] * qy[7  :  0]) >> SAR[5  :  0]
  qz[23 : 16] = (qx[23 : 16] * qy[23 : 16] + qx[31 : 24] * qy[31 : 24]) >> SAR[5  :  0]
  qz[31 : 24] = (qx[23 : 16] * qy[31 : 24] - qx[31 : 24] * qy[23 : 16]) >> SAR[5  :  0]
  qz[39 : 32] = (qx[39 : 32] * qy[39 : 32] + qx[47 : 40] * qy[47 : 40]) >> SAR[5  :  0]
  qz[47 : 40] = (qx[39 : 32] * qy[47 : 40] - qx[47 : 40] * qy[39 : 32]) >> SAR[5  :  0]
  qz[55 : 48] = (qx[55 : 48] * qy[55 : 48] + qx[63 : 56] * qy[63 : 56]) >> SAR[5  :  0]
  qz[63 : 56] = (qx[55 : 48] * qy[63 : 56] - qx[63 : 56] * qy[55 : 48]) >> SAR[5  :  0]
else if sel4 == 3:
  qz[71 : 64] = (qx[71 : 64] * qy[71 : 64] + qx[79 : 72] * qy[79 : 72]) >> SAR[5  :  0]
  qz[79 : 72] = (qx[71 : 64] * qy[79 : 72] - qx[79 : 72] * qy[71 : 64]) >> SAR[5  :  0]
  qz[87 : 80] = (qx[87 : 80] * qy[87 : 80] + qx[95 : 88] * qy[95 : 88]) >> SAR[5  :  0]
  qz[95 : 88] = (qx[87 : 80] * qy[95 : 88] - qx[95 : 88] * qy[87 : 80]) >> SAR[5  :  0]
  qz[103: 96] = (qx[103: 96] * qy[103: 96] + qx[111:104] * qy[111:104]) >> SAR[5  :  0]
  qz[111:104] = (qx[103: 96] * qy[111:104] - qx[111:104] * qy[103: 96]) >> SAR[5  :  0]
  qz[119:112] = (qx[119:112] * qy[119:112] + qx[127:120] * qy[127:120]) >> SAR[5  :  0]
  qz[127:120] = (qx[119:112] * qy[127:120] - qx[127:120] * qy[119:112]) >> SAR[5  :  0]

\end{lstlisting}

\textbf{Schedule}\\
\begin{longtable}[c]{|l|l|l|}
    \hline
    \rowcolor{lightgray}
    \textbf{Operands} & \textbf{Use} & \textbf{Def} \\\hline
    \endfirsthead
    \hline
    \rowcolor{lightgray}
     \textbf{Operands} & \textbf{Use} & \textbf{Def} \\\hline
    \endhead
    \hline
    \endfoot


    qz & - & 2 \\\hline
     qx & 1 & - \\\hline
     qy & 1 & -

\end{longtable}

\textbf{Format}\\

\begin{figure*}[!htbp]
 {\setlength{\tabcolsep}{1pt}
 \begin{center}
\begin{tabular}{@{}F@{}F@{}E@{}I@{}W@{}W@{}F@{}F@{}O}
\instbitrange{31}{29}&
\instbitrange{28}{26}&
\instbitrange{25}{18}&
\instbit{17}&
\instbitrange{16}{15}&
\instbitrange{14}{13}&
\instbitrange{12}{10}&
\instbitrange{9}{7}&
\instbitrange{6}{0} \\
\hline
\multicolumn{1}{|c|}{{\scriptsize qx[2:0]}} &
\multicolumn{1}{c|}{{\scriptsize qy[2:0]}} &
\multicolumn{1}{c|}{110000X0} &
\multicolumn{1}{c|}{{\scriptsize sat[0]}} &
\multicolumn{1}{c|}{{\scriptsize sel4[1:0]}} &
\multicolumn{1}{c|}{00} &
\multicolumn{1}{c|}{{\scriptsize rm[2:0]}} &
\multicolumn{1}{c|}{{\scriptsize qz[2:0]}} &
\multicolumn{1}{c|}{0011111}\\
\hline
3& 3& 8& 1& 2& 2& 3& 3& 7 \\
\end{tabular}
 \end{center}
 }
 \vspace{-0.1in}
 \caption[]{ESP.CMUL.U8} \end{figure*}


\newpage
\subsubsection{ESP.CMUL.U8.LD.INCP}\label{instruction:ESPCMULU8LDINCP}
\textbf{Syntax}\\
ESP.CMUL.U8.LD.INCP qu,rs1,qz,qx,qy,sel4,sat,rm

\textbf{Description}\\
This instruction performs an 8-bit unsigned complex multiplication. The range of the immediate number \lstinline{sel4} is 0 \verb+~+ 3, which specifies the 16 bits in the two QR registers \lstinline{qx} and \lstinline{qy} for complex multiplication. The real and imaginary parts of complex numbers are stored in the upper 8 bits and lower 8 bits of the 16 bits respectively. The calculated real part and imaginary part results are stored in the corresponding 16 bits of the register \lstinline{qz}.

Meanwhile, it also loads 16-byte data from the memory into the register \lstinline{qu}. After the access, the value in the register \lstinline{rs1} is incremented by 16.\\
\textbf{Operation}\\
\begin{lstlisting}
if sel4 == 0:
  qz[  7:  0] = (qx[  7:  0] * qy[  7:  0] - qx[ 15:  8] * qy[ 15:  8]) >> SAR[5  :  0]
  qz[ 15:  8] = (qx[  7:  0] * qy[ 15:  8] + qx[ 15:  8] * qy[  7:  0]) >> SAR[5  :  0]
  qz[ 23: 16] = (qx[ 23: 16] * qy[ 23: 16] - qx[ 31: 16] * qy[ 31: 24]) >> SAR[5  :  0]
  qz[ 31: 24] = (qx[ 23: 16] * qy[ 31: 24] + qx[ 31: 16] * qy[ 23: 16]) >> SAR[5  :  0]
  qz[ 39: 32] = (qx[ 39: 32] * qy[ 39: 32] - qx[ 47: 40] * qy[ 47: 40]) >> SAR[5  :  0]
  qz[47 : 40] = (qx[39 : 32] * qy[47 : 40] + qx[47 : 40] * qy[39 : 32]) >> SAR[5  :  0]
  qz[55 : 48] = (qx[55 : 48] * qy[55 : 48] - qx[63 : 56] * qy[63 : 56]) >> SAR[5  :  0]
  qz[63 : 56] = (qx[55 : 48] * qy[63 : 56] + qx[63 : 56] * qy[55 : 48]) >> SAR[5  :  0]
else if sel4 == 1:
  qz[71 : 64] = (qx[71 : 64] * qy[71 : 64] - qx[79 : 72] * qy[79 : 72]) >> SAR[5  :  0]
  qz[79 : 72] = (qx[71 : 64] * qy[79 : 72] + qx[79 : 72] * qy[71 : 64]) >> SAR[5  :  0]
  qz[87 : 80] = (qx[87 : 80] * qy[87 : 80] - qx[95 : 88] * qy[95 : 88]) >> SAR[5  :  0]
  qz[95 : 88] = (qx[87 : 80] * qy[95 : 88] + qx[95 : 88] * qy[87 : 80]) >> SAR[5  :  0]
  qz[103: 96] = (qx[103: 96] * qy[103: 96] - qx[111:104] * qy[111:104]) >> SAR[5  :  0]
  qz[111:104] = (qx[103: 96] * qy[111:104] + qx[111:104] * qy[103: 96]) >> SAR[5  :  0]
  qz[119:112] = (qx[119:112] * qy[119:112] - qx[127:120] * qy[127:120]) >> SAR[5  :  0]
  qz[127:120] = (qx[119:112] * qy[127:120] + qx[127:120] * qy[119:112]) >> SAR[5  :  0]
else if sel4 == 2:
  qz[7  :  0] = (qx[7  :  0] * qy[7  :  0] + qx[15 :  8] * qy[15 :  8]) >> SAR[5  :  0]
  qz[15 :  8] = (qx[7  :  0] * qy[15 :  8] - qx[15 :  8] * qy[7  :  0]) >> SAR[5  :  0]
  qz[23 : 16] = (qx[23 : 16] * qy[23 : 16] + qx[31 : 24] * qy[31 : 24]) >> SAR[5  :  0]
  qz[31 : 24] = (qx[23 : 16] * qy[31 : 24] - qx[31 : 24] * qy[23 : 16]) >> SAR[5  :  0]
  qz[39 : 32] = (qx[39 : 32] * qy[39 : 32] + qx[47 : 40] * qy[47 : 40]) >> SAR[5  :  0]
  qz[47 : 40] = (qx[39 : 32] * qy[47 : 40] - qx[47 : 40] * qy[39 : 32]) >> SAR[5  :  0]
  qz[55 : 48] = (qx[55 : 48] * qy[55 : 48] + qx[63 : 56] * qy[63 : 56]) >> SAR[5  :  0]
  qz[63 : 56] = (qx[55 : 48] * qy[63 : 56] - qx[63 : 56] * qy[55 : 48]) >> SAR[5  :  0]
else if sel4 == 3:
  qz[71 : 64] = (qx[71 : 64] * qy[71 : 64] + qx[79 : 72] * qy[79 : 72]) >> SAR[5  :  0]
  qz[79 : 72] = (qx[71 : 64] * qy[79 : 72] - qx[79 : 72] * qy[71 : 64]) >> SAR[5  :  0]
  qz[87 : 80] = (qx[87 : 80] * qy[87 : 80] + qx[95 : 88] * qy[95 : 88]) >> SAR[5  :  0]
  qz[95 : 88] = (qx[87 : 80] * qy[95 : 88] - qx[95 : 88] * qy[87 : 80]) >> SAR[5  :  0]
  qz[103: 96] = (qx[103: 96] * qy[103: 96] + qx[111:104] * qy[111:104]) >> SAR[5  :  0]
  qz[111:104] = (qx[103: 96] * qy[111:104] - qx[111:104] * qy[103: 96]) >> SAR[5  :  0]
  qz[119:112] = (qx[119:112] * qy[119:112] + qx[127:120] * qy[127:120]) >> SAR[5  :  0]
  qz[127:120] = (qx[119:112] * qy[127:120] - qx[127:120] * qy[119:112]) >> SAR[5  :  0]

qu[127:0] = load128(rs1[31:0])
rs1[31:0] = rs1[31:0] + 16

\end{lstlisting}


\textbf{Schedule}\\
\begin{longtable}[c]{|l|l|l|}
    \hline
    \rowcolor{lightgray}
    \textbf{Operands} & \textbf{Use} & \textbf{Def} \\\hline
    \endfirsthead
    \hline
    \rowcolor{lightgray}
     \textbf{Operands} & \textbf{Use} & \textbf{Def} \\\hline
    \endhead
    \hline
    \endfoot


    qz & - & 2 \\\hline
     qx & 1 & - \\\hline
     qy & 1 & - \\\hline
     qu & - & 2 \\\hline
     rs1 & 1 & 1

\end{longtable}

\textbf{Format}\\

\begin{figure*}[!htbp]
 {\setlength{\tabcolsep}{1pt}
 \begin{center}
\begin{tabular}{@{}F@{}F@{}F@{}I@{}W@{}I@{}I@{}F@{}W@{}F@{}F@{}O}
\instbitrange{31}{29}&
\instbitrange{28}{26}&
\instbitrange{25}{23}&
\instbit{22}&
\instbitrange{21}{20}&
\instbit{19}&
\instbit{18}&
\instbitrange{17}{15}&
\instbitrange{14}{13}&
\instbitrange{12}{10}&
\instbitrange{9}{7}&
\instbitrange{6}{0} \\
\hline
\multicolumn{1}{|c|}{{\scriptsize qx[2:0]}} &
\multicolumn{1}{c|}{{\scriptsize qy[2:0]}} &
\multicolumn{1}{c|}{000} &
\multicolumn{1}{c|}{{\scriptsize sat[0]}} &
\multicolumn{1}{c|}{{\scriptsize sel4[1:0]}} &
\multicolumn{1}{c|}{{\scriptsize rm[2]}} &
\multicolumn{1}{c|}{{\scriptsize rs1[4]}} &
\multicolumn{1}{c|}{{\scriptsize rs1[2:0]}} &
\multicolumn{1}{c|}{{\scriptsize rm[1:0]}} &
\multicolumn{1}{c|}{{\scriptsize qu[2:0]}} &
\multicolumn{1}{c|}{{\scriptsize qz[2:0]}} &
\multicolumn{1}{c|}{1111011}\\
\hline
3& 3& 3& 1& 2& 1& 1& 3& 2& 3& 3& 7 \\
\end{tabular}
 \end{center}
 }
 \vspace{-0.1in}
 \caption[]{ESP.CMUL.U8.LD.INCP} \end{figure*}


\newpage
\subsubsection{ESP.CMUL.U8.ST.INCP}\label{instruction:ESPCMULU8STINCP}
\textbf{Syntax}\\
ESP.CMUL.U8.ST.INCP qu,rs1,qz,qx,qy,sel4,sat,rm

\textbf{Description}\\
This instruction performs an 8-bit unsigned complex multiplication. The range of the immediate number \lstinline{sel4} is 0 \verb+~+ 3, which specifies the 16 bits in the two QR registers \lstinline{qx} and \lstinline{qy} for complex multiplication. The real and imaginary parts of complex numbers are stored in the upper 8 bits and lower 8 bits of the 16 bits respectively. The calculated real part and imaginary part results are stored in the corresponding 16 bits of the register \lstinline{qz}.

Meanwhile, it also stores the 16-byte data of \lstinline{qu} in memory. After the access, the value in the register \lstinline{rs1} is incremented by 16.\\
\textbf{Operation}\\
\begin{lstlisting}
if sel4 == 0:
  qz[  7:  0] = (qx[  7:  0] * qy[  7:  0] - qx[ 15:  8] * qy[ 15:  8]) >> SAR[5  :  0]
  qz[ 15:  8] = (qx[  7:  0] * qy[ 15:  8] + qx[ 15:  8] * qy[  7:  0]) >> SAR[5  :  0]
  qz[ 23: 16] = (qx[ 23: 16] * qy[ 23: 16] - qx[ 31: 16] * qy[ 31: 24]) >> SAR[5  :  0]
  qz[ 31: 24] = (qx[ 23: 16] * qy[ 31: 24] + qx[ 31: 16] * qy[ 23: 16]) >> SAR[5  :  0]
  qz[ 39: 32] = (qx[ 39: 32] * qy[ 39: 32] - qx[ 47: 40] * qy[ 47: 40]) >> SAR[5  :  0]
  qz[47 : 40] = (qx[39 : 32] * qy[47 : 40] + qx[47 : 40] * qy[39 : 32]) >> SAR[5  :  0]
  qz[55 : 48] = (qx[55 : 48] * qy[55 : 48] - qx[63 : 56] * qy[63 : 56]) >> SAR[5  :  0]
  qz[63 : 56] = (qx[55 : 48] * qy[63 : 56] + qx[63 : 56] * qy[55 : 48]) >> SAR[5  :  0]
else if sel4 == 1:
  qz[71 : 64] = (qx[71 : 64] * qy[71 : 64] - qx[79 : 72] * qy[79 : 72]) >> SAR[5  :  0]
  qz[79 : 72] = (qx[71 : 64] * qy[79 : 72] + qx[79 : 72] * qy[71 : 64]) >> SAR[5  :  0]
  qz[87 : 80] = (qx[87 : 80] * qy[87 : 80] - qx[95 : 88] * qy[95 : 88]) >> SAR[5  :  0]
  qz[95 : 88] = (qx[87 : 80] * qy[95 : 88] + qx[95 : 88] * qy[87 : 80]) >> SAR[5  :  0]
  qz[103: 96] = (qx[103: 96] * qy[103: 96] - qx[111:104] * qy[111:104]) >> SAR[5  :  0]
  qz[111:104] = (qx[103: 96] * qy[111:104] + qx[111:104] * qy[103: 96]) >> SAR[5  :  0]
  qz[119:112] = (qx[119:112] * qy[119:112] - qx[127:120] * qy[127:120]) >> SAR[5  :  0]
  qz[127:120] = (qx[119:112] * qy[127:120] + qx[127:120] * qy[119:112]) >> SAR[5  :  0]
else if sel4 == 2:
  qz[7  :  0] = (qx[7  :  0] * qy[7  :  0] + qx[15 :  8] * qy[15 :  8]) >> SAR[5  :  0]
  qz[15 :  8] = (qx[7  :  0] * qy[15 :  8] - qx[15 :  8] * qy[7  :  0]) >> SAR[5  :  0]
  qz[23 : 16] = (qx[23 : 16] * qy[23 : 16] + qx[31 : 24] * qy[31 : 24]) >> SAR[5  :  0]
  qz[31 : 24] = (qx[23 : 16] * qy[31 : 24] - qx[31 : 24] * qy[23 : 16]) >> SAR[5  :  0]
  qz[39 : 32] = (qx[39 : 32] * qy[39 : 32] + qx[47 : 40] * qy[47 : 40]) >> SAR[5  :  0]
  qz[47 : 40] = (qx[39 : 32] * qy[47 : 40] - qx[47 : 40] * qy[39 : 32]) >> SAR[5  :  0]
  qz[55 : 48] = (qx[55 : 48] * qy[55 : 48] + qx[63 : 56] * qy[63 : 56]) >> SAR[5  :  0]
  qz[63 : 56] = (qx[55 : 48] * qy[63 : 56] - qx[63 : 56] * qy[55 : 48]) >> SAR[5  :  0]
else if sel4 == 3:
  qz[71 : 64] = (qx[71 : 64] * qy[71 : 64] + qx[79 : 72] * qy[79 : 72]) >> SAR[5  :  0]
  qz[79 : 72] = (qx[71 : 64] * qy[79 : 72] - qx[79 : 72] * qy[71 : 64]) >> SAR[5  :  0]
  qz[87 : 80] = (qx[87 : 80] * qy[87 : 80] + qx[95 : 88] * qy[95 : 88]) >> SAR[5  :  0]
  qz[95 : 88] = (qx[87 : 80] * qy[95 : 88] - qx[95 : 88] * qy[87 : 80]) >> SAR[5  :  0]
  qz[103: 96] = (qx[103: 96] * qy[103: 96] + qx[111:104] * qy[111:104]) >> SAR[5  :  0]
  qz[111:104] = (qx[103: 96] * qy[111:104] - qx[111:104] * qy[103: 96]) >> SAR[5  :  0]
  qz[119:112] = (qx[119:112] * qy[119:112] + qx[127:120] * qy[127:120]) >> SAR[5  :  0]
  qz[127:120] = (qx[119:112] * qy[127:120] - qx[127:120] * qy[119:112]) >> SAR[5  :  0]

  qu[127:0] => store128(rs1[31:0])
  rs1[31:0] = rs1[31:0] + 16
\end{lstlisting}

\textbf{Schedule}\\
\begin{longtable}[c]{|l|l|l|}
    \hline
    \rowcolor{lightgray}
    \textbf{Operands} & \textbf{Use} & \textbf{Def} \\\hline
    \endfirsthead
    \hline
    \rowcolor{lightgray}
     \textbf{Operands} & \textbf{Use} & \textbf{Def} \\\hline
    \endhead
    \hline
    \endfoot


    qz & - & 2 \\\hline
     qx & 1 & - \\\hline
     qy & 1 & - \\\hline
     qu & 2 & - \\\hline
     rs1 & 1 & 1

\end{longtable}

\textbf{Format}\\

\begin{figure*}[!htbp]
 {\setlength{\tabcolsep}{1pt}
 \begin{center}
\begin{tabular}{@{}F@{}F@{}F@{}I@{}W@{}I@{}I@{}F@{}W@{}F@{}F@{}O}
\instbitrange{31}{29}&
\instbitrange{28}{26}&
\instbitrange{25}{23}&
\instbit{22}&
\instbitrange{21}{20}&
\instbit{19}&
\instbit{18}&
\instbitrange{17}{15}&
\instbitrange{14}{13}&
\instbitrange{12}{10}&
\instbitrange{9}{7}&
\instbitrange{6}{0} \\
\hline
\multicolumn{1}{|c|}{{\scriptsize qx[2:0]}} &
\multicolumn{1}{c|}{{\scriptsize qy[2:0]}} &
\multicolumn{1}{c|}{100} &
\multicolumn{1}{c|}{{\scriptsize sat[0]}} &
\multicolumn{1}{c|}{{\scriptsize sel4[1:0]}} &
\multicolumn{1}{c|}{{\scriptsize rm[2]}} &
\multicolumn{1}{c|}{{\scriptsize rs1[4]}} &
\multicolumn{1}{c|}{{\scriptsize rs1[2:0]}} &
\multicolumn{1}{c|}{{\scriptsize rm[1:0]}} &
\multicolumn{1}{c|}{{\scriptsize qu[2:0]}} &
\multicolumn{1}{c|}{{\scriptsize qz[2:0]}} &
\multicolumn{1}{c|}{1111011}\\
\hline
3& 3& 3& 1& 2& 1& 1& 3& 2& 3& 3& 7 \\
\end{tabular}
 \end{center}
 }
 \vspace{-0.1in}
 \caption[]{ESP.CMUL.U8.ST.INCP} \end{figure*}


\newpage
\subsubsection{ESP.MAX.S16.A}\label{instruction:ESPMAXS16A}
\textbf{Syntax}\\
ESP.MAX.S16.A qw,rd

\textbf{Description}\\
This instruction saves the maximum signed value of 8 elements of 16 bits from \lstinline{qw} to \lstinline{rd}.

\textbf{Operation}\\
\begin{lstlisting}
    max_value[15 :  0] = max(qw[15 :  0], qw[31 : 16], ..., qw[127:112])
    rd[31 :  0] = {16{max_value[15]}, max_value[15 :  0]}
\end{lstlisting}

\textbf{Schedule}\\
\begin{longtable}[c]{|l|l|l|}
    \hline
    \rowcolor{lightgray}
    \textbf{Operands} & \textbf{Use} & \textbf{Def} \\\hline
    \endfirsthead
    \hline
    \rowcolor{lightgray}
     \textbf{Operands} & \textbf{Use} & \textbf{Def} \\\hline
    \endhead
    \hline
    \endfoot


    rd & - & 2 \\\hline
     qw & 1 & -

\end{longtable}

\textbf{Format}\\

\begin{figure*}[!htbp]
 {\setlength{\tabcolsep}{1pt}
 \begin{center}
\begin{tabular}{@{}S@{}W@{}Y@{}I@{}E@{}I@{}F@{}O}
\instbitrange{31}{26}&
\instbitrange{25}{24}&
\instbitrange{23}{20}&
\instbit{19}&
\instbitrange{18}{11}&
\instbit{10}&
\instbitrange{9}{7}&
\instbitrange{6}{0} \\
\hline
\multicolumn{1}{|c|}{1100XX} &
\multicolumn{1}{c|}{{\scriptsize qw[1:0]}} &
\multicolumn{1}{c|}{00XX} &
\multicolumn{1}{c|}{{\scriptsize qw[2]}} &
\multicolumn{1}{c|}{0XXX01XX} &
\multicolumn{1}{c|}{{\scriptsize rd[4]}} &
\multicolumn{1}{c|}{{\scriptsize rd[2:0]}} &
\multicolumn{1}{c|}{0011011}\\
\hline
6& 2& 4& 1& 8& 1& 3& 7 \\
\end{tabular}
 \end{center}
 }
 \vspace{-0.1in}
 \caption[]{ESP.MAX.S16.A} \end{figure*}


\newpage
\subsubsection{ESP.MAX.S32.A}\label{instruction:ESPMAXS32A}
\textbf{Syntax}\\
ESP.MAX.S32.A qw,rd

\textbf{Description}\\
This instruction saves the maximum signed value of 4 elements of 32 bits from \lstinline{qw} to \lstinline{rd}.


\textbf{Operation}\\
\begin{lstlisting}
    max_value[31 :  0] = max(qw[31 :  0], qw[63 : 32], ..., qw[127:96])
    rd[31 :  0] = max_value[31 :  0]
\end{lstlisting}

\textbf{Schedule}\\
\begin{longtable}[c]{|l|l|l|}
    \hline
    \rowcolor{lightgray}
    \textbf{Operands} & \textbf{Use} & \textbf{Def} \\\hline
    \endfirsthead
    \hline
    \rowcolor{lightgray}
     \textbf{Operands} & \textbf{Use} & \textbf{Def} \\\hline
    \endhead
    \hline
    \endfoot


    rd & - & 2 \\\hline
     qw & 1 & -

\end{longtable}

\textbf{Format}\\

\begin{figure*}[!htbp]
 {\setlength{\tabcolsep}{1pt}
 \begin{center}
\begin{tabular}{@{}S@{}W@{}Y@{}I@{}E@{}I@{}F@{}O}
\instbitrange{31}{26}&
\instbitrange{25}{24}&
\instbitrange{23}{20}&
\instbit{19}&
\instbitrange{18}{11}&
\instbit{10}&
\instbitrange{9}{7}&
\instbitrange{6}{0} \\
\hline
\multicolumn{1}{|c|}{1X10XX} &
\multicolumn{1}{c|}{{\scriptsize qw[1:0]}} &
\multicolumn{1}{c|}{00XX} &
\multicolumn{1}{c|}{{\scriptsize qw[2]}} &
\multicolumn{1}{c|}{0XXX01XX} &
\multicolumn{1}{c|}{{\scriptsize rd[4]}} &
\multicolumn{1}{c|}{{\scriptsize rd[2:0]}} &
\multicolumn{1}{c|}{0011011}\\
\hline
6& 2& 4& 1& 8& 1& 3& 7 \\
\end{tabular}
 \end{center}
 }
 \vspace{-0.1in}
 \caption[]{ESP.MAX.S32.A} \end{figure*}


\newpage
\subsubsection{ESP.MAX.S8.A}\label{instruction:ESPMAXS8A}
\textbf{Syntax}\\
ESP.MAX.S8.A qw,rd

\textbf{Description}\\
This instruction saves the maximum signed value of 16 elements of 8 bits from \lstinline{qw} to \lstinline{rd}.

\textbf{Operation}\\
\begin{lstlisting}
    max_value[7 :  0] = max(qw[7 :  0], qw[16 : 8], ..., qw[127:120])
    rd[31 :  0] = {24{max_value[7]}, max_value[7 :  0]}
\end{lstlisting}

\textbf{Schedule}\\
\begin{longtable}[c]{|l|l|l|}
    \hline
    \rowcolor{lightgray}
    \textbf{Operands} & \textbf{Use} & \textbf{Def} \\\hline
    \endfirsthead
    \hline
    \rowcolor{lightgray}
     \textbf{Operands} & \textbf{Use} & \textbf{Def} \\\hline
    \endhead
    \hline
    \endfoot


    rd & - & 2 \\\hline
     qw & 1 & -

\end{longtable}

\textbf{Format}\\

\begin{figure*}[!htbp]
 {\setlength{\tabcolsep}{1pt}
 \begin{center}
\begin{tabular}{@{}S@{}W@{}Y@{}I@{}E@{}I@{}F@{}O}
\instbitrange{31}{26}&
\instbitrange{25}{24}&
\instbitrange{23}{20}&
\instbit{19}&
\instbitrange{18}{11}&
\instbit{10}&
\instbitrange{9}{7}&
\instbitrange{6}{0} \\
\hline
\multicolumn{1}{|c|}{0100XX} &
\multicolumn{1}{c|}{{\scriptsize qw[1:0]}} &
\multicolumn{1}{c|}{00XX} &
\multicolumn{1}{c|}{{\scriptsize qw[2]}} &
\multicolumn{1}{c|}{0XXX01XX} &
\multicolumn{1}{c|}{{\scriptsize rd[4]}} &
\multicolumn{1}{c|}{{\scriptsize rd[2:0]}} &
\multicolumn{1}{c|}{0011011}\\
\hline
6& 2& 4& 1& 8& 1& 3& 7 \\
\end{tabular}
 \end{center}
 }
 \vspace{-0.1in}
 \caption[]{ESP.MAX.S8.A} \end{figure*}


\newpage
\subsubsection{ESP.MAX.U16.A}\label{instruction:ESPMAXU16A}
\textbf{Syntax}\\
ESP.MAX.U16.A qw,rd

\textbf{Description}\\
This instruction saves the maximum unsigned value of 8 elements of 16 bits from \lstinline{qw} to \lstinline{rd}.

\textbf{Operation}\\
\begin{lstlisting}
    max_value[15 :  0] = max(qw[15 :  0], qw[31 : 16], ..., qw[127:112])
    rd[31 :  0] = {16'b0, max_value[15 :  0]}
\end{lstlisting}


\textbf{Schedule}\\
\begin{longtable}[c]{|l|l|l|}
    \hline
    \rowcolor{lightgray}
    \textbf{Operands} & \textbf{Use} & \textbf{Def} \\\hline
    \endfirsthead
    \hline
    \rowcolor{lightgray}
     \textbf{Operands} & \textbf{Use} & \textbf{Def} \\\hline
    \endhead
    \hline
    \endfoot


    rd & - & 2 \\\hline
     qw & 1 & -

\end{longtable}

\textbf{Format}\\

\begin{figure*}[!htbp]
 {\setlength{\tabcolsep}{1pt}
 \begin{center}
\begin{tabular}{@{}S@{}W@{}Y@{}I@{}E@{}I@{}F@{}O}
\instbitrange{31}{26}&
\instbitrange{25}{24}&
\instbitrange{23}{20}&
\instbit{19}&
\instbitrange{18}{11}&
\instbit{10}&
\instbitrange{9}{7}&
\instbitrange{6}{0} \\
\hline
\multicolumn{1}{|c|}{1000XX} &
\multicolumn{1}{c|}{{\scriptsize qw[1:0]}} &
\multicolumn{1}{c|}{00XX} &
\multicolumn{1}{c|}{{\scriptsize qw[2]}} &
\multicolumn{1}{c|}{0XXX01XX} &
\multicolumn{1}{c|}{{\scriptsize rd[4]}} &
\multicolumn{1}{c|}{{\scriptsize rd[2:0]}} &
\multicolumn{1}{c|}{0011011}\\
\hline
6& 2& 4& 1& 8& 1& 3& 7 \\
\end{tabular}
 \end{center}
 }
 \vspace{-0.1in}
 \caption[]{ESP.MAX.U16.A} \end{figure*}


\newpage
\subsubsection{ESP.MAX.U32.A}\label{instruction:ESPMAXU32A}
\textbf{Syntax}\\
ESP.MAX.U32.A qw,rd

\textbf{Description}\\
This instruction saves the maximum unsigned value of 4 elements of 32 bits from \lstinline{qw} to \lstinline{rd}.


\textbf{Operation}\\
\begin{lstlisting}
    max_value[31 :  0] = max(qw[31 :  0], qw[63 : 32], ..., qw[127:96])
    rd[31 :  0] = max_value[31 :  0]
\end{lstlisting}

\textbf{Schedule}\\
\begin{longtable}[c]{|l|l|l|}
    \hline
    \rowcolor{lightgray}
    \textbf{Operands} & \textbf{Use} & \textbf{Def} \\\hline
    \endfirsthead
    \hline
    \rowcolor{lightgray}
     \textbf{Operands} & \textbf{Use} & \textbf{Def} \\\hline
    \endhead
    \hline
    \endfoot


    rd & - & 2 \\\hline
     qw & 1 & -

\end{longtable}

\textbf{Format}\\

\begin{figure*}[!htbp]
 {\setlength{\tabcolsep}{1pt}
 \begin{center}
\begin{tabular}{@{}S@{}W@{}Y@{}I@{}E@{}I@{}F@{}O}
\instbitrange{31}{26}&
\instbitrange{25}{24}&
\instbitrange{23}{20}&
\instbit{19}&
\instbitrange{18}{11}&
\instbit{10}&
\instbitrange{9}{7}&
\instbitrange{6}{0} \\
\hline
\multicolumn{1}{|c|}{0X10XX} &
\multicolumn{1}{c|}{{\scriptsize qw[1:0]}} &
\multicolumn{1}{c|}{00XX} &
\multicolumn{1}{c|}{{\scriptsize qw[2]}} &
\multicolumn{1}{c|}{0XXX01XX} &
\multicolumn{1}{c|}{{\scriptsize rd[4]}} &
\multicolumn{1}{c|}{{\scriptsize rd[2:0]}} &
\multicolumn{1}{c|}{0011011}\\
\hline
6& 2& 4& 1& 8& 1& 3& 7 \\
\end{tabular}
 \end{center}
 }
 \vspace{-0.1in}
 \caption[]{ESP.MAX.U32.A} \end{figure*}


\newpage
\subsubsection{ESP.MAX.U8.A}\label{instruction:ESPMAXU8A}
\textbf{Syntax}\\
ESP.MAX.U8.A qw,rd

\textbf{Description}\\
This instruction saves the maximum unsigned value of 16 elements of 8 bits from \lstinline{qw} to \lstinline{rd}.

\textbf{Operation}\\
\begin{lstlisting}
    max_value[7 :  0] = max(qw[7 :  0], qw[16 : 8], ..., qw[127:120])
    rd[31 :  0] = {24{0}, max_value[7 :  0]}
\end{lstlisting}


\textbf{Schedule}\\
\begin{longtable}[c]{|l|l|l|}
    \hline
    \rowcolor{lightgray}
    \textbf{Operands} & \textbf{Use} & \textbf{Def} \\\hline
    \endfirsthead
    \hline
    \rowcolor{lightgray}
     \textbf{Operands} & \textbf{Use} & \textbf{Def} \\\hline
    \endhead
    \hline
    \endfoot


    rd & - & 2 \\\hline
     qw & 1 & -

\end{longtable}

\textbf{Format}\\

\begin{figure*}[!htbp]
 {\setlength{\tabcolsep}{1pt}
 \begin{center}
\begin{tabular}{@{}S@{}W@{}Y@{}I@{}E@{}I@{}F@{}O}
\instbitrange{31}{26}&
\instbitrange{25}{24}&
\instbitrange{23}{20}&
\instbit{19}&
\instbitrange{18}{11}&
\instbit{10}&
\instbitrange{9}{7}&
\instbitrange{6}{0} \\
\hline
\multicolumn{1}{|c|}{0000XX} &
\multicolumn{1}{c|}{{\scriptsize qw[1:0]}} &
\multicolumn{1}{c|}{00XX} &
\multicolumn{1}{c|}{{\scriptsize qw[2]}} &
\multicolumn{1}{c|}{0XXX01XX} &
\multicolumn{1}{c|}{{\scriptsize rd[4]}} &
\multicolumn{1}{c|}{{\scriptsize rd[2:0]}} &
\multicolumn{1}{c|}{0011011}\\
\hline
6& 2& 4& 1& 8& 1& 3& 7 \\
\end{tabular}
 \end{center}
 }
 \vspace{-0.1in}
 \caption[]{ESP.MAX.U8.A} \end{figure*}


\newpage
\subsubsection{ESP.MIN.S16.A}\label{instruction:ESPMINS16A}
\textbf{Syntax}\\
ESP.MIN.S16.A qw,rd

\textbf{Description}\\
This instruction saves the minimum signed value of 8 elements of 16 bits from \lstinline{qw} to \lstinline{rd}.

\textbf{Operation}\\
\begin{lstlisting}
    min_value[15 :  0] = min(qw[15 :  0], qw[31 : 16], ..., qw[127:112])
    rd[31 :  0] = {16{min_value[15]}, min_value[15 :  0]}
\end{lstlisting}


\textbf{Schedule}\\
\begin{longtable}[c]{|l|l|l|}
    \hline
    \rowcolor{lightgray}
    \textbf{Operands} & \textbf{Use} & \textbf{Def} \\\hline
    \endfirsthead
    \hline
    \rowcolor{lightgray}
     \textbf{Operands} & \textbf{Use} & \textbf{Def} \\\hline
    \endhead
    \hline
    \endfoot


    rd & - & 2 \\\hline
     qw & 1 & -

\end{longtable}

\textbf{Format}\\

\begin{figure*}[!htbp]
 {\setlength{\tabcolsep}{1pt}
 \begin{center}
\begin{tabular}{@{}S@{}W@{}Y@{}I@{}E@{}I@{}F@{}O}
\instbitrange{31}{26}&
\instbitrange{25}{24}&
\instbitrange{23}{20}&
\instbit{19}&
\instbitrange{18}{11}&
\instbit{10}&
\instbitrange{9}{7}&
\instbitrange{6}{0} \\
\hline
\multicolumn{1}{|c|}{1101XX} &
\multicolumn{1}{c|}{{\scriptsize qw[1:0]}} &
\multicolumn{1}{c|}{00XX} &
\multicolumn{1}{c|}{{\scriptsize qw[2]}} &
\multicolumn{1}{c|}{0XXX01XX} &
\multicolumn{1}{c|}{{\scriptsize rd[4]}} &
\multicolumn{1}{c|}{{\scriptsize rd[2:0]}} &
\multicolumn{1}{c|}{0011011}\\
\hline
6& 2& 4& 1& 8& 1& 3& 7 \\
\end{tabular}
 \end{center}
 }
 \vspace{-0.1in}
 \caption[]{ESP.MIN.S16.A} \end{figure*}


\newpage
\subsubsection{ESP.MIN.S32.A}\label{instruction:ESPMINS32A}
\textbf{Syntax}\\
ESP.MIN.S32.A qw,rd

\textbf{Description}\\
This instruction saves the minimum signed value of 4 elements of 32 bits of \lstinline{qw} to \lstinline{rd}.


\textbf{Operation}\\
\begin{lstlisting}
    min_value[31 :  0] = min(qw[31 :  0], qw[63 : 32], ..., qw[127:96])
    rd[31 :  0] = min_value[31 :  0]
\end{lstlisting}


\textbf{Schedule}\\
\begin{longtable}[c]{|l|l|l|}
    \hline
    \rowcolor{lightgray}
    \textbf{Operands} & \textbf{Use} & \textbf{Def} \\\hline
    \endfirsthead
    \hline
    \rowcolor{lightgray}
     \textbf{Operands} & \textbf{Use} & \textbf{Def} \\\hline
    \endhead
    \hline
    \endfoot


    rd & - & 2 \\\hline
     qw & 1 & -

\end{longtable}

\textbf{Format}\\

\begin{figure*}[!htbp]
 {\setlength{\tabcolsep}{1pt}
 \begin{center}
\begin{tabular}{@{}S@{}W@{}Y@{}I@{}E@{}I@{}F@{}O}
\instbitrange{31}{26}&
\instbitrange{25}{24}&
\instbitrange{23}{20}&
\instbit{19}&
\instbitrange{18}{11}&
\instbit{10}&
\instbitrange{9}{7}&
\instbitrange{6}{0} \\
\hline
\multicolumn{1}{|c|}{1X11XX} &
\multicolumn{1}{c|}{{\scriptsize qw[1:0]}} &
\multicolumn{1}{c|}{00XX} &
\multicolumn{1}{c|}{{\scriptsize qw[2]}} &
\multicolumn{1}{c|}{0XXX01XX} &
\multicolumn{1}{c|}{{\scriptsize rd[4]}} &
\multicolumn{1}{c|}{{\scriptsize rd[2:0]}} &
\multicolumn{1}{c|}{0011011}\\
\hline
6& 2& 4& 1& 8& 1& 3& 7 \\
\end{tabular}
 \end{center}
 }
 \vspace{-0.1in}
 \caption[]{ESP.MIN.S32.A} \end{figure*}


\newpage
\subsubsection{ESP.MIN.S8.A}\label{instruction:ESPMINS8A}
\textbf{Syntax}\\
ESP.MIN.S8.A qw,rd

\textbf{Description}\\
This instruction saves the minimum signed value of 16 elements of 8 bits from \lstinline{qw} to \lstinline{rd}.

\textbf{Operation}\\
\begin{lstlisting}
    min_value[7 :  0] = min(qw[7 :  0], qw[16 : 8], ..., qw[127:120])
    rd[31 :  0] = {24{min_value[7]}, min_value[7 :  0]}
\end{lstlisting}


\textbf{Schedule}\\
\begin{longtable}[c]{|l|l|l|}
    \hline
    \rowcolor{lightgray}
    \textbf{Operands} & \textbf{Use} & \textbf{Def} \\\hline
    \endfirsthead
    \hline
    \rowcolor{lightgray}
     \textbf{Operands} & \textbf{Use} & \textbf{Def} \\\hline
    \endhead
    \hline
    \endfoot


    rd & - & 2 \\\hline
     qw & 1 & -

\end{longtable}

\textbf{Format}\\

\begin{figure*}[!htbp]
 {\setlength{\tabcolsep}{1pt}
 \begin{center}
\begin{tabular}{@{}S@{}W@{}Y@{}I@{}E@{}I@{}F@{}O}
\instbitrange{31}{26}&
\instbitrange{25}{24}&
\instbitrange{23}{20}&
\instbit{19}&
\instbitrange{18}{11}&
\instbit{10}&
\instbitrange{9}{7}&
\instbitrange{6}{0} \\
\hline
\multicolumn{1}{|c|}{0101XX} &
\multicolumn{1}{c|}{{\scriptsize qw[1:0]}} &
\multicolumn{1}{c|}{00XX} &
\multicolumn{1}{c|}{{\scriptsize qw[2]}} &
\multicolumn{1}{c|}{0XXX01XX} &
\multicolumn{1}{c|}{{\scriptsize rd[4]}} &
\multicolumn{1}{c|}{{\scriptsize rd[2:0]}} &
\multicolumn{1}{c|}{0011011}\\
\hline
6& 2& 4& 1& 8& 1& 3& 7 \\
\end{tabular}
 \end{center}
 }
 \vspace{-0.1in}
 \caption[]{ESP.MIN.S8.A} \end{figure*}


\newpage
\subsubsection{ESP.MIN.U16.A}\label{instruction:ESPMINU16A}
\textbf{Syntax}\\
ESP.MIN.U16.A qw,rd

\textbf{Description}\\
This instruction saves the minimum unsigned value of 8 elements of 16 bits from \lstinline{qw} to \lstinline{rd}.

\textbf{Operation}\\
\begin{lstlisting}
    min_value[15 :  0] = min(qw[15 :  0], qw[31 : 16], ..., qw[127:112])
    rd[31 :  0] = {16'b0, min_value[15 :  0]}
\end{lstlisting}


\textbf{Schedule}\\
\begin{longtable}[c]{|l|l|l|}
    \hline
    \rowcolor{lightgray}
    \textbf{Operands} & \textbf{Use} & \textbf{Def} \\\hline
    \endfirsthead
    \hline
    \rowcolor{lightgray}
     \textbf{Operands} & \textbf{Use} & \textbf{Def} \\\hline
    \endhead
    \hline
    \endfoot


    rd & - & 2 \\\hline
     qw & 1 & -

\end{longtable}

\textbf{Format}\\

\begin{figure*}[!htbp]
 {\setlength{\tabcolsep}{1pt}
 \begin{center}
\begin{tabular}{@{}S@{}W@{}Y@{}I@{}E@{}I@{}F@{}O}
\instbitrange{31}{26}&
\instbitrange{25}{24}&
\instbitrange{23}{20}&
\instbit{19}&
\instbitrange{18}{11}&
\instbit{10}&
\instbitrange{9}{7}&
\instbitrange{6}{0} \\
\hline
\multicolumn{1}{|c|}{1001XX} &
\multicolumn{1}{c|}{{\scriptsize qw[1:0]}} &
\multicolumn{1}{c|}{00XX} &
\multicolumn{1}{c|}{{\scriptsize qw[2]}} &
\multicolumn{1}{c|}{0XXX01XX} &
\multicolumn{1}{c|}{{\scriptsize rd[4]}} &
\multicolumn{1}{c|}{{\scriptsize rd[2:0]}} &
\multicolumn{1}{c|}{0011011}\\
\hline
6& 2& 4& 1& 8& 1& 3& 7 \\
\end{tabular}
 \end{center}
 }
 \vspace{-0.1in}
 \caption[]{ESP.MIN.U16.A} \end{figure*}


\newpage
\subsubsection{ESP.MIN.U32.A}\label{instruction:ESPMINU32A}
\textbf{Syntax}\\
ESP.MIN.U32.A qw,rd

\textbf{Description}\\
This instruction saves the minimum unsigned value of 4 elements of 32 bits from \lstinline{qw} to \lstinline{rd}.


\textbf{Operation}\\
\begin{lstlisting}
    min_value[31 :  0] = min(qw[31 :  0], qw[63 : 32], ..., qw[127:96])
    rd[31 :  0] = min_value[31 :  0]
\end{lstlisting}


\textbf{Schedule}\\
\begin{longtable}[c]{|l|l|l|}
    \hline
    \rowcolor{lightgray}
    \textbf{Operands} & \textbf{Use} & \textbf{Def} \\\hline
    \endfirsthead
    \hline
    \rowcolor{lightgray}
     \textbf{Operands} & \textbf{Use} & \textbf{Def} \\\hline
    \endhead
    \hline
    \endfoot


    rd & - & 2 \\\hline
     qw & 1 & -

\end{longtable}

\textbf{Format}\\

\begin{figure*}[!htbp]
 {\setlength{\tabcolsep}{1pt}
 \begin{center}
\begin{tabular}{@{}S@{}W@{}Y@{}I@{}E@{}I@{}F@{}O}
\instbitrange{31}{26}&
\instbitrange{25}{24}&
\instbitrange{23}{20}&
\instbit{19}&
\instbitrange{18}{11}&
\instbit{10}&
\instbitrange{9}{7}&
\instbitrange{6}{0} \\
\hline
\multicolumn{1}{|c|}{0X11XX} &
\multicolumn{1}{c|}{{\scriptsize qw[1:0]}} &
\multicolumn{1}{c|}{00XX} &
\multicolumn{1}{c|}{{\scriptsize qw[2]}} &
\multicolumn{1}{c|}{0XXX01XX} &
\multicolumn{1}{c|}{{\scriptsize rd[4]}} &
\multicolumn{1}{c|}{{\scriptsize rd[2:0]}} &
\multicolumn{1}{c|}{0011011}\\
\hline
6& 2& 4& 1& 8& 1& 3& 7 \\
\end{tabular}
 \end{center}
 }
 \vspace{-0.1in}
 \caption[]{ESP.MIN.U32.A} \end{figure*}


\newpage
\subsubsection{ESP.MIN.U8.A}\label{instruction:ESPMINU8A}
\textbf{Syntax}\\
ESP.MIN.U8.A qw,rd

\textbf{Description}\\
This instruction saves the minimum unsigned value of 16 elements of 8 bits from \lstinline{qw} to \lstinline{rd}.

\textbf{Operation}\\
\begin{lstlisting}
    min_value[7 :  0] = min(qw[7 :  0], qw[16 : 8], ..., qw[127:120])
    rd[31 :  0] = {24{0}, min_value[7 :  0]}
\end{lstlisting}

\textbf{Schedule}\\
\begin{longtable}[c]{|l|l|l|}
    \hline
    \rowcolor{lightgray}
    \textbf{Operands} & \textbf{Use} & \textbf{Def} \\\hline
    \endfirsthead
    \hline
    \rowcolor{lightgray}
     \textbf{Operands} & \textbf{Use} & \textbf{Def} \\\hline
    \endhead
    \hline
    \endfoot


    rd & - & 2 \\\hline
     qw & 1 & -

\end{longtable}

\textbf{Format}\\

\begin{figure*}[!htbp]
 {\setlength{\tabcolsep}{1pt}
 \begin{center}
\begin{tabular}{@{}S@{}W@{}Y@{}I@{}E@{}I@{}F@{}O}
\instbitrange{31}{26}&
\instbitrange{25}{24}&
\instbitrange{23}{20}&
\instbit{19}&
\instbitrange{18}{11}&
\instbit{10}&
\instbitrange{9}{7}&
\instbitrange{6}{0} \\
\hline
\multicolumn{1}{|c|}{0001XX} &
\multicolumn{1}{c|}{{\scriptsize qw[1:0]}} &
\multicolumn{1}{c|}{00XX} &
\multicolumn{1}{c|}{{\scriptsize qw[2]}} &
\multicolumn{1}{c|}{0XXX01XX} &
\multicolumn{1}{c|}{{\scriptsize rd[4]}} &
\multicolumn{1}{c|}{{\scriptsize rd[2:0]}} &
\multicolumn{1}{c|}{0011011}\\
\hline
6& 2& 4& 1& 8& 1& 3& 7 \\
\end{tabular}
 \end{center}
 }
 \vspace{-0.1in}
 \caption[]{ESP.MIN.U8.A} \end{figure*}


\newpage
\subsubsection{ESP.VABS.16}\label{instruction:ESPVABS16}
\textbf{Syntax}\\
ESP.VABS.16 qv,qy

\textbf{Description}\\
This instruction takes the absolute value of the 8 16-bit data segments.

\textbf{Operation}\\
\begin{lstlisting}[language=Python, escapechar=@]
   qv[ 15:  0] = @(qy[  7:  0] < 0) ? -qy[  7:  0] : qy[  7:  0]@
   qv[ 31: 16] = @(qy[ 15: 16] < 0) ? -qy[ 15:  8] : qy[ 15:  8]@
   qv[ 47: 32] = @(qy[ 47: 32] < 0) ? -qy[ 47: 16] : qy[ 47: 16]@
   qv[ 63: 48] = @(qy[ 63: 48] < 0) ? -qy[ 63: 48] : qy[ 63: 48]@
   qv[ 79: 64] = @(qy[ 79: 64] < 0) ? -qy[ 79: 64] : qy[ 79: 64]@
   qv[ 95: 80] = @(qy[ 95: 80] < 0) ? -qy[ 95: 80] : qy[ 95: 80]@
   qv[111: 96] = @(qy[111: 96] < 0) ? -qy[111: 96] : qy[111: 96]@
   qv[127:112] = @(qy[127:112] < 0) ? -qy[127:112] : qy[127:112]@
\end{lstlisting}

\textbf{Schedule}\\
\begin{longtable}[c]{|l|l|l|}
    \hline
    \rowcolor{lightgray}
    \textbf{Operands} & \textbf{Use} & \textbf{Def} \\\hline
    \endfirsthead
    \hline
    \rowcolor{lightgray}
     \textbf{Operands} & \textbf{Use} & \textbf{Def} \\\hline
    \endhead
    \hline
    \endfoot

    qy & 1 & - \\\hline
     qv & - & 1

\end{longtable}

\textbf{Format}\\

\begin{figure*}[!htbp]
 {\setlength{\tabcolsep}{1pt}
 \begin{center}
\begin{tabular}{@{}F@{}F@{}F@{}F@{}K@{}O}
\instbitrange{31}{29}&
\instbitrange{28}{26}&
\instbitrange{25}{23}&
\instbitrange{22}{20}&
\instbitrange{19}{7}&
\instbitrange{6}{0} \\
\hline
\multicolumn{1}{|c|}{0XX} &
\multicolumn{1}{c|}{{\scriptsize qy[2:0]}} &
\multicolumn{1}{c|}{XXX} &
\multicolumn{1}{c|}{{\scriptsize qv[2:0]}} &
\multicolumn{1}{c|}{X10XX100XX1XX} &
\multicolumn{1}{c|}{0011011}\\
\hline
3& 3& 3& 3& 13& 7 \\
\end{tabular}
 \end{center}
 }
 \vspace{-0.1in}
 \caption[]{ESP.VABS.16} \end{figure*}


\newpage
\subsubsection{ESP.VABS.32}\label{instruction:ESPVABS32}
\textbf{Syntax}\\
ESP.VABS.32 qv,qy

\textbf{Description}\\
This instruction takes the absolute value of the 4 32-bit data segments.

\textbf{Operation}\\
\begin{lstlisting}[language=Python, escapechar=@]
   qv[ 31:  0] = @(qy[ 31:  0] < 0) ? -qy[ 31:  0] : qy[ 31:  0]@
   qv[ 63: 32] = @(qy[ 63: 32] < 0) ? -qy[ 63: 32] : qy[ 63: 32]@
   qv[ 95: 64] = @(qy[ 95: 64] < 0) ? -qy[ 95: 64] : qy[ 95: 64]@
   qv[127: 96] = @(qy[127: 96] < 0) ? -qy[127: 96] : qy[127: 96]@
\end{lstlisting}


\textbf{Schedule}\\
\begin{longtable}[c]{|l|l|l|}
    \hline
    \rowcolor{lightgray}
    \textbf{Operands} & \textbf{Use} & \textbf{Def} \\\hline
    \endfirsthead
    \hline
    \rowcolor{lightgray}
     \textbf{Operands} & \textbf{Use} & \textbf{Def} \\\hline
    \endhead
    \hline
    \endfoot

    qy & 1 & - \\\hline
     qv & - & 1

\end{longtable}

\textbf{Format}\\

\begin{figure*}[!htbp]
 {\setlength{\tabcolsep}{1pt}
 \begin{center}
\begin{tabular}{@{}F@{}F@{}F@{}F@{}K@{}O}
\instbitrange{31}{29}&
\instbitrange{28}{26}&
\instbitrange{25}{23}&
\instbitrange{22}{20}&
\instbitrange{19}{7}&
\instbitrange{6}{0} \\
\hline
\multicolumn{1}{|c|}{0XX} &
\multicolumn{1}{c|}{{\scriptsize qy[2:0]}} &
\multicolumn{1}{c|}{XXX} &
\multicolumn{1}{c|}{{\scriptsize qv[2:0]}} &
\multicolumn{1}{c|}{XX1XX100XX1XX} &
\multicolumn{1}{c|}{0011011}\\
\hline
3& 3& 3& 3& 13& 7 \\
\end{tabular}
 \end{center}
 }
 \vspace{-0.1in}
 \caption[]{ESP.VABS.32} \end{figure*}


\newpage
\subsubsection{ESP.VABS.8}\label{instruction:ESPVABS8}
\textbf{Syntax}\\
ESP.VABS.8 qv,qy

\textbf{Description}\\
This instruction takes the absolute value of the 16 8-bit data segments.

\textbf{Operation}\\
\begin{lstlisting}[language=Python, escapechar=@]
   qv[  7:  0] = @(qy[  7:  0] < 0) ? -qy[  7:  0] : qy[  7:  0]@
   qv[ 15:  8] = @(qy[ 15:  8] < 0) ? -qy[ 15:  8] : qy[ 15:  8]@
   qv[ 23: 16] = @(qy[ 23: 16] < 0) ? -qy[ 23: 16] : qy[ 23: 16]@
   qv[ 31: 24] = @(qy[ 31: 24] < 0) ? -qy[ 31: 24] : qy[ 31: 24]@
   qv[ 39: 32] = @(qy[ 39: 32] < 0) ? -qy[ 39: 32] : qy[ 39: 32]@
   qv[ 47: 40] = @(qy[ 47: 40] < 0) ? -qy[ 47: 40] : qy[ 47: 40]@
   qv[ 55: 48] = @(qy[ 55: 48] < 0) ? -qy[ 55: 48] : qy[ 55: 48]@
   qv[ 63: 56] = @(qy[ 63: 56] < 0) ? -qy[ 63: 56] : qy[ 63: 56]@
   qv[ 71: 64] = @(qy[ 71: 64] < 0) ? -qy[ 71: 64] : qy[ 71: 64]@
   qv[ 79: 72] = @(qy[ 79: 72] < 0) ? -qy[ 79: 72] : qy[ 79: 72]@
   qv[ 87: 80] = @(qy[ 87: 80] < 0) ? -qy[ 87: 80] : qy[ 87: 80]@
   qv[ 95: 88] = @(qy[ 95: 88] < 0) ? -qy[ 95: 88] : qy[ 95: 88]@
   qv[103: 96] = @(qy[103: 96] < 0) ? -qy[103: 96] : qy[103: 96]@
   qv[111:104] = @(qy[111:104] < 0) ? -qy[111:104] : qy[111:104]@
   qv[119:112] = @(qy[119:112] < 0) ? -qy[119:112] : qy[119:112]@
   qv[127:120] = @(qy[127:120] < 0) ? -qy[127:120] : qy[127:120]@
\end{lstlisting}


\textbf{Schedule}\\
\begin{longtable}[c]{|l|l|l|}
    \hline
    \rowcolor{lightgray}
    \textbf{Operands} & \textbf{Use} & \textbf{Def} \\\hline
    \endfirsthead
    \hline
    \rowcolor{lightgray}
     \textbf{Operands} & \textbf{Use} & \textbf{Def} \\\hline
    \endhead
    \hline
    \endfoot

    qy & 1 & - \\\hline
     qv & - & 1

\end{longtable}

\textbf{Format}\\

\begin{figure*}[!htbp]
 {\setlength{\tabcolsep}{1pt}
 \begin{center}
\begin{tabular}{@{}F@{}F@{}F@{}F@{}K@{}O}
\instbitrange{31}{29}&
\instbitrange{28}{26}&
\instbitrange{25}{23}&
\instbitrange{22}{20}&
\instbitrange{19}{7}&
\instbitrange{6}{0} \\
\hline
\multicolumn{1}{|c|}{0XX} &
\multicolumn{1}{c|}{{\scriptsize qy[2:0]}} &
\multicolumn{1}{c|}{XXX} &
\multicolumn{1}{c|}{{\scriptsize qv[2:0]}} &
\multicolumn{1}{c|}{X00XX100XX1XX} &
\multicolumn{1}{c|}{0011011}\\
\hline
3& 3& 3& 3& 13& 7 \\
\end{tabular}
 \end{center}
 }
 \vspace{-0.1in}
 \caption[]{ESP.VABS.8} \end{figure*}


\newpage
\subsubsection{ESP.VADD.S16}\label{instruction:ESPVADDS16}
\textbf{Syntax}\\
ESP.VADD.S16 qv,qx,qy,sat

\textbf{Description}\\
This instruction performs a vector signed addition on 16-bit data in the two registers \lstinline{qx} and \lstinline{qy}. Then, the 8 results obtained from the calculation are saturated, and the saturated results are written to the register \lstinline{qv}.

\textbf{Operation}\\
\begin{lstlisting}
  qv[ 15:  0] = min(max(qx[ 15:  0] + qy[ 15:  0], -2^{15}), 2^{15}-1)
  qv[ 31: 16] = min(max(qx[ 31: 16] + qy[ 31: 16], -2^{15}), 2^{15}-1)
  ...
  qv[127:112] = min(max(qx[127:112] + qy[127:112], -2^{15}), 2^{15}-1)
\end{lstlisting}

\textbf{Schedule}\\
\begin{longtable}[c]{|l|l|l|}
    \hline
    \rowcolor{lightgray}
    \textbf{Operands} & \textbf{Use} & \textbf{Def} \\\hline
    \endfirsthead
    \hline
    \rowcolor{lightgray}
     \textbf{Operands} & \textbf{Use} & \textbf{Def} \\\hline
    \endhead
    \hline
    \endfoot

    qx & 1 & - \\\hline
  qy & 1 & - \\\hline
     qv & - & 1

\end{longtable}

\textbf{Format}\\

\begin{figure*}[!htbp]
 {\setlength{\tabcolsep}{1pt}
 \begin{center}
\begin{tabular}{@{}F@{}F@{}F@{}F@{}F@{}I@{}T@{}O}
\instbitrange{31}{29}&
\instbitrange{28}{26}&
\instbitrange{25}{23}&
\instbitrange{22}{20}&
\instbitrange{19}{17}&
\instbit{16}&
\instbitrange{15}{7}&
\instbitrange{6}{0} \\
\hline
\multicolumn{1}{|c|}{{\scriptsize qx[2:0]}} &
\multicolumn{1}{c|}{{\scriptsize qy[2:0]}} &
\multicolumn{1}{c|}{1X1} &
\multicolumn{1}{c|}{{\scriptsize qv[2:0]}} &
\multicolumn{1}{c|}{X11} &
\multicolumn{1}{c|}{{\scriptsize sat[0]}} &
\multicolumn{1}{c|}{001XX011X} &
\multicolumn{1}{c|}{0011011}\\
\hline
3& 3& 3& 3& 3& 1& 9& 7 \\
\end{tabular}
 \end{center}
 }
 \vspace{-0.1in}
 \caption[]{ESP.VADD.S16} \end{figure*}


\newpage
\subsubsection{ESP.VADD.S16.LD.INCP}\label{instruction:ESPVADDS16LDINCP}
\textbf{Syntax}\\
ESP.VADD.S16.LD.INCP qu,rs1,qv,qx,qy,sat

\textbf{Description}\\
This instruction performs a vector signed addition on 16-bit data in the two registers \lstinline{qx} and \lstinline{qy}. Then, the 8 results obtained from the calculation are saturated, and the saturated results are written to the register \lstinline{qv}. \\ During the operation, 16 bytes are loaded from the memory to the register \lstinline{qu}. After the access, the value in the register \lstinline{rs1} is incremented by 16.

\textbf{Operation}\\
\begin{lstlisting}
  qv[ 15:  0] = min(max(qx[ 15:  0] + qy[ 15:  0], -2^{15}), 2^{15}-1)
  qv[ 31: 16] = min(max(qx[ 31: 16] + qy[ 31: 16], -2^{15}), 2^{15}-1)
  ...

  qu[127:0] = load128(rs1[31:0])
  rs1[31:0] = rs1[31:0] + 16
\end{lstlisting}

\textbf{Schedule}\\
\begin{longtable}[c]{|l|l|l|}
    \hline
    \rowcolor{lightgray}
    \textbf{Operands} & \textbf{Use} & \textbf{Def} \\\hline
    \endfirsthead
    \hline
    \rowcolor{lightgray}
     \textbf{Operands} & \textbf{Use} & \textbf{Def} \\\hline
    \endhead
    \hline
    \endfoot

    qx & 1 & - \\\hline
  qy & 1 & - \\\hline
     qv & - & 1 \\\hline
     qu & - & 2 \\\hline
     rs1 & 1 & 1

\end{longtable}

\textbf{Format}\\

\begin{figure*}[!htbp]
 {\setlength{\tabcolsep}{1pt}
 \begin{center}
\begin{tabular}{@{}F@{}F@{}F@{}F@{}I@{}I@{}F@{}W@{}F@{}F@{}O}
\instbitrange{31}{29}&
\instbitrange{28}{26}&
\instbitrange{25}{23}&
\instbitrange{22}{20}&
\instbit{19}&
\instbit{18}&
\instbitrange{17}{15}&
\instbitrange{14}{13}&
\instbitrange{12}{10}&
\instbitrange{9}{7}&
\instbitrange{6}{0} \\
\hline
\multicolumn{1}{|c|}{{\scriptsize qx[2:0]}} &
\multicolumn{1}{c|}{{\scriptsize qy[2:0]}} &
\multicolumn{1}{c|}{011} &
\multicolumn{1}{c|}{{\scriptsize qv[2:0]}} &
\multicolumn{1}{c|}{{\scriptsize sat[0]}} &
\multicolumn{1}{c|}{{\scriptsize rs1[4]}} &
\multicolumn{1}{c|}{{\scriptsize rs1[2:0]}} &
\multicolumn{1}{c|}{10} &
\multicolumn{1}{c|}{{\scriptsize qu[2:0]}} &
\multicolumn{1}{c|}{X10} &
\multicolumn{1}{c|}{0011111}\\
\hline
3& 3& 3& 3& 1& 1& 3& 2& 3& 3& 7 \\
\end{tabular}
 \end{center}
 }
 \vspace{-0.1in}
 \caption[]{ESP.VADD.S16.LD.INCP} \end{figure*}


\newpage
\subsubsection{ESP.VADD.S16.ST.INCP}\label{instruction:ESPVADDS16STINCP}
\textbf{Syntax}\\
ESP.VADD.S16.ST.INCP qu,rs1,qv,qx,qy,sat

\textbf{Description}\\
This instruction performs a vector signed addition on 16-bit data in the two registers \lstinline{qx} and \lstinline{qy}. Then, the 8 results obtained from the calculation are saturated, and the saturated results are written to the register \lstinline{qv}. \\ During the operation, the value in the register \lstinline{qu} is stored in memory. After the access, the value in the register \lstinline{rs1} is incremented by 16.

\textbf{Operation}\\
\begin{lstlisting}
  qv[ 15:  0] = min(max(qx[ 15:  0] + qy[ 15:  0], -2^{15}), 2^{15}-1)
  qv[ 31: 16] = min(max(qx[ 31: 16] + qy[ 31: 16], -2^{15}), 2^{15}-1)
  ...

  qu[127:0] => store128(rs1[31:0])
  rs1[31:0] = rs1[31:0] + 16
\end{lstlisting}

\textbf{Schedule}\\
\begin{longtable}[c]{|l|l|l|}
    \hline
    \rowcolor{lightgray}
    \textbf{Operands} & \textbf{Use} & \textbf{Def} \\\hline
    \endfirsthead
    \hline
    \rowcolor{lightgray}
     \textbf{Operands} & \textbf{Use} & \textbf{Def} \\\hline
    \endhead
    \hline
    \endfoot

    qx & 1 & - \\\hline
  qy & 1 & - \\\hline
     qv & - & 1 \\\hline
     qu & 2 & - \\\hline
     rs1 & 1 & 1

\end{longtable}

\textbf{Format}\\

\begin{figure*}[!htbp]
 {\setlength{\tabcolsep}{1pt}
 \begin{center}
\begin{tabular}{@{}F@{}F@{}F@{}F@{}I@{}I@{}F@{}W@{}F@{}F@{}O}
\instbitrange{31}{29}&
\instbitrange{28}{26}&
\instbitrange{25}{23}&
\instbitrange{22}{20}&
\instbit{19}&
\instbit{18}&
\instbitrange{17}{15}&
\instbitrange{14}{13}&
\instbitrange{12}{10}&
\instbitrange{9}{7}&
\instbitrange{6}{0} \\
\hline
\multicolumn{1}{|c|}{{\scriptsize qx[2:0]}} &
\multicolumn{1}{c|}{{\scriptsize qy[2:0]}} &
\multicolumn{1}{c|}{111} &
\multicolumn{1}{c|}{{\scriptsize qv[2:0]}} &
\multicolumn{1}{c|}{{\scriptsize sat[0]}} &
\multicolumn{1}{c|}{{\scriptsize rs1[4]}} &
\multicolumn{1}{c|}{{\scriptsize rs1[2:0]}} &
\multicolumn{1}{c|}{10} &
\multicolumn{1}{c|}{{\scriptsize qu[2:0]}} &
\multicolumn{1}{c|}{X10} &
\multicolumn{1}{c|}{0011111}\\
\hline
3& 3& 3& 3& 1& 1& 3& 2& 3& 3& 7 \\
\end{tabular}
 \end{center}
 }
 \vspace{-0.1in}
 \caption[]{ESP.VADD.S16.ST.INCP} \end{figure*}


\newpage
\subsubsection{ESP.VADD.S32}\label{instruction:ESPVADDS32}
\textbf{Syntax}\\
ESP.VADD.S32 qv,qx,qy,sat

\textbf{Description}\\
This instruction performs a vector signed addition on 32-bit data in the two registers \lstinline{qx} and \lstinline{qy}. Then, the 4 results obtained from the calculation are saturated, and the saturated results are written to the register \lstinline{qv}.

\textbf{Operation}\\
\begin{lstlisting}
  qv[ 31:  0] = min(max(qx[ 31:  0] + qy[ 31:  0], -2^{31}), 2^{31}-1)
  qv[ 63: 32] = min(max(qx[ 63: 32] + qy[ 63: 32], -2^{31}), 2^{31}-1)
  ...
\end{lstlisting}


\textbf{Schedule}\\
\begin{longtable}[c]{|l|l|l|}
    \hline
    \rowcolor{lightgray}
    \textbf{Operands} & \textbf{Use} & \textbf{Def} \\\hline
    \endfirsthead
    \hline
    \rowcolor{lightgray}
     \textbf{Operands} & \textbf{Use} & \textbf{Def} \\\hline
    \endhead
    \hline
    \endfoot

    qx & 1 & - \\\hline
  qy & 1 & - \\\hline
     qv & - & 1

\end{longtable}

\textbf{Format}\\

\begin{figure*}[!htbp]
 {\setlength{\tabcolsep}{1pt}
 \begin{center}
\begin{tabular}{@{}F@{}F@{}F@{}F@{}W@{}I@{}T@{}O}
\instbitrange{31}{29}&
\instbitrange{28}{26}&
\instbitrange{25}{23}&
\instbitrange{22}{20}&
\instbitrange{19}{18}&
\instbit{17}&
\instbitrange{16}{7}&
\instbitrange{6}{0} \\
\hline
\multicolumn{1}{|c|}{{\scriptsize qx[2:0]}} &
\multicolumn{1}{c|}{{\scriptsize qy[2:0]}} &
\multicolumn{1}{c|}{1X1} &
\multicolumn{1}{c|}{{\scriptsize qv[2:0]}} &
\multicolumn{1}{c|}{X1} &
\multicolumn{1}{c|}{{\scriptsize sat[0]}} &
\multicolumn{1}{c|}{X101XX011X} &
\multicolumn{1}{c|}{0011011}\\
\hline
3& 3& 3& 3& 2& 1& 10& 7 \\
\end{tabular}
 \end{center}
 }
 \vspace{-0.1in}
 \caption[]{ESP.VADD.S32} \end{figure*}


\newpage
\subsubsection{ESP.VADD.S32.LD.INCP}\label{instruction:ESPVADDS32LDINCP}
\textbf{Syntax}\\
ESP.VADD.S32.LD.INCP qu,rs1,qv,qx,qy,sat

\textbf{Description}\\
This instruction performs a vector signed addition on 32-bit data in the two registers \lstinline{qx} and \lstinline{qy}. Then, the 4 results obtained from the calculation are saturated, and the saturated results are written to the register \lstinline{qv}. \\ During the operation, 16 bytes are loaded from the memory to the register \lstinline{qu}. After the access, the value in the register \lstinline{rs1} is incremented by 16.

\textbf{Operation}\\
\begin{lstlisting}
  qv[ 31:  0] = min(max(qx[ 31:  0] + qy[ 31:  0], -2^{31}), 2^{31}-1)
  qv[ 63: 32] = min(max(qx[ 63: 32] + qy[ 63: 32], -2^{31}), 2^{31}-1)
  ...

  qu[127:0] = load128(rs1[31:0])
  rs1[31:0] = rs1[31:0] + 16
\end{lstlisting}


\textbf{Schedule}\\
\begin{longtable}[c]{|l|l|l|}
    \hline
    \rowcolor{lightgray}
    \textbf{Operands} & \textbf{Use} & \textbf{Def} \\\hline
    \endfirsthead
    \hline
    \rowcolor{lightgray}
     \textbf{Operands} & \textbf{Use} & \textbf{Def} \\\hline
    \endhead
    \hline
    \endfoot

    qx & 1 & - \\\hline
  qy & 1 & - \\\hline
     qv & - & 1 \\\hline
     qu & - & 2 \\\hline
     rs1 & 1 & 1

\end{longtable}

\textbf{Format}\\

\begin{figure*}[!htbp]
 {\setlength{\tabcolsep}{1pt}
 \begin{center}
\begin{tabular}{@{}F@{}F@{}F@{}F@{}I@{}I@{}F@{}W@{}F@{}F@{}O}
\instbitrange{31}{29}&
\instbitrange{28}{26}&
\instbitrange{25}{23}&
\instbitrange{22}{20}&
\instbit{19}&
\instbit{18}&
\instbitrange{17}{15}&
\instbitrange{14}{13}&
\instbitrange{12}{10}&
\instbitrange{9}{7}&
\instbitrange{6}{0} \\
\hline
\multicolumn{1}{|c|}{{\scriptsize qx[2:0]}} &
\multicolumn{1}{c|}{{\scriptsize qy[2:0]}} &
\multicolumn{1}{c|}{01X} &
\multicolumn{1}{c|}{{\scriptsize qv[2:0]}} &
\multicolumn{1}{c|}{{\scriptsize sat[0]}} &
\multicolumn{1}{c|}{{\scriptsize rs1[4]}} &
\multicolumn{1}{c|}{{\scriptsize rs1[2:0]}} &
\multicolumn{1}{c|}{10} &
\multicolumn{1}{c|}{{\scriptsize qu[2:0]}} &
\multicolumn{1}{c|}{001} &
\multicolumn{1}{c|}{0011111}\\
\hline
3& 3& 3& 3& 1& 1& 3& 2& 3& 3& 7 \\
\end{tabular}
 \end{center}
 }
 \vspace{-0.1in}
 \caption[]{ESP.VADD.S32.LD.INCP} \end{figure*}


\newpage
\subsubsection{ESP.VADD.S32.ST.INCP}\label{instruction:ESPVADDS32STINCP}
\textbf{Syntax}\\
ESP.VADD.S32.ST.INCP qu,rs1,qv,qx,qy,sat

\textbf{Description}\\
This instruction performs a vector signed addition on 32-bit data in the two registers \lstinline{qx} and \lstinline{qy}. Then, the 4 results obtained from the calculation are saturated, and the saturated results are written to the register \lstinline{qu}. \\ During the operation, the value in the register \lstinline{qu} is stored in memory. After the access, the value in the register \lstinline{rs1} is incremented by 16.

\textbf{Operation}\\
\begin{lstlisting}
  qv[ 31:  0] = min(max(qx[ 31:  0] + qy[ 31:  0], -2^{31}), 2^{31}-1)
  qv[ 63: 32] = min(max(qx[ 63: 32] + qy[ 63: 32], -2^{31}), 2^{31}-1)
  ...

  qu[127:0] => store128(rs1[31:0])
  rs1[31:0] = rs1[31:0] + 16
\end{lstlisting}


\textbf{Schedule}\\
\begin{longtable}[c]{|l|l|l|}
    \hline
    \rowcolor{lightgray}
    \textbf{Operands} & \textbf{Use} & \textbf{Def} \\\hline
    \endfirsthead
    \hline
    \rowcolor{lightgray}
     \textbf{Operands} & \textbf{Use} & \textbf{Def} \\\hline
    \endhead
    \hline
    \endfoot

    qx & 1 & - \\\hline
  qy & 1 & - \\\hline
     qv & - & 1 \\\hline
     qu & 2 & - \\\hline
     rs1 & 1 & 1

\end{longtable}

\textbf{Format}\\

\begin{figure*}[!htbp]
 {\setlength{\tabcolsep}{1pt}
 \begin{center}
\begin{tabular}{@{}F@{}F@{}F@{}F@{}I@{}I@{}F@{}W@{}F@{}F@{}O}
\instbitrange{31}{29}&
\instbitrange{28}{26}&
\instbitrange{25}{23}&
\instbitrange{22}{20}&
\instbit{19}&
\instbit{18}&
\instbitrange{17}{15}&
\instbitrange{14}{13}&
\instbitrange{12}{10}&
\instbitrange{9}{7}&
\instbitrange{6}{0} \\
\hline
\multicolumn{1}{|c|}{{\scriptsize qx[2:0]}} &
\multicolumn{1}{c|}{{\scriptsize qy[2:0]}} &
\multicolumn{1}{c|}{11X} &
\multicolumn{1}{c|}{{\scriptsize qv[2:0]}} &
\multicolumn{1}{c|}{{\scriptsize sat[0]}} &
\multicolumn{1}{c|}{{\scriptsize rs1[4]}} &
\multicolumn{1}{c|}{{\scriptsize rs1[2:0]}} &
\multicolumn{1}{c|}{10} &
\multicolumn{1}{c|}{{\scriptsize qu[2:0]}} &
\multicolumn{1}{c|}{001} &
\multicolumn{1}{c|}{0011111}\\
\hline
3& 3& 3& 3& 1& 1& 3& 2& 3& 3& 7 \\
\end{tabular}
 \end{center}
 }
 \vspace{-0.1in}
 \caption[]{ESP.VADD.S32.ST.INCP} \end{figure*}


\newpage
\subsubsection{ESP.VADD.S8}\label{instruction:ESPVADDS8}
\textbf{Syntax}\\
ESP.VADD.S8 qv,qx,qy,sat

\textbf{Description}\\
This instruction performs a vector signed addition on 8-bit data in the two registers \lstinline{qx} and \lstinline{qy}. Then, the 16 results obtained from the calculation are saturated, and the saturated results are written to the register \lstinline{qv}.

\textbf{Operation}\\
\begin{lstlisting}
  qv[  7:  0] = min(max(qx[  7:  0] + qy[  7:  0], -2^{7}), 2^{7}-1)
  qv[ 15:  8] = min(max(qx[ 15:  8] + qy[ 15:  8], -2^{7}), 2^{7}-1)
  ...
  qv[127:120] = min(max(qx[127:120] + qy[127:120], -2^{7}), 2^{7}-1)
\end{lstlisting}


\textbf{Schedule}\\
\begin{longtable}[c]{|l|l|l|}
    \hline
    \rowcolor{lightgray}
    \textbf{Operands} & \textbf{Use} & \textbf{Def} \\\hline
    \endfirsthead
    \hline
    \rowcolor{lightgray}
     \textbf{Operands} & \textbf{Use} & \textbf{Def} \\\hline
    \endhead
    \hline
    \endfoot

    qx & 1 & - \\\hline
  qy & 1 & - \\\hline
     qv & - & 1

\end{longtable}

\textbf{Format}\\

\begin{figure*}[!htbp]
 {\setlength{\tabcolsep}{1pt}
 \begin{center}
\begin{tabular}{@{}F@{}F@{}F@{}F@{}F@{}I@{}T@{}O}
\instbitrange{31}{29}&
\instbitrange{28}{26}&
\instbitrange{25}{23}&
\instbitrange{22}{20}&
\instbitrange{19}{17}&
\instbit{16}&
\instbitrange{15}{7}&
\instbitrange{6}{0} \\
\hline
\multicolumn{1}{|c|}{{\scriptsize qx[2:0]}} &
\multicolumn{1}{c|}{{\scriptsize qy[2:0]}} &
\multicolumn{1}{c|}{1X1} &
\multicolumn{1}{c|}{{\scriptsize qv[2:0]}} &
\multicolumn{1}{c|}{X01} &
\multicolumn{1}{c|}{{\scriptsize sat[0]}} &
\multicolumn{1}{c|}{001XX011X} &
\multicolumn{1}{c|}{0011011}\\
\hline
3& 3& 3& 3& 3& 1& 9& 7 \\
\end{tabular}
 \end{center}
 }
 \vspace{-0.1in}
 \caption[]{ESP.VADD.S8} \end{figure*}


\newpage
\subsubsection{ESP.VADD.S8.LD.INCP}\label{instruction:ESPVADDS8LDINCP}
\textbf{Syntax}\\
ESP.VADD.S8.LD.INCP qu,rs1,qv,qx,qy,sat

\textbf{Description}\\
This instruction performs a vector signed addition on 8-bit data in the two registers \lstinline{qx} and \lstinline{qy}. Then, the 16 results obtained from the calculation are saturated, and the saturated results are written to the register \lstinline{qa}. \\ During the operation, 16 bytes are loaded from the memory to the register \lstinline{qu}. After the access, the value in the register \lstinline{rs1} is incremented by 16.

\textbf{Operation}\\
\begin{lstlisting}
  qv[  7:  0] = min(max(qx[  7:  0] + qy[  7:  0], -2^{7}), 2^{7}-1)
  qv[ 15:  8] = min(max(qx[ 15:  8] + qy[ 15:  8], -2^{7}), 2^{7}-1)
  ...

  qu[127:0] = load128(rs1[31:0])
  rs1[31:0] = rs1[31:0] + 16
\end{lstlisting}


\textbf{Schedule}\\
\begin{longtable}[c]{|l|l|l|}
    \hline
    \rowcolor{lightgray}
    \textbf{Operands} & \textbf{Use} & \textbf{Def} \\\hline
    \endfirsthead
    \hline
    \rowcolor{lightgray}
     \textbf{Operands} & \textbf{Use} & \textbf{Def} \\\hline
    \endhead
    \hline
    \endfoot

    qx & 1 & - \\\hline
  qy & 1 & - \\\hline
     qv & - & 1 \\\hline
     qu & - & 2 \\\hline
     rs1 & 1 & 1

\end{longtable}

\textbf{Format}\\

\begin{figure*}[!htbp]
 {\setlength{\tabcolsep}{1pt}
 \begin{center}
\begin{tabular}{@{}F@{}F@{}F@{}F@{}I@{}I@{}F@{}W@{}F@{}F@{}O}
\instbitrange{31}{29}&
\instbitrange{28}{26}&
\instbitrange{25}{23}&
\instbitrange{22}{20}&
\instbit{19}&
\instbit{18}&
\instbitrange{17}{15}&
\instbitrange{14}{13}&
\instbitrange{12}{10}&
\instbitrange{9}{7}&
\instbitrange{6}{0} \\
\hline
\multicolumn{1}{|c|}{{\scriptsize qx[2:0]}} &
\multicolumn{1}{c|}{{\scriptsize qy[2:0]}} &
\multicolumn{1}{c|}{001} &
\multicolumn{1}{c|}{{\scriptsize qv[2:0]}} &
\multicolumn{1}{c|}{{\scriptsize sat[0]}} &
\multicolumn{1}{c|}{{\scriptsize rs1[4]}} &
\multicolumn{1}{c|}{{\scriptsize rs1[2:0]}} &
\multicolumn{1}{c|}{10} &
\multicolumn{1}{c|}{{\scriptsize qu[2:0]}} &
\multicolumn{1}{c|}{X10} &
\multicolumn{1}{c|}{0011111}\\
\hline
3& 3& 3& 3& 1& 1& 3& 2& 3& 3& 7 \\
\end{tabular}
 \end{center}
 }
 \vspace{-0.1in}
 \caption[]{ESP.VADD.S8.LD.INCP} \end{figure*}


\newpage
\subsubsection{ESP.VADD.S8.ST.INCP}\label{instruction:ESPVADDS8STINCP}
\textbf{Syntax}\\
ESP.VADD.S8.ST.INCP qu,rs1,qv,qx,qy,sat

\textbf{Description}\\
This instruction performs a vector signed addition on 8-bit data in the two registers \lstinline{qx} and \lstinline{qy}. Then, the 16 results obtained from the calculation are saturated, and the saturated results are written to the register \lstinline{qa}. \\ During the operation, the value in the register \lstinline{qu} is stored in memory. After the access, the value in the register \lstinline{rs1} is incremented by 16.

\textbf{Operation}\\
\begin{lstlisting}
  qv[  7:  0] = min(max(qx[  7:  0] + qy[  7:  0], -2^{7}), 2^{7}-1)
  qv[ 15:  8] = min(max(qx[ 15:  8] + qy[ 15:  8], -2^{7}), 2^{7}-1)
  ...

  qu[127:0] => store128(rs1[31:0])
  rs1 [31:0] = rs1[31:0] + 16
\end{lstlisting}

\textbf{Schedule}\\
\begin{longtable}[c]{|l|l|l|}
    \hline
    \rowcolor{lightgray}
    \textbf{Operands} & \textbf{Use} & \textbf{Def} \\\hline
    \endfirsthead
    \hline
    \rowcolor{lightgray}
     \textbf{Operands} & \textbf{Use} & \textbf{Def} \\\hline
    \endhead
    \hline
    \endfoot

    qx & 1 & - \\\hline
  qy & 1 & - \\\hline
     qv & - & 1 \\\hline
     qu & 2 & - \\\hline
     rs1 & 1 & 1

\end{longtable}

\textbf{Format}\\

\begin{figure*}[!htbp]
 {\setlength{\tabcolsep}{1pt}
 \begin{center}
\begin{tabular}{@{}F@{}F@{}F@{}F@{}I@{}I@{}F@{}W@{}F@{}F@{}O}
\instbitrange{31}{29}&
\instbitrange{28}{26}&
\instbitrange{25}{23}&
\instbitrange{22}{20}&
\instbit{19}&
\instbit{18}&
\instbitrange{17}{15}&
\instbitrange{14}{13}&
\instbitrange{12}{10}&
\instbitrange{9}{7}&
\instbitrange{6}{0} \\
\hline
\multicolumn{1}{|c|}{{\scriptsize qx[2:0]}} &
\multicolumn{1}{c|}{{\scriptsize qy[2:0]}} &
\multicolumn{1}{c|}{101} &
\multicolumn{1}{c|}{{\scriptsize qv[2:0]}} &
\multicolumn{1}{c|}{{\scriptsize sat[0]}} &
\multicolumn{1}{c|}{{\scriptsize rs1[4]}} &
\multicolumn{1}{c|}{{\scriptsize rs1[2:0]}} &
\multicolumn{1}{c|}{10} &
\multicolumn{1}{c|}{{\scriptsize qu[2:0]}} &
\multicolumn{1}{c|}{X10} &
\multicolumn{1}{c|}{0011111}\\
\hline
3& 3& 3& 3& 1& 1& 3& 2& 3& 3& 7 \\
\end{tabular}
 \end{center}
 }
 \vspace{-0.1in}
 \caption[]{ESP.VADD.S8.ST.INCP} \end{figure*}


\newpage
\subsubsection{ESP.VADD.U16}\label{instruction:ESPVADDU16}
\textbf{Syntax}\\
ESP.VADD.U16 qv,qx,qy,sat

\textbf{Description}\\
This instruction performs a vector unsigned addition on 16-bit data in the two registers \lstinline{qx} and \lstinline{qy}. Then, the 8 results obtained from the calculation are saturated, and the saturated results are written to the register \lstinline{qv}.

\textbf{Operation}\\
\begin{lstlisting}
  qv[ 15:  0] = min(qx[ 15:  0] + qy[ 15:  0], 2^{16}-1)
  qv[ 31: 16] = min(qx[ 31: 16] + qy[ 31: 16], 2^{16}-1)
  ...
  qv[127:112] = min(qx[127:112] + qy[127:112], 2^{16}-1)
\end{lstlisting}


\textbf{Schedule}\\
\begin{longtable}[c]{|l|l|l|}
    \hline
    \rowcolor{lightgray}
    \textbf{Operands} & \textbf{Use} & \textbf{Def} \\\hline
    \endfirsthead
    \hline
    \rowcolor{lightgray}
     \textbf{Operands} & \textbf{Use} & \textbf{Def} \\\hline
    \endhead
    \hline
    \endfoot

    qx & 1 & - \\\hline
  qy & 1 & - \\\hline
     qv & - & 1

\end{longtable}

\textbf{Format}\\

\begin{figure*}[!htbp]
 {\setlength{\tabcolsep}{1pt}
 \begin{center}
\begin{tabular}{@{}F@{}F@{}F@{}F@{}F@{}I@{}T@{}O}
\instbitrange{31}{29}&
\instbitrange{28}{26}&
\instbitrange{25}{23}&
\instbitrange{22}{20}&
\instbitrange{19}{17}&
\instbit{16}&
\instbitrange{15}{7}&
\instbitrange{6}{0} \\
\hline
\multicolumn{1}{|c|}{{\scriptsize qx[2:0]}} &
\multicolumn{1}{c|}{{\scriptsize qy[2:0]}} &
\multicolumn{1}{c|}{1X1} &
\multicolumn{1}{c|}{{\scriptsize qv[2:0]}} &
\multicolumn{1}{c|}{X10} &
\multicolumn{1}{c|}{{\scriptsize sat[0]}} &
\multicolumn{1}{c|}{001XX011X} &
\multicolumn{1}{c|}{0011011}\\
\hline
3& 3& 3& 3& 3& 1& 9& 7 \\
\end{tabular}
 \end{center}
 }
 \vspace{-0.1in}
 \caption[]{ESP.VADD.U16} \end{figure*}


\newpage
\subsubsection{ESP.VADD.U16.LD.INCP}\label{instruction:ESPVADDU16LDINCP}
\textbf{Syntax}\\
ESP.VADD.U16.LD.INCP qu,rs1,qv,qx,qy,sat

\textbf{Description}\\
This instruction performs a vector unsigned addition on 16-bit data in the two registers \lstinline{qx} and \lstinline{qy}. Then, the 8 results obtained from the calculation are saturated, and the saturated results are written to the register \lstinline{qv}.\\ During the operation, the value in the register \lstinline{qu} is stored in memory. After the access, the value in the register \lstinline{rs1} is incremented by 16.

\textbf{Operation}\\
\begin{lstlisting}
  qv[ 15:  0] = min(qx[ 15:  0] + qy[ 15:  0], 2^{16}-1)
  qv[ 31: 16] = min(qx[ 31: 16] + qy[ 31: 16], 2^{16}-1)
  ...
  qv[127:112] = min(qx[127:112] + qy[127:112], 2^{16}-1)

   qu[127:0] = load128(rs1[31:0])
  rs1[31:0] = rs1[31:0] + 16
\end{lstlisting}


\textbf{Schedule}\\
\begin{longtable}[c]{|l|l|l|}
    \hline
    \rowcolor{lightgray}
    \textbf{Operands} & \textbf{Use} & \textbf{Def} \\\hline
    \endfirsthead
    \hline
    \rowcolor{lightgray}
     \textbf{Operands} & \textbf{Use} & \textbf{Def} \\\hline
    \endhead
    \hline
    \endfoot

    qx & 1 & - \\\hline
  qy & 1 & - \\\hline
     qv & - & 1 \\\hline
     qu & - & 2 \\\hline
     rs1 & 1 & 1

\end{longtable}

\textbf{Format}\\

\begin{figure*}[!htbp]
 {\setlength{\tabcolsep}{1pt}
 \begin{center}
\begin{tabular}{@{}F@{}F@{}F@{}F@{}I@{}I@{}F@{}W@{}F@{}F@{}O}
\instbitrange{31}{29}&
\instbitrange{28}{26}&
\instbitrange{25}{23}&
\instbitrange{22}{20}&
\instbit{19}&
\instbit{18}&
\instbitrange{17}{15}&
\instbitrange{14}{13}&
\instbitrange{12}{10}&
\instbitrange{9}{7}&
\instbitrange{6}{0} \\
\hline
\multicolumn{1}{|c|}{{\scriptsize qx[2:0]}} &
\multicolumn{1}{c|}{{\scriptsize qy[2:0]}} &
\multicolumn{1}{c|}{010} &
\multicolumn{1}{c|}{{\scriptsize qv[2:0]}} &
\multicolumn{1}{c|}{{\scriptsize sat[0]}} &
\multicolumn{1}{c|}{{\scriptsize rs1[4]}} &
\multicolumn{1}{c|}{{\scriptsize rs1[2:0]}} &
\multicolumn{1}{c|}{10} &
\multicolumn{1}{c|}{{\scriptsize qu[2:0]}} &
\multicolumn{1}{c|}{X10} &
\multicolumn{1}{c|}{0011111}\\
\hline
3& 3& 3& 3& 1& 1& 3& 2& 3& 3& 7 \\
\end{tabular}
 \end{center}
 }
 \vspace{-0.1in}
 \caption[]{ESP.VADD.U16.LD.INCP} \end{figure*}


\newpage
\subsubsection{ESP.VADD.U16.ST.INCP}\label{instruction:ESPVADDU16STINCP}
\textbf{Syntax}\\
ESP.VADD.U16.ST.INCP qu,rs1,qv,qx,qy,sat

\textbf{Description}\\
This instruction performs a vector unsigned addition on 16-bit data in the two registers \lstinline{qx} and \lstinline{qy}. Then, the 8 results obtained from the calculation are saturated, and the saturated results are written to the register \lstinline{qv}.\\ During the operation, the value in the register \lstinline{qu} is stored to memory. After the access, the value in the register \lstinline{rs1} is incremented by 16.

\textbf{Operation}\\
\begin{lstlisting}
  qv[ 15:  0] = min(qx[ 15:  0] + qy[ 15:  0], 2^{16}-1)
  qv[ 31: 16] = min(qx[ 31: 16] + qy[ 31: 16], 2^{16}-1)
  ...
  qv[127:112] = min(qx[127:112] + qy[127:112], 2^{16}-1)

  qu[127:0] => store128(rs1[31:0])
  rs1[31:0] = rs1[31:0] + 16
\end{lstlisting}


\textbf{Schedule}\\
\begin{longtable}[c]{|l|l|l|}
    \hline
    \rowcolor{lightgray}
    \textbf{Operands} & \textbf{Use} & \textbf{Def} \\\hline
    \endfirsthead
    \hline
    \rowcolor{lightgray}
     \textbf{Operands} & \textbf{Use} & \textbf{Def} \\\hline
    \endhead
    \hline
    \endfoot

    qx & 1 & - \\\hline
  qy & 1 & - \\\hline
     qv & - & 1 \\\hline
     qu & 2 & - \\\hline
     rs1 & 1 & 1

\end{longtable}

\textbf{Format}\\

\begin{figure*}[!htbp]
 {\setlength{\tabcolsep}{1pt}
 \begin{center}
\begin{tabular}{@{}F@{}F@{}F@{}F@{}I@{}I@{}F@{}W@{}F@{}F@{}O}
\instbitrange{31}{29}&
\instbitrange{28}{26}&
\instbitrange{25}{23}&
\instbitrange{22}{20}&
\instbit{19}&
\instbit{18}&
\instbitrange{17}{15}&
\instbitrange{14}{13}&
\instbitrange{12}{10}&
\instbitrange{9}{7}&
\instbitrange{6}{0} \\
\hline
\multicolumn{1}{|c|}{{\scriptsize qx[2:0]}} &
\multicolumn{1}{c|}{{\scriptsize qy[2:0]}} &
\multicolumn{1}{c|}{110} &
\multicolumn{1}{c|}{{\scriptsize qv[2:0]}} &
\multicolumn{1}{c|}{{\scriptsize sat[0]}} &
\multicolumn{1}{c|}{{\scriptsize rs1[4]}} &
\multicolumn{1}{c|}{{\scriptsize rs1[2:0]}} &
\multicolumn{1}{c|}{10} &
\multicolumn{1}{c|}{{\scriptsize qu[2:0]}} &
\multicolumn{1}{c|}{X10} &
\multicolumn{1}{c|}{0011111}\\
\hline
3& 3& 3& 3& 1& 1& 3& 2& 3& 3& 7 \\
\end{tabular}
 \end{center}
 }
 \vspace{-0.1in}
 \caption[]{ESP.VADD.U16.ST.INCP} \end{figure*}


\newpage
\subsubsection{ESP.VADD.U32}\label{instruction:ESPVADDU32}
\textbf{Syntax}\\
ESP.VADD.U32 qv,qx,qy,sat

\textbf{Description}\\
This instruction performs a vector unsigned addition on 32-bit data in the two registers \lstinline{qx} and \lstinline{qy}. Then, the 4 results obtained from the calculation are saturated, and the saturated results are written to the register \lstinline{qv}.

\textbf{Operation}\\
\begin{lstlisting}
  qv[ 31:  0] = min(qx[ 31:  0] + qy[ 31:  0], 2^{32}-1)
  qv[ 63: 32] = min(qx[ 63: 32] + qy[ 63: 32], 2^{32}-1)
  ...
\end{lstlisting}

\textbf{Schedule}\\
\begin{longtable}[c]{|l|l|l|}
    \hline
    \rowcolor{lightgray}
    \textbf{Operands} & \textbf{Use} & \textbf{Def} \\\hline
    \endfirsthead
    \hline
    \rowcolor{lightgray}
     \textbf{Operands} & \textbf{Use} & \textbf{Def} \\\hline
    \endhead
    \hline
    \endfoot

    qx & 1 & - \\\hline
  qy & 1 & - \\\hline
     qv & - & 1

\end{longtable}

\textbf{Format}\\

\begin{figure*}[!htbp]
 {\setlength{\tabcolsep}{1pt}
 \begin{center}
\begin{tabular}{@{}F@{}F@{}F@{}F@{}W@{}I@{}T@{}O}
\instbitrange{31}{29}&
\instbitrange{28}{26}&
\instbitrange{25}{23}&
\instbitrange{22}{20}&
\instbitrange{19}{18}&
\instbit{17}&
\instbitrange{16}{7}&
\instbitrange{6}{0} \\
\hline
\multicolumn{1}{|c|}{{\scriptsize qx[2:0]}} &
\multicolumn{1}{c|}{{\scriptsize qy[2:0]}} &
\multicolumn{1}{c|}{1X1} &
\multicolumn{1}{c|}{{\scriptsize qv[2:0]}} &
\multicolumn{1}{c|}{X0} &
\multicolumn{1}{c|}{{\scriptsize sat[0]}} &
\multicolumn{1}{c|}{X101XX011X} &
\multicolumn{1}{c|}{0011011}\\
\hline
3& 3& 3& 3& 2& 1& 10& 7 \\
\end{tabular}
 \end{center}
 }
 \vspace{-0.1in}
 \caption[]{ESP.VADD.U32} \end{figure*}


\newpage
\subsubsection{ESP.VADD.U32.LD.INCP}\label{instruction:ESPVADDU32LDINCP}
\textbf{Syntax}\\
ESP.VADD.U32.LD.INCP qu,rs1,qv,qx,qy,sat

\textbf{Description}\\
This instruction performs a vector unsigned addition on 32-bit data in the two registers \lstinline{qx} and \lstinline{qy}. Then, the 4 results obtained from the calculation are saturated, and the saturated results are written to the register \lstinline{qv}.\\ During the operation, 16 bytes are loaded from the memory to the register \lstinline{qu}. After the access, the value in the register \lstinline{rs1} is incremented by 16.

\textbf{Operation}\\
\begin{lstlisting}
  qv[ 31:  0] = min(qx[ 31:  0] + qy[ 31:  0], 2^{32}-1)
  qv[ 63: 32] = min(qx[ 63: 32] + qy[ 63: 32], 2^{32}-1)
  ...

  qu[127:0] = load128(rs1[31:0])
  rs1[31:0] = rs1[31:0] + 16
\end{lstlisting}

\textbf{Schedule}\\
\begin{longtable}[c]{|l|l|l|}
    \hline
    \rowcolor{lightgray}
    \textbf{Operands} & \textbf{Use} & \textbf{Def} \\\hline
    \endfirsthead
    \hline
    \rowcolor{lightgray}
     \textbf{Operands} & \textbf{Use} & \textbf{Def} \\\hline
    \endhead
    \hline
    \endfoot

    qx & 1 & - \\\hline
  qy & 1 & - \\\hline
     qv & - & 1 \\\hline
     qu & - & 2 \\\hline
     rs1 & 1 & 1

\end{longtable}

\textbf{Format}\\

\begin{figure*}[!htbp]
 {\setlength{\tabcolsep}{1pt}
 \begin{center}
\begin{tabular}{@{}F@{}F@{}F@{}F@{}I@{}I@{}F@{}W@{}F@{}F@{}O}
\instbitrange{31}{29}&
\instbitrange{28}{26}&
\instbitrange{25}{23}&
\instbitrange{22}{20}&
\instbit{19}&
\instbit{18}&
\instbitrange{17}{15}&
\instbitrange{14}{13}&
\instbitrange{12}{10}&
\instbitrange{9}{7}&
\instbitrange{6}{0} \\
\hline
\multicolumn{1}{|c|}{{\scriptsize qx[2:0]}} &
\multicolumn{1}{c|}{{\scriptsize qy[2:0]}} &
\multicolumn{1}{c|}{00X} &
\multicolumn{1}{c|}{{\scriptsize qv[2:0]}} &
\multicolumn{1}{c|}{{\scriptsize sat[0]}} &
\multicolumn{1}{c|}{{\scriptsize rs1[4]}} &
\multicolumn{1}{c|}{{\scriptsize rs1[2:0]}} &
\multicolumn{1}{c|}{10} &
\multicolumn{1}{c|}{{\scriptsize qu[2:0]}} &
\multicolumn{1}{c|}{001} &
\multicolumn{1}{c|}{0011111}\\
\hline
3& 3& 3& 3& 1& 1& 3& 2& 3& 3& 7 \\
\end{tabular}
 \end{center}
 }
 \vspace{-0.1in}
 \caption[]{ESP.VADD.U32.LD.INCP} \end{figure*}


\newpage
\subsubsection{ESP.VADD.U32.ST.INCP}\label{instruction:ESPVADDU32STINCP}
\textbf{Syntax}\\
ESP.VADD.U32.ST.INCP qu,rs1,qv,qx,qy,sat

\textbf{Description}\\
This instruction performs a vector unsigned addition on 32-bit data in the two registers \lstinline{qx} and \lstinline{qy}. Then, the 4 results obtained from the calculation are saturated, and the saturated results are written to the register \lstinline{qv}.\\ During the operation, the value in the register \lstinline{qu} is stored in memory. After the access, the value in the register \lstinline{rs1} is incremented by 16.


\textbf{Operation}\\
\begin{lstlisting}
  qv[ 31:  0] = min(qx[ 31:  0] + qy[ 31:  0], 2^{32}-1)
  qv[ 63: 32] = min(qx[ 63: 32] + qy[ 63: 32], 2^{32}-1)
  ...
  qu[127:0] => store128(rs1[31:0])
  rs1[31:0] = rs1[31:0] + 16
\end{lstlisting}


\textbf{Schedule}\\
\begin{longtable}[c]{|l|l|l|}
    \hline
    \rowcolor{lightgray}
    \textbf{Operands} & \textbf{Use} & \textbf{Def} \\\hline
    \endfirsthead
    \hline
    \rowcolor{lightgray}
     \textbf{Operands} & \textbf{Use} & \textbf{Def} \\\hline
    \endhead
    \hline
    \endfoot

    qx & 1 & - \\\hline
  qy & 1 & - \\\hline
     qv & - & 1 \\\hline
     qu & 2 & - \\\hline
     rs1 & 1 & 1

\end{longtable}

\textbf{Format}\\

\begin{figure*}[!htbp]
 {\setlength{\tabcolsep}{1pt}
 \begin{center}
\begin{tabular}{@{}F@{}F@{}F@{}F@{}I@{}I@{}F@{}W@{}F@{}F@{}O}
\instbitrange{31}{29}&
\instbitrange{28}{26}&
\instbitrange{25}{23}&
\instbitrange{22}{20}&
\instbit{19}&
\instbit{18}&
\instbitrange{17}{15}&
\instbitrange{14}{13}&
\instbitrange{12}{10}&
\instbitrange{9}{7}&
\instbitrange{6}{0} \\
\hline
\multicolumn{1}{|c|}{{\scriptsize qx[2:0]}} &
\multicolumn{1}{c|}{{\scriptsize qy[2:0]}} &
\multicolumn{1}{c|}{10X} &
\multicolumn{1}{c|}{{\scriptsize qv[2:0]}} &
\multicolumn{1}{c|}{{\scriptsize sat[0]}} &
\multicolumn{1}{c|}{{\scriptsize rs1[4]}} &
\multicolumn{1}{c|}{{\scriptsize rs1[2:0]}} &
\multicolumn{1}{c|}{10} &
\multicolumn{1}{c|}{{\scriptsize qu[2:0]}} &
\multicolumn{1}{c|}{001} &
\multicolumn{1}{c|}{0011111}\\
\hline
3& 3& 3& 3& 1& 1& 3& 2& 3& 3& 7 \\
\end{tabular}
 \end{center}
 }
 \vspace{-0.1in}
 \caption[]{ESP.VADD.U32.ST.INCP} \end{figure*}


\newpage
\subsubsection{ESP.VADD.U8}\label{instruction:ESPVADDU8}
\textbf{Syntax}\\
ESP.VADD.U8 qv,qx,qy,sat

\textbf{Description}\\
This instruction performs a vector unsigned addition on 8-bit data in the two registers \lstinline{qx} and \lstinline{qy}. Then, the 16 results obtained from the calculation are saturated, and the saturated results are written to the register \lstinline{qv}.

\textbf{Operation}\\
\begin{lstlisting}
  qv[  7:  0] = min(qx[  7:  0] + qy[  7:  0], 2^{8}-1)
  qv[ 15:  8] = min(qx[ 15:  8] + qy[ 15:  8], 2^{8}-1)
  ...
  qv[127:120] = min(qx[127:120] + qy[127:120], 2^{8}-1)
\end{lstlisting}


\textbf{Schedule}\\
\begin{longtable}[c]{|l|l|l|}
    \hline
    \rowcolor{lightgray}
    \textbf{Operands} & \textbf{Use} & \textbf{Def} \\\hline
    \endfirsthead
    \hline
    \rowcolor{lightgray}
     \textbf{Operands} & \textbf{Use} & \textbf{Def} \\\hline
    \endhead
    \hline
    \endfoot

    qx & 1 & - \\\hline
  qy & 1 & - \\\hline
     qv & - & 1

\end{longtable}

\textbf{Format}\\

\begin{figure*}[!htbp]
 {\setlength{\tabcolsep}{1pt}
 \begin{center}
\begin{tabular}{@{}F@{}F@{}F@{}F@{}F@{}I@{}T@{}O}
\instbitrange{31}{29}&
\instbitrange{28}{26}&
\instbitrange{25}{23}&
\instbitrange{22}{20}&
\instbitrange{19}{17}&
\instbit{16}&
\instbitrange{15}{7}&
\instbitrange{6}{0} \\
\hline
\multicolumn{1}{|c|}{{\scriptsize qx[2:0]}} &
\multicolumn{1}{c|}{{\scriptsize qy[2:0]}} &
\multicolumn{1}{c|}{1X1} &
\multicolumn{1}{c|}{{\scriptsize qv[2:0]}} &
\multicolumn{1}{c|}{X00} &
\multicolumn{1}{c|}{{\scriptsize sat[0]}} &
\multicolumn{1}{c|}{001XX011X} &
\multicolumn{1}{c|}{0011011}\\
\hline
3& 3& 3& 3& 3& 1& 9& 7 \\
\end{tabular}
 \end{center}
 }
 \vspace{-0.1in}
 \caption[]{ESP.VADD.U8} \end{figure*}


\newpage
\subsubsection{ESP.VADD.U8.LD.INCP}\label{instruction:ESPVADDU8LDINCP}
\textbf{Syntax}\\
ESP.VADD.U8.LD.INCP qu,rs1,qv,qx,qy,sat

\textbf{Description}\\
This instruction performs a vector unsigned addition on 8-bit data in the two registers \lstinline{qx} and \lstinline{qy}. Then, the 16 results obtained from the calculation are saturated, and the saturated results are written to the register \lstinline{qv}. \\ During the operation, 16 bytes are loaded from the memory to the register \lstinline{qu}. After the access, the value in the register \lstinline{rs1} is incremented by 16.

\textbf{Operation}\\
\begin{lstlisting}
  qv[  7:  0] = min(qx[  7:  0] + qy[  7:  0], 2^{8}-1)
  qv[ 15:  8] = min(qx[ 15:  8] + qy[ 15:  8], 2^{8}-1)
  ...

  qu[127:0] = load128(rs1[31:0])
  rs1[31:0] = rs1[31:0] + 16
\end{lstlisting}


\textbf{Schedule}\\
\begin{longtable}[c]{|l|l|l|}
    \hline
    \rowcolor{lightgray}
    \textbf{Operands} & \textbf{Use} & \textbf{Def} \\\hline
    \endfirsthead
    \hline
    \rowcolor{lightgray}
     \textbf{Operands} & \textbf{Use} & \textbf{Def} \\\hline
    \endhead
    \hline
    \endfoot

    qx & 1 & - \\\hline
  qy & 1 & - \\\hline
     qv & - & 1 \\\hline
     qu & - & 2 \\\hline
     rs1 & 1 & 1

\end{longtable}

\textbf{Format}\\

\begin{figure*}[!htbp]
 {\setlength{\tabcolsep}{1pt}
 \begin{center}
\begin{tabular}{@{}F@{}F@{}F@{}F@{}I@{}I@{}F@{}W@{}F@{}F@{}O}
\instbitrange{31}{29}&
\instbitrange{28}{26}&
\instbitrange{25}{23}&
\instbitrange{22}{20}&
\instbit{19}&
\instbit{18}&
\instbitrange{17}{15}&
\instbitrange{14}{13}&
\instbitrange{12}{10}&
\instbitrange{9}{7}&
\instbitrange{6}{0} \\
\hline
\multicolumn{1}{|c|}{{\scriptsize qx[2:0]}} &
\multicolumn{1}{c|}{{\scriptsize qy[2:0]}} &
\multicolumn{1}{c|}{000} &
\multicolumn{1}{c|}{{\scriptsize qv[2:0]}} &
\multicolumn{1}{c|}{{\scriptsize sat[0]}} &
\multicolumn{1}{c|}{{\scriptsize rs1[4]}} &
\multicolumn{1}{c|}{{\scriptsize rs1[2:0]}} &
\multicolumn{1}{c|}{10} &
\multicolumn{1}{c|}{{\scriptsize qu[2:0]}} &
\multicolumn{1}{c|}{X10} &
\multicolumn{1}{c|}{0011111}\\
\hline
3& 3& 3& 3& 1& 1& 3& 2& 3& 3& 7 \\
\end{tabular}
 \end{center}
 }
 \vspace{-0.1in}
 \caption[]{ESP.VADD.U8.LD.INCP} \end{figure*}


\newpage
\subsubsection{ESP.VADD.U8.ST.INCP}\label{instruction:ESPVADDU8STINCP}
\textbf{Syntax}\\
ESP.VADD.U8.ST.INCP qu,rs1,qv,qx,qy,sat

\textbf{Description}\\
This instruction performs a vector unsigned addition on 8-bit data in the two registers \lstinline{qx} and \lstinline{qy}. Then, the 16 results obtained from the calculation are saturated, and the saturated results are written to the register \lstinline{qv}. \\ During the operation, the value in the register \lstinline{qu} is stored in memory. After the access, the value in the register \lstinline{rs1} is incremented by 16.

\textbf{Operation}\\
\begin{lstlisting}
  qv[  7:  0] = min(qx[  7:  0] + qy[  7:  0], 2^{8}-1)
  qv[ 15:  8] = min(qx[ 15:  8] + qy[ 15:  8], 2^{8}-1)
  ...

  qu[127:0] => store128(rs1[31:0])
  rs1 [31:0] = rs1[31:0] + 16
\end{lstlisting}

\textbf{Schedule}\\
\begin{longtable}[c]{|l|l|l|}
    \hline
    \rowcolor{lightgray}
    \textbf{Operands} & \textbf{Use} & \textbf{Def} \\\hline
    \endfirsthead
    \hline
    \rowcolor{lightgray}
     \textbf{Operands} & \textbf{Use} & \textbf{Def} \\\hline
    \endhead
    \hline
    \endfoot

    qx & 1 & - \\\hline
  qy & 1 & - \\\hline
     qv & - & 1 \\\hline
     qu & 2 & - \\\hline
     rs1 & 1 & 1

\end{longtable}

\textbf{Format}\\

\begin{figure*}[!htbp]
 {\setlength{\tabcolsep}{1pt}
 \begin{center}
\begin{tabular}{@{}F@{}F@{}F@{}F@{}I@{}I@{}F@{}W@{}F@{}F@{}O}
\instbitrange{31}{29}&
\instbitrange{28}{26}&
\instbitrange{25}{23}&
\instbitrange{22}{20}&
\instbit{19}&
\instbit{18}&
\instbitrange{17}{15}&
\instbitrange{14}{13}&
\instbitrange{12}{10}&
\instbitrange{9}{7}&
\instbitrange{6}{0} \\
\hline
\multicolumn{1}{|c|}{{\scriptsize qx[2:0]}} &
\multicolumn{1}{c|}{{\scriptsize qy[2:0]}} &
\multicolumn{1}{c|}{100} &
\multicolumn{1}{c|}{{\scriptsize qv[2:0]}} &
\multicolumn{1}{c|}{{\scriptsize sat[0]}} &
\multicolumn{1}{c|}{{\scriptsize rs1[4]}} &
\multicolumn{1}{c|}{{\scriptsize rs1[2:0]}} &
\multicolumn{1}{c|}{10} &
\multicolumn{1}{c|}{{\scriptsize qu[2:0]}} &
\multicolumn{1}{c|}{X10} &
\multicolumn{1}{c|}{0011111}\\
\hline
3& 3& 3& 3& 1& 1& 3& 2& 3& 3& 7 \\
\end{tabular}
 \end{center}
 }
 \vspace{-0.1in}
 \caption[]{ESP.VADD.U8.ST.INCP} \end{figure*}


\newpage
\subsubsection{ESP.VCLAMP.S16}\label{instruction:ESPVCLAMPS16}
\textbf{Syntax}\\
ESP.VCLAMP.S16 qz,qx,sel16

\textbf{Description}\\
This instruction performs clamping for 16-bit signed data.

\textbf{Operation}\\
\begin{lstlisting}[escapeinside=`?]
    qz[ 15:  0] = min( max(qx[ 15:  0], -2^{sel16}), 2^{sel16}-1 )
    qz[ 31: 16] = min( max(qx[ 31: 16], -2^{sel16}), 2^{sel16}-1 )
    qz[ 47: 32] = min( max(qx[ 47: 32], -2^{sel16}), 2^{sel16}-1 )
    qz[ 63: 48] = min( max(qx[ 63: 48], -2^{sel16}), 2^{sel16}-1 )
    qz[ 79: 64] = min( max(qx[ 79: 64], -2^{sel16}), 2^{sel16}-1 )
    qz[ 95: 80] = min( max(qx[ 95: 80], -2^{sel16}), 2^{sel16}-1 )
    qz[111: 96] = min( max(qx[111: 96], -2^{sel16}), 2^{sel16}-1 )
    qz[127:112] = min( max(qx[127:112], -2^{sel16}), 2^{sel16}-1 )
\end{lstlisting}


\textbf{Schedule}\\
\begin{longtable}[c]{|l|l|l|}
    \hline
    \rowcolor{lightgray}
    \textbf{Operands} & \textbf{Use} & \textbf{Def} \\\hline
    \endfirsthead
    \hline
    \rowcolor{lightgray}
     \textbf{Operands} & \textbf{Use} & \textbf{Def} \\\hline
    \endhead
    \hline
    \endfoot

    qx & 1 & - \\\hline
  qy & 1 & - \\\hline
     qz & - & 1

\end{longtable}

\textbf{Format}\\

\begin{figure*}[!htbp]
 {\setlength{\tabcolsep}{1pt}
 \begin{center}
\begin{tabular}{@{}F@{}T@{}Y@{}V@{}F@{}O}
\instbitrange{31}{29}&
\instbitrange{28}{19}&
\instbitrange{18}{15}&
\instbitrange{14}{10}&
\instbitrange{9}{7}&
\instbitrange{6}{0} \\
\hline
\multicolumn{1}{|c|}{{\scriptsize qx[2:0]}} &
\multicolumn{1}{c|}{0XXXXXXXXX} &
\multicolumn{1}{c|}{{\scriptsize sel16[3:0]}} &
\multicolumn{1}{c|}{101XX} &
\multicolumn{1}{c|}{{\scriptsize qz[2:0]}} &
\multicolumn{1}{c|}{0011011}\\
\hline
3& 10& 4& 5& 3& 7 \\
\end{tabular}
 \end{center}
 }
 \vspace{-0.1in}
 \caption[]{ESP.VCLAMP.S16} \end{figure*}


\newpage
\subsubsection{ESP.VMAX.S16}\label{instruction:ESPVMAXS16}
\textbf{Syntax}\\
ESP.VMAX.S16 qz,qx,qy

\textbf{Description}\\
This instruction compares the numerical values of the eight 16-bit signed vector data segments in the registers \lstinline{qx} and \lstinline{qy}. The data segment with the larger value is written into the corresponding 16-bit data segment in the register \lstinline{qz}.\\
\textbf{Operation}
\begin{lstlisting}[escapeinside=`?]
  qz[ 15:  0] = (qx[ 15:  0]>=qy[ 15:  0]) ? qx[ 15:  0] : qy[ 15:  0]
  qz[ 31: 16] = (qx[ 31: 16]>=qy[ 31: 16]) ? qx[ 31: 16] : qy[ 31: 16]
  ...
  qz[127:112] = (qx[127:112]>=qy[127:112]) ? qx[127:112] : qy[127:112]
\end{lstlisting}


\textbf{Schedule}\\
\begin{longtable}[c]{|l|l|l|}
    \hline
    \rowcolor{lightgray}
    \textbf{Operands} & \textbf{Use} & \textbf{Def} \\\hline
    \endfirsthead
    \hline
    \rowcolor{lightgray}
     \textbf{Operands} & \textbf{Use} & \textbf{Def} \\\hline
    \endhead
    \hline
    \endfoot

    qx & 1 & - \\\hline
  qy & 1 & - \\\hline
     qz & - & 1

\end{longtable}

\textbf{Format}\\

\begin{figure*}[!htbp]
 {\setlength{\tabcolsep}{1pt}
 \begin{center}
\begin{tabular}{@{}F@{}F@{}K@{}F@{}O}
\instbitrange{31}{29}&
\instbitrange{28}{26}&
\instbitrange{25}{10}&
\instbitrange{9}{7}&
\instbitrange{6}{0} \\
\hline
\multicolumn{1}{|c|}{{\scriptsize qx[2:0]}} &
\multicolumn{1}{c|}{{\scriptsize qy[2:0]}} &
\multicolumn{1}{c|}{110100X11000011X} &
\multicolumn{1}{c|}{{\scriptsize qz[2:0]}} &
\multicolumn{1}{c|}{0011111}\\
\hline
3& 3& 16& 3& 7 \\
\end{tabular}
 \end{center}
 }
 \vspace{-0.1in}
 \caption[]{ESP.VMAX.S16} \end{figure*}


\newpage
\subsubsection{ESP.VMAX.S16.LD.INCP}\label{instruction:ESPVMAXS16LDINCP}
\textbf{Syntax}\\
ESP.VMAX.S16.LD.INCP qu,rs1,qz,qx,qy

\textbf{Description}\\
This instruction compares the numerical values of the eight 16-bit signed vector data segments in the registers \lstinline{qx} and \lstinline{qy}. The data segment with the larger value is written into the corresponding 16-bit data segment in the register \lstinline{qz}.\\During the operation, 16 bytes are loaded from the memory into the register \lstinline{qu}. After the access, the value in the register \lstinline{rs1} is incremented by 16.\\
\textbf{Operation}
\begin{lstlisting}[escapeinside=`?]
  qz[ 15:  0] = (qx[ 15:  0]>=qy[ 15:  0]) ? qx[ 15:  0] : qy[ 15:  0]
  qz[ 31: 16] = (qx[ 31: 16]>=qy[ 31: 16]) ? qx[ 31: 16] : qy[ 31: 16]
  ...
  qz[127:112] = (qx[127:112]>=qy[127:112]) ? qx[127:112] : qy[127:112]

  qu[127:0] = load128(rs1[31:0])
  rs1[31:0] = rs1[31:0] + 16
\end{lstlisting}

\textbf{Schedule}\\
\begin{longtable}[c]{|l|l|l|}
    \hline
    \rowcolor{lightgray}
    \textbf{Operands} & \textbf{Use} & \textbf{Def} \\\hline
    \endfirsthead
    \hline
    \rowcolor{lightgray}
     \textbf{Operands} & \textbf{Use} & \textbf{Def} \\\hline
    \endhead
    \hline
    \endfoot

    qx & 1 & - \\\hline
  qy & 1 & - \\\hline
     qz & - & 1 \\\hline
     qu & - & 2 \\\hline
     rs1 & 1 & 1

\end{longtable}

\textbf{Format}\\

\begin{figure*}[!htbp]
 {\setlength{\tabcolsep}{1pt}
 \begin{center}
\begin{tabular}{@{}F@{}F@{}O@{}I@{}F@{}W@{}F@{}F@{}O}
\instbitrange{31}{29}&
\instbitrange{28}{26}&
\instbitrange{25}{19}&
\instbit{18}&
\instbitrange{17}{15}&
\instbitrange{14}{13}&
\instbitrange{12}{10}&
\instbitrange{9}{7}&
\instbitrange{6}{0} \\
\hline
\multicolumn{1}{|c|}{{\scriptsize qx[2:0]}} &
\multicolumn{1}{c|}{{\scriptsize qy[2:0]}} &
\multicolumn{1}{c|}{0X0110X} &
\multicolumn{1}{c|}{{\scriptsize rs1[4]}} &
\multicolumn{1}{c|}{{\scriptsize rs1[2:0]}} &
\multicolumn{1}{c|}{X0} &
\multicolumn{1}{c|}{{\scriptsize qu[2:0]}} &
\multicolumn{1}{c|}{{\scriptsize qz[2:0]}} &
\multicolumn{1}{c|}{1011111}\\
\hline
3& 3& 7& 1& 3& 2& 3& 3& 7 \\
\end{tabular}
 \end{center}
 }
 \vspace{-0.1in}
 \caption[]{ESP.VMAX.S16.LD.INCP} \end{figure*}


\newpage
\subsubsection{ESP.VMAX.S16.ST.INCP}\label{instruction:ESPVMAXS16STINCP}
\textbf{Syntax}\\
ESP.VMAX.S16.ST.INCP qu,rs1,qz,qx,qy

\textbf{Description}\\
This instruction compares the numerical values of the eight 16-bit signed vector data segments in the registers \lstinline{qx} and \lstinline{qy}. The data segment with the larger value is written into the corresponding 16-bit data segment in the register \lstinline{qz}.\\During the operation, the value in the register \lstinline{qu} is stored in memory. After the access, the value in the register \lstinline{rs1} is incremented by 16.\\
\textbf{Operation}
\begin{lstlisting}[escapeinside=`?]
  qz[ 15:  0] = (qx[ 15:  0]>=qy[ 15:  0]) ? qx[ 15:  0] : qy[ 15:  0]
  qz[ 31: 16] = (qx[ 31: 16]>=qy[ 31: 16]) ? qx[ 31: 16] : qy[ 31: 16]
  ...
  qz[127:112] = (qx[127:112]>=qy[127:112]) ? qx[127:112] : qy[127:112]

  qu[127:0] => store128(rs1[31:0])
  rs1 [31:0] = rs1[31:0] + 16
\end{lstlisting}

\textbf{Schedule}\\
\begin{longtable}[c]{|l|l|l|}
    \hline
    \rowcolor{lightgray}
    \textbf{Operands} & \textbf{Use} & \textbf{Def} \\\hline
    \endfirsthead
    \hline
    \rowcolor{lightgray}
     \textbf{Operands} & \textbf{Use} & \textbf{Def} \\\hline
    \endhead
    \hline
    \endfoot

    qx & 1 & - \\\hline
  qy & 1 & - \\\hline
     qz & - & 1 \\\hline
     qu & 2 & - \\\hline
     rs1 & 1 & 1

\end{longtable}

\textbf{Format}\\

\begin{figure*}[!htbp]
 {\setlength{\tabcolsep}{1pt}
 \begin{center}
\begin{tabular}{@{}F@{}F@{}O@{}I@{}F@{}W@{}F@{}F@{}O}
\instbitrange{31}{29}&
\instbitrange{28}{26}&
\instbitrange{25}{19}&
\instbit{18}&
\instbitrange{17}{15}&
\instbitrange{14}{13}&
\instbitrange{12}{10}&
\instbitrange{9}{7}&
\instbitrange{6}{0} \\
\hline
\multicolumn{1}{|c|}{{\scriptsize qx[2:0]}} &
\multicolumn{1}{c|}{{\scriptsize qy[2:0]}} &
\multicolumn{1}{c|}{0X1110X} &
\multicolumn{1}{c|}{{\scriptsize rs1[4]}} &
\multicolumn{1}{c|}{{\scriptsize rs1[2:0]}} &
\multicolumn{1}{c|}{X0} &
\multicolumn{1}{c|}{{\scriptsize qu[2:0]}} &
\multicolumn{1}{c|}{{\scriptsize qz[2:0]}} &
\multicolumn{1}{c|}{1011111}\\
\hline
3& 3& 7& 1& 3& 2& 3& 3& 7 \\
\end{tabular}
 \end{center}
 }
 \vspace{-0.1in}
 \caption[]{ESP.VMAX.S16.ST.INCP} \end{figure*}


\newpage
\subsubsection{ESP.VMAX.S32}\label{instruction:ESPVMAXS32}
\textbf{Syntax}\\
ESP.VMAX.S32 qz,qx,qy

\textbf{Description}\\
This instruction compares the numerical values of the four 32-bit signed vector data segments in the registers \lstinline{qx} and \lstinline{qy}. The data segment with the larger value is written into the corresponding 32-bit data segment in the register \lstinline{qz}.\\
\textbf{Operation}
\begin{lstlisting}[escapeinside=`?]
  qz[ 31:  0] = (qx[ 31:  0]>=qy[ 31:  0]) ? qx[ 31:  0] : qy[ 31:  0]
  qz[ 63: 32] = (qx[ 63: 32]>=qy[ 63: 32]) ? qx[ 63: 32] : qy[ 63: 32]
  ...
  qz[127: 96] = (qx[127: 96]>=qy[127: 96]) ? qx[127: 96] : qy[127: 96]
\end{lstlisting}

\textbf{Schedule}\\
\begin{longtable}[c]{|l|l|l|}
    \hline
    \rowcolor{lightgray}
    \textbf{Operands} & \textbf{Use} & \textbf{Def} \\\hline
    \endfirsthead
    \hline
    \rowcolor{lightgray}
     \textbf{Operands} & \textbf{Use} & \textbf{Def} \\\hline
    \endhead
    \hline
    \endfoot

    qx & 1 & - \\\hline
  qy & 1 & - \\\hline
     qz & - & 1

\end{longtable}

\textbf{Format}\\

\begin{figure*}[!htbp]
 {\setlength{\tabcolsep}{1pt}
 \begin{center}
\begin{tabular}{@{}F@{}F@{}K@{}F@{}O}
\instbitrange{31}{29}&
\instbitrange{28}{26}&
\instbitrange{25}{10}&
\instbitrange{9}{7}&
\instbitrange{6}{0} \\
\hline
\multicolumn{1}{|c|}{{\scriptsize qx[2:0]}} &
\multicolumn{1}{c|}{{\scriptsize qy[2:0]}} &
\multicolumn{1}{c|}{110100X1X100011X} &
\multicolumn{1}{c|}{{\scriptsize qz[2:0]}} &
\multicolumn{1}{c|}{0011111}\\
\hline
3& 3& 16& 3& 7 \\
\end{tabular}
 \end{center}
 }
 \vspace{-0.1in}
 \caption[]{ESP.VMAX.S32} \end{figure*}


\newpage
\subsubsection{ESP.VMAX.S32.LD.INCP}\label{instruction:ESPVMAXS32LDINCP}
\textbf{Syntax}\\
ESP.VMAX.S32.LD.INCP qu,rs1,qz,qx,qy

\textbf{Description}\\
This instruction compares the numerical values of the four 32-bit signed vector data segments in registers \lstinline{qx} and \lstinline{qy}. The data segment with the larger value is written into the corresponding 32-bit data segment in register \lstinline{qz}.\\During the operation, 16 bytes are loaded from the memory into register \lstinline{qu}. After the access, the value in register \lstinline{rs1} is incremented by 16.\\
\textbf{Operation}
\begin{lstlisting}[escapeinside=`?]
  qz[ 31:  0] = (qx[ 31:  0]>=qy[ 31:  0]) ? qx[ 31:  0] : qy[ 31:  0]
  qz[ 63: 32] = (qx[ 63: 32]>=qy[ 63: 32]) ? qx[ 63: 32] : qy[ 63: 32]
  ...
  qz[127: 96] = (qx[127: 96]>=qy[127: 96]) ? qx[127: 96] : qy[127: 96]

   qu[127:0] = load128(rs1[31:0])
  rs1[31:0] = rs1[31:0] + 16
\end{lstlisting}


\textbf{Schedule}\\
\begin{longtable}[c]{|l|l|l|}
    \hline
    \rowcolor{lightgray}
    \textbf{Operands} & \textbf{Use} & \textbf{Def} \\\hline
    \endfirsthead
    \hline
    \rowcolor{lightgray}
     \textbf{Operands} & \textbf{Use} & \textbf{Def} \\\hline
    \endhead
    \hline
    \endfoot

    qx & 1 & - \\\hline
  qy & 1 & - \\\hline
     qz & - & 1 \\\hline
     qu & - & 2 \\\hline
     rs1 & 1 & 1

\end{longtable}

\textbf{Format}\\

\begin{figure*}[!htbp]
 {\setlength{\tabcolsep}{1pt}
 \begin{center}
\begin{tabular}{@{}F@{}F@{}O@{}I@{}F@{}W@{}F@{}F@{}O}
\instbitrange{31}{29}&
\instbitrange{28}{26}&
\instbitrange{25}{19}&
\instbit{18}&
\instbitrange{17}{15}&
\instbitrange{14}{13}&
\instbitrange{12}{10}&
\instbitrange{9}{7}&
\instbitrange{6}{0} \\
\hline
\multicolumn{1}{|c|}{{\scriptsize qx[2:0]}} &
\multicolumn{1}{c|}{{\scriptsize qy[2:0]}} &
\multicolumn{1}{c|}{0X01X1X} &
\multicolumn{1}{c|}{{\scriptsize rs1[4]}} &
\multicolumn{1}{c|}{{\scriptsize rs1[2:0]}} &
\multicolumn{1}{c|}{X0} &
\multicolumn{1}{c|}{{\scriptsize qu[2:0]}} &
\multicolumn{1}{c|}{{\scriptsize qz[2:0]}} &
\multicolumn{1}{c|}{1011111}\\
\hline
3& 3& 7& 1& 3& 2& 3& 3& 7 \\
\end{tabular}
 \end{center}
 }
 \vspace{-0.1in}
 \caption[]{ESP.VMAX.S32.LD.INCP} \end{figure*}


\newpage
\subsubsection{ESP.VMAX.S32.ST.INCP}\label{instruction:ESPVMAXS32STINCP}
\textbf{Syntax}\\
ESP.VMAX.S32.ST.INCP qu,rs1,qz,qx,qy

\textbf{Description}\\
This instruction compares the numerical values of the four 32-bit signed vector data segments in the registers \lstinline{qx} and \lstinline{qy}. The data segment with the larger value is written into the corresponding 32-bit data segment in the register \lstinline{qz}.\\During the operation, the value in the register \lstinline{qu} is stored in memory. After the access, the value in the register \lstinline{rs1} is incremented by 16.\\
\textbf{Operation}
\begin{lstlisting}[escapeinside=`?]
  qz[ 31:  0] = (qx[ 31:  0]>=qy[ 31:  0]) ? qx[ 31:  0] : qy[ 31:  0]
  qz[ 63: 32] = (qx[ 63: 32]>=qy[ 63: 32]) ? qx[ 63: 32] : qy[ 63: 32]
  ...
  qz[127: 96] = (qx[127: 96]>=qy[127: 96]) ? qx[127: 96] : qy[127: 96]

  qu[127:0] => store128(rs1[31:0])
  rs1 [31:0] = rs1[31:0] + 16
\end{lstlisting}

\textbf{Schedule}\\
\begin{longtable}[c]{|l|l|l|}
    \hline
    \rowcolor{lightgray}
    \textbf{Operands} & \textbf{Use} & \textbf{Def} \\\hline
    \endfirsthead
    \hline
    \rowcolor{lightgray}
     \textbf{Operands} & \textbf{Use} & \textbf{Def} \\\hline
    \endhead
    \hline
    \endfoot

    qx & 1 & - \\\hline
  qy & 1 & - \\\hline
     qz & - & 1 \\\hline
     qu & 2 & - \\\hline
     rs1 & 1 & 1

\end{longtable}

\textbf{Format}\\

\begin{figure*}[!htbp]
 {\setlength{\tabcolsep}{1pt}
 \begin{center}
\begin{tabular}{@{}F@{}F@{}O@{}I@{}F@{}W@{}F@{}F@{}O}
\instbitrange{31}{29}&
\instbitrange{28}{26}&
\instbitrange{25}{19}&
\instbit{18}&
\instbitrange{17}{15}&
\instbitrange{14}{13}&
\instbitrange{12}{10}&
\instbitrange{9}{7}&
\instbitrange{6}{0} \\
\hline
\multicolumn{1}{|c|}{{\scriptsize qx[2:0]}} &
\multicolumn{1}{c|}{{\scriptsize qy[2:0]}} &
\multicolumn{1}{c|}{0X11X1X} &
\multicolumn{1}{c|}{{\scriptsize rs1[4]}} &
\multicolumn{1}{c|}{{\scriptsize rs1[2:0]}} &
\multicolumn{1}{c|}{X0} &
\multicolumn{1}{c|}{{\scriptsize qu[2:0]}} &
\multicolumn{1}{c|}{{\scriptsize qz[2:0]}} &
\multicolumn{1}{c|}{1011111}\\
\hline
3& 3& 7& 1& 3& 2& 3& 3& 7 \\
\end{tabular}
 \end{center}
 }
 \vspace{-0.1in}
 \caption[]{ESP.VMAX.S32.ST.INCP} \end{figure*}


\newpage
\subsubsection{ESP.VMAX.S8}\label{instruction:ESPVMAXS8}
\textbf{Syntax}\\
ESP.VMAX.S8 qz,qx,qy

\textbf{Description}\\
This instruction compares the numerical values of the 16 8-bit signed vector data segments in registers \lstinline{qx} and \lstinline{qy}. The data segment with the larger value is written into the corresponding 8-bit data segment in the register \lstinline{qz}.\\
\textbf{Operation}
\begin{lstlisting}[escapeinside=`?]
  qz[  7:  0] = (qx[  7:  0]>=qy[  7:  0]) ? qx[  7:  0] : qy[  7:  0]
  qz[ 15:  8] = (qx[ 15:  8]>=qy[ 15:  8]) ? qx[ 15:  8] : qy[ 15:  8]
  ...
  qz[127:120] = (qx[127:120]>=qy[127:120]) ? qx[127:120] : qy[127:120]
\end{lstlisting}

\textbf{Schedule}\\
\begin{longtable}[c]{|l|l|l|}
    \hline
    \rowcolor{lightgray}
    \textbf{Operands} & \textbf{Use} & \textbf{Def} \\\hline
    \endfirsthead
    \hline
    \rowcolor{lightgray}
     \textbf{Operands} & \textbf{Use} & \textbf{Def} \\\hline
    \endhead
    \hline
    \endfoot

    qx & 1 & - \\\hline
  qy & 1 & - \\\hline
     qz & - & 1

\end{longtable}

\textbf{Format}\\

\begin{figure*}[!htbp]
 {\setlength{\tabcolsep}{1pt}
 \begin{center}
\begin{tabular}{@{}F@{}F@{}K@{}F@{}O}
\instbitrange{31}{29}&
\instbitrange{28}{26}&
\instbitrange{25}{10}&
\instbitrange{9}{7}&
\instbitrange{6}{0} \\
\hline
\multicolumn{1}{|c|}{{\scriptsize qx[2:0]}} &
\multicolumn{1}{c|}{{\scriptsize qy[2:0]}} &
\multicolumn{1}{c|}{110100X01000011X} &
\multicolumn{1}{c|}{{\scriptsize qz[2:0]}} &
\multicolumn{1}{c|}{0011111}\\
\hline
3& 3& 16& 3& 7 \\
\end{tabular}
 \end{center}
 }
 \vspace{-0.1in}
 \caption[]{ESP.VMAX.S8} \end{figure*}


\newpage
\subsubsection{ESP.VMAX.S8.LD.INCP}\label{instruction:ESPVMAXS8LDINCP}
\textbf{Syntax}\\
ESP.VMAX.S8.LD.INCP qu,rs1,qz,qx,qy

\textbf{Description}\\
This instruction compares the numerical values of the 16 8-bit signed vector data segments in registers \lstinline{qx} and \lstinline{qy}. The data segment with the larger value is written into the corresponding 8-bit data segment in register \lstinline{qz}.\\During the operation, 16 bytes are loaded from the memory into register \lstinline{qu}. After the access, the value in register \lstinline{rs1} is incremented by 16.\\
\textbf{Operation}
\begin{lstlisting}[escapeinside=`?]
  qz[  7:  0] = (qx[  7:  0]>=qy[  7:  0]) ? qx[  7:  0] : qy[  7:  0]
  qz[ 15:  8] = (qx[ 15:  8]>=qy[ 15:  8]) ? qx[ 15:  8] : qy[ 15:  8]
  ...
  qz[127:120] = (qx[127:120]>=qy[127:120]) ? qx[127:120] : qy[127:120]

  qu[127:0] = load128(rs1[31:0])
  rs1[31:0] = rs1[31:0] + 16
\end{lstlisting}


\textbf{Schedule}\\
\begin{longtable}[c]{|l|l|l|}
    \hline
    \rowcolor{lightgray}
    \textbf{Operands} & \textbf{Use} & \textbf{Def} \\\hline
    \endfirsthead
    \hline
    \rowcolor{lightgray}
     \textbf{Operands} & \textbf{Use} & \textbf{Def} \\\hline
    \endhead
    \hline
    \endfoot

    qx & 1 & - \\\hline
  qy & 1 & - \\\hline
     qz & - & 1 \\\hline
     qu & - & 2 \\\hline
     rs1 & 1 & 1

\end{longtable}

\textbf{Format}\\

\begin{figure*}[!htbp]
 {\setlength{\tabcolsep}{1pt}
 \begin{center}
\begin{tabular}{@{}F@{}F@{}O@{}I@{}F@{}W@{}F@{}F@{}O}
\instbitrange{31}{29}&
\instbitrange{28}{26}&
\instbitrange{25}{19}&
\instbit{18}&
\instbitrange{17}{15}&
\instbitrange{14}{13}&
\instbitrange{12}{10}&
\instbitrange{9}{7}&
\instbitrange{6}{0} \\
\hline
\multicolumn{1}{|c|}{{\scriptsize qx[2:0]}} &
\multicolumn{1}{c|}{{\scriptsize qy[2:0]}} &
\multicolumn{1}{c|}{0X0010X} &
\multicolumn{1}{c|}{{\scriptsize rs1[4]}} &
\multicolumn{1}{c|}{{\scriptsize rs1[2:0]}} &
\multicolumn{1}{c|}{X0} &
\multicolumn{1}{c|}{{\scriptsize qu[2:0]}} &
\multicolumn{1}{c|}{{\scriptsize qz[2:0]}} &
\multicolumn{1}{c|}{1011111}\\
\hline
3& 3& 7& 1& 3& 2& 3& 3& 7 \\
\end{tabular}
 \end{center}
 }
 \vspace{-0.1in}
 \caption[]{ESP.VMAX.S8.LD.INCP} \end{figure*}


\newpage
\subsubsection{ESP.VMAX.S8.ST.INCP}\label{instruction:ESPVMAXS8STINCP}
\textbf{Syntax}\\
ESP.VMAX.S8.ST.INCP qu,rs1,qz,qx,qy

\textbf{Description}\\
This instruction compares the numerical values of the 16 8-bit vector data segments in registers \lstinline{qx} and \lstinline{qy}. The data segment with the larger value is written into the corresponding 8-bit data segment in register \lstinline{qz}.\\During the operation, the value in register \lstinline{qu} is stored to memory. After the access, the value in register \lstinline{rs1} is incremented by 16.\\
\textbf{Operation}
\begin{lstlisting}[escapeinside=`?]
  qz[  7:  0] = (qx[  7:  0]>=qy[  7:  0]) ? qx[  7:  0] : qy[  7:  0]
  qz[ 15:  8] = (qx[ 15:  8]>=qy[ 15:  8]) ? qx[ 15:  8] : qy[ 15:  8]
  ...
  qz[127:120] = (qx[127:120]>=qy[127:120]) ? qx[127:120] : qy[127:120]

  qu[127:0] => store128(rs1[31:0])
  rs1 [31:0] = rs1[31:0] + 16
\end{lstlisting}

\textbf{Schedule}\\
\begin{longtable}[c]{|l|l|l|}
    \hline
    \rowcolor{lightgray}
    \textbf{Operands} & \textbf{Use} & \textbf{Def} \\\hline
    \endfirsthead
    \hline
    \rowcolor{lightgray}
     \textbf{Operands} & \textbf{Use} & \textbf{Def} \\\hline
    \endhead
    \hline
    \endfoot

    qx & 1 & - \\\hline
  qy & 1 & - \\\hline
     qz & - & 1 \\\hline
     qu & 2 & - \\\hline
     rs1 & 1 & 1

\end{longtable}

\textbf{Format}\\

\begin{figure*}[!htbp]
 {\setlength{\tabcolsep}{1pt}
 \begin{center}
\begin{tabular}{@{}F@{}F@{}O@{}I@{}F@{}W@{}F@{}F@{}O}
\instbitrange{31}{29}&
\instbitrange{28}{26}&
\instbitrange{25}{19}&
\instbit{18}&
\instbitrange{17}{15}&
\instbitrange{14}{13}&
\instbitrange{12}{10}&
\instbitrange{9}{7}&
\instbitrange{6}{0} \\
\hline
\multicolumn{1}{|c|}{{\scriptsize qx[2:0]}} &
\multicolumn{1}{c|}{{\scriptsize qy[2:0]}} &
\multicolumn{1}{c|}{0X1010X} &
\multicolumn{1}{c|}{{\scriptsize rs1[4]}} &
\multicolumn{1}{c|}{{\scriptsize rs1[2:0]}} &
\multicolumn{1}{c|}{X0} &
\multicolumn{1}{c|}{{\scriptsize qu[2:0]}} &
\multicolumn{1}{c|}{{\scriptsize qz[2:0]}} &
\multicolumn{1}{c|}{1011111}\\
\hline
3& 3& 7& 1& 3& 2& 3& 3& 7 \\
\end{tabular}
 \end{center}
 }
 \vspace{-0.1in}
 \caption[]{ESP.VMAX.S8.ST.INCP} \end{figure*}


\newpage
\subsubsection{ESP.VMAX.U16}\label{instruction:ESPVMAXU16}
\textbf{Syntax}\\
ESP.VMAX.U16 qz,qx,qy

\textbf{Description}\\
This instruction compares the numerical values of the eight 16-bit unsigned vector data segments in the registers \lstinline{qx} and \lstinline{qy}. The data segment with the larger value is written into the corresponding 16-bit data segment in the register \lstinline{qz}.\\
\textbf{Operation}
\begin{lstlisting}[escapeinside=`?]
  qz[ 15:  0] = (qx[ 15:  0]>=qy[ 15:  0]) ? qx[ 15:  0] : qy[ 15:  0]
  qz[ 31: 16] = (qx[ 31: 16]>=qy[ 31: 16]) ? qx[ 31: 16] : qy[ 31: 16]
  ...
  qz[127:112] = (qx[127:112]>=qy[127:112]) ? qx[127:112] : qy[127:112]
\end{lstlisting}


\textbf{Schedule}\\
\begin{longtable}[c]{|l|l|l|}
    \hline
    \rowcolor{lightgray}
    \textbf{Operands} & \textbf{Use} & \textbf{Def} \\\hline
    \endfirsthead
    \hline
    \rowcolor{lightgray}
     \textbf{Operands} & \textbf{Use} & \textbf{Def} \\\hline
    \endhead
    \hline
    \endfoot

    qx & 1 & - \\\hline
  qy & 1 & - \\\hline
     qz & - & 1

\end{longtable}

\textbf{Format}\\

\begin{figure*}[!htbp]
 {\setlength{\tabcolsep}{1pt}
 \begin{center}
\begin{tabular}{@{}F@{}F@{}K@{}F@{}O}
\instbitrange{31}{29}&
\instbitrange{28}{26}&
\instbitrange{25}{10}&
\instbitrange{9}{7}&
\instbitrange{6}{0} \\
\hline
\multicolumn{1}{|c|}{{\scriptsize qx[2:0]}} &
\multicolumn{1}{c|}{{\scriptsize qy[2:0]}} &
\multicolumn{1}{c|}{110100X10000011X} &
\multicolumn{1}{c|}{{\scriptsize qz[2:0]}} &
\multicolumn{1}{c|}{0011111}\\
\hline
3& 3& 16& 3& 7 \\
\end{tabular}
 \end{center}
 }
 \vspace{-0.1in}
 \caption[]{ESP.VMAX.U16} \end{figure*}


\newpage
\subsubsection{ESP.VMAX.U16.LD.INCP}\label{instruction:ESPVMAXU16LDINCP}
\textbf{Syntax}\\
ESP.VMAX.U16.LD.INCP qu,rs1,qz,qx,qy

\textbf{Description}\\
This instruction compares the numerical values of the eight 16-bit unsigned vector data segments in the registers \lstinline{qx} and \lstinline{qy}. The data segment with the larger value is written into the corresponding 16-bit data segment in the register \lstinline{qz}.\\During the operation, 16 bytes are loaded from the memory into the register \lstinline{qu}. After the access, the value in the register \lstinline{rs1} is incremented by 16.\\
\textbf{Operation}
\begin{lstlisting}[escapeinside=`?]
  qz[ 15:  0] = (qx[ 15:  0]>=qy[ 15:  0]) ? qx[ 15:  0] : qy[ 15:  0]
  qz[ 31: 16] = (qx[ 31: 16]>=qy[ 31: 16]) ? qx[ 31: 16] : qy[ 31: 16]
  ...
  qz[127:112] = (qx[127:112]>=qy[127:112]) ? qx[127:112] : qy[127:112]

  qu[127:0] = load128(rs1[31:0])
  rs1[31:0] = rs1[31:0] + 16
\end{lstlisting}


\textbf{Schedule}\\
\begin{longtable}[c]{|l|l|l|}
    \hline
    \rowcolor{lightgray}
    \textbf{Operands} & \textbf{Use} & \textbf{Def} \\\hline
    \endfirsthead
    \hline
    \rowcolor{lightgray}
     \textbf{Operands} & \textbf{Use} & \textbf{Def} \\\hline
    \endhead
    \hline
    \endfoot

    qx & 1 & - \\\hline
  qy & 1 & - \\\hline
     qz & - & 1 \\\hline
     qu & - & 2 \\\hline
     rs1 & 1 & 1

\end{longtable}

\textbf{Format}\\

\begin{figure*}[!htbp]
 {\setlength{\tabcolsep}{1pt}
 \begin{center}
\begin{tabular}{@{}F@{}F@{}O@{}I@{}F@{}W@{}F@{}F@{}O}
\instbitrange{31}{29}&
\instbitrange{28}{26}&
\instbitrange{25}{19}&
\instbit{18}&
\instbitrange{17}{15}&
\instbitrange{14}{13}&
\instbitrange{12}{10}&
\instbitrange{9}{7}&
\instbitrange{6}{0} \\
\hline
\multicolumn{1}{|c|}{{\scriptsize qx[2:0]}} &
\multicolumn{1}{c|}{{\scriptsize qy[2:0]}} &
\multicolumn{1}{c|}{0X0100X} &
\multicolumn{1}{c|}{{\scriptsize rs1[4]}} &
\multicolumn{1}{c|}{{\scriptsize rs1[2:0]}} &
\multicolumn{1}{c|}{X0} &
\multicolumn{1}{c|}{{\scriptsize qu[2:0]}} &
\multicolumn{1}{c|}{{\scriptsize qz[2:0]}} &
\multicolumn{1}{c|}{1011111}\\
\hline
3& 3& 7& 1& 3& 2& 3& 3& 7 \\
\end{tabular}
 \end{center}
 }
 \vspace{-0.1in}
 \caption[]{ESP.VMAX.U16.LD.INCP} \end{figure*}


\newpage
\subsubsection{ESP.VMAX.U16.ST.INCP}\label{instruction:ESPVMAXU16STINCP}
\textbf{Syntax}\\
ESP.VMAX.U16.ST.INCP qu,rs1,qz,qx,qy

\textbf{Description}\\
This instruction compares the numerical values of the eight 16-bit unsigned vector data segments in the registers \lstinline{qx} and \lstinline{qy}. The data segment with the larger value is written into the corresponding 16-bit data segment in the register \lstinline{qz}.\\During the operation, the value in the register \lstinline{qu} is stored in memory. After the access, the value in the register \lstinline{rs1} is incremented by 16.\\
\textbf{Operation}
\begin{lstlisting}[escapeinside=`?]
  qz[ 15:  0] = (qx[ 15:  0]>=qy[ 15:  0]) ? qx[ 15:  0] : qy[ 15:  0]
  qz[ 31: 16] = (qx[ 31: 16]>=qy[ 31: 16]) ? qx[ 31: 16] : qy[ 31: 16]
  ...
  qz[127:112] = (qx[127:112]>=qy[127:112]) ? qx[127:112] : qy[127:112]

  qu[127:0] => store128(rs1[31:0])
  rs1 [31:0] = rs1[31:0] + 16
\end{lstlisting}


\textbf{Schedule}\\
\begin{longtable}[c]{|l|l|l|}
    \hline
    \rowcolor{lightgray}
    \textbf{Operands} & \textbf{Use} & \textbf{Def} \\\hline
    \endfirsthead
    \hline
    \rowcolor{lightgray}
     \textbf{Operands} & \textbf{Use} & \textbf{Def} \\\hline
    \endhead
    \hline
    \endfoot

    qx & 1 & - \\\hline
  qy & 1 & - \\\hline
     qz & - & 1 \\\hline
     qu & 2 & - \\\hline
     rs1 & 1 & 1

\end{longtable}

\textbf{Format}\\

\begin{figure*}[!htbp]
 {\setlength{\tabcolsep}{1pt}
 \begin{center}
\begin{tabular}{@{}F@{}F@{}O@{}I@{}F@{}W@{}F@{}F@{}O}
\instbitrange{31}{29}&
\instbitrange{28}{26}&
\instbitrange{25}{19}&
\instbit{18}&
\instbitrange{17}{15}&
\instbitrange{14}{13}&
\instbitrange{12}{10}&
\instbitrange{9}{7}&
\instbitrange{6}{0} \\
\hline
\multicolumn{1}{|c|}{{\scriptsize qx[2:0]}} &
\multicolumn{1}{c|}{{\scriptsize qy[2:0]}} &
\multicolumn{1}{c|}{0X1100X} &
\multicolumn{1}{c|}{{\scriptsize rs1[4]}} &
\multicolumn{1}{c|}{{\scriptsize rs1[2:0]}} &
\multicolumn{1}{c|}{X0} &
\multicolumn{1}{c|}{{\scriptsize qu[2:0]}} &
\multicolumn{1}{c|}{{\scriptsize qz[2:0]}} &
\multicolumn{1}{c|}{1011111}\\
\hline
3& 3& 7& 1& 3& 2& 3& 3& 7 \\
\end{tabular}
 \end{center}
 }
 \vspace{-0.1in}
 \caption[]{ESP.VMAX.U16.ST.INCP} \end{figure*}


\newpage
\subsubsection{ESP.VMAX.U32}\label{instruction:ESPVMAXU32}
\textbf{Syntax}\\
ESP.VMAX.U32 qz,qx,qy

\textbf{Description}\\
This instruction compares the numerical values of the four 32-bit unsigned vector data segments in the registers \lstinline{qx} and \lstinline{qy}. The data segment with the larger value is written into the corresponding 32-bit data segment in the register \lstinline{qz}.\\
\textbf{Operation}
\begin{lstlisting}[escapeinside=`?]
  qz[ 31:  0] = (qx[ 31:  0]>=qy[ 31:  0]) ? qx[ 31:  0] : qy[ 31:  0]
  qz[ 63: 32] = (qx[ 63: 32]>=qy[ 63: 32]) ? qx[ 63: 32] : qy[ 63: 32]
  ...
  qz[127: 96] = (qx[127: 96]>=qy[127: 96]) ? qx[127: 96] : qy[127: 96]
\end{lstlisting}


\textbf{Schedule}\\
\begin{longtable}[c]{|l|l|l|}
    \hline
    \rowcolor{lightgray}
    \textbf{Operands} & \textbf{Use} & \textbf{Def} \\\hline
    \endfirsthead
    \hline
    \rowcolor{lightgray}
     \textbf{Operands} & \textbf{Use} & \textbf{Def} \\\hline
    \endhead
    \hline
    \endfoot

    qx & 1 & - \\\hline
  qy & 1 & - \\\hline
     qz & - & 1

\end{longtable}

\textbf{Format}\\

\begin{figure*}[!htbp]
 {\setlength{\tabcolsep}{1pt}
 \begin{center}
\begin{tabular}{@{}F@{}F@{}K@{}F@{}O}
\instbitrange{31}{29}&
\instbitrange{28}{26}&
\instbitrange{25}{10}&
\instbitrange{9}{7}&
\instbitrange{6}{0} \\
\hline
\multicolumn{1}{|c|}{{\scriptsize qx[2:0]}} &
\multicolumn{1}{c|}{{\scriptsize qy[2:0]}} &
\multicolumn{1}{c|}{110100X0X100011X} &
\multicolumn{1}{c|}{{\scriptsize qz[2:0]}} &
\multicolumn{1}{c|}{0011111}\\
\hline
3& 3& 16& 3& 7 \\
\end{tabular}
 \end{center}
 }
 \vspace{-0.1in}
 \caption[]{ESP.VMAX.U32} \end{figure*}


\newpage
\subsubsection{ESP.VMAX.U32.LD.INCP}\label{instruction:ESPVMAXU32LDINCP}
\textbf{Syntax}\\
ESP.VMAX.U32.LD.INCP qu,rs1,qz,qx,qy

\textbf{Description}\\
This instruction compares the numerical values of the four 32-bit unsigned vector data segments in registers \lstinline{qx} and \lstinline{qy}. The data segment with the larger value is written into the corresponding 32-bit data segment in register \lstinline{qz}.\\During the operation, 16 bytes are loaded from the memory into register \lstinline{qu}. After the access, the value in register \lstinline{rs1} is incremented by 16.\\
\textbf{Operation}
\begin{lstlisting}[escapeinside=`?]
  qz[ 31:  0] = (qx[ 31:  0]>=qy[ 31:  0]) ? qx[ 31:  0] : qy[ 31:  0]
  qz[ 63: 32] = (qx[ 63: 32]>=qy[ 63: 32]) ? qx[ 63: 32] : qy[ 63: 32]
  ...
  qz[127: 96] = (qx[127: 96]>=qy[127: 96]) ? qx[127: 96] : qy[127: 96]

  qu[127:0] = load128(rs1[31:0])
  rs1[31:0] = rs1[31:0] + 16
\end{lstlisting}


\textbf{Schedule}\\
\begin{longtable}[c]{|l|l|l|}
    \hline
    \rowcolor{lightgray}
    \textbf{Operands} & \textbf{Use} & \textbf{Def} \\\hline
    \endfirsthead
    \hline
    \rowcolor{lightgray}
     \textbf{Operands} & \textbf{Use} & \textbf{Def} \\\hline
    \endhead
    \hline
    \endfoot

    qx & 1 & - \\\hline
  qy & 1 & - \\\hline
     qz & - & 1 \\\hline
     qu & - & 2 \\\hline
     rs1 & 1 & 1

\end{longtable}

\textbf{Format}\\

\begin{figure*}[!htbp]
 {\setlength{\tabcolsep}{1pt}
 \begin{center}
\begin{tabular}{@{}F@{}F@{}O@{}I@{}F@{}W@{}F@{}F@{}O}
\instbitrange{31}{29}&
\instbitrange{28}{26}&
\instbitrange{25}{19}&
\instbit{18}&
\instbitrange{17}{15}&
\instbitrange{14}{13}&
\instbitrange{12}{10}&
\instbitrange{9}{7}&
\instbitrange{6}{0} \\
\hline
\multicolumn{1}{|c|}{{\scriptsize qx[2:0]}} &
\multicolumn{1}{c|}{{\scriptsize qy[2:0]}} &
\multicolumn{1}{c|}{0X00X1X} &
\multicolumn{1}{c|}{{\scriptsize rs1[4]}} &
\multicolumn{1}{c|}{{\scriptsize rs1[2:0]}} &
\multicolumn{1}{c|}{X0} &
\multicolumn{1}{c|}{{\scriptsize qu[2:0]}} &
\multicolumn{1}{c|}{{\scriptsize qz[2:0]}} &
\multicolumn{1}{c|}{1011111}\\
\hline
3& 3& 7& 1& 3& 2& 3& 3& 7 \\
\end{tabular}
 \end{center}
 }
 \vspace{-0.1in}
 \caption[]{ESP.VMAX.U32.LD.INCP} \end{figure*}


\newpage
\subsubsection{ESP.VMAX.U32.ST.INCP}\label{instruction:ESPVMAXU32STINCP}
\textbf{Syntax}\\
ESP.VMAX.U32.ST.INCP qu,rs1,qz,qx,qy

\textbf{Description}\\
This instruction compares the numerical values of the four 32-bit unsigned vector data segments in the registers \lstinline{qx} and \lstinline{qy}. The data segment with the larger value is written into the corresponding 32-bit data segment in the register \lstinline{qz}.\\During the operation, the value in the register \lstinline{qu} is stored in memory. After the access, the value in the register \lstinline{rs1} is incremented by 16.\\
\textbf{Operation}
\begin{lstlisting}[escapeinside=`?]
  qz[ 31:  0] = (qx[ 31:  0]>=qy[ 31:  0]) ? qx[ 31:  0] : qy[ 31:  0]
  qz[ 63: 32] = (qx[ 63: 32]>=qy[ 63: 32]) ? qx[ 63: 32] : qy[ 63: 32]
  ...
  qz[127: 96] = (qx[127: 96]>=qy[127: 96]) ? qx[127: 96] : qy[127: 96]

  qu[127:0] => store128(rs1[31:0])
  rs1 [31:0] = rs1[31:0] + 16
\end{lstlisting}


\textbf{Schedule}\\
\begin{longtable}[c]{|l|l|l|}
    \hline
    \rowcolor{lightgray}
    \textbf{Operands} & \textbf{Use} & \textbf{Def} \\\hline
    \endfirsthead
    \hline
    \rowcolor{lightgray}
     \textbf{Operands} & \textbf{Use} & \textbf{Def} \\\hline
    \endhead
    \hline
    \endfoot

    qx & 1 & - \\\hline
  qy & 1 & - \\\hline
     qz & - & 1 \\\hline
     qu & 2 & - \\\hline
     rs1 & 1 & 1

\end{longtable}

\textbf{Format}\\

\begin{figure*}[!htbp]
 {\setlength{\tabcolsep}{1pt}
 \begin{center}
\begin{tabular}{@{}F@{}F@{}O@{}I@{}F@{}W@{}F@{}F@{}O}
\instbitrange{31}{29}&
\instbitrange{28}{26}&
\instbitrange{25}{19}&
\instbit{18}&
\instbitrange{17}{15}&
\instbitrange{14}{13}&
\instbitrange{12}{10}&
\instbitrange{9}{7}&
\instbitrange{6}{0} \\
\hline
\multicolumn{1}{|c|}{{\scriptsize qx[2:0]}} &
\multicolumn{1}{c|}{{\scriptsize qy[2:0]}} &
\multicolumn{1}{c|}{0X10X1X} &
\multicolumn{1}{c|}{{\scriptsize rs1[4]}} &
\multicolumn{1}{c|}{{\scriptsize rs1[2:0]}} &
\multicolumn{1}{c|}{X0} &
\multicolumn{1}{c|}{{\scriptsize qu[2:0]}} &
\multicolumn{1}{c|}{{\scriptsize qz[2:0]}} &
\multicolumn{1}{c|}{1011111}\\
\hline
3& 3& 7& 1& 3& 2& 3& 3& 7 \\
\end{tabular}
 \end{center}
 }
 \vspace{-0.1in}
 \caption[]{ESP.VMAX.U32.ST.INCP} \end{figure*}


\newpage
\subsubsection{ESP.VMAX.U8}\label{instruction:ESPVMAXU8}
\textbf{Syntax}\\
ESP.VMAX.U8 qz,qx,qy

\textbf{Description}\\
This instruction compares the numerical values of the 16 8-bit unsigned vector data segments in the registers \lstinline{qx} and \lstinline{qy}. The data segment with the larger value is written into the corresponding 8-bit data segment in the register \lstinline{qz}.\\
\textbf{Operation}
\begin{lstlisting}[escapeinside=`?]
  qz[  7:  0] = (qx[  7:  0]>=qy[  7:  0]) ? qx[  7:  0] : qy[  7:  0]
  qz[ 15:  8] = (qx[ 15:  8]>=qy[ 15:  8]) ? qx[ 15:  8] : qy[ 15:  8]
  ...
  qz[127:120] = (qx[127:120]>=qy[127:120]) ? qx[127:120] : qy[127:120]
\end{lstlisting}


\textbf{Schedule}\\
\begin{longtable}[c]{|l|l|l|}
    \hline
    \rowcolor{lightgray}
    \textbf{Operands} & \textbf{Use} & \textbf{Def} \\\hline
    \endfirsthead
    \hline
    \rowcolor{lightgray}
     \textbf{Operands} & \textbf{Use} & \textbf{Def} \\\hline
    \endhead
    \hline
    \endfoot

    qx & 1 & - \\\hline
  qy & 1 & - \\\hline
     qz & - & 1

\end{longtable}

\textbf{Format}\\

\begin{figure*}[!htbp]
 {\setlength{\tabcolsep}{1pt}
 \begin{center}
\begin{tabular}{@{}F@{}F@{}K@{}F@{}O}
\instbitrange{31}{29}&
\instbitrange{28}{26}&
\instbitrange{25}{10}&
\instbitrange{9}{7}&
\instbitrange{6}{0} \\
\hline
\multicolumn{1}{|c|}{{\scriptsize qx[2:0]}} &
\multicolumn{1}{c|}{{\scriptsize qy[2:0]}} &
\multicolumn{1}{c|}{110100X00000011X} &
\multicolumn{1}{c|}{{\scriptsize qz[2:0]}} &
\multicolumn{1}{c|}{0011111}\\
\hline
3& 3& 16& 3& 7 \\
\end{tabular}
 \end{center}
 }
 \vspace{-0.1in}
 \caption[]{ESP.VMAX.U8} \end{figure*}


\newpage
\subsubsection{ESP.VMAX.U8.LD.INCP}\label{instruction:ESPVMAXU8LDINCP}
\textbf{Syntax}\\
ESP.VMAX.U8.LD.INCP qu,rs1,qz,qx,qy

\textbf{Description}\\
This instruction compares the numerical values of the 16 8-bit unsigned vector data segments in registers \lstinline{qx} and \lstinline{qy}. The data segment with the larger value is written into the corresponding 8-bit data segment in register \lstinline{qz}.\\During the operation, 16 bytes are loaded from the memory into register \lstinline{qu}. After the access, the value in register \lstinline{rs1} is incremented by 16.\\
\textbf{Operation}
\begin{lstlisting}[escapeinside=`?]
  qz[  7:  0] = (qx[  7:  0]>=qy[  7:  0]) ? qx[  7:  0] : qy[  7:  0]
  qz[ 15:  8] = (qx[ 15:  8]>=qy[ 15:  8]) ? qx[ 15:  8] : qy[ 15:  8]
  ...
  qz[127:120] = (qx[127:120]>=qy[127:120]) ? qx[127:120] : qy[127:120]

 qu[127:0] = load128(rs1[31:0])
  rs1[31:0] = rs1[31:0] + 16
\end{lstlisting}


\textbf{Schedule}\\
\begin{longtable}[c]{|l|l|l|}
    \hline
    \rowcolor{lightgray}
    \textbf{Operands} & \textbf{Use} & \textbf{Def} \\\hline
    \endfirsthead
    \hline
    \rowcolor{lightgray}
     \textbf{Operands} & \textbf{Use} & \textbf{Def} \\\hline
    \endhead
    \hline
    \endfoot

    qx & 1 & - \\\hline
  qy & 1 & - \\\hline
     qz & - & 1 \\\hline
     qu & - & 2 \\\hline
     rs1 & 1 & 1

\end{longtable}

\textbf{Format}\\

\begin{figure*}[!htbp]
 {\setlength{\tabcolsep}{1pt}
 \begin{center}
\begin{tabular}{@{}F@{}F@{}O@{}I@{}F@{}W@{}F@{}F@{}O}
\instbitrange{31}{29}&
\instbitrange{28}{26}&
\instbitrange{25}{19}&
\instbit{18}&
\instbitrange{17}{15}&
\instbitrange{14}{13}&
\instbitrange{12}{10}&
\instbitrange{9}{7}&
\instbitrange{6}{0} \\
\hline
\multicolumn{1}{|c|}{{\scriptsize qx[2:0]}} &
\multicolumn{1}{c|}{{\scriptsize qy[2:0]}} &
\multicolumn{1}{c|}{0X0000X} &
\multicolumn{1}{c|}{{\scriptsize rs1[4]}} &
\multicolumn{1}{c|}{{\scriptsize rs1[2:0]}} &
\multicolumn{1}{c|}{X0} &
\multicolumn{1}{c|}{{\scriptsize qu[2:0]}} &
\multicolumn{1}{c|}{{\scriptsize qz[2:0]}} &
\multicolumn{1}{c|}{1011111}\\
\hline
3& 3& 7& 1& 3& 2& 3& 3& 7 \\
\end{tabular}
 \end{center}
 }
 \vspace{-0.1in}
 \caption[]{ESP.VMAX.U8.LD.INCP} \end{figure*}


\newpage
\subsubsection{ESP.VMAX.U8.ST.INCP}\label{instruction:ESPVMAXU8STINCP}
\textbf{Syntax}\\
ESP.VMAX.U8.ST.INCP qu,rs1,qz,qx,qy

\textbf{Description}\\
This instruction compares the numerical values of the 16 8-bit unsigned vector data segments in registers \lstinline{qx} and \lstinline{qy}. The data segment with the larger value is written into the corresponding 8-bit data segment in register \lstinline{qz}.\\During the operation, the value in register \lstinline{qu} is stored in memory. After the access, the value in register \lstinline{rs1} is incremented by 16.\\
\textbf{Operation}
\begin{lstlisting}[escapeinside=`?]
  qz[  7:  0] = (qx[  7:  0]>=qy[  7:  0]) ? qx[  7:  0] : qy[  7:  0]
  qz[ 15:  8] = (qx[ 15:  8]>=qy[ 15:  8]) ? qx[ 15:  8] : qy[ 15:  8]
  ...
  qz[127:120] = (qx[127:120]>=qy[127:120]) ? qx[127:120] : qy[127:120]

    qu[127:0] => store128(rs1[31:0])
  rs1 [31:0] = rs1[31:0] + 16
\end{lstlisting}


\textbf{Schedule}\\
\begin{longtable}[c]{|l|l|l|}
    \hline
    \rowcolor{lightgray}
    \textbf{Operands} & \textbf{Use} & \textbf{Def} \\\hline
    \endfirsthead
    \hline
    \rowcolor{lightgray}
     \textbf{Operands} & \textbf{Use} & \textbf{Def} \\\hline
    \endhead
    \hline
    \endfoot

    qx & 1 & - \\\hline
  qy & 1 & - \\\hline
     qz & - & 1 \\\hline
     qu & 2 & - \\\hline
     rs1 & 1 & 1

\end{longtable}

\textbf{Format}\\

\begin{figure*}[!htbp]
 {\setlength{\tabcolsep}{1pt}
 \begin{center}
\begin{tabular}{@{}F@{}F@{}O@{}I@{}F@{}W@{}F@{}F@{}O}
\instbitrange{31}{29}&
\instbitrange{28}{26}&
\instbitrange{25}{19}&
\instbit{18}&
\instbitrange{17}{15}&
\instbitrange{14}{13}&
\instbitrange{12}{10}&
\instbitrange{9}{7}&
\instbitrange{6}{0} \\
\hline
\multicolumn{1}{|c|}{{\scriptsize qx[2:0]}} &
\multicolumn{1}{c|}{{\scriptsize qy[2:0]}} &
\multicolumn{1}{c|}{0X1000X} &
\multicolumn{1}{c|}{{\scriptsize rs1[4]}} &
\multicolumn{1}{c|}{{\scriptsize rs1[2:0]}} &
\multicolumn{1}{c|}{X0} &
\multicolumn{1}{c|}{{\scriptsize qu[2:0]}} &
\multicolumn{1}{c|}{{\scriptsize qz[2:0]}} &
\multicolumn{1}{c|}{1011111}\\
\hline
3& 3& 7& 1& 3& 2& 3& 3& 7 \\
\end{tabular}
 \end{center}
 }
 \vspace{-0.1in}
 \caption[]{ESP.VMAX.U8.ST.INCP} \end{figure*}


\newpage
\subsubsection{ESP.VMIN.S16}\label{instruction:ESPVMINS16}
\textbf{Syntax}\\
ESP.VMIN.S16 qz,qx,qy

\textbf{Description}\\
This instruction compares the numerical values of the eight 16-bit signed vector data segments in the registers \lstinline{qx} and \lstinline{qy}. The data segment with the smaller value is written into the corresponding 16-bit data segment in the register \lstinline{qz}.\\
\textbf{Operation}
\begin{lstlisting}[escapeinside=`?]
  qz[ 15:  0] = (qx[ 15:  0]<=qy[ 15:  0]) ? qx[ 15:  0] : qy[ 15:  0]
  qz[ 31: 16] = (qx[ 31: 16]<=qy[ 31: 16]) ? qx[ 31: 16] : qy[ 31: 16]
  ...
  qz[127:112] = (qx[127:112]<=qy[127:112]) ? qx[127:112] : qy[127:112]
\end{lstlisting}


\textbf{Schedule}\\
\begin{longtable}[c]{|l|l|l|}
    \hline
    \rowcolor{lightgray}
    \textbf{Operands} & \textbf{Use} & \textbf{Def} \\\hline
    \endfirsthead
    \hline
    \rowcolor{lightgray}
     \textbf{Operands} & \textbf{Use} & \textbf{Def} \\\hline
    \endhead
    \hline
    \endfoot

    qx & 1 & - \\\hline
  qy & 1 & - \\\hline
     qz & - & 1

\end{longtable}

\textbf{Format}\\

\begin{figure*}[!htbp]
 {\setlength{\tabcolsep}{1pt}
 \begin{center}
\begin{tabular}{@{}F@{}F@{}K@{}F@{}O}
\instbitrange{31}{29}&
\instbitrange{28}{26}&
\instbitrange{25}{10}&
\instbitrange{9}{7}&
\instbitrange{6}{0} \\
\hline
\multicolumn{1}{|c|}{{\scriptsize qx[2:0]}} &
\multicolumn{1}{c|}{{\scriptsize qy[2:0]}} &
\multicolumn{1}{c|}{110100X11010011X} &
\multicolumn{1}{c|}{{\scriptsize qz[2:0]}} &
\multicolumn{1}{c|}{0011111}\\
\hline
3& 3& 16& 3& 7 \\
\end{tabular}
 \end{center}
 }
 \vspace{-0.1in}
 \caption[]{ESP.VMIN.S16} \end{figure*}


\newpage
\subsubsection{ESP.VMIN.S16.LD.INCP}\label{instruction:ESPVMINS16LDINCP}
\textbf{Syntax}\\
ESP.VMIN.S16.LD.INCP qu,rs1,qz,qx,qy

\textbf{Description}\\
This instruction compares the numerical values of the eight 16-bit signed vector data segments in the registers \lstinline{qx} and \lstinline{qy}. The data segment with the smaller value is written into the corresponding 16-bit data segment in the register \lstinline{qz}.\\During the operation, 16 bytes are loaded from the memory into the register \lstinline{qu}. After the access, the value in the register \lstinline{rs1} is incremented by 16.\\
\textbf{Operation}
\begin{lstlisting}[escapeinside=`?]
  qz[ 15:  0] = (qx[ 15:  0]<=qy[ 15:  0]) ? qx[ 15:  0] : qy[ 15:  0]
  qz[ 31: 16] = (qx[ 31: 16]<=qy[ 31: 16]) ? qx[ 31: 16] : qy[ 31: 16]
  ...
  qz[127:112] = (qx[127:112]<=qy[127:112]) ? qx[127:112] : qy[127:112]

  qu[127:0] = load128(rs1[31:0])
  rs1[31:0] = rs1[31:0] + 16
\end{lstlisting}


\textbf{Schedule}\\
\begin{longtable}[c]{|l|l|l|}
    \hline
    \rowcolor{lightgray}
    \textbf{Operands} & \textbf{Use} & \textbf{Def} \\\hline
    \endfirsthead
    \hline
    \rowcolor{lightgray}
     \textbf{Operands} & \textbf{Use} & \textbf{Def} \\\hline
    \endhead
    \hline
    \endfoot

    qx & 1 & - \\\hline
  qy & 1 & - \\\hline
     qz & - & 1 \\\hline
     qu & - & 2 \\\hline
     rs1 & 1 & 1

\end{longtable}

\textbf{Format}\\

\begin{figure*}[!htbp]
 {\setlength{\tabcolsep}{1pt}
 \begin{center}
\begin{tabular}{@{}F@{}F@{}O@{}I@{}F@{}W@{}F@{}F@{}O}
\instbitrange{31}{29}&
\instbitrange{28}{26}&
\instbitrange{25}{19}&
\instbit{18}&
\instbitrange{17}{15}&
\instbitrange{14}{13}&
\instbitrange{12}{10}&
\instbitrange{9}{7}&
\instbitrange{6}{0} \\
\hline
\multicolumn{1}{|c|}{{\scriptsize qx[2:0]}} &
\multicolumn{1}{c|}{{\scriptsize qy[2:0]}} &
\multicolumn{1}{c|}{1X0110X} &
\multicolumn{1}{c|}{{\scriptsize rs1[4]}} &
\multicolumn{1}{c|}{{\scriptsize rs1[2:0]}} &
\multicolumn{1}{c|}{X0} &
\multicolumn{1}{c|}{{\scriptsize qu[2:0]}} &
\multicolumn{1}{c|}{{\scriptsize qz[2:0]}} &
\multicolumn{1}{c|}{1011111}\\
\hline
3& 3& 7& 1& 3& 2& 3& 3& 7 \\
\end{tabular}
 \end{center}
 }
 \vspace{-0.1in}
 \caption[]{ESP.VMIN.S16.LD.INCP} \end{figure*}


\newpage
\subsubsection{ESP.VMIN.S16.ST.INCP}\label{instruction:ESPVMINS16STINCP}
\textbf{Syntax}\\
ESP.VMIN.S16.ST.INCP qu,rs1,qz,qx,qy

\textbf{Description}\\
This instruction compares the numerical values of the eight 16-bit signed vector data segments in the registers \lstinline{qx} and \lstinline{qy}. The data segment with the smaller value is written into the corresponding 16-bit data segment in the register \lstinline{qz}.\\During the operation, the value in the register \lstinline{qu} is stored in memory. After the access, the value in the register \lstinline{rs1} is incremented by 16.\\
\textbf{Operation}
\begin{lstlisting}[escapeinside=`?]
  qz[ 15:  0] = (qx[ 15:  0]<=qy[ 15:  0]) ? qx[ 15:  0] : qy[ 15:  0]
  qz[ 31: 16] = (qx[ 31: 16]<=qy[ 31: 16]) ? qx[ 31: 16] : qy[ 31: 16]
  ...
  qz[127:112] = (qx[127:112]<=qy[127:112]) ? qx[127:112] : qy[127:112]

  qu[127:0] => store128(rs1[31:0])
  rs1[31:0] = rs1[31:0] + 16
\end{lstlisting}


\textbf{Schedule}\\
\begin{longtable}[c]{|l|l|l|}
    \hline
    \rowcolor{lightgray}
    \textbf{Operands} & \textbf{Use} & \textbf{Def} \\\hline
    \endfirsthead
    \hline
    \rowcolor{lightgray}
     \textbf{Operands} & \textbf{Use} & \textbf{Def} \\\hline
    \endhead
    \hline
    \endfoot

    qx & 1 & - \\\hline
  qy & 1 & - \\\hline
     qz & - & 1 \\\hline
     qu & 2 & - \\\hline
     rs1 & 1 & 1

\end{longtable}

\textbf{Format}\\

\begin{figure*}[!htbp]
 {\setlength{\tabcolsep}{1pt}
 \begin{center}
\begin{tabular}{@{}F@{}F@{}O@{}I@{}F@{}W@{}F@{}F@{}O}
\instbitrange{31}{29}&
\instbitrange{28}{26}&
\instbitrange{25}{19}&
\instbit{18}&
\instbitrange{17}{15}&
\instbitrange{14}{13}&
\instbitrange{12}{10}&
\instbitrange{9}{7}&
\instbitrange{6}{0} \\
\hline
\multicolumn{1}{|c|}{{\scriptsize qx[2:0]}} &
\multicolumn{1}{c|}{{\scriptsize qy[2:0]}} &
\multicolumn{1}{c|}{1X1110X} &
\multicolumn{1}{c|}{{\scriptsize rs1[4]}} &
\multicolumn{1}{c|}{{\scriptsize rs1[2:0]}} &
\multicolumn{1}{c|}{X0} &
\multicolumn{1}{c|}{{\scriptsize qu[2:0]}} &
\multicolumn{1}{c|}{{\scriptsize qz[2:0]}} &
\multicolumn{1}{c|}{1011111}\\
\hline
3& 3& 7& 1& 3& 2& 3& 3& 7 \\
\end{tabular}
 \end{center}
 }
 \vspace{-0.1in}
 \caption[]{ESP.VMIN.S16.ST.INCP} \end{figure*}


\newpage
\subsubsection{ESP.VMIN.S32}\label{instruction:ESPVMINS32}
\textbf{Syntax}\\
ESP.VMIN.S32 qz,qx,qy

\textbf{Description}\\
This instruction compares the numerical values of the four 32-bit signed vector data segments in the registers \lstinline{qx} and \lstinline{qy}. The data segment with the smaller value is written into the corresponding 32-bit data segment in the register \lstinline{qz}.\\
\textbf{Operation}
\begin{lstlisting}[escapeinside=`?]
  qz[ 31:  0] = (qx[ 31:  0]<=qy[ 31:  0]) ? qx[ 31:  0] : qy[ 31:  0]
  qz[ 63: 32] = (qx[ 63: 32]<=qy[ 63: 32]) ? qx[ 63: 32] : qy[ 63: 32]
  ...
  qz[127: 96] = (qx[127: 96]<=qy[127: 96]) ? qx[127: 96] : qy[127: 96]
\end{lstlisting}


\textbf{Schedule}\\
\begin{longtable}[c]{|l|l|l|}
    \hline
    \rowcolor{lightgray}
    \textbf{Operands} & \textbf{Use} & \textbf{Def} \\\hline
    \endfirsthead
    \hline
    \rowcolor{lightgray}
     \textbf{Operands} & \textbf{Use} & \textbf{Def} \\\hline
    \endhead
    \hline
    \endfoot

    qx & 1 & - \\\hline
  qy & 1 & - \\\hline
     qz & - & 1

\end{longtable}

\textbf{Format}\\

\begin{figure*}[!htbp]
 {\setlength{\tabcolsep}{1pt}
 \begin{center}
\begin{tabular}{@{}F@{}F@{}K@{}F@{}O}
\instbitrange{31}{29}&
\instbitrange{28}{26}&
\instbitrange{25}{10}&
\instbitrange{9}{7}&
\instbitrange{6}{0} \\
\hline
\multicolumn{1}{|c|}{{\scriptsize qx[2:0]}} &
\multicolumn{1}{c|}{{\scriptsize qy[2:0]}} &
\multicolumn{1}{c|}{110100X1X110011X} &
\multicolumn{1}{c|}{{\scriptsize qz[2:0]}} &
\multicolumn{1}{c|}{0011111}\\
\hline
3& 3& 16& 3& 7 \\
\end{tabular}
 \end{center}
 }
 \vspace{-0.1in}
 \caption[]{ESP.VMIN.S32} \end{figure*}


\newpage
\subsubsection{ESP.VMIN.S32.LD.INCP}\label{instruction:ESPVMINS32LDINCP}
\textbf{Syntax}\\
ESP.VMIN.S32.LD.INCP qu,rs1,qz,qx,qy

\textbf{Description}\\
This instruction compares the numerical values of the four 32-bit signed vector data segments in registers \lstinline{qx} and \lstinline{qy}. The data segment with the smaller value is written into the corresponding 32-bit data segment in register \lstinline{qz}.\\During the operation, 16 bytes are loaded from the memory into register \lstinline{qu}. After the access, the value in register \lstinline{rs1} is incremented by 16.\\
\textbf{Operation}
\begin{lstlisting}[escapeinside=`?]
  qz[ 31:  0] = (qx[ 31:  0]<=qy[ 31:  0]) ? qx[ 31:  0] : qy[ 31:  0]
  qz[ 63: 32] = (qx[ 63: 32]<=qy[ 63: 32]) ? qx[ 63: 32] : qy[ 63: 32]
  ...
  qz[127: 96] = (qx[127: 96]<=qy[127: 96]) ? qx[127: 96] : qy[127: 96]

  qu[127:0] = load128(rs1[31:0])
  rs1[31:0] = rs1[31:0] + 16
\end{lstlisting}


\textbf{Schedule}\\
\begin{longtable}[c]{|l|l|l|}
    \hline
    \rowcolor{lightgray}
    \textbf{Operands} & \textbf{Use} & \textbf{Def} \\\hline
    \endfirsthead
    \hline
    \rowcolor{lightgray}
     \textbf{Operands} & \textbf{Use} & \textbf{Def} \\\hline
    \endhead
    \hline
    \endfoot

    qx & 1 & - \\\hline
  qy & 1 & - \\\hline
     qz & - & 1 \\\hline
     qu & - & 2 \\\hline
     rs1 & 1 & 1

\end{longtable}

\textbf{Format}\\

\begin{figure*}[!htbp]
 {\setlength{\tabcolsep}{1pt}
 \begin{center}
\begin{tabular}{@{}F@{}F@{}O@{}I@{}F@{}W@{}F@{}F@{}O}
\instbitrange{31}{29}&
\instbitrange{28}{26}&
\instbitrange{25}{19}&
\instbit{18}&
\instbitrange{17}{15}&
\instbitrange{14}{13}&
\instbitrange{12}{10}&
\instbitrange{9}{7}&
\instbitrange{6}{0} \\
\hline
\multicolumn{1}{|c|}{{\scriptsize qx[2:0]}} &
\multicolumn{1}{c|}{{\scriptsize qy[2:0]}} &
\multicolumn{1}{c|}{1X01X1X} &
\multicolumn{1}{c|}{{\scriptsize rs1[4]}} &
\multicolumn{1}{c|}{{\scriptsize rs1[2:0]}} &
\multicolumn{1}{c|}{X0} &
\multicolumn{1}{c|}{{\scriptsize qu[2:0]}} &
\multicolumn{1}{c|}{{\scriptsize qz[2:0]}} &
\multicolumn{1}{c|}{1011111}\\
\hline
3& 3& 7& 1& 3& 2& 3& 3& 7 \\
\end{tabular}
 \end{center}
 }
 \vspace{-0.1in}
 \caption[]{ESP.VMIN.S32.LD.INCP} \end{figure*}


\newpage
\subsubsection{ESP.VMIN.S32.ST.INCP}\label{instruction:ESPVMINS32STINCP}
\textbf{Syntax}\\
ESP.VMIN.S32.ST.INCP qu,rs1,qz,qx,qy

\textbf{Description}\\
This instruction compares the numerical values of the four 32-bit signed vector data segments in the registers \lstinline{qx} and \lstinline{qy}. The data segment with the smaller value is written into the corresponding 32-bit data segment in the register \lstinline{qz}.\\During the operation, the value in the register \lstinline{qu} is stored in memory. After the access, the value in the register \lstinline{rs1} is incremented by 16.\\
\textbf{Operation}
\begin{lstlisting}[escapeinside=`?]
  qz[ 31:  0] = (qx[ 31:  0]<=qy[ 31:  0]) ? qx[ 31:  0] : qy[ 31:  0]
  qz[ 63: 32] = (qx[ 63: 32]<=qy[ 63: 32]) ? qx[ 63: 32] : qy[ 63: 32]
  ...
  qz[127: 96] = (qx[127: 96]<=qy[127: 96]) ? qx[127: 96] : qy[127: 96]

  qu[127:0] = store128(rs1[31:0])
  rs1[31:0] = rs1[31:0] + 16
\end{lstlisting}


\textbf{Schedule}\\
\begin{longtable}[c]{|l|l|l|}
    \hline
    \rowcolor{lightgray}
    \textbf{Operands} & \textbf{Use} & \textbf{Def} \\\hline
    \endfirsthead
    \hline
    \rowcolor{lightgray}
     \textbf{Operands} & \textbf{Use} & \textbf{Def} \\\hline
    \endhead
    \hline
    \endfoot

    qx & 1 & - \\\hline
  qy & 1 & - \\\hline
     qz & - & 1 \\\hline
     qu & 2 & - \\\hline
     rs1 & 1 & 1

\end{longtable}

\textbf{Format}\\

\begin{figure*}[!htbp]
 {\setlength{\tabcolsep}{1pt}
 \begin{center}
\begin{tabular}{@{}F@{}F@{}O@{}I@{}F@{}W@{}F@{}F@{}O}
\instbitrange{31}{29}&
\instbitrange{28}{26}&
\instbitrange{25}{19}&
\instbit{18}&
\instbitrange{17}{15}&
\instbitrange{14}{13}&
\instbitrange{12}{10}&
\instbitrange{9}{7}&
\instbitrange{6}{0} \\
\hline
\multicolumn{1}{|c|}{{\scriptsize qx[2:0]}} &
\multicolumn{1}{c|}{{\scriptsize qy[2:0]}} &
\multicolumn{1}{c|}{1X11X1X} &
\multicolumn{1}{c|}{{\scriptsize rs1[4]}} &
\multicolumn{1}{c|}{{\scriptsize rs1[2:0]}} &
\multicolumn{1}{c|}{X0} &
\multicolumn{1}{c|}{{\scriptsize qu[2:0]}} &
\multicolumn{1}{c|}{{\scriptsize qz[2:0]}} &
\multicolumn{1}{c|}{1011111}\\
\hline
3& 3& 7& 1& 3& 2& 3& 3& 7 \\
\end{tabular}
 \end{center}
 }
 \vspace{-0.1in}
 \caption[]{ESP.VMIN.S32.ST.INCP} \end{figure*}


\newpage
\subsubsection{ESP.VMIN.S8}\label{instruction:ESPVMINS8}
\textbf{Syntax}\\
ESP.VMIN.S8 qz,qx,qy

\textbf{Description}\\
This instruction compares the numerical values of the 16 8-bit signed vector data segments in registers \lstinline{qx} and \lstinline{qy}. The data segment with the smaller value is written into the corresponding 8-bit data segment in the register \lstinline{qz}.\\
\textbf{Operation}
\begin{lstlisting}[escapeinside=`?]
  qz[  7:  0] = (qx[  7:  0]<=qy[  7:  0]) ? qx[  7:  0] : qy[  7:  0]
  qz[ 15:  8] = (qx[ 15:  8]<=qy[ 15:  8]) ? qx[ 15:  8] : qy[ 15:  8]
  ...
  qz[127:120] = (qx[127:120]<=qy[127:120]) ? qx[127:120] : qy[127:120]
\end{lstlisting}


\textbf{Schedule}\\
\begin{longtable}[c]{|l|l|l|}
    \hline
    \rowcolor{lightgray}
    \textbf{Operands} & \textbf{Use} & \textbf{Def} \\\hline
    \endfirsthead
    \hline
    \rowcolor{lightgray}
     \textbf{Operands} & \textbf{Use} & \textbf{Def} \\\hline
    \endhead
    \hline
    \endfoot

    qx & 1 & - \\\hline
  qy & 1 & - \\\hline
     qz & - & 1

\end{longtable}

\textbf{Format}\\

\begin{figure*}[!htbp]
 {\setlength{\tabcolsep}{1pt}
 \begin{center}
\begin{tabular}{@{}F@{}F@{}K@{}F@{}O}
\instbitrange{31}{29}&
\instbitrange{28}{26}&
\instbitrange{25}{10}&
\instbitrange{9}{7}&
\instbitrange{6}{0} \\
\hline
\multicolumn{1}{|c|}{{\scriptsize qx[2:0]}} &
\multicolumn{1}{c|}{{\scriptsize qy[2:0]}} &
\multicolumn{1}{c|}{110100X01010011X} &
\multicolumn{1}{c|}{{\scriptsize qz[2:0]}} &
\multicolumn{1}{c|}{0011111}\\
\hline
3& 3& 16& 3& 7 \\
\end{tabular}
 \end{center}
 }
 \vspace{-0.1in}
 \caption[]{ESP.VMIN.S8} \end{figure*}


\newpage
\subsubsection{ESP.VMIN.S8.LD.INCP}\label{instruction:ESPVMINS8LDINCP}
\textbf{Syntax}\\
ESP.VMIN.S8.LD.INCP qu,rs1,qz,qx,qy

\textbf{Description}\\
This instruction compares the numerical values of the 16 8-bit signed vector data segments in registers \lstinline{qx} and \lstinline{qy}. The data segment with the smaller value is written into the corresponding 8-bit data segment in register \lstinline{qz}.\\During the operation, 16 bytes are loaded from the memory into register \lstinline{qu}. After the access, the value in register \lstinline{rs1} is incremented by 16.\\
\textbf{Operation}
\begin{lstlisting}[escapeinside=`?]
  qz[  7:  0] = (qx[  7:  0]<=qy[  7:  0]) ? qx[  7:  0] : qy[  7:  0]
  qz[ 15:  8] = (qx[ 15:  8]<=qy[ 15:  8]) ? qx[ 15:  8] : qy[ 15:  8]
  ...
  qz[127:120] = (qx[127:120]<=qy[127:120]) ? qx[127:120] : qy[127:120]

  qu[127:0] = load128(rs1[31:0])
  rs1[31:0] = rs1[31:0] + 16
\end{lstlisting}


\textbf{Schedule}\\
\begin{longtable}[c]{|l|l|l|}
    \hline
    \rowcolor{lightgray}
    \textbf{Operands} & \textbf{Use} & \textbf{Def} \\\hline
    \endfirsthead
    \hline
    \rowcolor{lightgray}
     \textbf{Operands} & \textbf{Use} & \textbf{Def} \\\hline
    \endhead
    \hline
    \endfoot

    qx & 1 & - \\\hline
  qy & 1 & - \\\hline
     qz & - & 1 \\\hline
     qu & - & 2 \\\hline
     rs1 & 1 & 1

\end{longtable}

\textbf{Format}\\

\begin{figure*}[!htbp]
 {\setlength{\tabcolsep}{1pt}
 \begin{center}
\begin{tabular}{@{}F@{}F@{}O@{}I@{}F@{}W@{}F@{}F@{}O}
\instbitrange{31}{29}&
\instbitrange{28}{26}&
\instbitrange{25}{19}&
\instbit{18}&
\instbitrange{17}{15}&
\instbitrange{14}{13}&
\instbitrange{12}{10}&
\instbitrange{9}{7}&
\instbitrange{6}{0} \\
\hline
\multicolumn{1}{|c|}{{\scriptsize qx[2:0]}} &
\multicolumn{1}{c|}{{\scriptsize qy[2:0]}} &
\multicolumn{1}{c|}{1X0010X} &
\multicolumn{1}{c|}{{\scriptsize rs1[4]}} &
\multicolumn{1}{c|}{{\scriptsize rs1[2:0]}} &
\multicolumn{1}{c|}{X0} &
\multicolumn{1}{c|}{{\scriptsize qu[2:0]}} &
\multicolumn{1}{c|}{{\scriptsize qz[2:0]}} &
\multicolumn{1}{c|}{1011111}\\
\hline
3& 3& 7& 1& 3& 2& 3& 3& 7 \\
\end{tabular}
 \end{center}
 }
 \vspace{-0.1in}
 \caption[]{ESP.VMIN.S8.LD.INCP} \end{figure*}


\newpage
\subsubsection{ESP.VMIN.S8.ST.INCP}\label{instruction:ESPVMINS8STINCP}
\textbf{Syntax}\\
ESP.VMIN.S8.ST.INCP qu,rs1,qz,qx,qy

\textbf{Description}\\
This instruction compares the numerical values of the 16 8-bit signed vector data segments in registers \lstinline{qx} and \lstinline{qy}. The data segment with the smaller value is written into the corresponding 8-bit data segment in register \lstinline{qz}.\\During the operation, the value in register \lstinline{qu} is stored in memory. After the access, the value in register \lstinline{rs1} is incremented by 16.\\
\textbf{Operation}
\begin{lstlisting}[escapeinside=`?]
  qz[  7:  0] = (qx[  7:  0]<=qy[  7:  0]) ? qx[  7:  0] : qy[  7:  0]
  qz[ 15:  8] = (qx[ 15:  8]<=qy[ 15:  8]) ? qx[ 15:  8] : qy[ 15:  8]
  ...
  qz[127:120] = (qx[127:120]<=qy[127:120]) ? qx[127:120] : qy[127:120]

  qu[127:0] = store128(rs1[31:0])
  rs1[31:0] = rs1[31:0] + 16
\end{lstlisting}


\textbf{Schedule}\\
\begin{longtable}[c]{|l|l|l|}
    \hline
    \rowcolor{lightgray}
    \textbf{Operands} & \textbf{Use} & \textbf{Def} \\\hline
    \endfirsthead
    \hline
    \rowcolor{lightgray}
     \textbf{Operands} & \textbf{Use} & \textbf{Def} \\\hline
    \endhead
    \hline
    \endfoot

    qx & 1 & - \\\hline
  qy & 1 & - \\\hline
     qz & - & 1 \\\hline
     qu & 2 & - \\\hline
     rs1 & 1 & 1

\end{longtable}

\textbf{Format}\\

\begin{figure*}[!htbp]
 {\setlength{\tabcolsep}{1pt}
 \begin{center}
\begin{tabular}{@{}F@{}F@{}O@{}I@{}F@{}W@{}F@{}F@{}O}
\instbitrange{31}{29}&
\instbitrange{28}{26}&
\instbitrange{25}{19}&
\instbit{18}&
\instbitrange{17}{15}&
\instbitrange{14}{13}&
\instbitrange{12}{10}&
\instbitrange{9}{7}&
\instbitrange{6}{0} \\
\hline
\multicolumn{1}{|c|}{{\scriptsize qx[2:0]}} &
\multicolumn{1}{c|}{{\scriptsize qy[2:0]}} &
\multicolumn{1}{c|}{1X1010X} &
\multicolumn{1}{c|}{{\scriptsize rs1[4]}} &
\multicolumn{1}{c|}{{\scriptsize rs1[2:0]}} &
\multicolumn{1}{c|}{X0} &
\multicolumn{1}{c|}{{\scriptsize qu[2:0]}} &
\multicolumn{1}{c|}{{\scriptsize qz[2:0]}} &
\multicolumn{1}{c|}{1011111}\\
\hline
3& 3& 7& 1& 3& 2& 3& 3& 7 \\
\end{tabular}
 \end{center}
 }
 \vspace{-0.1in}
 \caption[]{ESP.VMIN.S8.ST.INCP} \end{figure*}


\newpage
\subsubsection{ESP.VMIN.U16}\label{instruction:ESPVMINU16}
\textbf{Syntax}\\
ESP.VMIN.U16 qz,qx,qy

\textbf{Description}\\
This instruction compares the numerical values of the eight 16-bit unsigned vector data segments in the registers \lstinline{qx} and \lstinline{qy}. The data segment with the smaller value is written into the corresponding 16-bit data segment in the register \lstinline{qz}.\\
\textbf{Operation}
\begin{lstlisting}[escapeinside=`?]
  qz[ 15:  0] = (qx[ 15:  0]<=qy[ 15:  0]) ? qx[ 15:  0] : qy[ 15:  0]
  qz[ 31: 16] = (qx[ 31: 16]<=qy[ 31: 16]) ? qx[ 31: 16] : qy[ 31: 16]
  ...
  qz[127:112] = (qx[127:112]<=qy[127:112]) ? qx[127:112] : qy[127:112]
\end{lstlisting}


\textbf{Schedule}\\
\begin{longtable}[c]{|l|l|l|}
    \hline
    \rowcolor{lightgray}
    \textbf{Operands} & \textbf{Use} & \textbf{Def} \\\hline
    \endfirsthead
    \hline
    \rowcolor{lightgray}
     \textbf{Operands} & \textbf{Use} & \textbf{Def} \\\hline
    \endhead
    \hline
    \endfoot

    qx & 1 & - \\\hline
  qy & 1 & - \\\hline
     qz & - & 1 \\\hline

\end{longtable}

\textbf{Format}\\

\begin{figure*}[!htbp]
 {\setlength{\tabcolsep}{1pt}
 \begin{center}
\begin{tabular}{@{}F@{}F@{}K@{}F@{}O}
\instbitrange{31}{29}&
\instbitrange{28}{26}&
\instbitrange{25}{10}&
\instbitrange{9}{7}&
\instbitrange{6}{0} \\
\hline
\multicolumn{1}{|c|}{{\scriptsize qx[2:0]}} &
\multicolumn{1}{c|}{{\scriptsize qy[2:0]}} &
\multicolumn{1}{c|}{110100X10010011X} &
\multicolumn{1}{c|}{{\scriptsize qz[2:0]}} &
\multicolumn{1}{c|}{0011111}\\
\hline
3& 3& 16& 3& 7 \\
\end{tabular}
 \end{center}
 }
 \vspace{-0.1in}
 \caption[]{ESP.VMIN.U16} \end{figure*}


\newpage
\subsubsection{ESP.VMIN.U16.LD.INCP}\label{instruction:ESPVMINU16LDINCP}
\textbf{Syntax}\\
ESP.VMIN.U16.LD.INCP qu,rs1,qz,qx,qy

\textbf{Description}\\
This instruction compares the numerical values of the eight 16-bit unsigned vector data segments in registers \lstinline{qx} and \lstinline{qy}. The data segment with the smaller value is written into the corresponding 16-bit data segment in register \lstinline{qz}.\\During the operation, 16 bytes are loaded from the memory into register \lstinline{qu}. After the access, the value in register \lstinline{rs1} is incremented by 16.\\
\textbf{Operation}
\begin{lstlisting}[escapeinside=`?]
  qz[ 15:  0] = (qx[ 15:  0]<=qy[ 15:  0]) ? qx[ 15:  0] : qy[ 15:  0]
  qz[ 31: 16] = (qx[ 31: 16]<=qy[ 31: 16]) ? qx[ 31: 16] : qy[ 31: 16]
  ...
  qz[127:112] = (qx[127:112]<=qy[127:112]) ? qx[127:112] : qy[127:112]

  qu[127:0] = load128(rs1[31:0])
  rs1[31:0] = rs1[31:0] + 16
\end{lstlisting}


\textbf{Schedule}\\
\begin{longtable}[c]{|l|l|l|}
    \hline
    \rowcolor{lightgray}
    \textbf{Operands} & \textbf{Use} & \textbf{Def} \\\hline
    \endfirsthead
    \hline
    \rowcolor{lightgray}
     \textbf{Operands} & \textbf{Use} & \textbf{Def} \\\hline
    \endhead
    \hline
    \endfoot

    qx & 1 & - \\\hline
  qy & 1 & - \\\hline
     qz & - & 1 \\\hline
     qu & - & 2 \\\hline
     rs1 & 1 & 1

\end{longtable}

\textbf{Format}\\

\begin{figure*}[!htbp]
 {\setlength{\tabcolsep}{1pt}
 \begin{center}
\begin{tabular}{@{}F@{}F@{}O@{}I@{}F@{}W@{}F@{}F@{}O}
\instbitrange{31}{29}&
\instbitrange{28}{26}&
\instbitrange{25}{19}&
\instbit{18}&
\instbitrange{17}{15}&
\instbitrange{14}{13}&
\instbitrange{12}{10}&
\instbitrange{9}{7}&
\instbitrange{6}{0} \\
\hline
\multicolumn{1}{|c|}{{\scriptsize qx[2:0]}} &
\multicolumn{1}{c|}{{\scriptsize qy[2:0]}} &
\multicolumn{1}{c|}{1X0100X} &
\multicolumn{1}{c|}{{\scriptsize rs1[4]}} &
\multicolumn{1}{c|}{{\scriptsize rs1[2:0]}} &
\multicolumn{1}{c|}{X0} &
\multicolumn{1}{c|}{{\scriptsize qu[2:0]}} &
\multicolumn{1}{c|}{{\scriptsize qz[2:0]}} &
\multicolumn{1}{c|}{1011111}\\
\hline
3& 3& 7& 1& 3& 2& 3& 3& 7 \\
\end{tabular}
 \end{center}
 }
 \vspace{-0.1in}
 \caption[]{ESP.VMIN.U16.LD.INCP} \end{figure*}


\newpage
\subsubsection{ESP.VMIN.U16.ST.INCP}\label{instruction:ESPVMINU16STINCP}
\textbf{Syntax}\\
ESP.VMIN.U16.ST.INCP qu,rs1,qz,qx,qy

\textbf{Description}\\
This instruction compares the numerical values of the eight 16-bit unsigned vector data segments in the registers \lstinline{qx} and \lstinline{qy}. The data segment with the smaller value is written into the corresponding 16-bit data segment in the register \lstinline{qz}.\\During the operation, the value in the register \lstinline{qu} is stored in memory. After the access, the value in the register \lstinline{rs1} is incremented by 16.\\
\textbf{Operation}
\begin{lstlisting}[escapeinside=`?]
  qz[ 15:  0] = (qx[ 15:  0]<=qy[ 15:  0]) ? qx[ 15:  0] : qy[ 15:  0]
  qz[ 31: 16] = (qx[ 31: 16]<=qy[ 31: 16]) ? qx[ 31: 16] : qy[ 31: 16]
  ...
  qz[127:112] = (qx[127:112]<=qy[127:112]) ? qx[127:112] : qy[127:112]

  qu[127:0] = store128(rs1[31:0])
  rs1[31:0] = rs1[31:0] + 16
\end{lstlisting}


\textbf{Schedule}\\
\begin{longtable}[c]{|l|l|l|}
    \hline
    \rowcolor{lightgray}
    \textbf{Operands} & \textbf{Use} & \textbf{Def} \\\hline
    \endfirsthead
    \hline
    \rowcolor{lightgray}
     \textbf{Operands} & \textbf{Use} & \textbf{Def} \\\hline
    \endhead
    \hline
    \endfoot

    qx & 1 & - \\\hline
  qy & 1 & - \\\hline
     qz & - & 1 \\\hline
     qu & 2 & - \\\hline
     rs1 & 1 & 1

\end{longtable}

\textbf{Format}\\

\begin{figure*}[!htbp]
 {\setlength{\tabcolsep}{1pt}
 \begin{center}
\begin{tabular}{@{}F@{}F@{}O@{}I@{}F@{}W@{}F@{}F@{}O}
\instbitrange{31}{29}&
\instbitrange{28}{26}&
\instbitrange{25}{19}&
\instbit{18}&
\instbitrange{17}{15}&
\instbitrange{14}{13}&
\instbitrange{12}{10}&
\instbitrange{9}{7}&
\instbitrange{6}{0} \\
\hline
\multicolumn{1}{|c|}{{\scriptsize qx[2:0]}} &
\multicolumn{1}{c|}{{\scriptsize qy[2:0]}} &
\multicolumn{1}{c|}{1X1100X} &
\multicolumn{1}{c|}{{\scriptsize rs1[4]}} &
\multicolumn{1}{c|}{{\scriptsize rs1[2:0]}} &
\multicolumn{1}{c|}{X0} &
\multicolumn{1}{c|}{{\scriptsize qu[2:0]}} &
\multicolumn{1}{c|}{{\scriptsize qz[2:0]}} &
\multicolumn{1}{c|}{1011111}\\
\hline
3& 3& 7& 1& 3& 2& 3& 3& 7 \\
\end{tabular}
 \end{center}
 }
 \vspace{-0.1in}
 \caption[]{ESP.VMIN.U16.ST.INCP} \end{figure*}


\newpage
\subsubsection{ESP.VMIN.U32}\label{instruction:ESPVMINU32}
\textbf{Syntax}\\
ESP.VMIN.U32 qz,qx,qy

\textbf{Description}\\
This instruction compares the numerical values of the four 32-bit unsigned vector data segments in the registers \lstinline{qx} and \lstinline{qy}. The data segment with the smaller value is written into the corresponding 32-bit data segment in the register \lstinline{qz}.\\
\textbf{Operation}
\begin{lstlisting}[escapeinside=`?]
  qz[ 31:  0] = (qx[ 31:  0]<=qy[ 31:  0]) ? qx[ 31:  0] : qy[ 31:  0]
  qz[ 63: 32] = (qx[ 63: 32]<=qy[ 63: 32]) ? qx[ 63: 32] : qy[ 63: 32]
  ...
  qz[127: 96] = (qx[127: 96]<=qy[127: 96]) ? qx[127: 96] : qy[127: 96]
\end{lstlisting}


\textbf{Schedule}\\
\begin{longtable}[c]{|l|l|l|}
    \hline
    \rowcolor{lightgray}
    \textbf{Operands} & \textbf{Use} & \textbf{Def} \\\hline
    \endfirsthead
    \hline
    \rowcolor{lightgray}
     \textbf{Operands} & \textbf{Use} & \textbf{Def} \\\hline
    \endhead
    \hline
    \endfoot

    qx & 1 & - \\\hline
  qy & 1 & - \\\hline
     qz & - & 1

\end{longtable}

\textbf{Format}\\

\begin{figure*}[!htbp]
 {\setlength{\tabcolsep}{1pt}
 \begin{center}
\begin{tabular}{@{}F@{}F@{}K@{}F@{}O}
\instbitrange{31}{29}&
\instbitrange{28}{26}&
\instbitrange{25}{10}&
\instbitrange{9}{7}&
\instbitrange{6}{0} \\
\hline
\multicolumn{1}{|c|}{{\scriptsize qx[2:0]}} &
\multicolumn{1}{c|}{{\scriptsize qy[2:0]}} &
\multicolumn{1}{c|}{110100X0X110011X} &
\multicolumn{1}{c|}{{\scriptsize qz[2:0]}} &
\multicolumn{1}{c|}{0011111}\\
\hline
3& 3& 16& 3& 7 \\
\end{tabular}
 \end{center}
 }
 \vspace{-0.1in}
 \caption[]{ESP.VMIN.U32} \end{figure*}


\newpage
\subsubsection{ESP.VMIN.U32.LD.INCP}\label{instruction:ESPVMINU32LDINCP}
\textbf{Syntax}\\
ESP.VMIN.U32.LD.INCP qu,rs1,qz,qx,qy

\textbf{Description}\\
This instruction compares the numerical values of the four 32-bit unsigned vector data segments in registers \lstinline{qx} and \lstinline{qy}. The data segment with the smaller value is written into the corresponding 32-bit data segment in register \lstinline{qz}.\\During the operation, 16 bytes are loaded from the memory into register \lstinline{qu}. After the access, the value in register \lstinline{rs1} is incremented by 16.\\
\textbf{Operation}
\begin{lstlisting}[escapeinside=`?]
  qz[ 31:  0] = (qx[ 31:  0]<=qy[ 31:  0]) ? qx[ 31:  0] : qy[ 31:  0]
  qz[ 63: 32] = (qx[ 63: 32]<=qy[ 63: 32]) ? qx[ 63: 32] : qy[ 63: 32]
  ...
  qz[127: 96] = (qx[127: 96]<=qy[127: 96]) ? qx[127: 96] : qy[127: 96]

  qu[127:0] = load128(rs1[31:0])
  rs1[31:0] = rs1[31:0] + 16
\end{lstlisting}


\textbf{Schedule}\\
\begin{longtable}[c]{|l|l|l|}
    \hline
    \rowcolor{lightgray}
    \textbf{Operands} & \textbf{Use} & \textbf{Def} \\\hline
    \endfirsthead
    \hline
    \rowcolor{lightgray}
     \textbf{Operands} & \textbf{Use} & \textbf{Def} \\\hline
    \endhead
    \hline
    \endfoot

    qx & 1 & - \\\hline
  qy & 1 & - \\\hline
     qz & - & 1 \\\hline
     qu & - & 2 \\\hline
     rs1 & 1 & 1

\end{longtable}

\textbf{Format}\\

\begin{figure*}[!htbp]
 {\setlength{\tabcolsep}{1pt}
 \begin{center}
\begin{tabular}{@{}F@{}F@{}O@{}I@{}F@{}W@{}F@{}F@{}O}
\instbitrange{31}{29}&
\instbitrange{28}{26}&
\instbitrange{25}{19}&
\instbit{18}&
\instbitrange{17}{15}&
\instbitrange{14}{13}&
\instbitrange{12}{10}&
\instbitrange{9}{7}&
\instbitrange{6}{0} \\
\hline
\multicolumn{1}{|c|}{{\scriptsize qx[2:0]}} &
\multicolumn{1}{c|}{{\scriptsize qy[2:0]}} &
\multicolumn{1}{c|}{1X00X1X} &
\multicolumn{1}{c|}{{\scriptsize rs1[4]}} &
\multicolumn{1}{c|}{{\scriptsize rs1[2:0]}} &
\multicolumn{1}{c|}{X0} &
\multicolumn{1}{c|}{{\scriptsize qu[2:0]}} &
\multicolumn{1}{c|}{{\scriptsize qz[2:0]}} &
\multicolumn{1}{c|}{1011111}\\
\hline
3& 3& 7& 1& 3& 2& 3& 3& 7 \\
\end{tabular}
 \end{center}
 }
 \vspace{-0.1in}
 \caption[]{ESP.VMIN.U32.LD.INCP} \end{figure*}


\newpage
\subsubsection{ESP.VMIN.U32.ST.INCP}\label{instruction:ESPVMINU32STINCP}
\textbf{Syntax}\\
ESP.VMIN.U32.ST.INCP qu,rs1,qz,qx,qy

\textbf{Description}\\
This instruction compares the numerical values of the four 32-bit unsigned vector data segments in the registers \lstinline{qx} and \lstinline{qy}. The data segment with the smaller value is written into the corresponding 32-bit data segment in the register \lstinline{qz}.\\During the operation, the value in the register \lstinline{qu} is stored in memory. After the access, the value in the register \lstinline{rs1} is incremented by 16.\\
\textbf{Operation}
\begin{lstlisting}[escapeinside=`?]
  qz[ 31:  0] = (qx[ 31:  0]<=qy[ 31:  0]) ? qx[ 31:  0] : qy[ 31:  0]
  qz[ 63: 32] = (qx[ 63: 32]<=qy[ 63: 32]) ? qx[ 63: 32] : qy[ 63: 32]
  ...
  qz[127: 96] = (qx[127: 96]<=qy[127: 96]) ? qx[127: 96] : qy[127: 96]

  qu[127:0] = store128(rs1[31:0])
  rs1[31:0] = rs1[31:0] + 16
\end{lstlisting}


\textbf{Schedule}\\
\begin{longtable}[c]{|l|l|l|}
    \hline
    \rowcolor{lightgray}
    \textbf{Operands} & \textbf{Use} & \textbf{Def} \\\hline
    \endfirsthead
    \hline
    \rowcolor{lightgray}
     \textbf{Operands} & \textbf{Use} & \textbf{Def} \\\hline
    \endhead
    \hline
    \endfoot

    qx & 1 & - \\\hline
  qy & 1 & - \\\hline
     qz & - & 1 \\\hline
     qu & 2 & - \\\hline
     rs1 & 1 & 1

\end{longtable}

\textbf{Format}\\

\begin{figure*}[!htbp]
 {\setlength{\tabcolsep}{1pt}
 \begin{center}
\begin{tabular}{@{}F@{}F@{}O@{}I@{}F@{}W@{}F@{}F@{}O}
\instbitrange{31}{29}&
\instbitrange{28}{26}&
\instbitrange{25}{19}&
\instbit{18}&
\instbitrange{17}{15}&
\instbitrange{14}{13}&
\instbitrange{12}{10}&
\instbitrange{9}{7}&
\instbitrange{6}{0} \\
\hline
\multicolumn{1}{|c|}{{\scriptsize qx[2:0]}} &
\multicolumn{1}{c|}{{\scriptsize qy[2:0]}} &
\multicolumn{1}{c|}{1X10X1X} &
\multicolumn{1}{c|}{{\scriptsize rs1[4]}} &
\multicolumn{1}{c|}{{\scriptsize rs1[2:0]}} &
\multicolumn{1}{c|}{X0} &
\multicolumn{1}{c|}{{\scriptsize qu[2:0]}} &
\multicolumn{1}{c|}{{\scriptsize qz[2:0]}} &
\multicolumn{1}{c|}{1011111}\\
\hline
3& 3& 7& 1& 3& 2& 3& 3& 7 \\
\end{tabular}
 \end{center}
 }
 \vspace{-0.1in}
 \caption[]{ESP.VMIN.U32.ST.INCP} \end{figure*}


\newpage
\subsubsection{ESP.VMIN.U8}\label{instruction:ESPVMINU8}
\textbf{Syntax}\\
ESP.VMIN.U8 qz,qx,qy

\textbf{Description}\\
This instruction compares the numerical values of the 16 8-bit unsigned vector data segments in registers \lstinline{qx} and \lstinline{qy}. The data segment with the smaller value is written into the corresponding 8-bit data segment in the register \lstinline{qz}.\\
\textbf{Operation}
\begin{lstlisting}[escapeinside=`?]
  qz[  7:  0] = (qx[  7:  0]<=qy[  7:  0]) ? qx[  7:  0] : qy[  7:  0]
  qz[ 15:  8] = (qx[ 15:  8]<=qy[ 15:  8]) ? qx[ 15:  8] : qy[ 15:  8]
  ...
  qz[127:120] = (qx[127:120]<=qy[127:120]) ? qx[127:120] : qy[127:120]
\end{lstlisting}


\textbf{Schedule}\\
\begin{longtable}[c]{|l|l|l|}
    \hline
    \rowcolor{lightgray}
    \textbf{Operands} & \textbf{Use} & \textbf{Def} \\\hline
    \endfirsthead
    \hline
    \rowcolor{lightgray}
     \textbf{Operands} & \textbf{Use} & \textbf{Def} \\\hline
    \endhead
    \hline
    \endfoot

    qx & 1 & - \\\hline
  qy & 1 & - \\\hline
     qz & - & 1

\end{longtable}

\textbf{Format}\\

\begin{figure*}[!htbp]
 {\setlength{\tabcolsep}{1pt}
 \begin{center}
\begin{tabular}{@{}F@{}F@{}K@{}F@{}O}
\instbitrange{31}{29}&
\instbitrange{28}{26}&
\instbitrange{25}{10}&
\instbitrange{9}{7}&
\instbitrange{6}{0} \\
\hline
\multicolumn{1}{|c|}{{\scriptsize qx[2:0]}} &
\multicolumn{1}{c|}{{\scriptsize qy[2:0]}} &
\multicolumn{1}{c|}{110100X00010011X} &
\multicolumn{1}{c|}{{\scriptsize qz[2:0]}} &
\multicolumn{1}{c|}{0011111}\\
\hline
3& 3& 16& 3& 7 \\
\end{tabular}
 \end{center}
 }
 \vspace{-0.1in}
 \caption[]{ESP.VMIN.U8} \end{figure*}


\newpage
\subsubsection{ESP.VMIN.U8.LD.INCP}\label{instruction:ESPVMINU8LDINCP}
\textbf{Syntax}\\
ESP.VMIN.U8.LD.INCP qu,rs1,qz,qx,qy

\textbf{Description}\\
This instruction compares the numerical values of the 16 8-bit unsigned vector data segments in registers \lstinline{qx} and \lstinline{qy}. The data segment with the smaller value is written into the corresponding 8-bit data segment in register \lstinline{qz}.\\During the operation, 16 bytes are loaded from the memory into register \lstinline{qu}. After the access, the value in register \lstinline{rs1} is incremented by 16.\\
\textbf{Operation}
\begin{lstlisting}[escapeinside=`?]
  qz[  7:  0] = (qx[  7:  0]<=qy[  7:  0]) ? qx[  7:  0] : qy[  7:  0]
  qz[ 15:  8] = (qx[ 15:  8]<=qy[ 15:  8]) ? qx[ 15:  8] : qy[ 15:  8]
  ...
  qz[127:120] = (qx[127:120]<=qy[127:120]) ? qx[127:120] : qy[127:120]

  qu[127:0] = load128(rs1[31:0])
  rs1[31:0] = rs1[31:0] + 16
\end{lstlisting}


\textbf{Schedule}\\
\begin{longtable}[c]{|l|l|l|}
    \hline
    \rowcolor{lightgray}
    \textbf{Operands} & \textbf{Use} & \textbf{Def} \\\hline
    \endfirsthead
    \hline
    \rowcolor{lightgray}
     \textbf{Operands} & \textbf{Use} & \textbf{Def} \\\hline
    \endhead
    \hline
    \endfoot

    qx & 1 & - \\\hline
  qy & 1 & - \\\hline
     qz & - & 1 \\\hline
     qu & - & 2 \\\hline
     rs1 & 1 & 1

\end{longtable}

\textbf{Format}\\

\begin{figure*}[!htbp]
 {\setlength{\tabcolsep}{1pt}
 \begin{center}
\begin{tabular}{@{}F@{}F@{}O@{}I@{}F@{}W@{}F@{}F@{}O}
\instbitrange{31}{29}&
\instbitrange{28}{26}&
\instbitrange{25}{19}&
\instbit{18}&
\instbitrange{17}{15}&
\instbitrange{14}{13}&
\instbitrange{12}{10}&
\instbitrange{9}{7}&
\instbitrange{6}{0} \\
\hline
\multicolumn{1}{|c|}{{\scriptsize qx[2:0]}} &
\multicolumn{1}{c|}{{\scriptsize qy[2:0]}} &
\multicolumn{1}{c|}{1X0000X} &
\multicolumn{1}{c|}{{\scriptsize rs1[4]}} &
\multicolumn{1}{c|}{{\scriptsize rs1[2:0]}} &
\multicolumn{1}{c|}{X0} &
\multicolumn{1}{c|}{{\scriptsize qu[2:0]}} &
\multicolumn{1}{c|}{{\scriptsize qz[2:0]}} &
\multicolumn{1}{c|}{1011111}\\
\hline
3& 3& 7& 1& 3& 2& 3& 3& 7 \\
\end{tabular}
 \end{center}
 }
 \vspace{-0.1in}
 \caption[]{ESP.VMIN.U8.LD.INCP} \end{figure*}


\newpage
\subsubsection{ESP.VMIN.U8.ST.INCP}\label{instruction:ESPVMINU8STINCP}
\textbf{Syntax}\\
ESP.VMIN.U8.ST.INCP qu,rs1,qz,qx,qy

\textbf{Description}\\
This instruction compares the numerical values of the 16 8-bit unsigned vector data segments in registers \lstinline{qx} and \lstinline{qy}. The data segment with the smaller value is written into the corresponding 8-bit data segment in register \lstinline{qz}.\\During the operation, the value in register \lstinline{qu} is stored in memory. After the access, the value in register \lstinline{rs1} is incremented by 16.\\
\textbf{Operation}
\begin{lstlisting}[escapeinside=`?]
  qz[  7:  0] = (qx[  7:  0]<=qy[  7:  0]) ? qx[  7:  0] : qy[  7:  0]
  qz[ 15:  8] = (qx[ 15:  8]<=qy[ 15:  8]) ? qx[ 15:  8] : qy[ 15:  8]
  ...
  qz[127:120] = (qx[127:120]<=qy[127:120]) ? qx[127:120] : qy[127:120]

  qu[127:0] = store128(rs1[31:0])
  rs1[31:0] = rs1[31:0] + 16
\end{lstlisting}


\textbf{Schedule}\\
\begin{longtable}[c]{|l|l|l|}
    \hline
    \rowcolor{lightgray}
    \textbf{Operands} & \textbf{Use} & \textbf{Def} \\\hline
    \endfirsthead
    \hline
    \rowcolor{lightgray}
     \textbf{Operands} & \textbf{Use} & \textbf{Def} \\\hline
    \endhead
    \hline
    \endfoot

    qx & 1 & - \\\hline
  qy & 1 & - \\\hline
     qz & - & 1 \\\hline
     qu & 2 & - \\\hline
     rs1 & 1 & 1

\end{longtable}

\textbf{Format}\\

\begin{figure*}[!htbp]
 {\setlength{\tabcolsep}{1pt}
 \begin{center}
\begin{tabular}{@{}F@{}F@{}O@{}I@{}F@{}W@{}F@{}F@{}O}
\instbitrange{31}{29}&
\instbitrange{28}{26}&
\instbitrange{25}{19}&
\instbit{18}&
\instbitrange{17}{15}&
\instbitrange{14}{13}&
\instbitrange{12}{10}&
\instbitrange{9}{7}&
\instbitrange{6}{0} \\
\hline
\multicolumn{1}{|c|}{{\scriptsize qx[2:0]}} &
\multicolumn{1}{c|}{{\scriptsize qy[2:0]}} &
\multicolumn{1}{c|}{1X1000X} &
\multicolumn{1}{c|}{{\scriptsize rs1[4]}} &
\multicolumn{1}{c|}{{\scriptsize rs1[2:0]}} &
\multicolumn{1}{c|}{X0} &
\multicolumn{1}{c|}{{\scriptsize qu[2:0]}} &
\multicolumn{1}{c|}{{\scriptsize qz[2:0]}} &
\multicolumn{1}{c|}{1011111}\\
\hline
3& 3& 7& 1& 3& 2& 3& 3& 7 \\
\end{tabular}
 \end{center}
 }
 \vspace{-0.1in}
 \caption[]{ESP.VMIN.U8.ST.INCP} \end{figure*}


\newpage
\subsubsection{ESP.VMUL.S16}\label{instruction:ESPVMULS16}
\textbf{Syntax}\\
ESP.VMUL.S16 qz,qx,qy,sat,rm

\textbf{Description}\\
This instruction performs a signed vector multiplication on 16-bit data. Registers \lstinline{qx} and \lstinline{qy} are the multiplier and the multiplicand respectively. The eight 32-bit data results obtained from the calculation are arithmetically right-shifted by the value in the special register \lstinline{SAR}. Then, the lower 16-bit data of the shift result are written into the corresponding segment of the register \lstinline{qz}.\\
\textbf{Operation}
\begin{lstlisting}
  qz[ 15:  0] = (qx[ 15:  0] * qy[ 15:  0]) >> SAR[5:0]
  qz[ 31: 16] = (qx[ 31: 16] * qy[ 31: 16]) >> SAR[5:0]
  qz[ 47: 32] = (qx[ 47: 32] * qy[ 47: 32]) >> SAR[5:0]
  qz[ 63: 48] = (qx[ 63: 48] * qy[ 63: 48]) >> SAR[5:0]
  qz[ 79: 64] = (qx[ 79: 64] * qy[ 79: 64]) >> SAR[5:0]
  qz[ 95: 80] = (qx[ 95: 80] * qy[ 95: 80]) >> SAR[5:0]
  qz[111: 96] = (qx[111: 96] * qy[111: 96]) >> SAR[5:0]
  qz[127:112] = (qx[127:112] * qy[127:112]) >> SAR[5:0]
\end{lstlisting}


\textbf{Schedule}\\
\begin{longtable}[c]{|l|l|l|}
    \hline
    \rowcolor{lightgray}
    \textbf{Operands} & \textbf{Use} & \textbf{Def} \\\hline
    \endfirsthead
    \hline
    \rowcolor{lightgray}
     \textbf{Operands} & \textbf{Use} & \textbf{Def} \\\hline
    \endhead
    \hline
    \endfoot

    qx & 1 & - \\\hline
  qy & 1 & - \\\hline
     qz & - & 2 \\\hline
    SAR & 1 & -
\end{longtable}

\textbf{Format}\\

\begin{figure*}[!htbp]
 {\setlength{\tabcolsep}{1pt}
 \begin{center}
\begin{tabular}{@{}F@{}F@{}T@{}I@{}I@{}W@{}W@{}I@{}F@{}O}
\instbitrange{31}{29}&
\instbitrange{28}{26}&
\instbitrange{25}{17}&
\instbit{16}&
\instbit{15}&
\instbitrange{14}{13}&
\instbitrange{12}{11}&
\instbit{10}&
\instbitrange{9}{7}&
\instbitrange{6}{0} \\
\hline
\multicolumn{1}{|c|}{{\scriptsize qx[2:0]}} &
\multicolumn{1}{c|}{{\scriptsize qy[2:0]}} &
\multicolumn{1}{c|}{111100X11} &
\multicolumn{1}{c|}{{\scriptsize sat[0]}} &
\multicolumn{1}{c|}{{\scriptsize rm[2]}} &
\multicolumn{1}{c|}{00} &
\multicolumn{1}{c|}{{\scriptsize rm[1:0]}} &
\multicolumn{1}{c|}{0} &
\multicolumn{1}{c|}{{\scriptsize qz[2:0]}} &
\multicolumn{1}{c|}{0011111}\\
\hline
3& 3& 9& 1& 1& 2& 2& 1& 3& 7 \\
\end{tabular}
 \end{center}
 }
 \vspace{-0.1in}
 \caption[]{ESP.VMUL.S16} \end{figure*}


\newpage
\subsubsection{ESP.VMUL.S16.LD.INCP}\label{instruction:ESPVMULS16LDINCP}
\textbf{Syntax}\\
ESP.VMUL.S16.LD.INCP qu,rs1,qz,qx,qy,sat,rm

\textbf{Description}\\
This instruction performs a signed vector multiplication on 16-bit data. Registers \lstinline{qx} and \lstinline{qy} are the multiplier and the multiplicand respectively. The eight 32-bit data results obtained from the calculation are arithmetically right-shifted by the value in the special register \lstinline{SAR}. Then, the lower 16-bit data of the shift result are written into the corresponding segment of the register \lstinline{qz}.\\
During the operation, 16 bytes are loaded from the memory into the register \lstinline{qu}. After the access, the value in the register \lstinline{rs1} is incremented by 16.\\
\textbf{Operation}
\begin{lstlisting}
  qz[ 15:  0] = (qx[ 15:  0] * qy[ 15:  0]) >> SAR[5:0]
  qz[ 31: 16] = (qx[ 31: 16] * qy[ 31: 16]) >> SAR[5:0]
  qz[ 47: 32] = (qx[ 47: 32] * qy[ 47: 32]) >> SAR[5:0]
  qz[ 63: 48] = (qx[ 63: 48] * qy[ 63: 48]) >> SAR[5:0]
  qz[ 79: 64] = (qx[ 79: 64] * qy[ 79: 64]) >> SAR[5:0]
  qz[ 95: 80] = (qx[ 95: 80] * qy[ 95: 80]) >> SAR[5:0]
  qz[111: 96] = (qx[111: 96] * qy[111: 96]) >> SAR[5:0]
  qz[127:112] = (qx[127:112] * qy[127:112]) >> SAR[5:0]

  qu[127:0] = load128(rs1[31:0])
  rs1[31:0] = rs1[31:0] + 16
\end{lstlisting}


\textbf{Schedule}\\
\begin{longtable}[c]{|l|l|l|}
    \hline
    \rowcolor{lightgray}
    \textbf{Operands} & \textbf{Use} & \textbf{Def} \\\hline
    \endfirsthead
    \hline
    \rowcolor{lightgray}
     \textbf{Operands} & \textbf{Use} & \textbf{Def} \\\hline
    \endhead
    \hline
    \endfoot

    qx & 1 & - \\\hline
  qy & 1 & - \\\hline
     qz & - & 2 \\\hline
    SAR & 1 & - \\\hline
    qu & - & 2 \\\hline
    rs1 & 1 & 1
\end{longtable}

\textbf{Format}\\

\begin{figure*}[!htbp]
 {\setlength{\tabcolsep}{1pt}
 \begin{center}
\begin{tabular}{@{}F@{}F@{}F@{}I@{}F@{}I@{}F@{}W@{}F@{}F@{}O}
\instbitrange{31}{29}&
\instbitrange{28}{26}&
\instbitrange{25}{23}&
\instbit{22}&
\instbitrange{21}{19}&
\instbit{18}&
\instbitrange{17}{15}&
\instbitrange{14}{13}&
\instbitrange{12}{10}&
\instbitrange{9}{7}&
\instbitrange{6}{0} \\
\hline
\multicolumn{1}{|c|}{{\scriptsize qx[2:0]}} &
\multicolumn{1}{c|}{{\scriptsize qy[2:0]}} &
\multicolumn{1}{c|}{011} &
\multicolumn{1}{c|}{{\scriptsize sat[0]}} &
\multicolumn{1}{c|}{{\scriptsize rm[2:0]}} &
\multicolumn{1}{c|}{{\scriptsize rs1[4]}} &
\multicolumn{1}{c|}{{\scriptsize rs1[2:0]}} &
\multicolumn{1}{c|}{01} &
\multicolumn{1}{c|}{{\scriptsize qu[2:0]}} &
\multicolumn{1}{c|}{{\scriptsize qz[2:0]}} &
\multicolumn{1}{c|}{1011111}\\
\hline
3& 3& 3& 1& 3& 1& 3& 2& 3& 3& 7 \\
\end{tabular}
 \end{center}
 }
 \vspace{-0.1in}
 \caption[]{ESP.VMUL.S16.LD.INCP} \end{figure*}


\newpage
\subsubsection{ESP.VMUL.S16.S8XS8}\label{instruction:ESPVMULS16S8XS8}
\textbf{Syntax}\\
ESP.VMUL.S16.S8XS8 qz,qv,qx,qy,rm

\textbf{Description}\\
This instruction performs a signed vector multiplication on 8-bit data. Registers \lstinline{qx} and \lstinline{qy} are the multiplier and the multiplicand respectively. The 16 16-bit data results obtained from the calculation are arithmetically right-shifted by the value in the special register \lstinline{SAR}. Then, the shift results are written into the corresponding segments of registers \lstinline{qz} and \lstinline{qv}.\\
\textbf{Operation}
\begin{lstlisting}
  qz[ 15:  0] = (qx[  7:  0] * qy[  7:  0]) >> SAR[5:0]
  qz[ 31: 16] = (qx[ 15:  8] * qy[ 15:  8]) >> SAR[5:0]
  qz[ 47: 32] = (qx[ 23: 16] * qy[ 23: 16]) >> SAR[5:0]
  qz[ 63: 48] = (qx[ 31: 24] * qy[ 31: 24]) >> SAR[5:0]
  qz[ 79: 64] = (qx[ 39: 32] * qy[ 39: 32]) >> SAR[5:0]
  qz[ 95: 80] = (qx[ 47: 40] * qy[ 47: 40]) >> SAR[5:0]
  qz[111: 96] = (qx[ 55: 48] * qy[ 55: 48]) >> SAR[5:0]
  qz[127:112] = (qx[ 63: 56] * qy[ 63: 56]) >> SAR[5:0]

  qv[ 15:  0] = (qx[ 71: 64] * qy[ 71: 64]) >> SAR[5:0]
  qv[ 31: 16] = (qx[ 79: 72] * qy[ 79: 72]) >> SAR[5:0]
  qv[ 47: 32] = (qx[ 87: 80] * qy[ 87: 80]) >> SAR[5:0]
  qv[ 63: 48] = (qx[ 95: 88] * qy[ 95: 88]) >> SAR[5:0]
  qv[ 79: 64] = (qx[103: 96] * qy[103: 96]) >> SAR[5:0]
  qv[ 95: 80] = (qx[111:104] * qy[111:104]) >> SAR[5:0]
  qv[111: 96] = (qx[119:112] * qy[119:112]) >> SAR[5:0]
  qv[127:112] = (qx[127:120] * qy[127:120]) >> SAR[5:0]
\end{lstlisting}


\textbf{Schedule}\\
\begin{longtable}[c]{|l|l|l|}
    \hline
    \rowcolor{lightgray}
    \textbf{Operands} & \textbf{Use} & \textbf{Def} \\\hline
    \endfirsthead
    \hline
    \rowcolor{lightgray}
     \textbf{Operands} & \textbf{Use} & \textbf{Def} \\\hline
    \endhead
    \hline
    \endfoot

    qx & 1 & - \\\hline
  qy & 1 & - \\\hline
     qz & - & 2 \\\hline
     qv & - & 2 \\\hline
    SAR & 1 & -

\end{longtable}

\textbf{Format}\\

\begin{figure*}[!htbp]
 {\setlength{\tabcolsep}{1pt}
 \begin{center}
\begin{tabular}{@{}F@{}F@{}F@{}F@{}I@{}F@{}S@{}F@{}O}
\instbitrange{31}{29}&
\instbitrange{28}{26}&
\instbitrange{25}{23}&
\instbitrange{22}{20}&
\instbit{19}&
\instbitrange{18}{16}&
\instbitrange{15}{10}&
\instbitrange{9}{7}&
\instbitrange{6}{0} \\
\hline
\multicolumn{1}{|c|}{{\scriptsize qx[2:0]}} &
\multicolumn{1}{c|}{{\scriptsize qy[2:0]}} &
\multicolumn{1}{c|}{101} &
\multicolumn{1}{c|}{{\scriptsize qv[2:0]}} &
\multicolumn{1}{c|}{X} &
\multicolumn{1}{c|}{{\scriptsize rm[2:0]}} &
\multicolumn{1}{c|}{0000XX} &
\multicolumn{1}{c|}{{\scriptsize qz[2:0]}} &
\multicolumn{1}{c|}{0011111}\\
\hline
3& 3& 3& 3& 1& 3& 6& 3& 7 \\
\end{tabular}
 \end{center}
 }
 \vspace{-0.1in}
 \caption[]{ESP.VMUL.S16.S8XS8} \end{figure*}


\newpage
\subsubsection{ESP.VMUL.S16.ST.INCP}\label{instruction:ESPVMULS16STINCP}
\textbf{Syntax}\\
ESP.VMUL.S16.ST.INCP qu,rs1,qz,qx,qy,sat,rm

\textbf{Description}\\
This instruction performs a signed vector multiplication on 16-bit data. Registers \lstinline{qx} and \lstinline{qy} are the multiplier and the multiplicand respectively. The eight 32-bit data results obtained from the calculation are arithmetically right-shifted by the value in the special register \lstinline{SAR}. Then, the lower 16-bit data of the shift result are written into the corresponding segment of the register \lstinline{qz}.
\\During the operation, the value in the register \lstinline{qu} is stored to memory. After the access, the value in the register \lstinline{rs1} is incremented by 16.\\
\textbf{Operation}
\begin{lstlisting}
  qz[ 15:  0] = (qx[ 15:  0] * qy[ 15:  0]) >> SAR[5:0]
  qz[ 31: 16] = (qx[ 31: 16] * qy[ 31: 16]) >> SAR[5:0]
  qz[ 47: 32] = (qx[ 47: 32] * qy[ 47: 32]) >> SAR[5:0]
  qz[ 63: 48] = (qx[ 63: 48] * qy[ 63: 48]) >> SAR[5:0]
  qz[ 79: 64] = (qx[ 79: 64] * qy[ 79: 64]) >> SAR[5:0]
  qz[ 95: 80] = (qx[ 95: 80] * qy[ 95: 80]) >> SAR[5:0]
  qz[111: 96] = (qx[111: 96] * qy[111: 96]) >> SAR[5:0]
  qz[127:112] = (qx[127:112] * qy[127:112]) >> SAR[5:0]

  qu[127:0] => store128(rs1[31:0])
  rs1[31:0] = rs1[31:0] + 16
\end{lstlisting}


\textbf{Schedule}\\
\begin{longtable}[c]{|l|l|l|}
    \hline
    \rowcolor{lightgray}
    \textbf{Operands} & \textbf{Use} & \textbf{Def} \\\hline
    \endfirsthead
    \hline
    \rowcolor{lightgray}
     \textbf{Operands} & \textbf{Use} & \textbf{Def} \\\hline
    \endhead
    \hline
    \endfoot

    qx & 1 & - \\\hline
  qy & 1 & - \\\hline
     qz & - & 2 \\\hline
    SAR & 1 & － \\\hline
    qu & 2 & - \\\hline
    rs1 & 1 & 1

\end{longtable}

\textbf{Format}\\

\begin{figure*}[!htbp]
 {\setlength{\tabcolsep}{1pt}
 \begin{center}
\begin{tabular}{@{}F@{}F@{}F@{}I@{}F@{}I@{}F@{}W@{}F@{}F@{}O}
\instbitrange{31}{29}&
\instbitrange{28}{26}&
\instbitrange{25}{23}&
\instbit{22}&
\instbitrange{21}{19}&
\instbit{18}&
\instbitrange{17}{15}&
\instbitrange{14}{13}&
\instbitrange{12}{10}&
\instbitrange{9}{7}&
\instbitrange{6}{0} \\
\hline
\multicolumn{1}{|c|}{{\scriptsize qx[2:0]}} &
\multicolumn{1}{c|}{{\scriptsize qy[2:0]}} &
\multicolumn{1}{c|}{111} &
\multicolumn{1}{c|}{{\scriptsize sat[0]}} &
\multicolumn{1}{c|}{{\scriptsize rm[2:0]}} &
\multicolumn{1}{c|}{{\scriptsize rs1[4]}} &
\multicolumn{1}{c|}{{\scriptsize rs1[2:0]}} &
\multicolumn{1}{c|}{01} &
\multicolumn{1}{c|}{{\scriptsize qu[2:0]}} &
\multicolumn{1}{c|}{{\scriptsize qz[2:0]}} &
\multicolumn{1}{c|}{1011111}\\
\hline
3& 3& 3& 1& 3& 1& 3& 2& 3& 3& 7 \\
\end{tabular}
 \end{center}
 }
 \vspace{-0.1in}
 \caption[]{ESP.VMUL.S16.ST.INCP} \end{figure*}


\newpage
\subsubsection{ESP.VMUL.S32.S16XS16}\label{instruction:ESPVMULS32S16XS16}
\textbf{Syntax}\\
ESP.VMUL.S32.S16XS16 qz,qv,qx,qy,rm

\textbf{Description}\\
This instruction performs a signed vector multiplication on 16-bit data. Registers \lstinline{qx} and \lstinline{qy} are the multiplier and the multiplicand respectively. The eight 32-bit data results obtained from the calculation are arithmetically right-shifted by the value in the special register \lstinline{SAR}. Then, the shift results are written into the corresponding segments of registers \lstinline{qz} and \lstinline{qv}.\\
\textbf{Operation}
\begin{lstlisting}
  qz[ 31:  0] = (qx[ 15:  0] * qy[ 15:  0]) >> SAR[5:0]
  qz[ 63:  0] = (qx[ 31: 16] * qy[ 31: 16]) >> SAR[5:0]
  qz[ 95: 64] = (qx[ 47: 32] * qy[ 47: 32]) >> SAR[5:0]
  qz[127: 96] = (qx[ 63: 48] * qy[ 63: 48]) >> SAR[5:0]

  qv[ 31:  0] = (qx[ 79: 64] * qy[ 79: 64]) >> SAR[5:0]
  qv[ 63: 32] = (qx[ 95: 80] * qy[ 95: 80]) >> SAR[5:0]
  qv[ 95: 64] = (qx[111: 96] * qy[111: 96]) >> SAR[5:0]
  qv[127: 96] = (qx[127:112] * qy[127:112]) >> SAR[5:0]
\end{lstlisting}


\textbf{Schedule}\\
\begin{longtable}[c]{|l|l|l|}
    \hline
    \rowcolor{lightgray}
    \textbf{Operands} & \textbf{Use} & \textbf{Def} \\\hline
    \endfirsthead
    \hline
    \rowcolor{lightgray}
     \textbf{Operands} & \textbf{Use} & \textbf{Def} \\\hline
    \endhead
    \hline
    \endfoot

    qx & 1 & - \\\hline
  qy & 1 & - \\\hline
     qz & - & 2 \\\hline
     qv & - & 2 \\\hline
    SAR & 1 & -

\end{longtable}

\textbf{Format}\\

\begin{figure*}[!htbp]
 {\setlength{\tabcolsep}{1pt}
 \begin{center}
\begin{tabular}{@{}F@{}F@{}F@{}F@{}I@{}F@{}S@{}F@{}O}
\instbitrange{31}{29}&
\instbitrange{28}{26}&
\instbitrange{25}{23}&
\instbitrange{22}{20}&
\instbit{19}&
\instbitrange{18}{16}&
\instbitrange{15}{10}&
\instbitrange{9}{7}&
\instbitrange{6}{0} \\
\hline
\multicolumn{1}{|c|}{{\scriptsize qx[2:0]}} &
\multicolumn{1}{c|}{{\scriptsize qy[2:0]}} &
\multicolumn{1}{c|}{101} &
\multicolumn{1}{c|}{{\scriptsize qv[2:0]}} &
\multicolumn{1}{c|}{X} &
\multicolumn{1}{c|}{{\scriptsize rm[2:0]}} &
\multicolumn{1}{c|}{1000XX} &
\multicolumn{1}{c|}{{\scriptsize qz[2:0]}} &
\multicolumn{1}{c|}{0011111}\\
\hline
3& 3& 3& 3& 1& 3& 6& 3& 7 \\
\end{tabular}
 \end{center}
 }
 \vspace{-0.1in}
 \caption[]{ESP.VMUL.S32.S16XS16} \end{figure*}


\newpage
\subsubsection{ESP.VMUL.S8}\label{instruction:ESPVMULS8}
\textbf{Syntax}\\
ESP.VMUL.S8 qz,qx,qy,sat,rm

\textbf{Description}\\
This instruction performs a signed vector multiplication on 8-bit data. Registers \lstinline{qx} and \lstinline{qy} are the multiplier and the multiplicand respectively. The eight 16-bit data results obtained from the calculation are arithmetically right-shifted by the value in the special register \lstinline{SAR}. Then, the lower 8-bit data of the shift result are written into the corresponding segment of the register \lstinline{qz}.\\
\textbf{Operation}
\begin{lstlisting}
  qz[  7:  0] = (qx[  7:  0] * qy[  7:  0]) >> SAR[4:0]
  qz[ 15:  8] = (qx[ 15:  8] * qy[ 15:  8]) >> SAR[4:0]
  qz[ 23: 16] = (qx[ 23: 16] * qy[ 23: 16]) >> SAR[4:0]
  qz[ 31: 24] = (qx[ 31: 24] * qy[ 31: 24]) >> SAR[4:0]
  qz[ 39: 32] = (qx[ 39: 32] * qy[ 39: 32]) >> SAR[4:0]
  qz[ 47: 40] = (qx[ 47: 40] * qy[ 47: 40]) >> SAR[4:0]
  qz[ 55: 48] = (qx[ 55: 48] * qy[ 55: 48]) >> SAR[4:0]
  qz[ 63: 56] = (qx[ 63: 56] * qy[ 63: 56]) >> SAR[4:0]
  qz[ 71: 64] = (qx[ 71: 64] * qy[ 71: 64]) >> SAR[4:0]
  qz[ 79: 71] = (qx[ 79: 72] * qy[ 79: 72]) >> SAR[4:0]
  qz[ 87: 80] = (qx[ 87: 80] * qy[ 87: 80]) >> SAR[4:0]
  qz[ 95: 88] = (qx[ 95: 88] * qy[ 95: 88]) >> SAR[4:0]
  qz[103: 96] = (qx[103: 96] * qy[103: 96]) >> SAR[4:0]
  qz[111:104] = (qx[111:104] * qy[111:104]) >> SAR[4:0]
  qz[119:112] = (qx[119:112] * qy[119:112]) >> SAR[4:0]
  qz[127:120] = (qx[127:120] * qy[127:120]) >> SAR[4:0]

\end{lstlisting}


\textbf{Schedule}\\
\begin{longtable}[c]{|l|l|l|}
    \hline
    \rowcolor{lightgray}
    \textbf{Operands} & \textbf{Use} & \textbf{Def} \\\hline
    \endfirsthead
    \hline
    \rowcolor{lightgray}
     \textbf{Operands} & \textbf{Use} & \textbf{Def} \\\hline
    \endhead
    \hline
    \endfoot

    qx & 1 & - \\\hline
  qy & 1 & - \\\hline
     qz & - & 2 \\\hline
    SAR & 1 & -

\end{longtable}

\textbf{Format}\\

\begin{figure*}[!htbp]
 {\setlength{\tabcolsep}{1pt}
 \begin{center}
\begin{tabular}{@{}F@{}F@{}T@{}I@{}I@{}W@{}W@{}I@{}F@{}O}
\instbitrange{31}{29}&
\instbitrange{28}{26}&
\instbitrange{25}{17}&
\instbit{16}&
\instbit{15}&
\instbitrange{14}{13}&
\instbitrange{12}{11}&
\instbit{10}&
\instbitrange{9}{7}&
\instbitrange{6}{0} \\
\hline
\multicolumn{1}{|c|}{{\scriptsize qx[2:0]}} &
\multicolumn{1}{c|}{{\scriptsize qy[2:0]}} &
\multicolumn{1}{c|}{111100X01} &
\multicolumn{1}{c|}{{\scriptsize sat[0]}} &
\multicolumn{1}{c|}{{\scriptsize rm[2]}} &
\multicolumn{1}{c|}{00} &
\multicolumn{1}{c|}{{\scriptsize rm[1:0]}} &
\multicolumn{1}{c|}{0} &
\multicolumn{1}{c|}{{\scriptsize qz[2:0]}} &
\multicolumn{1}{c|}{0011111}\\
\hline
3& 3& 9& 1& 1& 2& 2& 1& 3& 7 \\
\end{tabular}
 \end{center}
 }
 \vspace{-0.1in}
 \caption[]{ESP.VMUL.S8} \end{figure*}


\newpage
\subsubsection{ESP.VMUL.S8.LD.INCP}\label{instruction:ESPVMULS8LDINCP}
\textbf{Syntax}\\
ESP.VMUL.S8.LD.INCP qu,rs1,qz,qx,qy,sat,rm

\textbf{Description}\\
This instruction performs a signed vector multiplication on 8-bit data. Registers \lstinline{qx} and \lstinline{qy} are the multiplier and the multiplicand respectively. The eight 16-bit data results obtained from the calculation are arithmetically right-shifted by the value in the special register \lstinline{SAR}. Then, the lower 8-bit data of the shift result are written into the corresponding segment of the register \lstinline{qz}.\\
During the operation, 16 bytes are loaded from the memory into the register \lstinline{qu}. After the access, the value in the register \lstinline{rs1} is incremented by 16.\\
\textbf{Operation}
\begin{lstlisting}
  qz[  7:  0] = (qx[  7:  0] * qy[  7:  0]) >> SAR[4:0]
  qz[ 15:  8] = (qx[ 15:  8] * qy[ 15:  8]) >> SAR[4:0]
  qz[ 23: 16] = (qx[ 23: 16] * qy[ 23: 16]) >> SAR[4:0]
  qz[ 31: 24] = (qx[ 31: 24] * qy[ 31: 24]) >> SAR[4:0]
  qz[ 39: 32] = (qx[ 39: 32] * qy[ 39: 32]) >> SAR[4:0]
  qz[ 47: 40] = (qx[ 47: 40] * qy[ 47: 40]) >> SAR[4:0]
  qz[ 55: 48] = (qx[ 55: 48] * qy[ 55: 48]) >> SAR[4:0]
  qz[ 63: 56] = (qx[ 63: 56] * qy[ 63: 56]) >> SAR[4:0]
  qz[ 71: 64] = (qx[ 71: 64] * qy[ 71: 64]) >> SAR[4:0]
  qz[ 79: 71] = (qx[ 79: 72] * qy[ 79: 72]) >> SAR[4:0]
  qz[ 87: 80] = (qx[ 87: 80] * qy[ 87: 80]) >> SAR[4:0]
  qz[ 95: 88] = (qx[ 95: 88] * qy[ 95: 88]) >> SAR[4:0]
  qz[103: 96] = (qx[103: 96] * qy[103: 96]) >> SAR[4:0]
  qz[111:104] = (qx[111:104] * qy[111:104]) >> SAR[4:0]
  qz[119:112] = (qx[119:112] * qy[119:112]) >> SAR[4:0]
  qz[127:120] = (qx[127:120] * qy[127:120]) >> SAR[4:0]

  qu[127:0] = load128(rs1[31:0])
  rs1[31:0] = rs1[31:0] + 16
\end{lstlisting}


\textbf{Schedule}\\
\begin{longtable}[c]{|l|l|l|}
    \hline
    \rowcolor{lightgray}
    \textbf{Operands} & \textbf{Use} & \textbf{Def} \\\hline
    \endfirsthead
    \hline
    \rowcolor{lightgray}
     \textbf{Operands} & \textbf{Use} & \textbf{Def} \\\hline
    \endhead
    \hline
    \endfoot

    qx & 1 & - \\\hline
  qy & 1 & - \\\hline
     qz & - & 2 \\\hline
    SAR & 1 & - \\\hline
    qu & - & 2 \\\hline
    rs1 & 1 & 1

\end{longtable}

\textbf{Format}\\

\begin{figure*}[!htbp]
 {\setlength{\tabcolsep}{1pt}
 \begin{center}
\begin{tabular}{@{}F@{}F@{}F@{}I@{}F@{}I@{}F@{}W@{}F@{}F@{}O}
\instbitrange{31}{29}&
\instbitrange{28}{26}&
\instbitrange{25}{23}&
\instbit{22}&
\instbitrange{21}{19}&
\instbit{18}&
\instbitrange{17}{15}&
\instbitrange{14}{13}&
\instbitrange{12}{10}&
\instbitrange{9}{7}&
\instbitrange{6}{0} \\
\hline
\multicolumn{1}{|c|}{{\scriptsize qx[2:0]}} &
\multicolumn{1}{c|}{{\scriptsize qy[2:0]}} &
\multicolumn{1}{c|}{001} &
\multicolumn{1}{c|}{{\scriptsize sat[0]}} &
\multicolumn{1}{c|}{{\scriptsize rm[2:0]}} &
\multicolumn{1}{c|}{{\scriptsize rs1[4]}} &
\multicolumn{1}{c|}{{\scriptsize rs1[2:0]}} &
\multicolumn{1}{c|}{01} &
\multicolumn{1}{c|}{{\scriptsize qu[2:0]}} &
\multicolumn{1}{c|}{{\scriptsize qz[2:0]}} &
\multicolumn{1}{c|}{1011111}\\
\hline
3& 3& 3& 1& 3& 1& 3& 2& 3& 3& 7 \\
\end{tabular}
 \end{center}
 }
 \vspace{-0.1in}
 \caption[]{ESP.VMUL.S8.LD.INCP} \end{figure*}


\newpage
\subsubsection{ESP.VMUL.S8.ST.INCP}\label{instruction:ESPVMULS8STINCP}
\textbf{Syntax}\\
ESP.VMUL.S8.ST.INCP qu,rs1,qz,qx,qy,sat,rm

\textbf{Description}\\
This instruction performs a signed vector multiplication on 8-bit data. Registers \lstinline{qx} and \lstinline{qy} are the multiplier and the multiplicand respectively. The eight 16-bit data results obtained from the calculation are arithmetically right-shifted by the value in the special register \lstinline{SAR}. Then, the lower 8-bit data of the shift result are written into the corresponding segment of the register \lstinline{qz}.\\
During the operation, the value in the register \lstinline{qu} is stored to memory. After the access, the value in the register \lstinline{rs1} is incremented by 16.\\
\textbf{Operation}
\begin{lstlisting}
  qz[  7:  0] = (qx[  7:  0] * qy[  7:  0]) >> SAR[4:0]
  qz[ 15:  8] = (qx[ 15:  8] * qy[ 15:  8]) >> SAR[4:0]
  qz[ 23: 16] = (qx[ 23: 16] * qy[ 23: 16]) >> SAR[4:0]
  qz[ 31: 24] = (qx[ 31: 24] * qy[ 31: 24]) >> SAR[4:0]
  qz[ 39: 32] = (qx[ 39: 32] * qy[ 39: 32]) >> SAR[4:0]
  qz[ 47: 40] = (qx[ 47: 40] * qy[ 47: 40]) >> SAR[4:0]
  qz[ 55: 48] = (qx[ 55: 48] * qy[ 55: 48]) >> SAR[4:0]
  qz[ 63: 56] = (qx[ 63: 56] * qy[ 63: 56]) >> SAR[4:0]
  qz[ 71: 64] = (qx[ 71: 64] * qy[ 71: 64]) >> SAR[4:0]
  qz[ 79: 71] = (qx[ 79: 72] * qy[ 79: 72]) >> SAR[4:0]
  qz[ 87: 80] = (qx[ 87: 80] * qy[ 87: 80]) >> SAR[4:0]
  qz[ 95: 88] = (qx[ 95: 88] * qy[ 95: 88]) >> SAR[4:0]
  qz[103: 96] = (qx[103: 96] * qy[103: 96]) >> SAR[4:0]
  qz[111:104] = (qx[111:104] * qy[111:104]) >> SAR[4:0]
  qz[119:112] = (qx[119:112] * qy[119:112]) >> SAR[4:0]
  qz[127:120] = (qx[127:120] * qy[127:120]) >> SAR[4:0]

  qu[127:0] => store128(rs1[31:0])
  rs1[31:0] = rs1[31:0] + 16
\end{lstlisting}


\textbf{Schedule}\\
\begin{longtable}[c]{|l|l|l|}
    \hline
    \rowcolor{lightgray}
    \textbf{Operands} & \textbf{Use} & \textbf{Def} \\\hline
    \endfirsthead
    \hline
    \rowcolor{lightgray}
     \textbf{Operands} & \textbf{Use} & \textbf{Def} \\\hline
    \endhead
    \hline
    \endfoot

    qx & 1 & - \\\hline
  qy & 1 & - \\\hline
     qz & - & 2 \\\hline
    SAR & 1 & - \\\hline
    qu & 2 & - \\\hline
    rs1 & 1 & 1

\end{longtable}

\textbf{Format}\\

\begin{figure*}[!htbp]
 {\setlength{\tabcolsep}{1pt}
 \begin{center}
\begin{tabular}{@{}F@{}F@{}F@{}I@{}F@{}I@{}F@{}W@{}F@{}F@{}O}
\instbitrange{31}{29}&
\instbitrange{28}{26}&
\instbitrange{25}{23}&
\instbit{22}&
\instbitrange{21}{19}&
\instbit{18}&
\instbitrange{17}{15}&
\instbitrange{14}{13}&
\instbitrange{12}{10}&
\instbitrange{9}{7}&
\instbitrange{6}{0} \\
\hline
\multicolumn{1}{|c|}{{\scriptsize qx[2:0]}} &
\multicolumn{1}{c|}{{\scriptsize qy[2:0]}} &
\multicolumn{1}{c|}{101} &
\multicolumn{1}{c|}{{\scriptsize sat[0]}} &
\multicolumn{1}{c|}{{\scriptsize rm[2:0]}} &
\multicolumn{1}{c|}{{\scriptsize rs1[4]}} &
\multicolumn{1}{c|}{{\scriptsize rs1[2:0]}} &
\multicolumn{1}{c|}{01} &
\multicolumn{1}{c|}{{\scriptsize qu[2:0]}} &
\multicolumn{1}{c|}{{\scriptsize qz[2:0]}} &
\multicolumn{1}{c|}{1011111}\\
\hline
3& 3& 3& 1& 3& 1& 3& 2& 3& 3& 7 \\
\end{tabular}
 \end{center}
 }
 \vspace{-0.1in}
 \caption[]{ESP.VMUL.S8.ST.INCP} \end{figure*}


\newpage
\subsubsection{ESP.VMUL.U16}\label{instruction:ESPVMULU16}
\textbf{Syntax}\\
ESP.VMUL.U16 qz,qx,qy,sat,rm

\textbf{Description}\\
This instruction performs an unsigned vector multiplication on 16-bit data. Registers \lstinline{qx} and \lstinline{qy} are the multiplier and the multiplicand respectively. The eight 32-bit data results obtained from the calculation are logically right-shifted by the value in the special register \lstinline{SAR}. Then, the lower 16-bit data of the shift result are written into the corresponding segment of the register \lstinline{qz}.\\

\textbf{Operation}\\
\begin{lstlisting}
  qz[ 15:  0] = (qx[ 15:  0] * qy[ 15:  0]) >> SAR[5:0]
  qz[ 31: 16] = (qx[ 31: 16] * qy[ 31: 16]) >> SAR[5:0]
  qz[ 47: 32] = (qx[ 47: 32] * qy[ 47: 32]) >> SAR[5:0]
  qz[ 63: 48] = (qx[ 63: 48] * qy[ 63: 48]) >> SAR[5:0]
  qz[ 79: 64] = (qx[ 79: 64] * qy[ 79: 64]) >> SAR[5:0]
  qz[ 95: 80] = (qx[ 95: 80] * qy[ 95: 80]) >> SAR[5:0]
  qz[111: 96] = (qx[111: 96] * qy[111: 96]) >> SAR[5:0]
  qz[127:112] = (qx[127:112] * qy[127:112]) >> SAR[5:0]
\end{lstlisting}


\textbf{Schedule}\\
\begin{longtable}[c]{|l|l|l|}
    \hline
    \rowcolor{lightgray}
    \textbf{Operands} & \textbf{Use} & \textbf{Def} \\\hline
    \endfirsthead
    \hline
    \rowcolor{lightgray}
     \textbf{Operands} & \textbf{Use} & \textbf{Def} \\\hline
    \endhead
    \hline
    \endfoot

    qx & 1 & - \\\hline
  qy & 1 & - \\\hline
     qz & - & 2 \\\hline
    SAR & 1 & -

\end{longtable}

\textbf{Format}\\

\begin{figure*}[!htbp]
 {\setlength{\tabcolsep}{1pt}
 \begin{center}
\begin{tabular}{@{}F@{}F@{}T@{}I@{}I@{}W@{}W@{}I@{}F@{}O}
\instbitrange{31}{29}&
\instbitrange{28}{26}&
\instbitrange{25}{17}&
\instbit{16}&
\instbit{15}&
\instbitrange{14}{13}&
\instbitrange{12}{11}&
\instbit{10}&
\instbitrange{9}{7}&
\instbitrange{6}{0} \\
\hline
\multicolumn{1}{|c|}{{\scriptsize qx[2:0]}} &
\multicolumn{1}{c|}{{\scriptsize qy[2:0]}} &
\multicolumn{1}{c|}{111100X10} &
\multicolumn{1}{c|}{{\scriptsize sat[0]}} &
\multicolumn{1}{c|}{{\scriptsize rm[2]}} &
\multicolumn{1}{c|}{00} &
\multicolumn{1}{c|}{{\scriptsize rm[1:0]}} &
\multicolumn{1}{c|}{0} &
\multicolumn{1}{c|}{{\scriptsize qz[2:0]}} &
\multicolumn{1}{c|}{0011111}\\
\hline
3& 3& 9& 1& 1& 2& 2& 1& 3& 7 \\
\end{tabular}
 \end{center}
 }
 \vspace{-0.1in}
 \caption[]{ESP.VMUL.U16} \end{figure*}


\newpage
\subsubsection{ESP.VMUL.U16.LD.INCP}\label{instruction:ESPVMULU16LDINCP}
\textbf{Syntax}\\
ESP.VMUL.U16.LD.INCP qu,rs1,qz,qx,qy,sat,rm

\textbf{Description}\\
This instruction performs an unsigned vector multiplication on 16-bit data. Registers \lstinline{qx} and \lstinline{qy} are the multiplier and the multiplicand respectively. The eight 32-bit data results obtained from the calculation are logically right-shifted by the value in the special register \lstinline{SAR}. Then, the lower 16-bit data of the shift result are written into the corresponding segment of the register \lstinline{qz}. \\ During the operation, the 16-byte data are loaded from the memory to the register \lstinline{qu}. After the access, the value in the register \lstinline{rs1} is incremented by 16.

\textbf{Operation}\\
\begin{lstlisting}
  qz[ 15:  0] = (qx[ 15:  0] * qy[ 15:  0]) >> SAR[5:0]
  qz[ 31: 16] = (qx[ 31: 16] * qy[ 31: 16]) >> SAR[5:0]
  qz[ 47: 32] = (qx[ 47: 32] * qy[ 47: 32]) >> SAR[5:0]
  qz[ 63: 48] = (qx[ 63: 48] * qy[ 63: 48]) >> SAR[5:0]
  qz[ 79: 64] = (qx[ 79: 64] * qy[ 79: 64]) >> SAR[5:0]
  qz[ 95: 80] = (qx[ 95: 80] * qy[ 95: 80]) >> SAR[5:0]
  qz[111: 96] = (qx[111: 96] * qy[111: 96]) >> SAR[5:0]
  qz[127:112] = (qx[127:112] * qy[127:112]) >> SAR[5:0]

  qu[127:0] = load128(rs1[31:0])
  rs1[31:0] = rs1[31:0] + 16
\end{lstlisting}


\textbf{Schedule}\\
\begin{longtable}[c]{|l|l|l|}
    \hline
    \rowcolor{lightgray}
    \textbf{Operands} & \textbf{Use} & \textbf{Def} \\\hline
    \endfirsthead
    \hline
    \rowcolor{lightgray}
     \textbf{Operands} & \textbf{Use} & \textbf{Def} \\\hline
    \endhead
    \hline
    \endfoot

    qx & 1 & - \\\hline
  qy & 1 & - \\\hline
     qz & - & 2 \\\hline
    SAR & 1 & - \\\hline
    qu & - & 2 \\\hline
    rs1 & 1 & 1

\end{longtable}

\textbf{Format}\\

\begin{figure*}[!htbp]
 {\setlength{\tabcolsep}{1pt}
 \begin{center}
\begin{tabular}{@{}F@{}F@{}F@{}I@{}F@{}I@{}F@{}W@{}F@{}F@{}O}
\instbitrange{31}{29}&
\instbitrange{28}{26}&
\instbitrange{25}{23}&
\instbit{22}&
\instbitrange{21}{19}&
\instbit{18}&
\instbitrange{17}{15}&
\instbitrange{14}{13}&
\instbitrange{12}{10}&
\instbitrange{9}{7}&
\instbitrange{6}{0} \\
\hline
\multicolumn{1}{|c|}{{\scriptsize qx[2:0]}} &
\multicolumn{1}{c|}{{\scriptsize qy[2:0]}} &
\multicolumn{1}{c|}{010} &
\multicolumn{1}{c|}{{\scriptsize sat[0]}} &
\multicolumn{1}{c|}{{\scriptsize rm[2:0]}} &
\multicolumn{1}{c|}{{\scriptsize rs1[4]}} &
\multicolumn{1}{c|}{{\scriptsize rs1[2:0]}} &
\multicolumn{1}{c|}{01} &
\multicolumn{1}{c|}{{\scriptsize qu[2:0]}} &
\multicolumn{1}{c|}{{\scriptsize qz[2:0]}} &
\multicolumn{1}{c|}{1011111}\\
\hline
3& 3& 3& 1& 3& 1& 3& 2& 3& 3& 7 \\
\end{tabular}
 \end{center}
 }
 \vspace{-0.1in}
 \caption[]{ESP.VMUL.U16.LD.INCP} \end{figure*}


\newpage
\subsubsection{ESP.VMUL.U16.ST.INCP}\label{instruction:ESPVMULU16STINCP}
\textbf{Syntax}\\
ESP.VMUL.U16.ST.INCP qu,rs1,qz,qx,qy,sat,rm

\textbf{Description}\\
This instruction performs an unsigned vector multiplication on 16-bit data. Registers \lstinline{qx} and \lstinline{qy} are the multiplier and the multiplicand respectively. The eight 32-bit data results obtained from the calculation are logically right-shifted by the value in the special register \lstinline{SAR}. Then, the lower 16-bit data of the shift result are written into the corresponding segment of the register \lstinline{qz}. \\ During the operation, the value in the register \lstinline{qu} is stored to memory. After the access, the value in the register \lstinline{rs1} is incremented by 16.

\textbf{Operation}\\
\begin{lstlisting}
  qz[ 15:  0] = (qx[ 15:  0] * qy[ 15:  0]) >> SAR[5:0]
  qz[ 31: 16] = (qx[ 31: 16] * qy[ 31: 16]) >> SAR[5:0]
  qz[ 47: 32] = (qx[ 47: 32] * qy[ 47: 32]) >> SAR[5:0]
  qz[ 63: 48] = (qx[ 63: 48] * qy[ 63: 48]) >> SAR[5:0]
  qz[ 79: 64] = (qx[ 79: 64] * qy[ 79: 64]) >> SAR[5:0]
  qz[ 95: 80] = (qx[ 95: 80] * qy[ 95: 80]) >> SAR[5:0]
  qz[111: 96] = (qx[111: 96] * qy[111: 96]) >> SAR[5:0]
  qz[127:112] = (qx[127:112] * qy[127:112]) >> SAR[5:0]

  qu[127:0] => store128(rs1[31:0])
  rs1[31:0] = rs1[31:0] + 16
\end{lstlisting}


\textbf{Schedule}\\
\begin{longtable}[c]{|l|l|l|}
    \hline
    \rowcolor{lightgray}
    \textbf{Operands} & \textbf{Use} & \textbf{Def} \\\hline
    \endfirsthead
    \hline
    \rowcolor{lightgray}
     \textbf{Operands} & \textbf{Use} & \textbf{Def} \\\hline
    \endhead
    \hline
    \endfoot

    qx & 1 & - \\\hline
  qy & 1 & - \\\hline
     qz & - & 2 \\\hline
    SAR & 1 & - \\\hline
    qu & 2 & - \\\hline
    rs1 & 1 & 1

\end{longtable}

\textbf{Format}\\

\begin{figure*}[!htbp]
 {\setlength{\tabcolsep}{1pt}
 \begin{center}
\begin{tabular}{@{}F@{}F@{}F@{}I@{}F@{}I@{}F@{}W@{}F@{}F@{}O}
\instbitrange{31}{29}&
\instbitrange{28}{26}&
\instbitrange{25}{23}&
\instbit{22}&
\instbitrange{21}{19}&
\instbit{18}&
\instbitrange{17}{15}&
\instbitrange{14}{13}&
\instbitrange{12}{10}&
\instbitrange{9}{7}&
\instbitrange{6}{0} \\
\hline
\multicolumn{1}{|c|}{{\scriptsize qx[2:0]}} &
\multicolumn{1}{c|}{{\scriptsize qy[2:0]}} &
\multicolumn{1}{c|}{110} &
\multicolumn{1}{c|}{{\scriptsize sat[0]}} &
\multicolumn{1}{c|}{{\scriptsize rm[2:0]}} &
\multicolumn{1}{c|}{{\scriptsize rs1[4]}} &
\multicolumn{1}{c|}{{\scriptsize rs1[2:0]}} &
\multicolumn{1}{c|}{01} &
\multicolumn{1}{c|}{{\scriptsize qu[2:0]}} &
\multicolumn{1}{c|}{{\scriptsize qz[2:0]}} &
\multicolumn{1}{c|}{1011111}\\
\hline
3& 3& 3& 1& 3& 1& 3& 2& 3& 3& 7 \\
\end{tabular}
 \end{center}
 }
 \vspace{-0.1in}
 \caption[]{ESP.VMUL.U16.ST.INCP} \end{figure*}


\newpage
\subsubsection{ESP.VMUL.U8}\label{instruction:ESPVMULU8}
\textbf{Syntax}\\
ESP.VMUL.U8 qz,qx,qy,sat,rm

\textbf{Description}\\
This instruction performs an unsigned vector multiplication on 8-bit data. Registers \lstinline{qx} and \lstinline{qy} are the multiplier and the multiplicand respectively. The 16 16-bit data results obtained from the calculation are logically right-shifted by the value in the special register \lstinline{SAR}. Then, the lower 8-bit data of the shift result are written into the corresponding segment of the register \lstinline{qz}.

\textbf{Operation}\\
\begin{lstlisting}
  qz[  7:  0] = (qx[  7:  0] * qy[  7:  0]) >> SAR[5:0]
  qz[ 15:  8] = (qx[ 15:  8] * qy[ 15:  8]) >> SAR[5:0]
  qz[ 23: 16] = (qx[ 23: 16] * qy[ 23: 16]) >> SAR[5:0]
  qz[ 31: 24] = (qx[ 31: 24] * qy[ 31: 24]) >> SAR[5:0]
  qz[ 39: 32] = (qx[ 39: 32] * qy[ 39: 32]) >> SAR[5:0]
  qz[ 47: 40] = (qx[ 47: 40] * qy[ 47: 40]) >> SAR[5:0]
  qz[ 55: 48] = (qx[ 55: 48] * qy[ 55: 48]) >> SAR[5:0]
  qz[ 63: 56] = (qx[ 63: 56] * qy[ 63: 56]) >> SAR[5:0]
  qz[ 71: 64] = (qx[ 71: 64] * qy[ 71: 64]) >> SAR[5:0]
  qz[ 79: 72] = (qx[ 79: 72] * qy[ 79: 72]) >> SAR[5:0]
  qz[ 87: 80] = (qx[ 87: 80] * qy[ 87: 80]) >> SAR[5:0]
  qz[ 95: 88] = (qx[ 95: 88] * qy[ 95: 88]) >> SAR[5:0]
  qz[103: 96] = (qx[103: 96] * qy[103: 96]) >> SAR[5:0]
  qz[111:104] = (qx[111:104] * qy[111:104]) >> SAR[5:0]
  qz[119:112] = (qx[119:112] * qy[119:112]) >> SAR[5:0]
  qz[127:120] = (qx[127:120] * qy[127:120]) >> SAR[5:0]
\end{lstlisting}


\textbf{Schedule}\\
\begin{longtable}[c]{|l|l|l|}
    \hline
    \rowcolor{lightgray}
    \textbf{Operands} & \textbf{Use} & \textbf{Def} \\\hline
    \endfirsthead
    \hline
    \rowcolor{lightgray}
     \textbf{Operands} & \textbf{Use} & \textbf{Def} \\\hline
    \endhead
    \hline
    \endfoot

    qx & 1 & - \\\hline
  qy & 1 & - \\\hline
     qz & - & 2 \\\hline
    SAR & 1 & - \\\hline
    qu & 2 & - \\\hline
    rs1 & 1 & 1

\end{longtable}

\textbf{Format}\\

\begin{figure*}[!htbp]
 {\setlength{\tabcolsep}{1pt}
 \begin{center}
\begin{tabular}{@{}F@{}F@{}T@{}I@{}I@{}W@{}W@{}I@{}F@{}O}
\instbitrange{31}{29}&
\instbitrange{28}{26}&
\instbitrange{25}{17}&
\instbit{16}&
\instbit{15}&
\instbitrange{14}{13}&
\instbitrange{12}{11}&
\instbit{10}&
\instbitrange{9}{7}&
\instbitrange{6}{0} \\
\hline
\multicolumn{1}{|c|}{{\scriptsize qx[2:0]}} &
\multicolumn{1}{c|}{{\scriptsize qy[2:0]}} &
\multicolumn{1}{c|}{111100X00} &
\multicolumn{1}{c|}{{\scriptsize sat[0]}} &
\multicolumn{1}{c|}{{\scriptsize rm[2]}} &
\multicolumn{1}{c|}{00} &
\multicolumn{1}{c|}{{\scriptsize rm[1:0]}} &
\multicolumn{1}{c|}{0} &
\multicolumn{1}{c|}{{\scriptsize qz[2:0]}} &
\multicolumn{1}{c|}{0011111}\\
\hline
3& 3& 9& 1& 1& 2& 2& 1& 3& 7 \\
\end{tabular}
 \end{center}
 }
 \vspace{-0.1in}
 \caption[]{ESP.VMUL.U8} \end{figure*}


\newpage
\subsubsection{ESP.VMUL.U8.LD.INCP}\label{instruction:ESPVMULU8LDINCP}
\textbf{Syntax}\\
ESP.VMUL.U8.LD.INCP qu,rs1,qz,qx,qy,sat,rm

\textbf{Description}\\
This instruction performs an unsigned vector multiplication on 8-bit data. Registers \lstinline{qx} and \lstinline{qy} are the multiplier and the multiplicand respectively. The 16 32-bit data results obtained from the calculation are logically right-shifted by the value in the special register \lstinline{SAR}. Then, the lower 8-bit data of the shift result are written into the corresponding segment of the register \lstinline{qz}. \\ During the operation, the 16-byte data are loaded from the memory to the register \lstinline{qu}. After the access, the value in the register \lstinline{rs1} is incremented by 16.

\textbf{Operation}\\
\begin{lstlisting}
  qz[  7:  0] = (qx[  7:  0] * qy[  7:  0]) >> SAR[5:0]
  qz[ 15:  8] = (qx[ 15:  8] * qy[ 15:  8]) >> SAR[5:0]
  qz[ 23: 16] = (qx[ 23: 16] * qy[ 23: 16]) >> SAR[5:0]
  qz[ 31: 24] = (qx[ 31: 24] * qy[ 31: 24]) >> SAR[5:0]
  qz[ 39: 32] = (qx[ 39: 32] * qy[ 39: 32]) >> SAR[5:0]
  qz[ 47: 40] = (qx[ 47: 40] * qy[ 47: 40]) >> SAR[5:0]
  qz[ 55: 48] = (qx[ 55: 48] * qy[ 55: 48]) >> SAR[5:0]
  qz[ 63: 56] = (qx[ 63: 56] * qy[ 63: 56]) >> SAR[5:0]
  qz[ 71: 64] = (qx[ 71: 64] * qy[ 71: 64]) >> SAR[5:0]
  qz[ 79: 72] = (qx[ 79: 72] * qy[ 79: 72]) >> SAR[5:0]
  qz[ 87: 80] = (qx[ 87: 80] * qy[ 87: 80]) >> SAR[5:0]
  qz[ 95: 88] = (qx[ 95: 88] * qy[ 95: 88]) >> SAR[5:0]
  qz[103: 96] = (qx[103: 96] * qy[103: 96]) >> SAR[5:0]
  qz[111:104] = (qx[111:104] * qy[111:104]) >> SAR[5:0]
  qz[119:112] = (qx[119:112] * qy[119:112]) >> SAR[5:0]
  qz[127:120] = (qx[127:120] * qy[127:120]) >> SAR[5:0]

  qu[127:0] = load128(rs1[31:0])
  rs1[31:0] = rs1[31:0] + 16
\end{lstlisting}


\textbf{Schedule}\\
\begin{longtable}[c]{|l|l|l|}
    \hline
    \rowcolor{lightgray}
    \textbf{Operands} & \textbf{Use} & \textbf{Def} \\\hline
    \endfirsthead
    \hline
    \rowcolor{lightgray}
     \textbf{Operands} & \textbf{Use} & \textbf{Def} \\\hline
    \endhead
    \hline
    \endfoot

    qx & 1 & - \\\hline
  qy & 1 & - \\\hline
     qz & - & 2 \\\hline
    SAR & 1 & - \\\hline
    qu & - & 2 \\\hline
    rs1 & 1 & 1

\end{longtable}

\textbf{Format}\\

\begin{figure*}[!htbp]
 {\setlength{\tabcolsep}{1pt}
 \begin{center}
\begin{tabular}{@{}F@{}F@{}F@{}I@{}F@{}I@{}F@{}W@{}F@{}F@{}O}
\instbitrange{31}{29}&
\instbitrange{28}{26}&
\instbitrange{25}{23}&
\instbit{22}&
\instbitrange{21}{19}&
\instbit{18}&
\instbitrange{17}{15}&
\instbitrange{14}{13}&
\instbitrange{12}{10}&
\instbitrange{9}{7}&
\instbitrange{6}{0} \\
\hline
\multicolumn{1}{|c|}{{\scriptsize qx[2:0]}} &
\multicolumn{1}{c|}{{\scriptsize qy[2:0]}} &
\multicolumn{1}{c|}{000} &
\multicolumn{1}{c|}{{\scriptsize sat[0]}} &
\multicolumn{1}{c|}{{\scriptsize rm[2:0]}} &
\multicolumn{1}{c|}{{\scriptsize rs1[4]}} &
\multicolumn{1}{c|}{{\scriptsize rs1[2:0]}} &
\multicolumn{1}{c|}{01} &
\multicolumn{1}{c|}{{\scriptsize qu[2:0]}} &
\multicolumn{1}{c|}{{\scriptsize qz[2:0]}} &
\multicolumn{1}{c|}{1011111}\\
\hline
3& 3& 3& 1& 3& 1& 3& 2& 3& 3& 7 \\
\end{tabular}
 \end{center}
 }
 \vspace{-0.1in}
 \caption[]{ESP.VMUL.U8.LD.INCP} \end{figure*}


\newpage
\subsubsection{ESP.VMUL.U8.ST.INCP}\label{instruction:ESPVMULU8STINCP}
\textbf{Syntax}\\
ESP.VMUL.U8.ST.INCP qu,rs1,qz,qx,qy,sat,rm

\textbf{Description}\\
This instruction performs an unsigned vector multiplication on 8-bit data. Registers \lstinline{qx} and \lstinline{qy} are the multiplier and the multiplicand respectively. The 16 16-bit data results obtained from the calculation are logically right-shifted by the value in the special register \lstinline{SAR}. Then, the lower 8-bit data of the shift result are written into the corresponding segment of the register \lstinline{qz}. \\ During the operation, the value in the register \lstinline{qu} is stored to memory. After the access, the value in the register \lstinline{rs1} is incremented by 16.

\textbf{Operation}\\
\begin{lstlisting}
  qz[  7:  0] = (qx[  7:  0] * qy[  7:  0]) >> SAR[5:0]
  qz[ 15:  8] = (qx[ 15:  8] * qy[ 15:  8]) >> SAR[5:0]
  qz[ 23: 16] = (qx[ 23: 16] * qy[ 23: 16]) >> SAR[5:0]
  qz[ 31: 24] = (qx[ 31: 24] * qy[ 31: 24]) >> SAR[5:0]
  qz[ 39: 32] = (qx[ 39: 32] * qy[ 39: 32]) >> SAR[5:0]
  qz[ 47: 40] = (qx[ 47: 40] * qy[ 47: 40]) >> SAR[5:0]
  qz[ 55: 48] = (qx[ 55: 48] * qy[ 55: 48]) >> SAR[5:0]
  qz[ 63: 56] = (qx[ 63: 56] * qy[ 63: 56]) >> SAR[5:0]
  qz[ 71: 64] = (qx[ 71: 64] * qy[ 71: 64]) >> SAR[5:0]
  qz[ 79: 72] = (qx[ 79: 72] * qy[ 79: 72]) >> SAR[5:0]
  qz[ 87: 80] = (qx[ 87: 80] * qy[ 87: 80]) >> SAR[5:0]
  qz[ 95: 88] = (qx[ 95: 88] * qy[ 95: 88]) >> SAR[5:0]
  qz[103: 96] = (qx[103: 96] * qy[103: 96]) >> SAR[5:0]
  qz[111:104] = (qx[111:104] * qy[111:104]) >> SAR[5:0]
  qz[119:112] = (qx[119:112] * qy[119:112]) >> SAR[5:0]
  qz[127:120] = (qx[127:120] * qy[127:120]) >> SAR[5:0]

  qu[127:0] => store128(rs1[31:0])
  rs1[31:0] = rs1[31:0] + 16
\end{lstlisting}


\textbf{Schedule}\\
\begin{longtable}[c]{|l|l|l|}
    \hline
    \rowcolor{lightgray}
    \textbf{Operands} & \textbf{Use} & \textbf{Def} \\\hline
    \endfirsthead
    \hline
    \rowcolor{lightgray}
     \textbf{Operands} & \textbf{Use} & \textbf{Def} \\\hline
    \endhead
    \hline
    \endfoot

    qx & 1 & - \\\hline
  qy & 1 & - \\\hline
     qz & - & 2 \\\hline
    SAR & 1 & - \\\hline
    qu & 2 & - \\\hline
    rs1 & 1 & 1

\end{longtable}

\textbf{Format}\\

\begin{figure*}[!htbp]
 {\setlength{\tabcolsep}{1pt}
 \begin{center}
\begin{tabular}{@{}F@{}F@{}F@{}I@{}F@{}I@{}F@{}W@{}F@{}F@{}O}
\instbitrange{31}{29}&
\instbitrange{28}{26}&
\instbitrange{25}{23}&
\instbit{22}&
\instbitrange{21}{19}&
\instbit{18}&
\instbitrange{17}{15}&
\instbitrange{14}{13}&
\instbitrange{12}{10}&
\instbitrange{9}{7}&
\instbitrange{6}{0} \\
\hline
\multicolumn{1}{|c|}{{\scriptsize qx[2:0]}} &
\multicolumn{1}{c|}{{\scriptsize qy[2:0]}} &
\multicolumn{1}{c|}{100} &
\multicolumn{1}{c|}{{\scriptsize sat[0]}} &
\multicolumn{1}{c|}{{\scriptsize rm[2:0]}} &
\multicolumn{1}{c|}{{\scriptsize rs1[4]}} &
\multicolumn{1}{c|}{{\scriptsize rs1[2:0]}} &
\multicolumn{1}{c|}{01} &
\multicolumn{1}{c|}{{\scriptsize qu[2:0]}} &
\multicolumn{1}{c|}{{\scriptsize qz[2:0]}} &
\multicolumn{1}{c|}{1011111}\\
\hline
3& 3& 3& 1& 3& 1& 3& 2& 3& 3& 7 \\
\end{tabular}
 \end{center}
 }
 \vspace{-0.1in}
 \caption[]{ESP.VMUL.U8.ST.INCP} \end{figure*}


\newpage
\subsubsection{ESP.VPRELU.S16}\label{instruction:ESPVPRELUS16}
\textbf{Syntax}\\
ESP.VPRELU.S16 qz,qy,qx,rs1,sat,rm

\textbf{Description}\\
This instruction divides the register \lstinline{qx} into 8 data segments of 16 bits each. If the value of a segment is not greater than 0, it performs multiplication with the value of the corresponding 16-bit segment in the register \lstinline{qy}, right-shifted by the lower 6-bit value in the register \lstinline{rs1}. The result obtained is then assigned to the corresponding 16-bit data segment in the register \lstinline{qz}. Otherwise, the value of the segment in \lstinline{qx} is assigned to \lstinline{qz}. The data is signed. \\
\textbf{Operation}
\begin{lstlisting}[escapeinside=`?]
  qz[ 15:  0] = (qx[ 15:  0]<=0) ? (qx[ 15:  0] * qy[ 15:  0]) >> rs1[5:0] : qx[ 15:  0]
  qz[ 31: 16] = (qx[ 31: 16]<=0) ? (qx[ 31: 16] * qy[ 31: 16]) >> rs1[5:0] : qx[ 31: 16]
  ...
  qz[127:112] = (qx[127:112]<=0) ? (qx[127:112] * qy[127:112]) >> rs1[5:0] : qx[127:112]
\end{lstlisting}


\textbf{Schedule}\\
\begin{longtable}[c]{|l|l|l|}
    \hline
    \rowcolor{lightgray}
    \textbf{Operands} & \textbf{Use} & \textbf{Def} \\\hline
    \endfirsthead
    \hline
    \rowcolor{lightgray}
     \textbf{Operands} & \textbf{Use} & \textbf{Def} \\\hline
    \endhead
    \hline
    \endfoot

    qx & 1 & - \\\hline
  qy & 1 & - \\\hline
     qz & - & 2 \\\hline
    rs1 & 1 & -

\end{longtable}

\textbf{Format}\\

\begin{figure*}[!htbp]
 {\setlength{\tabcolsep}{1pt}
 \begin{center}
\begin{tabular}{@{}F@{}F@{}F@{}I@{}F@{}I@{}F@{}W@{}F@{}F@{}O}
\instbitrange{31}{29}&
\instbitrange{28}{26}&
\instbitrange{25}{23}&
\instbit{22}&
\instbitrange{21}{19}&
\instbit{18}&
\instbitrange{17}{15}&
\instbitrange{14}{13}&
\instbitrange{12}{10}&
\instbitrange{9}{7}&
\instbitrange{6}{0} \\
\hline
\multicolumn{1}{|c|}{{\scriptsize qx[2:0]}} &
\multicolumn{1}{c|}{{\scriptsize qy[2:0]}} &
\multicolumn{1}{c|}{111} &
\multicolumn{1}{c|}{{\scriptsize sat[0]}} &
\multicolumn{1}{c|}{11X} &
\multicolumn{1}{c|}{{\scriptsize rs1[4]}} &
\multicolumn{1}{c|}{{\scriptsize rs1[2:0]}} &
\multicolumn{1}{c|}{00} &
\multicolumn{1}{c|}{{\scriptsize rm[2:0]}} &
\multicolumn{1}{c|}{{\scriptsize qz[2:0]}} &
\multicolumn{1}{c|}{0011111}\\
\hline
3& 3& 3& 1& 3& 1& 3& 2& 3& 3& 7 \\
\end{tabular}
 \end{center}
 }
 \vspace{-0.1in}
 \caption[]{ESP.VPRELU.S16} \end{figure*}


\newpage
\subsubsection{ESP.VPRELU.S8}\label{instruction:ESPVPRELUS8}
\textbf{Syntax}\\
ESP.VPRELU.S8 qz,qy,qx,rs1,sat,rm

\textbf{Description}\\
This instruction divides the register \lstinline{qx} into 16 data segments of 8 bits each. If the value of a segment is not greater than 0, it performs multiplication with the value of the corresponding 8-bit segment in the register \lstinline{qy}, right-shifts by the lower 5-bit value in the register \lstinline{rs1}, and then assigns the calculated result to the corresponding 8-bit data segment in the register \lstinline{qz}. Otherwise, the value of the segment in \lstinline{qx} is assigned to \lstinline{qz}.\\
\textbf{Operation}
\begin{lstlisting}[escapeinside=`?]
  qz[  7:  0] = (qx[  7:  0]<=0) ? (qx[  7:  0] * qy[  7:  0]) >> rs1[4:0] : qx[  7:  0]
  qz[ 15:  8] = (qx[ 15:  8]<=0) ? (qx[ 15:  8] * qy[ 15:  8]) >> rs1[4:0] : qx[ 15:  8]
  ...
  qz[127:120] = (qx[127:120]<=0) ? (qx[127:120] * qy[127:120]) >> rs1[4:0] : qx[127:120]
\end{lstlisting}


\textbf{Schedule}\\
\begin{longtable}[c]{|l|l|l|}
    \hline
    \rowcolor{lightgray}
    \textbf{Operands} & \textbf{Use} & \textbf{Def} \\\hline
    \endfirsthead
    \hline
    \rowcolor{lightgray}
     \textbf{Operands} & \textbf{Use} & \textbf{Def} \\\hline
    \endhead
    \hline
    \endfoot

    qx & 1 & - \\\hline
  qy & 1 & - \\\hline
     qz & - & 2 \\\hline
    rs1 & 1 & -

\end{longtable}

\textbf{Format}\\

\begin{figure*}[!htbp]
 {\setlength{\tabcolsep}{1pt}
 \begin{center}
\begin{tabular}{@{}F@{}F@{}F@{}I@{}F@{}I@{}F@{}W@{}F@{}F@{}O}
\instbitrange{31}{29}&
\instbitrange{28}{26}&
\instbitrange{25}{23}&
\instbit{22}&
\instbitrange{21}{19}&
\instbit{18}&
\instbitrange{17}{15}&
\instbitrange{14}{13}&
\instbitrange{12}{10}&
\instbitrange{9}{7}&
\instbitrange{6}{0} \\
\hline
\multicolumn{1}{|c|}{{\scriptsize qx[2:0]}} &
\multicolumn{1}{c|}{{\scriptsize qy[2:0]}} &
\multicolumn{1}{c|}{110} &
\multicolumn{1}{c|}{{\scriptsize sat[0]}} &
\multicolumn{1}{c|}{11X} &
\multicolumn{1}{c|}{{\scriptsize rs1[4]}} &
\multicolumn{1}{c|}{{\scriptsize rs1[2:0]}} &
\multicolumn{1}{c|}{00} &
\multicolumn{1}{c|}{{\scriptsize rm[2:0]}} &
\multicolumn{1}{c|}{{\scriptsize qz[2:0]}} &
\multicolumn{1}{c|}{0011111}\\
\hline
3& 3& 3& 1& 3& 1& 3& 2& 3& 3& 7 \\
\end{tabular}
 \end{center}
 }
 \vspace{-0.1in}
 \caption[]{ESP.VPRELU.S8} \end{figure*}


\newpage
\subsubsection{ESP.VRELU.S16}\label{instruction:ESPVRELUS16}
\textbf{Syntax}\\
ESP.VRELU.S16 qy,rs2,rs1,sat,rm

\textbf{Description}\\
This instruction divides register \lstinline{qy} into 8 data segments by 16 bits. If the value of the segment is not greater than 0, it will be multiplied by the value of the lower 16 bits in register \lstinline{rs2} and right-shifted by the value of the lower 6 bits in register \lstinline{rs1}, and then the result obtained will overwrite the value of the segment. Otherwise, the value of the segment will remain unchanged.\\
\textbf{Operation}
\begin{lstlisting}[escapeinside=`?]
  qy[ 15:  0] = (qy[ 15:  0]<=0) ? (qy[ 15:  0] * rs2[15:0]) >> rs1[5:0] : qy[ 15:  0]
  qy[ 31: 16] = (qy[ 31: 16]<=0) ? (qy[ 31: 16] * rs2[15:0]) >> rs1[5:0] : qy[ 31: 16]
  ...
  qy[127:112] = (qy[127:112]<=0) ? (qy[127:112] * rs2[15:0]) >> rs1[5:0] : qy[127:112]
\end{lstlisting}


\textbf{Schedule}\\
\begin{longtable}[c]{|l|l|l|}
    \hline
    \rowcolor{lightgray}
    \textbf{Operands} & \textbf{Use} & \textbf{Def} \\\hline
    \endfirsthead
    \hline
    \rowcolor{lightgray}
     \textbf{Operands} & \textbf{Use} & \textbf{Def} \\\hline
    \endhead
    \hline
    \endfoot

  qy & 1 & 2 \\\hline
    rs2 & 1 & - \\\hline
    rs1 & 1 & -

\end{longtable}

\textbf{Format}\\

\begin{figure*}[!htbp]
 {\setlength{\tabcolsep}{1pt}
 \begin{center}
\begin{tabular}{@{}F@{}F@{}I@{}I@{}I@{}F@{}I@{}I@{}F@{}W@{}F@{}F@{}O}
\instbitrange{31}{29}&
\instbitrange{28}{26}&
\instbit{25}&
\instbit{24}&
\instbit{23}&
\instbitrange{22}{20}&
\instbit{19}&
\instbit{18}&
\instbitrange{17}{15}&
\instbitrange{14}{13}&
\instbitrange{12}{10}&
\instbitrange{9}{7}&
\instbitrange{6}{0} \\
\hline
\multicolumn{1}{|c|}{11X} &
\multicolumn{1}{c|}{{\scriptsize qy[2:0]}} &
\multicolumn{1}{c|}{1} &
\multicolumn{1}{c|}{{\scriptsize sat[0]}} &
\multicolumn{1}{c|}{{\scriptsize rs2[4]}} &
\multicolumn{1}{c|}{{\scriptsize rs2[2:0]}} &
\multicolumn{1}{c|}{X} &
\multicolumn{1}{c|}{{\scriptsize rs1[4]}} &
\multicolumn{1}{c|}{{\scriptsize rs1[2:0]}} &
\multicolumn{1}{c|}{11} &
\multicolumn{1}{c|}{{\scriptsize rm[2:0]}} &
\multicolumn{1}{c|}{X0X} &
\multicolumn{1}{c|}{0011011}\\
\hline
3& 3& 1& 1& 1& 3& 1& 1& 3& 2& 3& 3& 7 \\
\end{tabular}
 \end{center}
 }
 \vspace{-0.1in}
 \caption[]{ESP.VRELU.S16} \end{figure*}


\newpage
\subsubsection{ESP.VRELU.S8}\label{instruction:ESPVRELUS8}
\textbf{Syntax}\\
ESP.VRELU.S8 qy,rs2,rs1,sat,rm

\textbf{Description}\\
This instruction divides register \lstinline{qy} into 16 data segments by 8 bits. If the value of the segment is not greater than 0, it will be multiplied by the value of the lower 8 bits in register \lstinline{rs2} and right-shifted by the value of the lower 5 bits in register \lstinline{rs1}, and then the result obtained will overwrite the value of the segment. Otherwise, the value of the segment will remain unchanged.\\
\textbf{Operation}
\begin{lstlisting}[escapeinside=`?]
  qy[  7:  0] = (qy[  7:  0]<=0) ? (qy[  7:  0] * rs2[7:0]) >> rs1[4:0] : qy[  7:  0]
  qy[ 15:  8] = (qy[ 15:  8]<=0) ? (qy[ 15:  8] * rs2[7:0]) >> rs1[4:0] : qy[ 15:  8]
  ...
  qy[127:120] = (qy[127:120]<=0) ? (qy[127:120] * rs2[7:0]) >> rs1[4:0] : qy[127:120]
\end{lstlisting}


\textbf{Schedule}\\
\begin{longtable}[c]{|l|l|l|}
    \hline
    \rowcolor{lightgray}
    \textbf{Operands} & \textbf{Use} & \textbf{Def} \\\hline
    \endfirsthead
    \hline
    \rowcolor{lightgray}
     \textbf{Operands} & \textbf{Use} & \textbf{Def} \\\hline
    \endhead
    \hline
    \endfoot

  qy & 1 & 2 \\\hline
    rs2 & 1 & - \\\hline
    rs1 & 1 & -

\end{longtable}

\textbf{Format}\\

\begin{figure*}[!htbp]
 {\setlength{\tabcolsep}{1pt}
 \begin{center}
\begin{tabular}{@{}F@{}F@{}I@{}I@{}I@{}F@{}I@{}I@{}F@{}W@{}F@{}F@{}O}
\instbitrange{31}{29}&
\instbitrange{28}{26}&
\instbit{25}&
\instbit{24}&
\instbit{23}&
\instbitrange{22}{20}&
\instbit{19}&
\instbit{18}&
\instbitrange{17}{15}&
\instbitrange{14}{13}&
\instbitrange{12}{10}&
\instbitrange{9}{7}&
\instbitrange{6}{0} \\
\hline
\multicolumn{1}{|c|}{11X} &
\multicolumn{1}{c|}{{\scriptsize qy[2:0]}} &
\multicolumn{1}{c|}{0} &
\multicolumn{1}{c|}{{\scriptsize sat[0]}} &
\multicolumn{1}{c|}{{\scriptsize rs2[4]}} &
\multicolumn{1}{c|}{{\scriptsize rs2[2:0]}} &
\multicolumn{1}{c|}{X} &
\multicolumn{1}{c|}{{\scriptsize rs1[4]}} &
\multicolumn{1}{c|}{{\scriptsize rs1[2:0]}} &
\multicolumn{1}{c|}{11} &
\multicolumn{1}{c|}{{\scriptsize rm[2:0]}} &
\multicolumn{1}{c|}{X0X} &
\multicolumn{1}{c|}{0011011}\\
\hline
3& 3& 1& 1& 1& 3& 1& 1& 3& 2& 3& 3& 7 \\
\end{tabular}
 \end{center}
 }
 \vspace{-0.1in}
 \caption[]{ESP.VRELU.S8} \end{figure*}


\newpage
\subsubsection{ESP.VSADDS.S16}\label{instruction:ESPVSADDSS16}
\textbf{Syntax}\\
ESP.VSADDS.S16 qv,qx,rs1,sat

\textbf{Description}\\
This instruction performs a signed vector addition on 16-bit data.

\textbf{Operation}\\
\begin{lstlisting}[escapeinside=`?]
    qv[ 15:  0] = min( max( qx[ 15:  0] + rs1[15:0], -2^{15} ), 2^{15}-1 )
    qv[ 31: 16] = min( max( qx[ 31: 16] + rs1[15:0], -2^{15} ), 2^{15}-1 )
    qv[ 47: 32] = min( max( qx[ 47: 32] + rs1[15:0], -2^{15} ), 2^{15}-1 )
    qv[ 63: 48] = min( max( qx[ 63: 48] + rs1[15:0], -2^{15} ), 2^{15}-1 )
    qv[ 79: 64] = min( max( qx[ 79: 64] + rs1[15:0], -2^{15} ), 2^{15}-1 )
    qv[ 95: 80] = min( max( qx[ 95: 80] + rs1[15:0], -2^{15} ), 2^{15}-1 )
    qv[111: 96] = min( max( qx[111: 96] + rs1[15:0], -2^{15} ), 2^{15}-1 )
    qv[127:112] = min( max( qx[127:112] + rs1[15:0], -2^{15} ), 2^{15}-1 )
\end{lstlisting}


\textbf{Schedule}\\
\begin{longtable}[c]{|l|l|l|}
    \hline
    \rowcolor{lightgray}
    \textbf{Operands} & \textbf{Use} & \textbf{Def} \\\hline
    \endfirsthead
    \hline
    \rowcolor{lightgray}
     \textbf{Operands} & \textbf{Use} & \textbf{Def} \\\hline
    \endhead
    \hline
    \endfoot

  qv & - & 2 \\\hline
    qx & 1 & - \\\hline
    rs1 & 1 & -

\end{longtable}

\textbf{Format}\\

\begin{figure*}[!htbp]
 {\setlength{\tabcolsep}{1pt}
 \begin{center}
\begin{tabular}{@{}F@{}W@{}I@{}F@{}F@{}I@{}I@{}F@{}E@{}O}
\instbitrange{31}{29}&
\instbitrange{28}{27}&
\instbit{26}&
\instbitrange{25}{23}&
\instbitrange{22}{20}&
\instbit{19}&
\instbit{18}&
\instbitrange{17}{15}&
\instbitrange{14}{7}&
\instbitrange{6}{0} \\
\hline
\multicolumn{1}{|c|}{{\scriptsize qx[2:0]}} &
\multicolumn{1}{c|}{11} &
\multicolumn{1}{c|}{{\scriptsize sat[0]}} &
\multicolumn{1}{c|}{1X1} &
\multicolumn{1}{c|}{{\scriptsize qv[2:0]}} &
\multicolumn{1}{c|}{X} &
\multicolumn{1}{c|}{{\scriptsize rs1[4]}} &
\multicolumn{1}{c|}{{\scriptsize rs1[2:0]}} &
\multicolumn{1}{c|}{01XX001X} &
\multicolumn{1}{c|}{0011011}\\
\hline
3& 2& 1& 3& 3& 1& 1& 3& 8& 7 \\
\end{tabular}
 \end{center}
 }
 \vspace{-0.1in}
 \caption[]{ESP.VSADDS.S16} \end{figure*}


\newpage
\subsubsection{ESP.VSADDS.S8}\label{instruction:ESPVSADDSS8}
\textbf{Syntax}\\
ESP.VSADDS.S8 qv,qx,rs1,sat

\textbf{Description}\\
This instruction performs a signed vector addition on 8-bit data.

\textbf{Operation}\\
\begin{lstlisting}[escapeinside=`?]
    qv[  7:  0] = min( max( qx[  7:  0] + rs1[7:0], -2^{7} ), 2^{7}-1 )
    qv[ 15:  8] = min( max( qx[ 15:  8] + rs1[7:0], -2^{7} ), 2^{7}-1 )
    qv[ 23: 16] = min( max( qx[ 23: 16] + rs1[7:0], -2^{7} ), 2^{7}-1 )
    qv[ 31: 24] = min( max( qx[ 31: 24] + rs1[7:0], -2^{7} ), 2^{7}-1 )
    qv[ 39: 32] = min( max( qx[ 39: 32] + rs1[7:0], -2^{7} ), 2^{7}-1 )
    qv[ 47: 40] = min( max( qx[ 47: 40] + rs1[7:0], -2^{7} ), 2^{7}-1 )
    qv[ 55: 48] = min( max( qx[ 55: 48] + rs1[7:0], -2^{7} ), 2^{7}-1 )
    qv[ 63: 56] = min( max( qx[ 63: 56] + rs1[7:0], -2^{7} ), 2^{7}-1 )
    qv[ 71: 64] = min( max( qx[ 71: 64] + rs1[7:0], -2^{7} ), 2^{7}-1 )
    qv[ 79: 72] = min( max( qx[ 79: 72] + rs1[7:0], -2^{7} ), 2^{7}-1 )
    qv[ 87: 80] = min( max( qx[ 87: 80] + rs1[7:0], -2^{7} ), 2^{7}-1 )
    qv[ 95: 88] = min( max( qx[ 95: 88] + rs1[7:0], -2^{7} ), 2^{7}-1 )
    qv[103: 96] = min( max( qx[103: 96] + rs1[7:0], -2^{7} ), 2^{7}-1 )
    qv[111:104] = min( max( qx[111:104] + rs1[7:0], -2^{7} ), 2^{7}-1 )
    qv[119:112] = min( max( qx[119:112] + rs1[7:0], -2^{7} ), 2^{7}-1 )
    qv[127:120] = min( max( qx[127:120] + rs1[7:0], -2^{7} ), 2^{7}-1 )
\end{lstlisting}

\textbf{Schedule}\\
\begin{longtable}[c]{|l|l|l|}
    \hline
    \rowcolor{lightgray}
    \textbf{Operands} & \textbf{Use} & \textbf{Def} \\\hline
    \endfirsthead
    \hline
    \rowcolor{lightgray}
     \textbf{Operands} & \textbf{Use} & \textbf{Def} \\\hline
    \endhead
    \hline
    \endfoot

  qv & - & 2 \\\hline
    qx & 1 & - \\\hline
    rs1 & 1 & -

\end{longtable}

\textbf{Format}\\

\begin{figure*}[!htbp]
 {\setlength{\tabcolsep}{1pt}
 \begin{center}
\begin{tabular}{@{}F@{}W@{}I@{}F@{}F@{}I@{}I@{}F@{}E@{}O}
\instbitrange{31}{29}&
\instbitrange{28}{27}&
\instbit{26}&
\instbitrange{25}{23}&
\instbitrange{22}{20}&
\instbit{19}&
\instbit{18}&
\instbitrange{17}{15}&
\instbitrange{14}{7}&
\instbitrange{6}{0} \\
\hline
\multicolumn{1}{|c|}{{\scriptsize qx[2:0]}} &
\multicolumn{1}{c|}{01} &
\multicolumn{1}{c|}{{\scriptsize sat[0]}} &
\multicolumn{1}{c|}{1X1} &
\multicolumn{1}{c|}{{\scriptsize qv[2:0]}} &
\multicolumn{1}{c|}{X} &
\multicolumn{1}{c|}{{\scriptsize rs1[4]}} &
\multicolumn{1}{c|}{{\scriptsize rs1[2:0]}} &
\multicolumn{1}{c|}{01XX001X} &
\multicolumn{1}{c|}{0011011}\\
\hline
3& 2& 1& 3& 3& 1& 1& 3& 8& 7 \\
\end{tabular}
 \end{center}
 }
 \vspace{-0.1in}
 \caption[]{ESP.VSADDS.S8} \end{figure*}


\newpage
\subsubsection{ESP.VSADDS.U16}\label{instruction:ESPVSADDSU16}
\textbf{Syntax}\\
ESP.VSADDS.U16 qv,qx,rs1,sat

\textbf{Description}\\
This instruction performs an unsigned vector addition on 16-bit data.

\textbf{Operation}\\
\begin{lstlisting}[escapeinside=`?]
    qv[ 15:  0] = min( qx[ 15:  0] + rs1[15:0], 2^{16}-1 )
    qv[ 31: 16] = min( qx[ 31: 16] + rs1[15:0], 2^{16}-1 )
    qv[ 47: 32] = min( qx[ 47: 32] + rs1[15:0], 2^{16}-1 )
    qv[ 63: 48] = min( qx[ 63: 48] + rs1[15:0], 2^{16}-1 )
    qv[ 79: 64] = min( qx[ 79: 64] + rs1[15:0], 2^{16}-1 )
    qv[ 95: 80] = min( qx[ 95: 80] + rs1[15:0], 2^{16}-1 )
    qv[111: 96] = min( qx[111: 96] + rs1[15:0], 2^{16}-1 )
    qv[127:112] = min( qx[127:112] + rs1[15:0], 2^{16}-1 )
\end{lstlisting}


\textbf{Schedule}\\
\begin{longtable}[c]{|l|l|l|}
    \hline
    \rowcolor{lightgray}
    \textbf{Operands} & \textbf{Use} & \textbf{Def} \\\hline
    \endfirsthead
    \hline
    \rowcolor{lightgray}
     \textbf{Operands} & \textbf{Use} & \textbf{Def} \\\hline
    \endhead
    \hline
    \endfoot

  qv & - & 2 \\\hline
    qx & 1 & - \\\hline
    rs1 & 1 & -

\end{longtable}

\textbf{Format}\\

\begin{figure*}[!htbp]
 {\setlength{\tabcolsep}{1pt}
 \begin{center}
\begin{tabular}{@{}F@{}W@{}I@{}F@{}F@{}I@{}I@{}F@{}E@{}O}
\instbitrange{31}{29}&
\instbitrange{28}{27}&
\instbit{26}&
\instbitrange{25}{23}&
\instbitrange{22}{20}&
\instbit{19}&
\instbit{18}&
\instbitrange{17}{15}&
\instbitrange{14}{7}&
\instbitrange{6}{0} \\
\hline
\multicolumn{1}{|c|}{{\scriptsize qx[2:0]}} &
\multicolumn{1}{c|}{10} &
\multicolumn{1}{c|}{{\scriptsize sat[0]}} &
\multicolumn{1}{c|}{1X1} &
\multicolumn{1}{c|}{{\scriptsize qv[2:0]}} &
\multicolumn{1}{c|}{X} &
\multicolumn{1}{c|}{{\scriptsize rs1[4]}} &
\multicolumn{1}{c|}{{\scriptsize rs1[2:0]}} &
\multicolumn{1}{c|}{01XX001X} &
\multicolumn{1}{c|}{0011011}\\
\hline
3& 2& 1& 3& 3& 1& 1& 3& 8& 7 \\
\end{tabular}
 \end{center}
 }
 \vspace{-0.1in}
 \caption[]{ESP.VSADDS.U16} \end{figure*}


\newpage
\subsubsection{ESP.VSADDS.U8}\label{instruction:ESPVSADDSU8}
\textbf{Syntax}\\
ESP.VSADDS.U8 qv,qx,rs1,sat

\textbf{Description}\\
This instruction performs an unsigned vector addition on 8-bit data.

\textbf{Operation}\\
\begin{lstlisting}[escapeinside=`?]
    qv[  7:  0] = min( qx[  7:  0] + rs1[7:0], 2^{8}-1 )
    qv[ 15:  8] = min( qx[ 15:  8] + rs1[7:0], 2^{8}-1 )
    qv[ 23: 16] = min( qx[ 23: 16] + rs1[7:0], 2^{8}-1 )
    qv[ 31: 24] = min( qx[ 31: 24] + rs1[7:0], 2^{8}-1 )
    qv[ 39: 32] = min( qx[ 39: 32] + rs1[7:0], 2^{8}-1 )
    qv[ 47: 40] = min( qx[ 47: 40] + rs1[7:0], 2^{8}-1 )
    qv[ 55: 48] = min( qx[ 55: 48] + rs1[7:0], 2^{8}-1 )
    qv[ 63: 56] = min( qx[ 63: 56] + rs1[7:0], 2^{8}-1 )
    qv[ 71: 64] = min( qx[ 71: 64] + rs1[7:0], 2^{8}-1 )
    qv[ 79: 72] = min( qx[ 79: 72] + rs1[7:0], 2^{8}-1 )
    qv[ 87: 80] = min( qx[ 87: 80] + rs1[7:0], 2^{8}-1 )
    qv[ 95: 88] = min( qx[ 95: 88] + rs1[7:0], 2^{8}-1 )
    qv[103: 96] = min( qx[103: 96] + rs1[7:0], 2^{8}-1 )
    qv[111:104] = min( qx[111:104] + rs1[7:0], 2^{8}-1 )
    qv[119:112] = min( qx[119:112] + rs1[7:0], 2^{8}-1 )
    qv[127:120] = min( qx[127:120] + rs1[7:0], 2^{8}-1 )
\end{lstlisting}


\textbf{Schedule}\\
\begin{longtable}[c]{|l|l|l|}
    \hline
    \rowcolor{lightgray}
    \textbf{Operands} & \textbf{Use} & \textbf{Def} \\\hline
    \endfirsthead
    \hline
    \rowcolor{lightgray}
     \textbf{Operands} & \textbf{Use} & \textbf{Def} \\\hline
    \endhead
    \hline
    \endfoot

  qv & - & 2 \\\hline
    qx & 1 & - \\\hline
    rs1 & 1 & -

\end{longtable}

\textbf{Format}\\

\begin{figure*}[!htbp]
 {\setlength{\tabcolsep}{1pt}
 \begin{center}
\begin{tabular}{@{}F@{}W@{}I@{}F@{}F@{}I@{}I@{}F@{}E@{}O}
\instbitrange{31}{29}&
\instbitrange{28}{27}&
\instbit{26}&
\instbitrange{25}{23}&
\instbitrange{22}{20}&
\instbit{19}&
\instbit{18}&
\instbitrange{17}{15}&
\instbitrange{14}{7}&
\instbitrange{6}{0} \\
\hline
\multicolumn{1}{|c|}{{\scriptsize qx[2:0]}} &
\multicolumn{1}{c|}{00} &
\multicolumn{1}{c|}{{\scriptsize sat[0]}} &
\multicolumn{1}{c|}{1X1} &
\multicolumn{1}{c|}{{\scriptsize qv[2:0]}} &
\multicolumn{1}{c|}{X} &
\multicolumn{1}{c|}{{\scriptsize rs1[4]}} &
\multicolumn{1}{c|}{{\scriptsize rs1[2:0]}} &
\multicolumn{1}{c|}{01XX001X} &
\multicolumn{1}{c|}{0011011}\\
\hline
3& 2& 1& 3& 3& 1& 1& 3& 8& 7 \\
\end{tabular}
 \end{center}
 }
 \vspace{-0.1in}
 \caption[]{ESP.VSADDS.U8} \end{figure*}


\newpage
\subsubsection{ESP.VSAT.S16}\label{instruction:ESPVSATS16}
\textbf{Syntax}\\
ESP.VSAT.S16 qz,qx,rs1,rs2

\textbf{Description}\\
This instruction performs vector saturation for 16-bit signed segments.

\textbf{Operation}\\
\begin{lstlisting}[escapeinside=`?]
    min_t[15:0] = max( rs1[15:0], rs2[15:0] )
    max_t[15:0] = min( rs1[15:0], rs2[15:0] )

    qz[ 15:  0] = max( min( qx[ 15:  0], max_t[15:0]), min_t[15:0] )
    qz[ 31: 16] = max( min( qx[ 31: 16], maxt_[15:0]), min_t[15:0] )
    qz[ 47: 32] = max( min( qx[ 47: 32], max_t[15:0]), min_t[15:0] )
    qz[ 63: 48] = max( min( qx[ 63: 48], max_t[15:0]), min_t[15:0] )
    qz[ 79: 64] = max( min( qx[ 79: 64], max_t[15:0]), min_t[15:0] )
    qz[ 95: 80] = max( min( qx[ 95: 80], max_t[15:0]), min_t[15:0] )
    qz[111: 96] = max( min( qx[111: 96], max_t[15:0]), min_t[15:0] )
    qz[127:112] = max( min( qx[127:112], max_t[15:0]), min_t[15:0] )
\end{lstlisting}


\textbf{Schedule}\\
\begin{longtable}[c]{|l|l|l|}
    \hline
    \rowcolor{lightgray}
    \textbf{Operands} & \textbf{Use} & \textbf{Def} \\\hline
    \endfirsthead
    \hline
    \rowcolor{lightgray}
     \textbf{Operands} & \textbf{Use} & \textbf{Def} \\\hline
    \endhead
    \hline
    \endfoot

  qz & - & 1 \\\hline
    qx & 1 & - \\\hline
    rs1 & 1 & - \\\hline
    rs2 & 1 & -

\end{longtable}

\textbf{Format}\\

\begin{figure*}[!htbp]
 {\setlength{\tabcolsep}{1pt}
 \begin{center}
\begin{tabular}{@{}F@{}V@{}I@{}F@{}I@{}I@{}F@{}V@{}F@{}O}
\instbitrange{31}{29}&
\instbitrange{28}{24}&
\instbit{23}&
\instbitrange{22}{20}&
\instbit{19}&
\instbit{18}&
\instbitrange{17}{15}&
\instbitrange{14}{10}&
\instbitrange{9}{7}&
\instbitrange{6}{0} \\
\hline
\multicolumn{1}{|c|}{{\scriptsize qx[2:0]}} &
\multicolumn{1}{c|}{11010} &
\multicolumn{1}{c|}{{\scriptsize rs2[4]}} &
\multicolumn{1}{c|}{{\scriptsize rs2[2:0]}} &
\multicolumn{1}{c|}{X} &
\multicolumn{1}{c|}{{\scriptsize rs1[4]}} &
\multicolumn{1}{c|}{{\scriptsize rs1[2:0]}} &
\multicolumn{1}{c|}{001XX} &
\multicolumn{1}{c|}{{\scriptsize qz[2:0]}} &
\multicolumn{1}{c|}{0011111}\\
\hline
3& 5& 1& 3& 1& 1& 3& 5& 3& 7 \\
\end{tabular}
 \end{center}
 }
 \vspace{-0.1in}
 \caption[]{ESP.VSAT.S16} \end{figure*}


\newpage
\subsubsection{ESP.VSAT.S32}\label{instruction:ESPVSATS32}
\textbf{Syntax}\\
ESP.VSAT.S32 qz,qx,rs1,rs2

\textbf{Description}\\
This instruction performs vector saturation for 32-bit signed segments.

\textbf{Operation}\\
\begin{lstlisting}[escapeinside=`?]
    min_t[31:0] = max( rs1[31;0], rs2[31:0] )
    max_t[31:0] = min( rs1[31:0], rs2[31:0] )

    qz[ 31:  0] = max( min( qx[ 31:  0], max_t[31:0]), min_t[31:0] )
    qz[ 63: 32] = max( min( qx[ 63: 32], max_t[31:0]), min_t[31:0] )
    qz[ 95: 64] = max( min( qx[ 95: 64], max_t[31:0]), min_t[31:0] )
    qz[127: 96] = max( min( qx[127: 96], maxt_[31:0]), min_t[31:0] )

\end{lstlisting}


\textbf{Schedule}\\
\begin{longtable}[c]{|l|l|l|}
    \hline
    \rowcolor{lightgray}
    \textbf{Operands} & \textbf{Use} & \textbf{Def} \\\hline
    \endfirsthead
    \hline
    \rowcolor{lightgray}
     \textbf{Operands} & \textbf{Use} & \textbf{Def} \\\hline
    \endhead
    \hline
    \endfoot

  qz & - & 1 \\\hline
    qx & 1 & - \\\hline
    rs1 & 1 & - \\\hline
    rs2 & 1 & -

\end{longtable}

\textbf{Format}\\

\begin{figure*}[!htbp]
 {\setlength{\tabcolsep}{1pt}
 \begin{center}
\begin{tabular}{@{}F@{}V@{}I@{}F@{}I@{}I@{}F@{}V@{}F@{}O}
\instbitrange{31}{29}&
\instbitrange{28}{24}&
\instbit{23}&
\instbitrange{22}{20}&
\instbit{19}&
\instbit{18}&
\instbitrange{17}{15}&
\instbitrange{14}{10}&
\instbitrange{9}{7}&
\instbitrange{6}{0} \\
\hline
\multicolumn{1}{|c|}{{\scriptsize qx[2:0]}} &
\multicolumn{1}{c|}{1X110} &
\multicolumn{1}{c|}{{\scriptsize rs2[4]}} &
\multicolumn{1}{c|}{{\scriptsize rs2[2:0]}} &
\multicolumn{1}{c|}{X} &
\multicolumn{1}{c|}{{\scriptsize rs1[4]}} &
\multicolumn{1}{c|}{{\scriptsize rs1[2:0]}} &
\multicolumn{1}{c|}{001XX} &
\multicolumn{1}{c|}{{\scriptsize qz[2:0]}} &
\multicolumn{1}{c|}{0011111}\\
\hline
3& 5& 1& 3& 1& 1& 3& 5& 3& 7 \\
\end{tabular}
 \end{center}
 }
 \vspace{-0.1in}
 \caption[]{ESP.VSAT.S32} \end{figure*}


\newpage
\subsubsection{ESP.VSAT.S8}\label{instruction:ESPVSATS8}
\textbf{Syntax}\\
ESP.VSAT.S8 qz,qx,rs1,rs2

\textbf{Description}\\
This instruction performs vector saturation for 8-bit signed segments.

\textbf{Operation}\\
\begin{lstlisting}[escapeinside=`?]
    min_t[7:0] = max( rs1[7:0], rs2[7:0] )
    max_t[7:0] = min( rs1[7:0], rs2[7:0] )

    qz[  7:  0] = max( min( qx[  7:  0], max_t[7:0]), min_t[7:0] )
    qz[ 15:  8] = max( min( qx[ 15:  8], max_t[7:0]), min_t[7:0] )
    qz[ 23: 16] = max( min( qx[ 23: 16], max_t[7:0]), min_t[7:0] )
    qz[ 31: 24] = max( min( qx[ 31: 24], maxt_[7:0]), min_t[7:0] )
    qz[ 39: 32] = max( min( qx[ 39: 32], max_t[7:0]), min_t[7:0] )
    qz[ 47: 40] = max( min( qx[ 47: 40], max_t[7:0]), min_t[7:0] )
    qz[ 55: 48] = max( min( qx[ 55: 48], max_t[7:0]), min_t[7:0] )
    qz[ 63: 56] = max( min( qx[ 63: 56], max_t[7:0]), min_t[7:0] )
    qz[ 71: 64] = max( min( qx[ 71: 64], max_t[7:0]), min_t[7:0] )
    qz[ 79: 72] = max( min( qx[ 79: 72], max_t[7:0]), min_t[7:0] )
    qz[ 87: 80] = max( min( qx[ 87: 80], max_t[7:0]), min_t[7:0] )
    qz[ 95: 88] = max( min( qx[ 95: 88], max_t[7:0]), min_t[7:0] )
    qz[103: 96] = max( min( qx[103: 96], max_t[7:0]), min_t[7:0] )
    qz[111:104] = max( min( qx[111:104], max_t[7:0]), min_t[7:0] )
    qz[119:112] = max( min( qx[119:112], max_t[7:0]), min_t[7:0] )
    qz[127:120] = max( min( qx[127:120], max_t[7:0]), min_t[7:0] )
\end{lstlisting}

\textbf{Schedule}\\
\begin{longtable}[c]{|l|l|l|}
    \hline
    \rowcolor{lightgray}
    \textbf{Operands} & \textbf{Use} & \textbf{Def} \\\hline
    \endfirsthead
    \hline
    \rowcolor{lightgray}
     \textbf{Operands} & \textbf{Use} & \textbf{Def} \\\hline
    \endhead
    \hline
    \endfoot

  qz & - & 1 \\\hline
    qx & 1 & - \\\hline
    rs1 & 1 & - \\\hline
    rs2 & 1 & -

\end{longtable}

\textbf{Format}\\

\begin{figure*}[!htbp]
 {\setlength{\tabcolsep}{1pt}
 \begin{center}
\begin{tabular}{@{}F@{}V@{}I@{}F@{}I@{}I@{}F@{}V@{}F@{}O}
\instbitrange{31}{29}&
\instbitrange{28}{24}&
\instbit{23}&
\instbitrange{22}{20}&
\instbit{19}&
\instbit{18}&
\instbitrange{17}{15}&
\instbitrange{14}{10}&
\instbitrange{9}{7}&
\instbitrange{6}{0} \\
\hline
\multicolumn{1}{|c|}{{\scriptsize qx[2:0]}} &
\multicolumn{1}{c|}{01010} &
\multicolumn{1}{c|}{{\scriptsize rs2[4]}} &
\multicolumn{1}{c|}{{\scriptsize rs2[2:0]}} &
\multicolumn{1}{c|}{X} &
\multicolumn{1}{c|}{{\scriptsize rs1[4]}} &
\multicolumn{1}{c|}{{\scriptsize rs1[2:0]}} &
\multicolumn{1}{c|}{001XX} &
\multicolumn{1}{c|}{{\scriptsize qz[2:0]}} &
\multicolumn{1}{c|}{0011111}\\
\hline
3& 5& 1& 3& 1& 1& 3& 5& 3& 7 \\
\end{tabular}
 \end{center}
 }
 \vspace{-0.1in}
 \caption[]{ESP.VSAT.S8} \end{figure*}


\newpage
\subsubsection{ESP.VSAT.U16}\label{instruction:ESPVSATU16}
\textbf{Syntax}\\
ESP.VSAT.U16 qz,qx,rs1,rs2

\textbf{Description}\\
This instruction performs vector saturation for 16-bit unsigned segments.

\textbf{Operation}\\
\begin{lstlisting}[escapeinside=`?]
    min_t[15:0] = max( rs1[15:0], rs2[15:0] )
    max_t[15:0] = min( rs1[15:0], rs2[15:0] )

    qz[ 15:  0] = max( min( qx[ 15:  0], max_t[15:0]), min_t[15:0] )
    qz[ 31: 16] = max( min( qx[ 31: 16], maxt_[15:0]), min_t[15:0] )
    qz[ 47: 32] = max( min( qx[ 47: 32], max_t[15:0]), min_t[15:0] )
    qz[ 63: 48] = max( min( qx[ 63: 48], max_t[15:0]), min_t[15:0] )
    qz[ 79: 64] = max( min( qx[ 79: 64], max_t[15:0]), min_t[15:0] )
    qz[ 95: 80] = max( min( qx[ 95: 80], max_t[15:0]), min_t[15:0] )
    qz[111: 96] = max( min( qx[111: 96], max_t[15:0]), min_t[15:0] )
    qz[127:112] = max( min( qx[127:112], max_t[15:0]), min_t[15:0] )
\end{lstlisting}


\textbf{Schedule}\\
\begin{longtable}[c]{|l|l|l|}
    \hline
    \rowcolor{lightgray}
    \textbf{Operands} & \textbf{Use} & \textbf{Def} \\\hline
    \endfirsthead
    \hline
    \rowcolor{lightgray}
     \textbf{Operands} & \textbf{Use} & \textbf{Def} \\\hline
    \endhead
    \hline
    \endfoot

  qz & - & 1 \\\hline
    qx & 1 & - \\\hline
    rs1 & 1 & - \\\hline
    rs2 & 1 & -

\end{longtable}

\textbf{Format}\\

\begin{figure*}[!htbp]
 {\setlength{\tabcolsep}{1pt}
 \begin{center}
\begin{tabular}{@{}F@{}V@{}I@{}F@{}I@{}I@{}F@{}V@{}F@{}O}
\instbitrange{31}{29}&
\instbitrange{28}{24}&
\instbit{23}&
\instbitrange{22}{20}&
\instbit{19}&
\instbit{18}&
\instbitrange{17}{15}&
\instbitrange{14}{10}&
\instbitrange{9}{7}&
\instbitrange{6}{0} \\
\hline
\multicolumn{1}{|c|}{{\scriptsize qx[2:0]}} &
\multicolumn{1}{c|}{10010} &
\multicolumn{1}{c|}{{\scriptsize rs2[4]}} &
\multicolumn{1}{c|}{{\scriptsize rs2[2:0]}} &
\multicolumn{1}{c|}{X} &
\multicolumn{1}{c|}{{\scriptsize rs1[4]}} &
\multicolumn{1}{c|}{{\scriptsize rs1[2:0]}} &
\multicolumn{1}{c|}{001XX} &
\multicolumn{1}{c|}{{\scriptsize qz[2:0]}} &
\multicolumn{1}{c|}{0011111}\\
\hline
3& 5& 1& 3& 1& 1& 3& 5& 3& 7 \\
\end{tabular}
 \end{center}
 }
 \vspace{-0.1in}
 \caption[]{ESP.VSAT.U16} \end{figure*}


\newpage
\subsubsection{ESP.VSAT.U32}\label{instruction:ESPVSATU32}
\textbf{Syntax}\\
ESP.VSAT.U32 qz,qx,rs1,rs2

\textbf{Description}\\
This instruction performs vector saturation for 32-bit unsigned segments.

\textbf{Operation}\\
\begin{lstlisting}[escapeinside=`?]
    min_t[31:0] = max( rs1[31;0], rs2[31:0] )
    max_t[31:0] = min( rs1[31:0], rs2[31:0] )

    qz[ 31:  0] = max( min( qx[ 31:  0], max_t[31:0]), min_t[31:0] )
    qz[ 63: 32] = max( min( qx[ 63: 32], max_t[31:0]), min_t[31:0] )
    qz[ 95: 64] = max( min( qx[ 95: 64], max_t[31:0]), min_t[31:0] )
    qz[127: 96] = max( min( qx[127: 96], maxt_[31:0]), min_t[31:0] )

\end{lstlisting}


\textbf{Schedule}\\
\begin{longtable}[c]{|l|l|l|}
    \hline
    \rowcolor{lightgray}
    \textbf{Operands} & \textbf{Use} & \textbf{Def} \\\hline
    \endfirsthead
    \hline
    \rowcolor{lightgray}
     \textbf{Operands} & \textbf{Use} & \textbf{Def} \\\hline
    \endhead
    \hline
    \endfoot

  qz & - & 1 \\\hline
    qx & 1 & - \\\hline
    rs1 & 1 & - \\\hline
    rs2 & 1 & -

\end{longtable}

\textbf{Format}\\

\begin{figure*}[!htbp]
 {\setlength{\tabcolsep}{1pt}
 \begin{center}
\begin{tabular}{@{}F@{}V@{}I@{}F@{}I@{}I@{}F@{}V@{}F@{}O}
\instbitrange{31}{29}&
\instbitrange{28}{24}&
\instbit{23}&
\instbitrange{22}{20}&
\instbit{19}&
\instbit{18}&
\instbitrange{17}{15}&
\instbitrange{14}{10}&
\instbitrange{9}{7}&
\instbitrange{6}{0} \\
\hline
\multicolumn{1}{|c|}{{\scriptsize qx[2:0]}} &
\multicolumn{1}{c|}{0X110} &
\multicolumn{1}{c|}{{\scriptsize rs2[4]}} &
\multicolumn{1}{c|}{{\scriptsize rs2[2:0]}} &
\multicolumn{1}{c|}{X} &
\multicolumn{1}{c|}{{\scriptsize rs1[4]}} &
\multicolumn{1}{c|}{{\scriptsize rs1[2:0]}} &
\multicolumn{1}{c|}{001XX} &
\multicolumn{1}{c|}{{\scriptsize qz[2:0]}} &
\multicolumn{1}{c|}{0011111}\\
\hline
3& 5& 1& 3& 1& 1& 3& 5& 3& 7 \\
\end{tabular}
 \end{center}
 }
 \vspace{-0.1in}
 \caption[]{ESP.VSAT.U32} \end{figure*}


\newpage
\subsubsection{ESP.VSAT.U8}\label{instruction:ESPVSATU8}
\textbf{Syntax}\\
ESP.VSAT.U8 qz,qx,rs1,rs2

\textbf{Description}\\
This instruction performs vector saturation for 8-bit unsigned segments.

\textbf{Operation}\\
\begin{lstlisting}[escapeinside=`?]
    min_t[7:0] = max( rs1[7:0], rs2[7:0] )
    max_t[7:0] = min( rs1[7:0], rs2[7:0] )

    qz[  7:  0] = max( min( qx[  7:  0], max_t[7:0]), min_t[7:0] )
    qz[ 15:  8] = max( min( qx[ 15:  8], max_t[7:0]), min_t[7:0] )
    qz[ 23: 16] = max( min( qx[ 23: 16], max_t[7:0]), min_t[7:0] )
    qz[ 31: 24] = max( min( qx[ 31: 24], maxt_[7:0]), min_t[7:0] )
    qz[ 39: 32] = max( min( qx[ 39: 32], max_t[7:0]), min_t[7:0] )
    qz[ 47: 40] = max( min( qx[ 47: 40], max_t[7:0]), min_t[7:0] )
    qz[ 55: 48] = max( min( qx[ 55: 48], max_t[7:0]), min_t[7:0] )
    qz[ 63: 56] = max( min( qx[ 63: 56], max_t[7:0]), min_t[7:0] )
    qz[ 71: 64] = max( min( qx[ 71: 64], max_t[7:0]), min_t[7:0] )
    qz[ 79: 72] = max( min( qx[ 79: 72], max_t[7:0]), min_t[7:0] )
    qz[ 87: 80] = max( min( qx[ 87: 80], max_t[7:0]), min_t[7:0] )
    qz[ 95: 88] = max( min( qx[ 95: 88], max_t[7:0]), min_t[7:0] )
    qz[103: 96] = max( min( qx[103: 96], max_t[7:0]), min_t[7:0] )
    qz[111:104] = max( min( qx[111:104], max_t[7:0]), min_t[7:0] )
    qz[119:112] = max( min( qx[119:112], max_t[7:0]), min_t[7:0] )
    qz[127:120] = max( min( qx[127:120], max_t[7:0]), min_t[7:0] )
\end{lstlisting}


\textbf{Schedule}\\
\begin{longtable}[c]{|l|l|l|}
    \hline
    \rowcolor{lightgray}
    \textbf{Operands} & \textbf{Use} & \textbf{Def} \\\hline
    \endfirsthead
    \hline
    \rowcolor{lightgray}
     \textbf{Operands} & \textbf{Use} & \textbf{Def} \\\hline
    \endhead
    \hline
    \endfoot

  qz & - & 1 \\\hline
    qx & 1 & - \\\hline
    rs1 & 1 & - \\\hline
    rs2 & 1 & -

\end{longtable}

\textbf{Format}\\

\begin{figure*}[!htbp]
 {\setlength{\tabcolsep}{1pt}
 \begin{center}
\begin{tabular}{@{}F@{}V@{}I@{}F@{}I@{}I@{}F@{}V@{}F@{}O}
\instbitrange{31}{29}&
\instbitrange{28}{24}&
\instbit{23}&
\instbitrange{22}{20}&
\instbit{19}&
\instbit{18}&
\instbitrange{17}{15}&
\instbitrange{14}{10}&
\instbitrange{9}{7}&
\instbitrange{6}{0} \\
\hline
\multicolumn{1}{|c|}{{\scriptsize qx[2:0]}} &
\multicolumn{1}{c|}{00010} &
\multicolumn{1}{c|}{{\scriptsize rs2[4]}} &
\multicolumn{1}{c|}{{\scriptsize rs2[2:0]}} &
\multicolumn{1}{c|}{X} &
\multicolumn{1}{c|}{{\scriptsize rs1[4]}} &
\multicolumn{1}{c|}{{\scriptsize rs1[2:0]}} &
\multicolumn{1}{c|}{001XX} &
\multicolumn{1}{c|}{{\scriptsize qz[2:0]}} &
\multicolumn{1}{c|}{0011111}\\
\hline
3& 5& 1& 3& 1& 1& 3& 5& 3& 7 \\
\end{tabular}
 \end{center}
 }
 \vspace{-0.1in}
 \caption[]{ESP.VSAT.U8} \end{figure*}


\newpage
\subsubsection{ESP.VSSUBS.S16}\label{instruction:ESPVSSUBSS16}
\textbf{Syntax}\\
ESP.VSSUBS.S16 qv,qx,rs1,sat

\textbf{Description}\\
This instruction performs a signed vector subtraction on 16-bit data.

\textbf{Operation}\\
\begin{lstlisting}[escapeinside=`?]
    qv[ 15:  0] = min( max( qx[ 15:  0] － rs1[15:0], -2^{15} ), 2^{15}-1 )
    qv[ 31: 16] = min( max( qx[ 31: 16] － rs1[15:0], -2^{15} ), 2^{15}-1 )
    qv[ 47: 32] = min( max( qx[ 47: 32] － rs1[15:0], -2^{15} ), 2^{15}-1 )
    qv[ 63: 48] = min( max( qx[ 63: 48] － rs1[15:0], -2^{15} ), 2^{15}-1 )
    qv[ 79: 64] = min( max( qx[ 79: 64] － rs1[15:0], -2^{15} ), 2^{15}-1 )
    qv[ 95: 80] = min( max( qx[ 95: 80] － rs1[15:0], -2^{15} ), 2^{15}-1 )
    qv[111: 96] = min( max( qx[111: 96] － rs1[15:0], -2^{15} ), 2^{15}-1 )
    qv[127:112] = min( max( qx[127:112] － rs1[15:0], -2^{15} ), 2^{15}-1 )
\end{lstlisting}


\textbf{Schedule}\\
\begin{longtable}[c]{|l|l|l|}
    \hline
    \rowcolor{lightgray}
    \textbf{Operands} & \textbf{Use} & \textbf{Def} \\\hline
    \endfirsthead
    \hline
    \rowcolor{lightgray}
     \textbf{Operands} & \textbf{Use} & \textbf{Def} \\\hline
    \endhead
    \hline
    \endfoot

  qv & - & 1 \\\hline
    qx & 1 & - \\\hline
    rs1 & 1 & -

\end{longtable}

\textbf{Format}\\

\begin{figure*}[!htbp]
 {\setlength{\tabcolsep}{1pt}
 \begin{center}
\begin{tabular}{@{}F@{}W@{}I@{}F@{}F@{}I@{}I@{}F@{}E@{}O}
\instbitrange{31}{29}&
\instbitrange{28}{27}&
\instbit{26}&
\instbitrange{25}{23}&
\instbitrange{22}{20}&
\instbit{19}&
\instbit{18}&
\instbitrange{17}{15}&
\instbitrange{14}{7}&
\instbitrange{6}{0} \\
\hline
\multicolumn{1}{|c|}{{\scriptsize qx[2:0]}} &
\multicolumn{1}{c|}{11} &
\multicolumn{1}{c|}{{\scriptsize sat[0]}} &
\multicolumn{1}{c|}{1X1} &
\multicolumn{1}{c|}{{\scriptsize qv[2:0]}} &
\multicolumn{1}{c|}{X} &
\multicolumn{1}{c|}{{\scriptsize rs1[4]}} &
\multicolumn{1}{c|}{{\scriptsize rs1[2:0]}} &
\multicolumn{1}{c|}{01XX101X} &
\multicolumn{1}{c|}{0011011}\\
\hline
3& 2& 1& 3& 3& 1& 1& 3& 8& 7 \\
\end{tabular}
 \end{center}
 }
 \vspace{-0.1in}
 \caption[]{ESP.VSSUBS.S16} \end{figure*}


\newpage
\subsubsection{ESP.VSSUBS.S8}\label{instruction:ESPVSSUBSS8}
\textbf{Syntax}\\
ESP.VSSUBS.S8 qv,qx,rs1,sat

\textbf{Description}\\
This instruction performs a signed vector subtraction on 8-bit data.

\textbf{Operation}\\
\begin{lstlisting}[escapeinside=`?]
    qv[  7:  0] = min( max( qx[  7:  0] － rs1[7:0], -2^{7} ), 2^{7}-1 )
    qv[ 15:  8] = min( max( qx[ 15:  8] － rs1[7:0], -2^{7} ), 2^{7}-1 )
    qv[ 23: 16] = min( max( qx[ 23: 16] － rs1[7:0], -2^{7} ), 2^{7}-1 )
    qv[ 31: 24] = min( max( qx[ 31: 24] － rs1[7:0], -2^{7} ), 2^{7}-1 )
    qv[ 39: 32] = min( max( qx[ 39: 32] － rs1[7:0], -2^{7} ), 2^{7}-1 )
    qv[ 47: 40] = min( max( qx[ 47: 40] － rs1[7:0], -2^{7} ), 2^{7}-1 )
    qv[ 55: 48] = min( max( qx[ 55: 48] － rs1[7:0], -2^{7} ), 2^{7}-1 )
    qv[ 63: 56] = min( max( qx[ 63: 56] － rs1[7:0], -2^{7} ), 2^{7}-1 )
    qv[ 71: 64] = min( max( qx[ 71: 64] － rs1[7:0], -2^{7} ), 2^{7}-1 )
    qv[ 79: 72] = min( max( qx[ 79: 72] － rs1[7:0], -2^{7} ), 2^{7}-1 )
    qv[ 87: 80] = min( max( qx[ 87: 80] － rs1[7:0], -2^{7} ), 2^{7}-1 )
    qv[ 95: 88] = min( max( qx[ 95: 88] － rs1[7:0], -2^{7} ), 2^{7}-1 )
    qv[103: 96] = min( max( qx[103: 96] － rs1[7:0], -2^{7} ), 2^{7}-1 )
    qv[111:104] = min( max( qx[111:104] － rs1[7:0], -2^{7} ), 2^{7}-1 )
    qv[119:112] = min( max( qx[119:112] － rs1[7:0], -2^{7} ), 2^{7}-1 )
    qv[127:120] = min( max( qx[127:120] － rs1[7:0], -2^{7} ), 2^{7}-1 )
\end{lstlisting}


\textbf{Schedule}\\
\begin{longtable}[c]{|l|l|l|}
    \hline
    \rowcolor{lightgray}
    \textbf{Operands} & \textbf{Use} & \textbf{Def} \\\hline
    \endfirsthead
    \hline
    \rowcolor{lightgray}
     \textbf{Operands} & \textbf{Use} & \textbf{Def} \\\hline
    \endhead
    \hline
    \endfoot

  qv & - & 1 \\\hline
    qx & 1 & - \\\hline
    rs1 & 1 & -

\end{longtable}

\textbf{Format}\\

\begin{figure*}[!htbp]
 {\setlength{\tabcolsep}{1pt}
 \begin{center}
\begin{tabular}{@{}F@{}W@{}I@{}F@{}F@{}I@{}I@{}F@{}E@{}O}
\instbitrange{31}{29}&
\instbitrange{28}{27}&
\instbit{26}&
\instbitrange{25}{23}&
\instbitrange{22}{20}&
\instbit{19}&
\instbit{18}&
\instbitrange{17}{15}&
\instbitrange{14}{7}&
\instbitrange{6}{0} \\
\hline
\multicolumn{1}{|c|}{{\scriptsize qx[2:0]}} &
\multicolumn{1}{c|}{01} &
\multicolumn{1}{c|}{{\scriptsize sat[0]}} &
\multicolumn{1}{c|}{1X1} &
\multicolumn{1}{c|}{{\scriptsize qv[2:0]}} &
\multicolumn{1}{c|}{X} &
\multicolumn{1}{c|}{{\scriptsize rs1[4]}} &
\multicolumn{1}{c|}{{\scriptsize rs1[2:0]}} &
\multicolumn{1}{c|}{01XX101X} &
\multicolumn{1}{c|}{0011011}\\
\hline
3& 2& 1& 3& 3& 1& 1& 3& 8& 7 \\
\end{tabular}
 \end{center}
 }
 \vspace{-0.1in}
 \caption[]{ESP.VSSUBS.S8} \end{figure*}


\newpage
\subsubsection{ESP.VSSUBS.U16}\label{instruction:ESPVSSUBSU16}
\textbf{Syntax}\\
ESP.VSSUBS.U16 qv,qx,rs1,sat

\textbf{Description}\\
This instruction performs an unsigned vector subtraction on 16-bit data.

\textbf{Operation}\\
\begin{lstlisting}[escapeinside=`?]
    qv[ 15:  0] = min( qx[ 15:  0] － rs1[15:0], 2^{16}-1 )
    qv[ 31: 16] = min( qx[ 31: 16] － rs1[15:0], 2^{16}-1 )
    qv[ 47: 32] = min( qx[ 47: 32] － rs1[15:0], 2^{16}-1 )
    qv[ 63: 48] = min( qx[ 63: 48] － rs1[15:0], 2^{16}-1 )
    qv[ 79: 64] = min( qx[ 79: 64] － rs1[15:0], 2^{16}-1 )
    qv[ 95: 80] = min( qx[ 95: 80] － rs1[15:0], 2^{16}-1 )
    qv[111: 96] = min( qx[111: 96] － rs1[15:0], 2^{16}-1 )
    qv[127:112] = min( qx[127:112] － rs1[15:0], 2^{16}-1 )
\end{lstlisting}


\textbf{Schedule}\\
\begin{longtable}[c]{|l|l|l|}
    \hline
    \rowcolor{lightgray}
    \textbf{Operands} & \textbf{Use} & \textbf{Def} \\\hline
    \endfirsthead
    \hline
    \rowcolor{lightgray}
     \textbf{Operands} & \textbf{Use} & \textbf{Def} \\\hline
    \endhead
    \hline
    \endfoot

  qv & - & 1 \\\hline
    qx & 1 & - \\\hline
    rs1 & 1 & -

\end{longtable}

\textbf{Format}\\

\begin{figure*}[!htbp]
 {\setlength{\tabcolsep}{1pt}
 \begin{center}
\begin{tabular}{@{}F@{}W@{}I@{}F@{}F@{}I@{}I@{}F@{}E@{}O}
\instbitrange{31}{29}&
\instbitrange{28}{27}&
\instbit{26}&
\instbitrange{25}{23}&
\instbitrange{22}{20}&
\instbit{19}&
\instbit{18}&
\instbitrange{17}{15}&
\instbitrange{14}{7}&
\instbitrange{6}{0} \\
\hline
\multicolumn{1}{|c|}{{\scriptsize qx[2:0]}} &
\multicolumn{1}{c|}{10} &
\multicolumn{1}{c|}{{\scriptsize sat[0]}} &
\multicolumn{1}{c|}{1X1} &
\multicolumn{1}{c|}{{\scriptsize qv[2:0]}} &
\multicolumn{1}{c|}{X} &
\multicolumn{1}{c|}{{\scriptsize rs1[4]}} &
\multicolumn{1}{c|}{{\scriptsize rs1[2:0]}} &
\multicolumn{1}{c|}{01XX101X} &
\multicolumn{1}{c|}{0011011}\\
\hline
3& 2& 1& 3& 3& 1& 1& 3& 8& 7 \\
\end{tabular}
 \end{center}
 }
 \vspace{-0.1in}
 \caption[]{ESP.VSSUBS.U16} \end{figure*}


\newpage
\subsubsection{ESP.VSSUBS.U8}\label{instruction:ESPVSSUBSU8}
\textbf{Syntax}\\
ESP.VSSUBS.U8 qv,qx,rs1,sat

\textbf{Description}\\
This instruction performs a signed vector subtraction on 8-bit data.

\textbf{Operation}\\
\begin{lstlisting}[escapeinside=`?]
    qv[  7:  0] = min( qx[  7:  0] - rs1[7:0], 2^{8}-1 )
    qv[ 15:  8] = min( qx[ 15:  8] - rs1[7:0], 2^{8}-1 )
    qv[ 23: 16] = min( qx[ 23: 16] - rs1[7:0], 2^{8}-1 )
    qv[ 31: 24] = min( qx[ 31: 24] - rs1[7:0], 2^{8}-1 )
    qv[ 39: 32] = min( qx[ 39: 32] - rs1[7:0], 2^{8}-1 )
    qv[ 47: 40] = min( qx[ 47: 40] - rs1[7:0], 2^{8}-1 )
    qv[ 55: 48] = min( qx[ 55: 48] - rs1[7:0], 2^{8}-1 )
    qv[ 63: 56] = min( qx[ 63: 56] - rs1[7:0], 2^{8}-1 )
    qv[ 71: 64] = min( qx[ 71: 64] - rs1[7:0], 2^{8}-1 )
    qv[ 79: 72] = min( qx[ 79: 72] - rs1[7:0], 2^{8}-1 )
    qv[ 87: 80] = min( qx[ 87: 80] - rs1[7:0], 2^{8}-1 )
    qv[ 95: 88] = min( qx[ 95: 88] - rs1[7:0], 2^{8}-1 )
    qv[103: 96] = min( qx[103: 96] - rs1[7:0], 2^{8}-1 )
    qv[111:104] = min( qx[111:104] - rs1[7:0], 2^{8}-1 )
    qv[119:112] = min( qx[119:112] - rs1[7:0], 2^{8}-1 )
    qv[127:120] = min( qx[127:120] - rs1[7:0], 2^{8}-1 )
\end{lstlisting}


\textbf{Schedule}\\
\begin{longtable}[c]{|l|l|l|}
    \hline
    \rowcolor{lightgray}
    \textbf{Operands} & \textbf{Use} & \textbf{Def} \\\hline
    \endfirsthead
    \hline
    \rowcolor{lightgray}
     \textbf{Operands} & \textbf{Use} & \textbf{Def} \\\hline
    \endhead
    \hline
    \endfoot

  qv & - & 1 \\\hline
    qx & 1 & - \\\hline
    rs1 & 1 & -

\end{longtable}

\textbf{Format}\\

\begin{figure*}[!htbp]
 {\setlength{\tabcolsep}{1pt}
 \begin{center}
\begin{tabular}{@{}F@{}W@{}I@{}F@{}F@{}I@{}I@{}F@{}E@{}O}
\instbitrange{31}{29}&
\instbitrange{28}{27}&
\instbit{26}&
\instbitrange{25}{23}&
\instbitrange{22}{20}&
\instbit{19}&
\instbit{18}&
\instbitrange{17}{15}&
\instbitrange{14}{7}&
\instbitrange{6}{0} \\
\hline
\multicolumn{1}{|c|}{{\scriptsize qx[2:0]}} &
\multicolumn{1}{c|}{00} &
\multicolumn{1}{c|}{{\scriptsize sat[0]}} &
\multicolumn{1}{c|}{1X1} &
\multicolumn{1}{c|}{{\scriptsize qv[2:0]}} &
\multicolumn{1}{c|}{X} &
\multicolumn{1}{c|}{{\scriptsize rs1[4]}} &
\multicolumn{1}{c|}{{\scriptsize rs1[2:0]}} &
\multicolumn{1}{c|}{01XX101X} &
\multicolumn{1}{c|}{0011011}\\
\hline
3& 2& 1& 3& 3& 1& 1& 3& 8& 7 \\
\end{tabular}
 \end{center}
 }
 \vspace{-0.1in}
 \caption[]{ESP.VSSUBS.U8} \end{figure*}


\newpage
\subsubsection{ESP.VSUB.S16}\label{instruction:ESPVSUBS16}
\textbf{Syntax}\\
ESP.VSUB.S16 qv,qx,qy,sat

\textbf{Description}\\
This instruction performs a vector subtraction on 16-bit data. Registers \lstinline{qx} and \lstinline{qy} are the minuend and the subtrahend respectively. Then, the 8 results obtained from the calculation are saturated and then written into the \lstinline{qv} register.

\textbf{Operation}\\
\begin{lstlisting}
  qv[ 15:  0] = min(max(qx[ 15:  0] - qy[ 15:  0], -2^{15}), 2^{15}-1)
  qv[ 31: 16] = min(max(qx[ 31: 16] - qy[ 31: 16], -2^{15}), 2^{15}-1)
  ...
\end{lstlisting}


\textbf{Schedule}\\
\begin{longtable}[c]{|l|l|l|}
    \hline
    \rowcolor{lightgray}
    \textbf{Operands} & \textbf{Use} & \textbf{Def} \\\hline
    \endfirsthead
    \hline
    \rowcolor{lightgray}
     \textbf{Operands} & \textbf{Use} & \textbf{Def} \\\hline
    \endhead
    \hline
    \endfoot

  qv & - & 1 \\\hline
    qx & 1 & - \\\hline
    qy & 1 & -

\end{longtable}

\textbf{Format}\\

\begin{figure*}[!htbp]
 {\setlength{\tabcolsep}{1pt}
 \begin{center}
\begin{tabular}{@{}F@{}F@{}F@{}F@{}F@{}I@{}T@{}O}
\instbitrange{31}{29}&
\instbitrange{28}{26}&
\instbitrange{25}{23}&
\instbitrange{22}{20}&
\instbitrange{19}{17}&
\instbit{16}&
\instbitrange{15}{7}&
\instbitrange{6}{0} \\
\hline
\multicolumn{1}{|c|}{{\scriptsize qx[2:0]}} &
\multicolumn{1}{c|}{{\scriptsize qy[2:0]}} &
\multicolumn{1}{c|}{1X1} &
\multicolumn{1}{c|}{{\scriptsize qv[2:0]}} &
\multicolumn{1}{c|}{X11} &
\multicolumn{1}{c|}{{\scriptsize sat[0]}} &
\multicolumn{1}{c|}{001XX111X} &
\multicolumn{1}{c|}{0011011}\\
\hline
3& 3& 3& 3& 3& 1& 9& 7 \\
\end{tabular}
 \end{center}
 }
 \vspace{-0.1in}
 \caption[]{ESP.VSUB.S16} \end{figure*}


\newpage
\subsubsection{ESP.VSUB.S16.LD.INCP}\label{instruction:ESPVSUBS16LDINCP}
\textbf{Syntax}\\
ESP.VSUB.S16.LD.INCP qu,rs1,qv,qx,qy,sat

\textbf{Description}\\
This instruction performs a vector subtraction on 16-bit data. Registers \lstinline{qx} and \lstinline{qy} are the subtrahend and the minuend, respectively. Then, the 8 results obtained from the calculation are saturated and then written into the register \lstinline{qv}. \\ During the operation, the 16-byte data is loaded from the memory into the register \lstinline{qu}. After the access, the value in the register \lstinline{rs1} is incremented by 16.

\textbf{Operation}\\
\begin{lstlisting}
  qv[ 15:  0] = min(max(qx[ 15:  0] - qy[ 15:  0], -2^{15}), 2^{15}-1)
  qv[ 31: 16] = min(max(qx[ 31: 16] - qy[ 31: 16], -2^{15}), 2^{15}-1)
  ...

  qu[127:0] = load128(rs1[31:0])
  rs1[31:0] = rs1[31:0] + 16
\end{lstlisting}


\textbf{Schedule}\\
\begin{longtable}[c]{|l|l|l|}
    \hline
    \rowcolor{lightgray}
    \textbf{Operands} & \textbf{Use} & \textbf{Def} \\\hline
    \endfirsthead
    \hline
    \rowcolor{lightgray}
     \textbf{Operands} & \textbf{Use} & \textbf{Def} \\\hline
    \endhead
    \hline
    \endfoot

  qv & - & 1 \\\hline
    qx & 1 & - \\\hline
    qy & 1 & - \\\hline
    qu & - & 2 \\\hline
    rs1 & 1 & 1

\end{longtable}

\textbf{Format}\\

\begin{figure*}[!htbp]
 {\setlength{\tabcolsep}{1pt}
 \begin{center}
\begin{tabular}{@{}F@{}F@{}F@{}F@{}I@{}I@{}F@{}W@{}F@{}F@{}O}
\instbitrange{31}{29}&
\instbitrange{28}{26}&
\instbitrange{25}{23}&
\instbitrange{22}{20}&
\instbit{19}&
\instbit{18}&
\instbitrange{17}{15}&
\instbitrange{14}{13}&
\instbitrange{12}{10}&
\instbitrange{9}{7}&
\instbitrange{6}{0} \\
\hline
\multicolumn{1}{|c|}{{\scriptsize qx[2:0]}} &
\multicolumn{1}{c|}{{\scriptsize qy[2:0]}} &
\multicolumn{1}{c|}{011} &
\multicolumn{1}{c|}{{\scriptsize qv[2:0]}} &
\multicolumn{1}{c|}{{\scriptsize sat[0]}} &
\multicolumn{1}{c|}{{\scriptsize rs1[4]}} &
\multicolumn{1}{c|}{{\scriptsize rs1[2:0]}} &
\multicolumn{1}{c|}{10} &
\multicolumn{1}{c|}{{\scriptsize qu[2:0]}} &
\multicolumn{1}{c|}{101} &
\multicolumn{1}{c|}{0011111}\\
\hline
3& 3& 3& 3& 1& 1& 3& 2& 3& 3& 7 \\
\end{tabular}
 \end{center}
 }
 \vspace{-0.1in}
 \caption[]{ESP.VSUB.S16.LD.INCP} \end{figure*}


\newpage
\subsubsection{ESP.VSUB.S16.ST.INCP}\label{instruction:ESPVSUBS16STINCP}
\textbf{Syntax}\\
ESP.VSUB.S16.ST.INCP qu,rs1,qv,qx,qy,sat

\textbf{Description}\\
This instruction performs a vector subtraction on 16-bit data. Registers \lstinline{qx} and \lstinline{qy} are the subtrahend and the minuend respectively. Then, the 8 results obtained from the calculation are saturated and then written into the register \lstinline{qv}. \\ During the operation, the value in the register \lstinline{qu} is stored to memory. After the access, the value in the register \lstinline{rs1} is incremented by 16.

\textbf{Operation}\\
\begin{lstlisting}
  qv[ 15:  0] = min(max(qx[ 15:  0] - qy[ 15:  0], -2^{15}), 2^{15}-1)
  qv[ 31: 16] = min(max(qx[ 31: 16] - qy[ 31: 16], -2^{15}), 2^{15}-1)
  ...

  qu[127:0] => store128(rs1[31:0])
  rs1[31:0] = rs1[31:0] + 16
\end{lstlisting}


\textbf{Schedule}\\
\begin{longtable}[c]{|l|l|l|}
    \hline
    \rowcolor{lightgray}
    \textbf{Operands} & \textbf{Use} & \textbf{Def} \\\hline
    \endfirsthead
    \hline
    \rowcolor{lightgray}
     \textbf{Operands} & \textbf{Use} & \textbf{Def} \\\hline
    \endhead
    \hline
    \endfoot

  qv & - & 1 \\\hline
    qx & 1 & - \\\hline
    qy & 1 & - \\\hline
    qu & 2 & - \\\hline
    rs1 & 1 & 1

\end{longtable}

\textbf{Format}\\

\begin{figure*}[!htbp]
 {\setlength{\tabcolsep}{1pt}
 \begin{center}
\begin{tabular}{@{}F@{}F@{}F@{}F@{}I@{}I@{}F@{}W@{}F@{}F@{}O}
\instbitrange{31}{29}&
\instbitrange{28}{26}&
\instbitrange{25}{23}&
\instbitrange{22}{20}&
\instbit{19}&
\instbit{18}&
\instbitrange{17}{15}&
\instbitrange{14}{13}&
\instbitrange{12}{10}&
\instbitrange{9}{7}&
\instbitrange{6}{0} \\
\hline
\multicolumn{1}{|c|}{{\scriptsize qx[2:0]}} &
\multicolumn{1}{c|}{{\scriptsize qy[2:0]}} &
\multicolumn{1}{c|}{111} &
\multicolumn{1}{c|}{{\scriptsize qv[2:0]}} &
\multicolumn{1}{c|}{{\scriptsize sat[0]}} &
\multicolumn{1}{c|}{{\scriptsize rs1[4]}} &
\multicolumn{1}{c|}{{\scriptsize rs1[2:0]}} &
\multicolumn{1}{c|}{10} &
\multicolumn{1}{c|}{{\scriptsize qu[2:0]}} &
\multicolumn{1}{c|}{101} &
\multicolumn{1}{c|}{0011111}\\
\hline
3& 3& 3& 3& 1& 1& 3& 2& 3& 3& 7 \\
\end{tabular}
 \end{center}
 }
 \vspace{-0.1in}
 \caption[]{ESP.VSUB.S16.ST.INCP} \end{figure*}


\newpage
\subsubsection{ESP.VSUB.S32}\label{instruction:ESPVSUBS32}
\textbf{Syntax}\\
ESP.VSUB.S32 qv,qx,qy,sat

\textbf{Description}\\
This instruction performs a vector subtraction on 32-bit data. Registers \lstinline{qx} and \lstinline{qy} are the minuend and the subtrahend respectively. Then, the 4 results obtained from the calculation are saturated and then written into the register \lstinline{qv}.

\textbf{Operation}\\
\begin{lstlisting}
  qv[ 31:  0] = min(max(qx[ 31:  0] - qy[ 31:  0], -2^{31}), 2^{31}-1)
  qv[ 63: 32] = min(max(qx[ 63: 32] - qy[ 63: 32], -2^{31}), 2^{31}-1)
  ...
\end{lstlisting}


\textbf{Schedule}\\
\begin{longtable}[c]{|l|l|l|}
    \hline
    \rowcolor{lightgray}
    \textbf{Operands} & \textbf{Use} & \textbf{Def} \\\hline
    \endfirsthead
    \hline
    \rowcolor{lightgray}
     \textbf{Operands} & \textbf{Use} & \textbf{Def} \\\hline
    \endhead
    \hline
    \endfoot

  qv & - & 1 \\\hline
    qx & 1 & - \\\hline
    qy & 1 & -

\end{longtable}

\textbf{Format}\\

\begin{figure*}[!htbp]
 {\setlength{\tabcolsep}{1pt}
 \begin{center}
\begin{tabular}{@{}F@{}F@{}F@{}F@{}W@{}I@{}T@{}O}
\instbitrange{31}{29}&
\instbitrange{28}{26}&
\instbitrange{25}{23}&
\instbitrange{22}{20}&
\instbitrange{19}{18}&
\instbit{17}&
\instbitrange{16}{7}&
\instbitrange{6}{0} \\
\hline
\multicolumn{1}{|c|}{{\scriptsize qx[2:0]}} &
\multicolumn{1}{c|}{{\scriptsize qy[2:0]}} &
\multicolumn{1}{c|}{1X1} &
\multicolumn{1}{c|}{{\scriptsize qv[2:0]}} &
\multicolumn{1}{c|}{X1} &
\multicolumn{1}{c|}{{\scriptsize sat[0]}} &
\multicolumn{1}{c|}{X101XX111X} &
\multicolumn{1}{c|}{0011011}\\
\hline
3& 3& 3& 3& 2& 1& 10& 7 \\
\end{tabular}
 \end{center}
 }
 \vspace{-0.1in}
 \caption[]{ESP.VSUB.S32} \end{figure*}


\newpage
\subsubsection{ESP.VSUB.S32.LD.INCP}\label{instruction:ESPVSUBS32LDINCP}
\textbf{Syntax}\\
ESP.VSUB.S32.LD.INCP qu,rs1,qv,qx,qy,sat

\textbf{Description}\\
This instruction performs a vector subtraction on 32-bit data. Registers \lstinline{qx} and \lstinline{qy} are the subtrahend and the minuend, respectively. Then, the 4 results obtained from the calculation are saturated and then written into the register \lstinline{qv}. \\ During the operation, the 16-byte data is loaded from the memory into the register \lstinline{qu}. After the access, the value in the register \lstinline{rs1} is incremented by 16.

\textbf{Operation}\\
\begin{lstlisting}
  qv[31: 0] = min(max(qx[31: 0] - qy[31: 0], -2^{31}), 2^{31}-1)
  qv[63:32] = min(max(qx[63:32] - qy[63:32], -2^{31}), 2^{31}-1)
  ...

  qu[127:0] = load128(rs1[31:0])
  rs1[31:0] = rs1[31:0] + 16
\end{lstlisting}


\textbf{Schedule}\\
\begin{longtable}[c]{|l|l|l|}
    \hline
    \rowcolor{lightgray}
    \textbf{Operands} & \textbf{Use} & \textbf{Def} \\\hline
    \endfirsthead
    \hline
    \rowcolor{lightgray}
     \textbf{Operands} & \textbf{Use} & \textbf{Def} \\\hline
    \endhead
    \hline
    \endfoot

  qv & - & 1 \\\hline
    qx & 1 & - \\\hline
    qy & 1 & - \\\hline
    qu & - & 2 \\\hline
    rs1 & 1 & 1

\end{longtable}

\textbf{Format}\\

\begin{figure*}[!htbp]
 {\setlength{\tabcolsep}{1pt}
 \begin{center}
\begin{tabular}{@{}F@{}F@{}F@{}F@{}I@{}I@{}F@{}W@{}F@{}F@{}O}
\instbitrange{31}{29}&
\instbitrange{28}{26}&
\instbitrange{25}{23}&
\instbitrange{22}{20}&
\instbit{19}&
\instbit{18}&
\instbitrange{17}{15}&
\instbitrange{14}{13}&
\instbitrange{12}{10}&
\instbitrange{9}{7}&
\instbitrange{6}{0} \\
\hline
\multicolumn{1}{|c|}{{\scriptsize qx[2:0]}} &
\multicolumn{1}{c|}{{\scriptsize qy[2:0]}} &
\multicolumn{1}{c|}{01X} &
\multicolumn{1}{c|}{{\scriptsize qv[2:0]}} &
\multicolumn{1}{c|}{{\scriptsize sat[0]}} &
\multicolumn{1}{c|}{{\scriptsize rs1[4]}} &
\multicolumn{1}{c|}{{\scriptsize rs1[2:0]}} &
\multicolumn{1}{c|}{10} &
\multicolumn{1}{c|}{{\scriptsize qu[2:0]}} &
\multicolumn{1}{c|}{011} &
\multicolumn{1}{c|}{0011111}\\
\hline
3& 3& 3& 3& 1& 1& 3& 2& 3& 3& 7 \\
\end{tabular}
 \end{center}
 }
 \vspace{-0.1in}
 \caption[]{ESP.VSUB.S32.LD.INCP} \end{figure*}


\newpage
\subsubsection{ESP.VSUB.S32.ST.INCP}\label{instruction:ESPVSUBS32STINCP}
\textbf{Syntax}\\
ESP.VSUB.S32.ST.INCP qu,rs1,qv,qx,qy,sat

\textbf{Description}\\
This instruction performs a vector subtraction on 32-bit data. Registers \lstinline{qx} and \lstinline{qy} are the subtrahend and the minuend, respectively. Then, the 4 results obtained from the calculation are saturated and then written into the register \lstinline{qv}. \\ During the operation, the value in the register \lstinline{qu} is stored in memory. After the access, the value in the register \lstinline{rs1} is incremented by 16.

\textbf{Operation}\\
\begin{lstlisting}
  qv[ 31:  0] = min(max(qx[ 31:  0] - qy[ 31:  0], -2^{31}), 2^{31}-1)
  qv[ 63: 32] = min(max(qx[ 63: 32] - qy[ 63: 32], -2^{31}), 2^{31}-1)
  ...

  qu[127:0] => store128(rs1[31:0])
  rs1[31:0] = rs1[31:0] + 16
\end{lstlisting}


\textbf{Schedule}\\
\begin{longtable}[c]{|l|l|l|}
    \hline
    \rowcolor{lightgray}
    \textbf{Operands} & \textbf{Use} & \textbf{Def} \\\hline
    \endfirsthead
    \hline
    \rowcolor{lightgray}
     \textbf{Operands} & \textbf{Use} & \textbf{Def} \\\hline
    \endhead
    \hline
    \endfoot

  qv & - & 1 \\\hline
    qx & 1 & - \\\hline
    qy & 1 & - \\\hline
    qu & 2 & - \\\hline
    rs1 & 1 & 1

\end{longtable}

\textbf{Format}\\

\begin{figure*}[!htbp]
 {\setlength{\tabcolsep}{1pt}
 \begin{center}
\begin{tabular}{@{}F@{}F@{}F@{}F@{}I@{}I@{}F@{}W@{}F@{}F@{}O}
\instbitrange{31}{29}&
\instbitrange{28}{26}&
\instbitrange{25}{23}&
\instbitrange{22}{20}&
\instbit{19}&
\instbit{18}&
\instbitrange{17}{15}&
\instbitrange{14}{13}&
\instbitrange{12}{10}&
\instbitrange{9}{7}&
\instbitrange{6}{0} \\
\hline
\multicolumn{1}{|c|}{{\scriptsize qx[2:0]}} &
\multicolumn{1}{c|}{{\scriptsize qy[2:0]}} &
\multicolumn{1}{c|}{11X} &
\multicolumn{1}{c|}{{\scriptsize qv[2:0]}} &
\multicolumn{1}{c|}{{\scriptsize sat[0]}} &
\multicolumn{1}{c|}{{\scriptsize rs1[4]}} &
\multicolumn{1}{c|}{{\scriptsize rs1[2:0]}} &
\multicolumn{1}{c|}{10} &
\multicolumn{1}{c|}{{\scriptsize qu[2:0]}} &
\multicolumn{1}{c|}{011} &
\multicolumn{1}{c|}{0011111}\\
\hline
3& 3& 3& 3& 1& 1& 3& 2& 3& 3& 7 \\
\end{tabular}
 \end{center}
 }
 \vspace{-0.1in}
 \caption[]{ESP.VSUB.S32.ST.INCP} \end{figure*}


\newpage
\subsubsection{ESP.VSUB.S8}\label{instruction:ESPVSUBS8}
\textbf{Syntax}\\
ESP.VSUB.S8 qv,qx,qy,sat

\textbf{Description}\\
This instruction performs a vector subtraction on 8-bit data. Registers \lstinline{qx} and \lstinline{qy} are the minuend and the subtrahend respectively. Then, the 16 results obtained from the calculation are saturated and then written into the register \lstinline{qv}.

\textbf{Operation}\\
\begin{lstlisting}
  qv[  7:  0] = min(max(qx[  7:  0] - qy[  7:  0], -2^{7}), 2^{7}-1)
  qv[ 15:  8] = min(max(qx[ 15:  8] - qy[ 15:  8], -2^{7}), 2^{7}-1)
  ...
\end{lstlisting}


\textbf{Schedule}\\
\begin{longtable}[c]{|l|l|l|}
    \hline
    \rowcolor{lightgray}
    \textbf{Operands} & \textbf{Use} & \textbf{Def} \\\hline
    \endfirsthead
    \hline
    \rowcolor{lightgray}
     \textbf{Operands} & \textbf{Use} & \textbf{Def} \\\hline
    \endhead
    \hline
    \endfoot

  qv & - & 1 \\\hline
    qx & 1 & - \\\hline
    qy & 1 & - \\\hline

\end{longtable}

\textbf{Format}\\

\begin{figure*}[!htbp]
 {\setlength{\tabcolsep}{1pt}
 \begin{center}
\begin{tabular}{@{}F@{}F@{}F@{}F@{}F@{}I@{}T@{}O}
\instbitrange{31}{29}&
\instbitrange{28}{26}&
\instbitrange{25}{23}&
\instbitrange{22}{20}&
\instbitrange{19}{17}&
\instbit{16}&
\instbitrange{15}{7}&
\instbitrange{6}{0} \\
\hline
\multicolumn{1}{|c|}{{\scriptsize qx[2:0]}} &
\multicolumn{1}{c|}{{\scriptsize qy[2:0]}} &
\multicolumn{1}{c|}{1X1} &
\multicolumn{1}{c|}{{\scriptsize qv[2:0]}} &
\multicolumn{1}{c|}{X01} &
\multicolumn{1}{c|}{{\scriptsize sat[0]}} &
\multicolumn{1}{c|}{001XX111X} &
\multicolumn{1}{c|}{0011011}\\
\hline
3& 3& 3& 3& 3& 1& 9& 7 \\
\end{tabular}
 \end{center}
 }
 \vspace{-0.1in}
 \caption[]{ESP.VSUB.S8} \end{figure*}


\newpage
\subsubsection{ESP.VSUB.S8.LD.INCP}\label{instruction:ESPVSUBS8LDINCP}
\textbf{Syntax}\\
ESP.VSUB.S8.LD.INCP qu,rs1,qv,qx,qy,sat

\textbf{Description}\\
This instruction performs a vector subtraction on 8-bit data. Registers \lstinline{qx} and \lstinline{qy} are the subtrahend and the minuend, respectively. Then, the 16 results obtained from the calculation are saturated and then written into the register \lstinline{qv}. \\ During the operation, the 16-byte data is loaded from the memory into the register \lstinline{qu}. After the access, the value in the register \lstinline{rs1} is incremented by 16.

\textbf{Operation}\\
\begin{lstlisting}
  qv[  7:  0] = min(max(qx[  7:  0] - qy[  7:  0], -2^{7}), 2^{7}-1)
  qv[ 15:  8] = min(max(qx[ 15:  8] - qy[ 15:  8], -2^{7}), 2^{7}-1)
  ...

  qu[127:0] = load128(rs1[31:0])
  rs1[31:0] = rs1[31:0] + 16
\end{lstlisting}


\textbf{Schedule}\\
\begin{longtable}[c]{|l|l|l|}
    \hline
    \rowcolor{lightgray}
    \textbf{Operands} & \textbf{Use} & \textbf{Def} \\\hline
    \endfirsthead
    \hline
    \rowcolor{lightgray}
     \textbf{Operands} & \textbf{Use} & \textbf{Def} \\\hline
    \endhead
    \hline
    \endfoot

  qv & - & 1 \\\hline
    qx & 1 & - \\\hline
    qy & 1 & - \\\hline
    qu & - & 2 \\\hline
    rs1 & 1 & 1

\end{longtable}

\textbf{Format}\\

\begin{figure*}[!htbp]
 {\setlength{\tabcolsep}{1pt}
 \begin{center}
\begin{tabular}{@{}F@{}F@{}F@{}F@{}I@{}I@{}F@{}W@{}F@{}F@{}O}
\instbitrange{31}{29}&
\instbitrange{28}{26}&
\instbitrange{25}{23}&
\instbitrange{22}{20}&
\instbit{19}&
\instbit{18}&
\instbitrange{17}{15}&
\instbitrange{14}{13}&
\instbitrange{12}{10}&
\instbitrange{9}{7}&
\instbitrange{6}{0} \\
\hline
\multicolumn{1}{|c|}{{\scriptsize qx[2:0]}} &
\multicolumn{1}{c|}{{\scriptsize qy[2:0]}} &
\multicolumn{1}{c|}{001} &
\multicolumn{1}{c|}{{\scriptsize qv[2:0]}} &
\multicolumn{1}{c|}{{\scriptsize sat[0]}} &
\multicolumn{1}{c|}{{\scriptsize rs1[4]}} &
\multicolumn{1}{c|}{{\scriptsize rs1[2:0]}} &
\multicolumn{1}{c|}{10} &
\multicolumn{1}{c|}{{\scriptsize qu[2:0]}} &
\multicolumn{1}{c|}{101} &
\multicolumn{1}{c|}{0011111}\\
\hline
3& 3& 3& 3& 1& 1& 3& 2& 3& 3& 7 \\
\end{tabular}
 \end{center}
 }
 \vspace{-0.1in}
 \caption[]{ESP.VSUB.S8.LD.INCP} \end{figure*}


\newpage
\subsubsection{ESP.VSUB.S8.ST.INCP}\label{instruction:ESPVSUBS8STINCP}
\textbf{Syntax}\\
ESP.VSUB.S8.ST.INCP qu,rs1,qv,qx,qy,sat

\textbf{Description}\\
This instruction performs a vector subtraction on 8-bit data. Registers \lstinline{qx} and \lstinline{qy} are the subtrahend and the minuend respectively. Then, the 16 results obtained from the calculation are saturated and then written into the register \lstinline{qv}. \\ During the operation, the value in the register \lstinline{qu} is stored to memory. After the access, the value in the register \lstinline{rs1} is incremented by 16.

\textbf{Operation}\\
\begin{lstlisting}
  qv[  7:  0] = min(max(qx[  7:  0] - qy[  7:  0], -2^{7}), 2^{7}-1)
  qv[ 15:  8] = min(max(qx[ 15:  8] - qy[ 15:  8], -2^{7}), 2^{7}-1)
  ...

  qu[127:0] => store128(rs1[31:0])
  rs1[31:0] = rs1[31:0] + 16
\end{lstlisting}


\textbf{Schedule}\\
\begin{longtable}[c]{|l|l|l|}
    \hline
    \rowcolor{lightgray}
    \textbf{Operands} & \textbf{Use} & \textbf{Def} \\\hline
    \endfirsthead
    \hline
    \rowcolor{lightgray}
     \textbf{Operands} & \textbf{Use} & \textbf{Def} \\\hline
    \endhead
    \hline
    \endfoot

  qv & - & 1 \\\hline
    qx & 1 & - \\\hline
    qy & 1 & - \\\hline
    qu & 2 & - \\\hline
    rs1 & 1 & 1

\end{longtable}

\textbf{Format}\\

\begin{figure*}[!htbp]
 {\setlength{\tabcolsep}{1pt}
 \begin{center}
\begin{tabular}{@{}F@{}F@{}F@{}F@{}I@{}I@{}F@{}W@{}F@{}F@{}O}
\instbitrange{31}{29}&
\instbitrange{28}{26}&
\instbitrange{25}{23}&
\instbitrange{22}{20}&
\instbit{19}&
\instbit{18}&
\instbitrange{17}{15}&
\instbitrange{14}{13}&
\instbitrange{12}{10}&
\instbitrange{9}{7}&
\instbitrange{6}{0} \\
\hline
\multicolumn{1}{|c|}{{\scriptsize qx[2:0]}} &
\multicolumn{1}{c|}{{\scriptsize qy[2:0]}} &
\multicolumn{1}{c|}{101} &
\multicolumn{1}{c|}{{\scriptsize qv[2:0]}} &
\multicolumn{1}{c|}{{\scriptsize sat[0]}} &
\multicolumn{1}{c|}{{\scriptsize rs1[4]}} &
\multicolumn{1}{c|}{{\scriptsize rs1[2:0]}} &
\multicolumn{1}{c|}{10} &
\multicolumn{1}{c|}{{\scriptsize qu[2:0]}} &
\multicolumn{1}{c|}{101} &
\multicolumn{1}{c|}{0011111}\\
\hline
3& 3& 3& 3& 1& 1& 3& 2& 3& 3& 7 \\
\end{tabular}
 \end{center}
 }
 \vspace{-0.1in}
 \caption[]{ESP.VSUB.S8.ST.INCP} \end{figure*}


\newpage
\subsubsection{ESP.VSUB.U16}\label{instruction:ESPVSUBU16}
\textbf{Syntax}\\
ESP.VSUB.U16 qv,qx,qy,sat

\textbf{Description}\\
This instruction performs a vector subtraction on 16-bit unsigned data. Registers \lstinline{qx} and \lstinline{qy} are the minuend and the subtrahend respectively. Then, the 8 results obtained from the calculation are saturated and then written into the \lstinline{qv} register.

\textbf{Operation}\\
\begin{lstlisting}
  qv[ 15:  0] = min( qx[ 15:  0] - qy[ 15:  0], 2^{16}-1)
  qv[ 31: 16] = min( qx[ 31: 16] - qy[ 31: 16], 2^{16}-1)
  ...
\end{lstlisting}


\textbf{Schedule}\\
\begin{longtable}[c]{|l|l|l|}
    \hline
    \rowcolor{lightgray}
    \textbf{Operands} & \textbf{Use} & \textbf{Def} \\\hline
    \endfirsthead
    \hline
    \rowcolor{lightgray}
     \textbf{Operands} & \textbf{Use} & \textbf{Def} \\\hline
    \endhead
    \hline
    \endfoot

  qv & - & 1 \\\hline
    qx & 1 & - \\\hline
    qy & 1 & -

\end{longtable}

\textbf{Format}\\

\begin{figure*}[!htbp]
 {\setlength{\tabcolsep}{1pt}
 \begin{center}
\begin{tabular}{@{}F@{}F@{}F@{}F@{}F@{}I@{}T@{}O}
\instbitrange{31}{29}&
\instbitrange{28}{26}&
\instbitrange{25}{23}&
\instbitrange{22}{20}&
\instbitrange{19}{17}&
\instbit{16}&
\instbitrange{15}{7}&
\instbitrange{6}{0} \\
\hline
\multicolumn{1}{|c|}{{\scriptsize qx[2:0]}} &
\multicolumn{1}{c|}{{\scriptsize qy[2:0]}} &
\multicolumn{1}{c|}{1X1} &
\multicolumn{1}{c|}{{\scriptsize qv[2:0]}} &
\multicolumn{1}{c|}{X10} &
\multicolumn{1}{c|}{{\scriptsize sat[0]}} &
\multicolumn{1}{c|}{001XX111X} &
\multicolumn{1}{c|}{0011011}\\
\hline
3& 3& 3& 3& 3& 1& 9& 7 \\
\end{tabular}
 \end{center}
 }
 \vspace{-0.1in}
 \caption[]{ESP.VSUB.U16} \end{figure*}


\newpage
\subsubsection{ESP.VSUB.U16.LD.INCP}\label{instruction:ESPVSUBU16LDINCP}
\textbf{Syntax}\\
ESP.VSUB.U16.LD.INCP qu,rs1,qv,qx,qy,sat

\textbf{Description}\\
This instruction performs a vector subtraction on 16-bit unsigned data. Registers \lstinline{qx} and \lstinline{qy} are the subtrahend and the minuend, respectively. Then, the 8 results obtained from the calculation are saturated and then written into the register \lstinline{qv}. \\ During the operation, the 16-byte data is loaded from the memory into the register \lstinline{qu}. After the access, the value in the register \lstinline{rs1} is incremented by 16.

\textbf{Operation}\\
\begin{lstlisting}
  qv[ 15:  0] = min(qx[ 15:  0] - qy[ 15:  0], 2^{16}-1)
  qv[ 31: 16] = min(qx[ 31: 16] - qy[ 31: 16], 2^{16}-1)
  ...

  qu[127:0] = load128(rs1[31:0])
  rs1[31:0] = rs1[31:0] + 16
\end{lstlisting}


\textbf{Schedule}\\
\begin{longtable}[c]{|l|l|l|}
    \hline
    \rowcolor{lightgray}
    \textbf{Operands} & \textbf{Use} & \textbf{Def} \\\hline
    \endfirsthead
    \hline
    \rowcolor{lightgray}
     \textbf{Operands} & \textbf{Use} & \textbf{Def} \\\hline
    \endhead
    \hline
    \endfoot

  qv & - & 1 \\\hline
    qx & 1 & - \\\hline
    qy & 1 & - \\\hline
    qu & - & 2 \\\hline
    rs1 & 1 & 1

\end{longtable}

\textbf{Format}\\

\begin{figure*}[!htbp]
 {\setlength{\tabcolsep}{1pt}
 \begin{center}
\begin{tabular}{@{}F@{}F@{}F@{}F@{}I@{}I@{}F@{}W@{}F@{}F@{}O}
\instbitrange{31}{29}&
\instbitrange{28}{26}&
\instbitrange{25}{23}&
\instbitrange{22}{20}&
\instbit{19}&
\instbit{18}&
\instbitrange{17}{15}&
\instbitrange{14}{13}&
\instbitrange{12}{10}&
\instbitrange{9}{7}&
\instbitrange{6}{0} \\
\hline
\multicolumn{1}{|c|}{{\scriptsize qx[2:0]}} &
\multicolumn{1}{c|}{{\scriptsize qy[2:0]}} &
\multicolumn{1}{c|}{010} &
\multicolumn{1}{c|}{{\scriptsize qv[2:0]}} &
\multicolumn{1}{c|}{{\scriptsize sat[0]}} &
\multicolumn{1}{c|}{{\scriptsize rs1[4]}} &
\multicolumn{1}{c|}{{\scriptsize rs1[2:0]}} &
\multicolumn{1}{c|}{10} &
\multicolumn{1}{c|}{{\scriptsize qu[2:0]}} &
\multicolumn{1}{c|}{101} &
\multicolumn{1}{c|}{0011111}\\
\hline
3& 3& 3& 3& 1& 1& 3& 2& 3& 3& 7 \\
\end{tabular}
 \end{center}
 }
 \vspace{-0.1in}
 \caption[]{ESP.VSUB.U16.LD.INCP} \end{figure*}


\newpage
\subsubsection{ESP.VSUB.U16.ST.INCP}\label{instruction:ESPVSUBU16STINCP}
\textbf{Syntax}\\
ESP.VSUB.U16.ST.INCP qu,rs1,qv,qx,qy,sat

\textbf{Description}\\
This instruction performs a vector subtraction on 16-bit unsigned data. Registers \lstinline{qx} and \lstinline{qy} are the subtrahend and the minuend, respectively. Then, the 8 results obtained from the calculation are saturated and then written into the register \lstinline{qv}. \\ During the operation, the value in the register \lstinline{qu} is stored to memory. After the access, the value in the register \lstinline{rs1} is incremented by 16.

\textbf{Operation}\\
\begin{lstlisting}
  qv[ 15:  0] = min(qx[ 15:  0] - qy[ 15:  0], 2^{16}-1)
  qv[ 31: 16] = min(qx[ 31: 16] - qy[ 31: 16], 2^{16}-1)
  ...

  qu[127:0] => store128(rs1[31:0])
  rs1[31:0] = rs1[31:0] + 16
\end{lstlisting}


\textbf{Schedule}\\
\begin{longtable}[c]{|l|l|l|}
    \hline
    \rowcolor{lightgray}
    \textbf{Operands} & \textbf{Use} & \textbf{Def} \\\hline
    \endfirsthead
    \hline
    \rowcolor{lightgray}
     \textbf{Operands} & \textbf{Use} & \textbf{Def} \\\hline
    \endhead
    \hline
    \endfoot

  qv & - & 1 \\\hline
    qx & 1 & - \\\hline
    qy & 1 & - \\\hline
    qu & 2 & - \\\hline
    rs1 & 1 & 1

\end{longtable}

\textbf{Format}\\

\begin{figure*}[!htbp]
 {\setlength{\tabcolsep}{1pt}
 \begin{center}
\begin{tabular}{@{}F@{}F@{}F@{}F@{}I@{}I@{}F@{}W@{}F@{}F@{}O}
\instbitrange{31}{29}&
\instbitrange{28}{26}&
\instbitrange{25}{23}&
\instbitrange{22}{20}&
\instbit{19}&
\instbit{18}&
\instbitrange{17}{15}&
\instbitrange{14}{13}&
\instbitrange{12}{10}&
\instbitrange{9}{7}&
\instbitrange{6}{0} \\
\hline
\multicolumn{1}{|c|}{{\scriptsize qx[2:0]}} &
\multicolumn{1}{c|}{{\scriptsize qy[2:0]}} &
\multicolumn{1}{c|}{110} &
\multicolumn{1}{c|}{{\scriptsize qv[2:0]}} &
\multicolumn{1}{c|}{{\scriptsize sat[0]}} &
\multicolumn{1}{c|}{{\scriptsize rs1[4]}} &
\multicolumn{1}{c|}{{\scriptsize rs1[2:0]}} &
\multicolumn{1}{c|}{10} &
\multicolumn{1}{c|}{{\scriptsize qu[2:0]}} &
\multicolumn{1}{c|}{101} &
\multicolumn{1}{c|}{0011111}\\
\hline
3& 3& 3& 3& 1& 1& 3& 2& 3& 3& 7 \\
\end{tabular}
 \end{center}
 }
 \vspace{-0.1in}
 \caption[]{ESP.VSUB.U16.ST.INCP} \end{figure*}


\newpage
\subsubsection{ESP.VSUB.U32}\label{instruction:ESPVSUBU32}
\textbf{Syntax}\\
ESP.VSUB.U32 qv,qx,qy,sat

\textbf{Description}\\
This instruction performs a vector subtraction on 32-bit unsigned data. Registers \lstinline{qx} and \lstinline{qy} are the minuend and the subtrahend respectively. Then, the 4 results obtained from the calculation are saturated and then written into the register \lstinline{qv}.

\textbf{Operation}\\
\begin{lstlisting}
  qv[ 31:  0] = min(qx[ 31:  0] - qy[ 31:  0], 2^{32}-1)
  qv[ 63: 32] = min(qx[ 63: 32] - qy[ 63: 32], 2^{32}-1)
  ...
\end{lstlisting}

\textbf{Schedule}\\
\begin{longtable}[c]{|l|l|l|}
    \hline
    \rowcolor{lightgray}
    \textbf{Operands} & \textbf{Use} & \textbf{Def} \\\hline
    \endfirsthead
    \hline
    \rowcolor{lightgray}
     \textbf{Operands} & \textbf{Use} & \textbf{Def} \\\hline
    \endhead
    \hline
    \endfoot

  qv & - & 1 \\\hline
    qx & 1 & - \\\hline
    qy & 1 & -

\end{longtable}

\textbf{Format}\\

\begin{figure*}[!htbp]
 {\setlength{\tabcolsep}{1pt}
 \begin{center}
\begin{tabular}{@{}F@{}F@{}F@{}F@{}W@{}I@{}T@{}O}
\instbitrange{31}{29}&
\instbitrange{28}{26}&
\instbitrange{25}{23}&
\instbitrange{22}{20}&
\instbitrange{19}{18}&
\instbit{17}&
\instbitrange{16}{7}&
\instbitrange{6}{0} \\
\hline
\multicolumn{1}{|c|}{{\scriptsize qx[2:0]}} &
\multicolumn{1}{c|}{{\scriptsize qy[2:0]}} &
\multicolumn{1}{c|}{1X1} &
\multicolumn{1}{c|}{{\scriptsize qv[2:0]}} &
\multicolumn{1}{c|}{X0} &
\multicolumn{1}{c|}{{\scriptsize sat[0]}} &
\multicolumn{1}{c|}{X101XX111X} &
\multicolumn{1}{c|}{0011011}\\
\hline
3& 3& 3& 3& 2& 1& 10& 7 \\
\end{tabular}
 \end{center}
 }
 \vspace{-0.1in}
 \caption[]{ESP.VSUB.U32} \end{figure*}


\newpage
\subsubsection{ESP.VSUB.U32.LD.INCP}\label{instruction:ESPVSUBU32LDINCP}
\textbf{Syntax}\\
ESP.VSUB.U32.LD.INCP qu,rs1,qv,qx,qy,sat

\textbf{Description}\\
This instruction performs a vector subtraction on 32-bit unsigned data. Registers \lstinline{qx} and \lstinline{qy} are the subtrahend and the minuend, respectively. Then, the 4 results obtained from the calculation are saturated and then written into the register \lstinline{qv}. \\ During the operation, the 16-byte data is loaded from the memory into the register \lstinline{qu}. After the access, the value in the register \lstinline{rs1} is incremented by 16.

\textbf{Operation}\\
\begin{lstlisting}
  qv[31: 0] = min(qx[31: 0] - qy[31: 0], 2^{32}-1)
  qv[63:32] = min(qx[63:32] - qy[63:32], 2^{32}-1)
  ...

  qu[127:0] = load128(rs1[31:0])
  rs1[31:0] = rs1[31:0] + 16
\end{lstlisting}

\textbf{Schedule}\\
\begin{longtable}[c]{|l|l|l|}
    \hline
    \rowcolor{lightgray}
    \textbf{Operands} & \textbf{Use} & \textbf{Def} \\\hline
    \endfirsthead
    \hline
    \rowcolor{lightgray}
     \textbf{Operands} & \textbf{Use} & \textbf{Def} \\\hline
    \endhead
    \hline
    \endfoot

  qv & - & 1 \\\hline
    qx & 1 & - \\\hline
    qy & 1 & - \\\hline
    qu & - & 2 \\\hline
    rs1 & 1 & 1

\end{longtable}

\textbf{Format}\\

\begin{figure*}[!htbp]
 {\setlength{\tabcolsep}{1pt}
 \begin{center}
\begin{tabular}{@{}F@{}F@{}F@{}F@{}I@{}I@{}F@{}W@{}F@{}F@{}O}
\instbitrange{31}{29}&
\instbitrange{28}{26}&
\instbitrange{25}{23}&
\instbitrange{22}{20}&
\instbit{19}&
\instbit{18}&
\instbitrange{17}{15}&
\instbitrange{14}{13}&
\instbitrange{12}{10}&
\instbitrange{9}{7}&
\instbitrange{6}{0} \\
\hline
\multicolumn{1}{|c|}{{\scriptsize qx[2:0]}} &
\multicolumn{1}{c|}{{\scriptsize qy[2:0]}} &
\multicolumn{1}{c|}{00X} &
\multicolumn{1}{c|}{{\scriptsize qv[2:0]}} &
\multicolumn{1}{c|}{{\scriptsize sat[0]}} &
\multicolumn{1}{c|}{{\scriptsize rs1[4]}} &
\multicolumn{1}{c|}{{\scriptsize rs1[2:0]}} &
\multicolumn{1}{c|}{10} &
\multicolumn{1}{c|}{{\scriptsize qu[2:0]}} &
\multicolumn{1}{c|}{011} &
\multicolumn{1}{c|}{0011111}\\
\hline
3& 3& 3& 3& 1& 1& 3& 2& 3& 3& 7 \\
\end{tabular}
 \end{center}
 }
 \vspace{-0.1in}
 \caption[]{ESP.VSUB.U32.LD.INCP} \end{figure*}


\newpage
\subsubsection{ESP.VSUB.U32.ST.INCP}\label{instruction:ESPVSUBU32STINCP}
\textbf{Syntax}\\
ESP.VSUB.U32.ST.INCP qu,rs1,qv,qx,qy,sat

\textbf{Description}\\
This instruction performs a vector subtraction on 32-bit unsigned data. Registers \lstinline{qx} and \lstinline{qy} are the subtrahend and the minuend, respectively. Then, the 4 results obtained from the calculation are saturated and then written into the register \lstinline{qv}. \\ During the operation, the value in the register \lstinline{qu} is stored to memory. After the access, the value in the register \lstinline{rs1} is incremented by 16.

\textbf{Operation}\\
\begin{lstlisting}
  qv[ 31:  0] = min(qx[ 31:  0] - qy[ 31:  0], 2^{32}-1)
  qv[ 63: 32] = min(qx[ 63: 32] - qy[ 63: 32], 2^{32}-1)
  ...

  qu[127:0] => store128(rs1[31:0])
  rs1[31:0] = rs1[31:0] + 16
\end{lstlisting}

\textbf{Schedule}\\
\begin{longtable}[c]{|l|l|l|}
    \hline
    \rowcolor{lightgray}
    \textbf{Operands} & \textbf{Use} & \textbf{Def} \\\hline
    \endfirsthead
    \hline
    \rowcolor{lightgray}
     \textbf{Operands} & \textbf{Use} & \textbf{Def} \\\hline
    \endhead
    \hline
    \endfoot

  qv & - & 1 \\\hline
    qx & 1 & - \\\hline
    qy & 1 & - \\\hline
    qu & 2 & - \\\hline
    rs1 & 1 & 1

\end{longtable}

\textbf{Format}\\

\begin{figure*}[!htbp]
 {\setlength{\tabcolsep}{1pt}
 \begin{center}
\begin{tabular}{@{}F@{}F@{}F@{}F@{}I@{}I@{}F@{}W@{}F@{}F@{}O}
\instbitrange{31}{29}&
\instbitrange{28}{26}&
\instbitrange{25}{23}&
\instbitrange{22}{20}&
\instbit{19}&
\instbit{18}&
\instbitrange{17}{15}&
\instbitrange{14}{13}&
\instbitrange{12}{10}&
\instbitrange{9}{7}&
\instbitrange{6}{0} \\
\hline
\multicolumn{1}{|c|}{{\scriptsize qx[2:0]}} &
\multicolumn{1}{c|}{{\scriptsize qy[2:0]}} &
\multicolumn{1}{c|}{10X} &
\multicolumn{1}{c|}{{\scriptsize qv[2:0]}} &
\multicolumn{1}{c|}{{\scriptsize sat[0]}} &
\multicolumn{1}{c|}{{\scriptsize rs1[4]}} &
\multicolumn{1}{c|}{{\scriptsize rs1[2:0]}} &
\multicolumn{1}{c|}{10} &
\multicolumn{1}{c|}{{\scriptsize qu[2:0]}} &
\multicolumn{1}{c|}{011} &
\multicolumn{1}{c|}{0011111}\\
\hline
3& 3& 3& 3& 1& 1& 3& 2& 3& 3& 7 \\
\end{tabular}
 \end{center}
 }
 \vspace{-0.1in}
 \caption[]{ESP.VSUB.U32.ST.INCP} \end{figure*}


\newpage
\subsubsection{ESP.VSUB.U8}\label{instruction:ESPVSUBU8}
\textbf{Syntax}\\
ESP.VSUB.U8 qv,qx,qy,sat

\textbf{Description}\\
This instruction performs a vector subtraction on 8-bit unsigned data. Registers \lstinline{qx} and \lstinline{qy} are the minuend and the subtrahend respectively. Then, the 16 results obtained from the calculation are saturated and then written into the register \lstinline{qv}.

\textbf{Operation}\\
\begin{lstlisting}
  qv[  7:  0] = min(qx[  7:  0] - qy[  7:  0], 2^{8}-1)
  qv[ 15:  8] = min(qx[ 15:  8] - qy[ 15:  8], 2^{8}-1)
  ...
\end{lstlisting}

\textbf{Schedule}\\
\begin{longtable}[c]{|l|l|l|}
    \hline
    \rowcolor{lightgray}
    \textbf{Operands} & \textbf{Use} & \textbf{Def} \\\hline
    \endfirsthead
    \hline
    \rowcolor{lightgray}
     \textbf{Operands} & \textbf{Use} & \textbf{Def} \\\hline
    \endhead
    \hline
    \endfoot

  qv & - & 1 \\\hline
    qx & 1 & - \\\hline
    qy & 1 & -

\end{longtable}

\textbf{Format}\\

\begin{figure*}[!htbp]
 {\setlength{\tabcolsep}{1pt}
 \begin{center}
\begin{tabular}{@{}F@{}F@{}F@{}F@{}F@{}I@{}T@{}O}
\instbitrange{31}{29}&
\instbitrange{28}{26}&
\instbitrange{25}{23}&
\instbitrange{22}{20}&
\instbitrange{19}{17}&
\instbit{16}&
\instbitrange{15}{7}&
\instbitrange{6}{0} \\
\hline
\multicolumn{1}{|c|}{{\scriptsize qx[2:0]}} &
\multicolumn{1}{c|}{{\scriptsize qy[2:0]}} &
\multicolumn{1}{c|}{1X1} &
\multicolumn{1}{c|}{{\scriptsize qv[2:0]}} &
\multicolumn{1}{c|}{X00} &
\multicolumn{1}{c|}{{\scriptsize sat[0]}} &
\multicolumn{1}{c|}{001XX111X} &
\multicolumn{1}{c|}{0011011}\\
\hline
3& 3& 3& 3& 3& 1& 9& 7 \\
\end{tabular}
 \end{center}
 }
 \vspace{-0.1in}
 \caption[]{ESP.VSUB.U8} \end{figure*}


\newpage
\subsubsection{ESP.VSUB.U8.LD.INCP}\label{instruction:ESPVSUBU8LDINCP}
\textbf{Syntax}\\
ESP.VSUB.U8.LD.INCP qu,rs1,qv,qx,qy,sat

\textbf{Description}\\
This instruction performs a vector subtraction on 8-bit unsigned data. Registers \lstinline{qx} and \lstinline{qy} are the subtrahend and the minuend, respectively. Then, the 16 results obtained from the calculation are saturated and then written into register \lstinline{qv}. \\ During the operation, the 16-byte data is loaded from the memory into register \lstinline{qu}. After the access, the value in register \lstinline{rs1} is incremented by 16.

\textbf{Operation}\\
\begin{lstlisting}
  qv[  7:  0] = min(qx[  7:  0] - qy[  7:  0], 2^{8}-1)
  qv[ 15:  8] = min(qx[ 15:  8] - qy[ 15:  8], 2^{8}-1)
  ...

  qu[127:0] = load128(rs1[31:0])
  rs1[31:0] = rs1[31:0] + 16
\end{lstlisting}

\textbf{Schedule}\\
\begin{longtable}[c]{|l|l|l|}
    \hline
    \rowcolor{lightgray}
    \textbf{Operands} & \textbf{Use} & \textbf{Def} \\\hline
    \endfirsthead
    \hline
    \rowcolor{lightgray}
     \textbf{Operands} & \textbf{Use} & \textbf{Def} \\\hline
    \endhead
    \hline
    \endfoot

  qv & - & 1 \\\hline
    qx & 1 & - \\\hline
    qy & 1 & - \\\hline
    qu & - & 2 \\\hline
    rs1 & 1 & 1

\end{longtable}

\textbf{Format}\\

\begin{figure*}[!htbp]
 {\setlength{\tabcolsep}{1pt}
 \begin{center}
\begin{tabular}{@{}F@{}F@{}F@{}F@{}I@{}I@{}F@{}W@{}F@{}F@{}O}
\instbitrange{31}{29}&
\instbitrange{28}{26}&
\instbitrange{25}{23}&
\instbitrange{22}{20}&
\instbit{19}&
\instbit{18}&
\instbitrange{17}{15}&
\instbitrange{14}{13}&
\instbitrange{12}{10}&
\instbitrange{9}{7}&
\instbitrange{6}{0} \\
\hline
\multicolumn{1}{|c|}{{\scriptsize qx[2:0]}} &
\multicolumn{1}{c|}{{\scriptsize qy[2:0]}} &
\multicolumn{1}{c|}{000} &
\multicolumn{1}{c|}{{\scriptsize qv[2:0]}} &
\multicolumn{1}{c|}{{\scriptsize sat[0]}} &
\multicolumn{1}{c|}{{\scriptsize rs1[4]}} &
\multicolumn{1}{c|}{{\scriptsize rs1[2:0]}} &
\multicolumn{1}{c|}{10} &
\multicolumn{1}{c|}{{\scriptsize qu[2:0]}} &
\multicolumn{1}{c|}{101} &
\multicolumn{1}{c|}{0011111}\\
\hline
3& 3& 3& 3& 1& 1& 3& 2& 3& 3& 7 \\
\end{tabular}
 \end{center}
 }
 \vspace{-0.1in}
 \caption[]{ESP.VSUB.U8.LD.INCP} \end{figure*}


\newpage
\subsubsection{ESP.VSUB.U8.ST.INCP}\label{instruction:ESPVSUBU8STINCP}
\textbf{Syntax}\\
ESP.VSUB.U8.ST.INCP qu,rs1,qv,qx,qy,sat

\textbf{Description}\\
This instruction performs a vector subtraction on 8-bit unsigned data. Registers \lstinline{qx} and \lstinline{qy} are the subtrahend and the minuend, respectively. Then, the 16 results obtained from the calculation are saturated and then written into the register \lstinline{qv}. \\ During the operation, the value in the register \lstinline{qu} is stored to memory. After the access, the value in the register \lstinline{rs1} is incremented by 16.

\textbf{Operation}\\
\begin{lstlisting}
  qv[  7:  0] = min(qx[  7:  0] - qy[  7:  0], 2^{8}-1)
  qv[ 15:  8] = min(qx[ 15:  8] - qy[ 15:  8], 2^{8}-1)
  ...

  qu[127:0] => store128(rs1[31:0])
  rs1[31:0] = rs1[31:0] + 16
\end{lstlisting}

\textbf{Schedule}\\
\begin{longtable}[c]{|l|l|l|}
    \hline
    \rowcolor{lightgray}
    \textbf{Operands} & \textbf{Use} & \textbf{Def} \\\hline
    \endfirsthead
    \hline
    \rowcolor{lightgray}
     \textbf{Operands} & \textbf{Use} & \textbf{Def} \\\hline
    \endhead
    \hline
    \endfoot

  qv & - & 1 \\\hline
    qx & 1 & - \\\hline
    qy & 1 & - \\\hline
    qu & 2 & - \\\hline
    rs1 & 1 & 1

\end{longtable}

\textbf{Format}\\

\begin{figure*}[!htbp]
 {\setlength{\tabcolsep}{1pt}
 \begin{center}
\begin{tabular}{@{}F@{}F@{}F@{}F@{}I@{}I@{}F@{}W@{}F@{}F@{}O}
\instbitrange{31}{29}&
\instbitrange{28}{26}&
\instbitrange{25}{23}&
\instbitrange{22}{20}&
\instbit{19}&
\instbit{18}&
\instbitrange{17}{15}&
\instbitrange{14}{13}&
\instbitrange{12}{10}&
\instbitrange{9}{7}&
\instbitrange{6}{0} \\
\hline
\multicolumn{1}{|c|}{{\scriptsize qx[2:0]}} &
\multicolumn{1}{c|}{{\scriptsize qy[2:0]}} &
\multicolumn{1}{c|}{100} &
\multicolumn{1}{c|}{{\scriptsize qv[2:0]}} &
\multicolumn{1}{c|}{{\scriptsize sat[0]}} &
\multicolumn{1}{c|}{{\scriptsize rs1[4]}} &
\multicolumn{1}{c|}{{\scriptsize rs1[2:0]}} &
\multicolumn{1}{c|}{10} &
\multicolumn{1}{c|}{{\scriptsize qu[2:0]}} &
\multicolumn{1}{c|}{101} &
\multicolumn{1}{c|}{0011111}\\
\hline
3& 3& 3& 3& 1& 1& 3& 2& 3& 3& 7 \\
\end{tabular}
 \end{center}
 }
 \vspace{-0.1in}
 \caption[]{ESP.VSUB.U8.ST.INCP} \end{figure*}


\newpage
\subsubsection{ESP.ADDX2}\label{instruction:ESPADDX2}
\textbf{Syntax}\\
ESP.ADDX2 rd,rs1,rs2

\textbf{Description}\\
This instruction shifts \lstinline{rs2} to the left by 1 and adds it to \lstinline{rs1}.

\textbf{Operation}\\
\begin{lstlisting}
    rd = rs1 + rs2<<1
\end{lstlisting}

\textbf{Schedule}\\
\begin{longtable}[c]{|l|l|l|}
    \hline
    \rowcolor{lightgray}
    \textbf{Operands} & \textbf{Use} & \textbf{Def} \\\hline
    \endfirsthead
    \hline
    \rowcolor{lightgray}
     \textbf{Operands} & \textbf{Use} & \textbf{Def} \\\hline
    \endhead
    \hline
    \endfoot

  rd & - & 1 \\\hline
    rs1 & 1 & - \\\hline
    rs2 & 1 & -

\end{longtable}

\textbf{Format}\\

\begin{figure*}[!htbp]
 {\setlength{\tabcolsep}{1pt}
 \begin{center}
\begin{tabular}{@{}O@{}V@{}V@{}F@{}V@{}O}
\instbitrange{31}{25}&
\instbitrange{24}{20}&
\instbitrange{19}{15}&
\instbitrange{14}{12}&
\instbitrange{11}{7}&
\instbitrange{6}{0} \\
\hline
\multicolumn{1}{|c|}{0000010} &
\multicolumn{1}{c|}{{\scriptsize rs2[4:0]}} &
\multicolumn{1}{c|}{{\scriptsize rs1[4:0]}} &
\multicolumn{1}{c|}{000} &
\multicolumn{1}{c|}{{\scriptsize rd[4:0]}} &
\multicolumn{1}{c|}{0110011}\\
\hline
7& 5& 5& 3& 5& 7 \\
\end{tabular}
 \end{center}
 }
 \vspace{-0.1in}
 \caption[]{ESP.ADDX2} \end{figure*}


\newpage
\subsubsection{ESP.ADDX4}\label{instruction:ESPADDX4}
\textbf{Syntax}\\
ESP.ADDX4 rd,rs1,rs2

\textbf{Description}\\
This instruction shifts \lstinline{rs2} to the left by 2 and adds it to \lstinline{rs1}.

\textbf{Operation}\\
\begin{lstlisting}
    rd = rs1 + rs2<<2
\end{lstlisting}


\textbf{Schedule}\\
\begin{longtable}[c]{|l|l|l|}
    \hline
    \rowcolor{lightgray}
    \textbf{Operands} & \textbf{Use} & \textbf{Def} \\\hline
    \endfirsthead
    \hline
    \rowcolor{lightgray}
     \textbf{Operands} & \textbf{Use} & \textbf{Def} \\\hline
    \endhead
    \hline
    \endfoot

  rd & - & 1 \\\hline
    rs1 & 1 & - \\\hline
    rs2 & 1 & -

\end{longtable}

\textbf{Format}\\

\begin{figure*}[!htbp]
 {\setlength{\tabcolsep}{1pt}
 \begin{center}
\begin{tabular}{@{}O@{}V@{}V@{}F@{}V@{}O}
\instbitrange{31}{25}&
\instbitrange{24}{20}&
\instbitrange{19}{15}&
\instbitrange{14}{12}&
\instbitrange{11}{7}&
\instbitrange{6}{0} \\
\hline
\multicolumn{1}{|c|}{0000100} &
\multicolumn{1}{c|}{{\scriptsize rs2[4:0]}} &
\multicolumn{1}{c|}{{\scriptsize rs1[4:0]}} &
\multicolumn{1}{c|}{000} &
\multicolumn{1}{c|}{{\scriptsize rd[4:0]}} &
\multicolumn{1}{c|}{0110011}\\
\hline
7& 5& 5& 3& 5& 7 \\
\end{tabular}
 \end{center}
 }
 \vspace{-0.1in}
 \caption[]{ESP.ADDX4} \end{figure*}


\newpage
\subsubsection{ESP.SAT}\label{instruction:ESPSAT}
\textbf{Syntax}\\
ESP.SAT rsd,rs0,rs1

\textbf{Description}\\
This instruction performs saturation for \lstinline{rsd}.

\textbf{Operation}\\
\begin{lstlisting}[escapeinside=`?]
    min_t[31:0] = min( rs1[31:0], rs0[31:0] )
    max_t[31:0] = max( rs1[31:0], rs0[31:0] )

    rsd[31:0] = max( min( rsd[31:0], max_t[31:0]), min_t[31:0] )
\end{lstlisting}

\textbf{Schedule}\\
\begin{longtable}[c]{|l|l|l|}
    \hline
    \rowcolor{lightgray}
    \textbf{Operands} & \textbf{Use} & \textbf{Def} \\\hline
    \endfirsthead
    \hline
    \rowcolor{lightgray}
     \textbf{Operands} & \textbf{Use} & \textbf{Def} \\\hline
    \endhead
    \hline
    \endfoot

  rsd & 1 & 1 \\\hline
    rs0 & 1 & - \\\hline
    rs1 & 1 & -

\end{longtable}

\textbf{Format}\\

\begin{figure*}[!htbp]
 {\setlength{\tabcolsep}{1pt}
 \begin{center}
\begin{tabular}{@{}O@{}V@{}V@{}F@{}V@{}O}
\instbitrange{31}{25}&
\instbitrange{24}{20}&
\instbitrange{19}{15}&
\instbitrange{14}{12}&
\instbitrange{11}{7}&
\instbitrange{6}{0} \\
\hline
\multicolumn{1}{|c|}{0100000} &
\multicolumn{1}{c|}{{\scriptsize rsd[4:0]}} &
\multicolumn{1}{c|}{{\scriptsize rs1[4:0]}} &
\multicolumn{1}{c|}{010} &
\multicolumn{1}{c|}{{\scriptsize rs0[4:0]}} &
\multicolumn{1}{c|}{0110011}\\
\hline
7& 5& 5& 3& 5& 7 \\
\end{tabular}
 \end{center}
 }
 \vspace{-0.1in}
 \caption[]{ESP.SAT} \end{figure*}


\newpage
\subsubsection{ESP.SUBX2}\label{instruction:ESPSUBX2}
\textbf{Syntax}\\
ESP.SUBX2 rd,rs1,rs2

\textbf{Description}\\
This instruction shifts \lstinline{rs2} to the left by 1 and subtracts it from \lstinline{rs1}.

\textbf{Operation}\\
\begin{lstlisting}
    rd = rs1 - rs2<<1
\end{lstlisting}

\textbf{Schedule}\\
\begin{longtable}[c]{|l|l|l|}
    \hline
    \rowcolor{lightgray}
    \textbf{Operands} & \textbf{Use} & \textbf{Def} \\\hline
    \endfirsthead
    \hline
    \rowcolor{lightgray}
     \textbf{Operands} & \textbf{Use} & \textbf{Def} \\\hline
    \endhead
    \hline
    \endfoot

  rd & - & 1 \\\hline
    rs1 & 1 & - \\\hline
    rs2 & 1 & -

\end{longtable}

\textbf{Format}\\

\begin{figure*}[!htbp]
 {\setlength{\tabcolsep}{1pt}
 \begin{center}
\begin{tabular}{@{}O@{}V@{}V@{}F@{}V@{}O}
\instbitrange{31}{25}&
\instbitrange{24}{20}&
\instbitrange{19}{15}&
\instbitrange{14}{12}&
\instbitrange{11}{7}&
\instbitrange{6}{0} \\
\hline
\multicolumn{1}{|c|}{0100010} &
\multicolumn{1}{c|}{{\scriptsize rs2[4:0]}} &
\multicolumn{1}{c|}{{\scriptsize rs1[4:0]}} &
\multicolumn{1}{c|}{000} &
\multicolumn{1}{c|}{{\scriptsize rd[4:0]}} &
\multicolumn{1}{c|}{0110011}\\
\hline
7& 5& 5& 3& 5& 7 \\
\end{tabular}
 \end{center}
 }
 \vspace{-0.1in}
 \caption[]{ESP.SUBX2} \end{figure*}


\newpage
\subsubsection{ESP.SUBX4}\label{instruction:ESPSUBX4}
\textbf{Syntax}\\
ESP.SUBX4 rd,rs1,rs2

\textbf{Description}\\
This instruction shifts \lstinline{rs2} to the left by 2 and subtracts it from \lstinline{rs1}.

\textbf{Operation}\\
\begin{lstlisting}
    rd = rs1 - rs2<<2
\end{lstlisting}

\textbf{Schedule}\\
\begin{longtable}[c]{|l|l|l|}
    \hline
    \rowcolor{lightgray}
    \textbf{Operands} & \textbf{Use} & \textbf{Def} \\\hline
    \endfirsthead
    \hline
    \rowcolor{lightgray}
     \textbf{Operands} & \textbf{Use} & \textbf{Def} \\\hline
    \endhead
    \hline
    \endfoot

  rd & - & 1 \\\hline
    rs1 & 1 & - \\\hline
    rs2 & 1 & -

\end{longtable}

\textbf{Format}\\

\begin{figure*}[!htbp]
 {\setlength{\tabcolsep}{1pt}
 \begin{center}
\begin{tabular}{@{}O@{}V@{}V@{}F@{}V@{}O}
\instbitrange{31}{25}&
\instbitrange{24}{20}&
\instbitrange{19}{15}&
\instbitrange{14}{12}&
\instbitrange{11}{7}&
\instbitrange{6}{0} \\
\hline
\multicolumn{1}{|c|}{0100100} &
\multicolumn{1}{c|}{{\scriptsize rs2[4:0]}} &
\multicolumn{1}{c|}{{\scriptsize rs1[4:0]}} &
\multicolumn{1}{c|}{000} &
\multicolumn{1}{c|}{{\scriptsize rd[4:0]}} &
\multicolumn{1}{c|}{0110011}\\
\hline
7& 5& 5& 3& 5& 7 \\
\end{tabular}
 \end{center}
 }
 \vspace{-0.1in}
 \caption[]{ESP.SUBX4} \end{figure*}


\newpage
\subsubsection{ESP.ANDQ}\label{instruction:ESPANDQ}
\textbf{Syntax}\\
ESP.ANDQ qz,qx,qy

\textbf{Description}\\
This instruction performs a bitwise AND operation on registers \lstinline{qx} and \lstinline{qy}, and writes the result of the logical operation to the register \lstinline{qz}.\\

\textbf{Operation}\\
\begin{lstlisting}
    qz = qx & qy
  \end{lstlisting}

\textbf{Schedule}\\
\begin{longtable}[c]{|l|l|l|}
    \hline
    \rowcolor{lightgray}
    \textbf{Operands} & \textbf{Use} & \textbf{Def} \\\hline
    \endfirsthead
    \hline
    \rowcolor{lightgray}
     \textbf{Operands} & \textbf{Use} & \textbf{Def} \\\hline
    \endhead
    \hline
    \endfoot

  qz & - & 1 \\\hline
    qx & 1 & - \\\hline
    qy & 1 & -

\end{longtable}

\textbf{Format}\\

\begin{figure*}[!htbp]
 {\setlength{\tabcolsep}{1pt}
 \begin{center}
\begin{tabular}{@{}F@{}F@{}K@{}F@{}O}
\instbitrange{31}{29}&
\instbitrange{28}{26}&
\instbitrange{25}{10}&
\instbitrange{9}{7}&
\instbitrange{6}{0} \\
\hline
\multicolumn{1}{|c|}{{\scriptsize qx[2:0]}} &
\multicolumn{1}{c|}{{\scriptsize qy[2:0]}} &
\multicolumn{1}{c|}{011XXXX0XXX00XXX} &
\multicolumn{1}{c|}{{\scriptsize qz[2:0]}} &
\multicolumn{1}{c|}{0011111}\\
\hline
3& 3& 16& 3& 7 \\
\end{tabular}
 \end{center}
 }
 \vspace{-0.1in}
 \caption[]{ESP.ANDQ} \end{figure*}


\newpage
\subsubsection{ESP.NOTQ}\label{instruction:ESPNOTQ}
\textbf{Syntax}\\
ESP.NOTQ qz,qx

\textbf{Description}\\
This instruction performs a bitwise NOT operation on the register \lstinline{qx} and writes the result to the register \lstinline{qz}.\\
\textbf{Operation}
\begin{lstlisting}
  qz = ~qx
\end{lstlisting}

\textbf{Schedule}\\
\begin{longtable}[c]{|l|l|l|}
    \hline
    \rowcolor{lightgray}
    \textbf{Operands} & \textbf{Use} & \textbf{Def} \\\hline
    \endfirsthead
    \hline
    \rowcolor{lightgray}
     \textbf{Operands} & \textbf{Use} & \textbf{Def} \\\hline
    \endhead
    \hline
    \endfoot

  qz & - & 1 \\\hline
    qx & 1 & -

\end{longtable}

\textbf{Format}\\

\begin{figure*}[!htbp]
 {\setlength{\tabcolsep}{1pt}
 \begin{center}
\begin{tabular}{@{}F@{}U@{}F@{}O}
\instbitrange{31}{29}&
\instbitrange{28}{10}&
\instbitrange{9}{7}&
\instbitrange{6}{0} \\
\hline
\multicolumn{1}{|c|}{{\scriptsize qx[2:0]}} &
\multicolumn{1}{c|}{XXX011XXXX1XXX00XXX} &
\multicolumn{1}{c|}{{\scriptsize qz[2:0]}} &
\multicolumn{1}{c|}{0011111}\\
\hline
3& 19& 3& 7 \\
\end{tabular}
 \end{center}
 }
 \vspace{-0.1in}
 \caption[]{ESP.NOTQ} \end{figure*}


\newpage
\subsubsection{ESP.ORQ}\label{instruction:ESPORQ}
\textbf{Syntax}\\
ESP.ORQ qz,qx,qy

\textbf{Description}\\
This instruction performs a bitwise or operation on registers \lstinline{qx} and \lstinline{qy}, and writes the result of the logical operation to the register \lstinline{qz}.\\
\textbf{Operation}
\begin{lstlisting}
  qz = qx | qy
\end{lstlisting}

\textbf{Schedule}\\
\begin{longtable}[c]{|l|l|l|}
    \hline
    \rowcolor{lightgray}
    \textbf{Operands} & \textbf{Use} & \textbf{Def} \\\hline
    \endfirsthead
    \hline
    \rowcolor{lightgray}
     \textbf{Operands} & \textbf{Use} & \textbf{Def} \\\hline
    \endhead
    \hline
    \endfoot

  qz & - & 1 \\\hline
    qx & 1 & - \\\hline
    qy & 1 & -

\end{longtable}

\textbf{Format}\\

\begin{figure*}[!htbp]
 {\setlength{\tabcolsep}{1pt}
 \begin{center}
\begin{tabular}{@{}F@{}F@{}K@{}F@{}O}
\instbitrange{31}{29}&
\instbitrange{28}{26}&
\instbitrange{25}{10}&
\instbitrange{9}{7}&
\instbitrange{6}{0} \\
\hline
\multicolumn{1}{|c|}{{\scriptsize qx[2:0]}} &
\multicolumn{1}{c|}{{\scriptsize qy[2:0]}} &
\multicolumn{1}{c|}{010XXXX0XXX00XXX} &
\multicolumn{1}{c|}{{\scriptsize qz[2:0]}} &
\multicolumn{1}{c|}{0011111}\\
\hline
3& 3& 16& 3& 7 \\
\end{tabular}
 \end{center}
 }
 \vspace{-0.1in}
 \caption[]{ESP.ORQ} \end{figure*}


\newpage
\subsubsection{ESP.XORQ}\label{instruction:ESPXORQ}
\textbf{Syntax}\\
ESP.XORQ qz,qx,qy

\textbf{Description}\\
This instruction performs a bitwise xor operation on registers \lstinline{qx} and \lstinline{qy}, and writes the result of the logical operation to the register \lstinline{qz}.\\
\textbf{Operation}
\begin{lstlisting}
  qz = qx ^ qy
\end{lstlisting}

\textbf{Schedule}\\
\begin{longtable}[c]{|l|l|l|}
    \hline
    \rowcolor{lightgray}
    \textbf{Operands} & \textbf{Use} & \textbf{Def} \\\hline
    \endfirsthead
    \hline
    \rowcolor{lightgray}
     \textbf{Operands} & \textbf{Use} & \textbf{Def} \\\hline
    \endhead
    \hline
    \endfoot

  qz & - & 1 \\\hline
    qx & 1 & - \\\hline
    qy & 1 & -

\end{longtable}

\textbf{Format}\\

\begin{figure*}[!htbp]
 {\setlength{\tabcolsep}{1pt}
 \begin{center}
\begin{tabular}{@{}F@{}F@{}K@{}F@{}O}
\instbitrange{31}{29}&
\instbitrange{28}{26}&
\instbitrange{25}{10}&
\instbitrange{9}{7}&
\instbitrange{6}{0} \\
\hline
\multicolumn{1}{|c|}{{\scriptsize qx[2:0]}} &
\multicolumn{1}{c|}{{\scriptsize qy[2:0]}} &
\multicolumn{1}{c|}{010XXXX1XXX00XXX} &
\multicolumn{1}{c|}{{\scriptsize qz[2:0]}} &
\multicolumn{1}{c|}{0011111}\\
\hline
3& 3& 16& 3& 7 \\
\end{tabular}
 \end{center}
 }
 \vspace{-0.1in}
 \caption[]{ESP.XORQ} \end{figure*}


\newpage
\subsubsection{ESP.VCMP.EQ.S16}\label{instruction:ESPVCMPEqy6}
\textbf{Syntax}\\
ESP.VCMP.EQ.S16 qz,qx,qy

\textbf{Description}\\
This instruction compares 16-bit vector data. It compares the numerical values of the eight 16-bit data segments in registers \lstinline{qx} and \lstinline{qy}. If the values are equal, it writes 0xFFFF into the corresponding 16-bit data segment in register \lstinline{qz}. Otherwise, it writes 0 into the segment.\\
\textbf{Operation}
\begin{lstlisting}[escapeinside=`?]
  qz[ 15:  0] = (qx[ 15:  0]==qy[ 15:  0]) ? 0xFFFF : 0
  qz[ 31: 16] = (qx[ 31: 16]==qy[ 31: 16]) ? 0xFFFF : 0
  ...
  qz[127:112] = (qx[127:112]==qy[127:112]) ? 0xFFFF : 0
\end{lstlisting}

\textbf{Schedule}\\
\begin{longtable}[c]{|l|l|l|}
    \hline
    \rowcolor{lightgray}
    \textbf{Operands} & \textbf{Use} & \textbf{Def} \\\hline
    \endfirsthead
    \hline
    \rowcolor{lightgray}
     \textbf{Operands} & \textbf{Use} & \textbf{Def} \\\hline
    \endhead
    \hline
    \endfoot

  qz & - & 1 \\\hline
    qx & 1 & - \\\hline
    qy & 1 & -

\end{longtable}

\textbf{Format}\\

\begin{figure*}[!htbp]
 {\setlength{\tabcolsep}{1pt}
 \begin{center}
\begin{tabular}{@{}F@{}F@{}K@{}F@{}O}
\instbitrange{31}{29}&
\instbitrange{28}{26}&
\instbitrange{25}{10}&
\instbitrange{9}{7}&
\instbitrange{6}{0} \\
\hline
\multicolumn{1}{|c|}{{\scriptsize qx[2:0]}} &
\multicolumn{1}{c|}{{\scriptsize qy[2:0]}} &
\multicolumn{1}{c|}{110100X001100X0X} &
\multicolumn{1}{c|}{{\scriptsize qz[2:0]}} &
\multicolumn{1}{c|}{0011111}\\
\hline
3& 3& 16& 3& 7 \\
\end{tabular}
 \end{center}
 }
 \vspace{-0.1in}
 \caption[]{ESP.VCMP.EQ.S16} \end{figure*}


\newpage
\subsubsection{ESP.VCMP.EQ.S32}\label{instruction:ESPVCMPEQS32}
\textbf{Syntax}\\
ESP.VCMP.EQ.S32 qz,qx,qy

\textbf{Description}\\
This instruction compares 32-bit vector data. It compares the numerical values of the four 32-bit data segments in registers \lstinline{qx} and \lstinline{qy}. If the values are equal, it writes 0xFFFFFFFF into the corresponding 32-bit data segment in the \lstinline{qz} register. Otherwise, it writes 0 into the segment.\\
\textbf{Operation}
\begin{lstlisting}[escapeinside=`?]
  qz[ 31:  0] = (qx[ 31:  0]==qy[ 31:  0]) ? 0xFFFFFFFF : 0
  qz[ 63: 32] = (qx[ 63: 32]==qy[ 63: 32]) ? 0xFFFFFFFF : 0
  ...
  qz[127: 96] = (qx[127: 96]==qy[127: 96]) ? 0xFFFFFFFF : 0
\end{lstlisting}

\textbf{Schedule}\\
\begin{longtable}[c]{|l|l|l|}
    \hline
    \rowcolor{lightgray}
    \textbf{Operands} & \textbf{Use} & \textbf{Def} \\\hline
    \endfirsthead
    \hline
    \rowcolor{lightgray}
     \textbf{Operands} & \textbf{Use} & \textbf{Def} \\\hline
    \endhead
    \hline
    \endfoot

  qz & - & 1 \\\hline
    qx & 1 & - \\\hline
    qy & 1 & -

\end{longtable}

\textbf{Format}\\

\begin{figure*}[!htbp]
 {\setlength{\tabcolsep}{1pt}
 \begin{center}
\begin{tabular}{@{}F@{}F@{}K@{}F@{}O}
\instbitrange{31}{29}&
\instbitrange{28}{26}&
\instbitrange{25}{10}&
\instbitrange{9}{7}&
\instbitrange{6}{0} \\
\hline
\multicolumn{1}{|c|}{{\scriptsize qx[2:0]}} &
\multicolumn{1}{c|}{{\scriptsize qy[2:0]}} &
\multicolumn{1}{c|}{110100X001X0001X} &
\multicolumn{1}{c|}{{\scriptsize qz[2:0]}} &
\multicolumn{1}{c|}{0011111}\\
\hline
3& 3& 16& 3& 7 \\
\end{tabular}
 \end{center}
 }
 \vspace{-0.1in}
 \caption[]{ESP.VCMP.EQ.S32} \end{figure*}


\newpage
\subsubsection{ESP.VCMP.EQ.S8}\label{instruction:ESPVCMPEQS8}
\textbf{Syntax}\\
ESP.VCMP.EQ.S8 qz,qx,qy

\textbf{Description}\\
This instruction compares 8-bit vector data. It compares the numerical values of the 16 8-bit data segments in registers \lstinline{qx} and \lstinline{qy}. If the values are equal, it writes 0xFF into the corresponding 8-bit data segment in register \lstinline{qz}. Otherwise, it writes 0 into the segment.\\
\textbf{Operation}
\begin{lstlisting}[escapeinside=`?]
  qz[  7:  0] = (qx[  7:  0]==qy[  7:  0]) ? 0xFF : 0
  qz[ 15:  8] = (qx[ 15:  8]==qy[ 15:  8]) ? 0xFF : 0
  ...
  qz[127:120] = (qx[127:120]==qy[127:120]) ? 0xFF : 0
\end{lstlisting}

\textbf{Schedule}\\
\begin{longtable}[c]{|l|l|l|}
    \hline
    \rowcolor{lightgray}
    \textbf{Operands} & \textbf{Use} & \textbf{Def} \\\hline
    \endfirsthead
    \hline
    \rowcolor{lightgray}
     \textbf{Operands} & \textbf{Use} & \textbf{Def} \\\hline
    \endhead
    \hline
    \endfoot

  qz & - & 1 \\\hline
    qx & 1 & - \\\hline
    qy & 1 & -

\end{longtable}

\textbf{Format}\\

\begin{figure*}[!htbp]
 {\setlength{\tabcolsep}{1pt}
 \begin{center}
\begin{tabular}{@{}F@{}F@{}K@{}F@{}O}
\instbitrange{31}{29}&
\instbitrange{28}{26}&
\instbitrange{25}{10}&
\instbitrange{9}{7}&
\instbitrange{6}{0} \\
\hline
\multicolumn{1}{|c|}{{\scriptsize qx[2:0]}} &
\multicolumn{1}{c|}{{\scriptsize qy[2:0]}} &
\multicolumn{1}{c|}{110100X000100X0X} &
\multicolumn{1}{c|}{{\scriptsize qz[2:0]}} &
\multicolumn{1}{c|}{0011111}\\
\hline
3& 3& 16& 3& 7 \\
\end{tabular}
 \end{center}
 }
 \vspace{-0.1in}
 \caption[]{ESP.VCMP.EQ.S8} \end{figure*}


\newpage
\subsubsection{ESP.VCMP.EQ.U16}\label{instruction:ESPVCMPEqu16}
\textbf{Syntax}\\
ESP.VCMP.EQ.U16 qz,qx,qy

\textbf{Description}\\
This instruction compares 16-bit vector data. It compares the numerical values of the eight 16-bit data segments in the registers \lstinline{qx} and \lstinline{qy}. If the values are equal, it writes 0xFFFF into the corresponding 16-bit data segment in the register \lstinline{qz}. Otherwise, it writes 0 into the segment.\\
\textbf{Operation}
\begin{lstlisting}[escapeinside=`?]
  qz[ 15:  0] = (qx[ 15:  0]==qy[ 15:  0]) ? 0xFFFF : 0
  qz[ 31: 16] = (qx[ 31: 16]==qy[ 31: 16]) ? 0xFFFF : 0
  ...
  qz[127:112] = (qx[127:112]==qy[127:112]) ? 0xFFFF : 0
\end{lstlisting}

\textbf{Schedule}\\
\begin{longtable}[c]{|l|l|l|}
    \hline
    \rowcolor{lightgray}
    \textbf{Operands} & \textbf{Use} & \textbf{Def} \\\hline
    \endfirsthead
    \hline
    \rowcolor{lightgray}
     \textbf{Operands} & \textbf{Use} & \textbf{Def} \\\hline
    \endhead
    \hline
    \endfoot

  qz & - & 1 \\\hline
    qx & 1 & - \\\hline
    qy & 1 & -

\end{longtable}

\textbf{Format}\\

\begin{figure*}[!htbp]
 {\setlength{\tabcolsep}{1pt}
 \begin{center}
\begin{tabular}{@{}F@{}F@{}K@{}F@{}O}
\instbitrange{31}{29}&
\instbitrange{28}{26}&
\instbitrange{25}{10}&
\instbitrange{9}{7}&
\instbitrange{6}{0} \\
\hline
\multicolumn{1}{|c|}{{\scriptsize qx[2:0]}} &
\multicolumn{1}{c|}{{\scriptsize qy[2:0]}} &
\multicolumn{1}{c|}{110100X001000X0X} &
\multicolumn{1}{c|}{{\scriptsize qz[2:0]}} &
\multicolumn{1}{c|}{0011111}\\
\hline
3& 3& 16& 3& 7 \\
\end{tabular}
 \end{center}
 }
 \vspace{-0.1in}
 \caption[]{ESP.VCMP.EQ.U16} \end{figure*}


\newpage
\subsubsection{ESP.VCMP.EQ.U32}\label{instruction:ESPVCMPEqu32}
\textbf{Syntax}\\
ESP.VCMP.EQ.U32 qz,qx,qy

\textbf{Description}\\
This instruction compares 32-bit vector data. It compares the numerical values of the four 32-bit data segments in the registers \lstinline{qx} and \lstinline{qy}. If the values are equal, it writes 0xFFFFFFFF into the corresponding 32-bit data segment in the register \lstinline{qz}. Otherwise, it writes 0 into the segment.\\
\textbf{Operation}
\begin{lstlisting}[escapeinside=`?]
  qz[ 31:  0] = (qx[ 31:  0]==qy[ 31:  0]) ? 0xFFFFFFFF : 0
  qz[ 63: 32] = (qx[ 63: 32]==qy[ 63: 32]) ? 0xFFFFFFFF : 0
  ...
  qz[127: 96] = (qx[127: 96]==qy[127: 96]) ? 0xFFFFFFFF : 0
\end{lstlisting}

\textbf{Schedule}\\
\begin{longtable}[c]{|l|l|l|}
    \hline
    \rowcolor{lightgray}
    \textbf{Operands} & \textbf{Use} & \textbf{Def} \\\hline
    \endfirsthead
    \hline
    \rowcolor{lightgray}
     \textbf{Operands} & \textbf{Use} & \textbf{Def} \\\hline
    \endhead
    \hline
    \endfoot

  qz & - & 1 \\\hline
    qx & 1 & - \\\hline
    qy & 1 & -

\end{longtable}

\textbf{Format}\\

\begin{figure*}[!htbp]
 {\setlength{\tabcolsep}{1pt}
 \begin{center}
\begin{tabular}{@{}F@{}F@{}K@{}F@{}O}
\instbitrange{31}{29}&
\instbitrange{28}{26}&
\instbitrange{25}{10}&
\instbitrange{9}{7}&
\instbitrange{6}{0} \\
\hline
\multicolumn{1}{|c|}{{\scriptsize qx[2:0]}} &
\multicolumn{1}{c|}{{\scriptsize qy[2:0]}} &
\multicolumn{1}{c|}{110100X000X0001X} &
\multicolumn{1}{c|}{{\scriptsize qz[2:0]}} &
\multicolumn{1}{c|}{0011111}\\
\hline
3& 3& 16& 3& 7 \\
\end{tabular}
 \end{center}
 }
 \vspace{-0.1in}
 \caption[]{ESP.VCMP.EQ.U32} \end{figure*}


\newpage
\subsubsection{ESP.VCMP.EQ.U8}\label{instruction:ESPVCMPEqu8}
\textbf{Syntax}\\
ESP.VCMP.EQ.U8 qz,qx,qy

\textbf{Description}\\
This instruction compares 8-bit vector data. It compares the numerical values of the 16 8-bit data segments in registers \lstinline{qx} and \lstinline{qy}. If the values are equal, it writes 0xFF into the corresponding 8-bit data segment in the \lstinline{qz} register. Otherwise, it writes 0 into the segment.\\
\textbf{Operation}
\begin{lstlisting}[escapeinside=`?]
  qz[  7:  0] = (qx[  7:  0]==qy[  7:  0]) ? 0xFF : 0
  qz[ 15:  8] = (qx[ 15:  8]==qy[ 15:  8]) ? 0xFF : 0
  ...
  qz[127:120] = (qx[127:120]==qy[127:120]) ? 0xFF : 0
\end{lstlisting}


\textbf{Schedule}\\
\begin{longtable}[c]{|l|l|l|}
    \hline
    \rowcolor{lightgray}
    \textbf{Operands} & \textbf{Use} & \textbf{Def} \\\hline
    \endfirsthead
    \hline
    \rowcolor{lightgray}
     \textbf{Operands} & \textbf{Use} & \textbf{Def} \\\hline
    \endhead
    \hline
    \endfoot

  qz & - & 1 \\\hline
    qx & 1 & - \\\hline
    qy & 1 & -

\end{longtable}

\textbf{Format}\\

\begin{figure*}[!htbp]
 {\setlength{\tabcolsep}{1pt}
 \begin{center}
\begin{tabular}{@{}F@{}F@{}K@{}F@{}O}
\instbitrange{31}{29}&
\instbitrange{28}{26}&
\instbitrange{25}{10}&
\instbitrange{9}{7}&
\instbitrange{6}{0} \\
\hline
\multicolumn{1}{|c|}{{\scriptsize qx[2:0]}} &
\multicolumn{1}{c|}{{\scriptsize qy[2:0]}} &
\multicolumn{1}{c|}{110100X000000X0X} &
\multicolumn{1}{c|}{{\scriptsize qz[2:0]}} &
\multicolumn{1}{c|}{0011111}\\
\hline
3& 3& 16& 3& 7 \\
\end{tabular}
 \end{center}
 }
 \vspace{-0.1in}
 \caption[]{ESP.VCMP.EQ.U8} \end{figure*}


\newpage
\subsubsection{ESP.VCMP.GT.S16}\label{instruction:ESPVCMPGTS16}
\textbf{Syntax}\\
ESP.VCMP.GT.S16 qz,qx,qy

\textbf{Description}\\
This instruction compares 16-bit vector data. It compares the numerical values of the eight 16-bit data segments in registers \lstinline{qx} and \lstinline{qy}. If the former is larger than the latter, it writes 0xFFFF into the corresponding 16-bit data segment in register \lstinline{qz}. Otherwise, it writes 0 into the segment.\\
\textbf{Operation}
\begin{lstlisting}[escapeinside=`?]
  qz[ 15:  0] = (qx[ 15:  0]>qy[ 15:  0]) ? 0xFFFF : 0
  qz[ 31: 16] = (qx[ 31: 16]>qy[ 31: 16]) ? 0xFFFF : 0
  ...
  qz[127:112] = (qx[127:112]>qy[127:112]) ? 0xFFFF : 0
\end{lstlisting}

\textbf{Schedule}\\
\begin{longtable}[c]{|l|l|l|}
    \hline
    \rowcolor{lightgray}
    \textbf{Operands} & \textbf{Use} & \textbf{Def} \\\hline
    \endfirsthead
    \hline
    \rowcolor{lightgray}
     \textbf{Operands} & \textbf{Use} & \textbf{Def} \\\hline
    \endhead
    \hline
    \endfoot

  qz & - & 1 \\\hline
    qx & 1 & - \\\hline
    qy & 1 & -

\end{longtable}

\textbf{Format}\\

\begin{figure*}[!htbp]
 {\setlength{\tabcolsep}{1pt}
 \begin{center}
\begin{tabular}{@{}F@{}F@{}K@{}F@{}O}
\instbitrange{31}{29}&
\instbitrange{28}{26}&
\instbitrange{25}{10}&
\instbitrange{9}{7}&
\instbitrange{6}{0} \\
\hline
\multicolumn{1}{|c|}{{\scriptsize qx[2:0]}} &
\multicolumn{1}{c|}{{\scriptsize qy[2:0]}} &
\multicolumn{1}{c|}{110100X101100X0X} &
\multicolumn{1}{c|}{{\scriptsize qz[2:0]}} &
\multicolumn{1}{c|}{0011111}\\
\hline
3& 3& 16& 3& 7 \\
\end{tabular}
 \end{center}
 }
 \vspace{-0.1in}
 \caption[]{ESP.VCMP.GT.S16} \end{figure*}


\newpage
\subsubsection{ESP.VCMP.GT.S32}\label{instruction:ESPVCMPGTS32}
\textbf{Syntax}\\
ESP.VCMP.GT.S32 qz,qx,qy

\textbf{Description}\\
This instruction compares 32-bit signed vector data. It compares the numerical values of the four 32-bit data segments in the registers \lstinline{qx} and \lstinline{qy}. If the former is larger than the latter, it writes 0xFFFFFFFF into the corresponding 32-bit data segment in the register \lstinline{qz}. Otherwise, it writes 0 into the segment.\\
\textbf{Operation}
\begin{lstlisting}[escapeinside=`?]
  qz[ 31:  0] = (qx[ 31:  0]>qy[ 31:  0]) ? 0xFFFFFFFF : 0
  qz[ 63: 32] = (qx[ 63: 32]>qy[ 63: 32]) ? 0xFFFFFFFF : 0
  ...
  qz[127: 96] = (qx[127: 96]>qy[127: 96]) ? 0xFFFFFFFF : 0
\end{lstlisting}

\textbf{Schedule}\\
\begin{longtable}[c]{|l|l|l|}
    \hline
    \rowcolor{lightgray}
    \textbf{Operands} & \textbf{Use} & \textbf{Def} \\\hline
    \endfirsthead
    \hline
    \rowcolor{lightgray}
     \textbf{Operands} & \textbf{Use} & \textbf{Def} \\\hline
    \endhead
    \hline
    \endfoot

  qz & - & 1 \\\hline
    qx & 1 & - \\\hline
    qy & 1 & -

\end{longtable}

\textbf{Format}\\

\begin{figure*}[!htbp]
 {\setlength{\tabcolsep}{1pt}
 \begin{center}
\begin{tabular}{@{}F@{}F@{}K@{}F@{}O}
\instbitrange{31}{29}&
\instbitrange{28}{26}&
\instbitrange{25}{10}&
\instbitrange{9}{7}&
\instbitrange{6}{0} \\
\hline
\multicolumn{1}{|c|}{{\scriptsize qx[2:0]}} &
\multicolumn{1}{c|}{{\scriptsize qy[2:0]}} &
\multicolumn{1}{c|}{110100X101X0001X} &
\multicolumn{1}{c|}{{\scriptsize qz[2:0]}} &
\multicolumn{1}{c|}{0011111}\\
\hline
3& 3& 16& 3& 7 \\
\end{tabular}
 \end{center}
 }
 \vspace{-0.1in}
 \caption[]{ESP.VCMP.GT.S32} \end{figure*}


\newpage
\subsubsection{ESP.VCMP.GT.S8}\label{instruction:ESPVCMPGTS8}
\textbf{Syntax}\\
ESP.VCMP.GT.S8 qz,qx,qy

\textbf{Description}\\
This instruction compares 8-bit signed vector data. It compares the numerical values of the 16 8-bit data segments in the registers \lstinline{qx} and \lstinline{qy}. If the former is larger than the latter, it writes 0xFF into the corresponding 8-bit data segment in the register \lstinline{qz}. Otherwise, it writes 0 into the segment.\\
\textbf{Operation}
\begin{lstlisting}[escapeinside=`?]
  qz[  7:  0] = (qx[  7:  0]>qy[  7:  0]) ? 0xFF : 0
  qz[ 15:  8] = (qx[ 15:  8]>qy[ 15:  8]) ? 0xFF : 0
  ...
  qz[127:120] = (qx[127:120]>qy[127:120]) ? 0xFF : 0
\end{lstlisting}

\textbf{Schedule}\\
\begin{longtable}[c]{|l|l|l|}
    \hline
    \rowcolor{lightgray}
    \textbf{Operands} & \textbf{Use} & \textbf{Def} \\\hline
    \endfirsthead
    \hline
    \rowcolor{lightgray}
     \textbf{Operands} & \textbf{Use} & \textbf{Def} \\\hline
    \endhead
    \hline
    \endfoot

  qz & - & 1 \\\hline
    qx & 1 & - \\\hline
    qy & 1 & -

\end{longtable}

\textbf{Format}\\

\begin{figure*}[!htbp]
 {\setlength{\tabcolsep}{1pt}
 \begin{center}
\begin{tabular}{@{}F@{}F@{}K@{}F@{}O}
\instbitrange{31}{29}&
\instbitrange{28}{26}&
\instbitrange{25}{10}&
\instbitrange{9}{7}&
\instbitrange{6}{0} \\
\hline
\multicolumn{1}{|c|}{{\scriptsize qx[2:0]}} &
\multicolumn{1}{c|}{{\scriptsize qy[2:0]}} &
\multicolumn{1}{c|}{110100X100100X0X} &
\multicolumn{1}{c|}{{\scriptsize qz[2:0]}} &
\multicolumn{1}{c|}{0011111}\\
\hline
3& 3& 16& 3& 7 \\
\end{tabular}
 \end{center}
 }
 \vspace{-0.1in}
 \caption[]{ESP.VCMP.GT.S8} \end{figure*}


\newpage
\subsubsection{ESP.VCMP.GT.U16}\label{instruction:ESPVCMPGTU16}
\textbf{Syntax}\\
ESP.VCMP.GT.U16 qz,qx,qy

\textbf{Description}\\
This instruction compares 16-bit unsigned vector data. It compares the numerical values of the eight 16-bit data segments in the registers \lstinline{qx} and \lstinline{qy}. If the former is larger than the latter, it writes 0xFFFF into the corresponding 16-bit data segment in the register \lstinline{qz}. Otherwise, it writes 0 into the segment.\\
\textbf{Operation}
\begin{lstlisting}[escapeinside=`?]
  qz[ 15:  0] = (qx[ 15:  0]>qy[ 15:  0]) ? 0xFFFF : 0
  qz[ 31: 16] = (qx[ 31: 16]>qy[ 31: 16]) ? 0xFFFF : 0
  ...
  qz[127:112] = (qx[127:112]>qy[127:112]) ? 0xFFFF : 0
\end{lstlisting}

\textbf{Schedule}\\
\begin{longtable}[c]{|l|l|l|}
    \hline
    \rowcolor{lightgray}
    \textbf{Operands} & \textbf{Use} & \textbf{Def} \\\hline
    \endfirsthead
    \hline
    \rowcolor{lightgray}
     \textbf{Operands} & \textbf{Use} & \textbf{Def} \\\hline
    \endhead
    \hline
    \endfoot

  qz & - & 1 \\\hline
    qx & 1 & - \\\hline
    qy & 1 & -

\end{longtable}

\textbf{Format}\\

\begin{figure*}[!htbp]
 {\setlength{\tabcolsep}{1pt}
 \begin{center}
\begin{tabular}{@{}F@{}F@{}K@{}F@{}O}
\instbitrange{31}{29}&
\instbitrange{28}{26}&
\instbitrange{25}{10}&
\instbitrange{9}{7}&
\instbitrange{6}{0} \\
\hline
\multicolumn{1}{|c|}{{\scriptsize qx[2:0]}} &
\multicolumn{1}{c|}{{\scriptsize qy[2:0]}} &
\multicolumn{1}{c|}{110100X101000X0X} &
\multicolumn{1}{c|}{{\scriptsize qz[2:0]}} &
\multicolumn{1}{c|}{0011111}\\
\hline
3& 3& 16& 3& 7 \\
\end{tabular}
 \end{center}
 }
 \vspace{-0.1in}
 \caption[]{ESP.VCMP.GT.U16} \end{figure*}


\newpage
\subsubsection{ESP.VCMP.GT.U32}\label{instruction:ESPVCMPGTU32}
\textbf{Syntax}\\
ESP.VCMP.GT.U32 qz,qx,qy

\textbf{Description}\\
This instruction compares 32-bit unsigned vector data. It compares the numerical values of the four 32-bit data segments in registers \lstinline{qx} and \lstinline{qy}. If the former is larger than the latter, it writes 0xFFFFFFFF into the corresponding 32-bit data segment in the register \lstinline{qz}. Otherwise, it writes 0 into the segment.\\
\textbf{Operation}
\begin{lstlisting}[escapeinside=`?]
  qz[ 31:  0] = (qx[ 31:  0]>qy[ 31:  0]) ? 0xFFFFFFFF : 0
  qz[ 63: 32] = (qx[ 63: 32]>qy[ 63: 32]) ? 0xFFFFFFFF : 0
  ...
  qz[127: 96] = (qx[127: 96]>qy[127: 96]) ? 0xFFFFFFFF : 0
\end{lstlisting}

\textbf{Schedule}\\
\begin{longtable}[c]{|l|l|l|}
    \hline
    \rowcolor{lightgray}
    \textbf{Operands} & \textbf{Use} & \textbf{Def} \\\hline
    \endfirsthead
    \hline
    \rowcolor{lightgray}
     \textbf{Operands} & \textbf{Use} & \textbf{Def} \\\hline
    \endhead
    \hline
    \endfoot

  qz & - & 1 \\\hline
    qx & 1 & - \\\hline
    qy & 1 & -

\end{longtable}

\textbf{Format}\\

\begin{figure*}[!htbp]
 {\setlength{\tabcolsep}{1pt}
 \begin{center}
\begin{tabular}{@{}F@{}F@{}K@{}F@{}O}
\instbitrange{31}{29}&
\instbitrange{28}{26}&
\instbitrange{25}{10}&
\instbitrange{9}{7}&
\instbitrange{6}{0} \\
\hline
\multicolumn{1}{|c|}{{\scriptsize qx[2:0]}} &
\multicolumn{1}{c|}{{\scriptsize qy[2:0]}} &
\multicolumn{1}{c|}{110100X100X0001X} &
\multicolumn{1}{c|}{{\scriptsize qz[2:0]}} &
\multicolumn{1}{c|}{0011111}\\
\hline
3& 3& 16& 3& 7 \\
\end{tabular}
 \end{center}
 }
 \vspace{-0.1in}
 \caption[]{ESP.VCMP.GT.U32} \end{figure*}


\newpage
\subsubsection{ESP.VCMP.GT.U8}\label{instruction:ESPVCMPGTU8}
\textbf{Syntax}\\
ESP.VCMP.GT.U8 qz,qx,qy

\textbf{Description}\\
This instruction compares 8-bit unsigned vector data. It compares the numerical values of the 16 8-bit data segments in registers \lstinline{qx} and \lstinline{qy}. If the former is larger than the latter, it writes 0xFF into the corresponding 8-bit data segment in the register \lstinline{qz}. Otherwise, it writes 0 into the segment.\\
\textbf{Operation}
\begin{lstlisting}[escapeinside=`?]
  qz[  7:  0] = (qx[  7:  0]>qy[  7:  0]) ? 0xFF : 0
  qz[ 15:  8] = (qx[ 15:  8]>qy[ 15:  8]) ? 0xFF : 0
  ...
  qz[127:120] = (qx[127:120]>qy[127:120]) ? 0xFF : 0
\end{lstlisting}

\textbf{Schedule}\\
\begin{longtable}[c]{|l|l|l|}
    \hline
    \rowcolor{lightgray}
    \textbf{Operands} & \textbf{Use} & \textbf{Def} \\\hline
    \endfirsthead
    \hline
    \rowcolor{lightgray}
     \textbf{Operands} & \textbf{Use} & \textbf{Def} \\\hline
    \endhead
    \hline
    \endfoot

  qz & - & 1 \\\hline
    qx & 1 & - \\\hline
    qy & 1 & -

\end{longtable}

\textbf{Format}\\

\begin{figure*}[!htbp]
 {\setlength{\tabcolsep}{1pt}
 \begin{center}
\begin{tabular}{@{}F@{}F@{}K@{}F@{}O}
\instbitrange{31}{29}&
\instbitrange{28}{26}&
\instbitrange{25}{10}&
\instbitrange{9}{7}&
\instbitrange{6}{0} \\
\hline
\multicolumn{1}{|c|}{{\scriptsize qx[2:0]}} &
\multicolumn{1}{c|}{{\scriptsize qy[2:0]}} &
\multicolumn{1}{c|}{110100X100000X0X} &
\multicolumn{1}{c|}{{\scriptsize qz[2:0]}} &
\multicolumn{1}{c|}{0011111}\\
\hline
3& 3& 16& 3& 7 \\
\end{tabular}
 \end{center}
 }
 \vspace{-0.1in}
 \caption[]{ESP.VCMP.GT.U8} \end{figure*}


\newpage
\subsubsection{ESP.VCMP.LT.S16}\label{instruction:ESPVCMPLTS16}
\textbf{Syntax}\\
ESP.VCMP.LT.S16 qz,qx,qy

\textbf{Description}\\
This instruction compares 16-bit signed vector data. It compares the numerical values of the eight 16-bit data segments in registers \lstinline{qx} and \lstinline{qy}. If the former is smaller than the latter, it writes 0xFFFF into the corresponding 16-bit data segment in register \lstinline{qz}. Otherwise, it writes 0 into the segment.\\
\textbf{Operation}
\begin{lstlisting}[escapeinside=`?]
  qz[ 15:  0] = (qx[ 15:  0]<qy[ 15:  0]) ? 0xFFFF : 0
  qz[ 31: 16] = (qx[ 31: 16]<qy[ 31: 16]) ? 0xFFFF : 0
  ...
  qz[127:112] = (qx[127:112]<qy[127:112]) ? 0xFFFF : 0
\end{lstlisting}

\textbf{Schedule}\\
\begin{longtable}[c]{|l|l|l|}
    \hline
    \rowcolor{lightgray}
    \textbf{Operands} & \textbf{Use} & \textbf{Def} \\\hline
    \endfirsthead
    \hline
    \rowcolor{lightgray}
     \textbf{Operands} & \textbf{Use} & \textbf{Def} \\\hline
    \endhead
    \hline
    \endfoot

  qz & - & 1 \\\hline
    qx & 1 & - \\\hline
    qy & 1 & -

\end{longtable}

\textbf{Format}\\

\begin{figure*}[!htbp]
 {\setlength{\tabcolsep}{1pt}
 \begin{center}
\begin{tabular}{@{}F@{}F@{}K@{}F@{}O}
\instbitrange{31}{29}&
\instbitrange{28}{26}&
\instbitrange{25}{10}&
\instbitrange{9}{7}&
\instbitrange{6}{0} \\
\hline
\multicolumn{1}{|c|}{{\scriptsize qx[2:0]}} &
\multicolumn{1}{c|}{{\scriptsize qy[2:0]}} &
\multicolumn{1}{c|}{110100X011100X0X} &
\multicolumn{1}{c|}{{\scriptsize qz[2:0]}} &
\multicolumn{1}{c|}{0011111}\\
\hline
3& 3& 16& 3& 7 \\
\end{tabular}
 \end{center}
 }
 \vspace{-0.1in}
 \caption[]{ESP.VCMP.LT.S16} \end{figure*}


\newpage
\subsubsection{ESP.VCMP.LT.S32}\label{instruction:ESPVCMPLTS32}
\textbf{Syntax}\\
ESP.VCMP.LT.S32 qz,qx,qy

\textbf{Description}\\
This instruction compares 32-bit signed vector data. It performs a comparison of the numerical values of the four 32-bit data segments in registers \lstinline{qx} and \lstinline{qy}. If the former is smaller than the latter, it writes 0xFFFFFFFF into the corresponding 32-bit data segment in register \lstinline{qz}. Otherwise, it writes 0 into the segment.\\
\textbf{Operation}
\begin{lstlisting}[escapeinside=`?]
  qz[ 31:  0] = (qx[ 31:  0]<qy[ 31:  0]) ? 0xFFFFFFFF : 0
  qz[ 63: 32] = (qx[ 63: 32]<qy[ 63: 32]) ? 0xFFFFFFFF : 0
  ...
  qz[127: 96] = (qx[127: 96]<qy[127: 96]) ? 0xFFFFFFFF : 0
\end{lstlisting}

\textbf{Schedule}\\
\begin{longtable}[c]{|l|l|l|}
    \hline
    \rowcolor{lightgray}
    \textbf{Operands} & \textbf{Use} & \textbf{Def} \\\hline
    \endfirsthead
    \hline
    \rowcolor{lightgray}
     \textbf{Operands} & \textbf{Use} & \textbf{Def} \\\hline
    \endhead
    \hline
    \endfoot

  qz & - & 1 \\\hline
    qx & 1 & - \\\hline
    qy & 1 & -

\end{longtable}

\textbf{Format}\\

\begin{figure*}[!htbp]
 {\setlength{\tabcolsep}{1pt}
 \begin{center}
\begin{tabular}{@{}F@{}F@{}K@{}F@{}O}
\instbitrange{31}{29}&
\instbitrange{28}{26}&
\instbitrange{25}{10}&
\instbitrange{9}{7}&
\instbitrange{6}{0} \\
\hline
\multicolumn{1}{|c|}{{\scriptsize qx[2:0]}} &
\multicolumn{1}{c|}{{\scriptsize qy[2:0]}} &
\multicolumn{1}{c|}{110100X011X0001X} &
\multicolumn{1}{c|}{{\scriptsize qz[2:0]}} &
\multicolumn{1}{c|}{0011111}\\
\hline
3& 3& 16& 3& 7 \\
\end{tabular}
 \end{center}
 }
 \vspace{-0.1in}
 \caption[]{ESP.VCMP.LT.S32} \end{figure*}


\newpage
\subsubsection{ESP.VCMP.LT.S8}\label{instruction:ESPVCMPLTS8}
\textbf{Syntax}\\
ESP.VCMP.LT.S8 qz,qx,qy

\textbf{Description}\\
This instruction compares 8-bit signed vector data. It compares the numerical values of the 16 8-bit data segments in registers \lstinline{qx} and \lstinline{qy}. If the former is smaller than the latter, it writes 0xFF into the corresponding 8-bit data segment in register \lstinline{qz}. Otherwise, it writes 0 into the segment.\\
\textbf{Operation}
\begin{lstlisting}[escapeinside=`?]
  qz[  7:  0] = (qx[  7:  0]<qy[  7:  0]) ? 0xFF : 0
  qz[ 15:  8] = (qx[ 15:  8]<qy[ 15:  8]) ? 0xFF : 0
  ...
  qz[127:120] = (qx[127:120]<qy[127:120]) ? 0xFF : 0
\end{lstlisting}

\textbf{Schedule}\\
\begin{longtable}[c]{|l|l|l|}
    \hline
    \rowcolor{lightgray}
    \textbf{Operands} & \textbf{Use} & \textbf{Def} \\\hline
    \endfirsthead
    \hline
    \rowcolor{lightgray}
     \textbf{Operands} & \textbf{Use} & \textbf{Def} \\\hline
    \endhead
    \hline
    \endfoot

  qz & - & 1 \\\hline
    qx & 1 & - \\\hline
    qy & 1 & -

\end{longtable}

\textbf{Format}\\

\begin{figure*}[!htbp]
 {\setlength{\tabcolsep}{1pt}
 \begin{center}
\begin{tabular}{@{}F@{}F@{}K@{}F@{}O}
\instbitrange{31}{29}&
\instbitrange{28}{26}&
\instbitrange{25}{10}&
\instbitrange{9}{7}&
\instbitrange{6}{0} \\
\hline
\multicolumn{1}{|c|}{{\scriptsize qx[2:0]}} &
\multicolumn{1}{c|}{{\scriptsize qy[2:0]}} &
\multicolumn{1}{c|}{110100X010100X0X} &
\multicolumn{1}{c|}{{\scriptsize qz[2:0]}} &
\multicolumn{1}{c|}{0011111}\\
\hline
3& 3& 16& 3& 7 \\
\end{tabular}
 \end{center}
 }
 \vspace{-0.1in}
 \caption[]{ESP.VCMP.LT.S8} \end{figure*}


\newpage
\subsubsection{ESP.VCMP.LT.U16}\label{instruction:ESPVCMPLTU16}
\textbf{Syntax}\\
ESP.VCMP.LT.U16 qz,qx,qy

\textbf{Description}\\
This instruction compares 16-bit unsigned vector data. It compares the numerical values of the eight 16-bit data segments in the registers \lstinline{qx} and \lstinline{qy}. If the former is smaller than the latter, it writes 0xFFFF into the corresponding 16-bit data segment in the register \lstinline{qz}. Otherwise, it writes 0 into the segment.\\
\textbf{Operation}
\begin{lstlisting}[escapeinside=`?]
  qz[ 15:  0] = (qx[ 15:  0]<qy[ 15:  0]) ? 0xFFFF : 0
  qz[ 31: 16] = (qx[ 31: 16]<qy[ 31: 16]) ? 0xFFFF : 0
  ...
  qz[127:112] = (qx[127:112]<qy[127:112]) ? 0xFFFF : 0
\end{lstlisting}

\textbf{Schedule}\\
\begin{longtable}[c]{|l|l|l|}
    \hline
    \rowcolor{lightgray}
    \textbf{Operands} & \textbf{Use} & \textbf{Def} \\\hline
    \endfirsthead
    \hline
    \rowcolor{lightgray}
     \textbf{Operands} & \textbf{Use} & \textbf{Def} \\\hline
    \endhead
    \hline
    \endfoot

  qz & - & 1 \\\hline
    qx & 1 & - \\\hline
    qy & 1 & -

\end{longtable}

\textbf{Format}\\

\begin{figure*}[!htbp]
 {\setlength{\tabcolsep}{1pt}
 \begin{center}
\begin{tabular}{@{}F@{}F@{}K@{}F@{}O}
\instbitrange{31}{29}&
\instbitrange{28}{26}&
\instbitrange{25}{10}&
\instbitrange{9}{7}&
\instbitrange{6}{0} \\
\hline
\multicolumn{1}{|c|}{{\scriptsize qx[2:0]}} &
\multicolumn{1}{c|}{{\scriptsize qy[2:0]}} &
\multicolumn{1}{c|}{110100X011000X0X} &
\multicolumn{1}{c|}{{\scriptsize qz[2:0]}} &
\multicolumn{1}{c|}{0011111}\\
\hline
3& 3& 16& 3& 7 \\
\end{tabular}
 \end{center}
 }
 \vspace{-0.1in}
 \caption[]{ESP.VCMP.LT.U16} \end{figure*}


\newpage
\subsubsection{ESP.VCMP.LT.U32}\label{instruction:ESPVCMPLTU32}
\textbf{Syntax}\\
ESP.VCMP.LT.U32 qz,qx,qy

\textbf{Description}\\
This instruction compares 32-bit unsigned vector data. It compares the numerical values of the four 32-bit data segments in registers \lstinline{qx} and \lstinline{qy}. If the former is smaller than the latter, it writes 0xFFFFFFFF into the corresponding 32-bit data segment in register \lstinline{qz}. Otherwise, it writes 0 into the segment.\\
\textbf{Operation}
\begin{lstlisting}[escapeinside=`?]
  qz[ 31:  0] = (qx[ 31:  0]<qy[ 31:  0]) ? 0xFFFFFFFF : 0
  qz[ 63: 32] = (qx[ 63: 32]<qy[ 63: 32]) ? 0xFFFFFFFF : 0
  ...
  qz[127: 96] = (qx[127: 96]<qy[127: 96]) ? 0xFFFFFFFF : 0
\end{lstlisting}

\textbf{Schedule}\\
\begin{longtable}[c]{|l|l|l|}
    \hline
    \rowcolor{lightgray}
    \textbf{Operands} & \textbf{Use} & \textbf{Def} \\\hline
    \endfirsthead
    \hline
    \rowcolor{lightgray}
     \textbf{Operands} & \textbf{Use} & \textbf{Def} \\\hline
    \endhead
    \hline
    \endfoot

  qz & - & 1 \\\hline
    qx & 1 & - \\\hline
    qy & 1 & -

\end{longtable}

\textbf{Format}\\

\begin{figure*}[!htbp]
 {\setlength{\tabcolsep}{1pt}
 \begin{center}
\begin{tabular}{@{}F@{}F@{}K@{}F@{}O}
\instbitrange{31}{29}&
\instbitrange{28}{26}&
\instbitrange{25}{10}&
\instbitrange{9}{7}&
\instbitrange{6}{0} \\
\hline
\multicolumn{1}{|c|}{{\scriptsize qx[2:0]}} &
\multicolumn{1}{c|}{{\scriptsize qy[2:0]}} &
\multicolumn{1}{c|}{110100X010X0001X} &
\multicolumn{1}{c|}{{\scriptsize qz[2:0]}} &
\multicolumn{1}{c|}{0011111}\\
\hline
3& 3& 16& 3& 7 \\
\end{tabular}
 \end{center}
 }
 \vspace{-0.1in}
 \caption[]{ESP.VCMP.LT.U32} \end{figure*}


\newpage
\subsubsection{ESP.VCMP.LT.U8}\label{instruction:ESPVCMPLTU8}
\textbf{Syntax}\\
ESP.VCMP.LT.U8 qz,qx,qy

\textbf{Description}\\
This instruction compares 8-bit unsigned vector data. It compares the numerical values of the 16 8-bit data segments in registers \lstinline{qx} and \lstinline{qy}. If the former is smaller than the latter, it writes 0xFF into the corresponding 8-bit data segment in register \lstinline{qz}. Otherwise, it writes 0 into the segment.\\
\textbf{Operation}
\begin{lstlisting}[escapeinside=`?]
  qz[  7:  0] = (qx[  7:  0]<qy[  7:  0]) ? 0xFF : 0
  qz[ 15:  8] = (qx[ 15:  8]<qy[ 15:  8]) ? 0xFF : 0
  ...
  qz[127:120] = (qx[127:120]<qy[127:120]) ? 0xFF : 0
\end{lstlisting}

\textbf{Schedule}\\
\begin{longtable}[c]{|l|l|l|}
    \hline
    \rowcolor{lightgray}
    \textbf{Operands} & \textbf{Use} & \textbf{Def} \\\hline
    \endfirsthead
    \hline
    \rowcolor{lightgray}
     \textbf{Operands} & \textbf{Use} & \textbf{Def} \\\hline
    \endhead
    \hline
    \endfoot

  qz & - & 1 \\\hline
    qx & 1 & - \\\hline
    qy & 1 & -

\end{longtable}

\textbf{Format}\\

\begin{figure*}[!htbp]
 {\setlength{\tabcolsep}{1pt}
 \begin{center}
\begin{tabular}{@{}F@{}F@{}K@{}F@{}O}
\instbitrange{31}{29}&
\instbitrange{28}{26}&
\instbitrange{25}{10}&
\instbitrange{9}{7}&
\instbitrange{6}{0} \\
\hline
\multicolumn{1}{|c|}{{\scriptsize qx[2:0]}} &
\multicolumn{1}{c|}{{\scriptsize qy[2:0]}} &
\multicolumn{1}{c|}{110100X010000X0X} &
\multicolumn{1}{c|}{{\scriptsize qz[2:0]}} &
\multicolumn{1}{c|}{0011111}\\
\hline
3& 3& 16& 3& 7 \\
\end{tabular}
 \end{center}
 }
 \vspace{-0.1in}
 \caption[]{ESP.VCMP.LT.U8} \end{figure*}


\newpage
\subsubsection{ESP.MOV.S16.QACC}\label{instruction:ESPMOVS16QACC}
\textbf{Syntax}\\
ESP.MOV.S16.QACC qu

\textbf{Description}\\
This instruction signed-extends the 8 segments of 16-bit data in the register \lstinline{qs} to 64 bits and writes the results to the special registers \lstinline{QACC_H} and \lstinline{QACC_L}.

\textbf{Operation}\\
\begin{lstlisting}
    QACC_L[ 63:  0] = {48{qu[15]}, qu[ 15:  0]}
    QACC_L[127: 64] = {48{qu[31]}, qu[ 31: 16]}
    QACC_L[191:128] = {48{qu[47]}, qu[ 47: 32]}
    QACC_L[255:192] = {48{qu[63]}, qu[ 63: 48]}
    QACC_H[ 39:  0] = {48{qu[79]}, qu[ 79: 64]}
    QACC_H[ 79: 40] = {48{qu[95]}, qu[ 95: 80]}
    QACC_H[119: 80] = {48{qu[111]}, qu[111: 96]}
    QACC_H[159:120] = {48{qu[127]}, qu[127:112]}
\end{lstlisting}

\textbf{Schedule}\\
\begin{longtable}[c]{|l|l|l|}
    \hline
    \rowcolor{lightgray}
    \textbf{Operands} & \textbf{Use} & \textbf{Def} \\\hline
    \endfirsthead
    \hline
    \rowcolor{lightgray}
     \textbf{Operands} & \textbf{Use} & \textbf{Def} \\\hline
    \endhead
    \hline
    \endfoot

  qu & 2 & - \\\hline
    QACC\_H & - & 2 \\\hline
    QACC\_L & - & 2

\end{longtable}

\textbf{Format}\\

\begin{figure*}[!htbp]
 {\setlength{\tabcolsep}{1pt}
 \begin{center}
\begin{tabular}{@{}U@{}F@{}F@{}O}
\instbitrange{31}{13}&
\instbitrange{12}{10}&
\instbitrange{9}{7}&
\instbitrange{6}{0} \\
\hline
\multicolumn{1}{|c|}{1XXXXXXXXXXXXXXXX00} &
\multicolumn{1}{c|}{{\scriptsize qu[2:0]}} &
\multicolumn{1}{c|}{11X} &
\multicolumn{1}{c|}{0011011}\\
\hline
19& 3& 3& 7 \\
\end{tabular}
 \end{center}
 }
 \vspace{-0.1in}
 \caption[]{ESP.MOV.S16.QACC} \end{figure*}


\newpage
\subsubsection{ESP.MOV.S8.QACC}\label{instruction:ESPMOVS8QACC}
\textbf{Syntax}\\
ESP.MOV.S8.QACC qu

\textbf{Description}\\
This instruction sign-extends the 16 segments of 8-bit data in the register \lstinline{qs} to 32 bits and writes the result to the special registers \lstinline{QACC_H} and \lstinline{QACC_L}.

\textbf{Operation}\\
\begin{lstlisting}
  QACC_L[ 31:  0] = {24{qu[7]}, qu[  7:  0]}
  QACC_L[ 63: 32] = {24{qu[15]}, qu[ 15:  8]}
  QACC_L[ 95: 64] = {24{qu[23]}, qu[ 23: 16]}
  ...
  QACC_L[255:224] = {24{qu[63]}, qu[ 63: 56]}
  QACC_H[ 31:  0] = {24{qu[71]}, qu[ 71: 64]}
  QACC_H[ 63: 32] = {24{qu[79]}, qu[ 79: 72]}
  QACC_H[ 95: 64] = {24{qu[87]}, qu[ 87: 80]}
  ...
  QACC_H[255:224] = {24{qu[127]}, qu[127:120]}
\end{lstlisting}

\textbf{Schedule}\\
\begin{longtable}[c]{|l|l|l|}
    \hline
    \rowcolor{lightgray}
    \textbf{Operands} & \textbf{Use} & \textbf{Def} \\\hline
    \endfirsthead
    \hline
    \rowcolor{lightgray}
     \textbf{Operands} & \textbf{Use} & \textbf{Def} \\\hline
    \endhead
    \hline
    \endfoot

  qu & 2 & - \\\hline
    QACC\_H & - & 2 \\\hline
    QACC\_L & - & 2

\end{longtable}

\textbf{Format}\\

\begin{figure*}[!htbp]
 {\setlength{\tabcolsep}{1pt}
 \begin{center}
\begin{tabular}{@{}U@{}F@{}F@{}O}
\instbitrange{31}{13}&
\instbitrange{12}{10}&
\instbitrange{9}{7}&
\instbitrange{6}{0} \\
\hline
\multicolumn{1}{|c|}{1XXXXXXXXXXXXXXXX00} &
\multicolumn{1}{c|}{{\scriptsize qu[2:0]}} &
\multicolumn{1}{c|}{01X} &
\multicolumn{1}{c|}{0011011}\\
\hline
19& 3& 3& 7 \\
\end{tabular}
 \end{center}
 }
 \vspace{-0.1in}
 \caption[]{ESP.MOV.S8.QACC} \end{figure*}


\newpage
\subsubsection{ESP.MOV.U16.QACC}\label{instruction:ESPMOVU16QACC}
\textbf{Syntax}\\
ESP.MOV.U16.QACC qu

\textbf{Description}\\
This instruction zero-extends the 8 segments of 16-bit data in the register \lstinline{qs} to 64 bits and writes the results to the special registers \lstinline{QACC_H} and \lstinline{QACC_L}.

\textbf{Operation}\\
\begin{lstlisting}
    QACC_L[ 63:  0] = {48{0}, qu[ 15:  0]}
    QACC_L[127: 64] = {48{0}, qu[ 31: 16]}
    QACC_L[191:128] = {48{0}, qu[ 47: 32]}
    QACC_L[255:192] = {48{0}, qu[ 63: 48]}
    QACC_H[ 39:  0] = {48{0}, qu[ 79: 64]}
    QACC_H[ 79: 40] = {48{0}, qu[ 95: 80]}
    QACC_H[119: 80] = {48{0}, qu[111: 96]}
    QACC_H[159:120] = {48{0}, qu[127:112]}
\end{lstlisting}

\textbf{Schedule}\\
\begin{longtable}[c]{|l|l|l|}
    \hline
    \rowcolor{lightgray}
    \textbf{Operands} & \textbf{Use} & \textbf{Def} \\\hline
    \endfirsthead
    \hline
    \rowcolor{lightgray}
     \textbf{Operands} & \textbf{Use} & \textbf{Def} \\\hline
    \endhead
    \hline
    \endfoot

  qu & 2 & - \\\hline
    QACC\_H & - & 2 \\\hline
    QACC\_L & - & 2

\end{longtable}

\textbf{Format}\\

\begin{figure*}[!htbp]
 {\setlength{\tabcolsep}{1pt}
 \begin{center}
\begin{tabular}{@{}U@{}F@{}F@{}O}
\instbitrange{31}{13}&
\instbitrange{12}{10}&
\instbitrange{9}{7}&
\instbitrange{6}{0} \\
\hline
\multicolumn{1}{|c|}{1XXXXXXXXXXXXXXXX00} &
\multicolumn{1}{c|}{{\scriptsize qu[2:0]}} &
\multicolumn{1}{c|}{10X} &
\multicolumn{1}{c|}{0011011}\\
\hline
19& 3& 3& 7 \\
\end{tabular}
 \end{center}
 }
 \vspace{-0.1in}
 \caption[]{ESP.MOV.U16.QACC} \end{figure*}


\newpage
\subsubsection{ESP.MOV.U8.QACC}\label{instruction:ESPMOVU8QACC}
\textbf{Syntax}\\
ESP.MOV.U8.QACC qu

\textbf{Description}\\
This instruction zero-extends the 16 segments of 8-bit data in the register \lstinline{qs} to 32 bits and writes the result to the special registers \lstinline{QACC_H} and \lstinline{QACC_L}.

\textbf{Operation}\\
\begin{lstlisting}
  QACC_L[ 31:  0] = {24{0}, qu[  7:  0]}
  QACC_L[ 63: 32] = {24{0}, qu[ 15:  8]}
  QACC_L[ 95: 64] = {24{0}, qu[ 23: 16]}
  ...
  QACC_L[255:224] = {24{0}, qu[ 63: 56]}
  QACC_H[ 31:  0] = {24{0}, qu[ 71: 64]}
  QACC_H[ 63: 32] = {24{0}, qu[ 79: 72]}
  QACC_H[ 95: 64] = {24{0}, qu[ 87: 80]}
  ...
  QACC_H[255:224] = {24{0}, qu[127:120]}
\end{lstlisting}

\textbf{Schedule}\\
\begin{longtable}[c]{|l|l|l|}
    \hline
    \rowcolor{lightgray}
    \textbf{Operands} & \textbf{Use} & \textbf{Def} \\\hline
    \endfirsthead
    \hline
    \rowcolor{lightgray}
     \textbf{Operands} & \textbf{Use} & \textbf{Def} \\\hline
    \endhead
    \hline
    \endfoot

  qu & 2 & - \\\hline
    QACC\_H & - & 2 \\\hline
    QACC\_L & - & 2

\end{longtable}

\textbf{Format}\\

\begin{figure*}[!htbp]
 {\setlength{\tabcolsep}{1pt}
 \begin{center}
\begin{tabular}{@{}U@{}F@{}F@{}O}
\instbitrange{31}{13}&
\instbitrange{12}{10}&
\instbitrange{9}{7}&
\instbitrange{6}{0} \\
\hline
\multicolumn{1}{|c|}{1XXXXXXXXXXXXXXXX00} &
\multicolumn{1}{c|}{{\scriptsize qu[2:0]}} &
\multicolumn{1}{c|}{00X} &
\multicolumn{1}{c|}{0011011}\\
\hline
19& 3& 3& 7 \\
\end{tabular}
 \end{center}
 }
 \vspace{-0.1in}
 \caption[]{ESP.MOV.U8.QACC} \end{figure*}


\newpage
\subsubsection{ESP.MOVI.16.A}\label{instruction:ESPMOVI16A}
\textbf{Syntax}\\
ESP.MOVI.16.A qy,rd,sel16

\textbf{Description}\\
This instruction extracts 16 bits from the sel index of qy to the \lstinline{rd} register.

\textbf{Operation}\\
\begin{lstlisting}
    if sel == 0:    rd[31 :  0] = { 16’b0, qy[15 :  0] }
    else if sel==1: rd[31 :  0] = { 16’b0, qy[31 : 16] }
    else if sel==2: rd[31 :  0] = { 16’b0, qy[47 : 32] }
    else if sel==3: rd[31 :  0] = { 16’b0, qy[63 : 48] }
    else if sel==4: rd[31 :  0] = { 16’b0, qy[79 : 64] }
    else if sel==5: rd[31 :  0] = { 16’b0, qy[95 : 80] }
    else if sel==6: rd[31 :  0] = { 16’b0, qy[111: 96] }
    else if sel==7: rd[31 :  0] = { 16’b0, qy[127:112] }
\end{lstlisting}

\textbf{Schedule}\\
\begin{longtable}[c]{|l|l|l|}
    \hline
    \rowcolor{lightgray}
    \textbf{Operands} & \textbf{Use} & \textbf{Def} \\\hline
    \endfirsthead
    \hline
    \rowcolor{lightgray}
     \textbf{Operands} & \textbf{Use} & \textbf{Def} \\\hline
    \endhead
    \hline
    \endfoot

  rd & - & 2 \\\hline
    qy & 1 & -

\end{longtable}

\textbf{Format}\\

\begin{figure*}[!htbp]
 {\setlength{\tabcolsep}{1pt}
 \begin{center}
\begin{tabular}{@{}F@{}F@{}O@{}Y@{}Y@{}I@{}F@{}O}
\instbitrange{31}{29}&
\instbitrange{28}{26}&
\instbitrange{25}{19}&
\instbitrange{18}{15}&
\instbitrange{14}{11}&
\instbit{10}&
\instbitrange{9}{7}&
\instbitrange{6}{0} \\
\hline
\multicolumn{1}{|c|}{0XX} &
\multicolumn{1}{c|}{{\scriptsize qy[2:0]}} &
\multicolumn{1}{c|}{0X1100X} &
\multicolumn{1}{c|}{{\scriptsize sel16[3:0]}} &
\multicolumn{1}{c|}{01XX} &
\multicolumn{1}{c|}{{\scriptsize rd[4]}} &
\multicolumn{1}{c|}{{\scriptsize rd[2:0]}} &
\multicolumn{1}{c|}{0011011}\\
\hline
3& 3& 7& 4& 4& 1& 3& 7 \\
\end{tabular}
 \end{center}
 }
 \vspace{-0.1in}
 \caption[]{ESP.MOVI.16.A} \end{figure*}


\newpage
\subsubsection{ESP.MOVI.16.Q}\label{instruction:ESPMOVI16Q}
\textbf{Syntax}\\
ESP.MOVI.16.Q qy,rs1,sel16

\textbf{Description}\\
This instruction inserts 16 bits into the sel index of \lstinline{qy} from the \lstinline{rs1} register.

\textbf{Operation}\\
\begin{lstlisting}
    if sel == 0:        qy[31 :  0] = { qy[31 : 16], rs1[15:0] }
    else if sel == 1:   qy[31 :  0] = { rs1[15:0], qy[15 :  0] }
    else if sel == 2:   qy[63 : 32] = { qy[63 : 48], rs1[15:0] }
    else if sel == 3:   qy[63 : 48] = { rs1[15:0], qy[47 : 32] }
    else if sel == 4:   qy[95 : 64] = { qy[95 : 80], rs1[15:0] }
    else if sel == 5:   qy[95 : 64] = { rs1[15:0], qy[79 : 64] }
    else if sel == 6:   qy[127: 96] = { qy[127:112], rs1[15:0] }
    else if sel == 7:   qy[127: 96] = { rs1[15:0], qy[111: 96] }
\end{lstlisting}

\textbf{Schedule}\\
\begin{longtable}[c]{|l|l|l|}
    \hline
    \rowcolor{lightgray}
    \textbf{Operands} & \textbf{Use} & \textbf{Def} \\\hline
    \endfirsthead
    \hline
    \rowcolor{lightgray}
     \textbf{Operands} & \textbf{Use} & \textbf{Def} \\\hline
    \endhead
    \hline
    \endfoot

  qy & 1 & 1 \\\hline
    rs1 & 1 & -

\end{longtable}

\textbf{Format}\\

\begin{figure*}[!htbp]
 {\setlength{\tabcolsep}{1pt}
 \begin{center}
\begin{tabular}{@{}F@{}F@{}O@{}I@{}F@{}Y@{}Y@{}O}
\instbitrange{31}{29}&
\instbitrange{28}{26}&
\instbitrange{25}{19}&
\instbit{18}&
\instbitrange{17}{15}&
\instbitrange{14}{11}&
\instbitrange{10}{7}&
\instbitrange{6}{0} \\
\hline
\multicolumn{1}{|c|}{0XX} &
\multicolumn{1}{c|}{{\scriptsize qy[2:0]}} &
\multicolumn{1}{c|}{0X1101X} &
\multicolumn{1}{c|}{{\scriptsize rs1[4]}} &
\multicolumn{1}{c|}{{\scriptsize rs1[2:0]}} &
\multicolumn{1}{c|}{01XX} &
\multicolumn{1}{c|}{{\scriptsize sel16[3:0]}} &
\multicolumn{1}{c|}{0011011}\\
\hline
3& 3& 7& 1& 3& 4& 4& 7 \\
\end{tabular}
 \end{center}
 }
 \vspace{-0.1in}
 \caption[]{ESP.MOVI.16.Q} \end{figure*}


\newpage
\subsubsection{ESP.MOVI.32.A}\label{instruction:ESPMOVI32A}
\textbf{Syntax}\\
ESP.MOVI.32.A qy,rd,sel4

\textbf{Description}\\
This instruction extracts 32 bits from the sel index of qv to the \lstinline{rd} register.

\textbf{Operation}\\
\begin{lstlisting}
    if sel == 0:    rd[31 :  0] = qy[31 :  0]
    else if sel==1: rd[31 :  0] = qy[63 : 32]
    else if sel==2: rd[31 :  0] = qy[95 : 64]
    else if sel==3: rd[31 :  0] = qy[127: 96]
\end{lstlisting}

\textbf{Schedule}\\
\begin{longtable}[c]{|l|l|l|}
    \hline
    \rowcolor{lightgray}
    \textbf{Operands} & \textbf{Use} & \textbf{Def} \\\hline
    \endfirsthead
    \hline
    \rowcolor{lightgray}
     \textbf{Operands} & \textbf{Use} & \textbf{Def} \\\hline
    \endhead
    \hline
    \endfoot

  rd & - & 2 \\\hline
    qy & 1 & -

\end{longtable}

\textbf{Format}\\

\begin{figure*}[!htbp]
 {\setlength{\tabcolsep}{1pt}
 \begin{center}
\begin{tabular}{@{}F@{}F@{}F@{}I@{}F@{}I@{}O@{}I@{}F@{}O}
\instbitrange{31}{29}&
\instbitrange{28}{26}&
\instbitrange{25}{23}&
\instbit{22}&
\instbitrange{21}{19}&
\instbit{18}&
\instbitrange{17}{11}&
\instbit{10}&
\instbitrange{9}{7}&
\instbitrange{6}{0} \\
\hline
\multicolumn{1}{|c|}{0XX} &
\multicolumn{1}{c|}{{\scriptsize qy[2:0]}} &
\multicolumn{1}{c|}{0X1} &
\multicolumn{1}{c|}{{\scriptsize sel4[1]}} &
\multicolumn{1}{c|}{10X} &
\multicolumn{1}{c|}{{\scriptsize sel4[0]}} &
\multicolumn{1}{c|}{XXX01XX} &
\multicolumn{1}{c|}{{\scriptsize rd[4]}} &
\multicolumn{1}{c|}{{\scriptsize rd[2:0]}} &
\multicolumn{1}{c|}{0011011}\\
\hline
3& 3& 3& 1& 3& 1& 7& 1& 3& 7 \\
\end{tabular}
 \end{center}
 }
 \vspace{-0.1in}
 \caption[]{ESP.MOVI.32.A} \end{figure*}


\newpage
\subsubsection{ESP.MOVI.32.Q}\label{instruction:ESPMOVI32Q}
\textbf{Syntax}\\
ESP.MOVI.32.Q qy,rs1,sel4

\textbf{Description}\\
This instruction inserts 32 bits into the sel index of \lstinline{qy} from the \lstinline{rs1} register.

\textbf{Operation}\\
\begin{lstlisting}
    if sel == 0:        qy[31 :  0] = rs1[31:0]
    else if sel == 1:   qy[63 : 32] = rs1[31:0]
    else if sel == 2:   qy[95 : 64] = rs1[31:0]
    else if sel == 3:   qy[127: 96] = rs1[31:0]
\end{lstlisting}

\textbf{Schedule}\\
\begin{longtable}[c]{|l|l|l|}
    \hline
    \rowcolor{lightgray}
    \textbf{Operands} & \textbf{Use} & \textbf{Def} \\\hline
    \endfirsthead
    \hline
    \rowcolor{lightgray}
     \textbf{Operands} & \textbf{Use} & \textbf{Def} \\\hline
    \endhead
    \hline
    \endfoot

  qy & 1 & 1 \\\hline
    rs1 & 1 & -

\end{longtable}

\textbf{Format}\\

\begin{figure*}[!htbp]
 {\setlength{\tabcolsep}{1pt}
 \begin{center}
\begin{tabular}{@{}F@{}F@{}O@{}I@{}F@{}Y@{}W@{}W@{}O}
\instbitrange{31}{29}&
\instbitrange{28}{26}&
\instbitrange{25}{19}&
\instbit{18}&
\instbitrange{17}{15}&
\instbitrange{14}{11}&
\instbitrange{10}{9}&
\instbitrange{8}{7}&
\instbitrange{6}{0} \\
\hline
\multicolumn{1}{|c|}{0XX} &
\multicolumn{1}{c|}{{\scriptsize qy[2:0]}} &
\multicolumn{1}{c|}{0X1X11X} &
\multicolumn{1}{c|}{{\scriptsize rs1[4]}} &
\multicolumn{1}{c|}{{\scriptsize rs1[2:0]}} &
\multicolumn{1}{c|}{01XX} &
\multicolumn{1}{c|}{{\scriptsize sel4[1:0]}} &
\multicolumn{1}{c|}{XX} &
\multicolumn{1}{c|}{0011011}\\
\hline
3& 3& 7& 1& 3& 4& 2& 2& 7 \\
\end{tabular}
 \end{center}
 }
 \vspace{-0.1in}
 \caption[]{ESP.MOVI.32.Q} \end{figure*}


\newpage
\subsubsection{ESP.MOVI.8.A}\label{instruction:ESPMOVI8A}
\textbf{Syntax}\\
ESP.MOVI.8.A qy,rd,sel16

\textbf{Description}\\
This instruction extracts 8 bits from the sel index of qv to the \lstinline{rd} register.

\textbf{Operation}\\
\begin{lstlisting}
    if sel == 0:    rd[31 :  0] = { 24’b0, qy[7  :  0] }
    else if sel==1: rd[31 :  0] = { 24’b0, qy[15 :  8] }
    else if sel==2: rd[31 :  0] = { 24’b0, qy[23 : 16] }
    else if sel==3: rd[31 :  0] = { 24’b0, qy[31 : 24] }
    else if sel==4: rd[31 :  0] = { 24’b0, qy[39 : 32] }
    else if sel==5: rd[31 :  0] = { 24’b0, qy[47 : 40] }
    else if sel==6: rd[31 :  0] = { 24’b0, qy[55 : 48] }
    else if sel==7: rd[31 :  0] = { 24’b0, qy[63 : 56] }
    else if sel==8: rd[31 :  0] = { 24’b0, qy[71 : 64] }
    else if sel==9: rd[31 :  0] = { 24’b0, qy[79 : 72] }
    else if sel==10:rd[31 :  0] = { 24’b0, qy[87 : 80] }
    else if sel==11:rd[31 :  0] = { 24’b0, qy[95 : 88] }
    else if sel==12:rd[31 :  0] = { 24’b0, qy[103: 96] }
    else if sel==13:rd[31 :  0] = { 24’b0, qy[111:104] }
    else if sel==14:rd[31 :  0] = { 24’b0, qy[119:112] }
    else if sel==15:rd[31 :  0] = { 24’b0, qy[127:120] }
\end{lstlisting}

\textbf{Schedule}\\
\begin{longtable}[c]{|l|l|l|}
    \hline
    \rowcolor{lightgray}
    \textbf{Operands} & \textbf{Use} & \textbf{Def} \\\hline
    \endfirsthead
    \hline
    \rowcolor{lightgray}
     \textbf{Operands} & \textbf{Use} & \textbf{Def} \\\hline
    \endhead
    \hline
    \endfoot

  rd & - & 2 \\\hline
    qy & 1 & -

\end{longtable}

\textbf{Format}\\

\begin{figure*}[!htbp]
 {\setlength{\tabcolsep}{1pt}
 \begin{center}
\begin{tabular}{@{}F@{}F@{}O@{}Y@{}Y@{}I@{}F@{}O}
\instbitrange{31}{29}&
\instbitrange{28}{26}&
\instbitrange{25}{19}&
\instbitrange{18}{15}&
\instbitrange{14}{11}&
\instbit{10}&
\instbitrange{9}{7}&
\instbitrange{6}{0} \\
\hline
\multicolumn{1}{|c|}{0XX} &
\multicolumn{1}{c|}{{\scriptsize qy[2:0]}} &
\multicolumn{1}{c|}{0X1000X} &
\multicolumn{1}{c|}{{\scriptsize sel16[3:0]}} &
\multicolumn{1}{c|}{01XX} &
\multicolumn{1}{c|}{{\scriptsize rd[4]}} &
\multicolumn{1}{c|}{{\scriptsize rd[2:0]}} &
\multicolumn{1}{c|}{0011011}\\
\hline
3& 3& 7& 4& 4& 1& 3& 7 \\
\end{tabular}
 \end{center}
 }
 \vspace{-0.1in}
 \caption[]{ESP.MOVI.8.A} \end{figure*}


\newpage
\subsubsection{ESP.MOVI.8.Q}\label{instruction:ESPMOVI8Q}
\textbf{Syntax}\\
ESP.MOVI.8.Q qy,rs1,sel16

\textbf{Description}\\
This instruction inserts 8 bits into the sel index of \lstinline{qy} from the \lstinline{rs1} register.

\textbf{Operation}\\
\begin{lstlisting}
    if sel == 0:        qy[31 :  0] = { qy[31 :  8], rs1[7:0] }
    else if sel == 1:   qy[31 :  0] = { qy[31 : 16], rs1[7:0], qy[7  :  0] }
    else if sel == 2:   qy[31 :  0] = { qy[31 :  0], rs1[7:0], qy[15 :  0] }
    else if sel == 3:   qy[31 : 24] = { rs1[7:0], qy[23 :  0] }
    else if sel == 4:   qy[63 : 32] = { qy[63 : 40], rs1[7:0] }
    else if sel == 5:   qy[63 : 32] = { qy[63 : 48], rs1[7:0], qy[39 : 32] }
    else if sel == 6:   qy[63 : 32] = { qy[63 : 56], rs1[7:0], qy[47 : 32] }
    else if sel == 7:   qy[63 : 32] = { rs1[7:0], qy[55 : 32] }
    else if sel == 8:   qy[95 : 64] = { qy[95 : 72], rs1[7:0] }
    else if sel == 9:   qy[95 : 64] = { qy[95 : 80], rs1[7:0], qy[71 : 64] }
    else if sel == 10:  qy[95 : 64] = { qy[95 : 88], rs1[7:0], qy[79 : 64] }
    else if sel == 11:  qy[95 : 64] = { rs1[7:0], qy[87 : 64] }
    else if sel == 12:  qy[127: 96] = { qy[127:104], rs1[7:0] }
    else if sel == 13:  qy[127: 96] = { qy[127:112], rs1[7:0], qy[103: 96] }
    else if sel == 14:  qy[127: 96] = { qy[127:120], rs1[7:0], qy[111: 96] }
    else if sel == 15:  qy[127: 96] = { rs1[7:0], qy[119: 96] }
\end{lstlisting}

\textbf{Schedule}\\
\begin{longtable}[c]{|l|l|l|}
    \hline
    \rowcolor{lightgray}
    \textbf{Operands} & \textbf{Use} & \textbf{Def} \\\hline
    \endfirsthead
    \hline
    \rowcolor{lightgray}
     \textbf{Operands} & \textbf{Use} & \textbf{Def} \\\hline
    \endhead
    \hline
    \endfoot

  qy & 1 & 1 \\\hline
    rs1 & 1 & -

\end{longtable}

\textbf{Format}\\

\begin{figure*}[!htbp]
 {\setlength{\tabcolsep}{1pt}
 \begin{center}
\begin{tabular}{@{}F@{}F@{}O@{}I@{}F@{}Y@{}Y@{}O}
\instbitrange{31}{29}&
\instbitrange{28}{26}&
\instbitrange{25}{19}&
\instbit{18}&
\instbitrange{17}{15}&
\instbitrange{14}{11}&
\instbitrange{10}{7}&
\instbitrange{6}{0} \\
\hline
\multicolumn{1}{|c|}{0XX} &
\multicolumn{1}{c|}{{\scriptsize qy[2:0]}} &
\multicolumn{1}{c|}{0X1001X} &
\multicolumn{1}{c|}{{\scriptsize rs1[4]}} &
\multicolumn{1}{c|}{{\scriptsize rs1[2:0]}} &
\multicolumn{1}{c|}{01XX} &
\multicolumn{1}{c|}{{\scriptsize sel16[3:0]}} &
\multicolumn{1}{c|}{0011011}\\
\hline
3& 3& 7& 1& 3& 4& 4& 7 \\
\end{tabular}
 \end{center}
 }
 \vspace{-0.1in}
 \caption[]{ESP.MOVI.8.Q} \end{figure*}


\newpage
\subsubsection{ESP.MOVX.R.CFG}\label{instruction:ESPMOVXRCFG}
\textbf{Syntax}\\
ESP.MOVX.R.CFG rd

\textbf{Description}\\
This instruction moves data from the special register \lstinline{CFG} to the \lstinline{rd} register.

\textbf{Operation}\\
\begin{lstlisting}
    rd[31 :  0] = CFG[31:0]
\end{lstlisting}

\textbf{Schedule}\\
\begin{longtable}[c]{|l|l|l|}
    \hline
    \rowcolor{lightgray}
    \textbf{Operands} & \textbf{Use} & \textbf{Def} \\\hline
    \endfirsthead
    \hline
    \rowcolor{lightgray}
     \textbf{Operands} & \textbf{Use} & \textbf{Def} \\\hline
    \endhead
    \hline
    \endfoot

  rd & - & 2 \\\hline
    CFG & 1 & -

\end{longtable}

\textbf{Format}\\

\begin{figure*}[!htbp]
 {\setlength{\tabcolsep}{1pt}
 \begin{center}
\begin{tabular}{@{}J@{}I@{}F@{}O}
\instbitrange{31}{11}&
\instbit{10}&
\instbitrange{9}{7}&
\instbitrange{6}{0} \\
\hline
\multicolumn{1}{|c|}{1XX0X00X1000XXXXX01XX} &
\multicolumn{1}{c|}{{\scriptsize rd[4]}} &
\multicolumn{1}{c|}{{\scriptsize rd[2:0]}} &
\multicolumn{1}{c|}{0011011}\\
\hline
21& 1& 3& 7 \\
\end{tabular}
 \end{center}
 }
 \vspace{-0.1in}
 \caption[]{ESP.MOVX.R.CFG} \end{figure*}


\newpage
\subsubsection{ESP.MOVX.R.FFT.BIT.WIDTH}\label{instruction:ESPMOVXRFFT_BIT_WIDTH}
\textbf{Syntax}\\
ESP.MOVX.R.FFT.BIT.WIDTH rd

\textbf{Description}\\
This instruction moves data from the special register \lstinline{FFT_BIT_WIDTH} to the \lstinline{rd} register.

\textbf{Operation}\\
\begin{lstlisting}
    rd[31 :  0] = {28{0}, FFT_BIT_WIDTH[3  :  0]}
\end{lstlisting}

\textbf{Schedule}\\
\begin{longtable}[c]{|l|l|l|}
    \hline
    \rowcolor{lightgray}
    \textbf{Operands} & \textbf{Use} & \textbf{Def} \\\hline
    \endfirsthead
    \hline
    \rowcolor{lightgray}
     \textbf{Operands} & \textbf{Use} & \textbf{Def} \\\hline
    \endhead
    \hline
    \endfoot

  rd & - & 2 \\\hline
    FFT\_BIT\_WIDTH & 1 & -

\end{longtable}

\textbf{Format}\\

\begin{figure*}[!htbp]
 {\setlength{\tabcolsep}{1pt}
 \begin{center}
\begin{tabular}{@{}J@{}I@{}F@{}O}
\instbitrange{31}{11}&
\instbit{10}&
\instbitrange{9}{7}&
\instbitrange{6}{0} \\
\hline
\multicolumn{1}{|c|}{1XX0010X1000XXXXX01XX} &
\multicolumn{1}{c|}{{\scriptsize rd[4]}} &
\multicolumn{1}{c|}{{\scriptsize rd[2:0]}} &
\multicolumn{1}{c|}{0011011}\\
\hline
21& 1& 3& 7 \\
\end{tabular}
 \end{center}
 }
 \vspace{-0.1in}
 \caption[]{ESP.MOVX.R.FFT.BIT.WIDTH} \end{figure*}


\newpage
\subsubsection{ESP.MOVX.R.PERF}\label{instruction:ESPMOVXRPERF}
\textbf{Syntax}\\
ESP.MOVX.R.PERF rd,rs1

\textbf{Description}\\
This instruction moves data from the special register \lstinline{PERF} to the \lstinline{rd} register.

\textbf{Operation}\\
\begin{lstlisting}
    rd[31 :  0] = PERF[31 :  0]
\end{lstlisting}

\textbf{Schedule}\\
\begin{longtable}[c]{|l|l|l|}
    \hline
    \rowcolor{lightgray}
    \textbf{Operands} & \textbf{Use} & \textbf{Def} \\\hline
    \endfirsthead
    \hline
    \rowcolor{lightgray}
     \textbf{Operands} & \textbf{Use} & \textbf{Def} \\\hline
    \endhead
    \hline
    \endfoot

  rd & - & 2 \\\hline
    PERF & 1 & -

\end{longtable}

\textbf{Format}\\

\begin{figure*}[!htbp]
 {\setlength{\tabcolsep}{1pt}
 \begin{center}
\begin{tabular}{@{}K@{}I@{}F@{}Y@{}I@{}F@{}O}
\instbitrange{31}{19}&
\instbit{18}&
\instbitrange{17}{15}&
\instbitrange{14}{11}&
\instbit{10}&
\instbitrange{9}{7}&
\instbitrange{6}{0} \\
\hline
\multicolumn{1}{|c|}{1XX0110X1000X} &
\multicolumn{1}{c|}{{\scriptsize rs1[4]}} &
\multicolumn{1}{c|}{{\scriptsize rs1[2:0]}} &
\multicolumn{1}{c|}{01XX} &
\multicolumn{1}{c|}{{\scriptsize rd[4]}} &
\multicolumn{1}{c|}{{\scriptsize rd[2:0]}} &
\multicolumn{1}{c|}{0011011}\\
\hline
13& 1& 3& 4& 1& 3& 7 \\
\end{tabular}
 \end{center}
 }
 \vspace{-0.1in}
 \caption[]{ESP.MOVX.R.PERF} \end{figure*}


\newpage
\subsubsection{ESP.MOVX.R.SAR}\label{instruction:ESPMOVXRSAR}
\textbf{Syntax}\\
ESP.MOVX.R.SAR rd

\textbf{Description}\\
This instruction moves data from the special register \lstinline{SAR} to the \lstinline{rd} register.

\textbf{Operation}\\
\begin{lstlisting}
    rd[31 :  0] = {26’b0, SAR[5:0]}
\end{lstlisting}

\textbf{Schedule}\\
\begin{longtable}[c]{|l|l|l|}
    \hline
    \rowcolor{lightgray}
    \textbf{Operands} & \textbf{Use} & \textbf{Def} \\\hline
    \endfirsthead
    \hline
    \rowcolor{lightgray}
     \textbf{Operands} & \textbf{Use} & \textbf{Def} \\\hline
    \endhead
    \hline
    \endfoot

  rd & - & 2 \\\hline
    SAR & 1 & -

\end{longtable}

\textbf{Format}\\

\begin{figure*}[!htbp]
 {\setlength{\tabcolsep}{1pt}
 \begin{center}
\begin{tabular}{@{}J@{}I@{}F@{}O}
\instbitrange{31}{11}&
\instbit{10}&
\instbitrange{9}{7}&
\instbitrange{6}{0} \\
\hline
\multicolumn{1}{|c|}{1XX0000X1100XXXXX01XX} &
\multicolumn{1}{c|}{{\scriptsize rd[4]}} &
\multicolumn{1}{c|}{{\scriptsize rd[2:0]}} &
\multicolumn{1}{c|}{0011011}\\
\hline
21& 1& 3& 7 \\
\end{tabular}
 \end{center}
 }
 \vspace{-0.1in}
 \caption[]{ESP.MOVX.R.SAR} \end{figure*}


\newpage
\subsubsection{ESP.MOVX.R.SAR.BYTES}\label{instruction:ESPMOVXRSAR_BYTES}
\textbf{Syntax}\\
ESP.MOVX.R.SAR.BYTES rd

\textbf{Description}\\
This instruction moves data from the special register \lstinline{SAR_BYTES} to the \lstinline{rd} register.

\textbf{Operation}\\
\begin{lstlisting}
    rd[31 :  0] = {28’b0, SAR_BYTES[3:0]}
\end{lstlisting}

\textbf{Schedule}\\
\begin{longtable}[c]{|l|l|l|}
    \hline
    \rowcolor{lightgray}
    \textbf{Operands} & \textbf{Use} & \textbf{Def} \\\hline
    \endfirsthead
    \hline
    \rowcolor{lightgray}
     \textbf{Operands} & \textbf{Use} & \textbf{Def} \\\hline
    \endhead
    \hline
    \endfoot

  rd & - & 2 \\\hline
    SAR\_BYTES & 1 & -

\end{longtable}

\textbf{Format}\\

\begin{figure*}[!htbp]
 {\setlength{\tabcolsep}{1pt}
 \begin{center}
\begin{tabular}{@{}J@{}I@{}F@{}O}
\instbitrange{31}{11}&
\instbit{10}&
\instbitrange{9}{7}&
\instbitrange{6}{0} \\
\hline
\multicolumn{1}{|c|}{1XX0100X1100XXXXX01XX} &
\multicolumn{1}{c|}{{\scriptsize rd[4]}} &
\multicolumn{1}{c|}{{\scriptsize rd[2:0]}} &
\multicolumn{1}{c|}{0011011}\\
\hline
21& 1& 3& 7 \\
\end{tabular}
 \end{center}
 }
 \vspace{-0.1in}
 \caption[]{ESP.MOVX.R.SAR.BYTES} \end{figure*}


\newpage
\subsubsection{ESP.MOVX.R.XACC.H}\label{instruction:ESPMOVXRXACCH}
\textbf{Syntax}\\
ESP.MOVX.R.XACC.H rd

\textbf{Description}\\
This instruction moves 8-bit data from the special register \lstinline{XACC} to the \lstinline{rd} register.

\textbf{Operation}\\
\begin{lstlisting}
    rd[31 :  0] = {24’b0, XACC[39 : 32]}
\end{lstlisting}

\textbf{Schedule}\\
\begin{longtable}[c]{|l|l|l|}
    \hline
    \rowcolor{lightgray}
    \textbf{Operands} & \textbf{Use} & \textbf{Def} \\\hline
    \endfirsthead
    \hline
    \rowcolor{lightgray}
     \textbf{Operands} & \textbf{Use} & \textbf{Def} \\\hline
    \endhead
    \hline
    \endfoot

  rd & - & 2 \\\hline
    XACC & 2 & -

\end{longtable}

\textbf{Format}\\

\begin{figure*}[!htbp]
 {\setlength{\tabcolsep}{1pt}
 \begin{center}
\begin{tabular}{@{}J@{}I@{}F@{}O}
\instbitrange{31}{11}&
\instbit{10}&
\instbitrange{9}{7}&
\instbitrange{6}{0} \\
\hline
\multicolumn{1}{|c|}{1XX0110X1100XXXXX01XX} &
\multicolumn{1}{c|}{{\scriptsize rd[4]}} &
\multicolumn{1}{c|}{{\scriptsize rd[2:0]}} &
\multicolumn{1}{c|}{0011011}\\
\hline
21& 1& 3& 7 \\
\end{tabular}
 \end{center}
 }
 \vspace{-0.1in}
 \caption[]{ESP.MOVX.R.XACC.H} \end{figure*}


\newpage
\subsubsection{ESP.MOVX.R.XACC.L}\label{instruction:ESPMOVXRXACCL}
\textbf{Syntax}\\
ESP.MOVX.R.XACC.L rd

\textbf{Description}\\
This instruction moves 32 bits of the special register \lstinline{XACC} to the \lstinline{rd} register.

\textbf{Operation}\\
\begin{lstlisting}
    rd[31 :  0] = XACC[31 :  0]
\end{lstlisting}

\textbf{Schedule}\\
\begin{longtable}[c]{|l|l|l|}
    \hline
    \rowcolor{lightgray}
    \textbf{Operands} & \textbf{Use} & \textbf{Def} \\\hline
    \endfirsthead
    \hline
    \rowcolor{lightgray}
     \textbf{Operands} & \textbf{Use} & \textbf{Def} \\\hline
    \endhead
    \hline
    \endfoot

  rd & - & 2 \\\hline
    XACC & 2 & -

\end{longtable}

\textbf{Format}\\

\begin{figure*}[!htbp]
 {\setlength{\tabcolsep}{1pt}
 \begin{center}
\begin{tabular}{@{}J@{}I@{}F@{}O}
\instbitrange{31}{11}&
\instbit{10}&
\instbitrange{9}{7}&
\instbitrange{6}{0} \\
\hline
\multicolumn{1}{|c|}{1XX0010X1100XXXXX01XX} &
\multicolumn{1}{c|}{{\scriptsize rd[4]}} &
\multicolumn{1}{c|}{{\scriptsize rd[2:0]}} &
\multicolumn{1}{c|}{0011011}\\
\hline
21& 1& 3& 7 \\
\end{tabular}
 \end{center}
 }
 \vspace{-0.1in}
 \caption[]{ESP.MOVX.R.XACC.L} \end{figure*}


\newpage
\subsubsection{ESP.MOVX.W.CFG}\label{instruction:ESPMOVXWCFG}
\textbf{Syntax}\\
ESP.MOVX.W.CFG rs1

\textbf{Description}\\
This instruction writes to the special register \lstinline{CFG} with the value of the register \lstinline{rs1}.

\textbf{Operation}\\
\begin{lstlisting}
    CFG[31:0] = rs1[31:0]
\end{lstlisting}

\textbf{Schedule}\\
\begin{longtable}[c]{|l|l|l|}
    \hline
    \rowcolor{lightgray}
    \textbf{Operands} & \textbf{Use} & \textbf{Def} \\\hline
    \endfirsthead
    \hline
    \rowcolor{lightgray}
     \textbf{Operands} & \textbf{Use} & \textbf{Def} \\\hline
    \endhead
    \hline
    \endfoot

  rs1 & 1 & - \\\hline
    CFG & - & 1

\end{longtable}

\textbf{Format}\\

\begin{figure*}[!htbp]
 {\setlength{\tabcolsep}{1pt}
 \begin{center}
\begin{tabular}{@{}K@{}I@{}F@{}E@{}O}
\instbitrange{31}{19}&
\instbit{18}&
\instbitrange{17}{15}&
\instbitrange{14}{7}&
\instbitrange{6}{0} \\
\hline
\multicolumn{1}{|c|}{1XX1X00X1000X} &
\multicolumn{1}{c|}{{\scriptsize rs1[4]}} &
\multicolumn{1}{c|}{{\scriptsize rs1[2:0]}} &
\multicolumn{1}{c|}{01XXXXXX} &
\multicolumn{1}{c|}{0011011}\\
\hline
13& 1& 3& 8& 7 \\
\end{tabular}
 \end{center}
 }
 \vspace{-0.1in}
 \caption[]{ESP.MOVX.W.CFG} \end{figure*}


\newpage
\subsubsection{ESP.MOVX.W.FFT.BIT.WIDTH}\label{instruction:ESPMOVXWFFT_BIT_WIDTH}
\textbf{Syntax}\\
ESP.MOVX.W.FFT.BIT.WIDTH rs1

\textbf{Description}\\
This instruction writes the \lstinline{rs1} to \lstinline{FFT_BIT_WIDTH}.

\textbf{Operation}\\
\begin{lstlisting}
    FFT_BIT_WIDTH[3  :  0] = rs1[3:0]
\end{lstlisting}

\textbf{Schedule}\\
\begin{longtable}[c]{|l|l|l|}
    \hline
    \rowcolor{lightgray}
    \textbf{Operands} & \textbf{Use} & \textbf{Def} \\\hline
    \endfirsthead
    \hline
    \rowcolor{lightgray}
     \textbf{Operands} & \textbf{Use} & \textbf{Def} \\\hline
    \endhead
    \hline
    \endfoot

  rs1 & 1 & - \\\hline
    FFT\_BIT\_WIDTH & - & 1

\end{longtable}

\textbf{Format}\\

\begin{figure*}[!htbp]
 {\setlength{\tabcolsep}{1pt}
 \begin{center}
\begin{tabular}{@{}K@{}I@{}F@{}E@{}O}
\instbitrange{31}{19}&
\instbit{18}&
\instbitrange{17}{15}&
\instbitrange{14}{7}&
\instbitrange{6}{0} \\
\hline
\multicolumn{1}{|c|}{1XX1010X1000X} &
\multicolumn{1}{c|}{{\scriptsize rs1[4]}} &
\multicolumn{1}{c|}{{\scriptsize rs1[2:0]}} &
\multicolumn{1}{c|}{01XXXXXX} &
\multicolumn{1}{c|}{0011011}\\
\hline
13& 1& 3& 8& 7 \\
\end{tabular}
 \end{center}
 }
 \vspace{-0.1in}
 \caption[]{ESP.MOVX.W.FFT.BIT.WIDTH} \end{figure*}


\newpage
\subsubsection{ESP.MOVX.W.PERF}\label{instruction:ESPMOVXWPERF}
\textbf{Syntax}\\
ESP.MOVX.W.PERF rs1

\textbf{Description}\\
This instruction writes to the special register \lstinline{PERF} with the value of the register \lstinline{rs1}.

\textbf{Operation}\\
\begin{lstlisting}
    PERF[31:0] = rs1[31:0]
\end{lstlisting}

\textbf{Schedule}\\
\begin{longtable}[c]{|l|l|l|}
    \hline
    \rowcolor{lightgray}
    \textbf{Operands} & \textbf{Use} & \textbf{Def} \\\hline
    \endfirsthead
    \hline
    \rowcolor{lightgray}
     \textbf{Operands} & \textbf{Use} & \textbf{Def} \\\hline
    \endhead
    \hline
    \endfoot

  rs1 & 1 & - \\\hline
    PERF & - & 1

\end{longtable}

\textbf{Format}\\

\begin{figure*}[!htbp]
 {\setlength{\tabcolsep}{1pt}
 \begin{center}
\begin{tabular}{@{}K@{}I@{}F@{}E@{}O}
\instbitrange{31}{19}&
\instbit{18}&
\instbitrange{17}{15}&
\instbitrange{14}{7}&
\instbitrange{6}{0} \\
\hline
\multicolumn{1}{|c|}{1XX1110X1000X} &
\multicolumn{1}{c|}{{\scriptsize rs1[4]}} &
\multicolumn{1}{c|}{{\scriptsize rs1[2:0]}} &
\multicolumn{1}{c|}{01XXXXXX} &
\multicolumn{1}{c|}{0011011}\\
\hline
13& 1& 3& 8& 7 \\
\end{tabular}
 \end{center}
 }
 \vspace{-0.1in}
 \caption[]{ESP.MOVX.W.PERF} \end{figure*}


\newpage
\subsubsection{ESP.MOVX.W.SAR}\label{instruction:ESPMOVXWSAR}
\textbf{Syntax}\\
ESP.MOVX.W.SAR rs1

\textbf{Description}\\
This instruction moves data from \lstinline{rs1} to the special register \lstinline{SAR}.

\textbf{Operation}\\
\begin{lstlisting}
    SAR[5  :  0] = rs1[5:0]
\end{lstlisting}

\textbf{Schedule}\\
\begin{longtable}[c]{|l|l|l|}
    \hline
    \rowcolor{lightgray}
    \textbf{Operands} & \textbf{Use} & \textbf{Def} \\\hline
    \endfirsthead
    \hline
    \rowcolor{lightgray}
     \textbf{Operands} & \textbf{Use} & \textbf{Def} \\\hline
    \endhead
    \hline
    \endfoot

  rs1 & 1 & - \\\hline
    SAR & - & 1

\end{longtable}

\textbf{Format}\\

\begin{figure*}[!htbp]
 {\setlength{\tabcolsep}{1pt}
 \begin{center}
\begin{tabular}{@{}K@{}I@{}F@{}E@{}O}
\instbitrange{31}{19}&
\instbit{18}&
\instbitrange{17}{15}&
\instbitrange{14}{7}&
\instbitrange{6}{0} \\
\hline
\multicolumn{1}{|c|}{1XX1000X1100X} &
\multicolumn{1}{c|}{{\scriptsize rs1[4]}} &
\multicolumn{1}{c|}{{\scriptsize rs1[2:0]}} &
\multicolumn{1}{c|}{01XXXXXX} &
\multicolumn{1}{c|}{0011011}\\
\hline
13& 1& 3& 8& 7 \\
\end{tabular}
 \end{center}
 }
 \vspace{-0.1in}
 \caption[]{ESP.MOVX.W.SAR} \end{figure*}


\newpage
\subsubsection{ESP.MOVX.W.SAR.BYTES}\label{instruction:ESPMOVXWSAR_BYTES}
\textbf{Syntax}\\
ESP.MOVX.W.SAR.BYTES rs1

\textbf{Description}\\
This instruction writes the \lstinline{rs1} to \lstinline{SAR_BYTES}.

\textbf{Operation}\\
\begin{lstlisting}
    SAR_BYTES[3  :  0] = rs1[3:0]
\end{lstlisting}

\textbf{Schedule}\\
\begin{longtable}[c]{|l|l|l|}
    \hline
    \rowcolor{lightgray}
    \textbf{Operands} & \textbf{Use} & \textbf{Def} \\\hline
    \endfirsthead
    \hline
    \rowcolor{lightgray}
     \textbf{Operands} & \textbf{Use} & \textbf{Def} \\\hline
    \endhead
    \hline
    \endfoot

  rs1 & 1 & - \\\hline
    SAR\_BYTES & - & 1

\end{longtable}

\textbf{Format}\\

\begin{figure*}[!htbp]
 {\setlength{\tabcolsep}{1pt}
 \begin{center}
\begin{tabular}{@{}K@{}I@{}F@{}E@{}O}
\instbitrange{31}{19}&
\instbit{18}&
\instbitrange{17}{15}&
\instbitrange{14}{7}&
\instbitrange{6}{0} \\
\hline
\multicolumn{1}{|c|}{1XX1100X1100X} &
\multicolumn{1}{c|}{{\scriptsize rs1[4]}} &
\multicolumn{1}{c|}{{\scriptsize rs1[2:0]}} &
\multicolumn{1}{c|}{01XXXXXX} &
\multicolumn{1}{c|}{0011011}\\
\hline
13& 1& 3& 8& 7 \\
\end{tabular}
 \end{center}
 }
 \vspace{-0.1in}
 \caption[]{ESP.MOVX.W.SAR.BYTES} \end{figure*}


\newpage
\subsubsection{ESP.MOVX.W.XACC.H}\label{instruction:ESPMOVXWXACCH}
\textbf{Syntax}\\
ESP.MOVX.W.XACC.H rs1

\textbf{Description}\\
This instruction writes to the special register \lstinline{XACC[39:32]} with the value of \lstinline{rs1}.

\textbf{Operation}\\
\begin{lstlisting}
    XACC[39:32] = rs1[7:0]
\end{lstlisting}

\textbf{Schedule}\\
\begin{longtable}[c]{|l|l|l|}
    \hline
    \rowcolor{lightgray}
    \textbf{Operands} & \textbf{Use} & \textbf{Def} \\\hline
    \endfirsthead
    \hline
    \rowcolor{lightgray}
     \textbf{Operands} & \textbf{Use} & \textbf{Def} \\\hline
    \endhead
    \hline
    \endfoot

  rs1 & 1 & - \\\hline
    XACC & - & 2

\end{longtable}

\textbf{Format}\\

\begin{figure*}[!htbp]
 {\setlength{\tabcolsep}{1pt}
 \begin{center}
\begin{tabular}{@{}K@{}I@{}F@{}E@{}O}
\instbitrange{31}{19}&
\instbit{18}&
\instbitrange{17}{15}&
\instbitrange{14}{7}&
\instbitrange{6}{0} \\
\hline
\multicolumn{1}{|c|}{1XX1110X1100X} &
\multicolumn{1}{c|}{{\scriptsize rs1[4]}} &
\multicolumn{1}{c|}{{\scriptsize rs1[2:0]}} &
\multicolumn{1}{c|}{01XXXXXX} &
\multicolumn{1}{c|}{0011011}\\
\hline
13& 1& 3& 8& 7 \\
\end{tabular}
 \end{center}
 }
 \vspace{-0.1in}
 \caption[]{ESP.MOVX.W.XACC.H} \end{figure*}


\newpage
\subsubsection{ESP.MOVX.W.XACC.L}\label{instruction:ESPMOVXWXACCL}
\textbf{Syntax}\\
ESP.MOVX.W.XACC.L rs1

\textbf{Description}\\
This instruction writes to the special register \lstinline{XACC[31:0]} with the value of \lstinline{rs1}.

\textbf{Operation}\\
\begin{lstlisting}
    XACC[31:0] = rs1[31:0]
\end{lstlisting}

\textbf{Schedule}\\
\begin{longtable}[c]{|l|l|l|}
    \hline
    \rowcolor{lightgray}
    \textbf{Operands} & \textbf{Use} & \textbf{Def} \\\hline
    \endfirsthead
    \hline
    \rowcolor{lightgray}
     \textbf{Operands} & \textbf{Use} & \textbf{Def} \\\hline
    \endhead
    \hline
    \endfoot

  rs1 & 1 & - \\\hline
    XACC & - & 2

\end{longtable}

\textbf{Format}\\

\begin{figure*}[!htbp]
 {\setlength{\tabcolsep}{1pt}
 \begin{center}
\begin{tabular}{@{}K@{}I@{}F@{}E@{}O}
\instbitrange{31}{19}&
\instbit{18}&
\instbitrange{17}{15}&
\instbitrange{14}{7}&
\instbitrange{6}{0} \\
\hline
\multicolumn{1}{|c|}{1XX1010X1100X} &
\multicolumn{1}{c|}{{\scriptsize rs1[4]}} &
\multicolumn{1}{c|}{{\scriptsize rs1[2:0]}} &
\multicolumn{1}{c|}{01XXXXXX} &
\multicolumn{1}{c|}{0011011}\\
\hline
13& 1& 3& 8& 7 \\
\end{tabular}
 \end{center}
 }
 \vspace{-0.1in}
 \caption[]{ESP.MOVX.W.XACC.L} \end{figure*}


\newpage
\subsubsection{ESP.VEXT.S16}\label{instruction:ESPVEXTS16}
\textbf{Syntax}\\
ESP.VEXT.S16 qz,qv,qw

\textbf{Description}\\
This instruction signs-extends the load data from a 16-bit segment to a 32-bit segment.

\textbf{Operation}\\
\begin{lstlisting}
    qz[31 :  0] = { {16{qw[ 15]}}, qw[15 :  0] }
    qz[63 : 32] = { {16{qw[ 31]}}, qw[31 : 16] }
    qz[95 : 64] = { {16{qw[ 47]}}, qw[47 : 32] }
    qz[127: 96] = { {16{qw[ 63]}}, qw[63 : 48] }
    qv[31 :  0] = { {16{qw[ 79]}}, qw[79 : 64] }
    qv[63 : 32] = { {16{qw[ 95]}}, qw[95 : 80] }
    qv[95 : 64] = { {16{qw[111]}}, qw[111: 96] }
    qv[127: 96] = { {16{qw[127]}}, qw[127:112] }
\end{lstlisting}

\textbf{Schedule}\\
\begin{longtable}[c]{|l|l|l|}
    \hline
    \rowcolor{lightgray}
    \textbf{Operands} & \textbf{Use} & \textbf{Def} \\\hline
    \endfirsthead
    \hline
    \rowcolor{lightgray}
     \textbf{Operands} & \textbf{Use} & \textbf{Def} \\\hline
    \endhead
    \hline
    \endfoot

  qz & - & 1 \\\hline
  qv & - & 1 \\\hline
    qw & 1 & -

\end{longtable}

\textbf{Format}\\

\begin{figure*}[!htbp]
 {\setlength{\tabcolsep}{1pt}
 \begin{center}
\begin{tabular}{@{}S@{}W@{}I@{}F@{}I@{}T@{}F@{}O}
\instbitrange{31}{26}&
\instbitrange{25}{24}&
\instbit{23}&
\instbitrange{22}{20}&
\instbit{19}&
\instbitrange{18}{10}&
\instbitrange{9}{7}&
\instbitrange{6}{0} \\
\hline
\multicolumn{1}{|c|}{0XX11X} &
\multicolumn{1}{c|}{{\scriptsize qw[1:0]}} &
\multicolumn{1}{c|}{0} &
\multicolumn{1}{c|}{{\scriptsize qv[2:0]}} &
\multicolumn{1}{c|}{{\scriptsize qw[2]}} &
\multicolumn{1}{c|}{1XXX01XXX} &
\multicolumn{1}{c|}{{\scriptsize qz[2:0]}} &
\multicolumn{1}{c|}{0011011}\\
\hline
6& 2& 1& 3& 1& 9& 3& 7 \\
\end{tabular}
 \end{center}
 }
 \vspace{-0.1in}
 \caption[]{ESP.VEXT.S16} \end{figure*}


\newpage
\subsubsection{ESP.VEXT.S8}\label{instruction:ESPVEXTS8}
\textbf{Syntax}\\
ESP.VEXT.S8 qz,qv,qw

\textbf{Description}\\
This instruction signs-extends the load data from an 8-bit segment to a 16-bit segment.

\textbf{Operation}\\
\begin{lstlisting}
    qz[15 :  0] = { {8{qw[  7]}}, qw[7  :  0] }
    qz[31 : 16] = { {8{qw[ 15]}}, qw[15 :  8] }
    qz[47 : 32] = { {8{qw[ 23]}}, qw[23 : 16] }
    qz[63 : 48] = { {8{qw[ 31]}}, qw[31 : 24] }
    qz[79 : 64] = { {8{qw[ 39]}}, qw[39 : 32] }
    qz[95 : 80] = { {8{qw[ 47]}}, qw[47 : 40] }
    qz[111: 96] = { {8{qw[ 55]}}, qw[55 : 48] }
    qz[127:112] = { {8{qw[ 63]}}, qw[63 : 56] }
    qv[15 :  0] = { {8{qw[ 71]}}, qw[71 : 64] }
    qv[31 : 16] = { {8{qw[ 79]}}, qw[79 : 72] }
    qv[47 : 32] = { {8{qw[ 87]}}, qw[87 : 80] }
    qv[63 : 48] = { {8{qw[ 95]}}, qw[95 : 88] }
    qv[79 : 64] = { {8{qw[103]}}, qw[103: 96] }
    qv[95 : 80] = { {8{qw[111]}}, qw[111:104] }
    qv[111: 96] = { {8{qw[119]}}, qw[119:112] }
    qv[127:112] = { {8{qw[127]}}, qw[127:120] }
\end{lstlisting}

\textbf{Schedule}\\
\begin{longtable}[c]{|l|l|l|}
    \hline
    \rowcolor{lightgray}
    \textbf{Operands} & \textbf{Use} & \textbf{Def} \\\hline
    \endfirsthead
    \hline
    \rowcolor{lightgray}
     \textbf{Operands} & \textbf{Use} & \textbf{Def} \\\hline
    \endhead
    \hline
    \endfoot

  qz & - & 1 \\\hline
  qv & - & 1 \\\hline
    qw & 1 & -

\end{longtable}

\textbf{Format}\\

\begin{figure*}[!htbp]
 {\setlength{\tabcolsep}{1pt}
 \begin{center}
\begin{tabular}{@{}S@{}W@{}I@{}F@{}I@{}T@{}F@{}O}
\instbitrange{31}{26}&
\instbitrange{25}{24}&
\instbit{23}&
\instbitrange{22}{20}&
\instbit{19}&
\instbitrange{18}{10}&
\instbitrange{9}{7}&
\instbitrange{6}{0} \\
\hline
\multicolumn{1}{|c|}{0XX01X} &
\multicolumn{1}{c|}{{\scriptsize qw[1:0]}} &
\multicolumn{1}{c|}{0} &
\multicolumn{1}{c|}{{\scriptsize qv[2:0]}} &
\multicolumn{1}{c|}{{\scriptsize qw[2]}} &
\multicolumn{1}{c|}{1XXX01XXX} &
\multicolumn{1}{c|}{{\scriptsize qz[2:0]}} &
\multicolumn{1}{c|}{0011011}\\
\hline
6& 2& 1& 3& 1& 9& 3& 7 \\
\end{tabular}
 \end{center}
 }
 \vspace{-0.1in}
 \caption[]{ESP.VEXT.S8} \end{figure*}


\newpage
\subsubsection{ESP.VEXT.U16}\label{instruction:ESPVEXTU16}
\textbf{Syntax}\\
ESP.VEXT.U16 qz,qv,qw

\textbf{Description}\\
This instruction zero-extends the load data from a 16-bit segment to a 32-bit segment.

\textbf{Operation}\\
\begin{lstlisting}
    qz[31 :  0] = { {16{0}}, qw[15 :  0] }
    qz[63 : 32] = { {16{0}}, qw[31 : 16] }
    qz[95 : 64] = { {16{0}}, qw[47 : 32] }
    qz[127: 96] = { {16{0}}, qw[63 : 48] }
    qv[31 :  0] = { {16{0}}, qw[79 : 64] }
    qv[63 : 32] = { {16{0}}, qw[95 : 80] }
    qv[95 : 64] = { {16{0}}, qw[111: 96] }
    qv[127: 96] = { {16{0}}, qw[127:112] }
\end{lstlisting}

\textbf{Schedule}\\
\begin{longtable}[c]{|l|l|l|}
    \hline
    \rowcolor{lightgray}
    \textbf{Operands} & \textbf{Use} & \textbf{Def} \\\hline
    \endfirsthead
    \hline
    \rowcolor{lightgray}
     \textbf{Operands} & \textbf{Use} & \textbf{Def} \\\hline
    \endhead
    \hline
    \endfoot

  qz & - & 1 \\\hline
  qv & - & 1 \\\hline
    qw & 1 & -

\end{longtable}

\textbf{Format}\\

\begin{figure*}[!htbp]
 {\setlength{\tabcolsep}{1pt}
 \begin{center}
\begin{tabular}{@{}S@{}W@{}I@{}F@{}I@{}T@{}F@{}O}
\instbitrange{31}{26}&
\instbitrange{25}{24}&
\instbit{23}&
\instbitrange{22}{20}&
\instbit{19}&
\instbitrange{18}{10}&
\instbitrange{9}{7}&
\instbitrange{6}{0} \\
\hline
\multicolumn{1}{|c|}{0XX10X} &
\multicolumn{1}{c|}{{\scriptsize qw[1:0]}} &
\multicolumn{1}{c|}{0} &
\multicolumn{1}{c|}{{\scriptsize qv[2:0]}} &
\multicolumn{1}{c|}{{\scriptsize qw[2]}} &
\multicolumn{1}{c|}{1XXX01XXX} &
\multicolumn{1}{c|}{{\scriptsize qz[2:0]}} &
\multicolumn{1}{c|}{0011011}\\
\hline
6& 2& 1& 3& 1& 9& 3& 7 \\
\end{tabular}
 \end{center}
 }
 \vspace{-0.1in}
 \caption[]{ESP.VEXT.U16} \end{figure*}


\newpage
\subsubsection{ESP.VEXT.U8}\label{instruction:ESPVEXTU8}
\textbf{Syntax}\\
ESP.VEXT.U8 qz,qv,qw

\textbf{Description}\\
This instruction zero-extends the load data from an 8-bit segment to a 16-bit segment.

\textbf{Operation}\\
\begin{lstlisting}
    qz[15 :  0] = { {8{0}}, qw[7  :  0] }
    qz[31 : 16] = { {8{0}}, qw[15 :  8] }
    qz[47 : 32] = { {8{0}}, qw[23 : 16] }
    qz[63 : 48] = { {8{0}}, qw[31 : 24] }
    qz[79 : 64] = { {8{0}}, qw[39 : 32] }
    qz[95 : 80] = { {8{0}}, qw[47 : 40] }
    qz[111: 96] = { {8{0}}, qw[55 : 48] }
    qz[127:112] = { {8{0}}, qw[63 : 56] }
    qv[15 :  0] = { {8{0}}, qw[71 : 64] }
    qv[31 : 16] = { {8{0}}, qw[79 : 72] }
    qv[47 : 32] = { {8{0}}, qw[87 : 80] }
    qv[63 : 48] = { {8{0}}, qw[95 : 88] }
    qv[79 : 64] = { {8{0}}, qw[103: 96] }
    qv[95 : 80] = { {8{0}}, qw[111:104] }
    qv[111: 96] = { {8{0}}, qw[119:112] }
    qv[127:112] = { {8{0}}, qw[127:120] }
\end{lstlisting}

\textbf{Schedule}\\
\begin{longtable}[c]{|l|l|l|}
    \hline
    \rowcolor{lightgray}
    \textbf{Operands} & \textbf{Use} & \textbf{Def} \\\hline
    \endfirsthead
    \hline
    \rowcolor{lightgray}
     \textbf{Operands} & \textbf{Use} & \textbf{Def} \\\hline
    \endhead
    \hline
    \endfoot

  qz & - & 1 \\\hline
  qv & - & 1 \\\hline
    qw & 1 & -

\end{longtable}

\textbf{Format}\\

\begin{figure*}[!htbp]
 {\setlength{\tabcolsep}{1pt}
 \begin{center}
\begin{tabular}{@{}S@{}W@{}I@{}F@{}I@{}T@{}F@{}O}
\instbitrange{31}{26}&
\instbitrange{25}{24}&
\instbit{23}&
\instbitrange{22}{20}&
\instbit{19}&
\instbitrange{18}{10}&
\instbitrange{9}{7}&
\instbitrange{6}{0} \\
\hline
\multicolumn{1}{|c|}{0XX00X} &
\multicolumn{1}{c|}{{\scriptsize qw[1:0]}} &
\multicolumn{1}{c|}{0} &
\multicolumn{1}{c|}{{\scriptsize qv[2:0]}} &
\multicolumn{1}{c|}{{\scriptsize qw[2]}} &
\multicolumn{1}{c|}{1XXX01XXX} &
\multicolumn{1}{c|}{{\scriptsize qz[2:0]}} &
\multicolumn{1}{c|}{0011011}\\
\hline
6& 2& 1& 3& 1& 9& 3& 7 \\
\end{tabular}
 \end{center}
 }
 \vspace{-0.1in}
 \caption[]{ESP.VEXT.U8} \end{figure*}


\newpage
\subsubsection{ESP.VUNZIP.16}\label{instruction:ESPVUNZIP16}
\textbf{Syntax}\\
ESP.VUNZIP.16 qx,qy

\textbf{Description}\\
This instruction implements the unzip algorithm on 16-bit vector data.

\textbf{Operation}\\
\begin{lstlisting}
  qx[ 15:  0] = qx[ 15:  0]
  qx[ 31: 16] = qx[ 47: 32]
  qx[ 47: 32] = qx[ 79: 64]
  qx[ 63: 48] = qx[111: 96]
  qx[ 79: 64] = qy[ 15:  0]
  qx[ 95: 80] = qy[ 47: 32]
  qx[111: 96] = qy[ 79: 64]
  qx[127:112] = qy[111: 96]
  qy[ 15:  0] = qx[ 31: 16]
  qy[ 31: 16] = qx[ 63: 48]
  qy[ 47: 32] = qx[ 95: 80]
  qy[ 63: 48] = qx[127:112]
  qy[ 79: 64] = qy[ 31: 16]
  qy[ 95: 80] = qy[ 63: 48]
  qy[111: 96] = qy[ 95: 80]
  qy[127:112] = qy[127:112]
\end{lstlisting}

\textbf{Schedule}\\
\begin{longtable}[c]{|l|l|l|}
    \hline
    \rowcolor{lightgray}
    \textbf{Operands} & \textbf{Use} & \textbf{Def} \\\hline
    \endfirsthead
    \hline
    \rowcolor{lightgray}
     \textbf{Operands} & \textbf{Use} & \textbf{Def} \\\hline
    \endhead
    \hline
    \endfoot

  qx & 1 & 1 \\\hline
  qy & 1 & 1

\end{longtable}

\textbf{Format}\\

\begin{figure*}[!htbp]
 {\setlength{\tabcolsep}{1pt}
 \begin{center}
\begin{tabular}{@{}F@{}F@{}U@{}O}
\instbitrange{31}{29}&
\instbitrange{28}{26}&
\instbitrange{25}{7}&
\instbitrange{6}{0} \\
\hline
\multicolumn{1}{|c|}{{\scriptsize qx[2:0]}} &
\multicolumn{1}{c|}{{\scriptsize qy[2:0]}} &
\multicolumn{1}{c|}{1X1X0XX110001XX0X0X} &
\multicolumn{1}{c|}{0011011}\\
\hline
3& 3& 19& 7 \\
\end{tabular}
 \end{center}
 }
 \vspace{-0.1in}
 \caption[]{ESP.VUNZIP.16} \end{figure*}


\newpage
\subsubsection{ESP.VUNZIP.32}\label{instruction:ESPVUNZIP32}
\textbf{Syntax}\\
ESP.VUNZIP.32 qx,qy

\textbf{Description}\\
This instruction implements the unzip algorithm on 32-bit vector data.

\textbf{Operation}\\
\begin{lstlisting}
  qx[ 31:  0] = qx[ 31:  0]
  qx[ 63: 32] = qx[ 95: 64]
  qx[ 95: 64] = qy[ 31:  0]
  qx[127: 96] = qy[ 95: 64]
  qy[ 31:  0] = qx[ 63: 32]
  qy[ 63: 32] = qx[127: 96]
  qy[ 95: 64] = qy[ 63: 32]
  qy[127: 96] = qy[127: 96]
\end{lstlisting}

\textbf{Schedule}\\
\begin{longtable}[c]{|l|l|l|}
    \hline
    \rowcolor{lightgray}
    \textbf{Operands} & \textbf{Use} & \textbf{Def} \\\hline
    \endfirsthead
    \hline
    \rowcolor{lightgray}
     \textbf{Operands} & \textbf{Use} & \textbf{Def} \\\hline
    \endhead
    \hline
    \endfoot

  qx & 1 & 1 \\\hline
  qy & 1 & 1

\end{longtable}

\textbf{Format}\\

\begin{figure*}[!htbp]
 {\setlength{\tabcolsep}{1pt}
 \begin{center}
\begin{tabular}{@{}F@{}F@{}U@{}O}
\instbitrange{31}{29}&
\instbitrange{28}{26}&
\instbitrange{25}{7}&
\instbitrange{6}{0} \\
\hline
\multicolumn{1}{|c|}{{\scriptsize qx[2:0]}} &
\multicolumn{1}{c|}{{\scriptsize qy[2:0]}} &
\multicolumn{1}{c|}{1X1X0XX1X1001XX0X0X} &
\multicolumn{1}{c|}{0011011}\\
\hline
3& 3& 19& 7 \\
\end{tabular}
 \end{center}
 }
 \vspace{-0.1in}
 \caption[]{ESP.VUNZIP.32} \end{figure*}


\newpage
\subsubsection{ESP.VUNZIP.8}\label{instruction:ESPVUNZIP8}
\textbf{Syntax}\\
ESP.VUNZIP.8 qx,qy

\textbf{Description}\\
This instruction implements the unzip algorithm on 8-bit vector data.

\textbf{Operation}\\
\begin{lstlisting}
  qx[  7:  0] = qx[  7:  0]
  qx[ 15:  8] = qx[ 23: 16]
  qx[ 23: 16] = qx[ 39: 32]
  qx[ 31: 24] = qx[ 55: 48]
  qx[ 39: 32] = qx[ 71: 64]
  qx[ 47: 40] = qx[ 87: 80]
  qx[ 55: 48] = qx[103: 96]
  qx[ 63: 56] = qx[119:112]
  qx[ 71: 64] = qy[  7:  0]
  qx[ 79: 72] = qy[ 23: 16]
  qx[ 87: 80] = qy[ 39: 32]
  qx[ 95: 88] = qy[ 55: 48]
  qx[103: 96] = qy[ 71: 64]
  qx[111:104] = qy[ 87: 80]
  qx[119:112] = qy[103: 96]
  qx[127:120] = qy[119:112]
  qy[  7:  0] = qx[ 15:  8]
  qy[ 15:  8] = qx[ 31: 24]
  qy[ 23: 16] = qx[ 47: 40]
  qy[ 31: 24] = qx[ 63: 56]
  qy[ 39: 32] = qx[ 79: 72]
  qy[ 47: 40] = qx[ 95: 88]
  qy[ 55: 48] = qx[111:104]
  qy[ 63: 56] = qx[127:120]
  qy[ 71: 64] = qy[ 15:  8]
  qy[ 79: 72] = qy[ 31: 24]
  qy[ 87: 80] = qy[ 47: 40]
  qy[ 95: 88] = qy[ 63: 56]
  qy[103: 96] = qy[ 79: 72]
  qy[111:104] = qy[ 95: 88]
  qy[119:112] = qy[111:104]
  qy[127:120] = qy[127:120]
\end{lstlisting}

\textbf{Schedule}\\
\begin{longtable}[c]{|l|l|l|}
    \hline
    \rowcolor{lightgray}
    \textbf{Operands} & \textbf{Use} & \textbf{Def} \\\hline
    \endfirsthead
    \hline
    \rowcolor{lightgray}
     \textbf{Operands} & \textbf{Use} & \textbf{Def} \\\hline
    \endhead
    \hline
    \endfoot

  qx & 1 & 1 \\\hline
  qy & 1 & 1

\end{longtable}

\textbf{Format}\\

\begin{figure*}[!htbp]
 {\setlength{\tabcolsep}{1pt}
 \begin{center}
\begin{tabular}{@{}F@{}F@{}U@{}O}
\instbitrange{31}{29}&
\instbitrange{28}{26}&
\instbitrange{25}{7}&
\instbitrange{6}{0} \\
\hline
\multicolumn{1}{|c|}{{\scriptsize qx[2:0]}} &
\multicolumn{1}{c|}{{\scriptsize qy[2:0]}} &
\multicolumn{1}{c|}{1X1X0XX100001XX0X0X} &
\multicolumn{1}{c|}{0011011}\\
\hline
3& 3& 19& 7 \\
\end{tabular}
 \end{center}
 }
 \vspace{-0.1in}
 \caption[]{ESP.VUNZIP.8} \end{figure*}


\newpage
\subsubsection{ESP.VUNZIPT.16}\label{instruction:ESPVUNZIPT16}
\textbf{Syntax}\\
ESP.VUNZIPT.16 qx,qy,qw

\textbf{Description}\\
This instruction exchanges the values of \lstinline{qx}, \lstinline{qy}, and \lstinline{qw}; it can be used to convert RGB to YUV.

\textbf{Operation}\\
\begin{lstlisting}
  temp0[ 15:  0] = qx[ 15:  0]
  temp1[ 15:  0] = qx[ 31: 16]
  temp2[ 15:  0] = qx[ 47: 32]
  temp0[ 31: 16] = qx[ 63: 48]
  temp1[ 31: 16] = qx[ 79: 64]
  temp2[ 31: 16] = qx[ 95: 80]
  temp0[ 47: 32] = qx[111: 96]
  temp1[ 47: 32] = qx[127:112]

  temp2[ 47: 32] = qy[ 15:  0]
  temp0[ 63: 48] = qy[ 31: 16]
  temp1[ 63: 48] = qy[ 47: 32]
  temp2[ 63: 48] = qy[ 63: 48]
  temp0[ 79: 64] = qy[ 79: 64]
  temp1[ 79: 64] = qy[ 95: 80]
  temp2[ 79: 64] = qy[111: 96]
  temp0[ 95: 80] = qy[127:112]

  temp1[ 95: 80] = qw[ 15:  0]
  temp2[ 95: 80] = qw[ 31: 16]
  temp0[111: 96] = qw[ 47: 32]
  temp1[111: 96] = qw[ 63: 48]
  temp2[111: 96] = qw[ 79: 64]
  temp0[127:112] = qw[ 95: 80]
  temp1[127:112] = qw[111: 96]
  temp2[127:112] = qw[127:112]

  qx[127:0] = temp0[127:0]
  qy[127:0] = temp1[127:0]
  qw[127:0] = temp2[127:0]

\end{lstlisting}

\textbf{Schedule}\\
\begin{longtable}[c]{|l|l|l|}
    \hline
    \rowcolor{lightgray}
    \textbf{Operands} & \textbf{Use} & \textbf{Def} \\\hline
    \endfirsthead
    \hline
    \rowcolor{lightgray}
     \textbf{Operands} & \textbf{Use} & \textbf{Def} \\\hline
    \endhead
    \hline
    \endfoot

  qx & 1 & 1 \\\hline
  qy & 1 & 1 \\\hline
  qw & 1 & 1

\end{longtable}

\textbf{Format}\\

\begin{figure*}[!htbp]
 {\setlength{\tabcolsep}{1pt}
 \begin{center}
\begin{tabular}{@{}F@{}F@{}W@{}Y@{}I@{}M@{}O}
\instbitrange{31}{29}&
\instbitrange{28}{26}&
\instbitrange{25}{24}&
\instbitrange{23}{20}&
\instbit{19}&
\instbitrange{18}{7}&
\instbitrange{6}{0} \\
\hline
\multicolumn{1}{|c|}{{\scriptsize qx[2:0]}} &
\multicolumn{1}{c|}{{\scriptsize qy[2:0]}} &
\multicolumn{1}{c|}{{\scriptsize qw[1:0]}} &
\multicolumn{1}{c|}{01XX} &
\multicolumn{1}{c|}{{\scriptsize qw[2]}} &
\multicolumn{1}{c|}{0XXX01XX111X} &
\multicolumn{1}{c|}{0011011}\\
\hline
3& 3& 2& 4& 1& 12& 7 \\
\end{tabular}
 \end{center}
 }
 \vspace{-0.1in}
 \caption[]{ESP.VUNZIPT.16} \end{figure*}


\newpage
\subsubsection{ESP.VUNZIPT.8}\label{instruction:ESPVUNZIPT8}
\textbf{Syntax}\\
ESP.VUNZIPT.8 qx,qy,qw

\textbf{Description}\\
This instruction exchanges the values of \lstinline{qx}, \lstinline{qy}, and \lstinline{qw}; it can be used to convert RGB to YUV.

\textbf{Operation}\\
\begin{lstlisting}
  temp0[  7:  0] = qx[  7:  0]
  temp1[  7:  0] = qx[ 15:  8]
  temp2[  7:  0] = qx[ 23: 16]
  temp0[ 15:  8] = qx[ 31: 24]
  temp1[ 15:  8] = qx[ 39: 32]
  temp2[ 15:  8] = qx[ 47: 40]
  temp0[ 23: 16] = qx[ 55: 48]
  temp1[ 23: 16] = qx[ 63: 56]
  temp2[ 23: 16] = qx[ 71: 64]
  temp0[ 31: 24] = qx[ 79: 72]
  temp1[ 31: 24] = qx[ 87: 80]
  temp2[ 31: 24] = qx[ 95: 88]
  temp0[ 39: 32] = qx[103: 96]
  temp1[ 39: 32] = qx[111:104]
  temp2[ 39: 32] = qx[119:112]
  temp0[ 47: 40] = qx[127:120]

  temp1[ 47: 40] = qy[  7:  0]
  temp2[ 47: 40] = qy[ 15:  8]
  temp0[ 55: 48] = qy[ 23: 16]
  temp1[ 55: 48] = qy[ 31: 24]
  temp2[ 55: 48] = qy[ 39: 32]
  temp0[ 63: 56] = qy[ 47: 40]
  temp1[ 63: 56] = qy[ 55: 48]
  temp2[ 63: 56] = qy[ 63: 56]
  temp0[ 71: 64] = qy[ 71: 64]
  temp1[ 71: 64] = qy[ 79: 72]
  temp2[ 71: 64] = qy[ 87: 80]
  temp0[ 79: 72] = qy[ 95: 88]
  temp1[ 79: 72] = qy[103: 96]
  temp2[ 79: 72] = qy[111:104]
  temp0[ 87: 80] = qy[119:112]
  temp1[ 87: 80] = qy[127:120]

  temp2[ 87: 80] = qw[  7:  0]
  temp0[ 95: 88] = qw[ 15:  8]
  temp1[ 95: 88] = qw[ 23: 16]
  temp2[ 95: 88] = qw[ 31: 24]
  temp0[103: 96] = qw[ 39: 32]
  temp1[103: 96] = qw[ 47: 40]
  temp2[103: 96] = qw[ 55: 48]
  temp0[111:104] = qw[ 63: 56]
  temp1[111:104] = qw[ 71: 64]
  temp2[111:104] = qw[ 79: 72]
  temp0[119:112] = qw[ 87: 80]
  temp1[119:112] = qw[ 95: 88]
  temp2[119:112] = qw[103: 96]
  temp0[127:120] = qw[111:104]
  temp1[127:120] = qw[119:112]
  temp2[127:120] = qw[127:120]

  qx[127:0] = temp0[127:0]
  qy[127:0] = temp1[127:0]
  qw[127:0] = temp2[127:0]

\end{lstlisting}

\textbf{Schedule}\\
\begin{longtable}[c]{|l|l|l|}
    \hline
    \rowcolor{lightgray}
    \textbf{Operands} & \textbf{Use} & \textbf{Def} \\\hline
    \endfirsthead
    \hline
    \rowcolor{lightgray}
     \textbf{Operands} & \textbf{Use} & \textbf{Def} \\\hline
    \endhead
    \hline
    \endfoot

  qx & 1 & 1 \\\hline
  qy & 1 & 1 \\\hline
  qw & 1 & 1

\end{longtable}

\textbf{Format}\\

\begin{figure*}[!htbp]
 {\setlength{\tabcolsep}{1pt}
 \begin{center}
\begin{tabular}{@{}F@{}F@{}W@{}Y@{}I@{}M@{}O}
\instbitrange{31}{29}&
\instbitrange{28}{26}&
\instbitrange{25}{24}&
\instbitrange{23}{20}&
\instbit{19}&
\instbitrange{18}{7}&
\instbitrange{6}{0} \\
\hline
\multicolumn{1}{|c|}{{\scriptsize qx[2:0]}} &
\multicolumn{1}{c|}{{\scriptsize qy[2:0]}} &
\multicolumn{1}{c|}{{\scriptsize qw[1:0]}} &
\multicolumn{1}{c|}{01XX} &
\multicolumn{1}{c|}{{\scriptsize qw[2]}} &
\multicolumn{1}{c|}{0XXX01XX101X} &
\multicolumn{1}{c|}{0011011}\\
\hline
3& 3& 2& 4& 1& 12& 7 \\
\end{tabular}
 \end{center}
 }
 \vspace{-0.1in}
 \caption[]{ESP.VUNZIPT.8} \end{figure*}


\newpage
\subsubsection{ESP.VZIP.16}\label{instruction:ESPVZIP16}
\textbf{Syntax}\\
ESP.VZIP.16 qx,qy

\textbf{Description}\\
This instruction interleaves the values of \lstinline{qx} and \lstinline{qy}.

\textbf{Operation}\\
\begin{lstlisting}
  qx[ 15:  0] = qx[ 15:  0]
  qx[ 31: 16] = qy[ 15:  0]
  qx[ 47: 32] = qx[ 31: 16]
  qx[ 63: 48] = qy[ 31: 16]
  qx[ 79: 64] = qx[ 47: 32]
  qx[ 95: 80] = qy[ 47: 32]
  qx[111: 96] = qx[ 63: 48]
  qx[127:112] = qy[ 63: 48]
  qy[ 15:  0] = qx[ 79: 64]
  qy[ 31: 16] = qy[ 79: 64]
  qy[ 47: 32] = qx[ 95: 80]
  qy[ 63: 48] = qy[ 95: 80]
  qy[ 79: 64] = qx[111: 96]
  qy[ 95: 80] = qy[111: 96]
  qy[111: 96] = qx[127:112]
  qy[127:112] = qy[127:112]
\end{lstlisting}

\textbf{Schedule}\\
\begin{longtable}[c]{|l|l|l|}
    \hline
    \rowcolor{lightgray}
    \textbf{Operands} & \textbf{Use} & \textbf{Def} \\\hline
    \endfirsthead
    \hline
    \rowcolor{lightgray}
     \textbf{Operands} & \textbf{Use} & \textbf{Def} \\\hline
    \endhead
    \hline
    \endfoot

  qx & 1 & 1 \\\hline
  qy & 1 & 1

\end{longtable}

\textbf{Format}\\

\begin{figure*}[!htbp]
 {\setlength{\tabcolsep}{1pt}
 \begin{center}
\begin{tabular}{@{}F@{}F@{}U@{}O}
\instbitrange{31}{29}&
\instbitrange{28}{26}&
\instbitrange{25}{7}&
\instbitrange{6}{0} \\
\hline
\multicolumn{1}{|c|}{{\scriptsize qx[2:0]}} &
\multicolumn{1}{c|}{{\scriptsize qy[2:0]}} &
\multicolumn{1}{c|}{1X1X0XX010001XX0X0X} &
\multicolumn{1}{c|}{0011011}\\
\hline
3& 3& 19& 7 \\
\end{tabular}
 \end{center}
 }
 \vspace{-0.1in}
 \caption[]{ESP.VZIP.16} \end{figure*}


\newpage
\subsubsection{ESP.VZIP.32}\label{instruction:ESPVZIP32}
\textbf{Syntax}\\
ESP.VZIP.32 qx,qy

\textbf{Description}\\
This instruction interleaves the values of \lstinline{qx} and \lstinline{qy}.

\textbf{Operation}\\
\begin{lstlisting}
  qx[ 31:  0] = qx[ 31:  0]
  qx[ 63: 32] = qy[ 31:  0]
  qx[ 95: 64] = qx[ 63: 32]
  qx[127: 96] = qy[ 63: 32]
  qy[ 31:  0] = qx[ 95: 64]
  qy[ 63: 32] = qy[ 95: 64]
  qy[ 95: 64] = qx[127: 96]
  qy[127: 96] = qy[127: 96]
\end{lstlisting}

\textbf{Schedule}\\
\begin{longtable}[c]{|l|l|l|}
    \hline
    \rowcolor{lightgray}
    \textbf{Operands} & \textbf{Use} & \textbf{Def} \\\hline
    \endfirsthead
    \hline
    \rowcolor{lightgray}
     \textbf{Operands} & \textbf{Use} & \textbf{Def} \\\hline
    \endhead
    \hline
    \endfoot

  qx & 1 & 1 \\\hline
  qy & 1 & 1

\end{longtable}

\textbf{Format}\\

\begin{figure*}[!htbp]
 {\setlength{\tabcolsep}{1pt}
 \begin{center}
\begin{tabular}{@{}F@{}F@{}U@{}O}
\instbitrange{31}{29}&
\instbitrange{28}{26}&
\instbitrange{25}{7}&
\instbitrange{6}{0} \\
\hline
\multicolumn{1}{|c|}{{\scriptsize qx[2:0]}} &
\multicolumn{1}{c|}{{\scriptsize qy[2:0]}} &
\multicolumn{1}{c|}{1X1X0XX0X1001XX0X0X} &
\multicolumn{1}{c|}{0011011}\\
\hline
3& 3& 19& 7 \\
\end{tabular}
 \end{center}
 }
 \vspace{-0.1in}
 \caption[]{ESP.VZIP.32} \end{figure*}


\newpage
\subsubsection{ESP.VZIP.8}\label{instruction:ESPVZIP8}
\textbf{Syntax}\\
ESP.VZIP.8 qx,qy

\textbf{Description}\\
This instruction interleaves the values of \lstinline{qx} and \lstinline{qy}.

\textbf{Operation}\\
\begin{lstlisting}
  qx[  7:  0] = qx[  7:  0]
  qx[ 15:  8] = qy[  7:  0]
  qx[ 23: 16] = qx[ 15:  8]
  qx[ 31: 24] = qy[ 15:  8]
  qx[ 39: 32] = qx[ 23: 16]
  qx[ 47: 40] = qy[ 23: 16]
  qx[ 55: 48] = qx[ 31: 24]
  qx[ 63: 56] = qy[ 31: 24]
  qx[ 71: 64] = qx[ 39: 32]
  qx[ 79: 72] = qy[ 39: 32]
  qx[ 87: 80] = qx[ 47: 40]
  qx[ 95: 88] = qy[ 47: 40]
  qx[103: 96] = qx[ 55: 48]
  qx[111:104] = qy[ 55: 48]
  qx[119:112] = qx[ 63: 56]
  qx[127:120] = qy[ 63: 56]
  qy[  7:  0] = qx[ 71: 64]
  qy[ 15:  8] = qy[ 71: 64]
  qy[ 23: 16] = qx[ 79: 72]
  qy[ 31: 24] = qy[ 79: 72]
  qy[ 39: 32] = qx[ 87: 80]
  qy[ 47: 40] = qy[ 87: 80]
  qy[ 55: 48] = qx[ 95: 88]
  qy[ 63: 56] = qy[ 95: 88]
  qy[ 71: 64] = qx[103: 96]
  qy[ 79: 72] = qy[103: 96]
  qy[ 87: 80] = qx[111:104]
  qy[ 95: 88] = qy[111:104]
  qy[103: 96] = qx[119:112]
  qy[111:104] = qy[119:112]
  qy[119:112] = qx[127:120]
  qy[127:120] = qy[127:120]
\end{lstlisting}

\textbf{Schedule}\\
\begin{longtable}[c]{|l|l|l|}
    \hline
    \rowcolor{lightgray}
    \textbf{Operands} & \textbf{Use} & \textbf{Def} \\\hline
    \endfirsthead
    \hline
    \rowcolor{lightgray}
     \textbf{Operands} & \textbf{Use} & \textbf{Def} \\\hline
    \endhead
    \hline
    \endfoot

  qx & 1 & 1 \\\hline
  qy & 1 & 1

\end{longtable}

\textbf{Format}\\

\begin{figure*}[!htbp]
 {\setlength{\tabcolsep}{1pt}
 \begin{center}
\begin{tabular}{@{}F@{}F@{}U@{}O}
\instbitrange{31}{29}&
\instbitrange{28}{26}&
\instbitrange{25}{7}&
\instbitrange{6}{0} \\
\hline
\multicolumn{1}{|c|}{{\scriptsize qx[2:0]}} &
\multicolumn{1}{c|}{{\scriptsize qy[2:0]}} &
\multicolumn{1}{c|}{1X1X0XX000001XX0X0X} &
\multicolumn{1}{c|}{0011011}\\
\hline
3& 3& 19& 7 \\
\end{tabular}
 \end{center}
 }
 \vspace{-0.1in}
 \caption[]{ESP.VZIP.8} \end{figure*}


\newpage
\subsubsection{ESP.VZIPT.16}\label{instruction:ESPVZIPT16}
\textbf{Syntax}\\
ESP.VZIPT.16 qx,qy,qw

\textbf{Description}\\
This instruction interleaves the values of \lstinline{qx}, \lstinline{qy}, and \lstinline{qw}. It can be used for image color transmission, for example, converting YUV to RGB.

\textbf{Operation}\\
\begin{lstlisting}
  temp0[ 15:  0] = qx[ 15:  0]
  temp0[ 31: 16] = qy[ 15:  0]
  temp0[ 47: 32] = qw[ 15:  0]
  temp0[ 63: 48] = qx[ 31: 16]
  temp0[ 79: 64] = qy[ 31: 16]
  temp0[ 95: 80] = qw[ 31: 16]
  temp0[111: 96] = qx[ 47: 32]
  temp0[127:112] = qy[ 47: 32]

  temp1[ 15:  0] = qw[ 47: 32]
  temp1[ 31: 16] = qx[ 63: 48]
  temp1[ 47: 32] = qy[ 63: 48]
  temp1[ 63: 48] = qw[ 63: 48]
  temp1[ 79: 64] = qx[ 79: 64]
  temp1[ 95: 80] = qy[ 79: 64]
  temp1[111: 96] = qw[ 79: 64]
  temp1[127:112] = qx[ 95: 80]

  temp2[ 15:  0] = qy[ 95: 80]
  temp2[ 31: 16] = qw[ 95: 80]
  temp2[ 47: 32] = qx[111: 96]
  temp2[ 63: 48] = qy[111: 96]
  temp2[ 79: 64] = qw[111: 96]
  temp2[ 95: 80] = qx[127:112]
  temp2[111: 96] = qy[127:112]
  temp2[127:112] = qw[127:112]

  qx[127:0] = temp0[127:0]
  qy[127:0] = temp1[127:0]
  qw[127:0] = temp2[127:0]

\end{lstlisting}

\textbf{Schedule}\\
\begin{longtable}[c]{|l|l|l|}
    \hline
    \rowcolor{lightgray}
    \textbf{Operands} & \textbf{Use} & \textbf{Def} \\\hline
    \endfirsthead
    \hline
    \rowcolor{lightgray}
     \textbf{Operands} & \textbf{Use} & \textbf{Def} \\\hline
    \endhead
    \hline
    \endfoot

  qx & 1 & 1 \\\hline
  qy & 1 & 1 \\\hline
  qw & 1 & 1

\end{longtable}

\textbf{Format}\\

\begin{figure*}[!htbp]
 {\setlength{\tabcolsep}{1pt}
 \begin{center}
\begin{tabular}{@{}F@{}F@{}W@{}Y@{}I@{}M@{}O}
\instbitrange{31}{29}&
\instbitrange{28}{26}&
\instbitrange{25}{24}&
\instbitrange{23}{20}&
\instbit{19}&
\instbitrange{18}{7}&
\instbitrange{6}{0} \\
\hline
\multicolumn{1}{|c|}{{\scriptsize qx[2:0]}} &
\multicolumn{1}{c|}{{\scriptsize qy[2:0]}} &
\multicolumn{1}{c|}{{\scriptsize qw[1:0]}} &
\multicolumn{1}{c|}{01XX} &
\multicolumn{1}{c|}{{\scriptsize qw[2]}} &
\multicolumn{1}{c|}{0XXX01XX011X} &
\multicolumn{1}{c|}{0011011}\\
\hline
3& 3& 2& 4& 1& 12& 7 \\
\end{tabular}
 \end{center}
 }
 \vspace{-0.1in}
 \caption[]{ESP.VZIPT.16} \end{figure*}


\newpage
\subsubsection{ESP.VZIPT.8}\label{instruction:ESPVZIPT8}
\textbf{Syntax}\\
ESP.VZIPT.8 qx,qy,qw

\textbf{Description}\\
This instruction interleaves the values of \lstinline{qx}, \lstinline{qy}, and \lstinline{qw}. It can be used for image color transmission, for example, converting YUV to RGB.

\textbf{Operation}\\
\begin{lstlisting}
  temp0[  7:  0] = qx[  7:  0]
  temp0[ 15:  8] = qy[  7:  0]
  temp0[ 23: 16] = qw[  7:  0]
  temp0[ 31: 24] = qx[ 15:  8]
  temp0[ 39: 32] = qy[ 15:  8]
  temp0[ 47: 40] = qw[ 15:  8]
  temp0[ 55: 48] = qx[ 23: 16]
  temp0[ 63: 56] = qy[ 23: 16]
  temp0[ 71: 64] = qw[ 23: 16]
  temp0[ 79: 72] = qx[ 31: 24]
  temp0[ 87: 80] = qy[ 31: 24]
  temp0[ 95: 88] = qw[ 31: 24]
  temp0[103: 96] = qx[ 39: 32]
  temp0[111:104] = qy[ 39: 32]
  temp0[119:112] = qw[ 39: 32]
  temp0[127:120] = qx[ 47: 40]

  temp1[ 15:  0] = qy[ 47: 40]
  temp1[ 31: 16] = qw[ 47: 40]
  temp1[ 47: 32] = qx[ 55: 48]
  temp1[ 63: 48] = qy[ 55: 48]
  temp1[ 79: 64] = qw[ 55: 48]
  temp1[ 95: 80] = qx[ 63: 56]
  temp1[111: 96] = qy[ 63: 56]
  temp1[127:112] = qw[ 63: 56]
  temp1[ 15:  0] = qx[ 71: 64]
  temp1[ 31: 16] = qy[ 71: 64]
  temp1[ 47: 32] = qw[ 71: 64]
  temp1[ 63: 48] = qx[ 79: 72]
  temp1[ 79: 64] = qy[ 79: 72]
  temp1[ 95: 80] = qw[ 79: 72]
  temp1[111: 96] = qx[ 87: 80]
  temp1[127:112] = qy[ 87: 80]

  temp2[ 15:  0] = qw[ 87: 80]
  temp2[ 31: 16] = qx[ 95: 88]
  temp2[ 47: 32] = qy[ 95: 88]
  temp2[ 63: 48] = qw[ 95: 88]
  temp2[ 79: 64] = qx[103: 96]
  temp2[ 95: 80] = qy[103: 96]
  temp2[111: 96] = qw[103: 96]
  temp2[127:112] = qx[111:104]
  temp2[ 15:  0] = qy[111:104]
  temp2[ 31: 16] = qw[111:104]
  temp2[ 47: 32] = qx[119:112]
  temp2[ 63: 48] = qy[119:112]
  temp2[ 79: 64] = qw[119:112]
  temp2[ 95: 80] = qx[127:120]
  temp2[111: 96] = qy[127:120]
  temp2[127:112] = qw[127:120]

  qx[127:0] = temp0[127:0]
  qy[127:0] = temp1[127:0]
  qw[127:0] = temp2[127:0]
\end{lstlisting}

\textbf{Schedule}\\
\begin{longtable}[c]{|l|l|l|}
    \hline
    \rowcolor{lightgray}
    \textbf{Operands} & \textbf{Use} & \textbf{Def} \\\hline
    \endfirsthead
    \hline
    \rowcolor{lightgray}
     \textbf{Operands} & \textbf{Use} & \textbf{Def} \\\hline
    \endhead
    \hline
    \endfoot

  qx & 1 & 1 \\\hline
  qy & 1 & 1 \\\hline
  qw & 1 & 1

\end{longtable}

\textbf{Format}\\

\begin{figure*}[!htbp]
 {\setlength{\tabcolsep}{1pt}
 \begin{center}
\begin{tabular}{@{}F@{}F@{}W@{}Y@{}I@{}M@{}O}
\instbitrange{31}{29}&
\instbitrange{28}{26}&
\instbitrange{25}{24}&
\instbitrange{23}{20}&
\instbit{19}&
\instbitrange{18}{7}&
\instbitrange{6}{0} \\
\hline
\multicolumn{1}{|c|}{{\scriptsize qx[2:0]}} &
\multicolumn{1}{c|}{{\scriptsize qy[2:0]}} &
\multicolumn{1}{c|}{{\scriptsize qw[1:0]}} &
\multicolumn{1}{c|}{01XX} &
\multicolumn{1}{c|}{{\scriptsize qw[2]}} &
\multicolumn{1}{c|}{0XXX01XX001X} &
\multicolumn{1}{c|}{0011011}\\
\hline
3& 3& 2& 4& 1& 12& 7 \\
\end{tabular}
 \end{center}
 }
 \vspace{-0.1in}
 \caption[]{ESP.VZIPT.8} \end{figure*}


\newpage
\subsubsection{ESP.ZERO.Q}\label{instruction:ESPZEROQ}
\textbf{Syntax}\\
ESP.ZERO.Q qz

\textbf{Description}\\
This instruction sets the register \lstinline{qz} to zero.

\textbf{Operation}\\
\begin{lstlisting}
    qz[127:0] = 0
\end{lstlisting}

\textbf{Schedule}\\
\begin{longtable}[c]{|l|l|l|}
    \hline
    \rowcolor{lightgray}
    \textbf{Operands} & \textbf{Use} & \textbf{Def} \\\hline
    \endfirsthead
    \hline
    \rowcolor{lightgray}
     \textbf{Operands} & \textbf{Use} & \textbf{Def} \\\hline
    \endhead
    \hline
    \endfoot

  qz & - & 1

\end{longtable}

\textbf{Format}\\

\begin{figure*}[!htbp]
 {\setlength{\tabcolsep}{1pt}
 \begin{center}
\begin{tabular}{@{}J@{}F@{}O}
\instbitrange{31}{10}&
\instbitrange{9}{7}&
\instbitrange{6}{0} \\
\hline
\multicolumn{1}{|c|}{0XXXXXXXXXXXXXXXX001XX} &
\multicolumn{1}{c|}{{\scriptsize qz[2:0]}} &
\multicolumn{1}{c|}{0011011}\\
\hline
22& 3& 7 \\
\end{tabular}
 \end{center}
 }
 \vspace{-0.1in}
 \caption[]{ESP.ZERO.Q} \end{figure*}


\newpage
\subsubsection{ESP.ZERO.QACC}\label{instruction:ESPZEROQACC}
\textbf{Syntax}\\
ESP.ZERO.QACC 

\textbf{Description}\\
This instruction sets \lstinline{QACC_L} and \lstinline{QACC_H} to zero.

\textbf{Operation}\\
\begin{lstlisting}
    QACC_L = 0
    QACC_H = 0
\end{lstlisting}

\textbf{Schedule}\\
\begin{longtable}[c]{|l|l|l|}
    \hline
    \rowcolor{lightgray}
    \textbf{Operands} & \textbf{Use} & \textbf{Def} \\\hline
    \endfirsthead
    \hline
    \rowcolor{lightgray}
     \textbf{Operands} & \textbf{Use} & \textbf{Def} \\\hline
    \endhead
    \hline
    \endfoot

  QACC\_H & - & 2 \\\hline
  QACC\_L & - & 2

\end{longtable}

\textbf{Format}\\

\begin{figure*}[!htbp]
 {\setlength{\tabcolsep}{1pt}
 \begin{center}
\begin{tabular}{@{}J@{}O}
\instbitrange{31}{7}&
\instbitrange{6}{0} \\
\hline
\multicolumn{1}{|c|}{0XXXXXXXXXXXXXXXX000XX10X} &
\multicolumn{1}{c|}{0011011}\\
\hline
25& 7 \\
\end{tabular}
 \end{center}
 }
 \vspace{-0.1in}
 \caption[]{ESP.ZERO.QACC} \end{figure*}


\newpage
\subsubsection{ESP.ZERO.XACC}\label{instruction:ESPZEROXACC}
\textbf{Syntax}\\
ESP.ZERO.XACC 

\textbf{Description}\\
This instruction sets \lstinline{XACC} to zero.

\textbf{Operation}\\
\begin{lstlisting}
    XACC = 0
\end{lstlisting}

\textbf{Schedule}\\
\begin{longtable}[c]{|l|l|l|}
    \hline
    \rowcolor{lightgray}
    \textbf{Operands} & \textbf{Use} & \textbf{Def} \\\hline
    \endfirsthead
    \hline
    \rowcolor{lightgray}
     \textbf{Operands} & \textbf{Use} & \textbf{Def} \\\hline
    \endhead
    \hline
    \endfoot

  XACC & - & 2

\end{longtable}

\textbf{Format}\\

\begin{figure*}[!htbp]
 {\setlength{\tabcolsep}{1pt}
 \begin{center}
\begin{tabular}{@{}J@{}O}
\instbitrange{31}{7}&
\instbitrange{6}{0} \\
\hline
\multicolumn{1}{|c|}{0XXXXXXXXXXXXXXXX000XX00X} &
\multicolumn{1}{c|}{0011011}\\
\hline
25& 7 \\
\end{tabular}
 \end{center}
 }
 \vspace{-0.1in}
 \caption[]{ESP.ZERO.XACC} \end{figure*}


\newpage
\subsubsection{ESP.FFT.AMS.S16.LD.INCP}\label{instruction:ESPFFTAMSS16LDINCP}
\textbf{Syntax}\\
ESP.FFT.AMS.S16.LD.INCP qu,rs1,qz,qv,qx,qw,qy,sel2,sat

\textbf{Description}\\
It is a dedicated FFT instruction to perform addition, subtraction, multiplication, addition and subtraction, and shift operations on 16-bit data segments.\\
During the operation, the 16-byte data is loaded from the memory to register \lstinline{qu}. After the access, the value in register \lstinline{rs1} is incremented by 16.\\
\textbf{Operation}

\begin{lstlisting}
  temp0[16:0] = qx[ 47: 32] + qw[ 47: 32]
  temp1[16:0] = qx[ 63: 48] - qw[ 63: 48]

if sel2==0:
  temp2[33:0] = ((qx[ 47: 32] - qw[ 47: 32]) * qy[ 47: 32] - (qx[ 63: 48] + qw[ 63: 48]) * qy[ 63: 48]) >> SAR
  temp3[33:0] = ((qx[ 47: 32] - qw[ 47: 32]) * qy[ 63: 48] + (qx[ 63: 48] + qw[ 63: 48]) * qy[ 47: 32]) >> SAR
if sel2==1:
  temp2[33:0] = ((qx[ 63: 48] + qw[ 63: 48]) * qy[ 63: 48] + (qx[ 47: 32] - qw[ 47: 32]) * qy[ 47: 32]) >> SAR
  temp3[33:0] = ((qx[ 63: 48] + qw[ 63: 48]) * qy[ 47: 32] - (qx[ 47: 32] - qw[ 47: 32]) * qy[ 63: 48]) >> SAR

  add_res0[34: 0] = temp0[16:0] + temp2[33:0]
  add_res1[34: 0] = temp1[16:0] + temp3[33:0]
  sub_res0[34: 0] = temp0[16:0] - temp2[33:0]
  sub_res1[34: 0] = temp3[33:0] - temp1[16:0]

  qz[15:0] = sat_s16(add_res0[34:0])
  qz[31:16] = sat_s16(add_res1[34:0])
  qv[15:0] = sat_s16(sub_res0[34:0])
  qv[31:16] = sat_s16(sub_res1[34:0])

  qu = load128(as[31:0])
  rs1[31:0] = rs1[31:0] + 16
\end{lstlisting}

\textbf{Schedule}\\
\begin{longtable}[c]{|l|l|l|}
    \hline
    \rowcolor{lightgray}
    \textbf{Operands} & \textbf{Use} & \textbf{Def} \\\hline
    \endfirsthead
    \hline
    \rowcolor{lightgray}
     \textbf{Operands} & \textbf{Use} & \textbf{Def} \\\hline
    \endhead
    \hline
    \endfoot

  qu & - & 2 \\\hline
  rs1 & 1 & 1 \\\hline
  qz & - & 2 \\\hline
  qv & - & 2 \\\hline
  qx & 1 & - \\\hline
  qw & 1 & - \\\hline
  qy & 1 & - \\\hline
  SAR & 1 & -

\end{longtable}

\textbf{Format}\\

\begin{figure*}[!htbp]
 {\setlength{\tabcolsep}{1pt}
 \begin{center}
\begin{tabular}{@{}F@{}F@{}W@{}I@{}F@{}I@{}I@{}F@{}I@{}I@{}F@{}F@{}O}
\instbitrange{31}{29}&
\instbitrange{28}{26}&
\instbitrange{25}{24}&
\instbit{23}&
\instbitrange{22}{20}&
\instbit{19}&
\instbit{18}&
\instbitrange{17}{15}&
\instbit{14}&
\instbit{13}&
\instbitrange{12}{10}&
\instbitrange{9}{7}&
\instbitrange{6}{0} \\
\hline
\multicolumn{1}{|c|}{{\scriptsize qx[2:0]}} &
\multicolumn{1}{c|}{{\scriptsize qy[2:0]}} &
\multicolumn{1}{c|}{{\scriptsize qw[1:0]}} &
\multicolumn{1}{c|}{{\scriptsize sel2[0]}} &
\multicolumn{1}{c|}{{\scriptsize qv[2:0]}} &
\multicolumn{1}{c|}{{\scriptsize qw[2]}} &
\multicolumn{1}{c|}{{\scriptsize rs1[4]}} &
\multicolumn{1}{c|}{{\scriptsize rs1[2:0]}} &
\multicolumn{1}{c|}{{\scriptsize sat[0]}} &
\multicolumn{1}{c|}{0} &
\multicolumn{1}{c|}{{\scriptsize qu[2:0]}} &
\multicolumn{1}{c|}{{\scriptsize qz[2:0]}} &
\multicolumn{1}{c|}{1011011}\\
\hline
3& 3& 2& 1& 3& 1& 1& 3& 1& 1& 3& 3& 7 \\
\end{tabular}
 \end{center}
 }
 \vspace{-0.1in}
 \caption[]{ESP.FFT.AMS.S16.LD.INCP} \end{figure*}


\newpage
\subsubsection{ESP.FFT.AMS.S16.LD.INCP.UAUP}\label{instruction:ESPFFTAMSS16LDINCPUAUP}
\textbf{Syntax}\\
ESP.FFT.AMS.S16.LD.INCP.UAUP qu,rs1,qz,qv,qx,qw,qy,sel2,sat

\textbf{Description}\\
This is a dedicated FFT instruction designed to perform addition, subtraction, multiplication, addition and subtraction, and shift operations on 16-bit data segments.\\
During the operation, the 16-byte data are loaded from the memory. The instruction combines the loaded data and the data in the special register \lstinline{UA_STATE} into 32-byte data, right-shifts the 32-byte data by the result of the \lstinline{SAR_BYTE} value multiplied by 8, and assigns the lower 128 bits of the shifted result to the register \lstinline{qu}. Meanwhile, the register \lstinline{UA_STATE} is updated with the loaded 16-byte data. After the access, the value in the register \lstinline{rs1} is incremented by 16.\\

\textbf{Operation}
\begin{lstlisting}
  temp0[16:0] = qx[ 15:  0] + qw[ 15:  0]
  temp1[16:0] = qx[ 31: 16] - qw[ 31: 16]

if sel2==0:
  temp2[33:0] = ((qx[ 15:  0] - qw[ 15:  0]) * qy[ 15:  0] - (qx[ 31: 16] + qw[ 31: 16]) * qy[ 31: 16]) >> SAR
  temp3[33:0] = ((qx[ 15:  0] - qw[ 15:  0]) * qy[ 31: 16] + (qx[ 31: 16] + qw[ 31: 16]) * qy[ 15:  0]) >> SAR
if sel2==1:
  temp2[33:0] = ((qx[ 31: 16] + qw[ 31: 16]) * qy[ 31: 16] + (qx[ 15:  0] - qw[ 15:  0]) * qy[ 15:  0]) >> SAR
  temp3[33:0] = ((qx[ 31: 16] + qw[ 31: 16]) * qy[ 15:  0] - (qx[ 15:  0] - qw[ 15:  0]) * qy[ 31: 16]) >> SAR

add_res0[34: 0] = temp0[16:0] + temp2[33:0]
add_res1[34: 0] = temp1[16:0] + temp3[33:0]
sub_res0[34: 0] = temp0[16:0] - temp2[33:0]
sub_res1[34: 0] = temp3[33:0] - temp1[16:0]

qz[15:0] = sat_s16(add_res0[34:0])
qz[31:16] = sat_s16(add_res1[34:0])
qv[15:0] = sat_s16(sub_res0[34:0])
qv[31:16] = sat_s16(sub_res1[34:0])

  dataIn[127:0] = load128(as[31:0])

  qu[127: 0] = {dataIn[127:0], UA_STATE[127:0]} >> {SAR_BYTE[3:0] << 3}
  UA_STATE[127:0] = dataIn[127:0]
  rs1[31:0] = rs1[31:0] + 16
\end{lstlisting}

\textbf{Schedule}\\
\begin{longtable}[c]{|l|l|l|}
    \hline
    \rowcolor{lightgray}
    \textbf{Operands} & \textbf{Use} & \textbf{Def} \\\hline
    \endfirsthead
    \hline
    \rowcolor{lightgray}
     \textbf{Operands} & \textbf{Use} & \textbf{Def} \\\hline
    \endhead
    \hline
    \endfoot

  qu & - & 2 \\\hline
  rs1 & 1 & 1 \\\hline
  qz & - & 2 \\\hline
  qv & - & 2 \\\hline
  qx & 1 & - \\\hline
  qw & 1 & - \\\hline
  qy & 1 & - \\\hline
  SAR & 1 & - \\\hline
  SAR\_BYTES & 1 & -

\end{longtable}

\textbf{Format}\\

\begin{figure*}[!htbp]
 {\setlength{\tabcolsep}{1pt}
 \begin{center}
\begin{tabular}{@{}F@{}F@{}W@{}I@{}F@{}I@{}I@{}F@{}I@{}I@{}F@{}F@{}O}
\instbitrange{31}{29}&
\instbitrange{28}{26}&
\instbitrange{25}{24}&
\instbit{23}&
\instbitrange{22}{20}&
\instbit{19}&
\instbit{18}&
\instbitrange{17}{15}&
\instbit{14}&
\instbit{13}&
\instbitrange{12}{10}&
\instbitrange{9}{7}&
\instbitrange{6}{0} \\
\hline
\multicolumn{1}{|c|}{{\scriptsize qx[2:0]}} &
\multicolumn{1}{c|}{{\scriptsize qy[2:0]}} &
\multicolumn{1}{c|}{{\scriptsize qw[1:0]}} &
\multicolumn{1}{c|}{{\scriptsize sel2[0]}} &
\multicolumn{1}{c|}{{\scriptsize qv[2:0]}} &
\multicolumn{1}{c|}{{\scriptsize qw[2]}} &
\multicolumn{1}{c|}{{\scriptsize rs1[4]}} &
\multicolumn{1}{c|}{{\scriptsize rs1[2:0]}} &
\multicolumn{1}{c|}{{\scriptsize sat[0]}} &
\multicolumn{1}{c|}{0} &
\multicolumn{1}{c|}{{\scriptsize qu[2:0]}} &
\multicolumn{1}{c|}{{\scriptsize qz[2:0]}} &
\multicolumn{1}{c|}{0111011}\\
\hline
3& 3& 2& 1& 3& 1& 1& 3& 1& 1& 3& 3& 7 \\
\end{tabular}
 \end{center}
 }
 \vspace{-0.1in}
 \caption[]{ESP.FFT.AMS.S16.LD.INCP.UAUP} \end{figure*}


\newpage
\subsubsection{ESP.FFT.AMS.S16.LD.R32.DECP}\label{instruction:ESPFFTAMSS16LDR32DECP}
\textbf{Syntax}\\
ESP.FFT.AMS.S16.LD.R32.DECP qu,rs1,qz,qv,qx,qw,qy,sel2,sat

\textbf{Description}\\
It is a dedicated FFT instruction to perform addition, subtraction, multiplication, addition and subtraction, and shift operations on 16-bit data segments.\\
During the operation, the instruction forces the lower 4 bits of the access address in register \lstinline{rs1} to 0 and loads the 16-byte data from the memory to register \lstinline{qu} in the big-endian word order, namely, loads the segment [127: 96] of the data to [31:0] of \lstinline{qu}. After the access is completed, the value in register \lstinline{rs1} is decreased by 16.\\
\textbf{Operation}
\begin{lstlisting}
  temp0[15:0] = qx[ 79: 64] + qy[ 79: 64]
  temp1[15:0] = qx[ 95: 80] - qy[ 95: 80]

if sel2==0:
  temp2[15:0] = ((qx[ 79: 64] - qw[ 79: 64]) * qy[ 79: 64] - (qx[ 95: 80] + qw[ 95: 80]) * qy[ 95: 80]) >> SAR
  temp3[15:0] = ((qx[ 79: 64] - qw[ 79: 64]) * qy[ 95: 80] + (qx[ 95: 80] + qw[ 95: 80]) * qy[ 79: 64]) >> SAR
if sel2==1:
  temp2[15:0] = ((qx[ 95: 80] + qw[ 95: 80]) * qy[ 95: 80] + (qx[ 79: 64] - qw[ 79: 64]) * qy[ 79: 64]) >> SAR
  temp3[15:0] = ((qx[ 95: 80] + qw[ 95: 80]) * qy[ 79: 64] - (qx[ 79: 64] - qw[ 79: 64]) * qy[ 95: 80]) >> SAR

  qz[79: 64] = temp0[15:0] + temp2[15:0]
  qz[95: 80] = temp1[15:0] + temp3[15:0]
  qv[79:64] = temp0[15:0] - temp2[15:0]
  qv[95:80] = temp3[15:0] - temp1[15:0]
  {qu[31: 0], qu[63: 32], qu[95: 64], qu[127: 96]} = load128(rs1[31:0])
  rs1[31:0] = rs1[31:0] - 16
\end{lstlisting}

\textbf{Schedule}\\
\begin{longtable}[c]{|l|l|l|}
    \hline
    \rowcolor{lightgray}
    \textbf{Operands} & \textbf{Use} & \textbf{Def} \\\hline
    \endfirsthead
    \hline
    \rowcolor{lightgray}
     \textbf{Operands} & \textbf{Use} & \textbf{Def} \\\hline
    \endhead
    \hline
    \endfoot

  qu & - & 2 \\\hline
  rs1 & 1 & 1 \\\hline
  qz & - & 2 \\\hline
  qv & - & 2 \\\hline
  qx & 1 & - \\\hline
  qw & 1 & - \\\hline
  qy & 1 & - \\\hline
  SAR & 1 & -

\end{longtable}

\textbf{Format}\\

\begin{figure*}[!htbp]
 {\setlength{\tabcolsep}{1pt}
 \begin{center}
\begin{tabular}{@{}F@{}F@{}W@{}I@{}F@{}I@{}I@{}F@{}I@{}I@{}F@{}F@{}O}
\instbitrange{31}{29}&
\instbitrange{28}{26}&
\instbitrange{25}{24}&
\instbit{23}&
\instbitrange{22}{20}&
\instbit{19}&
\instbit{18}&
\instbitrange{17}{15}&
\instbit{14}&
\instbit{13}&
\instbitrange{12}{10}&
\instbitrange{9}{7}&
\instbitrange{6}{0} \\
\hline
\multicolumn{1}{|c|}{{\scriptsize qx[2:0]}} &
\multicolumn{1}{c|}{{\scriptsize qy[2:0]}} &
\multicolumn{1}{c|}{{\scriptsize qw[1:0]}} &
\multicolumn{1}{c|}{{\scriptsize sel2[0]}} &
\multicolumn{1}{c|}{{\scriptsize qv[2:0]}} &
\multicolumn{1}{c|}{{\scriptsize qw[2]}} &
\multicolumn{1}{c|}{{\scriptsize rs1[4]}} &
\multicolumn{1}{c|}{{\scriptsize rs1[2:0]}} &
\multicolumn{1}{c|}{{\scriptsize sat[0]}} &
\multicolumn{1}{c|}{1} &
\multicolumn{1}{c|}{{\scriptsize qu[2:0]}} &
\multicolumn{1}{c|}{{\scriptsize qz[2:0]}} &
\multicolumn{1}{c|}{0111011}\\
\hline
3& 3& 2& 1& 3& 1& 1& 3& 1& 1& 3& 3& 7 \\
\end{tabular}
 \end{center}
 }
 \vspace{-0.1in}
 \caption[]{ESP.FFT.AMS.S16.LD.R32.DECP} \end{figure*}


\newpage
\subsubsection{ESP.FFT.AMS.S16.ST.INCP}\label{instruction:ESPFFTAMSS16STINCP}
\textbf{Syntax}\\
ESP.FFT.AMS.S16.ST.INCP qu,qz,rs2,rs1,qx,qw,qy,sel2,sat

\textbf{Description}\\
It is a dedicated FFT instruction to perform addition, subtraction, multiplication, addition and subtraction, and shift operations on 16-bit data segments.\\
During the operation, splices the values in registers \lstinline{qu} and \lstinline{rs2} into 16-byte data, and stores the spliced result to memory. After the access is completed, the value in register \lstinline{rs1} is incremented by 16. Besides, the operation result is updated to register \lstinline{rs2}.\\
\textbf{Operation}
\begin{lstlisting}
  temp0[15:0] = qx[111: 96] + qw[111: 96]
  temp1[15:0] = qx[127:112] - qw[127:112]

if sel2==0:
  temp2[15:0] = ((qx[111: 96] - qw[111: 96]) * qy[111: 96] - (qx[127:112] + qw[127:112]) * qy[127:112]) >> SAR[5:0]
  temp3[15:0] = ((qx[111: 96] - qw[111: 96]) * qy[127:112] + (qx[127:112] + qw[127:112]) * qy[111: 96]) >> SAR[5:0]
  {qv[ 95: 80] >> 1, qv[ 79: 64] >> 1, qv[ 63: 48] >> 1, qv[ 47: 32] >> 1, qv[ 31: 16] >> 1, qv[ 15:  0] >> 1, rs2[31: 16] >> 1, rs2[15:0] >> 1} => store128(rs2[31:0])
if sel2==1:
  temp2[15:0] = ((qx[127:112] + qw[127:112]) * qy[127:112] + (qx[111: 96] - qw[111: 96]) * qy[111: 96]) >> SAR[5:0]
  temp3[15:0] = ((qx[127:112] + qw[127:112]) * qy[111: 96] - (qx[111: 96] - qw[111: 96]) * qy[127:112]) >> SAR[5:0]
  {qu[ 95: 64], qu[ 63: 32], qu[ 31:  0], rs2[31:0]} => store128(rs1[31:0])

  temp4[16:0] = temp1[15:0] + temp3[15:0]
  temp5[16:0] = temp0[15:0] + temp2[15:0]
  qz[111: 96] = temp0[15:0] - temp2[15:0]
  qz[127:112] = temp3[15:0] - temp1[15:0]
  rs2 = {temp4[15:0], temp5[15:0]}
  rs1[31:0] = rs1[31:0] +16
\end{lstlisting}

\textbf{Schedule}\\
\begin{longtable}[c]{|l|l|l|}
    \hline
    \rowcolor{lightgray}
    \textbf{Operands} & \textbf{Use} & \textbf{Def} \\\hline
    \endfirsthead
    \hline
    \rowcolor{lightgray}
     \textbf{Operands} & \textbf{Use} & \textbf{Def} \\\hline
    \endhead
    \hline
    \endfoot

  qu & 2 & - \\\hline
  rs1 & 1 & 1 \\\hline
  qz & - & 2 \\\hline
  qx & 1 & - \\\hline
  qw & 1 & - \\\hline
  qy & 1 & - \\\hline
  SAR & 1 & -

\end{longtable}

\textbf{Format}\\

\begin{figure*}[!htbp]
 {\setlength{\tabcolsep}{1pt}
 \begin{center}
\begin{tabular}{@{}F@{}F@{}W@{}I@{}F@{}I@{}I@{}F@{}I@{}I@{}F@{}F@{}O}
\instbitrange{31}{29}&
\instbitrange{28}{26}&
\instbitrange{25}{24}&
\instbit{23}&
\instbitrange{22}{20}&
\instbit{19}&
\instbit{18}&
\instbitrange{17}{15}&
\instbit{14}&
\instbit{13}&
\instbitrange{12}{10}&
\instbitrange{9}{7}&
\instbitrange{6}{0} \\
\hline
\multicolumn{1}{|c|}{{\scriptsize qx[2:0]}} &
\multicolumn{1}{c|}{{\scriptsize qy[2:0]}} &
\multicolumn{1}{c|}{{\scriptsize qw[1:0]}} &
\multicolumn{1}{c|}{{\scriptsize rs2[4]}} &
\multicolumn{1}{c|}{{\scriptsize rs2[2:0]}} &
\multicolumn{1}{c|}{{\scriptsize qw[2]}} &
\multicolumn{1}{c|}{{\scriptsize rs1[4]}} &
\multicolumn{1}{c|}{{\scriptsize rs1[2:0]}} &
\multicolumn{1}{c|}{{\scriptsize sel2[0]}} &
\multicolumn{1}{c|}{{\scriptsize sat[0]}} &
\multicolumn{1}{c|}{{\scriptsize qu[2:0]}} &
\multicolumn{1}{c|}{{\scriptsize qz[2:0]}} &
\multicolumn{1}{c|}{0111111}\\
\hline
3& 3& 2& 1& 3& 1& 1& 3& 1& 1& 3& 3& 7 \\
\end{tabular}
 \end{center}
 }
 \vspace{-0.1in}
 \caption[]{ESP.FFT.AMS.S16.ST.INCP} \end{figure*}


\newpage
\subsubsection{ESP.FFT.CMUL.S16.LD.XP}\label{instruction:ESPFFTCMULS16LDXP}
\textbf{Syntax}\\
ESP.FFT.CMUL.S16.LD.XP qu,rs1,rs2,qz,qy,qx,sel8,sat

\textbf{Description}\\
This instruction performs a 16-bit signed complex multiplication. It is similar to \hyperref[instruction:ESPCMULS16]{ESP.CMUL.S16}, except that the order of registers \lstinline{qx} and \lstinline{qy} is reversed.\\
During the operation, the 16-byte data are loaded from the memory to the register \lstinline{qu}. After the access is completed, the value in the register \lstinline{rs1} is incremented by the value in the register \lstinline{rs2}.\\
\textbf{Operation}
\begin{lstlisting}
if sel8 == 0:
  qz[ 15:  0] = (qy[ 15:  0] * qx[ 15:  0] + qy[ 31: 16] * qx[ 31: 16]) >> SAR[5:0]
  qz[ 31: 16] = (qy[ 31: 16] * qx[ 15:  0] - qy[ 15:  0] * qx[ 31: 16]) >> SAR[5:0]
if sel8 == 1:
  qz[ 15:  0] = (qy[ 15:  0] * qx[ 15:  0] - qy[ 31: 16] * qx[ 31: 16]) >> SAR[5:0]
  qz[ 31: 16] = (qy[ 31: 16] * qx[ 15:  0] + qy[ 15:  0] * qx[ 31: 16]) >> SAR[5:0]
if sel8 == 2:
  qz[ 47: 32] = (qy[ 47: 32] * qx[ 47: 32] + qy[ 63: 48] * qx[ 63: 48]) >> SAR[5:0]
  qz[ 63: 48] = (qy[ 63: 48] * qx[ 47: 32] - qy[ 47: 32] * qx[ 63: 48]) >> SAR[5:0]
if sel8 == 3:
  qz[ 47: 32] = (qy[ 47: 32] * qx[ 47: 32] - qy[ 63: 48] * qx[ 63: 48]) >> SAR[5:0]
  qz[ 63: 48] = (qy[ 63: 48] * qx[ 47: 32] + qy[ 47: 32] * qx[ 63: 48]) >> SAR[5:0]
if sel8 == 4:
  qz[ 79: 64] = (qy[ 79: 64] * qx[ 79: 64] + qy[ 95: 80] * qx[ 95: 80]) >> SAR[5:0]
  qz[ 95: 80] = (qy[ 95: 80] * qx[ 79: 64] - qy[ 79: 64] * qx[ 95: 80]) >> SAR[5:0]
if sel8 == 5:
  qz[ 79: 64] = (qy[ 79: 64] * qx[ 79: 64] - qy[ 95: 80] * qx[ 95: 80]) >> SAR[5:0]
  qz[ 95: 80] = (qy[ 95: 80] * qx[ 79: 64] + qy[ 79: 64] * qx[ 95: 80]) >> SAR[5:0]

  qu[127:0] = load128(rs1[31:0])
  rs1[31:0] = rs1[31:0] + rs2[31:0]
\end{lstlisting}

\textbf{Schedule}\\
\begin{longtable}[c]{|l|l|l|}
    \hline
    \rowcolor{lightgray}
    \textbf{Operands} & \textbf{Use} & \textbf{Def} \\\hline
    \endfirsthead
    \hline
    \rowcolor{lightgray}
     \textbf{Operands} & \textbf{Use} & \textbf{Def} \\\hline
    \endhead
    \hline
    \endfoot

  qu & - & 2 \\\hline
  rs1 & 1 & 1 \\\hline
  rs2 & 1 & - \\\hline
  qz & - & 2 \\\hline
  qx & 1 & - \\\hline
  qy & 1 & - \\\hline
  SAR & 1 & -

\end{longtable}

\textbf{Format}\\

\begin{figure*}[!htbp]
 {\setlength{\tabcolsep}{1pt}
 \begin{center}
\begin{tabular}{@{}F@{}F@{}I@{}I@{}I@{}F@{}I@{}I@{}F@{}I@{}I@{}F@{}F@{}O}
\instbitrange{31}{29}&
\instbitrange{28}{26}&
\instbit{25}&
\instbit{24}&
\instbit{23}&
\instbitrange{22}{20}&
\instbit{19}&
\instbit{18}&
\instbitrange{17}{15}&
\instbit{14}&
\instbit{13}&
\instbitrange{12}{10}&
\instbitrange{9}{7}&
\instbitrange{6}{0} \\
\hline
\multicolumn{1}{|c|}{{\scriptsize qx[2:0]}} &
\multicolumn{1}{c|}{{\scriptsize qy[2:0]}} &
\multicolumn{1}{c|}{{\scriptsize sat[0]}} &
\multicolumn{1}{c|}{{\scriptsize sel8[2]}} &
\multicolumn{1}{c|}{{\scriptsize rs2[4]}} &
\multicolumn{1}{c|}{{\scriptsize rs2[2:0]}} &
\multicolumn{1}{c|}{{\scriptsize sel8[1]}} &
\multicolumn{1}{c|}{{\scriptsize rs1[4]}} &
\multicolumn{1}{c|}{{\scriptsize rs1[2:0]}} &
\multicolumn{1}{c|}{{\scriptsize sel8[0]}} &
\multicolumn{1}{c|}{1} &
\multicolumn{1}{c|}{{\scriptsize qu[2:0]}} &
\multicolumn{1}{c|}{{\scriptsize qz[2:0]}} &
\multicolumn{1}{c|}{1011011}\\
\hline
3& 3& 1& 1& 1& 3& 1& 1& 3& 1& 1& 3& 3& 7 \\
\end{tabular}
 \end{center}
 }
 \vspace{-0.1in}
 \caption[]{ESP.FFT.CMUL.S16.LD.XP} \end{figure*}


\newpage
\subsubsection{ESP.FFT.CMUL.S16.ST.XP}\label{instruction:ESPFFTCMULS16STXP}
\textbf{Syntax}\\
ESP.FFT.CMUL.S16.ST.XP qy,qx,qu,rs1,rs2,sel8,upd4,sel4,sat

\textbf{Description}\\
This instruction performs a 16-bit signed complex multiplication. It is similar to \hyperref[instruction:ESPCMULS16]{ESP.CMUL.S16}, except that the order of registers \lstinline{qx} and \lstinline{qy} is reversed.\\
The result of the operation and the data segments in registers \lstinline{qx} and \lstinline{qu}, specified by the immediate data \lstinline{upd4}, are concatenated into 128 bits. These are then written into memory. After the access is completed, the value in register \lstinline{rs1} is incremented by the value in register \lstinline{rs2}.\\
\textbf{Operation}
\begin{lstlisting}
if sel8 == 6:
  temp[15: 0] = (qy[111: 96] * qx[111: 96] + qy[127:112] * qx[127:112]) >> SAR[5:0]
  temp[31:16] = (qy[127:112] * qx[111: 96] - qy[111: 96] * qx[127:112]) >> SAR[5:0]
if sel8 == 7:
  temp[15: 0] = (qy[111: 96] * qx[111: 96] - qy[127:112] * qx[127:112]) >> SAR[5:0]
  temp[31:16] = (qy[127:112] * qx[111: 96] + qy[111: 96] * qx[127:112]) >> SAR[5:0]

if upd4 == 0:  // normal radix 2
  {temp[31:0], qu[ 95:  0]} => store128(rs1[31:0])
if upd4 == 1:  // radix2 last second stage
  {
       temp[31:0],
       qu[ 95: 64],
       qy[ 63: 48] >> sel4,
       qy[ 47: 32] >> sel4,
       qy[ 31: 16] >> sel4,
       qy[ 15:  0] >> sel4
  } => store128(rs1[31:0])
if upd4 == 2:  // radix2 last stage
  {
       temp[31:0],
       qy[ 63: 48] >> sel4,
       qy[ 47: 32] >> sel4,
       qu[ 95: 64],
       qy[ 31: 16] >> sel4,
       qy[ 15:  0] >> sel4
  } => store128(rs1[31:0])

  rs1[31:0] = rs1[31:0] + rs2[31:0]
\end{lstlisting}

\textbf{Schedule}\\
\begin{longtable}[c]{|l|l|l|}
    \hline
    \rowcolor{lightgray}
    \textbf{Operands} & \textbf{Use} & \textbf{Def} \\\hline
    \endfirsthead
    \hline
    \rowcolor{lightgray}
     \textbf{Operands} & \textbf{Use} & \textbf{Def} \\\hline
    \endhead
    \hline
    \endfoot

  qu & 2 & - \\\hline
  rs1 & 1 & 1 \\\hline
  rs2 & 1 & - \\\hline
  qx & 1 & - \\\hline
  qy & 1 & - \\\hline
  SAR & 1 & -

\end{longtable}

\textbf{Format}\\

\begin{figure*}[!htbp]
 {\setlength{\tabcolsep}{1pt}
 \begin{center}
\begin{tabular}{@{}F@{}F@{}I@{}I@{}I@{}F@{}I@{}I@{}F@{}W@{}F@{}F@{}O}
\instbitrange{31}{29}&
\instbitrange{28}{26}&
\instbit{25}&
\instbit{24}&
\instbit{23}&
\instbitrange{22}{20}&
\instbit{19}&
\instbit{18}&
\instbitrange{17}{15}&
\instbitrange{14}{13}&
\instbitrange{12}{10}&
\instbitrange{9}{7}&
\instbitrange{6}{0} \\
\hline
\multicolumn{1}{|c|}{{\scriptsize qx[2:0]}} &
\multicolumn{1}{c|}{{\scriptsize qy[2:0]}} &
\multicolumn{1}{c|}{{\scriptsize sat[0]}} &
\multicolumn{1}{c|}{{\scriptsize sel4[1]}} &
\multicolumn{1}{c|}{{\scriptsize rs2[4]}} &
\multicolumn{1}{c|}{{\scriptsize rs2[2:0]}} &
\multicolumn{1}{c|}{{\scriptsize sel4[0]}} &
\multicolumn{1}{c|}{{\scriptsize rs1[4]}} &
\multicolumn{1}{c|}{{\scriptsize rs1[2:0]}} &
\multicolumn{1}{c|}{{\scriptsize upd4[1:0]}} &
\multicolumn{1}{c|}{{\scriptsize qu[2:0]}} &
\multicolumn{1}{c|}{{\scriptsize sel8[2:0]}} &
\multicolumn{1}{c|}{1111111}\\
\hline
3& 3& 1& 1& 1& 3& 1& 1& 3& 2& 3& 3& 7 \\
\end{tabular}
 \end{center}
 }
 \vspace{-0.1in}
 \caption[]{ESP.FFT.CMUL.S16.ST.XP} \end{figure*}


\newpage
\subsubsection{ESP.FFT.R2BF.S16}\label{instruction:ESPFFTR2BFS16}
\textbf{Syntax}\\
ESP.FFT.R2BF.S16 qz,qv,qx,qy,sel2,sat

\textbf{Description}\\
This instruction performs the radix-2 butterfly operation on the 16-byte values in registers \lstinline{qx} and \lstinline{qy}, with a data bit width of 16 bits. Some of the calculation results are written to the register \lstinline{qz}, while others are written to the register \lstinline{qv}.\\
\textbf{Operation}
\begin{lstlisting}
if sel2==0:
  op_a[127:  0] = {qy[ 63:  0], qx[ 63:  0]}
  op_b[127:  0] = {qy[127: 64], qx[127: 64]}
if sel2==1:
  op_a[127:  0] = {qy[ 95: 64], qy[ 31:  0], qx[ 95: 64], qx[ 31:  0]}
  op_b[127:  0] = {qy[127: 96], qy[ 63: 32], qx[127: 96], qx[ 63: 32]}

  qz[ 15:  0] = op_a[ 15:  0] + op_b[ 15:  0]
  qz[ 31: 16] = op_a[ 31: 16] + op_b[ 31: 16]
  qz[ 47: 32] = op_a[ 47: 32] + op_b[ 47: 32]
  qz[ 63: 48] = op_a[ 63: 48] + op_b[ 63: 48]
  qz[ 79: 64] = op_a[ 15:  0] - op_b[ 15:  0]
  qz[ 95: 80] = op_a[ 31: 16] - op_b[ 31: 16]
  qz[111: 96] = op_a[ 47: 32] - op_b[ 47: 32]
  qz[127:112] = op_a[ 63: 48] - op_b[ 63: 48]

  qv[ 15:  0] = op_a[ 79: 64] + op_b[ 79: 64]
  qv[ 31: 16] = op_a[ 95: 80] + op_b[ 95: 80]
  qv[ 47: 32] = op_a[111: 96] + op_b[111: 96]
  qv[ 63: 48] = op_a[127:112] + op_b[127:112]
  qv[ 79: 64] = op_a[ 79: 64] - op_b[ 79: 64]
  qv[ 95: 80] = op_a[ 95: 80] - op_b[ 95: 80]
  qv[111: 96] = op_a[111: 96] - op_b[111: 96]
  qv[127:112] = op_a[127:112] - op_b[127:112]
\end{lstlisting}

\textbf{Schedule}\\
\begin{longtable}[c]{|l|l|l|}
    \hline
    \rowcolor{lightgray}
    \textbf{Operands} & \textbf{Use} & \textbf{Def} \\\hline
    \endfirsthead
    \hline
    \rowcolor{lightgray}
     \textbf{Operands} & \textbf{Use} & \textbf{Def} \\\hline
    \endhead
    \hline
    \endfoot

  qz & - & 1 \\\hline
  qv & - & 1 \\\hline
  qx & 1 & - \\\hline
  qy & 1 & -

\end{longtable}

\textbf{Format}\\

\begin{figure*}[!htbp]
 {\setlength{\tabcolsep}{1pt}
 \begin{center}
\begin{tabular}{@{}F@{}F@{}F@{}F@{}I@{}I@{}I@{}O@{}F@{}O}
\instbitrange{31}{29}&
\instbitrange{28}{26}&
\instbitrange{25}{23}&
\instbitrange{22}{20}&
\instbit{19}&
\instbit{18}&
\instbit{17}&
\instbitrange{16}{10}&
\instbitrange{9}{7}&
\instbitrange{6}{0} \\
\hline
\multicolumn{1}{|c|}{{\scriptsize qx[2:0]}} &
\multicolumn{1}{c|}{{\scriptsize qy[2:0]}} &
\multicolumn{1}{c|}{100} &
\multicolumn{1}{c|}{{\scriptsize qv[2:0]}} &
\multicolumn{1}{c|}{X} &
\multicolumn{1}{c|}{{\scriptsize sel2[0]}} &
\multicolumn{1}{c|}{{\scriptsize sat[0]}} &
\multicolumn{1}{c|}{XX000XX} &
\multicolumn{1}{c|}{{\scriptsize qz[2:0]}} &
\multicolumn{1}{c|}{0011111}\\
\hline
3& 3& 3& 3& 1& 1& 1& 7& 3& 7 \\
\end{tabular}
 \end{center}
 }
 \vspace{-0.1in}
 \caption[]{ESP.FFT.R2BF.S16} \end{figure*}


\newpage
\subsubsection{ESP.FFT.R2BF.S16.ST.INCP}\label{instruction:ESPFFTR2BFS16STINCP}
\textbf{Syntax}\\
ESP.FFT.R2BF.S16.ST.INCP qz,qx,qy,rs1,sel4,sat

\textbf{Description}\\
This instruction performs the radix-2 butterfly operation on the 16-byte values in registers \lstinline{qx} and \lstinline{qy}, and the data bit width is 16 bits. Some of the calculation results are written to the register \lstinline{qz}, while others perform an arithmetic shift and are written into the memory address indicated by \lstinline{rs1}. After the access is completed, the value in the register \lstinline{rs1} is incremented by 16.\\
\textbf{Operation}
\begin{lstlisting}
  qz[ 15:  0] = qx[ 15:  0] - qy[ 15:  0]
  qz[ 31: 16] = qx[ 31: 16] - qy[ 31: 16]
  qz[ 47: 32] = qx[ 47: 32] - qy[ 47: 32]
  qz[ 63: 48] = qx[ 63: 48] - qy[ 63: 48]
  qz[ 79: 64] = qx[ 79: 64] - qy[ 79: 64]
  qz[ 95: 80] = qx[ 95: 80] - qy[ 95: 80]
  qz[111: 96] = qx[111: 96] - qy[111: 96]
  qz[127:112] = qx[127:112] - qy[127:112]

  {
    (qx[127:112] + qy[127:112]) >> sel4,
    (qx[111: 96] + qy[111: 96]) >> sel4,
    (qx[ 95: 80] + qy[ 95: 80]) >> sel4,
    (qx[ 79: 64] + qy[ 79: 64]) >> sel4,
    (qx[ 63: 48] + qy[ 63: 48]) >> sel4,
    (qx[ 47: 32] + qy[ 47: 32]) >> sel4,
    (qx[ 31: 16] + qy[ 31: 16]) >> sel4,
    (qx[ 15:  0] + qy[ 15:  0]) >> sel4
  } => store128(rs1[31:0])
  rs1[31:0] = rs1[31:0] + 16
\end{lstlisting}

\textbf{Schedule}\\
\begin{longtable}[c]{|l|l|l|}
    \hline
    \rowcolor{lightgray}
    \textbf{Operands} & \textbf{Use} & \textbf{Def} \\\hline
    \endfirsthead
    \hline
    \rowcolor{lightgray}
     \textbf{Operands} & \textbf{Use} & \textbf{Def} \\\hline
    \endhead
    \hline
    \endfoot

  qz & - & 1 \\\hline
  rs1 & 1 & 1 \\\hline
  qx & 1 & - \\\hline
  qy & 1 & -

\end{longtable}

\textbf{Format}\\

\begin{figure*}[!htbp]
 {\setlength{\tabcolsep}{1pt}
 \begin{center}
\begin{tabular}{@{}F@{}F@{}O@{}I@{}F@{}W@{}I@{}W@{}F@{}O}
\instbitrange{31}{29}&
\instbitrange{28}{26}&
\instbitrange{25}{19}&
\instbit{18}&
\instbitrange{17}{15}&
\instbitrange{14}{13}&
\instbit{12}&
\instbitrange{11}{10}&
\instbitrange{9}{7}&
\instbitrange{6}{0} \\
\hline
\multicolumn{1}{|c|}{{\scriptsize qx[2:0]}} &
\multicolumn{1}{c|}{{\scriptsize qy[2:0]}} &
\multicolumn{1}{c|}{111101X} &
\multicolumn{1}{c|}{{\scriptsize rs1[4]}} &
\multicolumn{1}{c|}{{\scriptsize rs1[2:0]}} &
\multicolumn{1}{c|}{00} &
\multicolumn{1}{c|}{{\scriptsize sat[0]}} &
\multicolumn{1}{c|}{{\scriptsize sel4[1:0]}} &
\multicolumn{1}{c|}{{\scriptsize qz[2:0]}} &
\multicolumn{1}{c|}{0011111}\\
\hline
3& 3& 7& 1& 3& 2& 1& 2& 3& 7 \\
\end{tabular}
 \end{center}
 }
 \vspace{-0.1in}
 \caption[]{ESP.FFT.R2BF.S16.ST.INCP} \end{figure*}


\newpage
\subsubsection{ESP.FFT.VST.R32.DECP}\label{instruction:ESPFFTVSTR32DECP}
\textbf{Syntax}\\
ESP.FFT.VST.R32.DECP qu,rs1,sel2

\textbf{Description}\\
It is a dedicated FFT instruction. This instruction divides data in register \lstinline {qu} into 8 segments of 16-bit data and performs an arithmetic right shift on them by 0 or 1 depending on the immediate number \lstinline{sel2}, and finally writes the result to the memory address indicated by register \lstinline{rs1} in word big-endian order. After the access is completed, the value in register \lstinline{rs1} is decreased by 16.\\
\textbf{Operation}
\begin{lstlisting}
  {
    qu[ 31: 16] >> sel2,
    qu[ 15:  0] >> sel2,
    qu[ 63: 48] >> sel2,
    qu[ 47: 32] >> sel2,
    qu[ 95: 80] >> sel2,
    qu[ 79: 64] >> sel2,
    qu[127:112] >> sel2,
    qu[111: 96] >> sel2
  } => store128(rs1[31:0])
  rs1[31:0] = rs1[31:0] - 16
\end{lstlisting}

\textbf{Schedule}\\
\begin{longtable}[c]{|l|l|l|}
    \hline
    \rowcolor{lightgray}
    \textbf{Operands} & \textbf{Use} & \textbf{Def} \\\hline
    \endfirsthead
    \hline
    \rowcolor{lightgray}
     \textbf{Operands} & \textbf{Use} & \textbf{Def} \\\hline
    \endhead
    \hline
    \endfoot

  qu & 2 & - \\\hline
  rs1 & 1 & 1

\end{longtable}

\textbf{Format}\\

\begin{figure*}[!htbp]
 {\setlength{\tabcolsep}{1pt}
 \begin{center}
\begin{tabular}{@{}I@{}M@{}I@{}F@{}W@{}F@{}F@{}O}
\instbit{31}&
\instbitrange{30}{19}&
\instbit{18}&
\instbitrange{17}{15}&
\instbitrange{14}{13}&
\instbitrange{12}{10}&
\instbitrange{9}{7}&
\instbitrange{6}{0} \\
\hline
\multicolumn{1}{|c|}{{\scriptsize sel2[0]}} &
\multicolumn{1}{c|}{XXXXX110010X} &
\multicolumn{1}{c|}{{\scriptsize rs1[4]}} &
\multicolumn{1}{c|}{{\scriptsize rs1[2:0]}} &
\multicolumn{1}{c|}{00} &
\multicolumn{1}{c|}{{\scriptsize qu[2:0]}} &
\multicolumn{1}{c|}{XXX} &
\multicolumn{1}{c|}{0011111}\\
\hline
1& 12& 1& 3& 2& 3& 3& 7 \\
\end{tabular}
 \end{center}
 }
 \vspace{-0.1in}
 \caption[]{ESP.FFT.VST.R32.DECP} \end{figure*}


\newpage
\subsubsection{ESP.LD.128.USAR.IP}\label{instruction:ESPLD128USARIP}
\textbf{Syntax}\\
ESP.LD.128.USAR.IP qu,rs1,imm16

\textbf{Description}\\
This instruction loads 16-byte data from memory to the register \lstinline{qu}. Meanwhile, it saves the value of the lower 4 bits in the register \lstinline{rs1} to the special register \lstinline{SAR_BYTE}. After the access is completed, the value in the register \lstinline{rs1} is incremented by an 8-bit sign-extended constant in the instruction code segment, left-shifted by 4.\\

imm starts at -2048, ends at 2032, and steps by 16.

\textbf{Operation}
\begin{lstlisting}
  qu[127:0] = load128(rs1[31:0])
  SAR_BYTE = rs1[3:0]
  rs1[31:0] = rs1[31:0] + imm
\end{lstlisting}

\textbf{Schedule}\\
\begin{longtable}[c]{|l|l|l|}
    \hline
    \rowcolor{lightgray}
    \textbf{Operands} & \textbf{Use} & \textbf{Def} \\\hline
    \endfirsthead
    \hline
    \rowcolor{lightgray}
     \textbf{Operands} & \textbf{Use} & \textbf{Def} \\\hline
    \endhead
    \hline
    \endfoot

  qu & - & 2 \\\hline
  rs1 & 1 & 1 \\\hline
  SAR\_BYTE & - & 1

\end{longtable}

\textbf{Format}\\

\begin{figure*}[!htbp]
 {\setlength{\tabcolsep}{1pt}
 \begin{center}
\begin{tabular}{@{}S@{}O@{}I@{}F@{}W@{}F@{}W@{}I@{}O}
\instbitrange{31}{26}&
\instbitrange{25}{19}&
\instbit{18}&
\instbitrange{17}{15}&
\instbitrange{14}{13}&
\instbitrange{12}{10}&
\instbitrange{9}{8}&
\instbit{7}&
\instbitrange{6}{0} \\
\hline
\multicolumn{1}{|c|}{{\scriptsize imm16[7:2]}} &
\multicolumn{1}{c|}{111010X} &
\multicolumn{1}{c|}{{\scriptsize rs1[4]}} &
\multicolumn{1}{c|}{{\scriptsize rs1[2:0]}} &
\multicolumn{1}{c|}{00} &
\multicolumn{1}{c|}{{\scriptsize qu[2:0]}} &
\multicolumn{1}{c|}{{\scriptsize imm16[1:0]}} &
\multicolumn{1}{c|}{0} &
\multicolumn{1}{c|}{0011111}\\
\hline
6& 7& 1& 3& 2& 3& 2& 1& 7 \\
\end{tabular}
 \end{center}
 }
 \vspace{-0.1in}
 \caption[]{ESP.LD.128.USAR.IP} \end{figure*}


\newpage
\subsubsection{ESP.LD.128.USAR.XP}\label{instruction:ESPLD128USARXP}
\textbf{Syntax}\\
ESP.LD.128.USAR.XP qu,rs1,rs2

\textbf{Description}\\
This instruction loads 16-byte data from memory into the register \lstinline{qu}. Meanwhile, it saves the value of the lower 4 bits in \lstinline{rs1} into the special register \lstinline{SAR_BYTE}. After the access is completed, the value in the register \lstinline{rs1} is incremented by the value in the register \lstinline{rs2}.\\
\textbf{Operation}
\begin{lstlisting}
  qu[127:0] = load128(rs1[31:0])
  SAR_BYTE = rs1[3:0]
  rs1[31:0] = rs1[31:0] + rs2[31:0]
\end{lstlisting}

\textbf{Schedule}\\
\begin{longtable}[c]{|l|l|l|}
    \hline
    \rowcolor{lightgray}
    \textbf{Operands} & \textbf{Use} & \textbf{Def} \\\hline
    \endfirsthead
    \hline
    \rowcolor{lightgray}
     \textbf{Operands} & \textbf{Use} & \textbf{Def} \\\hline
    \endhead
    \hline
    \endfoot

  qu & - & 2 \\\hline
  rs1 & 1 & 1 \\\hline
  rs2 & 1 & - \\\hline
  SAR\_BYTE & - & 1

\end{longtable}

\textbf{Format}\\

\begin{figure*}[!htbp]
 {\setlength{\tabcolsep}{1pt}
 \begin{center}
\begin{tabular}{@{}E@{}I@{}F@{}I@{}I@{}F@{}W@{}F@{}F@{}O}
\instbitrange{31}{24}&
\instbit{23}&
\instbitrange{22}{20}&
\instbit{19}&
\instbit{18}&
\instbitrange{17}{15}&
\instbitrange{14}{13}&
\instbitrange{12}{10}&
\instbitrange{9}{7}&
\instbitrange{6}{0} \\
\hline
\multicolumn{1}{|c|}{01X0XXXX} &
\multicolumn{1}{c|}{{\scriptsize rs2[4]}} &
\multicolumn{1}{c|}{{\scriptsize rs2[2:0]}} &
\multicolumn{1}{c|}{0} &
\multicolumn{1}{c|}{{\scriptsize rs1[4]}} &
\multicolumn{1}{c|}{{\scriptsize rs1[2:0]}} &
\multicolumn{1}{c|}{11} &
\multicolumn{1}{c|}{{\scriptsize qu[2:0]}} &
\multicolumn{1}{c|}{0X0} &
\multicolumn{1}{c|}{0011011}\\
\hline
8& 1& 3& 1& 1& 3& 2& 3& 3& 7 \\
\end{tabular}
 \end{center}
 }
 \vspace{-0.1in}
 \caption[]{ESP.LD.128.USAR.XP} \end{figure*}


\newpage
\subsubsection{ESP.LD.XACC.IP}\label{instruction:ESPLDXACCIP}
\textbf{Syntax}\\
ESP.LD.XACC.IP rs1,imm8

\textbf{Description}\\
Loads 64-bit data from memory, and saves its lower 40 bits to the special register \lstinline{XACC}. After the access is completed, the value in the register \lstinline{rs1} is incremented by an 8-bit sign-extended constant in the instruction code segment left-shifted by 3.\\

imm starts at -1024, ends at 1016, and the step is 8.

\textbf{Operation}
\begin{lstlisting}
  DataIn[63:0] = load64(rs1[31:0])
  XACC[39:0] = DataIn[39:0]
  rs1[31:0] = rs1[31:0] + imm
\end{lstlisting}

\textbf{Schedule}\\
\begin{longtable}[c]{|l|l|l|}
    \hline
    \rowcolor{lightgray}
    \textbf{Operands} & \textbf{Use} & \textbf{Def} \\\hline
    \endfirsthead
    \hline
    \rowcolor{lightgray}
     \textbf{Operands} & \textbf{Use} & \textbf{Def} \\\hline
    \endhead
    \hline
    \endfoot

  XACC & - & 2 \\\hline
  rs1 & 1 & 1

\end{longtable}

\textbf{Format}\\

\begin{figure*}[!htbp]
 {\setlength{\tabcolsep}{1pt}
 \begin{center}
\begin{tabular}{@{}F@{}F@{}W@{}F@{}W@{}I@{}F@{}W@{}W@{}Y@{}O}
\instbitrange{31}{29}&
\instbitrange{28}{26}&
\instbitrange{25}{24}&
\instbitrange{23}{21}&
\instbitrange{20}{19}&
\instbit{18}&
\instbitrange{17}{15}&
\instbitrange{14}{13}&
\instbitrange{12}{11}&
\instbitrange{10}{7}&
\instbitrange{6}{0} \\
\hline
\multicolumn{1}{|c|}{0XX} &
\multicolumn{1}{c|}{{\scriptsize imm8[7:5]}} &
\multicolumn{1}{c|}{00} &
\multicolumn{1}{c|}{{\scriptsize imm8[4:2]}} &
\multicolumn{1}{c|}{1X} &
\multicolumn{1}{c|}{{\scriptsize rs1[4]}} &
\multicolumn{1}{c|}{{\scriptsize rs1[2:0]}} &
\multicolumn{1}{c|}{00} &
\multicolumn{1}{c|}{{\scriptsize imm8[1:0]}} &
\multicolumn{1}{c|}{X110} &
\multicolumn{1}{c|}{0011111}\\
\hline
3& 3& 2& 3& 2& 1& 3& 2& 2& 4& 7 \\
\end{tabular}
 \end{center}
 }
 \vspace{-0.1in}
 \caption[]{ESP.LD.XACC.IP} \end{figure*}


\newpage
\subsubsection{ESP.LDQA.S16.128.IP}\label{instruction:ESPLDQAS16128IP}
\textbf{Syntax}\\
ESP.LDQA.S16.128.IP rs1,imm16

\textbf{Description}\\
This instruction loads 16-byte data from memory, divides it into 8 segments of 16 bits, sign-extends each segment to 64 bits, and then stores the results to the 256-bit special registers \lstinline{QACC_L} and \lstinline{QACC_H} respectively. After the access is completed, the value in the register \lstinline{rs1} is incremented by an 8-bit sign-extended constant in the instruction code segment, left-shifted by 4.\\

imm starts at -2048, ends at 2032, with a step of 16.

\textbf{Operation}
\begin{lstlisting}
  dataIn[127:0] = load128(rs1[31:0])
  QACC_L[ 63:  0] = {48{dataIn[15]}, dataIn[ 15:  0]}
  QACC_L[127: 64] = {48{dataIn[31]}, dataIn[ 31: 16]}
  QACC_L[191:128] = {48{dataIn[47]}, dataIn[ 47: 32]}
  QACC_L[255:192] = {48{dataIn[63]}, dataIn[ 63: 48]}

  QACC_H[ 63:  0] = {48{dataIn[79]}, dataIn[ 79: 64]}
  QACC_H[127: 64] = {48{dataIn[95]}, dataIn[ 95: 80]}
  QACC_H[191:128] = {48{dataIn[111]}, dataIn[111: 96]}
  QACC_H[159:192] = {48{dataIn[127]}, dataIn[127:112]}

  rs1[31:0] = rs1[31:0] + imm
\end{lstlisting}

\textbf{Schedule}\\
\begin{longtable}[c]{|l|l|l|}
    \hline
    \rowcolor{lightgray}
    \textbf{Operands} & \textbf{Use} & \textbf{Def} \\\hline
    \endfirsthead
    \hline
    \rowcolor{lightgray}
     \textbf{Operands} & \textbf{Use} & \textbf{Def} \\\hline
    \endhead
    \hline
    \endfoot

  QACC\_H & - & 2 \\\hline
    QACC\_L & - & 2 \\\hline
  rs1 & 1 & 1

\end{longtable}

\textbf{Format}\\

\begin{figure*}[!htbp]
 {\setlength{\tabcolsep}{1pt}
 \begin{center}
\begin{tabular}{@{}V@{}I@{}W@{}F@{}W@{}I@{}F@{}W@{}Y@{}W@{}O}
\instbitrange{31}{27}&
\instbit{26}&
\instbitrange{25}{24}&
\instbitrange{23}{21}&
\instbitrange{20}{19}&
\instbit{18}&
\instbitrange{17}{15}&
\instbitrange{14}{13}&
\instbitrange{12}{9}&
\instbitrange{8}{7}&
\instbitrange{6}{0} \\
\hline
\multicolumn{1}{|c|}{0XX11} &
\multicolumn{1}{c|}{{\scriptsize imm16[7]}} &
\multicolumn{1}{c|}{00} &
\multicolumn{1}{c|}{{\scriptsize imm16[6:4]}} &
\multicolumn{1}{c|}{1X} &
\multicolumn{1}{c|}{{\scriptsize rs1[4]}} &
\multicolumn{1}{c|}{{\scriptsize rs1[2:0]}} &
\multicolumn{1}{c|}{00} &
\multicolumn{1}{c|}{{\scriptsize imm16[3:0]}} &
\multicolumn{1}{c|}{01} &
\multicolumn{1}{c|}{0011111}\\
\hline
5& 1& 2& 3& 2& 1& 3& 2& 4& 2& 7 \\
\end{tabular}
 \end{center}
 }
 \vspace{-0.1in}
 \caption[]{ESP.LDQA.S16.128.IP} \end{figure*}


\newpage
\subsubsection{ESP.LDQA.S16.128.XP}\label{instruction:ESPLDQAS16128XP}
\textbf{Syntax}\\
ESP.LDQA.S16.128.XP rs1,rs2

\textbf{Description}\\
This instruction loads 16-byte data from memory, divides it into 8 segments of 16 bits, sign-extends each segment to 64 bits, and then stores the results in the 256-bit special registers \lstinline{QACC_L} and \lstinline{QACC_H} respectively. After the access is completed, the value in the register \lstinline{rs1} is incremented by the value in the register \lstinline{rs2}.\\
\textbf{Operation}
\begin{lstlisting}
  dataIn[127:0] = load128(rs1[31:0])
  QACC_L[ 63:  0] = {48{dataIn[15]}, dataIn[ 15:  0]}
  QACC_L[127: 64] = {48{dataIn[31]}, dataIn[ 31: 16]}
  QACC_L[191:128] = {48{dataIn[47]}, dataIn[ 47: 32]}
  QACC_L[255:192] = {48{dataIn[63]}, dataIn[ 63: 48]}

  QACC_H[ 63:  0] = {48{dataIn[79]}, dataIn[ 79: 64]}
  QACC_H[127: 64] = {48{dataIn[95]}, dataIn[ 95: 80]}
  QACC_H[191:128] = {48{dataIn[111]}, dataIn[111: 96]}
  QACC_H[159:192] = {48{dataIn[127]}, dataIn[127:112]}

  rs1[31:0] = rs1[31:0] + rs2[31:0]
\end{lstlisting}

\textbf{Schedule}\\
\begin{longtable}[c]{|l|l|l|}
    \hline
    \rowcolor{lightgray}
    \textbf{Operands} & \textbf{Use} & \textbf{Def} \\\hline
    \endfirsthead
    \hline
    \rowcolor{lightgray}
     \textbf{Operands} & \textbf{Use} & \textbf{Def} \\\hline
    \endhead
    \hline
    \endfoot

  QACC\_H & - & 2 \\\hline
    QACC\_L & - & 2 \\\hline
  rs1 & 1 & 1 \\\hline
  rs2 & 1 & -

\end{longtable}

\textbf{Format}\\

\begin{figure*}[!htbp]
 {\setlength{\tabcolsep}{1pt}
 \begin{center}
\begin{tabular}{@{}E@{}I@{}F@{}I@{}I@{}F@{}E@{}O}
\instbitrange{31}{24}&
\instbit{23}&
\instbitrange{22}{20}&
\instbit{19}&
\instbit{18}&
\instbitrange{17}{15}&
\instbitrange{14}{7}&
\instbitrange{6}{0} \\
\hline
\multicolumn{1}{|c|}{1101XXXX} &
\multicolumn{1}{c|}{{\scriptsize rs2[4]}} &
\multicolumn{1}{c|}{{\scriptsize rs2[2:0]}} &
\multicolumn{1}{c|}{X} &
\multicolumn{1}{c|}{{\scriptsize rs1[4]}} &
\multicolumn{1}{c|}{{\scriptsize rs1[2:0]}} &
\multicolumn{1}{c|}{101XXXXX} &
\multicolumn{1}{c|}{0011011}\\
\hline
8& 1& 3& 1& 1& 3& 8& 7 \\
\end{tabular}
 \end{center}
 }
 \vspace{-0.1in}
 \caption[]{ESP.LDQA.S16.128.XP} \end{figure*}


\newpage
\subsubsection{ESP.LDQA.S8.128.IP}\label{instruction:ESPLDQAS8128IP}
\textbf{Syntax}\\
ESP.LDQA.S8.128.IP rs1,imm16

\textbf{Description}\\
This instruction loads 16-byte data from memory, divides it into 16 segments of 8 bits, sign-extends each segment to 32 bits, and then stores the results in the 256-bit special registers \lstinline{QACC_L} and \lstinline{QACC_H} respectively. After the access is completed, the value in the register \lstinline{rs1} is incremented by the 8-bit sign-extended constant in the instruction code segment, left-shifted by 4.\\

imm starts at -2048, ends at 2032, with a step of 16.

\textbf{Operation}
\begin{lstlisting}
DataIn[127:0] = load128(rs1[31:0])
QACC_L[31 :  0] = {24{DataIn[ 7]}, DataIn[7  :  0]}
QACC_L[63 : 32] = {24{DataIn[15]}, DataIn[15 :  8]}
QACC_L[95 : 64] = {24{DataIn[23]}, DataIn[23 : 16]}
...
QACC_L[255:224] = {24{DataIn[63]}, DataIn[63 : 56]}
QACC_H[31 :  0] = {24{DataIn[71]}, DataIn[71 : 64]}
QACC_H[63 : 32] = {24{DataIn[79]}, DataIn[79 : 72]}
QACC_H[95 : 64] = {24{DataIn[87]}, DataIn[87 : 80]}
...
QACC_H[255:224] = {24{DataIn[127]}, DataIn[127:120]}

rs1[31:0] = rs1[31:0] + imm
\end{lstlisting}

\textbf{Schedule}\\
\begin{longtable}[c]{|l|l|l|}
    \hline
    \rowcolor{lightgray}
    \textbf{Operands} & \textbf{Use} & \textbf{Def} \\\hline
    \endfirsthead
    \hline
    \rowcolor{lightgray}
     \textbf{Operands} & \textbf{Use} & \textbf{Def} \\\hline
    \endhead
    \hline
    \endfoot

  QACC\_H & - & 2 \\\hline
    QACC\_L & - & 2 \\\hline
  rs1 & 1 & 1

\end{longtable}

\textbf{Format}\\

\begin{figure*}[!htbp]
 {\setlength{\tabcolsep}{1pt}
 \begin{center}
\begin{tabular}{@{}V@{}I@{}W@{}F@{}W@{}I@{}F@{}W@{}Y@{}W@{}O}
\instbitrange{31}{27}&
\instbit{26}&
\instbitrange{25}{24}&
\instbitrange{23}{21}&
\instbitrange{20}{19}&
\instbit{18}&
\instbitrange{17}{15}&
\instbitrange{14}{13}&
\instbitrange{12}{9}&
\instbitrange{8}{7}&
\instbitrange{6}{0} \\
\hline
\multicolumn{1}{|c|}{0XX01} &
\multicolumn{1}{c|}{{\scriptsize imm16[7]}} &
\multicolumn{1}{c|}{00} &
\multicolumn{1}{c|}{{\scriptsize imm16[6:4]}} &
\multicolumn{1}{c|}{1X} &
\multicolumn{1}{c|}{{\scriptsize rs1[4]}} &
\multicolumn{1}{c|}{{\scriptsize rs1[2:0]}} &
\multicolumn{1}{c|}{00} &
\multicolumn{1}{c|}{{\scriptsize imm16[3:0]}} &
\multicolumn{1}{c|}{01} &
\multicolumn{1}{c|}{0011111}\\
\hline
5& 1& 2& 3& 2& 1& 3& 2& 4& 2& 7 \\
\end{tabular}
 \end{center}
 }
 \vspace{-0.1in}
 \caption[]{ESP.LDQA.S8.128.IP} \end{figure*}


\newpage
\subsubsection{ESP.LDQA.S8.128.XP}\label{instruction:ESPLDQAS8128XP}
\textbf{Syntax}\\
ESP.LDQA.S8.128.XP rs1,rs2

\textbf{Description}\\
This instruction loads 16-byte data from memory, divides it into 16 segments of 8 bits, sign-extends each segment to 32 bits, and then stores the results in the 256-bit special registers \lstinline{QACC_L} and \lstinline{QACC_H} respectively. After the access is completed, the value in the register \lstinline{rs1} is incremented by the value in the register \lstinline{rs2}.\\
\textbf{Operation}
\begin{lstlisting}
DataIn[127:0] = load128(rs1[31:0])
QACC_L[31 :  0] = {24{DataIn[ 7]}, DataIn[7  :  0]}
QACC_L[63 : 32] = {24{DataIn[15]}, DataIn[15 :  8]}
QACC_L[95 : 64] = {24{DataIn[23]}, DataIn[23 : 16]}
...
QACC_L[255:224] = {24{DataIn[63]}, DataIn[63 : 56]}
QACC_H[31 :  0] = {24{DataIn[71]}, DataIn[71 : 64]}
QACC_H[63 : 32] = {24{DataIn[79]}, DataIn[79 : 72]}
QACC_H[95 : 64] = {24{DataIn[87]}, DataIn[87 : 80]}
...
QACC_H[255:224] = {24{DataIn[127]}, DataIn[127:120]}

rs1[31:0] = rs1[31:0] + rs2[31:0]
\end{lstlisting}

\textbf{Schedule}\\
\begin{longtable}[c]{|l|l|l|}
    \hline
    \rowcolor{lightgray}
    \textbf{Operands} & \textbf{Use} & \textbf{Def} \\\hline
    \endfirsthead
    \hline
    \rowcolor{lightgray}
     \textbf{Operands} & \textbf{Use} & \textbf{Def} \\\hline
    \endhead
    \hline
    \endfoot

  QACC\_H & - & 2 \\\hline
    QACC\_L & - & 2 \\\hline
  rs1 & 1 & 1 \\\hline
  rs2 & 1 & -

\end{longtable}

\textbf{Format}\\

\begin{figure*}[!htbp]
 {\setlength{\tabcolsep}{1pt}
 \begin{center}
\begin{tabular}{@{}E@{}I@{}F@{}I@{}I@{}F@{}E@{}O}
\instbitrange{31}{24}&
\instbit{23}&
\instbitrange{22}{20}&
\instbit{19}&
\instbit{18}&
\instbitrange{17}{15}&
\instbitrange{14}{7}&
\instbitrange{6}{0} \\
\hline
\multicolumn{1}{|c|}{0101XXXX} &
\multicolumn{1}{c|}{{\scriptsize rs2[4]}} &
\multicolumn{1}{c|}{{\scriptsize rs2[2:0]}} &
\multicolumn{1}{c|}{X} &
\multicolumn{1}{c|}{{\scriptsize rs1[4]}} &
\multicolumn{1}{c|}{{\scriptsize rs1[2:0]}} &
\multicolumn{1}{c|}{101XXXXX} &
\multicolumn{1}{c|}{0011011}\\
\hline
8& 1& 3& 1& 1& 3& 8& 7 \\
\end{tabular}
 \end{center}
 }
 \vspace{-0.1in}
 \caption[]{ESP.LDQA.S8.128.XP} \end{figure*}


\newpage
\subsubsection{ESP.LDQA.U16.128.IP}\label{instruction:ESPLDQAU16128IP}
\textbf{Syntax}\\
ESP.LDQA.U16.128.IP rs1,imm16

\textbf{Description}\\
This instruction loads 16-byte data from memory, divides it into 8 segments of 16 bits, zero-extends each segment to 64 bits, and then stores the results in the 256-bit special registers \lstinline{QACC_L} and \lstinline{QACC_H} respectively. After the access is completed, the value in the register \lstinline{rs1} is incremented by an 8-bit sign-extended constant in the instruction code segment, left-shifted by 4.\\

imm starts at -2048, ends at 2032, with a step of 16.

\textbf{Operation}
\begin{lstlisting}
  dataIn[127:0] = load128(rs1[31:0])
  QACC_L[ 63:  0] = {48{0}, dataIn[ 15:  0]}
  QACC_L[127: 64] = {48{0}, dataIn[ 31: 16]}
  QACC_L[191:128] = {48{0}, dataIn[ 47: 32]}
  QACC_L[255:192] = {48{0}, dataIn[ 63: 48]}

  QACC_H[ 63:  0] = {48{0}, dataIn[ 79: 64]}
  QACC_H[127: 64] = {48{0}, dataIn[ 95: 80]}
  QACC_H[191:128] = {48{0}, dataIn[111: 96]}
  QACC_H[159:192] = {48{0}, dataIn[127:112]}

  rs1[31:0] = rs1[31:0] + imm
\end{lstlisting}

\textbf{Schedule}\\
\begin{longtable}[c]{|l|l|l|}
    \hline
    \rowcolor{lightgray}
    \textbf{Operands} & \textbf{Use} & \textbf{Def} \\\hline
    \endfirsthead
    \hline
    \rowcolor{lightgray}
     \textbf{Operands} & \textbf{Use} & \textbf{Def} \\\hline
    \endhead
    \hline
    \endfoot

  QACC\_H & - & 2 \\\hline
    QACC\_L & - & 2 \\\hline
  rs1 & 1 & 1

\end{longtable}

\textbf{Format}\\

\begin{figure*}[!htbp]
 {\setlength{\tabcolsep}{1pt}
 \begin{center}
\begin{tabular}{@{}V@{}I@{}W@{}F@{}W@{}I@{}F@{}W@{}Y@{}W@{}O}
\instbitrange{31}{27}&
\instbit{26}&
\instbitrange{25}{24}&
\instbitrange{23}{21}&
\instbitrange{20}{19}&
\instbit{18}&
\instbitrange{17}{15}&
\instbitrange{14}{13}&
\instbitrange{12}{9}&
\instbitrange{8}{7}&
\instbitrange{6}{0} \\
\hline
\multicolumn{1}{|c|}{0XX10} &
\multicolumn{1}{c|}{{\scriptsize imm16[7]}} &
\multicolumn{1}{c|}{00} &
\multicolumn{1}{c|}{{\scriptsize imm16[6:4]}} &
\multicolumn{1}{c|}{1X} &
\multicolumn{1}{c|}{{\scriptsize rs1[4]}} &
\multicolumn{1}{c|}{{\scriptsize rs1[2:0]}} &
\multicolumn{1}{c|}{00} &
\multicolumn{1}{c|}{{\scriptsize imm16[3:0]}} &
\multicolumn{1}{c|}{01} &
\multicolumn{1}{c|}{0011111}\\
\hline
5& 1& 2& 3& 2& 1& 3& 2& 4& 2& 7 \\
\end{tabular}
 \end{center}
 }
 \vspace{-0.1in}
 \caption[]{ESP.LDQA.U16.128.IP} \end{figure*}


\newpage
\subsubsection{ESP.LDQA.U16.128.XP}\label{instruction:ESPLDQAU16128XP}
\textbf{Syntax}\\
ESP.LDQA.U16.128.XP rs1,rs2

\textbf{Description}\\
This instruction loads 16-byte data from memory, divides it into 8 segments of 16 bits, zero-extends each segment to 64 bits, and then stores the results in the 256-bit special registers \lstinline{QACC_L} and \lstinline{QACC_H} respectively. After the access is completed, the value in the register \lstinline{rs1} is incremented by the value in the register \lstinline{rs2}.\\
\textbf{Operation}
\begin{lstlisting}
  dataIn[127:0] = load128(rs1[31:0])
  QACC_L[ 63:  0] = {48{0}, dataIn[ 15:  0]}
  QACC_L[127: 64] = {48{0}, dataIn[ 31: 16]}
  QACC_L[191:128] = {48{0}, dataIn[ 47: 32]}
  QACC_L[255:192] = {48{0}, dataIn[ 63: 48]}

  QACC_H[ 63:  0] = {48{0}, dataIn[ 79: 64]}
  QACC_H[127: 64] = {48{0}, dataIn[ 95: 80]}
  QACC_H[191:128] = {48{0}, dataIn[111: 96]}
  QACC_H[159:192] = {48{0}, dataIn[127:112]}

  rs1[31:0] = rs1[31:0] + rs2[31:0]
\end{lstlisting}

\textbf{Schedule}\\
\begin{longtable}[c]{|l|l|l|}
    \hline
    \rowcolor{lightgray}
    \textbf{Operands} & \textbf{Use} & \textbf{Def} \\\hline
    \endfirsthead
    \hline
    \rowcolor{lightgray}
     \textbf{Operands} & \textbf{Use} & \textbf{Def} \\\hline
    \endhead
    \hline
    \endfoot

  QACC\_H & - & 2 \\\hline
    QACC\_L & - & 2 \\\hline
  rs1 & 1 & 1 \\\hline
  rs2 & 1 & -

\end{longtable}

\textbf{Format}\\

\begin{figure*}[!htbp]
 {\setlength{\tabcolsep}{1pt}
 \begin{center}
\begin{tabular}{@{}E@{}I@{}F@{}I@{}I@{}F@{}E@{}O}
\instbitrange{31}{24}&
\instbit{23}&
\instbitrange{22}{20}&
\instbit{19}&
\instbit{18}&
\instbitrange{17}{15}&
\instbitrange{14}{7}&
\instbitrange{6}{0} \\
\hline
\multicolumn{1}{|c|}{1001XXXX} &
\multicolumn{1}{c|}{{\scriptsize rs2[4]}} &
\multicolumn{1}{c|}{{\scriptsize rs2[2:0]}} &
\multicolumn{1}{c|}{X} &
\multicolumn{1}{c|}{{\scriptsize rs1[4]}} &
\multicolumn{1}{c|}{{\scriptsize rs1[2:0]}} &
\multicolumn{1}{c|}{101XXXXX} &
\multicolumn{1}{c|}{0011011}\\
\hline
8& 1& 3& 1& 1& 3& 8& 7 \\
\end{tabular}
 \end{center}
 }
 \vspace{-0.1in}
 \caption[]{ESP.LDQA.U16.128.XP} \end{figure*}


\newpage
\subsubsection{ESP.LDQA.U8.128.IP}\label{instruction:ESPLDQAU8128IP}
\textbf{Syntax}\\
ESP.LDQA.U8.128.IP rs1,imm16

\textbf{Description}\\
This instruction loads 16-byte data from memory, divides it into 16 segments of 8 bits, zero-extends each segment to 32 bits, and then stores the results in the 256-bit special registers \lstinline{QACC_L} and \lstinline{QACC_H} respectively. After the access is completed, the value in the register \lstinline{rs1} is incremented by the 8-bit sign-extended constant in the instruction code segment, which is left-shifted by 4.\\

imm starts at -2048, ends at 2032, and steps by 16.

\textbf{Operation}
\begin{lstlisting}
DataIn[127:0] = load128(rs1[31:0])
QACC_L[31 :  0] = {24{0}, DataIn[7  :  0]}
QACC_L[63 : 32] = {24{0}, DataIn[15 :  8]}
QACC_L[95 : 64] = {24{0}, DataIn[23 : 16]}
...
QACC_L[255:224] = {24{0}, DataIn[63 : 56]}
QACC_H[31 :  0] = {24{0}, DataIn[71 : 64]}
QACC_H[63 : 32] = {24{0}, DataIn[79 : 72]}
QACC_H[95 : 64] = {24{0}, DataIn[87 : 80]}
...
QACC_H[255:224] = {24{0}, DataIn[127:120]}

rs1[31:0] = rs1[31:0] + imm
\end{lstlisting}

\textbf{Schedule}\\
\begin{longtable}[c]{|l|l|l|}
    \hline
    \rowcolor{lightgray}
    \textbf{Operands} & \textbf{Use} & \textbf{Def} \\\hline
    \endfirsthead
    \hline
    \rowcolor{lightgray}
     \textbf{Operands} & \textbf{Use} & \textbf{Def} \\\hline
    \endhead
    \hline
    \endfoot

  QACC\_H & - & 2 \\\hline
    QACC\_L & - & 2 \\\hline
  rs1 & 1 & 1

\end{longtable}

\textbf{Format}\\

\begin{figure*}[!htbp]
 {\setlength{\tabcolsep}{1pt}
 \begin{center}
\begin{tabular}{@{}V@{}I@{}W@{}F@{}W@{}I@{}F@{}W@{}Y@{}W@{}O}
\instbitrange{31}{27}&
\instbit{26}&
\instbitrange{25}{24}&
\instbitrange{23}{21}&
\instbitrange{20}{19}&
\instbit{18}&
\instbitrange{17}{15}&
\instbitrange{14}{13}&
\instbitrange{12}{9}&
\instbitrange{8}{7}&
\instbitrange{6}{0} \\
\hline
\multicolumn{1}{|c|}{0XX00} &
\multicolumn{1}{c|}{{\scriptsize imm16[7]}} &
\multicolumn{1}{c|}{00} &
\multicolumn{1}{c|}{{\scriptsize imm16[6:4]}} &
\multicolumn{1}{c|}{1X} &
\multicolumn{1}{c|}{{\scriptsize rs1[4]}} &
\multicolumn{1}{c|}{{\scriptsize rs1[2:0]}} &
\multicolumn{1}{c|}{00} &
\multicolumn{1}{c|}{{\scriptsize imm16[3:0]}} &
\multicolumn{1}{c|}{01} &
\multicolumn{1}{c|}{0011111}\\
\hline
5& 1& 2& 3& 2& 1& 3& 2& 4& 2& 7 \\
\end{tabular}
 \end{center}
 }
 \vspace{-0.1in}
 \caption[]{ESP.LDQA.U8.128.IP} \end{figure*}


\newpage
\subsubsection{ESP.LDQA.U8.128.XP}\label{instruction:ESPLDQAU8128XP}
\textbf{Syntax}\\
ESP.LDQA.U8.128.XP rs1,rs2

\textbf{Description}\\
This instruction loads 16-byte data from memory, divides it into 16 segments of 8 bits, zero-extends each segment to 32 bits, and then stores the results in the 256-bit special registers \lstinline{QACC_L} and \lstinline{QACC_H} respectively. After the access is completed, the value in the register \lstinline{rs1} is incremented by the value in the register \lstinline{rs2}.\\
\textbf{Operation}
\begin{lstlisting}
DataIn[127:0] = load128(rs1[31:0])
QACC_L[31 :  0] = {24{0}, DataIn[7  :  0]}
QACC_L[63 : 32] = {24{0}, DataIn[15 :  8]}
QACC_L[95 : 64] = {24{0}, DataIn[23 : 16]}
...
QACC_L[255:224] = {24{0}, DataIn[63 : 56]}
QACC_H[31 :  0] = {24{0}, DataIn[71 : 64]}
QACC_H[63 : 32] = {24{0}, DataIn[79 : 72]}
QACC_H[95 : 64] = {24{0}, DataIn[87 : 80]}
...
QACC_H[255:224] = {24{0}, DataIn[127:120]}

rs1[31:0] = rs1[31:0] + rs2[31:0]
\end{lstlisting}

\textbf{Schedule}\\
\begin{longtable}[c]{|l|l|l|}
    \hline
    \rowcolor{lightgray}
    \textbf{Operands} & \textbf{Use} & \textbf{Def} \\\hline
    \endfirsthead
    \hline
    \rowcolor{lightgray}
     \textbf{Operands} & \textbf{Use} & \textbf{Def} \\\hline
    \endhead
    \hline
    \endfoot

  QACC\_H & - & 2 \\\hline
    QACC\_L & - & 2 \\\hline
  rs1 & 1 & 1 \\\hline
  rs2 & 1 & -

\end{longtable}

\textbf{Format}\\

\begin{figure*}[!htbp]
 {\setlength{\tabcolsep}{1pt}
 \begin{center}
\begin{tabular}{@{}E@{}I@{}F@{}I@{}I@{}F@{}E@{}O}
\instbitrange{31}{24}&
\instbit{23}&
\instbitrange{22}{20}&
\instbit{19}&
\instbit{18}&
\instbitrange{17}{15}&
\instbitrange{14}{7}&
\instbitrange{6}{0} \\
\hline
\multicolumn{1}{|c|}{0001XXXX} &
\multicolumn{1}{c|}{{\scriptsize rs2[4]}} &
\multicolumn{1}{c|}{{\scriptsize rs2[2:0]}} &
\multicolumn{1}{c|}{X} &
\multicolumn{1}{c|}{{\scriptsize rs1[4]}} &
\multicolumn{1}{c|}{{\scriptsize rs1[2:0]}} &
\multicolumn{1}{c|}{101XXXXX} &
\multicolumn{1}{c|}{0011011}\\
\hline
8& 1& 3& 1& 1& 3& 8& 7 \\
\end{tabular}
 \end{center}
 }
 \vspace{-0.1in}
 \caption[]{ESP.LDQA.U8.128.XP} \end{figure*}


\newpage
\subsubsection{ESP.VLDBC.16.IP}\label{instruction:ESPVLDBC16IP}
\textbf{Syntax}\\
ESP.VLDBC.16.IP qu,rs1,imm2

\textbf{Description}\\
These instructions load 16-bit data from memory  and broadcast the data to register \lstinline{qu}, with the pointer updated post-loading by an immediate value.

imm starts at -256, ends at 254, and steps by 2.

\textbf{Operation}\\
\begin{lstlisting}
    qu[127:0] = {8{load16({rs1[31:0]})}}
    rs1[31:0] = rs1[31:0] + imm
\end{lstlisting}

\textbf{Schedule}\\
\begin{longtable}[c]{|l|l|l|}
    \hline
    \rowcolor{lightgray}
    \textbf{Operands} & \textbf{Use} & \textbf{Def} \\\hline
    \endfirsthead
    \hline
    \rowcolor{lightgray}
     \textbf{Operands} & \textbf{Use} & \textbf{Def} \\\hline
    \endhead
    \hline
    \endfoot

  qu & - & 2 \\\hline
  rs1 & 1 & 1

\end{longtable}

\textbf{Format}\\

\begin{figure*}[!htbp]
 {\setlength{\tabcolsep}{1pt}
 \begin{center}
\begin{tabular}{@{}S@{}O@{}I@{}F@{}W@{}F@{}W@{}I@{}O}
\instbitrange{31}{26}&
\instbitrange{25}{19}&
\instbit{18}&
\instbitrange{17}{15}&
\instbitrange{14}{13}&
\instbitrange{12}{10}&
\instbitrange{9}{8}&
\instbit{7}&
\instbitrange{6}{0} \\
\hline
\multicolumn{1}{|c|}{{\scriptsize imm2[7:2]}} &
\multicolumn{1}{c|}{111001X} &
\multicolumn{1}{c|}{{\scriptsize rs1[4]}} &
\multicolumn{1}{c|}{{\scriptsize rs1[2:0]}} &
\multicolumn{1}{c|}{00} &
\multicolumn{1}{c|}{{\scriptsize qu[2:0]}} &
\multicolumn{1}{c|}{{\scriptsize imm2[1:0]}} &
\multicolumn{1}{c|}{X} &
\multicolumn{1}{c|}{0011111}\\
\hline
6& 7& 1& 3& 2& 3& 2& 1& 7 \\
\end{tabular}
 \end{center}
 }
 \vspace{-0.1in}
 \caption[]{ESP.VLDBC.16.IP} \end{figure*}


\newpage
\subsubsection{ESP.VLDBC.16.XP}\label{instruction:ESPVLDBC16XP}
\textbf{Syntax}\\
ESP.VLDBC.16.XP qu,rs1,rs2

\textbf{Description}\\
These instructions load 16-bit data from memory and broadcast to the register \lstinline{qu}, with the pointer updated post-loading by the value in the register \lstinline{rs2}.

\textbf{Operation}\\
\begin{lstlisting}
    qu[127:0] = {8{load16({rs1[31:0]})}}
    rs1[31:0] = rs1[31:0] + rs2[31:0]
\end{lstlisting}

\textbf{Schedule}\\
\begin{longtable}[c]{|l|l|l|}
    \hline
    \rowcolor{lightgray}
    \textbf{Operands} & \textbf{Use} & \textbf{Def} \\\hline
    \endfirsthead
    \hline
    \rowcolor{lightgray}
     \textbf{Operands} & \textbf{Use} & \textbf{Def} \\\hline
    \endhead
    \hline
    \endfoot

  qu & - & 2 \\\hline
  rs1 & 1 & 1 \\\hline
    rs2 & 1 & -

\end{longtable}

\textbf{Format}\\

\begin{figure*}[!htbp]
 {\setlength{\tabcolsep}{1pt}
 \begin{center}
\begin{tabular}{@{}E@{}I@{}F@{}I@{}I@{}F@{}W@{}F@{}F@{}O}
\instbitrange{31}{24}&
\instbit{23}&
\instbitrange{22}{20}&
\instbit{19}&
\instbit{18}&
\instbitrange{17}{15}&
\instbitrange{14}{13}&
\instbitrange{12}{10}&
\instbitrange{9}{7}&
\instbitrange{6}{0} \\
\hline
\multicolumn{1}{|c|}{01X0XX10} &
\multicolumn{1}{c|}{{\scriptsize rs2[4]}} &
\multicolumn{1}{c|}{{\scriptsize rs2[2:0]}} &
\multicolumn{1}{c|}{1} &
\multicolumn{1}{c|}{{\scriptsize rs1[4]}} &
\multicolumn{1}{c|}{{\scriptsize rs1[2:0]}} &
\multicolumn{1}{c|}{11} &
\multicolumn{1}{c|}{{\scriptsize qu[2:0]}} &
\multicolumn{1}{c|}{1X0} &
\multicolumn{1}{c|}{0011011}\\
\hline
8& 1& 3& 1& 1& 3& 2& 3& 3& 7 \\
\end{tabular}
 \end{center}
 }
 \vspace{-0.1in}
 \caption[]{ESP.VLDBC.16.XP} \end{figure*}


\newpage
\subsubsection{ESP.VLDBC.32.IP}\label{instruction:ESPVLDBC32IP}
\textbf{Syntax}\\
ESP.VLDBC.32.IP qu,rs1,imm4

\textbf{Description}\\
These instructions load 32-bit data from memory and broadcast to the register \lstinline{qu}, with the pointer post updated by an immediate value.

imm starts at -512, ends at 508, and steps by 4.

\textbf{Operation}\\
\begin{lstlisting}
    qu[127:0] = {4{load32({rs1[31:0]})}}
    rs1[31:0] = rs1[31:0] + imm
\end{lstlisting}

\textbf{Schedule}\\
\begin{longtable}[c]{|l|l|l|}
    \hline
    \rowcolor{lightgray}
    \textbf{Operands} & \textbf{Use} & \textbf{Def} \\\hline
    \endfirsthead
    \hline
    \rowcolor{lightgray}
     \textbf{Operands} & \textbf{Use} & \textbf{Def} \\\hline
    \endhead
    \hline
    \endfoot

  qu & - & 2 \\\hline
  rs1 & 1 & 1

\end{longtable}

\textbf{Format}\\

\begin{figure*}[!htbp]
 {\setlength{\tabcolsep}{1pt}
 \begin{center}
\begin{tabular}{@{}S@{}O@{}I@{}F@{}W@{}F@{}W@{}I@{}O}
\instbitrange{31}{26}&
\instbitrange{25}{19}&
\instbit{18}&
\instbitrange{17}{15}&
\instbitrange{14}{13}&
\instbitrange{12}{10}&
\instbitrange{9}{8}&
\instbit{7}&
\instbitrange{6}{0} \\
\hline
\multicolumn{1}{|c|}{{\scriptsize imm4[7:2]}} &
\multicolumn{1}{c|}{110101X} &
\multicolumn{1}{c|}{{\scriptsize rs1[4]}} &
\multicolumn{1}{c|}{{\scriptsize rs1[2:0]}} &
\multicolumn{1}{c|}{00} &
\multicolumn{1}{c|}{{\scriptsize qu[2:0]}} &
\multicolumn{1}{c|}{{\scriptsize imm4[1:0]}} &
\multicolumn{1}{c|}{0} &
\multicolumn{1}{c|}{0011111}\\
\hline
6& 7& 1& 3& 2& 3& 2& 1& 7 \\
\end{tabular}
 \end{center}
 }
 \vspace{-0.1in}
 \caption[]{ESP.VLDBC.32.IP} \end{figure*}


\newpage
\subsubsection{ESP.VLDBC.32.XP}\label{instruction:ESPVLDBC32XP}
\textbf{Syntax}\\
ESP.VLDBC.32.XP qu,rs1,rs2

\textbf{Description}\\
These instructions load 32 bits of data from memory and broadcast to the register \lstinline{qu}, with the pointer updated post-loading by the value in the register \lstinline{rs2}.

\textbf{Operation}\\
\begin{lstlisting}
    qu[127:0] = {4{load32({rs1[31:0]})}}
    rs1[31:0] = rs1[31:0] + rs2[31:0]
\end{lstlisting}

\textbf{Schedule}\\
\begin{longtable}[c]{|l|l|l|}
    \hline
    \rowcolor{lightgray}
    \textbf{Operands} & \textbf{Use} & \textbf{Def} \\\hline
    \endfirsthead
    \hline
    \rowcolor{lightgray}
     \textbf{Operands} & \textbf{Use} & \textbf{Def} \\\hline
    \endhead
    \hline
    \endfoot

  qu & - & 2 \\\hline
  rs1 & 1 & 1 \\\hline
    rs2 & 1 & -

\end{longtable}

\textbf{Format}\\

\begin{figure*}[!htbp]
 {\setlength{\tabcolsep}{1pt}
 \begin{center}
\begin{tabular}{@{}E@{}I@{}F@{}I@{}I@{}F@{}W@{}F@{}F@{}O}
\instbitrange{31}{24}&
\instbit{23}&
\instbitrange{22}{20}&
\instbit{19}&
\instbit{18}&
\instbitrange{17}{15}&
\instbitrange{14}{13}&
\instbitrange{12}{10}&
\instbitrange{9}{7}&
\instbitrange{6}{0} \\
\hline
\multicolumn{1}{|c|}{01X0XXX1} &
\multicolumn{1}{c|}{{\scriptsize rs2[4]}} &
\multicolumn{1}{c|}{{\scriptsize rs2[2:0]}} &
\multicolumn{1}{c|}{1} &
\multicolumn{1}{c|}{{\scriptsize rs1[4]}} &
\multicolumn{1}{c|}{{\scriptsize rs1[2:0]}} &
\multicolumn{1}{c|}{11} &
\multicolumn{1}{c|}{{\scriptsize qu[2:0]}} &
\multicolumn{1}{c|}{1X0} &
\multicolumn{1}{c|}{0011011}\\
\hline
8& 1& 3& 1& 1& 3& 2& 3& 3& 7 \\
\end{tabular}
 \end{center}
 }
 \vspace{-0.1in}
 \caption[]{ESP.VLDBC.32.XP} \end{figure*}


\newpage
\subsubsection{ESP.VLDBC.8.IP}\label{instruction:ESPVLDBC8IP}
\textbf{Syntax}\\
ESP.VLDBC.8.IP qu,rs1,imm1

\textbf{Description}\\
These instructions load 8 bits of data from memory and broadcast to the register \lstinline{qu} with a pointer post updated by an immediate value.

imm starts at -128, ends at 127, and the step is 1.

\textbf{Operation}\\
\begin{lstlisting}
    qu[127:0] = {16{load8({rs1[31:0]})}}
    rs1[31:0] = rs1[31:0] + imm
\end{lstlisting}

\textbf{Schedule}\\
\begin{longtable}[c]{|l|l|l|}
    \hline
    \rowcolor{lightgray}
    \textbf{Operands} & \textbf{Use} & \textbf{Def} \\\hline
    \endfirsthead
    \hline
    \rowcolor{lightgray}
     \textbf{Operands} & \textbf{Use} & \textbf{Def} \\\hline
    \endhead
    \hline
    \endfoot

  qu & - & 2 \\\hline
  rs1 & 1 & 1

\end{longtable}

\textbf{Format}\\

\begin{figure*}[!htbp]
 {\setlength{\tabcolsep}{1pt}
 \begin{center}
\begin{tabular}{@{}S@{}O@{}I@{}F@{}W@{}F@{}W@{}I@{}O}
\instbitrange{31}{26}&
\instbitrange{25}{19}&
\instbit{18}&
\instbitrange{17}{15}&
\instbitrange{14}{13}&
\instbitrange{12}{10}&
\instbitrange{9}{8}&
\instbit{7}&
\instbitrange{6}{0} \\
\hline
\multicolumn{1}{|c|}{{\scriptsize imm1[7:2]}} &
\multicolumn{1}{c|}{110101X} &
\multicolumn{1}{c|}{{\scriptsize rs1[4]}} &
\multicolumn{1}{c|}{{\scriptsize rs1[2:0]}} &
\multicolumn{1}{c|}{00} &
\multicolumn{1}{c|}{{\scriptsize qu[2:0]}} &
\multicolumn{1}{c|}{{\scriptsize imm1[1:0]}} &
\multicolumn{1}{c|}{1} &
\multicolumn{1}{c|}{0011111}\\
\hline
6& 7& 1& 3& 2& 3& 2& 1& 7 \\
\end{tabular}
 \end{center}
 }
 \vspace{-0.1in}
 \caption[]{ESP.VLDBC.8.IP} \end{figure*}


\newpage
\subsubsection{ESP.VLDBC.8.XP}\label{instruction:ESPVLDBC8XP}
\textbf{Syntax}\\
ESP.VLDBC.8.XP qu,rs1,rs2

\textbf{Description}\\
These instructions load 8 bits of data from memory and broadcast to the register \lstinline{qu}, with the pointer updated post-loading by the value in the register \lstinline{rs2}.

\textbf{Operation}\\
\begin{lstlisting}
    qu[127:0] = {16{load8({rs1[31:0]})}}
    rs1[31:0] = rs1[31:0] + rs2[31:0]
\end{lstlisting}

\textbf{Schedule}\\
\begin{longtable}[c]{|l|l|l|}
    \hline
    \rowcolor{lightgray}
    \textbf{Operands} & \textbf{Use} & \textbf{Def} \\\hline
    \endfirsthead
    \hline
    \rowcolor{lightgray}
     \textbf{Operands} & \textbf{Use} & \textbf{Def} \\\hline
    \endhead
    \hline
    \endfoot

  qu & - & 2 \\\hline
  rs1 & 1 & 1 \\\hline
    rs2 & 1 & -

\end{longtable}

\textbf{Format}\\

\begin{figure*}[!htbp]
 {\setlength{\tabcolsep}{1pt}
 \begin{center}
\begin{tabular}{@{}E@{}I@{}F@{}I@{}I@{}F@{}W@{}F@{}F@{}O}
\instbitrange{31}{24}&
\instbit{23}&
\instbitrange{22}{20}&
\instbit{19}&
\instbit{18}&
\instbitrange{17}{15}&
\instbitrange{14}{13}&
\instbitrange{12}{10}&
\instbitrange{9}{7}&
\instbitrange{6}{0} \\
\hline
\multicolumn{1}{|c|}{01X0XX00} &
\multicolumn{1}{c|}{{\scriptsize rs2[4]}} &
\multicolumn{1}{c|}{{\scriptsize rs2[2:0]}} &
\multicolumn{1}{c|}{1} &
\multicolumn{1}{c|}{{\scriptsize rs1[4]}} &
\multicolumn{1}{c|}{{\scriptsize rs1[2:0]}} &
\multicolumn{1}{c|}{11} &
\multicolumn{1}{c|}{{\scriptsize qu[2:0]}} &
\multicolumn{1}{c|}{1X0} &
\multicolumn{1}{c|}{0011011}\\
\hline
8& 1& 3& 1& 1& 3& 2& 3& 3& 7 \\
\end{tabular}
 \end{center}
 }
 \vspace{-0.1in}
 \caption[]{ESP.VLDBC.8.XP} \end{figure*}


\newpage
\subsubsection{ESP.VLDEXT.S16.IP}\label{instruction:ESPVLDEXTS16IP}
\textbf{Syntax}\\
ESP.VLDEXT.S16.IP qu,qz,rs1,imm16

\textbf{Description}\\
This instruction loads 128 bit of data from memory, sign-extends each 16-bit segment to 32 bits, with pointer post-incremented by an immediate value.

imm is start at -2048, end at 2032, step is 16.

\textbf{Operation}\\
\begin{lstlisting}
    DataIn[127:0] = load128(rs1[31:0])
    qu[ 31:  0] = { {16{DataIn[ 15]}}, DataIn[ 15:  0] }
    qu[ 63: 32] = { {16{DataIn[ 31]}}, DataIn[ 31: 16] }
    qu[ 95: 64] = { {16{DataIn[ 47]}}, DataIn[ 47: 32] }
    qu[127: 96] = { {16{DataIn[ 63]}}, DataIn[ 63: 48] }
    qz[ 31:  0] = { {16{DataIn[ 79]}}, DataIn[ 79: 64] }
    qz[ 63: 32] = { {16{DataIn[ 95]}}, DataIn[ 95: 80] }
    qz[ 95: 64] = { {16{DataIn[111]}}, DataIn[111: 96] }
    qz[127: 96] = { {16{DataIn[127]}}, DataIn[127:112] }
    rs1[31:0] = rs1[31:0] + imm
\end{lstlisting}

\textbf{Schedule}\\
\begin{longtable}[c]{|l|l|l|}
    \hline
    \rowcolor{lightgray}
    \textbf{Operands} & \textbf{Use} & \textbf{Def} \\\hline
    \endfirsthead
    \hline
    \rowcolor{lightgray}
     \textbf{Operands} & \textbf{Use} & \textbf{Def} \\\hline
    \endhead
    \hline
    \endfoot

  qu & - & 2 \\\hline
  qv & - & 2 \\\hline
  rs1 & 1 & 1

\end{longtable}

\textbf{Format}\\

\begin{figure*}[!htbp]
 {\setlength{\tabcolsep}{1pt}
 \begin{center}
\begin{tabular}{@{}V@{}I@{}W@{}F@{}W@{}I@{}F@{}W@{}F@{}F@{}O}
\instbitrange{31}{27}&
\instbit{26}&
\instbitrange{25}{24}&
\instbitrange{23}{21}&
\instbitrange{20}{19}&
\instbit{18}&
\instbitrange{17}{15}&
\instbitrange{14}{13}&
\instbitrange{12}{10}&
\instbitrange{9}{7}&
\instbitrange{6}{0} \\
\hline
\multicolumn{1}{|c|}{1XX11} &
\multicolumn{1}{c|}{{\scriptsize imm16[3]}} &
\multicolumn{1}{c|}{00} &
\multicolumn{1}{c|}{{\scriptsize imm16[2:0]}} &
\multicolumn{1}{c|}{0X} &
\multicolumn{1}{c|}{{\scriptsize rs1[4]}} &
\multicolumn{1}{c|}{{\scriptsize rs1[2:0]}} &
\multicolumn{1}{c|}{00} &
\multicolumn{1}{c|}{{\scriptsize qu[2:0]}} &
\multicolumn{1}{c|}{{\scriptsize qz[2:0]}} &
\multicolumn{1}{c|}{0011111}\\
\hline
5& 1& 2& 3& 2& 1& 3& 2& 3& 3& 7 \\
\end{tabular}
 \end{center}
 }
 \vspace{-0.1in}
 \caption[]{ESP.VLDEXT.S16.IP} \end{figure*}


\newpage
\subsubsection{ESP.VLDEXT.S16.XP}\label{instruction:ESPVLDEXTS16XP}
\textbf{Syntax}\\
ESP.VLDEXT.S16.XP qu,qz,rs1,rs2

\textbf{Description}\\
This instruction loads 128-bit data from memory and sign-extends the data from a 16-bit segment to a 32-bit segment.

\textbf{Operation}\\
\begin{lstlisting}
    DataIn[127:0] = load128(rs1[31:0])
    qu[ 31:  0] = { {16{DataIn[ 15]}}, DataIn[ 15:  0] }
    qu[ 63: 32] = { {16{DataIn[ 31]}}, DataIn[ 31: 16] }
    qu[ 95: 64] = { {16{DataIn[ 47]}}, DataIn[ 47: 32] }
    qu[127: 96] = { {16{DataIn[ 63]}}, DataIn[ 63: 48] }
    qz[ 31:  0] = { {16{DataIn[ 79]}}, DataIn[ 79: 64] }
    qz[ 63: 32] = { {16{DataIn[ 95]}}, DataIn[ 95: 80] }
    qz[ 95: 64] = { {16{DataIn[111]}}, DataIn[111: 96] }
    qz[127: 96] = { {16{DataIn[127]}}, DataIn[127:112] }
    rs1[31:0] = rs1[31:0] + rs2[31:0]
\end{lstlisting}

\textbf{Schedule}\\
\begin{longtable}[c]{|l|l|l|}
    \hline
    \rowcolor{lightgray}
    \textbf{Operands} & \textbf{Use} & \textbf{Def} \\\hline
    \endfirsthead
    \hline
    \rowcolor{lightgray}
     \textbf{Operands} & \textbf{Use} & \textbf{Def} \\\hline
    \endhead
    \hline
    \endfoot

  qu & - & 2 \\\hline
  qv & - & 2 \\\hline
  rs1 & 1 & 1 \\\hline
    rs2 & 1 & -

\end{longtable}

\textbf{Format}\\

\begin{figure*}[!htbp]
 {\setlength{\tabcolsep}{1pt}
 \begin{center}
\begin{tabular}{@{}E@{}I@{}F@{}I@{}I@{}F@{}W@{}F@{}F@{}O}
\instbitrange{31}{24}&
\instbit{23}&
\instbitrange{22}{20}&
\instbit{19}&
\instbit{18}&
\instbitrange{17}{15}&
\instbitrange{14}{13}&
\instbitrange{12}{10}&
\instbitrange{9}{7}&
\instbitrange{6}{0} \\
\hline
\multicolumn{1}{|c|}{01X1XX11} &
\multicolumn{1}{c|}{{\scriptsize rs2[4]}} &
\multicolumn{1}{c|}{{\scriptsize rs2[2:0]}} &
\multicolumn{1}{c|}{X} &
\multicolumn{1}{c|}{{\scriptsize rs1[4]}} &
\multicolumn{1}{c|}{{\scriptsize rs1[2:0]}} &
\multicolumn{1}{c|}{11} &
\multicolumn{1}{c|}{{\scriptsize qu[2:0]}} &
\multicolumn{1}{c|}{{\scriptsize qz[2:0]}} &
\multicolumn{1}{c|}{0011011}\\
\hline
8& 1& 3& 1& 1& 3& 2& 3& 3& 7 \\
\end{tabular}
 \end{center}
 }
 \vspace{-0.1in}
 \caption[]{ESP.VLDEXT.S16.XP} \end{figure*}


\newpage
\subsubsection{ESP.VLDEXT.S8.IP}\label{instruction:ESPVLDEXTS8IP}
\textbf{Syntax}\\
ESP.VLDEXT.S8.IP qu,qz,rs1,imm16

\textbf{Description}\\
This instruction loads 128 bit data from memory, and sign-extend the data as 8-bit segment to 16-bit, with pointer post update by immediate value.

imm starts at -2048, ends at 2032, with a step of step 16.

\textbf{Operation}\\
\begin{lstlisting}
    DataIn[127:0] = load128(rs1[31:0])
    qu[ 15:  0] = { {8{DataIn[  7]}}, DataIn[  7:  0] }
    qu[ 31: 16] = { {8{DataIn[ 15]}}, DataIn[ 15:  8] }
    qu[ 47: 32] = { {8{DataIn[ 23]}}, DataIn[ 23: 16] }
    qu[ 63: 48] = { {8{DataIn[ 31]}}, DataIn[ 31: 24] }
    qu[ 79: 64] = { {8{DataIn[ 39]}}, DataIn[ 39: 32] }
    qu[ 95: 80] = { {8{DataIn[ 47]}}, DataIn[ 47: 40] }
    qu[111: 96] = { {8{DataIn[ 55]}}, DataIn[ 55: 48] }
    qu[127:112] = { {8{DataIn[ 63]}}, DataIn[ 63: 56] }
    qz[ 15:  0] = { {8{DataIn[ 71]}}, DataIn[ 71: 64] }
    qz[ 31: 16] = { {8{DataIn[ 79]}}, DataIn[ 79: 72] }
    qz[ 47: 32] = { {8{DataIn[ 87]}}, DataIn[ 87: 80] }
    qz[ 63: 48] = { {8{DataIn[ 95]}}, DataIn[ 95: 88] }
    qz[ 79: 64] = { {8{DataIn[103]}}, DataIn[103: 96] }
    qz[ 95: 80] = { {8{DataIn[111]}}, DataIn[111:104] }
    qz[111: 96] = { {8{DataIn[119]}}, DataIn[119:112] }
    qz[127:112] = { {8{DataIn[127]}}, DataIn[127:120] }
    rs1[31:0] = rs1[31:0] + imm
\end{lstlisting}

\textbf{Schedule}\\
\begin{longtable}[c]{|l|l|l|}
    \hline
    \rowcolor{lightgray}
    \textbf{Operands} & \textbf{Use} & \textbf{Def} \\\hline
    \endfirsthead
    \hline
    \rowcolor{lightgray}
     \textbf{Operands} & \textbf{Use} & \textbf{Def} \\\hline
    \endhead
    \hline
    \endfoot

  qu & - & 2 \\\hline
  qv & - & 2 \\\hline
  rs1 & 1 & 1

\end{longtable}

\textbf{Format}\\

\begin{figure*}[!htbp]
 {\setlength{\tabcolsep}{1pt}
 \begin{center}
\begin{tabular}{@{}V@{}I@{}W@{}F@{}W@{}I@{}F@{}W@{}F@{}F@{}O}
\instbitrange{31}{27}&
\instbit{26}&
\instbitrange{25}{24}&
\instbitrange{23}{21}&
\instbitrange{20}{19}&
\instbit{18}&
\instbitrange{17}{15}&
\instbitrange{14}{13}&
\instbitrange{12}{10}&
\instbitrange{9}{7}&
\instbitrange{6}{0} \\
\hline
\multicolumn{1}{|c|}{1XX01} &
\multicolumn{1}{c|}{{\scriptsize imm16[3]}} &
\multicolumn{1}{c|}{00} &
\multicolumn{1}{c|}{{\scriptsize imm16[2:0]}} &
\multicolumn{1}{c|}{0X} &
\multicolumn{1}{c|}{{\scriptsize rs1[4]}} &
\multicolumn{1}{c|}{{\scriptsize rs1[2:0]}} &
\multicolumn{1}{c|}{00} &
\multicolumn{1}{c|}{{\scriptsize qu[2:0]}} &
\multicolumn{1}{c|}{{\scriptsize qz[2:0]}} &
\multicolumn{1}{c|}{0011111}\\
\hline
5& 1& 2& 3& 2& 1& 3& 2& 3& 3& 7 \\
\end{tabular}
 \end{center}
 }
 \vspace{-0.1in}
 \caption[]{ESP.VLDEXT.S8.IP} \end{figure*}


\newpage
\subsubsection{ESP.VLDEXT.S8.XP}\label{instruction:ESPVLDEXTS8XP}
\textbf{Syntax}\\
ESP.VLDEXT.S8.XP qu,qz,rs1,rs2

\textbf{Description}\\
This instruction loads 128-bit data from memory and sign-extends the data from an 8-bit segment to a 16-bit segment.

\textbf{Operation}\\
\begin{lstlisting}
    DataIn[127:0] = load128(rs1[31:0])
    qu[ 15:  0] = { {8{DataIn[  7]}}, DataIn[  7:  0] }
    qu[ 31: 16] = { {8{DataIn[ 15]}}, DataIn[ 15:  8] }
    qu[ 47: 32] = { {8{DataIn[ 23]}}, DataIn[ 23: 16] }
    qu[ 63: 48] = { {8{DataIn[ 31]}}, DataIn[ 31: 24] }
    qu[ 79: 64] = { {8{DataIn[ 39]}}, DataIn[ 39: 32] }
    qu[ 95: 80] = { {8{DataIn[ 47]}}, DataIn[ 47: 40] }
    qu[111: 96] = { {8{DataIn[ 55]}}, DataIn[ 55: 48] }
    qu[127:112] = { {8{DataIn[ 63]}}, DataIn[ 63: 56] }
    qv[ 15:  0] = { {8{DataIn[ 71]}}, DataIn[ 71: 64] }
    qv[ 31: 16] = { {8{DataIn[ 79]}}, DataIn[ 79: 72] }
    qv[ 47: 32] = { {8{DataIn[ 87]}}, DataIn[ 87: 80] }
    qv[ 63: 48] = { {8{DataIn[ 95]}}, DataIn[ 95: 88] }
    qv[ 79: 64] = { {8{DataIn[103]}}, DataIn[103: 96] }
    qv[ 95: 80] = { {8{DataIn[111]}}, DataIn[111:104] }
    qv[111: 96] = { {8{DataIn[119]}}, DataIn[119:112] }
    qv[127:112] = { {8{DataIn[127]}}, DataIn[127:120] }
    rs1[31:0] = rs1[31:0] + rs2[31:0]
\end{lstlisting}

\textbf{Schedule}\\
\begin{longtable}[c]{|l|l|l|}
    \hline
    \rowcolor{lightgray}
    \textbf{Operands} & \textbf{Use} & \textbf{Def} \\\hline
    \endfirsthead
    \hline
    \rowcolor{lightgray}
     \textbf{Operands} & \textbf{Use} & \textbf{Def} \\\hline
    \endhead
    \hline
    \endfoot

  qu & - & 2 \\\hline
  qv & - & 2 \\\hline
  rs1 & 1 & 1 \\\hline
    rs2 & 1 & -

\end{longtable}

\textbf{Format}\\

\begin{figure*}[!htbp]
 {\setlength{\tabcolsep}{1pt}
 \begin{center}
\begin{tabular}{@{}E@{}I@{}F@{}I@{}I@{}F@{}W@{}F@{}F@{}O}
\instbitrange{31}{24}&
\instbit{23}&
\instbitrange{22}{20}&
\instbit{19}&
\instbit{18}&
\instbitrange{17}{15}&
\instbitrange{14}{13}&
\instbitrange{12}{10}&
\instbitrange{9}{7}&
\instbitrange{6}{0} \\
\hline
\multicolumn{1}{|c|}{01X1XX01} &
\multicolumn{1}{c|}{{\scriptsize rs2[4]}} &
\multicolumn{1}{c|}{{\scriptsize rs2[2:0]}} &
\multicolumn{1}{c|}{X} &
\multicolumn{1}{c|}{{\scriptsize rs1[4]}} &
\multicolumn{1}{c|}{{\scriptsize rs1[2:0]}} &
\multicolumn{1}{c|}{11} &
\multicolumn{1}{c|}{{\scriptsize qu[2:0]}} &
\multicolumn{1}{c|}{{\scriptsize qz[2:0]}} &
\multicolumn{1}{c|}{0011011}\\
\hline
8& 1& 3& 1& 1& 3& 2& 3& 3& 7 \\
\end{tabular}
 \end{center}
 }
 \vspace{-0.1in}
 \caption[]{ESP.VLDEXT.S8.XP} \end{figure*}


\newpage
\subsubsection{ESP.VLDEXT.U16.IP}\label{instruction:ESPVLDEXTU16IP}
\textbf{Syntax}\\
ESP.VLDEXT.U16.IP qu,qz,rs1,imm16

\textbf{Description}\\
This instruction loads 128-bit data from memory and zero-extends the data from a 16-bit segment to a 32-bit segment.

\lstinline{imm} starts at -2048, ends at 2032, and the step is 16.

\textbf{Operation}\\
\begin{lstlisting}
    DataIn[127:0] = load128(rs1[31:0])
    qu[ 31:  0] = { {16{0}}, DataIn[ 15:  0] }
    qu[ 63: 32] = { {16{0}}, DataIn[ 31: 16] }
    qu[ 95: 64] = { {16{0}}, DataIn[ 47: 32] }
    qu[127: 96] = { {16{0}}, DataIn[ 63: 48] }
    qz[ 31:  0] = { {16{0}}, DataIn[ 79: 64] }
    qz[ 63: 32] = { {16{0}}, DataIn[ 95: 80] }
    qz[ 95: 64] = { {16{0}}, DataIn[111: 96] }
    qz[127: 96] = { {16{0}}, DataIn[127:112] }
    rs1[31:0] = rs1[31:0] + imm
\end{lstlisting}

\textbf{Schedule}\\
\begin{longtable}[c]{|l|l|l|}
    \hline
    \rowcolor{lightgray}
    \textbf{Operands} & \textbf{Use} & \textbf{Def} \\\hline
    \endfirsthead
    \hline
    \rowcolor{lightgray}
     \textbf{Operands} & \textbf{Use} & \textbf{Def} \\\hline
    \endhead
    \hline
    \endfoot

  qu & - & 2 \\\hline
  qv & - & 2 \\\hline
  rs1 & 1 & 1

\end{longtable}

\textbf{Format}\\

\begin{figure*}[!htbp]
 {\setlength{\tabcolsep}{1pt}
 \begin{center}
\begin{tabular}{@{}V@{}I@{}W@{}F@{}W@{}I@{}F@{}W@{}F@{}F@{}O}
\instbitrange{31}{27}&
\instbit{26}&
\instbitrange{25}{24}&
\instbitrange{23}{21}&
\instbitrange{20}{19}&
\instbit{18}&
\instbitrange{17}{15}&
\instbitrange{14}{13}&
\instbitrange{12}{10}&
\instbitrange{9}{7}&
\instbitrange{6}{0} \\
\hline
\multicolumn{1}{|c|}{1XX10} &
\multicolumn{1}{c|}{{\scriptsize imm16[3]}} &
\multicolumn{1}{c|}{00} &
\multicolumn{1}{c|}{{\scriptsize imm16[2:0]}} &
\multicolumn{1}{c|}{0X} &
\multicolumn{1}{c|}{{\scriptsize rs1[4]}} &
\multicolumn{1}{c|}{{\scriptsize rs1[2:0]}} &
\multicolumn{1}{c|}{00} &
\multicolumn{1}{c|}{{\scriptsize qu[2:0]}} &
\multicolumn{1}{c|}{{\scriptsize qz[2:0]}} &
\multicolumn{1}{c|}{0011111}\\
\hline
5& 1& 2& 3& 2& 1& 3& 2& 3& 3& 7 \\
\end{tabular}
 \end{center}
 }
 \vspace{-0.1in}
 \caption[]{ESP.VLDEXT.U16.IP} \end{figure*}


\newpage
\subsubsection{ESP.VLDEXT.U16.XP}\label{instruction:ESPVLDEXTU16XP}
\textbf{Syntax}\\
ESP.VLDEXT.U16.XP qu,qz,rs1,rs2

\textbf{Description}\\
This instruction loads 128-bit data from memory and zero-extends the data from a 16-bit segment to a 32-bit segment.

\textbf{Operation}\\
\begin{lstlisting}
    DataIn[127:0] = load128(rs1[31:0])
    qu[ 31:  0] = { {16{0}}, DataIn[ 15:  0] }
    qu[ 63: 32] = { {16{0}}, DataIn[ 31: 16] }
    qu[ 95: 64] = { {16{0}}, DataIn[ 47: 32] }
    qu[127: 96] = { {16{0}}, DataIn[ 63: 48] }
    qz[ 31:  0] = { {16{0}}, DataIn[ 79: 64] }
    qz[ 63: 32] = { {16{0}}, DataIn[ 95: 80] }
    qz[ 95: 64] = { {16{0}}, DataIn[111: 96] }
    qz[127: 96] = { {16{0}}, DataIn[127:112] }
    rs1[31:0] = rs1[31:0] + rs2[31:0]
\end{lstlisting}

\textbf{Schedule}\\
\begin{longtable}[c]{|l|l|l|}
    \hline
    \rowcolor{lightgray}
    \textbf{Operands} & \textbf{Use} & \textbf{Def} \\\hline
    \endfirsthead
    \hline
    \rowcolor{lightgray}
     \textbf{Operands} & \textbf{Use} & \textbf{Def} \\\hline
    \endhead
    \hline
    \endfoot

  qu & - & 2 \\\hline
  qv & - & 2 \\\hline
  rs1 & 1 & 1 \\\hline
    rs2 & 1 & -

\end{longtable}

\textbf{Format}\\

\begin{figure*}[!htbp]
 {\setlength{\tabcolsep}{1pt}
 \begin{center}
\begin{tabular}{@{}E@{}I@{}F@{}I@{}I@{}F@{}W@{}F@{}F@{}O}
\instbitrange{31}{24}&
\instbit{23}&
\instbitrange{22}{20}&
\instbit{19}&
\instbit{18}&
\instbitrange{17}{15}&
\instbitrange{14}{13}&
\instbitrange{12}{10}&
\instbitrange{9}{7}&
\instbitrange{6}{0} \\
\hline
\multicolumn{1}{|c|}{01X1XX10} &
\multicolumn{1}{c|}{{\scriptsize rs2[4]}} &
\multicolumn{1}{c|}{{\scriptsize rs2[2:0]}} &
\multicolumn{1}{c|}{X} &
\multicolumn{1}{c|}{{\scriptsize rs1[4]}} &
\multicolumn{1}{c|}{{\scriptsize rs1[2:0]}} &
\multicolumn{1}{c|}{11} &
\multicolumn{1}{c|}{{\scriptsize qu[2:0]}} &
\multicolumn{1}{c|}{{\scriptsize qz[2:0]}} &
\multicolumn{1}{c|}{0011011}\\
\hline
8& 1& 3& 1& 1& 3& 2& 3& 3& 7 \\
\end{tabular}
 \end{center}
 }
 \vspace{-0.1in}
 \caption[]{ESP.VLDEXT.U16.XP} \end{figure*}


\newpage
\subsubsection{ESP.VLDEXT.U8.IP}\label{instruction:ESPVLDEXTU8IP}
\textbf{Syntax}\\
ESP.VLDEXT.U8.IP qu,qz,rs1,imm16

\textbf{Description}\\
This instruction loads 128-bit data from memory and zero-extends the data from an 8-bit segment to a 16-bit segment.

\lstinline{imm} starts at -2048, ends at 2032, and steps in increments of 16.

\textbf{Operation}\\
\begin{lstlisting}
    DataIn[127:0] = load128(rs1[31:0])
    qu[ 15:  0] = { {8{0}}, DataIn[  7:  0] }
    qu[ 31: 16] = { {8{0}}, DataIn[ 15:  8] }
    qu[ 47: 32] = { {8{0}}, DataIn[ 23: 16] }
    qu[ 63: 48] = { {8{0}}, DataIn[ 31: 24] }
    qu[ 79: 64] = { {8{0}}, DataIn[ 39: 32] }
    qu[ 95: 80] = { {8{0}}, DataIn[ 47: 40] }
    qu[111: 96] = { {8{0}}, DataIn[ 55: 48] }
    qu[127:112] = { {8{0}}, DataIn[ 63: 56] }
    qz[ 15:  0] = { {8{0}}, DataIn[ 71: 64] }
    qz[ 31: 16] = { {8{0}}, DataIn[ 79: 72] }
    qz[ 47: 32] = { {8{0}}, DataIn[ 87: 80] }
    qz[ 63: 48] = { {8{0}}, DataIn[ 95: 88] }
    qz[ 79: 64] = { {8{0}}, DataIn[103: 96] }
    qz[ 95: 80] = { {8{0}}, DataIn[111:104] }
    qz[111: 96] = { {8{0}}, DataIn[119:112] }
    qz[127:112] = { {8{0}}, DataIn[127:120] }
    rs1[31:0] = rs1[31:0] + imm
\end{lstlisting}

\textbf{Schedule}\\
\begin{longtable}[c]{|l|l|l|}
    \hline
    \rowcolor{lightgray}
    \textbf{Operands} & \textbf{Use} & \textbf{Def} \\\hline
    \endfirsthead
    \hline
    \rowcolor{lightgray}
     \textbf{Operands} & \textbf{Use} & \textbf{Def} \\\hline
    \endhead
    \hline
    \endfoot

  qu & - & 2 \\\hline
  qv & - & 2 \\\hline
  rs1 & 1 & 1

\end{longtable}

\textbf{Format}\\

\begin{figure*}[!htbp]
 {\setlength{\tabcolsep}{1pt}
 \begin{center}
\begin{tabular}{@{}V@{}I@{}W@{}F@{}W@{}I@{}F@{}W@{}F@{}F@{}O}
\instbitrange{31}{27}&
\instbit{26}&
\instbitrange{25}{24}&
\instbitrange{23}{21}&
\instbitrange{20}{19}&
\instbit{18}&
\instbitrange{17}{15}&
\instbitrange{14}{13}&
\instbitrange{12}{10}&
\instbitrange{9}{7}&
\instbitrange{6}{0} \\
\hline
\multicolumn{1}{|c|}{1XX00} &
\multicolumn{1}{c|}{{\scriptsize imm16[3]}} &
\multicolumn{1}{c|}{00} &
\multicolumn{1}{c|}{{\scriptsize imm16[2:0]}} &
\multicolumn{1}{c|}{0X} &
\multicolumn{1}{c|}{{\scriptsize rs1[4]}} &
\multicolumn{1}{c|}{{\scriptsize rs1[2:0]}} &
\multicolumn{1}{c|}{00} &
\multicolumn{1}{c|}{{\scriptsize qu[2:0]}} &
\multicolumn{1}{c|}{{\scriptsize qz[2:0]}} &
\multicolumn{1}{c|}{0011111}\\
\hline
5& 1& 2& 3& 2& 1& 3& 2& 3& 3& 7 \\
\end{tabular}
 \end{center}
 }
 \vspace{-0.1in}
 \caption[]{ESP.VLDEXT.U8.IP} \end{figure*}


\newpage
\subsubsection{ESP.VLDEXT.U8.XP}\label{instruction:ESPVLDEXTU8XP}
\textbf{Syntax}\\
ESP.VLDEXT.U8.XP qu,qz,rs1,rs2

\textbf{Description}\\
This instruction loads 128-bit data from memory and zero-extends the data from an 8-bit segment to a 16-bit segment.

\textbf{Operation}\\
\begin{lstlisting}
    DataIn[127:0] = load128(rs1[31:0])
    qu[ 15:  0] = { {8{0}}, DataIn[  7:  0] }
    qu[ 31: 16] = { {8{0}}, DataIn[ 15:  8] }
    qu[ 47: 32] = { {8{0}}, DataIn[ 23: 16] }
    qu[ 63: 48] = { {8{0}}, DataIn[ 31: 24] }
    qu[ 79: 64] = { {8{0}}, DataIn[ 39: 32] }
    qu[ 95: 80] = { {8{0}}, DataIn[ 47: 40] }
    qu[111: 96] = { {8{0}}, DataIn[ 55: 48] }
    qu[127:112] = { {8{0}}, DataIn[ 63: 56] }
    qz[ 15:  0] = { {8{0}}, DataIn[ 71: 64] }
    qz[ 31: 16] = { {8{0}}, DataIn[ 79: 72] }
    qz[ 47: 32] = { {8{0}}, DataIn[ 87: 80] }
    qz[ 63: 48] = { {8{0}}, DataIn[ 95: 88] }
    qz[ 79: 64] = { {8{0}}, DataIn[103: 96] }
    qz[ 95: 80] = { {8{0}}, DataIn[111:104] }
    qz[111: 96] = { {8{0}}, DataIn[119:112] }
    qz[127:112] = { {8{0}}, DataIn[127:120] }
    rs1[31:0] = rs1[31:0] + rs2[31:0]
\end{lstlisting}

\textbf{Schedule}\\
\begin{longtable}[c]{|l|l|l|}
    \hline
    \rowcolor{lightgray}
    \textbf{Operands} & \textbf{Use} & \textbf{Def} \\\hline
    \endfirsthead
    \hline
    \rowcolor{lightgray}
     \textbf{Operands} & \textbf{Use} & \textbf{Def} \\\hline
    \endhead
    \hline
    \endfoot

  qu & - & 2 \\\hline
  qv & - & 2 \\\hline
  rs1 & 1 & 1 \\\hline
    rs2 & 1 & -

\end{longtable}

\textbf{Format}\\

\begin{figure*}[!htbp]
 {\setlength{\tabcolsep}{1pt}
 \begin{center}
\begin{tabular}{@{}E@{}I@{}F@{}I@{}I@{}F@{}W@{}F@{}F@{}O}
\instbitrange{31}{24}&
\instbit{23}&
\instbitrange{22}{20}&
\instbit{19}&
\instbit{18}&
\instbitrange{17}{15}&
\instbitrange{14}{13}&
\instbitrange{12}{10}&
\instbitrange{9}{7}&
\instbitrange{6}{0} \\
\hline
\multicolumn{1}{|c|}{01X1XX00} &
\multicolumn{1}{c|}{{\scriptsize rs2[4]}} &
\multicolumn{1}{c|}{{\scriptsize rs2[2:0]}} &
\multicolumn{1}{c|}{X} &
\multicolumn{1}{c|}{{\scriptsize rs1[4]}} &
\multicolumn{1}{c|}{{\scriptsize rs1[2:0]}} &
\multicolumn{1}{c|}{11} &
\multicolumn{1}{c|}{{\scriptsize qu[2:0]}} &
\multicolumn{1}{c|}{{\scriptsize qz[2:0]}} &
\multicolumn{1}{c|}{0011011}\\
\hline
8& 1& 3& 1& 1& 3& 2& 3& 3& 7 \\
\end{tabular}
 \end{center}
 }
 \vspace{-0.1in}
 \caption[]{ESP.VLDEXT.U8.XP} \end{figure*}


\newpage
\subsubsection{ESP.VLDHBC.16.INCP}\label{instruction:ESPVLDHBC16INCP}
\textbf{Syntax}\\
ESP.VLDHBC.16.INCP qu,qz,rs1

\textbf{Description}\\
This instruction loads 16-byte data from memory, extends it to 256 bits in the following way, and assigns the result to the registers \lstinline{qu} and \lstinline{qz}. After the access, the value in the register \lstinline{rs1} is incremented by 16.\\
\textbf{Operation}
\begin{lstlisting}
  dataIn[127:0] = load128(rs1[31:0])
  qu  = {2{dataIn[ 63: 48]}, 2{dataIn[ 47: 32]}, 2{dataIn[ 31: 16]}, 2{dataIn[ 15:  0]} }
  qz = {2{dataIn[127:112]}, 2{dataIn[111: 96]}, 2{dataIn[ 95: 80]}, 2{dataIn[ 79: 64]} }
  rs1[31:0] = rs1[31:0] + 16
\end{lstlisting}

\textbf{Schedule}\\
\begin{longtable}[c]{|l|l|l|}
    \hline
    \rowcolor{lightgray}
    \textbf{Operands} & \textbf{Use} & \textbf{Def} \\\hline
    \endfirsthead
    \hline
    \rowcolor{lightgray}
     \textbf{Operands} & \textbf{Use} & \textbf{Def} \\\hline
    \endhead
    \hline
    \endfoot

  qu & - & 2 \\\hline
  qz & - & 2 \\\hline
  rs1 & 1 & 1

\end{longtable}

\textbf{Format}\\

\begin{figure*}[!htbp]
 {\setlength{\tabcolsep}{1pt}
 \begin{center}
\begin{tabular}{@{}K@{}I@{}F@{}W@{}F@{}F@{}O}
\instbitrange{31}{19}&
\instbit{18}&
\instbitrange{17}{15}&
\instbitrange{14}{13}&
\instbitrange{12}{10}&
\instbitrange{9}{7}&
\instbitrange{6}{0} \\
\hline
\multicolumn{1}{|c|}{1XXXXX00XXX1X} &
\multicolumn{1}{c|}{{\scriptsize rs1[4]}} &
\multicolumn{1}{c|}{{\scriptsize rs1[2:0]}} &
\multicolumn{1}{c|}{00} &
\multicolumn{1}{c|}{{\scriptsize qu[2:0]}} &
\multicolumn{1}{c|}{{\scriptsize qz[2:0]}} &
\multicolumn{1}{c|}{0011111}\\
\hline
13& 1& 3& 2& 3& 3& 7 \\
\end{tabular}
 \end{center}
 }
 \vspace{-0.1in}
 \caption[]{ESP.VLDHBC.16.INCP} \end{figure*}


\newpage
\subsubsection{ESP.LD.QACC.H.H.128.IP}\label{instruction:ESPLDQACC_HH128IP}
\textbf{Syntax}\\
ESP.LD.QACC.H.H.128.IP rs1,imm16

\textbf{Description}\\
This instruction loads 16-byte data from memory into the special register \lstinline{QACC_H[255:128]}. After the access is completed, the value in the register \lstinline{rs1} is incremented by an 8-bit sign-extended constant in the instruction code segment, left-shifted by 4.\\

imm starts at -2048, ends at 2032, and steps by 16.

\textbf{Operation}
\begin{lstlisting}
  QACC_H[255:  128] = load128(rs1[31:0])
  rs1[31:0] = rs1[31:0] + imm
\end{lstlisting}

\textbf{Schedule}\\
\begin{longtable}[c]{|l|l|l|}
    \hline
    \rowcolor{lightgray}
    \textbf{Operands} & \textbf{Use} & \textbf{Def} \\\hline
    \endfirsthead
    \hline
    \rowcolor{lightgray}
     \textbf{Operands} & \textbf{Use} & \textbf{Def} \\\hline
    \endhead
    \hline
    \endfoot

  QACC\_H & - & 2 \\\hline
  rs1 & 1 & 1

\end{longtable}

\textbf{Format}\\

\begin{figure*}[!htbp]
 {\setlength{\tabcolsep}{1pt}
 \begin{center}
\begin{tabular}{@{}V@{}I@{}W@{}F@{}W@{}I@{}F@{}W@{}Y@{}W@{}O}
\instbitrange{31}{27}&
\instbit{26}&
\instbitrange{25}{24}&
\instbitrange{23}{21}&
\instbitrange{20}{19}&
\instbit{18}&
\instbitrange{17}{15}&
\instbitrange{14}{13}&
\instbitrange{12}{9}&
\instbitrange{8}{7}&
\instbitrange{6}{0} \\
\hline
\multicolumn{1}{|c|}{0XX01} &
\multicolumn{1}{c|}{{\scriptsize imm16[7]}} &
\multicolumn{1}{c|}{00} &
\multicolumn{1}{c|}{{\scriptsize imm16[6:4]}} &
\multicolumn{1}{c|}{0X} &
\multicolumn{1}{c|}{{\scriptsize rs1[4]}} &
\multicolumn{1}{c|}{{\scriptsize rs1[2:0]}} &
\multicolumn{1}{c|}{00} &
\multicolumn{1}{c|}{{\scriptsize imm16[3:0]}} &
\multicolumn{1}{c|}{XX} &
\multicolumn{1}{c|}{0011111}\\
\hline
5& 1& 2& 3& 2& 1& 3& 2& 4& 2& 7 \\
\end{tabular}
 \end{center}
 }
 \vspace{-0.1in}
 \caption[]{ESP.LD.QACC.H.H.128.IP} \end{figure*}


\newpage
\subsubsection{ESP.LD.QACC.H.L.128.IP}\label{instruction:ESPLDQACC_HL128IP}
\textbf{Syntax}\\
ESP.LD.QACC.H.L.128.IP rs1,imm16

\textbf{Description}\\
This instruction loads 16-byte data from memory into the special register \lstinline{QACC_H[127:0]}. After the access is completed, the value in the register \lstinline{rs1} is incremented by an 8-bit sign-extended constant in the instruction code segment, left-shifted by 4.\\

imm starts at -2048, ends at 2032, and steps by 16.

\textbf{Operation}
\begin{lstlisting}
  QACC_H[127:  0] = load128(rs1[31:0])
  rs1[31:0] = rs1[31:0] + imm
\end{lstlisting}

\textbf{Schedule}\\
\begin{longtable}[c]{|l|l|l|}
    \hline
    \rowcolor{lightgray}
    \textbf{Operands} & \textbf{Use} & \textbf{Def} \\\hline
    \endfirsthead
    \hline
    \rowcolor{lightgray}
     \textbf{Operands} & \textbf{Use} & \textbf{Def} \\\hline
    \endhead
    \hline
    \endfoot

  QACC\_H & - & 2 \\\hline
  rs1 & 1 & 1

\end{longtable}

\textbf{Format}\\

\begin{figure*}[!htbp]
 {\setlength{\tabcolsep}{1pt}
 \begin{center}
\begin{tabular}{@{}V@{}I@{}W@{}F@{}W@{}I@{}F@{}W@{}Y@{}W@{}O}
\instbitrange{31}{27}&
\instbit{26}&
\instbitrange{25}{24}&
\instbitrange{23}{21}&
\instbitrange{20}{19}&
\instbit{18}&
\instbitrange{17}{15}&
\instbitrange{14}{13}&
\instbitrange{12}{9}&
\instbitrange{8}{7}&
\instbitrange{6}{0} \\
\hline
\multicolumn{1}{|c|}{0XX01} &
\multicolumn{1}{c|}{{\scriptsize imm16[7]}} &
\multicolumn{1}{c|}{00} &
\multicolumn{1}{c|}{{\scriptsize imm16[6:4]}} &
\multicolumn{1}{c|}{1X} &
\multicolumn{1}{c|}{{\scriptsize rs1[4]}} &
\multicolumn{1}{c|}{{\scriptsize rs1[2:0]}} &
\multicolumn{1}{c|}{00} &
\multicolumn{1}{c|}{{\scriptsize imm16[3:0]}} &
\multicolumn{1}{c|}{00} &
\multicolumn{1}{c|}{0011111}\\
\hline
5& 1& 2& 3& 2& 1& 3& 2& 4& 2& 7 \\
\end{tabular}
 \end{center}
 }
 \vspace{-0.1in}
 \caption[]{ESP.LD.QACC.H.L.128.IP} \end{figure*}


\newpage
\subsubsection{ESP.LD.QACC.L.H.128.IP}\label{instruction:ESPLDQACC_LH128IP}
\textbf{Syntax}\\
ESP.LD.QACC.L.H.128.IP rs1,imm16

\textbf{Description}\\
This instruction loads 16-byte data from memory into the special register \lstinline{QACC_L[255:128]}. After the access is completed, the value in the register \lstinline{rs1} is incremented by an 8-bit sign-extended constant in the instruction code segment, left-shifted by 4.\\

imm starts at -2048, ends at 2032, and steps by 16.

\textbf{Operation}
\begin{lstlisting}
  QACC_L[255:  128] = load128(rs1[31:0])
  rs1[31:0] = rs1[31:0] + imm
\end{lstlisting}

\textbf{Schedule}\\
\begin{longtable}[c]{|l|l|l|}
    \hline
    \rowcolor{lightgray}
    \textbf{Operands} & \textbf{Use} & \textbf{Def} \\\hline
    \endfirsthead
    \hline
    \rowcolor{lightgray}
     \textbf{Operands} & \textbf{Use} & \textbf{Def} \\\hline
    \endhead
    \hline
    \endfoot

  QACC\_L & - & 2 \\\hline
  rs1 & 1 & 1

\end{longtable}

\textbf{Format}\\

\begin{figure*}[!htbp]
 {\setlength{\tabcolsep}{1pt}
 \begin{center}
\begin{tabular}{@{}V@{}I@{}W@{}F@{}W@{}I@{}F@{}W@{}Y@{}W@{}O}
\instbitrange{31}{27}&
\instbit{26}&
\instbitrange{25}{24}&
\instbitrange{23}{21}&
\instbitrange{20}{19}&
\instbit{18}&
\instbitrange{17}{15}&
\instbitrange{14}{13}&
\instbitrange{12}{9}&
\instbitrange{8}{7}&
\instbitrange{6}{0} \\
\hline
\multicolumn{1}{|c|}{0XX00} &
\multicolumn{1}{c|}{{\scriptsize imm16[7]}} &
\multicolumn{1}{c|}{00} &
\multicolumn{1}{c|}{{\scriptsize imm16[6:4]}} &
\multicolumn{1}{c|}{0X} &
\multicolumn{1}{c|}{{\scriptsize rs1[4]}} &
\multicolumn{1}{c|}{{\scriptsize rs1[2:0]}} &
\multicolumn{1}{c|}{00} &
\multicolumn{1}{c|}{{\scriptsize imm16[3:0]}} &
\multicolumn{1}{c|}{XX} &
\multicolumn{1}{c|}{0011111}\\
\hline
5& 1& 2& 3& 2& 1& 3& 2& 4& 2& 7 \\
\end{tabular}
 \end{center}
 }
 \vspace{-0.1in}
 \caption[]{ESP.LD.QACC.L.H.128.IP} \end{figure*}


\newpage
\subsubsection{ESP.LD.QACC.L.L.128.IP}\label{instruction:ESPLDQACC_LL128IP}
\textbf{Syntax}\\
ESP.LD.QACC.L.L.128.IP rs1,imm16

\textbf{Description}\\
This instruction loads 16-byte data from memory into the special register \lstinline{QACC_L[127:0]}. After the access is completed, the value in the register \lstinline{rs1} is incremented by an 8-bit sign-extended constant in the instruction code segment, left-shifted by 4.\\

imm starts at -2048, ends at 2032, and steps by 16.

\textbf{Operation}
\begin{lstlisting}
  QACC_L[127:  0] = load128(rs1[31:0])
  rs1[31:0] = rs1[31:0] + imm
\end{lstlisting}

\textbf{Schedule}\\
\begin{longtable}[c]{|l|l|l|}
    \hline
    \rowcolor{lightgray}
    \textbf{Operands} & \textbf{Use} & \textbf{Def} \\\hline
    \endfirsthead
    \hline
    \rowcolor{lightgray}
     \textbf{Operands} & \textbf{Use} & \textbf{Def} \\\hline
    \endhead
    \hline
    \endfoot

  QACC\_L & - & 2 \\\hline
  rs1 & 1 & 1

\end{longtable}

\textbf{Format}\\

\begin{figure*}[!htbp]
 {\setlength{\tabcolsep}{1pt}
 \begin{center}
\begin{tabular}{@{}V@{}I@{}W@{}F@{}W@{}I@{}F@{}W@{}Y@{}W@{}O}
\instbitrange{31}{27}&
\instbit{26}&
\instbitrange{25}{24}&
\instbitrange{23}{21}&
\instbitrange{20}{19}&
\instbit{18}&
\instbitrange{17}{15}&
\instbitrange{14}{13}&
\instbitrange{12}{9}&
\instbitrange{8}{7}&
\instbitrange{6}{0} \\
\hline
\multicolumn{1}{|c|}{0XX00} &
\multicolumn{1}{c|}{{\scriptsize imm16[7]}} &
\multicolumn{1}{c|}{00} &
\multicolumn{1}{c|}{{\scriptsize imm16[6:4]}} &
\multicolumn{1}{c|}{1X} &
\multicolumn{1}{c|}{{\scriptsize rs1[4]}} &
\multicolumn{1}{c|}{{\scriptsize rs1[2:0]}} &
\multicolumn{1}{c|}{00} &
\multicolumn{1}{c|}{{\scriptsize imm16[3:0]}} &
\multicolumn{1}{c|}{00} &
\multicolumn{1}{c|}{0011111}\\
\hline
5& 1& 2& 3& 2& 1& 3& 2& 4& 2& 7 \\
\end{tabular}
 \end{center}
 }
 \vspace{-0.1in}
 \caption[]{ESP.LD.QACC.L.L.128.IP} \end{figure*}


\newpage
\subsubsection{ESP.LD.UA.STATE.IP}\label{instruction:ESPLDUA_STATEIP}
\textbf{Syntax}\\
ESP.LD.UA.STATE.IP rs1,imm16

\textbf{Description}\\
This instruction loads 16-byte data from memory into the special register \lstinline{UA_STATE[127:0]}. After the access is completed, the value in the register \lstinline{rs1} is incremented by an 8-bit sign-extended constant in the instruction code segment, left-shifted by 4.\\

imm starts at -2048, ends at 2032, and steps by 16.

\textbf{Operation}
\begin{lstlisting}
  UA_STATE[127:  0] = load128(rs1[31:0])
  rs1[31:0] = rs1[31:0] + imm
\end{lstlisting}

\textbf{Schedule}\\
\begin{longtable}[c]{|l|l|l|}
    \hline
    \rowcolor{lightgray}
    \textbf{Operands} & \textbf{Use} & \textbf{Def} \\\hline
    \endfirsthead
    \hline
    \rowcolor{lightgray}
     \textbf{Operands} & \textbf{Use} & \textbf{Def} \\\hline
    \endhead
    \hline
    \endfoot

  UA\_STATE & - & 2 \\\hline
  rs1 & 1 & 1

\end{longtable}

\textbf{Format}\\

\begin{figure*}[!htbp]
 {\setlength{\tabcolsep}{1pt}
 \begin{center}
\begin{tabular}{@{}Y@{}W@{}W@{}F@{}W@{}I@{}F@{}W@{}F@{}F@{}O}
\instbitrange{31}{28}&
\instbitrange{27}{26}&
\instbitrange{25}{24}&
\instbitrange{23}{21}&
\instbitrange{20}{19}&
\instbit{18}&
\instbitrange{17}{15}&
\instbitrange{14}{13}&
\instbitrange{12}{10}&
\instbitrange{9}{7}&
\instbitrange{6}{0} \\
\hline
\multicolumn{1}{|c|}{0XX0} &
\multicolumn{1}{c|}{{\scriptsize imm16[7:6]}} &
\multicolumn{1}{c|}{00} &
\multicolumn{1}{c|}{{\scriptsize imm16[5:3]}} &
\multicolumn{1}{c|}{1X} &
\multicolumn{1}{c|}{{\scriptsize rs1[4]}} &
\multicolumn{1}{c|}{{\scriptsize rs1[2:0]}} &
\multicolumn{1}{c|}{00} &
\multicolumn{1}{c|}{{\scriptsize imm16[2:0]}} &
\multicolumn{1}{c|}{010} &
\multicolumn{1}{c|}{0011111}\\
\hline
4& 2& 2& 3& 2& 1& 3& 2& 3& 3& 7 \\
\end{tabular}
 \end{center}
 }
 \vspace{-0.1in}
 \caption[]{ESP.LD.UA.STATE.IP} \end{figure*}


\newpage
\subsubsection{ESP.LDXQ.32}\label{instruction:ESPLDXQ32}
\textbf{Syntax}\\
ESP.LDXQ.32 qu,qw,rs1,sel4,sel8

\textbf{Description}\\
This instruction selects one of the 8 segments of 16-bit data in \lstinline{qw} as the addend according to the immediate number \lstinline{sel8}. It performs the addition of the addend left-shifted by 2 bits to the value of the address register \lstinline{rs1}, uses the result as the access address, and stores the loaded data in a 32-bit data segment in the register \lstinline{qu} according to the value of the immediate number \lstinline{sel4}. \\
\textbf{Operation}
\begin{lstlisting}
  vaddr0[31:0] = rs1[31:0] + qw[ 15:  0] * 4
  vaddr1[31:0] = rs1[31:0] + qw[ 31: 16] * 4
  vaddr2[31:0] = rs1[31:0] + qw[ 47: 32] * 4
  vaddr3[31:0] = rs1[31:0] + qw[ 63: 47] * 4
  vaddr4[31:0] = rs1[31:0] + qw[ 79: 64] * 4
  vaddr5[31:0] = rs1[31:0] + qw[ 95: 80] * 4
  vaddr6[31:0] = rs1[31:0] + qw[111: 96] * 4
  vaddr7[31:0] = rs1[31:0] + qw[127:112] * 4

if sel8 == 0:
  dataIn[31:0] = load32(vaddr0[31:0])
if sel8 == 1:
  dataIn[31:0] = load32(vaddr1[31:0])
if sel8 == 2:
  dataIn[31:0] = load32(vaddr2[31:0])
if sel8 == 3:
  dataIn[31:0] = load32(vaddr3[31:0])
if sel8 == 4:
  dataIn[31:0] = load32(vaddr4[31:0])
if sel8 == 5:
  dataIn[31:0] = load32(vaddr5[31:0])
if sel8 == 6:
  dataIn[31:0] = load32(vaddr6[31:0])
if sel8 == 7:
  dataIn[31:0] = load32(vaddr7[31:0])

  qu[32*sel4+31:32*sel4] = dataIn[31:0]
\end{lstlisting}

\textbf{Schedule}\\
\begin{longtable}[c]{|l|l|l|}
    \hline
    \rowcolor{lightgray}
    \textbf{Operands} & \textbf{Use} & \textbf{Def} \\\hline
    \endfirsthead
    \hline
    \rowcolor{lightgray}
     \textbf{Operands} & \textbf{Use} & \textbf{Def} \\\hline
    \endhead
    \hline
    \endfoot

  qu & - & 2 \\\hline
  qw & 0 & - \\\hline
  rs1 & 1 & -

\end{longtable}

\textbf{Format}\\

\begin{figure*}[!htbp]
 {\setlength{\tabcolsep}{1pt}
 \begin{center}
\begin{tabular}{@{}S@{}W@{}I@{}W@{}I@{}I@{}I@{}F@{}W@{}F@{}W@{}I@{}O}
\instbitrange{31}{26}&
\instbitrange{25}{24}&
\instbit{23}&
\instbitrange{22}{21}&
\instbit{20}&
\instbit{19}&
\instbit{18}&
\instbitrange{17}{15}&
\instbitrange{14}{13}&
\instbitrange{12}{10}&
\instbitrange{9}{8}&
\instbit{7}&
\instbitrange{6}{0} \\
\hline
\multicolumn{1}{|c|}{01X0XX} &
\multicolumn{1}{c|}{{\scriptsize qw[1:0]}} &
\multicolumn{1}{c|}{0} &
\multicolumn{1}{c|}{{\scriptsize sel4[1:0]}} &
\multicolumn{1}{c|}{{\scriptsize sel8[2]}} &
\multicolumn{1}{c|}{{\scriptsize qw[2]}} &
\multicolumn{1}{c|}{{\scriptsize rs1[4]}} &
\multicolumn{1}{c|}{{\scriptsize rs1[2:0]}} &
\multicolumn{1}{c|}{11} &
\multicolumn{1}{c|}{{\scriptsize qu[2:0]}} &
\multicolumn{1}{c|}{{\scriptsize sel8[1:0]}} &
\multicolumn{1}{c|}{1} &
\multicolumn{1}{c|}{0011011}\\
\hline
6& 2& 1& 2& 1& 1& 1& 3& 2& 3& 2& 1& 7 \\
\end{tabular}
 \end{center}
 }
 \vspace{-0.1in}
 \caption[]{ESP.LDXQ.32} \end{figure*}


\newpage
\subsubsection{ESP.ST.QACC.H.H.128.IP}\label{instruction:ESPSTQACC_HH128IP}
\textbf{Syntax}\\
ESP.ST.QACC.H.H.128.IP rs1,imm16

\textbf{Description}\\
This instruction stores the upper 128 bits in the special register \lstinline{QACC_H} to memory. After the access is completed, the value in the register \lstinline{rs1} is incremented by an 8-bit sign-extended constant in the instruction code segment, left-shifted by 4.\\

imm starts at -2048, ends at 2032, and steps by 16.

\textbf{Operation}
\begin{lstlisting}
  MEM(rs1[31:0]) = QACC_H[255:128]
  rs1[31:0] = rs1[31:0] + imm
\end{lstlisting}

\textbf{Schedule}\\
\begin{longtable}[c]{|l|l|l|}
    \hline
    \rowcolor{lightgray}
    \textbf{Operands} & \textbf{Use} & \textbf{Def} \\\hline
    \endfirsthead
    \hline
    \rowcolor{lightgray}
     \textbf{Operands} & \textbf{Use} & \textbf{Def} \\\hline
    \endhead
    \hline
    \endfoot

  QACC\_H & 2 & - \\\hline
  rs1 & 1 & 1

\end{longtable}

\textbf{Format}\\

\begin{figure*}[!htbp]
 {\setlength{\tabcolsep}{1pt}
 \begin{center}
\begin{tabular}{@{}V@{}I@{}W@{}F@{}W@{}I@{}F@{}W@{}Y@{}W@{}O}
\instbitrange{31}{27}&
\instbit{26}&
\instbitrange{25}{24}&
\instbitrange{23}{21}&
\instbitrange{20}{19}&
\instbit{18}&
\instbitrange{17}{15}&
\instbitrange{14}{13}&
\instbitrange{12}{9}&
\instbitrange{8}{7}&
\instbitrange{6}{0} \\
\hline
\multicolumn{1}{|c|}{0XX11} &
\multicolumn{1}{c|}{{\scriptsize imm16[7]}} &
\multicolumn{1}{c|}{00} &
\multicolumn{1}{c|}{{\scriptsize imm16[6:4]}} &
\multicolumn{1}{c|}{0X} &
\multicolumn{1}{c|}{{\scriptsize rs1[4]}} &
\multicolumn{1}{c|}{{\scriptsize rs1[2:0]}} &
\multicolumn{1}{c|}{00} &
\multicolumn{1}{c|}{{\scriptsize imm16[3:0]}} &
\multicolumn{1}{c|}{XX} &
\multicolumn{1}{c|}{0011111}\\
\hline
5& 1& 2& 3& 2& 1& 3& 2& 4& 2& 7 \\
\end{tabular}
 \end{center}
 }
 \vspace{-0.1in}
 \caption[]{ESP.ST.QACC.H.H.128.IP} \end{figure*}


\newpage
\subsubsection{ESP.ST.QACC.H.L.128.IP}\label{instruction:ESPSTQACC_HL128IP}
\textbf{Syntax}\\
ESP.ST.QACC.H.L.128.IP rs1,imm16

\textbf{Description}\\
This instruction stores the lower 128 bits in the special register \lstinline{QACC_H} to memory. After the access is completed, the value in the register \lstinline{rs1} is incremented by an 8-bit sign-extended constant in the instruction code segment, left-shifted by 4.\\

imm starts at -2048, ends at 2032, and steps by 16.

\textbf{Operation}
\begin{lstlisting}
  MEM(rs1[31:0]) = QACC_H[127:0]
  rs1[31:0] = rs1[31:0] + imm
\end{lstlisting}

\textbf{Schedule}\\
\begin{longtable}[c]{|l|l|l|}
    \hline
    \rowcolor{lightgray}
    \textbf{Operands} & \textbf{Use} & \textbf{Def} \\\hline
    \endfirsthead
    \hline
    \rowcolor{lightgray}
     \textbf{Operands} & \textbf{Use} & \textbf{Def} \\\hline
    \endhead
    \hline
    \endfoot

  QACC\_H & 2 & - \\\hline
  rs1 & 1 & 1

\end{longtable}

\textbf{Format}\\

\begin{figure*}[!htbp]
 {\setlength{\tabcolsep}{1pt}
 \begin{center}
\begin{tabular}{@{}V@{}I@{}W@{}F@{}W@{}I@{}F@{}W@{}Y@{}W@{}O}
\instbitrange{31}{27}&
\instbit{26}&
\instbitrange{25}{24}&
\instbitrange{23}{21}&
\instbitrange{20}{19}&
\instbit{18}&
\instbitrange{17}{15}&
\instbitrange{14}{13}&
\instbitrange{12}{9}&
\instbitrange{8}{7}&
\instbitrange{6}{0} \\
\hline
\multicolumn{1}{|c|}{0XX11} &
\multicolumn{1}{c|}{{\scriptsize imm16[7]}} &
\multicolumn{1}{c|}{00} &
\multicolumn{1}{c|}{{\scriptsize imm16[6:4]}} &
\multicolumn{1}{c|}{1X} &
\multicolumn{1}{c|}{{\scriptsize rs1[4]}} &
\multicolumn{1}{c|}{{\scriptsize rs1[2:0]}} &
\multicolumn{1}{c|}{00} &
\multicolumn{1}{c|}{{\scriptsize imm16[3:0]}} &
\multicolumn{1}{c|}{00} &
\multicolumn{1}{c|}{0011111}\\
\hline
5& 1& 2& 3& 2& 1& 3& 2& 4& 2& 7 \\
\end{tabular}
 \end{center}
 }
 \vspace{-0.1in}
 \caption[]{ESP.ST.QACC.H.L.128.IP} \end{figure*}


\newpage
\subsubsection{ESP.ST.QACC.L.H.128.IP}\label{instruction:ESPSTQACC_LH128IP}
\textbf{Syntax}\\
ESP.ST.QACC.L.H.128.IP rs1,imm16

\textbf{Description}\\
This instruction stores the upper 128 bits in the special register \lstinline{QACC_L} to memory. After the access is completed, the value in the register \lstinline{rs1} is incremented by an 8-bit sign-extended constant in the instruction code segment, left-shifted by 4.\\

imm starts at -2048, ends at 2032, and steps by 16.

\textbf{Operation}
\begin{lstlisting}
  MEM(rs1[31:0]) = QACC_L[255:128]
  rs1[31:0] = rs1[31:0] + imm
\end{lstlisting}

\textbf{Schedule}\\
\begin{longtable}[c]{|l|l|l|}
    \hline
    \rowcolor{lightgray}
    \textbf{Operands} & \textbf{Use} & \textbf{Def} \\\hline
    \endfirsthead
    \hline
    \rowcolor{lightgray}
     \textbf{Operands} & \textbf{Use} & \textbf{Def} \\\hline
    \endhead
    \hline
    \endfoot

  QACC\_L & 2 & - \\\hline
  rs1 & 1 & 1

\end{longtable}

\textbf{Format}\\

\begin{figure*}[!htbp]
 {\setlength{\tabcolsep}{1pt}
 \begin{center}
\begin{tabular}{@{}V@{}I@{}W@{}F@{}W@{}I@{}F@{}W@{}Y@{}W@{}O}
\instbitrange{31}{27}&
\instbit{26}&
\instbitrange{25}{24}&
\instbitrange{23}{21}&
\instbitrange{20}{19}&
\instbit{18}&
\instbitrange{17}{15}&
\instbitrange{14}{13}&
\instbitrange{12}{9}&
\instbitrange{8}{7}&
\instbitrange{6}{0} \\
\hline
\multicolumn{1}{|c|}{0XX10} &
\multicolumn{1}{c|}{{\scriptsize imm16[7]}} &
\multicolumn{1}{c|}{00} &
\multicolumn{1}{c|}{{\scriptsize imm16[6:4]}} &
\multicolumn{1}{c|}{0X} &
\multicolumn{1}{c|}{{\scriptsize rs1[4]}} &
\multicolumn{1}{c|}{{\scriptsize rs1[2:0]}} &
\multicolumn{1}{c|}{00} &
\multicolumn{1}{c|}{{\scriptsize imm16[3:0]}} &
\multicolumn{1}{c|}{XX} &
\multicolumn{1}{c|}{0011111}\\
\hline
5& 1& 2& 3& 2& 1& 3& 2& 4& 2& 7 \\
\end{tabular}
 \end{center}
 }
 \vspace{-0.1in}
 \caption[]{ESP.ST.QACC.L.H.128.IP} \end{figure*}


\newpage
\subsubsection{ESP.ST.QACC.L.L.128.IP}\label{instruction:ESPSTQACC_LL128IP}
\textbf{Syntax}\\
ESP.ST.QACC.L.L.128.IP rs1,imm16

\textbf{Description}\\
This instruction stores the lower 128 bits in the special register \lstinline{QACC_L} to memory. After the access is completed, the value in the register \lstinline{rs1} is incremented by an 8-bit sign-extended constant in the instruction code segment, left-shifted by 4.\\

imm starts at -2048, ends at 2032, and steps by 16.

\textbf{Operation}
\begin{lstlisting}
  MEM(rs1[31:0]) = QACC_L[127:0]
  rs1[31:0] = rs1[31:0] + imm
\end{lstlisting}

\textbf{Schedule}\\
\begin{longtable}[c]{|l|l|l|}
    \hline
    \rowcolor{lightgray}
    \textbf{Operands} & \textbf{Use} & \textbf{Def} \\\hline
    \endfirsthead
    \hline
    \rowcolor{lightgray}
     \textbf{Operands} & \textbf{Use} & \textbf{Def} \\\hline
    \endhead
    \hline
    \endfoot

  QACC\_L & 2 & - \\\hline
  rs1 & 1 & 1

\end{longtable}

\textbf{Format}\\

\begin{figure*}[!htbp]
 {\setlength{\tabcolsep}{1pt}
 \begin{center}
\begin{tabular}{@{}V@{}I@{}W@{}F@{}W@{}I@{}F@{}W@{}Y@{}W@{}O}
\instbitrange{31}{27}&
\instbit{26}&
\instbitrange{25}{24}&
\instbitrange{23}{21}&
\instbitrange{20}{19}&
\instbit{18}&
\instbitrange{17}{15}&
\instbitrange{14}{13}&
\instbitrange{12}{9}&
\instbitrange{8}{7}&
\instbitrange{6}{0} \\
\hline
\multicolumn{1}{|c|}{0XX10} &
\multicolumn{1}{c|}{{\scriptsize imm16[7]}} &
\multicolumn{1}{c|}{00} &
\multicolumn{1}{c|}{{\scriptsize imm16[6:4]}} &
\multicolumn{1}{c|}{1X} &
\multicolumn{1}{c|}{{\scriptsize rs1[4]}} &
\multicolumn{1}{c|}{{\scriptsize rs1[2:0]}} &
\multicolumn{1}{c|}{00} &
\multicolumn{1}{c|}{{\scriptsize imm16[3:0]}} &
\multicolumn{1}{c|}{00} &
\multicolumn{1}{c|}{0011111}\\
\hline
5& 1& 2& 3& 2& 1& 3& 2& 4& 2& 7 \\
\end{tabular}
 \end{center}
 }
 \vspace{-0.1in}
 \caption[]{ESP.ST.QACC.L.L.128.IP} \end{figure*}


\newpage
\subsubsection{ESP.ST.UA.STATE.IP}\label{instruction:ESPSTUA_STATEIP}
\textbf{Syntax}\\
ESP.ST.UA.STATE.IP rs1,imm16

\textbf{Description}\\
This instruction stores the 128-bit data in the special register \lstinline{UA_STATE} to memory. After the access is completed, the value in the register \lstinline{rs1} is incremented by an 8-bit sign-extended constant in the instruction code segment, left-shifted by 4.\\

imm starts at -2048, ends at 2032, and steps by 16.

\textbf{Operation}
\begin{lstlisting}
  MEM(rs1[31:0]) = UA_STATE[127:0]
  rs1[31:0] = rs1[31:0] + imm
\end{lstlisting}

\textbf{Schedule}\\
\begin{longtable}[c]{|l|l|l|}
    \hline
    \rowcolor{lightgray}
    \textbf{Operands} & \textbf{Use} & \textbf{Def} \\\hline
    \endfirsthead
    \hline
    \rowcolor{lightgray}
     \textbf{Operands} & \textbf{Use} & \textbf{Def} \\\hline
    \endhead
    \hline
    \endfoot

  UA\_STATE & 2 & - \\\hline
  rs1 & 1 & 1

\end{longtable}

\textbf{Format}\\

\begin{figure*}[!htbp]
 {\setlength{\tabcolsep}{1pt}
 \begin{center}
\begin{tabular}{@{}Y@{}W@{}W@{}F@{}W@{}I@{}F@{}W@{}F@{}F@{}O}
\instbitrange{31}{28}&
\instbitrange{27}{26}&
\instbitrange{25}{24}&
\instbitrange{23}{21}&
\instbitrange{20}{19}&
\instbit{18}&
\instbitrange{17}{15}&
\instbitrange{14}{13}&
\instbitrange{12}{10}&
\instbitrange{9}{7}&
\instbitrange{6}{0} \\
\hline
\multicolumn{1}{|c|}{0XX1} &
\multicolumn{1}{c|}{{\scriptsize imm16[7:6]}} &
\multicolumn{1}{c|}{00} &
\multicolumn{1}{c|}{{\scriptsize imm16[5:3]}} &
\multicolumn{1}{c|}{1X} &
\multicolumn{1}{c|}{{\scriptsize rs1[4]}} &
\multicolumn{1}{c|}{{\scriptsize rs1[2:0]}} &
\multicolumn{1}{c|}{00} &
\multicolumn{1}{c|}{{\scriptsize imm16[2:0]}} &
\multicolumn{1}{c|}{010} &
\multicolumn{1}{c|}{0011111}\\
\hline
4& 2& 2& 3& 2& 1& 3& 2& 3& 3& 7 \\
\end{tabular}
 \end{center}
 }
 \vspace{-0.1in}
 \caption[]{ESP.ST.UA.STATE.IP} \end{figure*}


\newpage
\subsubsection{ESP.STXQ.32}\label{instruction:ESPSTXQ32}
\textbf{Syntax}\\
ESP.STXQ.32 qu,qw,rs1,sel4,sel8

\textbf{Description}\\
This instruction selects one of the four 32-bit data segments in the register \lstinline{qw} according to the immediate number \lstinline{sel4}. Then, it selects an addend from the eight 16-bit data segments in the register \lstinline{qw} according to the immediate number \lstinline{sel8}, adds the addend left-shifted by 2 bits to the value of the address register \lstinline{rs1}, and uses the result as the write address.\\
\textbf{Operation}
\begin{lstlisting}
  vaddr0[31:0] = rs1[31:0] + qw[ 15:  0] * 4
  vaddr1[31:0] = rs1[31:0] + qw[ 31: 16] * 4
  vaddr2[31:0] = rs1[31:0] + qw[ 47: 32] * 4
  vaddr3[31:0] = rs1[31:0] + qw[ 63: 47] * 4
  vaddr4[31:0] = rs1[31:0] + qw[ 79: 64] * 4
  vaddr5[31:0] = rs1[31:0] + qw[ 95: 80] * 4
  vaddr6[31:0] = rs1[31:0] + qw[111: 96] * 4
  vaddr7[31:0] = rs1[31:0] + qw[127:112] * 4

if sel8 == 0:
  qu[32*sel4+31:32*sel4] =>store32(vaddr0[31:0])
if sel8 == 1:
  qu[32*sel4+31:32*sel4] =>store32(vaddr1[31:0])
if sel8 == 2:
  qu[32*sel4+31:32*sel4] =>store32(vaddr2[31:0])
if sel8 == 3:
  qu[32*sel4+31:32*sel4] =>store32(vaddr3[31:0])
if sel8 == 4:
  qu[32*sel4+31:32*sel4] =>store32(vaddr4[31:0])
if sel8 == 5:
  qu[32*sel4+31:32*sel4] =>store32(vaddr5[31:0])
if sel8 == 6:
  qu[32*sel4+31:32*sel4] =>store32(vaddr6[31:0])
if sel8 == 7:
  qu[32*sel4+31:32*sel4] =>store32(vaddr7[31:0])
\end{lstlisting}

\textbf{Schedule}\\
\begin{longtable}[c]{|l|l|l|}
    \hline
    \rowcolor{lightgray}
    \textbf{Operands} & \textbf{Use} & \textbf{Def} \\\hline
    \endfirsthead
    \hline
    \rowcolor{lightgray}
     \textbf{Operands} & \textbf{Use} & \textbf{Def} \\\hline
    \endhead
    \hline
    \endfoot

  qu & 2 & - \\\hline
  qw & 0 & - \\\hline
  rs1 & 1 & 1

\end{longtable}

\textbf{Format}\\

\begin{figure*}[!htbp]
 {\setlength{\tabcolsep}{1pt}
 \begin{center}
\begin{tabular}{@{}S@{}W@{}I@{}W@{}I@{}I@{}I@{}F@{}W@{}F@{}W@{}I@{}O}
\instbitrange{31}{26}&
\instbitrange{25}{24}&
\instbit{23}&
\instbitrange{22}{21}&
\instbit{20}&
\instbit{19}&
\instbit{18}&
\instbitrange{17}{15}&
\instbitrange{14}{13}&
\instbitrange{12}{10}&
\instbitrange{9}{8}&
\instbit{7}&
\instbitrange{6}{0} \\
\hline
\multicolumn{1}{|c|}{01X0XX} &
\multicolumn{1}{c|}{{\scriptsize qw[1:0]}} &
\multicolumn{1}{c|}{1} &
\multicolumn{1}{c|}{{\scriptsize sel4[1:0]}} &
\multicolumn{1}{c|}{{\scriptsize sel8[2]}} &
\multicolumn{1}{c|}{{\scriptsize qw[2]}} &
\multicolumn{1}{c|}{{\scriptsize rs1[4]}} &
\multicolumn{1}{c|}{{\scriptsize rs1[2:0]}} &
\multicolumn{1}{c|}{11} &
\multicolumn{1}{c|}{{\scriptsize qu[2:0]}} &
\multicolumn{1}{c|}{{\scriptsize sel8[1:0]}} &
\multicolumn{1}{c|}{1} &
\multicolumn{1}{c|}{0011011}\\
\hline
6& 2& 1& 2& 1& 1& 1& 3& 2& 3& 2& 1& 7 \\
\end{tabular}
 \end{center}
 }
 \vspace{-0.1in}
 \caption[]{ESP.STXQ.32} \end{figure*}


\newpage
\subsubsection{ESP.VLD.128.IP}\label{instruction:ESPVLD128IP}
\textbf{Syntax}\\
ESP.VLD.128.IP qu,rs1,imm16

\textbf{Description}\\
This instruction loads 128-bit data from memory to the register \lstinline{qu}, with the pointer incremented by an immediate.\\

imm starts at -2048, ends at 2032, and steps by 16.

\textbf{Operation}
\begin{lstlisting}
  qu[127:  0] = load128(rs1[31:0])
  rs1[31:0] = rs1[31:0] + imm
\end{lstlisting}

\textbf{Schedule}\\
\begin{longtable}[c]{|l|l|l|}
    \hline
    \rowcolor{lightgray}
    \textbf{Operands} & \textbf{Use} & \textbf{Def} \\\hline
    \endfirsthead
    \hline
    \rowcolor{lightgray}
     \textbf{Operands} & \textbf{Use} & \textbf{Def} \\\hline
    \endhead
    \hline
    \endfoot

  qu & - & 2 \\\hline
  rs1 & 1 & 1

\end{longtable}

\textbf{Format}\\

\begin{figure*}[!htbp]
 {\setlength{\tabcolsep}{1pt}
 \begin{center}
\begin{tabular}{@{}I@{}V@{}O@{}I@{}F@{}W@{}F@{}F@{}O}
\instbit{31}&
\instbitrange{30}{26}&
\instbitrange{25}{19}&
\instbit{18}&
\instbitrange{17}{15}&
\instbitrange{14}{13}&
\instbitrange{12}{10}&
\instbitrange{9}{7}&
\instbitrange{6}{0} \\
\hline
\multicolumn{1}{|c|}{0} &
\multicolumn{1}{c|}{{\scriptsize imm16[7:3]}} &
\multicolumn{1}{c|}{110001X} &
\multicolumn{1}{c|}{{\scriptsize rs1[4]}} &
\multicolumn{1}{c|}{{\scriptsize rs1[2:0]}} &
\multicolumn{1}{c|}{00} &
\multicolumn{1}{c|}{{\scriptsize qu[2:0]}} &
\multicolumn{1}{c|}{{\scriptsize imm16[2:0]}} &
\multicolumn{1}{c|}{0011111}\\
\hline
1& 5& 7& 1& 3& 2& 3& 3& 7 \\
\end{tabular}
 \end{center}
 }
 \vspace{-0.1in}
 \caption[]{ESP.VLD.128.IP} \end{figure*}


\newpage
\subsubsection{ESP.VLD.128.XP}\label{instruction:ESPVLD128XP}
\textbf{Syntax}\\
ESP.VLD.128.XP qu,rs1,rs2

\textbf{Description}\\
This instruction loads 128-bit data from memory into the register \lstinline{qu}, with the pointer incremented by the value of the register \lstinline{rs2}.\\

\textbf{Operation}
\begin{lstlisting}
  qu[127:  0] = load128(rs1[31:0])
  rs1[31:0] = rs1[31:0] + rs2[31:0]
\end{lstlisting}

\textbf{Schedule}\\
\begin{longtable}[c]{|l|l|l|}
    \hline
    \rowcolor{lightgray}
    \textbf{Operands} & \textbf{Use} & \textbf{Def} \\\hline
    \endfirsthead
    \hline
    \rowcolor{lightgray}
     \textbf{Operands} & \textbf{Use} & \textbf{Def} \\\hline
    \endhead
    \hline
    \endfoot

  qu & - & 2 \\\hline
  rs1 & 1 & 1 \\\hline
  rs2 & 1 & -

\end{longtable}

\textbf{Format}\\

\begin{figure*}[!htbp]
 {\setlength{\tabcolsep}{1pt}
 \begin{center}
\begin{tabular}{@{}E@{}I@{}F@{}I@{}I@{}F@{}W@{}F@{}F@{}O}
\instbitrange{31}{24}&
\instbit{23}&
\instbitrange{22}{20}&
\instbit{19}&
\instbit{18}&
\instbitrange{17}{15}&
\instbitrange{14}{13}&
\instbitrange{12}{10}&
\instbitrange{9}{7}&
\instbitrange{6}{0} \\
\hline
\multicolumn{1}{|c|}{01X0XX0X} &
\multicolumn{1}{c|}{{\scriptsize rs2[4]}} &
\multicolumn{1}{c|}{{\scriptsize rs2[2:0]}} &
\multicolumn{1}{c|}{0} &
\multicolumn{1}{c|}{{\scriptsize rs1[4]}} &
\multicolumn{1}{c|}{{\scriptsize rs1[2:0]}} &
\multicolumn{1}{c|}{11} &
\multicolumn{1}{c|}{{\scriptsize qu[2:0]}} &
\multicolumn{1}{c|}{1X0} &
\multicolumn{1}{c|}{0011011}\\
\hline
8& 1& 3& 1& 1& 3& 2& 3& 3& 7 \\
\end{tabular}
 \end{center}
 }
 \vspace{-0.1in}
 \caption[]{ESP.VLD.128.XP} \end{figure*}


\newpage
\subsubsection{ESP.VLD.H.64.IP}\label{instruction:ESPVLDH64IP}
\textbf{Syntax}\\
ESP.VLD.H.64.IP qu,rs1,imm8

\textbf{Description}\\

This instruction loads 64-bit data from memory to the upper 64-bit segment in the register \lstinline{qu}, with the pointer incremented by an immediate.\\

imm starts at -1024, ends at 1016, and steps by 8.

\textbf{Operation}
\begin{lstlisting}
  qu[ 127:  64] = load64(rs1[31:0])
  rs1[31:0] = rs1[31:0] + imm
\end{lstlisting}

\textbf{Schedule}\\
\begin{longtable}[c]{|l|l|l|}
    \hline
    \rowcolor{lightgray}
    \textbf{Operands} & \textbf{Use} & \textbf{Def} \\\hline
    \endfirsthead
    \hline
    \rowcolor{lightgray}
     \textbf{Operands} & \textbf{Use} & \textbf{Def} \\\hline
    \endhead
    \hline
    \endfoot

  qu & - & 2 \\\hline
  rs1 & 1 & 1

\end{longtable}

\textbf{Format}\\

\begin{figure*}[!htbp]
 {\setlength{\tabcolsep}{1pt}
 \begin{center}
\begin{tabular}{@{}W@{}Y@{}W@{}I@{}Y@{}I@{}F@{}W@{}F@{}F@{}O}
\instbitrange{31}{30}&
\instbitrange{29}{26}&
\instbitrange{25}{24}&
\instbit{23}&
\instbitrange{22}{19}&
\instbit{18}&
\instbitrange{17}{15}&
\instbitrange{14}{13}&
\instbitrange{12}{10}&
\instbitrange{9}{7}&
\instbitrange{6}{0} \\
\hline
\multicolumn{1}{|c|}{01} &
\multicolumn{1}{c|}{{\scriptsize imm8[7:4]}} &
\multicolumn{1}{c|}{11} &
\multicolumn{1}{c|}{{\scriptsize imm8[3]}} &
\multicolumn{1}{c|}{110X} &
\multicolumn{1}{c|}{{\scriptsize rs1[4]}} &
\multicolumn{1}{c|}{{\scriptsize rs1[2:0]}} &
\multicolumn{1}{c|}{00} &
\multicolumn{1}{c|}{{\scriptsize qu[2:0]}} &
\multicolumn{1}{c|}{{\scriptsize imm8[2:0]}} &
\multicolumn{1}{c|}{0011111}\\
\hline
2& 4& 2& 1& 4& 1& 3& 2& 3& 3& 7 \\
\end{tabular}
 \end{center}
 }
 \vspace{-0.1in}
 \caption[]{ESP.VLD.H.64.IP} \end{figure*}


\newpage
\subsubsection{ESP.VLD.H.64.XP}\label{instruction:ESPVLDH64XP}
\textbf{Syntax}\\
ESP.VLD.H.64.XP qu,rs1,rs2

\textbf{Description}\\

This instruction loads 64-bit data from memory into the upper 64-bit segment of the register \lstinline{qu}, with the pointer incremented by the value in the register \lstinline{rs2}.\\

\textbf{Operation}
\begin{lstlisting}
  qu[ 127:  64] = load64(rs1[31:0])
  rs1[31:0] = rs1[31:0] + rs2[31:0]
\end{lstlisting}

\textbf{Schedule}\\
\begin{longtable}[c]{|l|l|l|}
    \hline
    \rowcolor{lightgray}
    \textbf{Operands} & \textbf{Use} & \textbf{Def} \\\hline
    \endfirsthead
    \hline
    \rowcolor{lightgray}
     \textbf{Operands} & \textbf{Use} & \textbf{Def} \\\hline
    \endhead
    \hline
    \endfoot

  qu & - & 2 \\\hline
  rs1 & 1 & 1 \\\hline
  rs2 & 1 & -

\end{longtable}

\textbf{Format}\\

\begin{figure*}[!htbp]
 {\setlength{\tabcolsep}{1pt}
 \begin{center}
\begin{tabular}{@{}E@{}I@{}F@{}I@{}I@{}F@{}W@{}F@{}F@{}O}
\instbitrange{31}{24}&
\instbit{23}&
\instbitrange{22}{20}&
\instbit{19}&
\instbit{18}&
\instbitrange{17}{15}&
\instbitrange{14}{13}&
\instbitrange{12}{10}&
\instbitrange{9}{7}&
\instbitrange{6}{0} \\
\hline
\multicolumn{1}{|c|}{01X0XX01} &
\multicolumn{1}{c|}{{\scriptsize rs2[4]}} &
\multicolumn{1}{c|}{{\scriptsize rs2[2:0]}} &
\multicolumn{1}{c|}{1} &
\multicolumn{1}{c|}{{\scriptsize rs1[4]}} &
\multicolumn{1}{c|}{{\scriptsize rs1[2:0]}} &
\multicolumn{1}{c|}{11} &
\multicolumn{1}{c|}{{\scriptsize qu[2:0]}} &
\multicolumn{1}{c|}{0X0} &
\multicolumn{1}{c|}{0011011}\\
\hline
8& 1& 3& 1& 1& 3& 2& 3& 3& 7 \\
\end{tabular}
 \end{center}
 }
 \vspace{-0.1in}
 \caption[]{ESP.VLD.H.64.XP} \end{figure*}


\newpage
\subsubsection{ESP.VLD.L.64.IP}\label{instruction:ESPVLDL64IP}
\textbf{Syntax}\\
ESP.VLD.L.64.IP qu,rs1,imm8

\textbf{Description}\\
This instruction loads 64-bit data from memory into the lower 64-bit segment of the register \lstinline{qu}, with the pointer incremented by an immediate.\\

imm starts at -1024, ends at 1016, and steps by 8.

\textbf{Operation}
\begin{lstlisting}
  qu[ 63:  0] = load64(rs1[31:0])
  rs1[31:0] = rs1[31:0] + imm
\end{lstlisting}

\textbf{Schedule}\\
\begin{longtable}[c]{|l|l|l|}
    \hline
    \rowcolor{lightgray}
    \textbf{Operands} & \textbf{Use} & \textbf{Def} \\\hline
    \endfirsthead
    \hline
    \rowcolor{lightgray}
     \textbf{Operands} & \textbf{Use} & \textbf{Def} \\\hline
    \endhead
    \hline
    \endfoot

  qu & - & 2 \\\hline
  rs1 & 1 & 1

\end{longtable}

\textbf{Format}\\

\begin{figure*}[!htbp]
 {\setlength{\tabcolsep}{1pt}
 \begin{center}
\begin{tabular}{@{}W@{}Y@{}W@{}I@{}Y@{}I@{}F@{}W@{}F@{}F@{}O}
\instbitrange{31}{30}&
\instbitrange{29}{26}&
\instbitrange{25}{24}&
\instbit{23}&
\instbitrange{22}{19}&
\instbit{18}&
\instbitrange{17}{15}&
\instbitrange{14}{13}&
\instbitrange{12}{10}&
\instbitrange{9}{7}&
\instbitrange{6}{0} \\
\hline
\multicolumn{1}{|c|}{00} &
\multicolumn{1}{c|}{{\scriptsize imm8[7:4]}} &
\multicolumn{1}{c|}{11} &
\multicolumn{1}{c|}{{\scriptsize imm8[3]}} &
\multicolumn{1}{c|}{110X} &
\multicolumn{1}{c|}{{\scriptsize rs1[4]}} &
\multicolumn{1}{c|}{{\scriptsize rs1[2:0]}} &
\multicolumn{1}{c|}{00} &
\multicolumn{1}{c|}{{\scriptsize qu[2:0]}} &
\multicolumn{1}{c|}{{\scriptsize imm8[2:0]}} &
\multicolumn{1}{c|}{0011111}\\
\hline
2& 4& 2& 1& 4& 1& 3& 2& 3& 3& 7 \\
\end{tabular}
 \end{center}
 }
 \vspace{-0.1in}
 \caption[]{ESP.VLD.L.64.IP} \end{figure*}


\newpage
\subsubsection{ESP.VLD.L.64.XP}\label{instruction:ESPVLDL64XP}
\textbf{Syntax}\\
ESP.VLD.L.64.XP qu,rs1,rs2

\textbf{Description}\\

This instruction loads 64-bit data from memory into the lower 64-bit segment of the \lstinline{qu} register, with the pointer incremented by the value in the \lstinline{rs2} register.\\

\textbf{Operation}
\begin{lstlisting}
  qu[ 63:  0] = load64(rs1[31:0])
  rs1[31:0] = rs1[31:0] + rs2[31:0]
\end{lstlisting}

\textbf{Schedule}\\
\begin{longtable}[c]{|l|l|l|}
    \hline
    \rowcolor{lightgray}
    \textbf{Operands} & \textbf{Use} & \textbf{Def} \\\hline
    \endfirsthead
    \hline
    \rowcolor{lightgray}
     \textbf{Operands} & \textbf{Use} & \textbf{Def} \\\hline
    \endhead
    \hline
    \endfoot

  qu & - & 2 \\\hline
  rs1 & 1 & 1 \\\hline
  rs2 & 1 & -

\end{longtable}

\textbf{Format}\\

\begin{figure*}[!htbp]
 {\setlength{\tabcolsep}{1pt}
 \begin{center}
\begin{tabular}{@{}E@{}I@{}F@{}I@{}I@{}F@{}W@{}F@{}F@{}O}
\instbitrange{31}{24}&
\instbit{23}&
\instbitrange{22}{20}&
\instbit{19}&
\instbit{18}&
\instbitrange{17}{15}&
\instbitrange{14}{13}&
\instbitrange{12}{10}&
\instbitrange{9}{7}&
\instbitrange{6}{0} \\
\hline
\multicolumn{1}{|c|}{01X0XX00} &
\multicolumn{1}{c|}{{\scriptsize rs2[4]}} &
\multicolumn{1}{c|}{{\scriptsize rs2[2:0]}} &
\multicolumn{1}{c|}{1} &
\multicolumn{1}{c|}{{\scriptsize rs1[4]}} &
\multicolumn{1}{c|}{{\scriptsize rs1[2:0]}} &
\multicolumn{1}{c|}{11} &
\multicolumn{1}{c|}{{\scriptsize qu[2:0]}} &
\multicolumn{1}{c|}{0X0} &
\multicolumn{1}{c|}{0011011}\\
\hline
8& 1& 3& 1& 1& 3& 2& 3& 3& 7 \\
\end{tabular}
 \end{center}
 }
 \vspace{-0.1in}
 \caption[]{ESP.VLD.L.64.XP} \end{figure*}


\newpage
\subsubsection{ESP.VST.128.IP}\label{instruction:ESPVST128IP}
\textbf{Syntax}\\
ESP.VST.128.IP qu,rs1,imm16

\textbf{Description}\\
This instruction stores the 128 bits in the register \lstinline{qu} to memory, with the pointer incremented by an immediate.\\

imm starts at -2048, ends at 2032, and steps in increments of 16.

\textbf{Operation}
\begin{lstlisting}
  MEM(rs1[31:0]) = qu[127:  0]
  rs1[31:0] = rs1[31:0] +imm
\end{lstlisting}

\textbf{Schedule}\\
\begin{longtable}[c]{|l|l|l|}
    \hline
    \rowcolor{lightgray}
    \textbf{Operands} & \textbf{Use} & \textbf{Def} \\\hline
    \endfirsthead
    \hline
    \rowcolor{lightgray}
     \textbf{Operands} & \textbf{Use} & \textbf{Def} \\\hline
    \endhead
    \hline
    \endfoot

  qu & 2 & - \\\hline
  rs1 & 1 & 1

\end{longtable}

\textbf{Format}\\

\begin{figure*}[!htbp]
 {\setlength{\tabcolsep}{1pt}
 \begin{center}
\begin{tabular}{@{}I@{}V@{}O@{}I@{}F@{}W@{}F@{}F@{}O}
\instbit{31}&
\instbitrange{30}{26}&
\instbitrange{25}{19}&
\instbit{18}&
\instbitrange{17}{15}&
\instbitrange{14}{13}&
\instbitrange{12}{10}&
\instbitrange{9}{7}&
\instbitrange{6}{0} \\
\hline
\multicolumn{1}{|c|}{1} &
\multicolumn{1}{c|}{{\scriptsize imm16[7:3]}} &
\multicolumn{1}{c|}{110001X} &
\multicolumn{1}{c|}{{\scriptsize rs1[4]}} &
\multicolumn{1}{c|}{{\scriptsize rs1[2:0]}} &
\multicolumn{1}{c|}{00} &
\multicolumn{1}{c|}{{\scriptsize qu[2:0]}} &
\multicolumn{1}{c|}{{\scriptsize imm16[2:0]}} &
\multicolumn{1}{c|}{0011111}\\
\hline
1& 5& 7& 1& 3& 2& 3& 3& 7 \\
\end{tabular}
 \end{center}
 }
 \vspace{-0.1in}
 \caption[]{ESP.VST.128.IP} \end{figure*}


\newpage
\subsubsection{ESP.VST.128.XP}\label{instruction:ESPVST128XP}
\textbf{Syntax}\\
ESP.VST.128.XP qu,rs1,rs2

\textbf{Description}\\
This instruction stores the 128 bits in the register \lstinline{qu} to memory, with the pointer incremented by the value in the register \lstinline{rs2}.\\

\textbf{Operation}
\begin{lstlisting}
  MEM(rs1[31:0]) = qu[127:  0]
  rs1[31:0] = rs1[31:0] + rs2[31:0]
\end{lstlisting}

\textbf{Schedule}\\
\begin{longtable}[c]{|l|l|l|}
    \hline
    \rowcolor{lightgray}
    \textbf{Operands} & \textbf{Use} & \textbf{Def} \\\hline
    \endfirsthead
    \hline
    \rowcolor{lightgray}
     \textbf{Operands} & \textbf{Use} & \textbf{Def} \\\hline
    \endhead
    \hline
    \endfoot

  qu & 2 & - \\\hline
  rs1 & 1 & 1

\end{longtable}

\textbf{Format}\\

\begin{figure*}[!htbp]
 {\setlength{\tabcolsep}{1pt}
 \begin{center}
\begin{tabular}{@{}E@{}I@{}F@{}I@{}I@{}F@{}W@{}F@{}F@{}O}
\instbitrange{31}{24}&
\instbit{23}&
\instbitrange{22}{20}&
\instbit{19}&
\instbit{18}&
\instbitrange{17}{15}&
\instbitrange{14}{13}&
\instbitrange{12}{10}&
\instbitrange{9}{7}&
\instbitrange{6}{0} \\
\hline
\multicolumn{1}{|c|}{01X0XX1X} &
\multicolumn{1}{c|}{{\scriptsize rs2[4]}} &
\multicolumn{1}{c|}{{\scriptsize rs2[2:0]}} &
\multicolumn{1}{c|}{0} &
\multicolumn{1}{c|}{{\scriptsize rs1[4]}} &
\multicolumn{1}{c|}{{\scriptsize rs1[2:0]}} &
\multicolumn{1}{c|}{11} &
\multicolumn{1}{c|}{{\scriptsize qu[2:0]}} &
\multicolumn{1}{c|}{1X0} &
\multicolumn{1}{c|}{0011011}\\
\hline
8& 1& 3& 1& 1& 3& 2& 3& 3& 7 \\
\end{tabular}
 \end{center}
 }
 \vspace{-0.1in}
 \caption[]{ESP.VST.128.XP} \end{figure*}


\newpage
\subsubsection{ESP.VST.H.64.IP}\label{instruction:ESPVSTH64IP}
\textbf{Syntax}\\
ESP.VST.H.64.IP qu,rs1,imm8

\textbf{Description}\\
This instruction stores the 64 bits in the upper 64 bits of the register \lstinline{qu} to memory, with the pointer incremented by an immediate.\\

imm starts at -1024, ends at 1016, and steps by 8.

\textbf{Operation}
\begin{lstlisting}
  MEM(rs1[31:0]) = qu[127:  64]
  rs1[31:0] = rs1[31:0] + imm
\end{lstlisting}

\textbf{Schedule}\\
\begin{longtable}[c]{|l|l|l|}
    \hline
    \rowcolor{lightgray}
    \textbf{Operands} & \textbf{Use} & \textbf{Def} \\\hline
    \endfirsthead
    \hline
    \rowcolor{lightgray}
     \textbf{Operands} & \textbf{Use} & \textbf{Def} \\\hline
    \endhead
    \hline
    \endfoot

  qu & 2 & - \\\hline
  rs1 & 1 & 1

\end{longtable}

\textbf{Format}\\

\begin{figure*}[!htbp]
 {\setlength{\tabcolsep}{1pt}
 \begin{center}
\begin{tabular}{@{}W@{}Y@{}W@{}I@{}Y@{}I@{}F@{}W@{}F@{}F@{}O}
\instbitrange{31}{30}&
\instbitrange{29}{26}&
\instbitrange{25}{24}&
\instbit{23}&
\instbitrange{22}{19}&
\instbit{18}&
\instbitrange{17}{15}&
\instbitrange{14}{13}&
\instbitrange{12}{10}&
\instbitrange{9}{7}&
\instbitrange{6}{0} \\
\hline
\multicolumn{1}{|c|}{11} &
\multicolumn{1}{c|}{{\scriptsize imm8[7:4]}} &
\multicolumn{1}{c|}{11} &
\multicolumn{1}{c|}{{\scriptsize imm8[3]}} &
\multicolumn{1}{c|}{110X} &
\multicolumn{1}{c|}{{\scriptsize rs1[4]}} &
\multicolumn{1}{c|}{{\scriptsize rs1[2:0]}} &
\multicolumn{1}{c|}{00} &
\multicolumn{1}{c|}{{\scriptsize qu[2:0]}} &
\multicolumn{1}{c|}{{\scriptsize imm8[2:0]}} &
\multicolumn{1}{c|}{0011111}\\
\hline
2& 4& 2& 1& 4& 1& 3& 2& 3& 3& 7 \\
\end{tabular}
 \end{center}
 }
 \vspace{-0.1in}
 \caption[]{ESP.VST.H.64.IP} \end{figure*}


\newpage
\subsubsection{ESP.VST.H.64.XP}\label{instruction:ESPVSTH64XP}
\textbf{Syntax}\\
ESP.VST.H.64.XP qu,rs1,rs2

\textbf{Description}\\
This instruction stores the 64 bits in the upper 64 bits of the register \lstinline{qu} to memory, with the pointer incremented by the value in the register \lstinline{rs2}.\\

\textbf{Operation}
\begin{lstlisting}
  MEM(rs1[31:0]) = qu[127:  64]
  rs1[31:0] = rs1[31:0] + rs2[31:0]
\end{lstlisting}

\textbf{Schedule}\\
\begin{longtable}[c]{|l|l|l|}
    \hline
    \rowcolor{lightgray}
    \textbf{Operands} & \textbf{Use} & \textbf{Def} \\\hline
    \endfirsthead
    \hline
    \rowcolor{lightgray}
     \textbf{Operands} & \textbf{Use} & \textbf{Def} \\\hline
    \endhead
    \hline
    \endfoot

  qu & 2 & - \\\hline
  rs1 & 1 & 1 \\\hline
  rs2 & 1 & -

\end{longtable}

\textbf{Format}\\

\begin{figure*}[!htbp]
 {\setlength{\tabcolsep}{1pt}
 \begin{center}
\begin{tabular}{@{}E@{}I@{}F@{}I@{}I@{}F@{}W@{}F@{}F@{}O}
\instbitrange{31}{24}&
\instbit{23}&
\instbitrange{22}{20}&
\instbit{19}&
\instbit{18}&
\instbitrange{17}{15}&
\instbitrange{14}{13}&
\instbitrange{12}{10}&
\instbitrange{9}{7}&
\instbitrange{6}{0} \\
\hline
\multicolumn{1}{|c|}{01X0XX11} &
\multicolumn{1}{c|}{{\scriptsize rs2[4]}} &
\multicolumn{1}{c|}{{\scriptsize rs2[2:0]}} &
\multicolumn{1}{c|}{1} &
\multicolumn{1}{c|}{{\scriptsize rs1[4]}} &
\multicolumn{1}{c|}{{\scriptsize rs1[2:0]}} &
\multicolumn{1}{c|}{11} &
\multicolumn{1}{c|}{{\scriptsize qu[2:0]}} &
\multicolumn{1}{c|}{0X0} &
\multicolumn{1}{c|}{0011011}\\
\hline
8& 1& 3& 1& 1& 3& 2& 3& 3& 7 \\
\end{tabular}
 \end{center}
 }
 \vspace{-0.1in}
 \caption[]{ESP.VST.H.64.XP} \end{figure*}


\newpage
\subsubsection{ESP.VST.L.64.IP}\label{instruction:ESPVSTL64IP}
\textbf{Syntax}\\
ESP.VST.L.64.IP qu,rs1,imm8

\textbf{Description}\\
This instruction stores the 64 bits in the lower 64 bits of the register \lstinline{qu} to memory, with the pointer incremented by an immediate.\\

imm starts at -1024, ends at 1016, and steps by 8.

\textbf{Operation}
\begin{lstlisting}
  MEM(rs1[31:0]) = qu[63:  0]
  rs1[31:0] = rs1[31:0] + imm
\end{lstlisting}

\textbf{Schedule}\\
\begin{longtable}[c]{|l|l|l|}
    \hline
    \rowcolor{lightgray}
    \textbf{Operands} & \textbf{Use} & \textbf{Def} \\\hline
    \endfirsthead
    \hline
    \rowcolor{lightgray}
     \textbf{Operands} & \textbf{Use} & \textbf{Def} \\\hline
    \endhead
    \hline
    \endfoot

  qu & 2 & - \\\hline
  rs1 & 1 & 1

\end{longtable}

\textbf{Format}\\

\begin{figure*}[!htbp]
 {\setlength{\tabcolsep}{1pt}
 \begin{center}
\begin{tabular}{@{}W@{}Y@{}W@{}I@{}Y@{}I@{}F@{}W@{}F@{}F@{}O}
\instbitrange{31}{30}&
\instbitrange{29}{26}&
\instbitrange{25}{24}&
\instbit{23}&
\instbitrange{22}{19}&
\instbit{18}&
\instbitrange{17}{15}&
\instbitrange{14}{13}&
\instbitrange{12}{10}&
\instbitrange{9}{7}&
\instbitrange{6}{0} \\
\hline
\multicolumn{1}{|c|}{10} &
\multicolumn{1}{c|}{{\scriptsize imm8[7:4]}} &
\multicolumn{1}{c|}{11} &
\multicolumn{1}{c|}{{\scriptsize imm8[3]}} &
\multicolumn{1}{c|}{110X} &
\multicolumn{1}{c|}{{\scriptsize rs1[4]}} &
\multicolumn{1}{c|}{{\scriptsize rs1[2:0]}} &
\multicolumn{1}{c|}{00} &
\multicolumn{1}{c|}{{\scriptsize qu[2:0]}} &
\multicolumn{1}{c|}{{\scriptsize imm8[2:0]}} &
\multicolumn{1}{c|}{0011111}\\
\hline
2& 4& 2& 1& 4& 1& 3& 2& 3& 3& 7 \\
\end{tabular}
 \end{center}
 }
 \vspace{-0.1in}
 \caption[]{ESP.VST.L.64.IP} \end{figure*}


\newpage
\subsubsection{ESP.VST.L.64.XP}\label{instruction:ESPVSTL64XP}
\textbf{Syntax}\\
ESP.VST.L.64.XP qu,rs1,rs2

\textbf{Description}\\
This instruction stores the 64 bits in the lower 64 bits of the register \lstinline{qu} to memory, with the pointer incremented by the value in the register \lstinline{rs2}.\\

\textbf{Operation}
\begin{lstlisting}
  MEM(rs1[31:0]) = qu[63:  0]
  rs1[31:0] = rs1[31:0] + rs2[31:0]
\end{lstlisting}

\textbf{Schedule}\\
\begin{longtable}[c]{|l|l|l|}
    \hline
    \rowcolor{lightgray}
    \textbf{Operands} & \textbf{Use} & \textbf{Def} \\\hline
    \endfirsthead
    \hline
    \rowcolor{lightgray}
     \textbf{Operands} & \textbf{Use} & \textbf{Def} \\\hline
    \endhead
    \hline
    \endfoot

  qu & 2 & - \\\hline
  rs1 & 1 & 1 \\\hline
  rs2 & 1 & -

\end{longtable}

\textbf{Format}\\

\begin{figure*}[!htbp]
 {\setlength{\tabcolsep}{1pt}
 \begin{center}
\begin{tabular}{@{}E@{}I@{}F@{}I@{}I@{}F@{}W@{}F@{}F@{}O}
\instbitrange{31}{24}&
\instbit{23}&
\instbitrange{22}{20}&
\instbit{19}&
\instbit{18}&
\instbitrange{17}{15}&
\instbitrange{14}{13}&
\instbitrange{12}{10}&
\instbitrange{9}{7}&
\instbitrange{6}{0} \\
\hline
\multicolumn{1}{|c|}{01X0XX10} &
\multicolumn{1}{c|}{{\scriptsize rs2[4]}} &
\multicolumn{1}{c|}{{\scriptsize rs2[2:0]}} &
\multicolumn{1}{c|}{1} &
\multicolumn{1}{c|}{{\scriptsize rs1[4]}} &
\multicolumn{1}{c|}{{\scriptsize rs1[2:0]}} &
\multicolumn{1}{c|}{11} &
\multicolumn{1}{c|}{{\scriptsize qu[2:0]}} &
\multicolumn{1}{c|}{0X0} &
\multicolumn{1}{c|}{0011011}\\
\hline
8& 1& 3& 1& 1& 3& 2& 3& 3& 7 \\
\end{tabular}
 \end{center}
 }
 \vspace{-0.1in}
 \caption[]{ESP.VST.L.64.XP} \end{figure*}


\newpage
\subsubsection{ESP.SLCI.2Q}\label{instruction:ESPSLCI2Q}
\textbf{Syntax}\\
ESP.SLCI.2Q qy,qw,sel16

\textbf{Description}\\
This instruction performs a left shift on the 32-byte concatenation of registers \lstinline{qw} and \lstinline{qy}, and pads the lower bits with 0. The upper 128 bits of the shift result are written to register \lstinline{qy}, and the lower 128 bits are written to \lstinline{qw}. The left shift amount is 8 times the sum of \lstinline{sel16} and 1.\\
\textbf{Operation}
\begin{lstlisting}
  {qy[127:  0], qw[127:  0]} = {qy[127:  0], qw[127:  0]} << ((sel16[3:0]+1)*8)
\end{lstlisting}

\textbf{Schedule}\\
\begin{longtable}[c]{|l|l|l|}
    \hline
    \rowcolor{lightgray}
    \textbf{Operands} & \textbf{Use} & \textbf{Def} \\\hline
    \endfirsthead
    \hline
    \rowcolor{lightgray}
     \textbf{Operands} & \textbf{Use} & \textbf{Def} \\\hline
    \endhead
    \hline
    \endfoot

  qy & 1 & 1 \\\hline
  qw & 1 & 1

\end{longtable}

\textbf{Format}\\

\begin{figure*}[!htbp]
 {\setlength{\tabcolsep}{1pt}
 \begin{center}
\begin{tabular}{@{}I@{}W@{}F@{}W@{}Y@{}I@{}E@{}W@{}W@{}O}
\instbit{31}&
\instbitrange{30}{29}&
\instbitrange{28}{26}&
\instbitrange{25}{24}&
\instbitrange{23}{20}&
\instbit{19}&
\instbitrange{18}{11}&
\instbitrange{10}{9}&
\instbitrange{8}{7}&
\instbitrange{6}{0} \\
\hline
\multicolumn{1}{|c|}{0} &
\multicolumn{1}{c|}{{\scriptsize sel16[3:2]}} &
\multicolumn{1}{c|}{{\scriptsize qy[2:0]}} &
\multicolumn{1}{c|}{{\scriptsize qw[1:0]}} &
\multicolumn{1}{c|}{01XX} &
\multicolumn{1}{c|}{{\scriptsize qw[2]}} &
\multicolumn{1}{c|}{0XXX01XX} &
\multicolumn{1}{c|}{{\scriptsize sel16[1:0]}} &
\multicolumn{1}{c|}{0X} &
\multicolumn{1}{c|}{0011011}\\
\hline
1& 2& 3& 2& 4& 1& 8& 2& 2& 7 \\
\end{tabular}
 \end{center}
 }
 \vspace{-0.1in}
 \caption[]{ESP.SLCI.2Q} \end{figure*}


\newpage
\subsubsection{ESP.SLCXXP.2Q}\label{instruction:ESPSLCXXP2Q}
\textbf{Syntax}\\
ESP.SLCXXP.2Q qy,qw,rs1,rs2

\textbf{Description}\\
This instruction performs a left shift on the 32-byte concatenation of registers \lstinline{qy} and \lstinline{qw}, and pads the lower bits with 0. The upper 128 bits of the shift result are written to register \lstinline{qy}, and the lower 128 bits are written to \lstinline{qw}. The left shift amount is 8, performed by the sum of 1 plus the lower 4-bit value of register \lstinline{rs1}. After the above operations, the value in \lstinline{rs1} is incremented by the value in \lstinline{rs2}.\\
\textbf{Operation}
\begin{lstlisting}
  {qy[127:  0], qw[127:  0]} = {qy[127:  0], qw[127:  0]} << ((rs1[3:0]+1)*8)
  rs1[31:0] = rs1[31:0] + rs2[31:0]
\end{lstlisting}

\textbf{Schedule}\\
\begin{longtable}[c]{|l|l|l|}
    \hline
    \rowcolor{lightgray}
    \textbf{Operands} & \textbf{Use} & \textbf{Def} \\\hline
    \endfirsthead
    \hline
    \rowcolor{lightgray}
     \textbf{Operands} & \textbf{Use} & \textbf{Def} \\\hline
    \endhead
    \hline
    \endfoot

  qy & 1 & 1 \\\hline
  qw & 1 & 1 \\\hline
  rs1 & 1 & 1 \\\hline
  rs2 & 1 & -

\end{longtable}

\textbf{Format}\\

\begin{figure*}[!htbp]
 {\setlength{\tabcolsep}{1pt}
 \begin{center}
\begin{tabular}{@{}F@{}F@{}W@{}I@{}F@{}I@{}I@{}F@{}E@{}O}
\instbitrange{31}{29}&
\instbitrange{28}{26}&
\instbitrange{25}{24}&
\instbit{23}&
\instbitrange{22}{20}&
\instbit{19}&
\instbit{18}&
\instbitrange{17}{15}&
\instbitrange{14}{7}&
\instbitrange{6}{0} \\
\hline
\multicolumn{1}{|c|}{X0X} &
\multicolumn{1}{c|}{{\scriptsize qy[2:0]}} &
\multicolumn{1}{c|}{{\scriptsize qw[1:0]}} &
\multicolumn{1}{c|}{{\scriptsize rs2[4]}} &
\multicolumn{1}{c|}{{\scriptsize rs2[2:0]}} &
\multicolumn{1}{c|}{{\scriptsize qw[2]}} &
\multicolumn{1}{c|}{{\scriptsize rs1[4]}} &
\multicolumn{1}{c|}{{\scriptsize rs1[2:0]}} &
\multicolumn{1}{c|}{11XX0XXX} &
\multicolumn{1}{c|}{0011011}\\
\hline
3& 3& 2& 1& 3& 1& 1& 3& 8& 7 \\
\end{tabular}
 \end{center}
 }
 \vspace{-0.1in}
 \caption[]{ESP.SLCXXP.2Q} \end{figure*}


\newpage
\subsubsection{ESP.SRC.Q}\label{instruction:ESPSRCQ}
\textbf{Syntax}\\
ESP.SRC.Q qz,qw,qy

\textbf{Description}\\
This instruction performs an arithmetic right shift on the 32-byte concatenation of registers \lstinline{qy} and \lstinline{qw}, which hold the loaded data from two consecutive aligned addresses. In this way, you can obtain unaligned 16-byte data, which will then be written to the register \lstinline{qw}. The right shift amount is \lstinline{SAR_BYTE} performed by 8.\\
\textbf{Operation}
\begin{lstlisting}
  qz[127:0] = {qy[127: 0], qw[127: 0]} >> {SAR_BYTE[3:0] << 3}
\end{lstlisting}

\textbf{Schedule}\\
\begin{longtable}[c]{|l|l|l|}
    \hline
    \rowcolor{lightgray}
    \textbf{Operands} & \textbf{Use} & \textbf{Def} \\\hline
    \endfirsthead
    \hline
    \rowcolor{lightgray}
     \textbf{Operands} & \textbf{Use} & \textbf{Def} \\\hline
    \endhead
    \hline
    \endfoot

  qz & - & 1 \\\hline
  qw & 1 & - \\\hline
  qy & 1 & - \\\hline
  SAR\_BYTES & 1 & -

\end{longtable}

\textbf{Format}\\

\begin{figure*}[!htbp]
 {\setlength{\tabcolsep}{1pt}
 \begin{center}
\begin{tabular}{@{}F@{}F@{}W@{}Y@{}I@{}T@{}F@{}O}
\instbitrange{31}{29}&
\instbitrange{28}{26}&
\instbitrange{25}{24}&
\instbitrange{23}{20}&
\instbit{19}&
\instbitrange{18}{10}&
\instbitrange{9}{7}&
\instbitrange{6}{0} \\
\hline
\multicolumn{1}{|c|}{1XX} &
\multicolumn{1}{c|}{{\scriptsize qy[2:0]}} &
\multicolumn{1}{c|}{{\scriptsize qw[1:0]}} &
\multicolumn{1}{c|}{X0XX} &
\multicolumn{1}{c|}{{\scriptsize qw[2]}} &
\multicolumn{1}{c|}{0XXX100XX} &
\multicolumn{1}{c|}{{\scriptsize qz[2:0]}} &
\multicolumn{1}{c|}{0011011}\\
\hline
3& 3& 2& 4& 1& 9& 3& 7 \\
\end{tabular}
 \end{center}
 }
 \vspace{-0.1in}
 \caption[]{ESP.SRC.Q} \end{figure*}


\newpage
\subsubsection{ESP.SRC.Q.LD.IP}\label{instruction:ESPSRCQLDIP}
\textbf{Syntax}\\
ESP.SRC.Q.LD.IP qu,rs1,imm16,qw,qy

\textbf{Description}\\
This instruction performs an arithmetic right shift on the 32-byte concatenation of registers \lstinline{qy} and \lstinline{qw}, which hold the loaded data from two consecutive aligned addresses. In this way, you can obtain unaligned 16-byte data, which will be written to the register \lstinline{qw}. The right shift amount is \lstinline{SAR_BYTE} performed by 8.\\
Simultaneously, the 16-byte data is loaded from the memory into the register \lstinline{qu}. After the access is completed, the value in the register \lstinline{rs1} is incremented by \lstinline{imm}.\\

\lstinline{imm} starts at -2048, ends at 2032, with a step of 16.

\textbf{Operation}
\begin{lstlisting}
  qw[127:0] = {qy[127: 0], qw[127: 0]} >> {SAR_BYTE[3:0] << 3}
  qu[127:0] = load128(rs1[31:0])
  rs1[31:0] = rs1[31:0] + imm
\end{lstlisting}

\textbf{Schedule}\\
\begin{longtable}[c]{|l|l|l|}
    \hline
    \rowcolor{lightgray}
    \textbf{Operands} & \textbf{Use} & \textbf{Def} \\\hline
    \endfirsthead
    \hline
    \rowcolor{lightgray}
     \textbf{Operands} & \textbf{Use} & \textbf{Def} \\\hline
    \endhead
    \hline
    \endfoot

  qz & - & 1 \\\hline
  qw & 1 & - \\\hline
  qy & 1 & - \\\hline
  SAR\_BYTES & 1 & - \\\hline
  qu & - & 2 \\\hline
  rs1 & 1 & 1

\end{longtable}

\textbf{Format}\\

\begin{figure*}[!htbp]
 {\setlength{\tabcolsep}{1pt}
 \begin{center}
\begin{tabular}{@{}F@{}F@{}W@{}Y@{}I@{}I@{}F@{}W@{}F@{}I@{}W@{}O}
\instbitrange{31}{29}&
\instbitrange{28}{26}&
\instbitrange{25}{24}&
\instbitrange{23}{20}&
\instbit{19}&
\instbit{18}&
\instbitrange{17}{15}&
\instbitrange{14}{13}&
\instbitrange{12}{10}&
\instbit{9}&
\instbitrange{8}{7}&
\instbitrange{6}{0} \\
\hline
\multicolumn{1}{|c|}{{\scriptsize imm16[7:5]}} &
\multicolumn{1}{c|}{{\scriptsize qy[2:0]}} &
\multicolumn{1}{c|}{{\scriptsize qw[1:0]}} &
\multicolumn{1}{c|}{{\scriptsize imm16[4:1]}} &
\multicolumn{1}{c|}{{\scriptsize qw[2]}} &
\multicolumn{1}{c|}{{\scriptsize rs1[4]}} &
\multicolumn{1}{c|}{{\scriptsize rs1[2:0]}} &
\multicolumn{1}{c|}{10} &
\multicolumn{1}{c|}{{\scriptsize qu[2:0]}} &
\multicolumn{1}{c|}{{\scriptsize imm16[0]}} &
\multicolumn{1}{c|}{00} &
\multicolumn{1}{c|}{0011111}\\
\hline
3& 3& 2& 4& 1& 1& 3& 2& 3& 1& 2& 7 \\
\end{tabular}
 \end{center}
 }
 \vspace{-0.1in}
 \caption[]{ESP.SRC.Q.LD.IP} \end{figure*}


\newpage
\subsubsection{ESP.SRC.Q.LD.XP}\label{instruction:ESPSRCQLDXP}
\textbf{Syntax}\\
ESP.SRC.Q.LD.XP qu,rs1,rs2,qw,qy

\textbf{Description}\\
This instruction performs an arithmetic right shift on the 32-byte concatenation of registers \lstinline{qy} and \lstinline{qw}, which hold the loaded data from two consecutive aligned addresses. In this way, you can obtain unaligned 16-byte data, which will be written to the register \lstinline{qw}. The right shift amount is the product of \lstinline{SAR_BYTE} and 8.\\
Simultaneously, the 16-byte data is loaded from memory into the register \lstinline{qu}. After the access is completed, the value in the register \lstinline{rs1} is incremented by the value in the register \lstinline{rs2}.\\
\textbf{Operation}
\begin{lstlisting}
  qw[127:0] = {qy[127: 0], qw[127: 0]} >> {SAR_BYTE[3:0] << 3}
  qu[127:0] = load128(rs1[31:0])
  rs1[31:0] = rs1[31:0] + rs2[31:0]
\end{lstlisting}

\textbf{Schedule}\\
\begin{longtable}[c]{|l|l|l|}
    \hline
    \rowcolor{lightgray}
    \textbf{Operands} & \textbf{Use} & \textbf{Def} \\\hline
    \endfirsthead
    \hline
    \rowcolor{lightgray}
     \textbf{Operands} & \textbf{Use} & \textbf{Def} \\\hline
    \endhead
    \hline
    \endfoot

  qz & - & 1 \\\hline
  qw & 1 & - \\\hline
  qy & 1 & - \\\hline
  SAR\_BYTES & 1 & - \\\hline
  qu & - & 2 \\\hline
  rs1 & 1 & 1 \\\hline
  rs2 & 1 & -

\end{longtable}

\textbf{Format}\\

\begin{figure*}[!htbp]
 {\setlength{\tabcolsep}{1pt}
 \begin{center}
\begin{tabular}{@{}F@{}F@{}W@{}I@{}F@{}I@{}I@{}F@{}W@{}F@{}F@{}O}
\instbitrange{31}{29}&
\instbitrange{28}{26}&
\instbitrange{25}{24}&
\instbit{23}&
\instbitrange{22}{20}&
\instbit{19}&
\instbit{18}&
\instbitrange{17}{15}&
\instbitrange{14}{13}&
\instbitrange{12}{10}&
\instbitrange{9}{7}&
\instbitrange{6}{0} \\
\hline
\multicolumn{1}{|c|}{11X} &
\multicolumn{1}{c|}{{\scriptsize qy[2:0]}} &
\multicolumn{1}{c|}{{\scriptsize qw[1:0]}} &
\multicolumn{1}{c|}{{\scriptsize rs2[4]}} &
\multicolumn{1}{c|}{{\scriptsize rs2[2:0]}} &
\multicolumn{1}{c|}{{\scriptsize qw[2]}} &
\multicolumn{1}{c|}{{\scriptsize rs1[4]}} &
\multicolumn{1}{c|}{{\scriptsize rs1[2:0]}} &
\multicolumn{1}{c|}{11} &
\multicolumn{1}{c|}{{\scriptsize qu[2:0]}} &
\multicolumn{1}{c|}{11X} &
\multicolumn{1}{c|}{0011011}\\
\hline
3& 3& 2& 1& 3& 1& 1& 3& 2& 3& 3& 7 \\
\end{tabular}
 \end{center}
 }
 \vspace{-0.1in}
 \caption[]{ESP.SRC.Q.LD.XP} \end{figure*}


\newpage
\subsubsection{ESP.SRCI.2Q}\label{instruction:ESPSRCI2Q}
\textbf{Syntax}\\
ESP.SRCI.2Q qy,qw,sel16

\textbf{Description}\\
This instruction performs a logical right shift on the 32-byte concatenation of registers \lstinline{qw} and \lstinline{qy}, and pads the higher bits with 0. The upper 128 bits of the shift result are written to the register \lstinline{qy}, and the lower 128 bits are written to \lstinline{qw}. The right shift amount is 8 times the sum of \lstinline{sel16} and 1.\\
\textbf{Operation}
\begin{lstlisting}
  {qy[127:  0], qw[127:  0]} = {qy[127:  0], qw[127:  0]} >> ((sel16+1)*8)
  qy[127:127-8*sel16] = 0
\end{lstlisting}

\textbf{Schedule}\\
\begin{longtable}[c]{|l|l|l|}
    \hline
    \rowcolor{lightgray}
    \textbf{Operands} & \textbf{Use} & \textbf{Def} \\\hline
    \endfirsthead
    \hline
    \rowcolor{lightgray}
     \textbf{Operands} & \textbf{Use} & \textbf{Def} \\\hline
    \endhead
    \hline
    \endfoot

  qz & 1 & 1 \\\hline
  qw & 1 & 1

\end{longtable}

\textbf{Format}\\

\begin{figure*}[!htbp]
 {\setlength{\tabcolsep}{1pt}
 \begin{center}
\begin{tabular}{@{}I@{}W@{}F@{}W@{}Y@{}I@{}E@{}W@{}W@{}O}
\instbit{31}&
\instbitrange{30}{29}&
\instbitrange{28}{26}&
\instbitrange{25}{24}&
\instbitrange{23}{20}&
\instbit{19}&
\instbitrange{18}{11}&
\instbitrange{10}{9}&
\instbitrange{8}{7}&
\instbitrange{6}{0} \\
\hline
\multicolumn{1}{|c|}{1} &
\multicolumn{1}{c|}{{\scriptsize sel16[3:2]}} &
\multicolumn{1}{c|}{{\scriptsize qy[2:0]}} &
\multicolumn{1}{c|}{{\scriptsize qw[1:0]}} &
\multicolumn{1}{c|}{01XX} &
\multicolumn{1}{c|}{{\scriptsize qw[2]}} &
\multicolumn{1}{c|}{0XXX01XX} &
\multicolumn{1}{c|}{{\scriptsize sel16[1:0]}} &
\multicolumn{1}{c|}{0X} &
\multicolumn{1}{c|}{0011011}\\
\hline
1& 2& 3& 2& 4& 1& 8& 2& 2& 7 \\
\end{tabular}
 \end{center}
 }
 \vspace{-0.1in}
 \caption[]{ESP.SRCI.2Q} \end{figure*}


\newpage
\subsubsection{ESP.SRCMB.S16.Q.QACC}\label{instruction:ESPSRCMBS16QQACC}
\textbf{Syntax}\\
ESP.SRCMB.S16.Q.QACC qu,qw,sat,rm

\textbf{Description}\\
This instruction extracts 8 data segments of 64 bits from special registers \lstinline{QACC_H} and \lstinline{QACC_L} and performs arithmetic right shift operations respectively, the shift amounts are from \lstinline{qw}. While writing the shift results back to \lstinline{QACC_H} and \lstinline{QACC_L}, the instruction saturates the results to 16-bit signed numbers and writes the 8 16-bit data obtained after saturation into the register \lstinline{qu}.\\
\textbf{Operation}
\begin{lstlisting}[language=C, escapechar=@]
  shf_value0[5:0] = @(qw[ 15:  0] >= 0x3f) ? 0x3f : qw[ 15:  0]@
  shf_value1[5:0] = @(qw[ 31: 16] >= 0x3f) ? 0x3f : qw[ 31: 16]@
  shf_value2[5:0] = @(qw[ 47: 32] >= 0x3f) ? 0x3f : qw[ 47: 32]@
  shf_value3[5:0] = @(qw[ 63: 48] >= 0x3f) ? 0x3f : qw[ 63: 48]@
  shf_value4[5:0] = @(qw[ 79: 64] >= 0x3f) ? 0x3f : qw[ 79: 64]@
  shf_value5[5:0] = @(qw[ 95: 80] >= 0x3f) ? 0x3f : qw[ 95: 80]@
  shf_value6[5:0] = @(qw[111: 96] >= 0x3f) ? 0x3f : qw[111: 96]@
  shf_value7[5:0] = @(qw[127:112] >= 0x3f) ? 0x3f : qw[127:112]@

  DATA_L[ 63:  0] = QACC_L[ 63:  0] >> shf_value0[5:0]
  DATA_L[127: 64] = QACC_L[127: 64] >> shf_value1[5:0]
  DATA_L[191:128] = QACC_L[191:128] >> shf_value2[5:0]
  DATA_L[255:192] = QACC_L[255:192] >> shf_value3[5:0]
  DATA_H[ 63:  0] = QACC_H[ 63:  0] >> shf_value4[5:0]
  DATA_H[127: 64] = QACC_H[127: 64] >> shf_value5[5:0]
  DATA_H[191:128] = QACC_H[191:128] >> shf_value6[5:0]
  DATA_H[255:192] = QACC_H[255:192] >> shf_value7[5:0]

  QACC_L[255:0] = DATA_L[255:0]
  QACC_H[255:0] = DATA_H[255:0]

  if sel2 == 0:
    qu[15: 0] = DATA_L[15:0]
    qu[31:16] = DATA_L[79:64]
    qu[47:32] = DATA_L[143:128]
    qu[63:48] = DATA_L[207:192]
    qu[79:64] = DATA_H[15:0]
    qu[95:80] = DATA_H[79:64]
    qu[111:96] = DATA_H[143:128]
    qu[127:112] = DATA_H[207:192]
  else if sel2 == 1:
    qu[ 15:  0] = min(max(DATA_L[63:0], -2^{15}), 2^{15}-1)
    qu[ 31: 16] = min(max(DATA_L[127:64], -2^{15}), 2^{15}-1)
    qu[ 47: 32] = min(max(DATA_L[191:128], -2^{15}), 2^{15}-1)
    qu[ 63: 48] = min(max(DATA_L[255:192], -2^{15}), 2^{15}-1)
    qu[ 79: 64] = min(max(DATA_H[63:0], -2^{15}), 2^{15}-1)
    qu[ 95: 80] = min(max(DATA_H[127:64], -2^{15}), 2^{15}-1)
    qu[111: 96] = min(max(DATA_H[191:128], -2^{15}), 2^{15}-1)
    qu[127:112] = min(max(DATA_H[255:192], -2^{15}), 2^{15}-1)
\end{lstlisting}

\textbf{Schedule}\\
\begin{longtable}[c]{|l|l|l|}
    \hline
    \rowcolor{lightgray}
    \textbf{Operands} & \textbf{Use} & \textbf{Def} \\\hline
    \endfirsthead
    \hline
    \rowcolor{lightgray}
     \textbf{Operands} & \textbf{Use} & \textbf{Def} \\\hline
    \endhead
    \hline
    \endfoot

  qu & - & 2 \\\hline
  rs1 & 1 & - \\\hline
  QACC\_H & 1 & 2 \\\hline
  QACC\_L & 1 & 2

\end{longtable}

\textbf{Format}\\

\begin{figure*}[!htbp]
 {\setlength{\tabcolsep}{1pt}
 \begin{center}
\begin{tabular}{@{}V@{}I@{}W@{}I@{}F@{}I@{}S@{}F@{}F@{}O}
\instbitrange{31}{27}&
\instbit{26}&
\instbitrange{25}{24}&
\instbit{23}&
\instbitrange{22}{20}&
\instbit{19}&
\instbitrange{18}{13}&
\instbitrange{12}{10}&
\instbitrange{9}{7}&
\instbitrange{6}{0} \\
\hline
\multicolumn{1}{|c|}{1XX11} &
\multicolumn{1}{c|}{{\scriptsize sat[0]}} &
\multicolumn{1}{c|}{{\scriptsize qw[1:0]}} &
\multicolumn{1}{c|}{0} &
\multicolumn{1}{c|}{{\scriptsize rm[2:0]}} &
\multicolumn{1}{c|}{{\scriptsize qw[2]}} &
\multicolumn{1}{c|}{1XXX01} &
\multicolumn{1}{c|}{{\scriptsize qu[2:0]}} &
\multicolumn{1}{c|}{1XX} &
\multicolumn{1}{c|}{0011011}\\
\hline
5& 1& 2& 1& 3& 1& 6& 3& 3& 7 \\
\end{tabular}
 \end{center}
 }
 \vspace{-0.1in}
 \caption[]{ESP.SRCMB.S16.Q.QACC} \end{figure*}


\newpage
\subsubsection{ESP.SRCMB.S16.QACC}\label{instruction:ESPSRCMBS16QACC}
\textbf{Syntax}\\
ESP.SRCMB.S16.QACC qu,rs1,sat,rm

\textbf{Description}\\
This instruction extracts 8 data segments of 64 bits from special registers \lstinline{QACC_H} and \lstinline{QACC_L} and performs arithmetic right shift operations respectively. While writing the shift results back to \lstinline{QACC_H} and \lstinline{QACC_L}, the instruction saturates the results to 16-bit signed numbers and writes the 8 16-bit data segments obtained after saturation into the \lstinline{qu} register.\\
\textbf{Operation}
\begin{lstlisting}
  temp0[63:0] = QACC_L[ 63:  0]
  ...
  temp3[63:0] = QACC_L[255:192]
  temp4[63:0] = QACC_H[ 63:  0]
  ...
  temp7[63:0] = QACC_H[255:192]

  temp_shf0[63:0] = temp0[63:0] >> rs1[5:0]
  temp_shf1[63:0] = temp1[63:0] >> rs1[5:0]
  ...
  temp_shf7[63:0] = temp7[63:0] >> rs1[5:0]

  QACC_L[ 63:  0] = temp_shf0[63:0]
  ...
  QACC_L[255:192] = temp_shf3[63:0]
  QACC_H[ 63:  0] = temp_shf4[63:0]
  ...
  QACC_H[255:192] = temp_shf7[63:0]

  if sel2 == 0:
    qu[ 15:  0] = temp0[15:0]
    qu[ 31: 16] = temp1[15:0]
    qu[ 47: 32] = temp2[15:0]
    qu[ 63: 48] = temp3[15:0]
    qu[ 79: 64] = temp4[15:0]
    qu[ 95: 80] = temp5[15:0]
    qu[111: 96] = temp6[15:0]
    qu[127:112] = temp7[15:0]
  else if sel2==1:
    qu[ 15:  0] = min(max(temp_shf0[63:0], -2^{15}), 2^{15}-1)
    ...
    qu[ 63: 48] = min(max(temp_shf3[63:0], -2^{15}), 2^{15}-1)
    qu[ 79: 64] = min(max(temp_shf4[63:0], -2^{15}), 2^{15}-1)
    ...
    qu[127:112] = min(max(temp_shf7[63:0], -2^{15}), 2^{15}-1)
\end{lstlisting}

\textbf{Schedule}\\
\begin{longtable}[c]{|l|l|l|}
    \hline
    \rowcolor{lightgray}
    \textbf{Operands} & \textbf{Use} & \textbf{Def} \\\hline
    \endfirsthead
    \hline
    \rowcolor{lightgray}
     \textbf{Operands} & \textbf{Use} & \textbf{Def} \\\hline
    \endhead
    \hline
    \endfoot

  qu & - & 2 \\\hline
  rs1 & 1 & - \\\hline
  QACC\_H & 1 & 2 \\\hline
  QACC\_L & 1 & 2

\end{longtable}

\textbf{Format}\\

\begin{figure*}[!htbp]
 {\setlength{\tabcolsep}{1pt}
 \begin{center}
\begin{tabular}{@{}W@{}I@{}F@{}O@{}I@{}F@{}W@{}F@{}F@{}O}
\instbitrange{31}{30}&
\instbit{29}&
\instbitrange{28}{26}&
\instbitrange{25}{19}&
\instbit{18}&
\instbitrange{17}{15}&
\instbitrange{14}{13}&
\instbitrange{12}{10}&
\instbitrange{9}{7}&
\instbitrange{6}{0} \\
\hline
\multicolumn{1}{|c|}{11} &
\multicolumn{1}{c|}{{\scriptsize sat[0]}} &
\multicolumn{1}{c|}{{\scriptsize rm[2:0]}} &
\multicolumn{1}{c|}{111010X} &
\multicolumn{1}{c|}{{\scriptsize rs1[4]}} &
\multicolumn{1}{c|}{{\scriptsize rs1[2:0]}} &
\multicolumn{1}{c|}{00} &
\multicolumn{1}{c|}{{\scriptsize qu[2:0]}} &
\multicolumn{1}{c|}{XX1} &
\multicolumn{1}{c|}{0011111}\\
\hline
2& 1& 3& 7& 1& 3& 2& 3& 3& 7 \\
\end{tabular}
 \end{center}
 }
 \vspace{-0.1in}
 \caption[]{ESP.SRCMB.S16.QACC} \end{figure*}


\newpage
\subsubsection{ESP.SRCMB.S8.Q.QACC}\label{instruction:ESPSRCMBS8QQACC}
\textbf{Syntax}\\
ESP.SRCMB.S8.Q.QACC qu,qw,sat,rm

\textbf{Description}\\
This instruction extracts 16 data segments of 32 bits from special registers \lstinline{QACC_H} and \lstinline{QACC_L} and performs arithmetic right shift operations respectively, the shift amounts are from \lstinline{qw}. While writing the shift results back to \lstinline{QACC_H} and \lstinline{QACC_L}, the instruction saturates the results to 8-bit signed numbers and writes the 16 8-bit data obtained after saturation into the register \lstinline{qu}.\\
\textbf{Operation}
\begin{lstlisting}[language=C, escapechar=@]
  shf_value0[5:0] = @(qw[  7:  0] >= 0x1f) ? 0x1f : qw[  7:  0]@
  shf_value1[5:0] = @(qw[ 15:  8] >= 0x1f) ? 0x1f : qw[ 15:  8]@
  shf_value2[5:0] = @(qw[ 23: 16] >= 0x1f) ? 0x1f : qw[ 23: 16]@
  shf_value3[5:0] = @(qw[ 31: 24] >= 0x1f) ? 0x1f : qw[ 31: 24]@
  shf_value4[5:0] = @(qw[ 39: 32] >= 0x1f) ? 0x1f : qw[ 39: 32]@
  shf_value5[5:0] = @(qw[ 47: 40] >= 0x1f) ? 0x1f : qw[ 47: 40]@
  shf_value6[5:0] = @(qw[ 55: 48] >= 0x1f) ? 0x1f : qw[ 55: 48]@
  shf_value7[5:0] = @(qw[ 63: 56] >= 0x1f) ? 0x1f : qw[ 63: 56]@
  shf_value8[5:0] = @(qw[ 71: 64] >= 0x1f) ? 0x1f : qw[ 71: 64]@
  shf_value9[5:0] = @(qw[ 79: 72] >= 0x1f) ? 0x1f : qw[ 79: 72]@
  shf_value10[5:0] = @(qw[ 87: 80] >= 0x1f) ? 0x1f : qw[ 87: 80]@
  shf_value11[5:0] = @(qw[ 95: 88] >= 0x1f) ? 0x1f : qw[ 95: 88]@
  shf_value12[5:0] = @(qw[103: 96] >= 0x1f) ? 0x1f : qw[103: 96]@
  shf_value13[5:0] = @(qw[111:104] >= 0x1f) ? 0x1f : qw[111:104]@
  shf_value14[5:0] = @(qw[119:112] >= 0x1f) ? 0x1f : qw[119:112]@
  shf_value15[5:0] = @(qw[127:120] >= 0x1f) ? 0x1f : qw[127:120]@

  DATA_L[ 31:  0] = QACC_L[ 31:  0] >> shf_value0[4:0]
  DATA_L[ 63: 32] = QACC_L[ 63: 32] >> shf_value1[4:0]
  DATA_L[ 95: 64] = QACC_L[ 95: 64] >> shf_value2[4:0]
  DATA_L[127: 96] = QACC_L[127: 96] >> shf_value3[4:0]
  DATA_L[159:128] = QACC_L[159:128] >> shf_value4[4:0]
  DATA_L[191:160] = QACC_L[191:160] >> shf_value5[4:0]
  DATA_L[223:192] = QACC_L[223:192] >> shf_value6[4:0]
  DATA_L[255:224] = QACC_L[255:224] >> shf_value7[4:0]
  DATA_H[ 31:  0] = QACC_H[ 31:  0] >> shf_value8[4:0]
  DATA_H[ 63: 32] = QACC_H[ 63: 32] >> shf_value9[4:0]
  DATA_H[ 95: 64] = QACC_H[ 95: 64] >> shf_value10[4:0]
  DATA_H[127: 96] = QACC_H[127: 96] >> shf_value11[4:0]
  DATA_H[159:128] = QACC_H[159:128] >> shf_value12[4:0]
  DATA_H[191:160] = QACC_H[191:160] >> shf_value13[4:0]
  DATA_H[223:192] = QACC_H[223:192] >> shf_value14[4:0]
  DATA_H[255:224] = QACC_H[255:224] >> shf_value15[4:0]


  QACC_L[255:0] = DATA_L[255:0]
  QACC_H[255:0] = DATA_H[255:0]

  if sel2 == 0:
    qu[  7:  0] = DATA_L[  7:  0]
    qu[ 15:  8] = DATA_L[ 39: 32]
    qu[ 23: 16] = DATA_L[ 71: 64]
    qu[ 31: 24] = DATA_L[103: 96]
    qu[ 39: 32] = DATA_L[135:128]
    qu[ 47: 40] = DATA_L[167:160]
    qu[ 55: 48] = DATA_L[199:192]
    qu[ 63: 56] = DATA_L[231:224]
    qu[ 71: 64] = DATA_H[  7:  0]
    qu[ 79: 72] = DATA_H[ 39: 32]
    qu[ 87: 80] = DATA_H[ 71: 64]
    qu[ 95: 88] = DATA_H[103: 96]
    qu[103: 96] = DATA_H[135:128]
    qu[111:104] = DATA_H[167:160]
    qu[119:112] = DATA_H[199:192]
    qu[127:120] = DATA_H[231:224]
  else if sel2 == 1:
    qu[  7:  0] = min(max(DATA_L[ 31:  0], -2^{7}), 2^{7}-1)
    qu[ 15:  8] = min(max(DATA_L[ 63: 32], -2^{7}), 2^{7}-1)
    qu[ 23: 16] = min(max(DATA_L[ 95: 64], -2^{7}), 2^{7}-1)
    qu[ 31: 24] = min(max(DATA_L[127: 96], -2^{7}), 2^{7}-1)
    qu[ 39: 32] = min(max(DATA_L[159:128], -2^{7}), 2^{7}-1)
    qu[ 47: 40] = min(max(DATA_L[191:160], -2^{7}), 2^{7}-1)
    qu[ 55: 48] = min(max(DATA_L[223:192], -2^{7}), 2^{7}-1)
    qu[ 63: 56] = min(max(DATA_L[255:224], -2^{7}), 2^{7}-1)
    qu[ 71: 64] = min(max(DATA_L[ 31:  0], -2^{7}), 2^{7}-1)
    qu[ 79: 72] = min(max(DATA_H[ 63: 32], -2^{7}), 2^{7}-1)
    qu[ 87: 80] = min(max(DATA_H[ 95: 64], -2^{7}), 2^{7}-1)
    qu[ 95: 88] = min(max(DATA_H[127: 96], -2^{7}), 2^{7}-1)
    qu[103: 96] = min(max(DATA_H[159:128], -2^{7}), 2^{7}-1)
    qu[111:104] = min(max(DATA_H[191:160], -2^{7}), 2^{7}-1)
    qu[119:112] = min(max(DATA_H[223:192], -2^{7}), 2^{7}-1)
    qu[127:120] = min(max(DATA_H[255:224], -2^{7}), 2^{7}-1)

\end{lstlisting}

\textbf{Schedule}\\
\begin{longtable}[c]{|l|l|l|}
    \hline
    \rowcolor{lightgray}
    \textbf{Operands} & \textbf{Use} & \textbf{Def} \\\hline
    \endfirsthead
    \hline
    \rowcolor{lightgray}
     \textbf{Operands} & \textbf{Use} & \textbf{Def} \\\hline
    \endhead
    \hline
    \endfoot

  qu & - & 2 \\\hline
  rs1 & 1 & - \\\hline
  QACC\_H & 1 & 2 \\\hline
  QACC\_L & 1 & 2

\end{longtable}

\textbf{Format}\\

\begin{figure*}[!htbp]
 {\setlength{\tabcolsep}{1pt}
 \begin{center}
\begin{tabular}{@{}V@{}I@{}W@{}I@{}F@{}I@{}S@{}F@{}F@{}O}
\instbitrange{31}{27}&
\instbit{26}&
\instbitrange{25}{24}&
\instbit{23}&
\instbitrange{22}{20}&
\instbit{19}&
\instbitrange{18}{13}&
\instbitrange{12}{10}&
\instbitrange{9}{7}&
\instbitrange{6}{0} \\
\hline
\multicolumn{1}{|c|}{1XX01} &
\multicolumn{1}{c|}{{\scriptsize sat[0]}} &
\multicolumn{1}{c|}{{\scriptsize qw[1:0]}} &
\multicolumn{1}{c|}{0} &
\multicolumn{1}{c|}{{\scriptsize rm[2:0]}} &
\multicolumn{1}{c|}{{\scriptsize qw[2]}} &
\multicolumn{1}{c|}{1XXX01} &
\multicolumn{1}{c|}{{\scriptsize qu[2:0]}} &
\multicolumn{1}{c|}{1XX} &
\multicolumn{1}{c|}{0011011}\\
\hline
5& 1& 2& 1& 3& 1& 6& 3& 3& 7 \\
\end{tabular}
 \end{center}
 }
 \vspace{-0.1in}
 \caption[]{ESP.SRCMB.S8.Q.QACC} \end{figure*}


\newpage
\subsubsection{ESP.SRCMB.S8.QACC}\label{instruction:ESPSRCMBS8QACC}
\textbf{Syntax}\\
ESP.SRCMB.S8.QACC qu,rs1,sat,rm

\textbf{Description}\\
This instruction extracts 16 data segments of 32 bits from special registers \lstinline{QACC_H} and \lstinline{QACC_L} and performs arithmetic right shift operations respectively. While writing the shift result back to \lstinline{QACC_H} and \lstinline{QACC_L}, the instruction saturates the result to 8-bit signed numbers and writes the 16 8-bit data obtained after saturation into the register \lstinline{qu}.\\
\textbf{Operation}
\begin{lstlisting}

  DATA_L[ 31:  0] = QACC_L[ 31:  0] >> rs1[4:0]
  DATA_L[ 63: 32] = QACC_L[ 63: 32] >> rs1[4:0]
  DATA_L[ 95: 64] = QACC_L[ 95: 64] >> rs1[4:0]
  DATA_L[127: 96] = QACC_L[127: 96] >> rs1[4:0]
  DATA_L[159:128] = QACC_L[159:128] >> rs1[4:0]
  DATA_L[191:160] = QACC_L[191:160] >> rs1[4:0]
  DATA_L[223:192] = QACC_L[223:192] >> rs1[4:0]
  DATA_L[255:224] = QACC_L[255:224] >> rs1[4:0]
  DATA_H[ 31:  0] = QACC_H[ 31:  0] >> rs1[4:0]
  DATA_H[ 63: 32] = QACC_H[ 63: 32] >> rs1[4:0]
  DATA_H[ 95: 64] = QACC_H[ 95: 64] >> rs1[4:0]
  DATA_H[127: 96] = QACC_H[127: 96] >> rs1[4:0]
  DATA_H[159:128] = QACC_H[159:128] >> rs1[4:0]
  DATA_H[191:160] = QACC_H[191:160] >> rs1[4:0]
  DATA_H[223:192] = QACC_H[223:192] >> rs1[4:0]
  DATA_H[255:224] = QACC_H[255:224] >> rs1[4:0]

  QACC_L[255:0] = DATA_L[255:0]
  QACC_H[255:0] = DATA_H[255:0]

  if sel2 == 0:
    qu[  7:  0] = DATA_L[  7:  0]
    qu[ 15:  8] = DATA_L[ 39: 32]
    qu[ 23: 16] = DATA_L[ 71: 64]
    qu[ 31: 24] = DATA_L[103: 96]
    qu[ 39: 32] = DATA_L[135:128]
    qu[ 47: 40] = DATA_L[167:160]
    qu[ 55: 48] = DATA_L[199:192]
    qu[ 63: 56] = DATA_L[231:224]
    qu[ 71: 64] = DATA_H[  7:  0]
    qu[ 79: 72] = DATA_H[ 39: 32]
    qu[ 87: 80] = DATA_H[ 71: 64]
    qu[ 95: 88] = DATA_H[103: 96]
    qu[103: 96] = DATA_H[135:128]
    qu[111:104] = DATA_H[167:160]
    qu[119:112] = DATA_H[199:192]
    qu[127:120] = DATA_H[231:224]
  else if sel2 == 1:
    qu[  7:  0] = min(max(DATA_L[ 31:  0], -2^{7}), 2^{7}-1)
    qu[ 15:  8] = min(max(DATA_L[ 63: 32], -2^{7}), 2^{7}-1)
    qu[ 23: 16] = min(max(DATA_L[ 95: 64], -2^{7}), 2^{7}-1)
    qu[ 31: 24] = min(max(DATA_L[127: 96], -2^{7}), 2^{7}-1)
    qu[ 39: 32] = min(max(DATA_L[159:128], -2^{7}), 2^{7}-1)
    qu[ 47: 40] = min(max(DATA_L[191:160], -2^{7}), 2^{7}-1)
    qu[ 55: 48] = min(max(DATA_L[223:192], -2^{7}), 2^{7}-1)
    qu[ 63: 56] = min(max(DATA_L[255:224], -2^{7}), 2^{7}-1)
    qu[ 71: 64] = min(max(DATA_L[ 31:  0], -2^{7}), 2^{7}-1)
    qu[ 79: 72] = min(max(DATA_H[ 63: 32], -2^{7}), 2^{7}-1)
    qu[ 87: 80] = min(max(DATA_H[ 95: 64], -2^{7}), 2^{7}-1)
    qu[ 95: 88] = min(max(DATA_H[127: 96], -2^{7}), 2^{7}-1)
    qu[103: 96] = min(max(DATA_H[159:128], -2^{7}), 2^{7}-1)
    qu[111:104] = min(max(DATA_H[191:160], -2^{7}), 2^{7}-1)
    qu[119:112] = min(max(DATA_H[223:192], -2^{7}), 2^{7}-1)
    qu[127:120] = min(max(DATA_H[255:224], -2^{7}), 2^{7}-1)

\end{lstlisting}

\textbf{Schedule}\\
\begin{longtable}[c]{|l|l|l|}
    \hline
    \rowcolor{lightgray}
    \textbf{Operands} & \textbf{Use} & \textbf{Def} \\\hline
    \endfirsthead
    \hline
    \rowcolor{lightgray}
     \textbf{Operands} & \textbf{Use} & \textbf{Def} \\\hline
    \endhead
    \hline
    \endfoot

  qu & - & 2 \\\hline
  rs1 & 1 & - \\\hline
  QACC\_H & 1 & 2 \\\hline
  QACC\_L & 1 & 2

\end{longtable}

\textbf{Format}\\

\begin{figure*}[!htbp]
 {\setlength{\tabcolsep}{1pt}
 \begin{center}
\begin{tabular}{@{}W@{}I@{}F@{}O@{}I@{}F@{}W@{}F@{}F@{}O}
\instbitrange{31}{30}&
\instbit{29}&
\instbitrange{28}{26}&
\instbitrange{25}{19}&
\instbit{18}&
\instbitrange{17}{15}&
\instbitrange{14}{13}&
\instbitrange{12}{10}&
\instbitrange{9}{7}&
\instbitrange{6}{0} \\
\hline
\multicolumn{1}{|c|}{01} &
\multicolumn{1}{c|}{{\scriptsize sat[0]}} &
\multicolumn{1}{c|}{{\scriptsize rm[2:0]}} &
\multicolumn{1}{c|}{111010X} &
\multicolumn{1}{c|}{{\scriptsize rs1[4]}} &
\multicolumn{1}{c|}{{\scriptsize rs1[2:0]}} &
\multicolumn{1}{c|}{00} &
\multicolumn{1}{c|}{{\scriptsize qu[2:0]}} &
\multicolumn{1}{c|}{XX1} &
\multicolumn{1}{c|}{0011111}\\
\hline
2& 1& 3& 7& 1& 3& 2& 3& 3& 7 \\
\end{tabular}
 \end{center}
 }
 \vspace{-0.1in}
 \caption[]{ESP.SRCMB.S8.QACC} \end{figure*}


\newpage
\subsubsection{ESP.SRCMB.U16.Q.QACC}\label{instruction:ESPSRCMBU16QQACC}
\textbf{Syntax}\\
ESP.SRCMB.U16.Q.QACC qu,qw,sat,rm

\textbf{Description}\\
This instruction extracts 8 data segments of 64 bits from special registers \lstinline{QACC_H} and \lstinline{QACC_L} and performs arithmetic right shift operations respectively, the shift amounts are from \lstinline{qw}. While writing the shift results back to \lstinline{QACC_H} and \lstinline{QACC_L}, the instruction saturates the results to 16-bit unsigned numbers and writes the 8 16-bit data obtained after saturation into the register \lstinline{qu}.\\
\textbf{Operation}
\begin{lstlisting}[language=C, escapechar=@]
  shf_value0[5:0] = @(qw[ 15:  0] >= 0x3f) ? 0x3f : qw[ 15:  0]@
  shf_value1[5:0] = @(qw[ 31: 16] >= 0x3f) ? 0x3f : qw[ 31: 16]@
  shf_value2[5:0] = @(qw[ 47: 32] >= 0x3f) ? 0x3f : qw[ 47: 32]@
  shf_value3[5:0] = @(qw[ 63: 48] >= 0x3f) ? 0x3f : qw[ 63: 48]@
  shf_value4[5:0] = @(qw[ 79: 64] >= 0x3f) ? 0x3f : qw[ 79: 64]@
  shf_value5[5:0] = @(qw[ 95: 80] >= 0x3f) ? 0x3f : qw[ 95: 80]@
  shf_value6[5:0] = @(qw[111: 96] >= 0x3f) ? 0x3f : qw[111: 96]@
  shf_value7[5:0] = @(qw[127:112] >= 0x3f) ? 0x3f : qw[127:112]@

  DATA_L[ 63:  0] = QACC_L[ 63:  0] >> shf_value0[5:0]
  DATA_L[127: 64] = QACC_L[127: 64] >> shf_value1[5:0]
  DATA_L[191:128] = QACC_L[191:128] >> shf_value2[5:0]
  DATA_L[255:192] = QACC_L[255:192] >> shf_value3[5:0]
  DATA_H[ 63:  0] = QACC_H[ 63:  0] >> shf_value4[5:0]
  DATA_H[127: 64] = QACC_H[127: 64] >> shf_value5[5:0]
  DATA_H[191:128] = QACC_H[191:128] >> shf_value6[5:0]
  DATA_H[255:192] = QACC_H[255:192] >> shf_value7[5:0]

  QACC_L[255:0] = DATA_L[255:0]
  QACC_H[255:0] = DATA_H[255:0]

  if sel2 == 0:
    qu[15: 0] = DATA_L[15:0]
    qu[31:16] = DATA_L[79:64]
    qu[47:32] = DATA_L[143:128]
    qu[63:48] = DATA_L[207:192]
    qu[79:64] = DATA_H[15:0]
    qu[95:80] = DATA_H[79:64]
    qu[111:96] = DATA_H[143:128]
    qu[127:112] = DATA_H[207:192]
  else if sel2 == 1:
    qu[ 15:  0] = min(DATA_L[ 63:  0], 2^{16}-1)
    qu[ 31: 16] = min(DATA_L[127: 64], 2^{16}-1)
    qu[ 47: 32] = min(DATA_L[191:128], 2^{16}-1)
    qu[ 63: 48] = min(DATA_L[255:192], 2^{16}-1)
    qu[ 79: 64] = min(DATA_H[ 63:  0], 2^{16}-1)
    qu[ 95: 80] = min(DATA_H[127: 64], 2^{16}-1)
    qu[111: 96] = min(DATA_H[191:128], 2^{16}-1)
    qu[127:112] = min(DATA_H[255:192], 2^{16}-1)
\end{lstlisting}

\textbf{Schedule}\\
\begin{longtable}[c]{|l|l|l|}
    \hline
    \rowcolor{lightgray}
    \textbf{Operands} & \textbf{Use} & \textbf{Def} \\\hline
    \endfirsthead
    \hline
    \rowcolor{lightgray}
     \textbf{Operands} & \textbf{Use} & \textbf{Def} \\\hline
    \endhead
    \hline
    \endfoot

  qu & - & 2 \\\hline
  rs1 & 1 & - \\\hline
  QACC\_H & 1 & 2 \\\hline
  QACC\_L & 1 & 2

\end{longtable}

\textbf{Format}\\

\begin{figure*}[!htbp]
 {\setlength{\tabcolsep}{1pt}
 \begin{center}
\begin{tabular}{@{}V@{}I@{}W@{}I@{}F@{}I@{}S@{}F@{}F@{}O}
\instbitrange{31}{27}&
\instbit{26}&
\instbitrange{25}{24}&
\instbit{23}&
\instbitrange{22}{20}&
\instbit{19}&
\instbitrange{18}{13}&
\instbitrange{12}{10}&
\instbitrange{9}{7}&
\instbitrange{6}{0} \\
\hline
\multicolumn{1}{|c|}{1XX10} &
\multicolumn{1}{c|}{{\scriptsize sat[0]}} &
\multicolumn{1}{c|}{{\scriptsize qw[1:0]}} &
\multicolumn{1}{c|}{0} &
\multicolumn{1}{c|}{{\scriptsize rm[2:0]}} &
\multicolumn{1}{c|}{{\scriptsize qw[2]}} &
\multicolumn{1}{c|}{1XXX01} &
\multicolumn{1}{c|}{{\scriptsize qu[2:0]}} &
\multicolumn{1}{c|}{1XX} &
\multicolumn{1}{c|}{0011011}\\
\hline
5& 1& 2& 1& 3& 1& 6& 3& 3& 7 \\
\end{tabular}
 \end{center}
 }
 \vspace{-0.1in}
 \caption[]{ESP.SRCMB.U16.Q.QACC} \end{figure*}


\newpage
\subsubsection{ESP.SRCMB.U16.QACC}\label{instruction:ESPSRCMBU16QACC}
\textbf{Syntax}\\
ESP.SRCMB.U16.QACC qu,rs1,sat,rm

\textbf{Description}\\
This instruction extracts 8 data segments of 64 bits from special registers \lstinline{QACC_H} and \lstinline{QACC_L} and performs arithmetic right shift operations respectively. While writing the shift results back to \lstinline{QACC_H} and \lstinline{QACC_L}, the instruction saturates the results to 16-bit unsigned numbers and writes the 8 16-bit data segments obtained after saturation into the \lstinline{qu} register.\\
\textbf{Operation}
\begin{lstlisting}
  temp0[63:0] = QACC_L[ 63:  0]
  ...
  temp3[63:0] = QACC_L[255:192]
  temp4[63:0] = QACC_H[ 63:  0]
  ...
  temp7[63:0] = QACC_H[255:192]

  temp_shf0[63:0] = temp0[63:0] >> rs1[5:0]
  temp_shf1[63:0] = temp1[63:0] >> rs1[5:0]
  ...
  temp_shf7[63:0] = temp7[63:0] >> rs1[5:0]

  QACC_L[ 63:  0] = temp_shf0[63:0]
  ...
  QACC_L[255:192] = temp_shf3[63:0]
  QACC_H[ 63:  0] = temp_shf4[63:0]
  ...
  QACC_H[255:192] = temp_shf7[63:0]

  if sel2 == 0:
    qu[ 15:  0] = temp0[15:0]
    qu[ 31: 16] = temp1[15:0]
    qu[ 47: 32] = temp2[15:0]
    qu[ 63: 48] = temp3[15:0]
    qu[ 79: 64] = temp4[15:0]
    qu[ 95: 80] = temp5[15:0]
    qu[111: 96] = temp6[15:0]
    qu[127:112] = temp7[15:0]
  else if sel2==1:
    qu[ 15:  0] = min(temp_shf0[63:0], 2^{16}-1)
    ...
    qu[ 63: 48] = min(temp_shf3[63:0], 2^{16}-1)
    qu[ 79: 64] = min(temp_shf4[63:0], 2^{16}-1)
    ...
    qu[127:112] = min(temp_shf7[63:0], 2^{16}-1)
\end{lstlisting}

\textbf{Schedule}\\
\begin{longtable}[c]{|l|l|l|}
    \hline
    \rowcolor{lightgray}
    \textbf{Operands} & \textbf{Use} & \textbf{Def} \\\hline
    \endfirsthead
    \hline
    \rowcolor{lightgray}
     \textbf{Operands} & \textbf{Use} & \textbf{Def} \\\hline
    \endhead
    \hline
    \endfoot

  qu & - & 2 \\\hline
  rs1 & 1 & - \\\hline
  QACC\_H & 1 & 2 \\\hline
  QACC\_L & 1 & 2

\end{longtable}

\textbf{Format}\\

\begin{figure*}[!htbp]
 {\setlength{\tabcolsep}{1pt}
 \begin{center}
\begin{tabular}{@{}W@{}I@{}F@{}O@{}I@{}F@{}W@{}F@{}F@{}O}
\instbitrange{31}{30}&
\instbit{29}&
\instbitrange{28}{26}&
\instbitrange{25}{19}&
\instbit{18}&
\instbitrange{17}{15}&
\instbitrange{14}{13}&
\instbitrange{12}{10}&
\instbitrange{9}{7}&
\instbitrange{6}{0} \\
\hline
\multicolumn{1}{|c|}{10} &
\multicolumn{1}{c|}{{\scriptsize sat[0]}} &
\multicolumn{1}{c|}{{\scriptsize rm[2:0]}} &
\multicolumn{1}{c|}{111010X} &
\multicolumn{1}{c|}{{\scriptsize rs1[4]}} &
\multicolumn{1}{c|}{{\scriptsize rs1[2:0]}} &
\multicolumn{1}{c|}{00} &
\multicolumn{1}{c|}{{\scriptsize qu[2:0]}} &
\multicolumn{1}{c|}{XX1} &
\multicolumn{1}{c|}{0011111}\\
\hline
2& 1& 3& 7& 1& 3& 2& 3& 3& 7 \\
\end{tabular}
 \end{center}
 }
 \vspace{-0.1in}
 \caption[]{ESP.SRCMB.U16.QACC} \end{figure*}


\newpage
\subsubsection{ESP.SRCMB.U8.Q.QACC}\label{instruction:ESPSRCMBU8QQACC}
\textbf{Syntax}\\
ESP.SRCMB.U8.Q.QACC qu,qw,sat,rm

\textbf{Description}\\
This instruction extracts 16 data segments of 32 bits from special registers \lstinline{QACC_H} and \lstinline{QACC_L} and performs arithmetic right shift operations respectively, the shift amounts are from \lstinline{qw}. While writing the shift results back to \lstinline{QACC_H} and \lstinline{QACC_L}, the instruction saturates the results to 8-bit unsigned numbers and writes the 16 8-bit data obtained after saturation into the register \lstinline{qu}.\\

\textbf{Operation}

\begin{lstlisting}[language=C, escapechar=@]
  shf_value0[5:0] = @(qw[  7:  0] >= 0x1f) ? 0x1f : qw[  7:  0]@
  shf_value1[5:0] = @(qw[ 15:  8] >= 0x1f) ? 0x1f : qw[ 15:  8]@
  shf_value2[5:0] = @(qw[ 23: 16] >= 0x1f) ? 0x1f : qw[ 23: 16]@
  shf_value3[5:0] = @(qw[ 31: 24] >= 0x1f) ? 0x1f : qw[ 31: 24]@
  shf_value4[5:0] = @(qw[ 39: 32] >= 0x1f) ? 0x1f : qw[ 39: 32]@
  shf_value5[5:0] = @(qw[ 47: 40] >= 0x1f) ? 0x1f : qw[ 47: 40]@
  shf_value6[5:0] = @(qw[ 55: 48] >= 0x1f) ? 0x1f : qw[ 55: 48]@
  shf_value7[5:0] = @(qw[ 63: 56] >= 0x1f) ? 0x1f : qw[ 63: 56]@
  shf_value8[5:0] = @(qw[ 71: 64] >= 0x1f) ? 0x1f : qw[ 71: 64]@
  shf_value9[5:0] = @(qw[ 79: 72] >= 0x1f) ? 0x1f : qw[ 79: 72]@
  shf_value10[5:0] = @(qw[ 87: 80] >= 0x1f) ? 0x1f : qw[ 87: 80]@
  shf_value11[5:0] = @(qw[ 95: 88] >= 0x1f) ? 0x1f : qw[ 95: 88]@
  shf_value12[5:0] = @(qw[103: 96] >= 0x1f) ? 0x1f : qw[103: 96]@
  shf_value13[5:0] = @(qw[111:104] >= 0x1f) ? 0x1f : qw[111:104]@
  shf_value14[5:0] = @(qw[119:112] >= 0x1f) ? 0x1f : qw[119:112]@
  shf_value15[5:0] = @(qw[127:120] >= 0x1f) ? 0x1f : qw[127:120]@

  DATA_L[ 31:  0] = QACC_L[ 31:  0] >> shf_value0[4:0]
  DATA_L[ 63: 32] = QACC_L[ 63: 32] >> shf_value1[4:0]
  DATA_L[ 95: 64] = QACC_L[ 95: 64] >> shf_value2[4:0]
  DATA_L[127: 96] = QACC_L[127: 96] >> shf_value3[4:0]
  DATA_L[159:128] = QACC_L[159:128] >> shf_value4[4:0]
  DATA_L[191:160] = QACC_L[191:160] >> shf_value5[4:0]
  DATA_L[223:192] = QACC_L[223:192] >> shf_value6[4:0]
  DATA_L[255:224] = QACC_L[255:224] >> shf_value7[4:0]
  DATA_H[ 31:  0] = QACC_H[ 31:  0] >> shf_value8[4:0]
  DATA_H[ 63: 32] = QACC_H[ 63: 32] >> shf_value9[4:0]
  DATA_H[ 95: 64] = QACC_H[ 95: 64] >> shf_value10[4:0]
  DATA_H[127: 96] = QACC_H[127: 96] >> shf_value11[4:0]
  DATA_H[159:128] = QACC_H[159:128] >> shf_value12[4:0]
  DATA_H[191:160] = QACC_H[191:160] >> shf_value13[4:0]
  DATA_H[223:192] = QACC_H[223:192] >> shf_value14[4:0]
  DATA_H[255:224] = QACC_H[255:224] >> shf_value15[4:0]

  QACC_L[255:0] = DATA_L[255:0]
  QACC_H[255:0] = DATA_H[255:0]

  if sel2 == 0:
    qu[  7:  0] = DATA_L[  7:  0]
    qu[ 15:  8] = DATA_L[ 39: 32]
    qu[ 23: 16] = DATA_L[ 71: 64]
    qu[ 31: 24] = DATA_L[103: 96]
    qu[ 39: 32] = DATA_L[135:128]
    qu[ 47: 40] = DATA_L[167:160]
    qu[ 55: 48] = DATA_L[199:192]
    qu[ 63: 56] = DATA_L[231:224]
    qu[ 71: 64] = DATA_H[  7:  0]
    qu[ 79: 72] = DATA_H[ 39: 32]
    qu[ 87: 80] = DATA_H[ 71: 64]
    qu[ 95: 88] = DATA_H[103: 96]
    qu[103: 96] = DATA_H[135:128]
    qu[111:104] = DATA_H[167:160]
    qu[119:112] = DATA_H[199:192]
    qu[127:120] = DATA_H[231:224]
  else if sel2 == 1:
    qu[  7:  0] = min(DATA_L[ 31:  0], 2^{8}-1)
    qu[ 15:  8] = min(DATA_L[ 63: 32], 2^{8}-1)
    qu[ 23: 16] = min(DATA_L[ 95: 64], 2^{8}-1)
    qu[ 31: 24] = min(DATA_L[127: 96], 2^{8}-1)
    qu[ 39: 32] = min(DATA_L[159:128], 2^{8}-1)
    qu[ 47: 40] = min(DATA_L[191:160], 2^{8}-1)
    qu[ 55: 48] = min(DATA_L[223:192], 2^{8}-1)
    qu[ 63: 56] = min(DATA_L[255:224], 2^{8}-1)
    qu[ 71: 64] = min(DATA_L[ 31:  0], 2^{8}-1)
    qu[ 79: 72] = min(DATA_H[ 63: 32], 2^{8}-1)
    qu[ 87: 80] = min(DATA_H[ 95: 64], 2^{8}-1)
    qu[ 95: 88] = min(DATA_H[127: 96], 2^{8}-1)
    qu[103: 96] = min(DATA_H[159:128], 2^{8}-1)
    qu[111:104] = min(DATA_H[191:160], 2^{8}-1)
    qu[119:112] = min(DATA_H[223:192], 2^{8}-1)
    qu[127:120] = min(DATA_H[255:224], 2^{8}-1)

\end{lstlisting}

\textbf{Schedule}\\
\begin{longtable}[c]{|l|l|l|}
    \hline
    \rowcolor{lightgray}
    \textbf{Operands} & \textbf{Use} & \textbf{Def} \\\hline
    \endfirsthead
    \hline
    \rowcolor{lightgray}
     \textbf{Operands} & \textbf{Use} & \textbf{Def} \\\hline
    \endhead
    \hline
    \endfoot

  qu & - & 2 \\\hline
  rs1 & 1 & - \\\hline
  QACC\_H & 1 & 2 \\\hline
  QACC\_L & 1 & 2

\end{longtable}

\textbf{Format}\\

\begin{figure*}[!htbp]
 {\setlength{\tabcolsep}{1pt}
 \begin{center}
\begin{tabular}{@{}V@{}I@{}W@{}I@{}F@{}I@{}S@{}F@{}F@{}O}
\instbitrange{31}{27}&
\instbit{26}&
\instbitrange{25}{24}&
\instbit{23}&
\instbitrange{22}{20}&
\instbit{19}&
\instbitrange{18}{13}&
\instbitrange{12}{10}&
\instbitrange{9}{7}&
\instbitrange{6}{0} \\
\hline
\multicolumn{1}{|c|}{1XX00} &
\multicolumn{1}{c|}{{\scriptsize sat[0]}} &
\multicolumn{1}{c|}{{\scriptsize qw[1:0]}} &
\multicolumn{1}{c|}{0} &
\multicolumn{1}{c|}{{\scriptsize rm[2:0]}} &
\multicolumn{1}{c|}{{\scriptsize qw[2]}} &
\multicolumn{1}{c|}{1XXX01} &
\multicolumn{1}{c|}{{\scriptsize qu[2:0]}} &
\multicolumn{1}{c|}{1XX} &
\multicolumn{1}{c|}{0011011}\\
\hline
5& 1& 2& 1& 3& 1& 6& 3& 3& 7 \\
\end{tabular}
 \end{center}
 }
 \vspace{-0.1in}
 \caption[]{ESP.SRCMB.U8.Q.QACC} \end{figure*}


\newpage
\subsubsection{ESP.SRCMB.U8.QACC}\label{instruction:ESPSRCMBU8QACC}
\textbf{Syntax}\\
ESP.SRCMB.U8.QACC qu,rs1,sat,rm

\textbf{Description}\\
This instruction extracts 16 data segments of 32 bits from special registers \lstinline{QACC_H} and \lstinline{QACC_L} and performs arithmetic right shift operations respectively. While writing the shift result back to \lstinline{QACC_H} and \lstinline{QACC_L}, the instruction saturates the result to 8-bit unsigned numbers and writes the 16 8-bit data obtained after saturation into the register \lstinline{qu}.\\
\textbf{Operation}
\begin{lstlisting}

  DATA_L[ 31:  0] = QACC_L[ 31:  0] >> rs1[4:0]
  DATA_L[ 63: 32] = QACC_L[ 63: 32] >> rs1[4:0]
  DATA_L[ 95: 64] = QACC_L[ 95: 64] >> rs1[4:0]
  DATA_L[127: 96] = QACC_L[127: 96] >> rs1[4:0]
  DATA_L[159:128] = QACC_L[159:128] >> rs1[4:0]
  DATA_L[191:160] = QACC_L[191:160] >> rs1[4:0]
  DATA_L[223:192] = QACC_L[223:192] >> rs1[4:0]
  DATA_L[255:224] = QACC_L[255:224] >> rs1[4:0]
  DATA_H[ 31:  0] = QACC_H[ 31:  0] >> rs1[4:0]
  DATA_H[ 63: 32] = QACC_H[ 63: 32] >> rs1[4:0]
  DATA_H[ 95: 64] = QACC_H[ 95: 64] >> rs1[4:0]
  DATA_H[127: 96] = QACC_H[127: 96] >> rs1[4:0]
  DATA_H[159:128] = QACC_H[159:128] >> rs1[4:0]
  DATA_H[191:160] = QACC_H[191:160] >> rs1[4:0]
  DATA_H[223:192] = QACC_H[223:192] >> rs1[4:0]
  DATA_H[255:224] = QACC_H[255:224] >> rs1[4:0]

  QACC_L[255:0] = DATA_L[255:0]
  QACC_H[255:0] = DATA_H[255:0]

  if sel2 == 0:
    qu[  7:  0] = DATA_L[  7:  0]
    qu[ 15:  8] = DATA_L[ 39: 32]
    qu[ 23: 16] = DATA_L[ 71: 64]
    qu[ 31: 24] = DATA_L[103: 96]
    qu[ 39: 32] = DATA_L[135:128]
    qu[ 47: 40] = DATA_L[167:160]
    qu[ 55: 48] = DATA_L[199:192]
    qu[ 63: 56] = DATA_L[231:224]
    qu[ 71: 64] = DATA_H[  7:  0]
    qu[ 79: 72] = DATA_H[ 39: 32]
    qu[ 87: 80] = DATA_H[ 71: 64]
    qu[ 95: 88] = DATA_H[103: 96]
    qu[103: 96] = DATA_H[135:128]
    qu[111:104] = DATA_H[167:160]
    qu[119:112] = DATA_H[199:192]
    qu[127:120] = DATA_H[231:224]
  else if sel2 == 1:
    qu[  7:  0] = min(DATA_L[ 31:  0], 2^{8}-1)
    qu[ 15:  8] = min(DATA_L[ 63: 32], 2^{8}-1)
    qu[ 23: 16] = min(DATA_L[ 95: 64], 2^{8}-1)
    qu[ 31: 24] = min(DATA_L[127: 96], 2^{8}-1)
    qu[ 39: 32] = min(DATA_L[159:128], 2^{8}-1)
    qu[ 47: 40] = min(DATA_L[191:160], 2^{8}-1)
    qu[ 55: 48] = min(DATA_L[223:192], 2^{8}-1)
    qu[ 63: 56] = min(DATA_L[255:224], 2^{8}-1)
    qu[ 71: 64] = min(DATA_L[ 31:  0], 2^{8}-1)
    qu[ 79: 72] = min(DATA_H[ 63: 32], 2^{8}-1)
    qu[ 87: 80] = min(DATA_H[ 95: 64], 2^{8}-1)
    qu[ 95: 88] = min(DATA_H[127: 96], 2^{8}-1)
    qu[103: 96] = min(DATA_H[159:128], 2^{8}-1)
    qu[111:104] = min(DATA_H[191:160], 2^{8}-1)
    qu[119:112] = min(DATA_H[223:192], 2^{8}-1)
    qu[127:120] = min(DATA_H[255:224], 2^{8}-1)

\end{lstlisting}

\textbf{Schedule}\\
\begin{longtable}[c]{|l|l|l|}
    \hline
    \rowcolor{lightgray}
    \textbf{Operands} & \textbf{Use} & \textbf{Def} \\\hline
    \endfirsthead
    \hline
    \rowcolor{lightgray}
     \textbf{Operands} & \textbf{Use} & \textbf{Def} \\\hline
    \endhead
    \hline
    \endfoot

  qu & - & 2 \\\hline
  rs1 & 1 & - \\\hline
  QACC\_H & 1 & 2 \\\hline
  QACC\_L & 1 & 2

\end{longtable}

\textbf{Format}\\

\begin{figure*}[!htbp]
 {\setlength{\tabcolsep}{1pt}
 \begin{center}
\begin{tabular}{@{}W@{}I@{}F@{}O@{}I@{}F@{}W@{}F@{}F@{}O}
\instbitrange{31}{30}&
\instbit{29}&
\instbitrange{28}{26}&
\instbitrange{25}{19}&
\instbit{18}&
\instbitrange{17}{15}&
\instbitrange{14}{13}&
\instbitrange{12}{10}&
\instbitrange{9}{7}&
\instbitrange{6}{0} \\
\hline
\multicolumn{1}{|c|}{00} &
\multicolumn{1}{c|}{{\scriptsize sat[0]}} &
\multicolumn{1}{c|}{{\scriptsize rm[2:0]}} &
\multicolumn{1}{c|}{111010X} &
\multicolumn{1}{c|}{{\scriptsize rs1[4]}} &
\multicolumn{1}{c|}{{\scriptsize rs1[2:0]}} &
\multicolumn{1}{c|}{00} &
\multicolumn{1}{c|}{{\scriptsize qu[2:0]}} &
\multicolumn{1}{c|}{XX1} &
\multicolumn{1}{c|}{0011111}\\
\hline
2& 1& 3& 7& 1& 3& 2& 3& 3& 7 \\
\end{tabular}
 \end{center}
 }
 \vspace{-0.1in}
 \caption[]{ESP.SRCMB.U8.QACC} \end{figure*}


\newpage
\subsubsection{ESP.SRCQ.128.ST.INCP}\label{instruction:ESPSRCQ128STINCP}
\textbf{Syntax}\\
ESP.SRCQ.128.ST.INCP qw,qy,rs1

\textbf{Description}\\
This instruction performs an arithmetic right shift on the 32-byte concatenation of registers \lstinline{qw} and \lstinline{qy}. Then, it writes the lower 128 bits of the shift result to memory. After the access, the value in the register \lstinline{rs1} is incremented by 16.\\
\textbf{Operation}
\begin{lstlisting}
  {qy[127:  0], qw[127:  0]} >> {SAR_BYTE[3:0] << 3}  => store128(rs1[31:0])
  rs1[31:0] = rs1[31:0] + 16
\end{lstlisting}

\textbf{Schedule}\\
\begin{longtable}[c]{|l|l|l|}
    \hline
    \rowcolor{lightgray}
    \textbf{Operands} & \textbf{Use} & \textbf{Def} \\\hline
    \endfirsthead
    \hline
    \rowcolor{lightgray}
     \textbf{Operands} & \textbf{Use} & \textbf{Def} \\\hline
    \endhead
    \hline
    \endfoot

  qz & - & 1 \\\hline
  qu & 1 & 1 \\\hline
  qv & 1 & 2 \\\hline
  rs1 & 1 & 1 \\\hline
  SAR\_BYTES & 1 & -

\end{longtable}

\textbf{Format}\\

\begin{figure*}[!htbp]
 {\setlength{\tabcolsep}{1pt}
 \begin{center}
\begin{tabular}{@{}F@{}F@{}W@{}Y@{}I@{}I@{}F@{}E@{}O}
\instbitrange{31}{29}&
\instbitrange{28}{26}&
\instbitrange{25}{24}&
\instbitrange{23}{20}&
\instbit{19}&
\instbit{18}&
\instbitrange{17}{15}&
\instbitrange{14}{7}&
\instbitrange{6}{0} \\
\hline
\multicolumn{1}{|c|}{1XX} &
\multicolumn{1}{c|}{{\scriptsize qy[2:0]}} &
\multicolumn{1}{c|}{{\scriptsize qw[1:0]}} &
\multicolumn{1}{c|}{X1XX} &
\multicolumn{1}{c|}{{\scriptsize qw[2]}} &
\multicolumn{1}{c|}{{\scriptsize rs1[4]}} &
\multicolumn{1}{c|}{{\scriptsize rs1[2:0]}} &
\multicolumn{1}{c|}{100XXXXX} &
\multicolumn{1}{c|}{0011011}\\
\hline
3& 3& 2& 4& 1& 1& 3& 8& 7 \\
\end{tabular}
 \end{center}
 }
 \vspace{-0.1in}
 \caption[]{ESP.SRCQ.128.ST.INCP} \end{figure*}


\newpage
\subsubsection{ESP.SRCXXP.2Q}\label{instruction:ESPSRCXXP2Q}
\textbf{Syntax}\\
ESP.SRCXXP.2Q qy,qw,rs1,rs2

\textbf{Description}\\
This instruction performs a logical right shift on the 32-byte concatenation of registers \lstinline{qw} and \lstinline{qy}, and pads the higher bits with 0. The upper 128 bits of the shift result are written to register \lstinline{qy}, and the lower 128 bits are written to \lstinline{qw}. The right shift amount is 8 times the sum of 1 plus the lower 4-bit value of register \lstinline{rs1}. After the above operations, the value in \lstinline{rs1} is incremented by the value in \lstinline{rs2}.\\
\textbf{Operation}
\begin{lstlisting}
  {qy[127:  0], qw[127:  0]} = {qy[127:  0], qw[127:  0]} >> ((rs1[3:0]+1)*8)
  qy[127:127-8*rs1[3:  0]] = 0
  rs1[31:0] = rs1[31:0] + rs2[31:0]
\end{lstlisting}

\textbf{Schedule}\\
\begin{longtable}[c]{|l|l|l|}
    \hline
    \rowcolor{lightgray}
    \textbf{Operands} & \textbf{Use} & \textbf{Def} \\\hline
    \endfirsthead
    \hline
    \rowcolor{lightgray}
     \textbf{Operands} & \textbf{Use} & \textbf{Def} \\\hline
    \endhead
    \hline
    \endfoot

  qy & 1 & 1 \\\hline
  qw & 1 & 1 \\\hline
  rs1 & 1 & 1 \\\hline
  rs2 & 1 & -

\end{longtable}

\textbf{Format}\\

\begin{figure*}[!htbp]
 {\setlength{\tabcolsep}{1pt}
 \begin{center}
\begin{tabular}{@{}F@{}F@{}W@{}I@{}F@{}I@{}I@{}F@{}E@{}O}
\instbitrange{31}{29}&
\instbitrange{28}{26}&
\instbitrange{25}{24}&
\instbit{23}&
\instbitrange{22}{20}&
\instbit{19}&
\instbit{18}&
\instbitrange{17}{15}&
\instbitrange{14}{7}&
\instbitrange{6}{0} \\
\hline
\multicolumn{1}{|c|}{X0X} &
\multicolumn{1}{c|}{{\scriptsize qy[2:0]}} &
\multicolumn{1}{c|}{{\scriptsize qw[1:0]}} &
\multicolumn{1}{c|}{{\scriptsize rs2[4]}} &
\multicolumn{1}{c|}{{\scriptsize rs2[2:0]}} &
\multicolumn{1}{c|}{{\scriptsize qw[2]}} &
\multicolumn{1}{c|}{{\scriptsize rs1[4]}} &
\multicolumn{1}{c|}{{\scriptsize rs1[2:0]}} &
\multicolumn{1}{c|}{11XX1XXX} &
\multicolumn{1}{c|}{0011011}\\
\hline
3& 3& 2& 1& 3& 1& 1& 3& 8& 7 \\
\end{tabular}
 \end{center}
 }
 \vspace{-0.1in}
 \caption[]{ESP.SRCXXP.2Q} \end{figure*}


\newpage
\subsubsection{ESP.SRS.S.XACC}\label{instruction:ESPSRSSXACC}
\textbf{Syntax}\\
ESP.SRS.S.XACC rd,rs1,sat,rm

\textbf{Description}\\
This instruction performs an arithmetic right shift on the special register \lstinline{XACC}. While writing the shift result back to \lstinline{XACC}, the instruction saturates the shift result to a 32-bit signed number and writes the saturated result into the register \lstinline{rd}.\\
\textbf{Operation}
\begin{lstlisting}
  temp_shf[39:0] = XACC[39:0] >> rs1[5:0]
  XACC = temp_shf[39:0]
  rd = min(max(temp_shf[39:0], -2^{31}), 2^{31}-1)
\end{lstlisting}

\textbf{Schedule}\\
\begin{longtable}[c]{|l|l|l|}
    \hline
    \rowcolor{lightgray}
    \textbf{Operands} & \textbf{Use} & \textbf{Def} \\\hline
    \endfirsthead
    \hline
    \rowcolor{lightgray}
     \textbf{Operands} & \textbf{Use} & \textbf{Def} \\\hline
    \endhead
    \hline
    \endfoot

  rd & - & 2 \\\hline
  rs1 & 1 & - \\\hline
  XACC & 1 & 2

\end{longtable}

\textbf{Format}\\

\begin{figure*}[!htbp]
 {\setlength{\tabcolsep}{1pt}
 \begin{center}
\begin{tabular}{@{}Y@{}I@{}I@{}F@{}W@{}W@{}I@{}F@{}Y@{}I@{}F@{}O}
\instbitrange{31}{28}&
\instbit{27}&
\instbit{26}&
\instbitrange{25}{23}&
\instbitrange{22}{21}&
\instbitrange{20}{19}&
\instbit{18}&
\instbitrange{17}{15}&
\instbitrange{14}{11}&
\instbit{10}&
\instbitrange{9}{7}&
\instbitrange{6}{0} \\
\hline
\multicolumn{1}{|c|}{1XX1} &
\multicolumn{1}{c|}{{\scriptsize sat[0]}} &
\multicolumn{1}{c|}{{\scriptsize rm[2]}} &
\multicolumn{1}{c|}{0X1} &
\multicolumn{1}{c|}{{\scriptsize rm[1:0]}} &
\multicolumn{1}{c|}{1X} &
\multicolumn{1}{c|}{{\scriptsize rs1[4]}} &
\multicolumn{1}{c|}{{\scriptsize rs1[2:0]}} &
\multicolumn{1}{c|}{01XX} &
\multicolumn{1}{c|}{{\scriptsize rd[4]}} &
\multicolumn{1}{c|}{{\scriptsize rd[2:0]}} &
\multicolumn{1}{c|}{0011011}\\
\hline
4& 1& 1& 3& 2& 2& 1& 3& 4& 1& 3& 7 \\
\end{tabular}
 \end{center}
 }
 \vspace{-0.1in}
 \caption[]{ESP.SRS.S.XACC} \end{figure*}


\newpage
\subsubsection{ESP.SRS.U.XACC}\label{instruction:ESPSRSUXACC}
\textbf{Syntax}\\
ESP.SRS.U.XACC rd,rs1,sat,rm

\textbf{Description}\\
This instruction performs an arithmetic right shift on the special register \lstinline{XACC}. While writing the shift result back to \lstinline{XACC}, the instruction saturates the shift result to a 32-bit unsigned number and writes the saturated result into the register \lstinline{rd}.\\
\textbf{Operation}
\begin{lstlisting}
  temp_shf[39:0] = XACC[39:0] >> rs1[5:0]
  XACC = temp_shf[39:0]
  rd = min(temp_shf[39:0], 2^{32}-1)
\end{lstlisting}

\textbf{Schedule}\\
\begin{longtable}[c]{|l|l|l|}
    \hline
    \rowcolor{lightgray}
    \textbf{Operands} & \textbf{Use} & \textbf{Def} \\\hline
    \endfirsthead
    \hline
    \rowcolor{lightgray}
     \textbf{Operands} & \textbf{Use} & \textbf{Def} \\\hline
    \endhead
    \hline
    \endfoot

  rd & - & 2 \\\hline
  rs1 & 1 & - \\\hline
  XACC & 1 & 2

\end{longtable}

\textbf{Format}\\

\begin{figure*}[!htbp]
 {\setlength{\tabcolsep}{1pt}
 \begin{center}
\begin{tabular}{@{}Y@{}I@{}I@{}F@{}W@{}W@{}I@{}F@{}Y@{}I@{}F@{}O}
\instbitrange{31}{28}&
\instbit{27}&
\instbit{26}&
\instbitrange{25}{23}&
\instbitrange{22}{21}&
\instbitrange{20}{19}&
\instbit{18}&
\instbitrange{17}{15}&
\instbitrange{14}{11}&
\instbit{10}&
\instbitrange{9}{7}&
\instbitrange{6}{0} \\
\hline
\multicolumn{1}{|c|}{1XX0} &
\multicolumn{1}{c|}{{\scriptsize sat[0]}} &
\multicolumn{1}{c|}{{\scriptsize rm[2]}} &
\multicolumn{1}{c|}{0X1} &
\multicolumn{1}{c|}{{\scriptsize rm[1:0]}} &
\multicolumn{1}{c|}{1X} &
\multicolumn{1}{c|}{{\scriptsize rs1[4]}} &
\multicolumn{1}{c|}{{\scriptsize rs1[2:0]}} &
\multicolumn{1}{c|}{01XX} &
\multicolumn{1}{c|}{{\scriptsize rd[4]}} &
\multicolumn{1}{c|}{{\scriptsize rd[2:0]}} &
\multicolumn{1}{c|}{0011011}\\
\hline
4& 1& 1& 3& 2& 2& 1& 3& 4& 1& 3& 7 \\
\end{tabular}
 \end{center}
 }
 \vspace{-0.1in}
 \caption[]{ESP.SRS.U.XACC} \end{figure*}


\newpage
\subsubsection{ESP.VSLD.16}\label{instruction:ESPVSLD16}
\textbf{Syntax}\\
ESP.VSLD.16 qu,qy,qw,sat,rm

\textbf{Description}\\
This instruction performs a 16-bit vector architecture shift and the shift amount comes from segments of \lstinline{qw}. If the shift amount is a negative number, it performs a right shift; otherwise, it performs a left shift.

\textbf{Operation}\\
\begin{lstlisting}[language=Python, escapechar=@]
    qu[ 15:  0] = @(qw[  4:  0] < 0) ? (qy[ 15:  0] >> (-qw[  4:  0]) ) : (qy[ 15:  0] << qw[  4:  0])@
    qu[ 31: 16] = @(qw[ 20: 16] < 0) ? (qy[ 31: 16] >> (-qw[ 20: 16]) ) : (qy[ 31: 16] << qw[ 20: 16])@
    qu[ 47: 32] = @(qw[ 36: 32] < 0) ? (qy[ 47: 32] >> (-qw[ 36: 32]) ) : (qy[ 47: 32] << qw[ 36: 32])@
    qu[ 63: 48] = @(qw[ 52: 48] < 0) ? (qy[ 63: 48] >> (-qw[ 52: 48]) ) : (qy[ 63: 48] << qw[ 52: 48])@
    qu[ 79: 64] = @(qw[ 68: 64] < 0) ? (qy[ 79: 64] >> (-qw[ 68: 64]) ) : (qy[ 79: 64] << qw[ 68: 64])@
    qu[ 95: 80] = @(qw[ 84: 80] < 0) ? (qy[ 95: 80] >> (-qw[ 84: 80]) ) : (qy[ 95: 80] << qw[ 84: 80])@
    qu[111: 96] = @(qw[100: 96] < 0) ? (qy[111: 96] >> (-qw[100: 96]) ) : (qy[111: 96] << qw[100: 96])@
    qu[127:112] = @(qw[116:112] < 0) ? (qy[127:112] >> (-qw[116:112]) ) : (qy[127:112] << qw[116:112])@

\end{lstlisting}

\textbf{Schedule}\\
\begin{longtable}[c]{|l|l|l|}
    \hline
    \rowcolor{lightgray}
    \textbf{Operands} & \textbf{Use} & \textbf{Def} \\\hline
    \endfirsthead
    \hline
    \rowcolor{lightgray}
     \textbf{Operands} & \textbf{Use} & \textbf{Def} \\\hline
    \endhead
    \hline
    \endfoot

  qu & - & 1 \\\hline
  qy & 1 & - \\\hline
  qw & 1 & -

\end{longtable}

\textbf{Format}\\

\begin{figure*}[!htbp]
 {\setlength{\tabcolsep}{1pt}
 \begin{center}
\begin{tabular}{@{}F@{}F@{}W@{}W@{}I@{}I@{}I@{}W@{}Y@{}F@{}F@{}O}
\instbitrange{31}{29}&
\instbitrange{28}{26}&
\instbitrange{25}{24}&
\instbitrange{23}{22}&
\instbit{21}&
\instbit{20}&
\instbit{19}&
\instbitrange{18}{17}&
\instbitrange{16}{13}&
\instbitrange{12}{10}&
\instbitrange{9}{7}&
\instbitrange{6}{0} \\
\hline
\multicolumn{1}{|c|}{11X} &
\multicolumn{1}{c|}{{\scriptsize qy[2:0]}} &
\multicolumn{1}{c|}{{\scriptsize qw[1:0]}} &
\multicolumn{1}{c|}{01} &
\multicolumn{1}{c|}{{\scriptsize sat[0]}} &
\multicolumn{1}{c|}{{\scriptsize rm[2]}} &
\multicolumn{1}{c|}{{\scriptsize qw[2]}} &
\multicolumn{1}{c|}{{\scriptsize rm[1:0]}} &
\multicolumn{1}{c|}{0X11} &
\multicolumn{1}{c|}{{\scriptsize qu[2:0]}} &
\multicolumn{1}{c|}{01X} &
\multicolumn{1}{c|}{0011011}\\
\hline
3& 3& 2& 2& 1& 1& 1& 2& 4& 3& 3& 7 \\
\end{tabular}
 \end{center}
 }
 \vspace{-0.1in}
 \caption[]{ESP.VSLD.16} \end{figure*}


\newpage
\subsubsection{ESP.VSLD.32}\label{instruction:ESPVSLD32}
\textbf{Syntax}\\
ESP.VSLD.32 qu,qy,qw,sat,rm

\textbf{Description}\\
This instruction performs a 32-bit vector architecture shift and the shift amount comes from segments of \lstinline{qw}. If the shift amount is a negative number, it performs a right shift; otherwise, it performs a left shift.

\textbf{Operation}\\
\begin{lstlisting}[language=Python, escapechar=@]
    qu[ 31:  0] = @(qw[  5:  0] < 0) ? (qy[ 31:  0] >> (-qw[  5:  0]) ) : (qy[ 31:  0] << qw[  5:  0])@
    qu[ 63: 32] = @(qw[ 37: 32] < 0) ? (qy[ 63: 32] >> (-qw[ 37: 32]) ) : (qy[ 63: 32] << qw[ 37: 32])@
    qu[ 95: 64] = @(qw[ 69: 64] < 0) ? (qy[ 95: 64] >> (-qw[ 69: 64]) ) : (qy[ 95: 64] << qw[ 69: 64])@
    qu[127: 96] = @(qw[101: 96] < 0) ? (qy[127: 96] >> (-qw[101: 96]) ) : (qy[127: 96] << qw[101: 96])@

\end{lstlisting}

\textbf{Schedule}\\
\begin{longtable}[c]{|l|l|l|}
    \hline
    \rowcolor{lightgray}
    \textbf{Operands} & \textbf{Use} & \textbf{Def} \\\hline
    \endfirsthead
    \hline
    \rowcolor{lightgray}
     \textbf{Operands} & \textbf{Use} & \textbf{Def} \\\hline
    \endhead
    \hline
    \endfoot

  qu & - & 1 \\\hline
  qy & 1 & - \\\hline
  qw & 1 & -

\end{longtable}

\textbf{Format}\\

\begin{figure*}[!htbp]
 {\setlength{\tabcolsep}{1pt}
 \begin{center}
\begin{tabular}{@{}F@{}F@{}W@{}I@{}I@{}W@{}I@{}I@{}V@{}F@{}F@{}O}
\instbitrange{31}{29}&
\instbitrange{28}{26}&
\instbitrange{25}{24}&
\instbit{23}&
\instbit{22}&
\instbitrange{21}{20}&
\instbit{19}&
\instbit{18}&
\instbitrange{17}{13}&
\instbitrange{12}{10}&
\instbitrange{9}{7}&
\instbitrange{6}{0} \\
\hline
\multicolumn{1}{|c|}{11X} &
\multicolumn{1}{c|}{{\scriptsize qy[2:0]}} &
\multicolumn{1}{c|}{{\scriptsize qw[1:0]}} &
\multicolumn{1}{c|}{0} &
\multicolumn{1}{c|}{{\scriptsize sat[0]}} &
\multicolumn{1}{c|}{{\scriptsize rm[2:1]}} &
\multicolumn{1}{c|}{{\scriptsize qw[2]}} &
\multicolumn{1}{c|}{{\scriptsize rm[0]}} &
\multicolumn{1}{c|}{X1X11} &
\multicolumn{1}{c|}{{\scriptsize qu[2:0]}} &
\multicolumn{1}{c|}{01X} &
\multicolumn{1}{c|}{0011011}\\
\hline
3& 3& 2& 1& 1& 2& 1& 1& 5& 3& 3& 7 \\
\end{tabular}
 \end{center}
 }
 \vspace{-0.1in}
 \caption[]{ESP.VSLD.32} \end{figure*}


\newpage
\subsubsection{ESP.VSLD.8}\label{instruction:ESPVSLD8}
\textbf{Syntax}\\
ESP.VSLD.8 qu,qy,qw,sat,rm

\textbf{Description}\\
This instruction performs an 8-bit vector architecture shift and the shift amount comes from segments of \lstinline{qw}. If the shift amount is a negative number, it performs a right shift; otherwise, it performs a left shift.

\textbf{Operation}\\
\begin{lstlisting}[language=Python, escapechar=@]
    qu[  7:  0] = @(qw[  3:  0] < 0) ? (qy[  7:  0] >> (-qw[  3:  0]) ) : (qy[  7:  0] << qw[  3:  0])@
    qu[ 15:  8] = @(qw[ 11:  8] < 0) ? (qy[ 15:  8] >> (-qw[ 11:  8]) ) : (qy[ 15:  8] << qw[ 11:  8])@
    qu[ 23: 16] = @(qw[ 19: 16] < 0) ? (qy[ 23: 16] >> (-qw[ 19: 16]) ) : (qy[ 23: 16] << qw[ 19: 16])@
    qu[ 31: 24] = @(qw[ 27: 24] < 0) ? (qy[ 31: 24] >> (-qw[ 27: 24]) ) : (qy[ 31: 24] << qw[ 27: 24])@
    qu[ 39: 32] = @(qw[ 35: 32] < 0) ? (qy[ 39: 32] >> (-qw[ 35: 32]) ) : (qy[ 39: 32] << qw[ 35: 32])@
    qu[ 47: 40] = @(qw[ 43: 40] < 0) ? (qy[ 47: 40] >> (-qw[ 43: 40]) ) : (qy[ 47: 40] << qw[ 43: 40])@
    qu[ 55: 48] = @(qw[ 51: 48] < 0) ? (qy[ 55: 48] >> (-qw[ 51: 48]) ) : (qy[ 55: 48] << qw[ 51: 48])@
    qu[ 63: 56] = @(qw[ 59: 56] < 0) ? (qy[ 63: 56] >> (-qw[ 59: 56]) ) : (qy[ 63: 56] << qw[ 59: 56])@
    qu[ 71: 64] = @(qw[ 67: 64] < 0) ? (qy[ 71: 64] >> (-qw[ 67: 64]) ) : (qy[ 71: 64] << qw[ 67: 64])@
    qu[ 79: 72] = @(qw[ 75: 72] < 0) ? (qy[ 79: 72] >> (-qw[ 75: 72]) ) : (qy[ 79: 72] << qw[ 75: 72])@
    qu[ 87: 80] = @(qw[ 83: 80] < 0) ? (qy[ 87: 80] >> (-qw[ 83: 80]) ) : (qy[ 87: 80] << qw[ 83: 80])@
    qu[ 95: 88] = @(qw[ 91: 88] < 0) ? (qy[ 95: 88] >> (-qw[ 91: 88]) ) : (qy[ 95: 88] << qw[ 91: 88])@
    qu[103: 96] = @(qw[ 99: 96] < 0) ? (qy[103: 96] >> (-qw[ 99: 96]) ) : (qy[103: 96] << qw[ 99: 96])@
    qu[111:104] = @(qw[107:104] < 0) ? (qy[111:104] >> (-qw[107:104]) ) : (qy[111:104] << qw[107:104])@
    qu[119:112] = @(qw[115:112] < 0) ? (qy[119:112] >> (-qw[115:112]) ) : (qy[119:112] << qw[115:112])@
    qu[127:120] = @(qw[123:120] < 0) ? (qy[127:120] >> (-qw[123:120]) ) : (qy[127:120] << qw[123:120])@
\end{lstlisting}

\textbf{Schedule}\\
\begin{longtable}[c]{|l|l|l|}
    \hline
    \rowcolor{lightgray}
    \textbf{Operands} & \textbf{Use} & \textbf{Def} \\\hline
    \endfirsthead
    \hline
    \rowcolor{lightgray}
     \textbf{Operands} & \textbf{Use} & \textbf{Def} \\\hline
    \endhead
    \hline
    \endfoot

  qu & - & 1 \\\hline
  qy & 1 & - \\\hline
  qw & 1 & -

\end{longtable}

\textbf{Format}\\

\begin{figure*}[!htbp]
 {\setlength{\tabcolsep}{1pt}
 \begin{center}
\begin{tabular}{@{}F@{}F@{}W@{}W@{}I@{}I@{}I@{}W@{}Y@{}F@{}F@{}O}
\instbitrange{31}{29}&
\instbitrange{28}{26}&
\instbitrange{25}{24}&
\instbitrange{23}{22}&
\instbit{21}&
\instbit{20}&
\instbit{19}&
\instbitrange{18}{17}&
\instbitrange{16}{13}&
\instbitrange{12}{10}&
\instbitrange{9}{7}&
\instbitrange{6}{0} \\
\hline
\multicolumn{1}{|c|}{11X} &
\multicolumn{1}{c|}{{\scriptsize qy[2:0]}} &
\multicolumn{1}{c|}{{\scriptsize qw[1:0]}} &
\multicolumn{1}{c|}{00} &
\multicolumn{1}{c|}{{\scriptsize sat[0]}} &
\multicolumn{1}{c|}{{\scriptsize rm[2]}} &
\multicolumn{1}{c|}{{\scriptsize qw[2]}} &
\multicolumn{1}{c|}{{\scriptsize rm[1:0]}} &
\multicolumn{1}{c|}{0X11} &
\multicolumn{1}{c|}{{\scriptsize qu[2:0]}} &
\multicolumn{1}{c|}{01X} &
\multicolumn{1}{c|}{0011011}\\
\hline
3& 3& 2& 2& 1& 1& 1& 2& 4& 3& 3& 7 \\
\end{tabular}
 \end{center}
 }
 \vspace{-0.1in}
 \caption[]{ESP.VSLD.8} \end{figure*}


\newpage
\subsubsection{ESP.VSR.S32}\label{instruction:ESPVSRS32}
\textbf{Syntax}\\
ESP.VSR.S32 qu,qy,rm

\textbf{Description}\\
This instruction performs an arithmetic right shift on the four 32-bit data segments in the register \lstinline{qy} respectively. The shift amount is the value in the 6-bit special register \lstinline{SAR}. During the shift process, the higher bits are padded with the signed bit. The lower 32 bits of the shift result are stored in the corresponding data segment in the register \lstinline{qu}.\\
\textbf{Operation}
\begin{lstlisting}
  qu[ 31:  0] = (qy[ 31:  0] >> SAR[5:0])
  qu[ 63: 32] = (qy[ 63: 32] >> SAR[5:0])
  qu[ 95: 64] = (qy[ 95: 64] >> SAR[5:0])
  qu[127: 96] = (qy[127: 96] >> SAR[5:0])
\end{lstlisting}

\textbf{Schedule}\\
\begin{longtable}[c]{|l|l|l|}
    \hline
    \rowcolor{lightgray}
    \textbf{Operands} & \textbf{Use} & \textbf{Def} \\\hline
    \endfirsthead
    \hline
    \rowcolor{lightgray}
     \textbf{Operands} & \textbf{Use} & \textbf{Def} \\\hline
    \endhead
    \hline
    \endfoot

  qu & - & 1 \\\hline
  qy & 1 & - \\\hline
  SAR & 1 & -

\end{longtable}

\textbf{Format}\\

\begin{figure*}[!htbp]
 {\setlength{\tabcolsep}{1pt}
 \begin{center}
\begin{tabular}{@{}F@{}F@{}I@{}I@{}I@{}W@{}E@{}F@{}F@{}O}
\instbitrange{31}{29}&
\instbitrange{28}{26}&
\instbit{25}&
\instbit{24}&
\instbit{23}&
\instbitrange{22}{21}&
\instbitrange{20}{13}&
\instbitrange{12}{10}&
\instbitrange{9}{7}&
\instbitrange{6}{0} \\
\hline
\multicolumn{1}{|c|}{1XX} &
\multicolumn{1}{c|}{{\scriptsize qy[2:0]}} &
\multicolumn{1}{c|}{1} &
\multicolumn{1}{c|}{{\scriptsize rm[2]}} &
\multicolumn{1}{c|}{0} &
\multicolumn{1}{c|}{{\scriptsize rm[1:0]}} &
\multicolumn{1}{c|}{1X1XXX01} &
\multicolumn{1}{c|}{{\scriptsize qu[2:0]}} &
\multicolumn{1}{c|}{0XX} &
\multicolumn{1}{c|}{0011011}\\
\hline
3& 3& 1& 1& 1& 2& 8& 3& 3& 7 \\
\end{tabular}
 \end{center}
 }
 \vspace{-0.1in}
 \caption[]{ESP.VSR.S32} \end{figure*}


\newpage
\subsubsection{ESP.VSR.U32}\label{instruction:ESPVSRU32}
\textbf{Syntax}\\
ESP.VSR.U32 qu,qy,rm

\textbf{Description}\\
This instruction performs an arithmetic right shift on the four 32-bit unsigned data segments in the register \lstinline{qy} respectively. The shift amount is the value in the 6-bit special register \lstinline{SAR}. During the shift process, the higher bits are padded with the signed bit. The lower 32 bits of the shift result are stored in the corresponding data segment in the register \lstinline{qu}.\\
\textbf{Operation}
\begin{lstlisting}
  qu[ 31:  0] = (qy[ 31:  0] >> SAR[5:0])
  qu[ 63: 32] = (qy[ 63: 32] >> SAR[5:0])
  qu[ 95: 64] = (qy[ 95: 64] >> SAR[5:0])
  qu[127: 96] = (qy[127: 96] >> SAR[5:0])
\end{lstlisting}

\textbf{Schedule}\\
\begin{longtable}[c]{|l|l|l|}
    \hline
    \rowcolor{lightgray}
    \textbf{Operands} & \textbf{Use} & \textbf{Def} \\\hline
    \endfirsthead
    \hline
    \rowcolor{lightgray}
     \textbf{Operands} & \textbf{Use} & \textbf{Def} \\\hline
    \endhead
    \hline
    \endfoot

  qu & - & 1 \\\hline
  qy & 1 & - \\\hline
  SAR & 1 & -

\end{longtable}

\textbf{Format}\\

\begin{figure*}[!htbp]
 {\setlength{\tabcolsep}{1pt}
 \begin{center}
\begin{tabular}{@{}F@{}F@{}I@{}I@{}I@{}W@{}E@{}F@{}F@{}O}
\instbitrange{31}{29}&
\instbitrange{28}{26}&
\instbit{25}&
\instbit{24}&
\instbit{23}&
\instbitrange{22}{21}&
\instbitrange{20}{13}&
\instbitrange{12}{10}&
\instbitrange{9}{7}&
\instbitrange{6}{0} \\
\hline
\multicolumn{1}{|c|}{1XX} &
\multicolumn{1}{c|}{{\scriptsize qy[2:0]}} &
\multicolumn{1}{c|}{0} &
\multicolumn{1}{c|}{{\scriptsize rm[2]}} &
\multicolumn{1}{c|}{0} &
\multicolumn{1}{c|}{{\scriptsize rm[1:0]}} &
\multicolumn{1}{c|}{1X1XXX01} &
\multicolumn{1}{c|}{{\scriptsize qu[2:0]}} &
\multicolumn{1}{c|}{0XX} &
\multicolumn{1}{c|}{0011011}\\
\hline
3& 3& 1& 1& 1& 2& 8& 3& 3& 7 \\
\end{tabular}
 \end{center}
 }
 \vspace{-0.1in}
 \caption[]{ESP.VSR.U32} \end{figure*}


\newpage
\subsubsection{ESP.VSRD.16}\label{instruction:ESPVSrd16}
\textbf{Syntax}\\
ESP.VSRD.16 qu,qy,qw,sat,rm

\textbf{Description}\\
This instruction performs a 16-bit vector architecture shift and the shift amount comes from segments of \lstinline{qw}. If the shift amount is a negative number, it performs a left shift; otherwise, it performs a right shift.

\textbf{Operation}\\
\begin{lstlisting}[language=Python, escapechar=@]
    qu[ 15:  0] = @(qw[  4:  0] < 0) ? (qy[ 15:  0] << (-qw[  4:  0]) ) : (qy[ 15:  0] >> qw[  4:  0])@
    qu[ 31: 16] = @(qw[ 20: 16] < 0) ? (qy[ 31: 16] << (-qw[ 20: 16]) ) : (qy[ 31: 16] >> qw[ 20: 16])@
    qu[ 47: 32] = @(qw[ 36: 32] < 0) ? (qy[ 47: 32] << (-qw[ 36: 32]) ) : (qy[ 47: 32] >> qw[ 36: 32])@
    qu[ 63: 48] = @(qw[ 52: 48] < 0) ? (qy[ 63: 48] << (-qw[ 52: 48]) ) : (qy[ 63: 48] >> qw[ 52: 48])@
    qu[ 79: 64] = @(qw[ 68: 64] < 0) ? (qy[ 79: 64] << (-qw[ 68: 64]) ) : (qy[ 79: 64] >> qw[ 68: 64])@
    qu[ 95: 80] = @(qw[ 84: 80] < 0) ? (qy[ 95: 80] << (-qw[ 84: 80]) ) : (qy[ 95: 80] >> qw[ 84: 80])@
    qu[111: 96] = @(qw[100: 96] < 0) ? (qy[111: 96] << (-qw[100: 96]) ) : (qy[111: 96] >> qw[100: 96])@
    qu[127:112] = @(qw[116:112] < 0) ? (qy[127:112] << (-qw[116:112]) ) : (qy[127:112] >> qw[116:112])@

\end{lstlisting}

\textbf{Schedule}\\
\begin{longtable}[c]{|l|l|l|}
    \hline
    \rowcolor{lightgray}
    \textbf{Operands} & \textbf{Use} & \textbf{Def} \\\hline
    \endfirsthead
    \hline
    \rowcolor{lightgray}
     \textbf{Operands} & \textbf{Use} & \textbf{Def} \\\hline
    \endhead
    \hline
    \endfoot

  qu & - & 1 \\\hline
  qy & 1 & - \\\hline
  qw & 1 & -

\end{longtable}


\textbf{Format}\\

\begin{figure*}[!htbp]
 {\setlength{\tabcolsep}{1pt}
 \begin{center}
\begin{tabular}{@{}F@{}F@{}W@{}W@{}I@{}I@{}I@{}W@{}Y@{}F@{}F@{}O}
\instbitrange{31}{29}&
\instbitrange{28}{26}&
\instbitrange{25}{24}&
\instbitrange{23}{22}&
\instbit{21}&
\instbit{20}&
\instbit{19}&
\instbitrange{18}{17}&
\instbitrange{16}{13}&
\instbitrange{12}{10}&
\instbitrange{9}{7}&
\instbitrange{6}{0} \\
\hline
\multicolumn{1}{|c|}{11X} &
\multicolumn{1}{c|}{{\scriptsize qy[2:0]}} &
\multicolumn{1}{c|}{{\scriptsize qw[1:0]}} &
\multicolumn{1}{c|}{11} &
\multicolumn{1}{c|}{{\scriptsize sat[0]}} &
\multicolumn{1}{c|}{{\scriptsize rm[2]}} &
\multicolumn{1}{c|}{{\scriptsize qw[2]}} &
\multicolumn{1}{c|}{{\scriptsize rm[1:0]}} &
\multicolumn{1}{c|}{0X11} &
\multicolumn{1}{c|}{{\scriptsize qu[2:0]}} &
\multicolumn{1}{c|}{01X} &
\multicolumn{1}{c|}{0011011}\\
\hline
3& 3& 2& 2& 1& 1& 1& 2& 4& 3& 3& 7 \\
\end{tabular}
 \end{center}
 }
 \vspace{-0.1in}
 \caption[]{ESP.VSRD.16} \end{figure*}


\newpage
\subsubsection{ESP.VSRD.32}\label{instruction:ESPVSrd32}
\textbf{Syntax}\\
ESP.VSRD.32 qu,qy,qw,sat,rm

\textbf{Description}\\
This instruction performs a 32-bit vector architecture shift and the shift amount comes from segments of \lstinline{qw}. If the shift amount is a negative number, it performs a left shift; otherwise, it performs a right shift.

\textbf{Operation}\\
\begin{lstlisting}[language=Python, escapechar=@]
    qu[ 31:  0] = @(qw[  5:  0] < 0) ? (qy[ 31:  0] << (-qw[  5:  0]) ) : (qy[ 31:  0] >> qw[  5:  0])@
    qu[ 63: 32] = @(qw[ 37: 32] < 0) ? (qy[ 63: 32] << (-qw[ 37: 32]) ) : (qy[ 63: 32] >> qw[ 37: 32])@
    qu[ 95: 64] = @(qw[ 69: 64] < 0) ? (qy[ 95: 64] << (-qw[ 69: 64]) ) : (qy[ 95: 64] >> qw[ 69: 64])@
    qu[127: 96] = @(qw[101: 96] < 0) ? (qy[127: 96] << (-qw[101: 96]) ) : (qy[127: 96] >> qw[101: 96])@

\end{lstlisting}

\textbf{Schedule}\\
\begin{longtable}[c]{|l|l|l|}
    \hline
    \rowcolor{lightgray}
    \textbf{Operands} & \textbf{Use} & \textbf{Def} \\\hline
    \endfirsthead
    \hline
    \rowcolor{lightgray}
     \textbf{Operands} & \textbf{Use} & \textbf{Def} \\\hline
    \endhead
    \hline
    \endfoot

  qu & - & 1 \\\hline
  qy & 1 & - \\\hline
  qw & 1 & -

\end{longtable}


\textbf{Format}\\

\begin{figure*}[!htbp]
 {\setlength{\tabcolsep}{1pt}
 \begin{center}
\begin{tabular}{@{}F@{}F@{}W@{}I@{}I@{}W@{}I@{}I@{}V@{}F@{}F@{}O}
\instbitrange{31}{29}&
\instbitrange{28}{26}&
\instbitrange{25}{24}&
\instbit{23}&
\instbit{22}&
\instbitrange{21}{20}&
\instbit{19}&
\instbit{18}&
\instbitrange{17}{13}&
\instbitrange{12}{10}&
\instbitrange{9}{7}&
\instbitrange{6}{0} \\
\hline
\multicolumn{1}{|c|}{11X} &
\multicolumn{1}{c|}{{\scriptsize qy[2:0]}} &
\multicolumn{1}{c|}{{\scriptsize qw[1:0]}} &
\multicolumn{1}{c|}{1} &
\multicolumn{1}{c|}{{\scriptsize sat[0]}} &
\multicolumn{1}{c|}{{\scriptsize rm[2:1]}} &
\multicolumn{1}{c|}{{\scriptsize qw[2]}} &
\multicolumn{1}{c|}{{\scriptsize rm[0]}} &
\multicolumn{1}{c|}{X1X11} &
\multicolumn{1}{c|}{{\scriptsize qu[2:0]}} &
\multicolumn{1}{c|}{01X} &
\multicolumn{1}{c|}{0011011}\\
\hline
3& 3& 2& 1& 1& 2& 1& 1& 5& 3& 3& 7 \\
\end{tabular}
 \end{center}
 }
 \vspace{-0.1in}
 \caption[]{ESP.VSRD.32} \end{figure*}


\newpage
\subsubsection{ESP.VSRD.8}\label{instruction:ESPVSrd8}
\textbf{Syntax}\\
ESP.VSRD.8 qu,qy,qw,sat,rm

\textbf{Description}\\
This instruction performs an 8-bit vector architecture shift, and the shift amount comes from segments of \lstinline{qw}. If the shift amount is a negative number, it performs a left shift; otherwise, it performs a right shift.

\textbf{Operation}\\
\begin{lstlisting}[language=Python, escapechar=@]
    qu[  7:  0] = @(qw[  3:  0] < 0) ? (qy[  7:  0] << (-qw[  3:  0]) ) : (qy[  7:  0] >> qw[  3:  0])@
    qu[ 15:  8] = @(qw[ 11:  8] < 0) ? (qy[ 15:  8] << (-qw[ 11:  8]) ) : (qy[ 15:  8] >> qw[ 11:  8])@
    qu[ 23: 16] = @(qw[ 19: 16] < 0) ? (qy[ 23: 16] << (-qw[ 19: 16]) ) : (qy[ 23: 16] >> qw[ 19: 16])@
    qu[ 31: 24] = @(qw[ 27: 24] < 0) ? (qy[ 31: 24] << (-qw[ 27: 24]) ) : (qy[ 31: 24] >> qw[ 27: 24])@
    qu[ 39: 32] = @(qw[ 35: 32] < 0) ? (qy[ 39: 32] << (-qw[ 35: 32]) ) : (qy[ 39: 32] >> qw[ 35: 32])@
    qu[ 47: 40] = @(qw[ 43: 40] < 0) ? (qy[ 47: 40] << (-qw[ 43: 40]) ) : (qy[ 47: 40] >> qw[ 43: 40])@
    qu[ 55: 48] = @(qw[ 51: 48] < 0) ? (qy[ 55: 48] << (-qw[ 51: 48]) ) : (qy[ 55: 48] >> qw[ 51: 48])@
    qu[ 63: 56] = @(qw[ 59: 56] < 0) ? (qy[ 63: 56] << (-qw[ 59: 56]) ) : (qy[ 63: 56] >> qw[ 59: 56])@
    qu[ 71: 64] = @(qw[ 67: 64] < 0) ? (qy[ 71: 64] << (-qw[ 67: 64]) ) : (qy[ 71: 64] >> qw[ 67: 64])@
    qu[ 79: 72] = @(qw[ 75: 72] < 0) ? (qy[ 79: 72] << (-qw[ 75: 72]) ) : (qy[ 79: 72] >> qw[ 75: 72])@
    qu[ 87: 80] = @(qw[ 83: 80] < 0) ? (qy[ 87: 80] << (-qw[ 83: 80]) ) : (qy[ 87: 80] >> qw[ 83: 80])@
    qu[ 95: 88] = @(qw[ 91: 88] < 0) ? (qy[ 95: 88] << (-qw[ 91: 88]) ) : (qy[ 95: 88] >> qw[ 91: 88])@
    qu[103: 96] = @(qw[ 99: 96] < 0) ? (qy[103: 96] << (-qw[ 99: 96]) ) : (qy[103: 96] >> qw[ 99: 96])@
    qu[111:104] = @(qw[107:104] < 0) ? (qy[111:104] << (-qw[107:104]) ) : (qy[111:104] >> qw[107:104])@
    qu[119:112] = @(qw[115:112] < 0) ? (qy[119:112] << (-qw[115:112]) ) : (qy[119:112] >> qw[115:112])@
    qu[127:120] = @(qw[123:120] < 0) ? (qy[127:120] << (-qw[123:120]) ) : (qy[127:120] >> qw[123:120])@
\end{lstlisting}

\begin{longtable}[c]{|l|l|l|}
    \hline
    \rowcolor{lightgray}
    \textbf{Operands} & \textbf{Use} & \textbf{Def} \\\hline
    \endfirsthead
    \rowcolor{lightgray}
     \textbf{Operands} & \textbf{Use} & \textbf{Def} \\\hline
    \endhead
    \endfoot

  qu & - & 1 \\\hline
  qy & 1 & - \\\hline
  qw & 1 & - \\\hline

\end{longtable}


\textbf{Format}\\

\begin{figure*}[!htbp]
 {\setlength{\tabcolsep}{1pt}
 \begin{center}
\begin{tabular}{@{}F@{}F@{}W@{}W@{}I@{}I@{}I@{}W@{}Y@{}F@{}F@{}O}
\instbitrange{31}{29}&
\instbitrange{28}{26}&
\instbitrange{25}{24}&
\instbitrange{23}{22}&
\instbit{21}&
\instbit{20}&
\instbit{19}&
\instbitrange{18}{17}&
\instbitrange{16}{13}&
\instbitrange{12}{10}&
\instbitrange{9}{7}&
\instbitrange{6}{0} \\
\hline
\multicolumn{1}{|c|}{11X} &
\multicolumn{1}{c|}{{\scriptsize qy[2:0]}} &
\multicolumn{1}{c|}{{\scriptsize qw[1:0]}} &
\multicolumn{1}{c|}{10} &
\multicolumn{1}{c|}{{\scriptsize sat[0]}} &
\multicolumn{1}{c|}{{\scriptsize rm[2]}} &
\multicolumn{1}{c|}{{\scriptsize qw[2]}} &
\multicolumn{1}{c|}{{\scriptsize rm[1:0]}} &
\multicolumn{1}{c|}{0X11} &
\multicolumn{1}{c|}{{\scriptsize qu[2:0]}} &
\multicolumn{1}{c|}{01X} &
\multicolumn{1}{c|}{0011011}\\
\hline
3& 3& 2& 2& 1& 1& 1& 2& 4& 3& 3& 7 \\
\end{tabular}
 \end{center}
 }
 \vspace{-0.1in}
 \caption[]{ESP.VSRD.8} \end{figure*}


\newpage
\subsubsection{ESP.ST.S.XACC.IP}\label{instruction:ESPSTSXACCIP}
\textbf{Syntax}\\
ESP.ST.S.XACC.IP rs1,imm8

\textbf{Description}\\
This instruction sign-extends the special register \lstinline{XACC} to 64 bits and stores the result in memory. After the access is completed, the value in the register \lstinline{rs1} is incremented by an 8-bit sign-extended constant in the instruction code segment, left-shifted by 3.\\

imm starts at -1024, ends at 1016, and steps by 8.

\textbf{Operation}
\begin{lstlisting}
  MEM(rs1[31:0]) = {24{XACC[39]}, XACC[39:0]}
  rs1[31:0] = rs1[31:0] + imm
\end{lstlisting}

\textbf{Schedule}\\
\begin{longtable}[c]{|l|l|l|}
    \hline
    \rowcolor{lightgray}
    \textbf{Operands} & \textbf{Use} & \textbf{Def} \\\hline
    \endfirsthead
    \hline
    \rowcolor{lightgray}
     \textbf{Operands} & \textbf{Use} & \textbf{Def} \\\hline
    \endhead
    \hline
    \endfoot

  XACC & 2 & - \\\hline
  rs1 & 1 & 1

\end{longtable}


\textbf{Format}\\

\begin{figure*}[!htbp]
 {\setlength{\tabcolsep}{1pt}
 \begin{center}
\begin{tabular}{@{}Y@{}W@{}W@{}F@{}W@{}I@{}F@{}W@{}F@{}F@{}O}
\instbitrange{31}{28}&
\instbitrange{27}{26}&
\instbitrange{25}{24}&
\instbitrange{23}{21}&
\instbitrange{20}{19}&
\instbit{18}&
\instbitrange{17}{15}&
\instbitrange{14}{13}&
\instbitrange{12}{10}&
\instbitrange{9}{7}&
\instbitrange{6}{0} \\
\hline
\multicolumn{1}{|c|}{0XX1} &
\multicolumn{1}{c|}{{\scriptsize imm8[7:6]}} &
\multicolumn{1}{c|}{00} &
\multicolumn{1}{c|}{{\scriptsize imm8[5:3]}} &
\multicolumn{1}{c|}{1X} &
\multicolumn{1}{c|}{{\scriptsize rs1[4]}} &
\multicolumn{1}{c|}{{\scriptsize rs1[2:0]}} &
\multicolumn{1}{c|}{00} &
\multicolumn{1}{c|}{{\scriptsize imm8[2:0]}} &
\multicolumn{1}{c|}{X11} &
\multicolumn{1}{c|}{0011111}\\
\hline
4& 2& 2& 3& 2& 1& 3& 2& 3& 3& 7 \\
\end{tabular}
 \end{center}
 }
 \vspace{-0.1in}
 \caption[]{ESP.ST.S.XACC.IP} \end{figure*}


\newpage
\subsubsection{ESP.ST.U.XACC.IP}\label{instruction:ESPSTUXACCIP}
\textbf{Syntax}\\
ESP.ST.U.XACC.IP rs1,imm8

\textbf{Description}\\
This instruction zero-extends the special register \lstinline{XACC} to 64 bits and stores the result in memory. After the access is completed, the value in the register \lstinline{rs1} is incremented by an 8-bit sign-extended constant in the instruction code segment, left-shifted by 3.\\

imm starts at -1024, ends at 1016, and steps by 8.

\textbf{Operation}
\begin{lstlisting}
  MEM(rs1[31:0]) = {24{0}, XACC[39:0]}
  rs1[31:0] = rs1[31:0] + imm
\end{lstlisting}

\textbf{Schedule}\\
\begin{longtable}[c]{|l|l|l|}
    \hline
    \rowcolor{lightgray}
    \textbf{Operands} & \textbf{Use} & \textbf{Def} \\\hline
    \endfirsthead
    \hline
    \rowcolor{lightgray}
     \textbf{Operands} & \textbf{Use} & \textbf{Def} \\\hline
    \endhead
    \hline
    \endfoot

  XACC & 2 & - \\\hline
  rs1 & 1 & 1

\end{longtable}

\textbf{Format}\\

\begin{figure*}[!htbp]
 {\setlength{\tabcolsep}{1pt}
 \begin{center}
\begin{tabular}{@{}Y@{}W@{}W@{}F@{}W@{}I@{}F@{}W@{}F@{}F@{}O}
\instbitrange{31}{28}&
\instbitrange{27}{26}&
\instbitrange{25}{24}&
\instbitrange{23}{21}&
\instbitrange{20}{19}&
\instbit{18}&
\instbitrange{17}{15}&
\instbitrange{14}{13}&
\instbitrange{12}{10}&
\instbitrange{9}{7}&
\instbitrange{6}{0} \\
\hline
\multicolumn{1}{|c|}{0XX0} &
\multicolumn{1}{c|}{{\scriptsize imm8[7:6]}} &
\multicolumn{1}{c|}{00} &
\multicolumn{1}{c|}{{\scriptsize imm8[5:3]}} &
\multicolumn{1}{c|}{1X} &
\multicolumn{1}{c|}{{\scriptsize rs1[4]}} &
\multicolumn{1}{c|}{{\scriptsize rs1[2:0]}} &
\multicolumn{1}{c|}{00} &
\multicolumn{1}{c|}{{\scriptsize imm8[2:0]}} &
\multicolumn{1}{c|}{X11} &
\multicolumn{1}{c|}{0011111}\\
\hline
4& 2& 2& 3& 2& 1& 3& 2& 3& 3& 7 \\
\end{tabular}
 \end{center}
 }
 \vspace{-0.1in}
 \caption[]{ESP.ST.U.XACC.IP} \end{figure*}




\newpage
\subsubsection{ESP.FFT.BITREV}\label{instruction:ESPFFTBITREV}
\textbf{Syntax}\\
ESP.FFT.BITREV qv,rs1

\textbf{Description}\\
This instruction swaps the bit order of data of different bit widths according to the value of the special register \lstinline{FFT_BIT_WIDTH}. Then, it compares the data before and after the swap, takes the larger value, pads the higher bits with 0 until it reaches 16 bits, and writes the result into the corresponding data segment of the register \lstinline{qa}. In the following, Switchx is a function that represents the bit inversion of the x-bit. Switch5(0b10100) = 0b00101. Here, 0b10100 represents 5-bit binary data.\\

\textbf{Operation}\\
\begin{lstlisting}
  temp0[15:0] = rs1[15:0]
  temp1[15:0] = rs1[15:0] + 1
  temp2[15:0] = rs1[15:0] + 2
  temp3[15:0] = rs1[15:0] + 3
  temp4[15:0] = rs1[15:0] + 4
  temp5[15:0] = rs1[15:0] + 5
  temp6[15:0] = rs1[15:0] + 6
  temp7[15:0] = rs1[15:0] + 7
if FFT_BIT_WIDTH==3:
  Switch3(X[2:0]) = X[0:2]
  qv[ 15:  0] = {13'h0, max(tmp0[2:0], Switch3(tmp0[2:0]))}
  qv[ 31: 16] = {13'h0, max(tmp1[2:0], Switch3(tmp1[2:0]))}
  qv[ 47: 32] = {13'h0, max(tmp2[2:0], Switch3(tmp2[2:0]))}
  qv[ 63: 48] = {13'h0, max(tmp3[2:0], Switch3(tmp3[2:0]))}
  qv[ 79: 64] = {13'h0, max(tmp4[2:0], Switch3(tmp4[2:0]))}
  qv[ 95: 80] = {13'h0, max(tmp5[2:0], Switch3(tmp5[2:0]))}
  qv[127: 96] = 0
if FFT_BIT_WIDTH==4:
  Switch4(X[3:0]) = X[0:3]
  qv[ 15:  0] = {12'h0, max(tmp0[3:0], Switch4(tmp0[3:0]))}
  qv[ 31: 16] = {12'h0, max(tmp1[3:0], Switch4(tmp1[3:0]))}
  qv[ 47: 32] = {12'h0, max(tmp2[3:0], Switch4(tmp2[3:0]))}
  qv[ 63: 48] = {12'h0, max(tmp3[3:0], Switch4(tmp3[3:0]))}
  qv[ 79: 64] = {12'h0, max(tmp4[3:0], Switch4(tmp4[3:0]))}
  qv[ 95: 80] = {12'h0, max(tmp5[3:0], Switch4(tmp5[3:0]))}
  qv[111: 96] = {12'h0, max(tmp6[3:0], Switch4(tmp6[3:0]))}
  qv[127:112] = {12'h0, max(tmp7[3:0], Switch4(tmp7[3:0]))}
  ...
if FFT_BIT_WIDTH==10:
  Switch10(X[9:0]) = X[0:9]
  qv[ 15:  0] = {6'h0, max(tmp0[9:0], Switch10(tmp0[9:0]))}
  qv[ 31: 16] = {6'h0, max(tmp1[9:0], Switch10(tmp1[9:0]))}
  qv[ 47: 32] = {6'h0, max(tmp2[9:0], Switch10(tmp2[9:0]))}
  qv[ 63: 48] = {6'h0, max(tmp3[9:0], Switch10(tmp3[9:0]))}
  qv[ 79: 64] = {6'h0, max(tmp4[9:0], Switch10(tmp4[9:0]))}
  qv[ 95: 80] = {6'h0, max(tmp5[9:0], Switch10(tmp5[9:0]))}
  qv[111: 96] = {6'h0, max(tmp6[9:0], Switch10(tmp6[9:0]))}
  qv[127:112] = {6'h0, max(tmp7[9:0], Switch10(tmp7[9:0]))}

  rs1[31:0] = rs1[31:0] + 8
\end{lstlisting}

\textbf{Schedule}\\
\begin{longtable}[c]{|l|l|l|}
    \hline
    \rowcolor{lightgray}
    \textbf{Operands} & \textbf{Use} & \textbf{Def} \\\hline
    \endfirsthead
    \hline
    \rowcolor{lightgray}
     \textbf{Operands} & \textbf{Use} & \textbf{Def} \\\hline
    \endhead
    \hline
    \endfoot


  qv & - & 1 \\\hline
  rs1 & 1 & 1
\end{longtable}

\textbf{Format}\\

\begin{figure*}[!htbp]
 {\setlength{\tabcolsep}{1pt}
 \begin{center}
\begin{tabular}{@{}T@{}F@{}I@{}I@{}F@{}E@{}O}
\instbitrange{31}{23}&
\instbitrange{22}{20}&
\instbit{19}&
\instbit{18}&
\instbitrange{17}{15}&
\instbitrange{14}{7}&
\instbitrange{6}{0} \\
\hline
\multicolumn{1}{|c|}{0XXXXXXXX} &
\multicolumn{1}{c|}{{\scriptsize qv[2:0]}} &
\multicolumn{1}{c|}{X} &
\multicolumn{1}{c|}{{\scriptsize rs1[4]}} &
\multicolumn{1}{c|}{{\scriptsize rs1[2:0]}} &
\multicolumn{1}{c|}{100XX0XX} &
\multicolumn{1}{c|}{0011011}\\
\hline
9& 3& 1& 1& 3& 8& 7 \\
\end{tabular}
 \end{center}
 }
 \vspace{-0.1in}
 \caption[]{ESP.FFT.BITREV} \end{figure*}



\newpage
\subsubsection{ESP.SRC.Q.QUP}\label{instruction:ESPSRCQQUP}
\textbf{Syntax}\\
ESP.SRC.Q.QUP qz,qw,qy

\textbf{Description}\\
This instruction performs an arithmetic right shift on the 32-byte concatenation of registers \lstinline{qw} and \lstinline{qy}, which hold the loaded data from two consecutive aligned addresses. In this way, it is possible to obtain unaligned 16-byte data, which will then be written to the register \lstinline{qz}. The right shift amount is \lstinline{SAR_BYTE} performed by 8. Simultaneously, the value of the register \lstinline{qy} is updated to \lstinline{qw}.\\

\textbf{Operation}\\
\begin{lstlisting}
  qz[127:  0] = {qy[127:  0], qw[127:  0]} >> {SAR_BYTE[3:0] << 3}
  qw = qy
\end{lstlisting}

\textbf{Schedule}\\
\begin{longtable}[c]{|l|l|l|}
    \hline
    \rowcolor{lightgray}
    \textbf{Operands} & \textbf{Use} & \textbf{Def} \\\hline
    \endfirsthead
    \hline
    \rowcolor{lightgray}
     \textbf{Operands} & \textbf{Use} & \textbf{Def} \\\hline
    \endhead
    \hline
    \endfoot

  qz & - & 1 \\\hline
  qw & 1 & - \\\hline
  qy & 1 & - \\\hline
  SAR\_BYTES & 1 & -

\end{longtable}

\textbf{Format}\\

\begin{figure*}[!htbp]
 {\setlength{\tabcolsep}{1pt}
 \begin{center}
\begin{tabular}{@{}F@{}F@{}W@{}Y@{}I@{}T@{}F@{}O}
\instbitrange{31}{29}&
\instbitrange{28}{26}&
\instbitrange{25}{24}&
\instbitrange{23}{20}&
\instbit{19}&
\instbitrange{18}{10}&
\instbitrange{9}{7}&
\instbitrange{6}{0} \\
\hline
\multicolumn{1}{|c|}{1XX} &
\multicolumn{1}{c|}{{\scriptsize qy[2:0]}} &
\multicolumn{1}{c|}{{\scriptsize qw[1:0]}} &
\multicolumn{1}{c|}{X0XX} &
\multicolumn{1}{c|}{{\scriptsize qw[2]}} &
\multicolumn{1}{c|}{1XXX100XX} &
\multicolumn{1}{c|}{{\scriptsize qz[2:0]}} &
\multicolumn{1}{c|}{0011011}\\
\hline
3& 3& 2& 4& 1& 9& 3& 7 \\
\end{tabular}
 \end{center}
 }
 \vspace{-0.1in}
 \caption[]{ESP.SRC.Q.QUP} \end{figure*}



\newpage
\subsubsection{ESP.VSL.S32}\label{instruction:ESPVSLS32}
\textbf{Syntax}\\
ESP.VSL.S32 qu,qy,sat

\textbf{Description}\\
This instruction performs a left shift on the four 32-bit signed data segments in the register \lstinline{qy} respectively. The left shift amount is the value in the 6-bit special register \lstinline{SAR}. During the shift process, the lower bits are padded with 0. The lower 32 bits of the shift result are stored in the corresponding data segment in the register \lstinline{qu}.\\

\textbf{Operation}\\
\begin{lstlisting}
  qu[ 31:  0] = (qy[ 31:  0] << SAR[5:0])
  qu[ 63: 32] = (qy[ 63: 32] << SAR[5:0])
  qu[ 95: 64] = (qy[ 95: 64] << SAR[5:0])
  qu[127: 96] = (qy[127: 96] << SAR[5:0])
\end{lstlisting}

\textbf{Schedule}\\
\begin{longtable}[c]{|l|l|l|}
    \hline
    \rowcolor{lightgray}
    \textbf{Operands} & \textbf{Use} & \textbf{Def} \\\hline
    \endfirsthead
    \hline
    \rowcolor{lightgray}
     \textbf{Operands} & \textbf{Use} & \textbf{Def} \\\hline
    \endhead
    \hline
    \endfoot

 qu & - & 1 \\\hline
  qy & 1 & - \\\hline
  SAR & 1 & -
\end{longtable}

\textbf{Format}\\

\begin{figure*}[!htbp]
 {\setlength{\tabcolsep}{1pt}
 \begin{center}
\begin{tabular}{@{}F@{}F@{}I@{}I@{}M@{}F@{}F@{}O}
\instbitrange{31}{29}&
\instbitrange{28}{26}&
\instbit{25}&
\instbit{24}&
\instbitrange{23}{13}&
\instbitrange{12}{10}&
\instbitrange{9}{7}&
\instbitrange{6}{0} \\
\hline
\multicolumn{1}{|c|}{1XX} &
\multicolumn{1}{c|}{{\scriptsize qy[2:0]}} &
\multicolumn{1}{c|}{1} &
\multicolumn{1}{c|}{{\scriptsize sat[0]}} &
\multicolumn{1}{c|}{0XX0X1XXX01} &
\multicolumn{1}{c|}{{\scriptsize qu[2:0]}} &
\multicolumn{1}{c|}{0XX} &
\multicolumn{1}{c|}{0011011}\\
\hline
3& 3& 1& 1& 11& 3& 3& 7 \\
\end{tabular}
 \end{center}
 }
 \vspace{-0.1in}
 \caption[]{ESP.VSL.S32} \end{figure*}



\newpage
\subsubsection{ESP.VSL.U32}\label{instruction:ESPVSLU32}
\textbf{Syntax}\\
ESP.VSL.U32 qu,qy,sat

\textbf{Description}\\
This instruction performs a left shift on the four 32-bit unsigned data segments in the register \lstinline{qy} respectively. The left shift amount is the value in the 6-bit special register \lstinline{SAR}. During the shift process, the lower bits are padded with 0. The lower 32 bits of the shift result are stored in the corresponding data segment in the register \lstinline{qu}.\\

\textbf{Operation}\\
\begin{lstlisting}
  qu[ 31:  0] = (qy[ 31:  0] << SAR[5:0])
  qu[ 63: 32] = (qy[ 63: 32] << SAR[5:0])
  qu[ 95: 64] = (qy[ 95: 64] << SAR[5:0])
  qu[127: 96] = (qy[127: 96] << SAR[5:0])
\end{lstlisting}

\textbf{Schedule}\\
\begin{longtable}[c]{|l|l|l|}
    \hline
    \rowcolor{lightgray}
    \textbf{Operands} & \textbf{Use} & \textbf{Def} \\\hline
    \endfirsthead
    \hline
    \rowcolor{lightgray}
     \textbf{Operands} & \textbf{Use} & \textbf{Def} \\\hline
    \endhead
    \hline
    \endfoot

  qu & - & 1 \\\hline
  qy & 1 & - \\\hline
  SAR & 1 & -
\end{longtable}

\textbf{Format}\\

\begin{figure*}[!htbp]
 {\setlength{\tabcolsep}{1pt}
 \begin{center}
\begin{tabular}{@{}F@{}F@{}I@{}I@{}M@{}F@{}F@{}O}
\instbitrange{31}{29}&
\instbitrange{28}{26}&
\instbit{25}&
\instbit{24}&
\instbitrange{23}{13}&
\instbitrange{12}{10}&
\instbitrange{9}{7}&
\instbitrange{6}{0} \\
\hline
\multicolumn{1}{|c|}{1XX} &
\multicolumn{1}{c|}{{\scriptsize qy[2:0]}} &
\multicolumn{1}{c|}{0} &
\multicolumn{1}{c|}{{\scriptsize sat[0]}} &
\multicolumn{1}{c|}{0XX0X1XXX01} &
\multicolumn{1}{c|}{{\scriptsize qu[2:0]}} &
\multicolumn{1}{c|}{0XX} &
\multicolumn{1}{c|}{0011011}\\
\hline
3& 3& 1& 1& 11& 3& 3& 7 \\
\end{tabular}
 \end{center}
 }
 \vspace{-0.1in}
 \caption[]{ESP.VSL.U32} \end{figure*}

